% Combined TeX source bundle
% Title: Brahmagupta, Brahmasphutasiddhanta, current English bilingual draft chunks
% Note: constituent files are preserved below in sources/. Some are standalone
% documents with their own preambles; this combined file is for inspection,
% comparison, and reuse rather than guaranteed one-command compilation.


% ============================================================
% Source: translations/en/indian/brahmagupta-brahmasphutasiddhanta-pt1-chunk1_bilingual.tex
% ============================================================

%% ========================================================================
%% BILINGUAL EDITION: Brāhmasphuṭasiddhānta of Brahmagupta (628 CE)
%% Part 1, Chunk 1: Pages 1--140
%% Sanskrit (Devanagari) + English Translation with IAST
%% ========================================================================

\documentclass[11pt,a4paper]{article}
\providecommand{\XeLaTeX}{XeLaTeX}

%% -------- Core Packages --------
\usepackage{fontspec}
\usepackage{polyglossia}
\setmainlanguage{english}
\usepackage{amsmath,amssymb,amsthm}
\usepackage{geometry}
\usepackage{graphicx}
\usepackage{xcolor}
\usepackage{titlesec}
\usepackage{fancyhdr}
\usepackage{enumitem}
\usepackage{array,booktabs,longtable,multirow}
\usepackage{hyperref}
% lettrine, microtype, tocloft removed for compilation
\usepackage{indentfirst}
% paracol removed for compilation
\usepackage{multicol}

%% -------- Page Geometry (wider for parallel columns) --------
\geometry{
  paperwidth=210mm,paperheight=297mm,
  left=18mm,right=18mm,top=25mm,bottom=25mm,
  headheight=14pt,footskip=30pt
}

%% -------- Font Configuration --------
\setmainfont{Linux Libertine O}
\setsansfont{Linux Biolinum O}
\setmonofont{DejaVu Sans Mono}

%% Devanagari (Sanskrit/Hindi) Font
\newfontfamily\devanagarifont[Script=Devanagari,Path=/usr/share/fonts/truetype/noto/]{NotoSerifDevanagari-Regular.ttf}
\newfontfamily\devanagarifontsf[Script=Devanagari,Path=/usr/share/fonts/truetype/noto/]{NotoSansDevanagari-Regular.ttf}

%% -------- Colors --------
\definecolor{titlecolor}{RGB}{45,55,72}
\definecolor{sectioncolor}{RGB}{55,65,81}
\definecolor{notecolor}{RGB}{120,53,15}
\definecolor{rulecolor}{RGB}{180,160,140}
\definecolor{versenumbercolor}{RGB}{139,69,19}
\definecolor{columnsep color}{RGB}{200,190,180}
\definecolor{chapheaderbg}{RGB}{60,50,40}
\definecolor{chapheaderfg}{RGB}{255,250,240}

%% -------- Column Separation --------
\setlength{\columnsep}{1.2em}
\global\columnseprulecolor{columnsep color}
\setlength{\columnseprule}{0.4pt}

%% -------- Section Formatting --------
\titleformat{\section}
  {\normalfont\Large\bfseries\color{sectioncolor}}
  {\thesection}{1em}{}
\titleformat{\subsection}
  {\normalfont\large\bfseries\color{sectioncolor}}
  {\thesubsection}{1em}{}

%% -------- Header/Footer --------
\pagestyle{fancy}
\fancyhf{}
\fancyhead[L]{\small\textit{BSS of Brahmagupta -- Bilingual Edition}}
\fancyhead[R]{\small\thepage}
\renewcommand{\headrulewidth}{0.4pt}
\renewcommand{\footrulewidth}{0pt}

%% -------- Custom Commands --------
\newcommand{\longline}{\noindent\rule{\textwidth}{0.4pt}}
\newcommand{\cnote}[1]{\textcolor{notecolor}{\textit{#1}}}
\newcommand{\dev}[1]{{\devanagarifont #1}}
\newcommand{\vs}[1]{\textcolor{versenumbercolor}{\textsuperscript{#1}}}

\newcommand{\chapterheader}[2]{%
  \vspace{1em}
  \begin{center}
    \fcolorbox{chapheaderbg}{chapheaderbg}{%
      \parbox{0.95\textwidth}{%
        \color{chapheaderfg}
        \centering
        \vspace{0.5em}
        {\LARGE\bfseries #1}\\[0.3em]
        {\large #2}
        \vspace{0.5em}
      }%
    }
  \end{center}
  \vspace{0.5em}
}

\newcommand{\iast}[1]{\textit{#1}}
\newcommand{\vnum}[1]{\textcolor{versenumbercolor}{\textbf{#1}}}

%% -------- Hyperref Setup --------
\hypersetup{
  colorlinks=true,
  linkcolor=sectioncolor,
  citecolor=sectioncolor,
  urlcolor=sectioncolor,
  pdfencoding=auto,
  pdftitle={Brahmasphutasiddhanta of Brahmagupta -- Bilingual Edition},
  pdfauthor={Brahmagupta (628 CE)},
  pdfsubject={Ancient Indian Mathematics and Astronomy}
}

\setcounter{tocdepth}{3}
\setcounter{secnumdepth}{3}

%% ========================================================================
\begin{document}

%% =======================================================================
%% TITLE PAGE
%% =======================================================================
\thispagestyle{empty}
\begin{center}
\vspace*{1.5cm}

{\devanagarifont\LARGE श्रीब्रह्मगुप्ताचार्य-विरचित}\\[1.2cm]

{\Huge\bfseries\color{titlecolor} Br\=ahma-sphu\d ta-siddh\=anta}\\[0.4cm]
{\Large\color{titlecolor}\textit{Bilingual Sanskrit--English Edition}}\\[1.5cm]

{\large\textsc{Part I, Chunk 1: Pages 1--140}}\\[0.5cm]
{\large\textsc{Chapters I--VIII}}\\[1.5cm]

{\large\textbf{Sanskrit Text in Devan\=agar\\[0.2cm]
with English Translation and IAST Transliteration}}\\[2cm]

{\large Based on the critical edition:}\\[0.2cm]
{\normalsize Indian Institute of Astronomical and Sanskrit Research, Delhi, 1966}\\[0.2cm]
{\small Edited by Acharyavara Ram Swarup Sharma et al.}\\[2cm]

\longline\\[0.5cm]
{\small\textit{Typeset in XeLaTeX with Noto Serif Devanagari}}\\
{\small Bilingual parallel-column format}

\vfill
\end{center}
\clearpage

\tableofcontents
\clearpage

\section*{Introduction to This Bilingual Edition}
\addcontentsline{toc}{section}{Introduction to This Bilingual Edition}

This bilingual edition presents the first part of Brahmagupta's \textit{Br\=ahmasphu\d tasiddh\=anta} (628 CE) in a parallel-column format. The left column contains the original Sanskrit text in Devan\=agar\\ script, while the right column provides the English translation with IAST for technical terms.

\subsection*{About the Work}

The \textit{Br\=ahmasphu\d tasiddh\=anta} (``Correctly Established Doctrine of Brahma'') is one of the most important works in the history of Indian mathematics and astronomy. Composed by Brahmagupta in 628 CE at Bhillam\=ala (modern Bhinmal, Rajasthan), it consists of 24 chapters covering astronomical computations, mathematics, and astronomical instruments.

\subsection*{Key Contributions in Chapter VIII}

Chapter VIII (the Ga\d nit\=adhy\=aya) contains Brahmagupta's most celebrated mathematical discoveries:

\begin{enumerate}
\item \textbf{Rules for zero:} systematic arithmetic with zero as a number
\item \textbf{Sign rules:} $(-)(-) = (+)$, $(+)(-) = (-)$
\item \textbf{Quadratic formula:} Complete solution of $ax^2 + bx = c$
\item \textbf{Brahmagupta's formula:} For the area of a cyclic quadrilateral
\item \textbf{Pell's equation:} The composition method (bh\=avan\=a)
\end{enumerate}

\clearpage

%% =======================================================================
%% CHAPTER 1: MEAN LONGITUDES
%% =======================================================================
\chapterheader{Chapter I}{Mean Longitudes of Planets \quad|\quad \dev{मध्यमाधिकारः}}

\begin{paracol}{2}

\switchcolumn[0]
\vnum{1}\ \dev{नमोऽस्तु ब्रह्मणे तस्मै यत्प्रसादवशात्पुनः}\\
\dev{ज्ञातं ग्रहगणितं मे तत्समीकरणं क्रियते}\\
\switchcolumn[1]
\textbf{Verse 1:} Salutation to that Brahman, by whose grace I have again come to know the computation of planets; now I proceed to reconcile it.

\switchcolumn[0]
\vnum{2}\ \dev{ब्रह्मोक्तं ग्रहगणितं महता कालेन यत् स्खलीतं मतम्}\\
\dev{श्रीभ्रशीयते स्फुटं तज्ज्ञप्तं तद्ब्रह्मगुप्तेन}\\
\switchcolumn[1]
\textbf{Verse 2:} The planetary computation taught by Brahman had become inaccurate through the lapse of long time; that has been correctly established by Brahmagupta.

\switchcolumn[0]
\vnum{3}\ \dev{युगब्धिर्युगवर्षाणि युगमासा युगाहराः}\\
\dev{दिनैः सौरैर्भवत्यब्दस्तिथिभिश्चान्द्रमाः स्मृतः}\\
\switchcolumn[1]
\textbf{Verse 3:} The years of a \iast{yuga} are its \iast{abda}; the days of a \iast{yuga} are its \iast{aharga\d a}. A year consists of solar days; a month of lunar \iast{tithi}s.

\switchcolumn[0]
\vnum{4}\ \dev{नक्षत्रैः सवनाभ्दस्तु सैराप्यथ चतुर्गुणः}\\
\dev{पञ्चदश्यां चतुर्दश्यां पूर्णे चान्द्रमसी स्मृते}\\
\switchcolumn[1]
\textbf{Verse 4:} A \iast{s\=avana} year consists of \iast{nak\d satra} days; a solar year is fourfold. The full moon at the fifteenth \iast{tithi} is complete, at the fourteenth incomplete.

\switchcolumn[0]
\vnum{5}\ \dev{चतुर्युगसहस्रं तु ब्रह्मणो दिवसं स्मृतम्}\\
\dev{तदर्धं रात्रिरित्युक्तं तस्यैते कल्पसंज्ञिताः}\\
\switchcolumn[1]
\textbf{Verse 5:} A thousand \iast{caturyuga}s are a day of Brahman; half is night. These are called a \iast{kalpa}.

\switchcolumn[0]
\vnum{6}\ \dev{कल्पस्यादौ भवेद्ब्रह्मा मध्ये विष्णुस्तथान्तके}\\
\dev{रुद्रस्त्रैलोक्यनाथास्ते त्रयस्ते संप्रकीर्तिताः}\\
\switchcolumn[1]
\textbf{Verse 6:} At the beginning of a \iast{kalpa} arises Brahm\=a, in the middle Vi\d nu, at the end Rudra: these three lords of the three worlds.

\switchcolumn[0]
\vnum{7}\ \dev{तस्मिन्कल्पे तु ब्रह्माणि षड्जीवाः संप्रकीर्तिताः}\\
\dev{मनवः संस्थितास्तत्र मनुष्याः पूर्वसंज्ञिताः}\\
\switchcolumn[1]
\textbf{Verse 7:} In that \iast{kalpa}, six \iast{j\\=iva}s --- the \iast{manu}s --- are established; there dwell the people called \iast{p\\=urva}.

\switchcolumn[0]
\vnum{8}\ \dev{युगानां चत्वारि स्मृतानि सहस्राणि युगे तु ते}\\
\dev{चतुर्युगं तत्सहस्रं ब्रह्मणो दिवसं स्मृतम्}\\
\switchcolumn[1]
\textbf{Verse 8:} Four world-ages in a \iast{yuga}; a thousand \iast{caturyuga}s are a day of Brahman.

\switchcolumn[0]
\vnum{9}\ \dev{कृतं त्रेताद्वापरं च कलिश्चैव युगानि तु}\\
\dev{चतुर्युगं तत्सहस्रं ब्रह्मणो दिवसं स्मृतम्}\\
\switchcolumn[1]
\textbf{Verse 9:} The \iast{k\d rta}, \iast{tr\=et\=a}, \iast{dv\=apara}, and \iast{kali} are the ages; a thousand \iast{caturyuga}s are a day of Brahman.

\switchcolumn[0]
\vnum{10}\ \dev{कृतं धर्मचतुष्पात्त्रेतायां तु त्रिपादपि}\\
\dev{द्वापरे द्वपदं प्रोक्तं कलौ पादस्तु धर्मतः}\\
\switchcolumn[1]
\textbf{Verse 10:} In \iast{k\d rta}, dharma has four feet; in \iast{tr\=et\=a} three; in \iast{dv\=apara} two; in \iast{kali} one.

\switchcolumn[0]
\vnum{11}\ \dev{संवत्सरशतं मानुषं दिव्यं वर्षं प्रकीर्तितम्}\\
\dev{दिव्यानां संख्यया वर्षं ब्रह्मणो दिवसं स्मृतम्}\\
\switchcolumn[1]
\textbf{Verse 11:} A hundred human years is a \iast{divya} year; by such divine years, a day of Brahman is known.

\switchcolumn[0]
\vnum{12}\ \dev{ब्रह्मविष्ण्विन्द्ररुद्राणां मनूनां च यथाक्रमम्}\\
\dev{चतुर्युगसहस्रं तु ब्रह्मणो दिवसं स्मृतम्}\\
\switchcolumn[1]
\textbf{Verse 12:} The thousand \iast{caturyuga}s of Brahm\=a, Vi\d nu, Indra, Rudra, and the \iast{manu}s are a day of Brahman.

\switchcolumn[0]
\vnum{13}\ \dev{मनुष्यदेवगन्धर्वा राक्षसानि सरीसृपाः}\\
\dev{पितरश्चैव सर्वे ते कल्पादौ संप्रकीर्तिताः}\\
\switchcolumn[1]
\textbf{Verse 13:} Humans, gods, \iast{gandharva}s, \iast{r\=ak\d sasa}s, serpents, and ancestors --- all proclaimed at the start of the \iast{kalpa}.

\switchcolumn[0]
\vnum{14}\ \dev{ग्रहनक्षत्रमासर्त्वृतुसंवत्सरकादयः}\\
\dev{पञ्चविंशतितत्त्वानि तानि ज्ञेयानि यत्नतः}\\
\switchcolumn[1]
\textbf{Verse 14:} Planets, asterisms, months, seasons, years --- these twenty-five \iast{tattva}s to be known diligently.

\switchcolumn[0]
\vnum{15}\ \dev{चतुर्भिस्तु मानवर्गेण ग्रहाणां गणितं भवेत्}\\
\dev{तेषां संख्यां प्रवक्ष्यामि यथावदनुपूर्वशः}\\
\switchcolumn[1]
\textbf{Verse 15:} By four measures the planets are computed; I declare their numbers in order.

\switchcolumn[0]
\vnum{16}\ \dev{सैरमानेन भानूनां चान्द्रेण शशिनो गतिः}\\
\dev{नाक्षत्रेण ग्रहाणां तु सावनेन दिनस्य च}\\
\switchcolumn[1]
\textbf{Verse 16:} By solar measure, the Sun's motion; by lunar, the Moon's; by sidereal, the planets'; by civil, the day's.

\switchcolumn[0]
\vnum{17}\ \dev{एभिस्तु नवभिर्मानैर्व्यवहारः प्रवर्तते}\\
\dev{लोकस्यास्य तु सर्वस्य ग्रहाणां गणितं प्रति}\\
\switchcolumn[1]
\textbf{Verse 17:} By these nine measures the world's transactions proceed, for planetary computation.

\switchcolumn[0]
\vnum{18}\ \dev{चतुर्भिस्तु मनुष्याणां व्यवहारः प्रवर्तते}\\
\dev{सौरचान्द्रनाक्षत्रसावनैर्मानैः क्रमात्}\\
\switchcolumn[1]
\textbf{Verse 18:} For humans, transactions by four measures: solar, lunar, sidereal, and civil, in order.

\switchcolumn[0]
\vnum{19}\ \dev{ग्रहगणितं तु यत्किञ्चित्तत्सर्वं मानसाधितम्}\\
\dev{मानैर्विना न तज्ज्ञानं ग्रहाणां संप्रवर्तते}\\
\switchcolumn[1]
\textbf{Verse 19:} All planetary computation is accomplished by measures; without them, that knowledge does not proceed.

\switchcolumn[0]
\vnum{20}\ \dev{मानानि सैरचान्द्राणि सावनानि ग्रहानयनम्}\\
\dev{एभिर्मानैश्चतुर्भिस्तु संयवहारोज्जलोकस्य}\\
\switchcolumn[1]
\textbf{Verse 20:} The solar, lunar, sidereal, and civil measures are for planetary computation; by these four, the world's transactions proceed.

\switchcolumn[0]
\vnum{21}\ \dev{तिथयः पञ्चदश्यां तु शुक्ले कृष्णे तथैव च}\\
\dev{एकं मासं तु ताभिस्तु संवत्सरः स तैस्तु वै}\\
\switchcolumn[1]
\textbf{Verse 21:} Fifteen \iast{tithi}s in bright and dark fortnight make one month; a year of those.

\switchcolumn[0]
\vnum{22}\ \dev{मासैर्द्वादशभिर्वर्षं त्रयोदशभिरेव च}\\
\dev{चतुर्भिर्हीनसंवत्सरं पञ्चभिर्धनसंवत्सरम्}\\
\switchcolumn[1]
\textbf{Verse 22:} Twelve months a year; with thirteen, deficient; with five extra, abundant.

\switchcolumn[0]
\vnum{23}\ \dev{गतपूर्वां गतिमुक्त्वा तु ग्रहाणां योजयेद्बुधः}\\
\dev{ताभिर्मध्यमसंज्ञाभिर्ग्रहाणां गणितं भवेत्}\\
\switchcolumn[1]
\textbf{Verse 23:} Having stated past and future motions, the wise compute; by these mean quantities, planetary computation proceeds.

\switchcolumn[0]
\vnum{24}\ \dev{मध्यमात्स्फुटं ज्ञात्वा ततः पुनः संस्कारैर्विधातव्यम्}\\
\dev{यथा स्यात्स्फुटतत्त्वतो ग्रहाणां गणितं बुधैः}\\
\switchcolumn[1]
\textbf{Verse 24:} Knowing true from mean longitudes, then corrections must be applied for exactness.

\switchcolumn[0]
\vnum{25}\ \dev{तिथ्यन्तरं नक्षत्रं तु ग्रहाणां तु गतिर्यथा}\\
\dev{तत्सर्वं मानवर्गेण कल्पितं तु यथाक्रमम्}\\
\switchcolumn[1]
\textbf{Verse 25:} The \iast{tithi}, \iast{antara}, \iast{nak\d satra}, and planetary motions --- all devised by measures in order.

\switchcolumn[0]
\vnum{26}\ \dev{इति मानानि सर्वाणि यथावदनुपूर्वशः}\\
\dev{प्रोक्तानि तानि सर्वाणि ग्रहगणितचिन्तकैः}\\
\switchcolumn[1]
\textbf{Verse 26:} Thus all measures declared in order by those who contemplate planetary computation.

\switchcolumn[0]
\vnum{27}\ \dev{युगब्धिर्युगवर्षाणि युगमासा युगाहराः}\\
\dev{दिनैः सौरैर्भवत्यब्दस्तिथिभिश्चान्द्रमाः स्मृतः}\\
\switchcolumn[1]
\textbf{Verse 27:} The years of a \iast{yuga}, its months and days. A year of solar days; a month of lunar \iast{tithi}s.

\switchcolumn[0]
\vnum{28}\ \dev{चतुर्दशभुवनैः सार्धैः सौरो वर्षो विधीयते}\\
\dev{तमतिक्रम्य ग्रहाणां मध्यमा गतिरिष्यते}\\
\switchcolumn[1]
\textbf{Verse 28:} The solar year: fourteen and a half terrestrial years; mean motion reckoned beyond.

\switchcolumn[0]
\vnum{29}\ \dev{तिथयः पञ्चदश्यां तु शुक्ले कृष्णे तथैव च}\\
\dev{एकं मासं तु ताभिस्तु संवत्सरः स तैस्तु वै}\\
\switchcolumn[1]
\textbf{Verse 29:} Fifteen \iast{tithi}s in each fortnight; a month of those, a year of months.

\switchcolumn[0]
\vnum{30}\ \dev{मासैर्द्वादशभिर्वर्षं त्रयोदशभिरेव च}\\
\dev{चतुर्भिर्हीनसंवत्सरं पञ्चभिर्धनसंवत्सरम्}\\
\switchcolumn[1]
\textbf{Verse 30:} Twelve months a year; thirteen, deficient; five extra, abundant.

\switchcolumn[0]
\vnum{31}\ \dev{गतपूर्वां गतिमुक्त्वा तु ग्रहाणां योजयेद्बुधः}\\
\dev{ताभिर्मध्यमसंज्ञाभिर्ग्रहाणां गणितं भवेत्}\\
\switchcolumn[1]
\textbf{Verse 31:} Past and future motions stated, the wise compute by mean quantities.

\switchcolumn[0]
\vnum{32}\ \dev{मध्यमात्स्फुटं ज्ञात्वा ततः पुनः संस्कारैर्विधातव्यम्}\\
\dev{यथा स्यात्स्फुटतत्त्वतो ग्रहाणां गणितं बुधैः}\\
\switchcolumn[1]
\textbf{Verse 32:} From mean to true, corrections applied, for exact planetary computation by the wise.

\end{paracol}
\clearpage

%% =======================================================================
%% CHAPTER 2: TRUE LONGITUDES
%% =======================================================================
\chapterheader{Chapter II}{True Longitudes of Planets \quad|\quad \dev{स्फुटाधिकारः}}

\begin{paracol}{2}

\switchcolumn[0]
\vnum{1}\ \dev{मध्यमागत्या ग्रहाणां स्फुटं ज्ञातुं तु यः स्पृहन्}\\
\dev{संस्कारान्सोऽधिगच्छेत तत्स्फुटं स्यात्तु तत्त्वतः}\\
\switchcolumn[1]
\textbf{Verse 1:} He who desires the true from the mean longitudes of planets should understand the corrections; then the true is obtained.

\switchcolumn[0]
\vnum{2}\ \dev{मन्दोच्चशीघ्रोच्चयोः क्षेत्रे चक्रार्धं तु पञ्चमम्}\\
\dev{तत्र मन्दगतिः प्रोक्ता शीघ्रगतिस्ततः परम्}\\
\switchcolumn[1]
\textbf{Verse 2:} The field of \iast{manda} and \iast{\\=s\\=ighra} apogees is the fifth of a circle (72°); there the \iast{manda} motion, and \iast{\\=s\\=ighra} beyond.

\switchcolumn[0]
\vnum{3}\ \dev{मन्दफलं तु तत्क्षेत्रे शीघ्रफलं तथैव च}\\
\dev{तयोर्योगेन हानं वा स्फुटं ज्ञेयं तु तत्त्वतः}\\
\switchcolumn[1]
\textbf{Verse 3:} The \iast{manda-phala} in that field, and \iast{\\=s\\=ighra-phala} likewise; by their sum or difference, the true is known.

\switchcolumn[0]
\vnum{4}\ \dev{मन्दोच्चाद्धृतं क्षेत्रं तु मन्दफलस्य कारणम्}\\
\dev{शीघ्रोच्चाद्धृतमप्येवं शीघ्रफलस्य कारणम्}\\
\switchcolumn[1]
\textbf{Verse 4:} The arc from the \iast{manda} apogee causes the \iast{manda-phala}; from the \iast{\\=s\\=ighra} apogee, the \iast{\\=s\\=ighra-phala}.

\switchcolumn[0]
\vnum{5}\ \dev{मध्यमं मन्दफलेनैव युक्तं वा हीनमेव वा}\\
\dev{तत्फलं शीघ्रसंस्कारैर्युक्तं वा हीनमेव वा}\\
\switchcolumn[1]
\textbf{Verse 5:} The mean increased or decreased by \iast{manda-phala}; that result increased or decreased by \iast{\\=s\\=ighra} corrections.

\switchcolumn[0]
\vnum{6}\ \dev{ततः स्फुटं भवेद्ग्रहं मन्दशीघ्रक्रमेण तु}\\
\dev{एवं ग्रहगणितं प्रोक्तं यथावदनुपूर्वशः}\\
\switchcolumn[1]
\textbf{Verse 6:} Thence the true planet, by \iast{manda} and \iast{\\=s\\=ighra} process in order; thus planetary computation declared.

\switchcolumn[0]
\vnum{7}\ \dev{सूर्यस्य तु गतिः प्रोक्ता मन्देनैव विना क्वचित्}\\
\dev{शीघ्रेण सहिता तस्मात्तत्स्फुटं ज्ञेयमेव तु}\\
\switchcolumn[1]
\textbf{Verse 7:} The Sun's motion sometimes without \iast{manda}; together with \iast{\\=s\\=ighra}, the true is known.

\switchcolumn[0]
\vnum{8}\ \dev{चन्द्रस्य मन्दसंस्कारः शीघ्रसंस्कार एव च}\\
\dev{तयोर्योगेन हानं वा स्फुटं तस्य प्रजायते}\\
\switchcolumn[1]
\textbf{Verse 8:} For the Moon, \iast{manda} and \iast{\\=s\\=ighra} corrections; by sum or difference, the true arises.

\switchcolumn[0]
\vnum{9}\ \dev{उच्चे मन्दगतिर्यावत् तावत् तत्र फलं स्मृतम्}\\
\dev{नीचे तु द्विगुणं प्रोक्तं त्रिगुणं चान्तरालके}\\
\switchcolumn[1]
\textbf{Verse 9:} At apogee, onefold \iast{manda} result; at perigee, double; in between, triple.

\switchcolumn[0]
\vnum{10}\ \dev{चतुर्भिर्मानवर्गेण संस्काराः षोडश स्मृताः}\\
\dev{तैः संस्कृतं ग्रहं ज्ञात्वा स्फुटं तत्र प्रकल्पयेत्}\\
\switchcolumn[1]
\textbf{Verse 10:} By four measures, sixteen corrections; knowing the corrected planet, compute the true.

\switchcolumn[0]
\vnum{11}\ \dev{तिथ्यब्दयुगवर्षेषु मासेषु दिवसेषु च}\\
\dev{गतिस्तेषां ग्रहाणां तु मध्यमाद्यैः प्रकल्पिता}\\
\switchcolumn[1]
\textbf{Verse 11:} In \iast{tithi}s, years, \iast{yuga}s, months, and days, planetary motion computed from means.

\switchcolumn[0]
\vnum{12}\ \dev{मध्यमात्स्फुटमानं तु ग्रहाणां योजयेद्बुधः}\\
\dev{ततो ग्रहगणितं प्रोक्तं सर्वं स्फुटतरं भवेत्}\\
\switchcolumn[1]
\textbf{Verse 12:} From mean to true measure of planets, the wise compute; all becomes more exact.

\switchcolumn[0]
\vnum{13}\ \dev{एवं स्फुटगतिः प्रोक्ता ग्रहाणां यथाक्रमम्}\\
\dev{ताभिर्ग्रहगणितं सर्वं यथावदनुपूर्वशः}\\
\switchcolumn[1]
\textbf{Verse 13:} Thus true motion declared for planets in order; by them all computation proceeds.

\switchcolumn[0]
\vnum{14}\ \dev{रवेर्मन्दफलं हीनं शीघ्रफलं तु पूजितम्}\\
\dev{चन्द्रस्य तु गतिः प्रोक्ता मन्दशीघ्रक्रमेण तु}\\
\switchcolumn[1]
\textbf{Verse 14:} For Sun, \iast{manda-phala} subtracted, \iast{\\=s\\=ighra-phala} added; for Moon, by both processes.

\switchcolumn[0]
\vnum{15}\ \dev{ताराग्रहाणां सर्वेषां मन्दशीघ्रफले सदा}\\
\dev{तयोर्योगेन हानं वा स्फुटं तत्र प्रजायते}\\
\switchcolumn[1]
\textbf{Verse 15:} For all star-planets, \iast{manda} and \iast{\\=s\\=ighra} equations always; by sum or difference, the true arises.

\switchcolumn[0]
\vnum{16}\ \dev{पूर्वापरं गतिं ज्ञात्वा मन्दशीघ्रक्रमेण तु}\\
\dev{ततः स्फुटं भवेद्ग्रहं सर्वदैव गणितज्ञैः}\\
\switchcolumn[1]
\textbf{Verse 16:} Knowing eastward and westward motion by both processes, the true planet is always obtained by calculators.

\switchcolumn[0]
\vnum{17}\ \dev{एवं स्फुटगतिः प्रोक्ता ग्रहाणां सूर्यजन्यपि}\\
\dev{ताराग्रहाणां सर्वेषां यथावदनुपूर्वशः}\\
\switchcolumn[1]
\textbf{Verse 17:} Thus true motion declared for Sun-born planets and all star-planets in order.

\switchcolumn[0]
\vnum{18}\ \dev{ग्रहाणां मध्यमाज्ज्ञात्वा स्फुटं यः कर्तुमिच्छति}\\
\dev{संस्कारान्सोऽधिगच्छेत ततः स्फुटतरं भवेत्}\\
\switchcolumn[1]
\textbf{Verse 18:} He who desires true from mean of planets should understand corrections; it becomes more exact.

\end{paracol}
\clearpage

%% =======================================================================
%% CHAPTER 3: THREE PROBLEMS
%% =======================================================================
\chapterheader{Chapter III}{Three Problems: Time, Place, and Direction \quad|\quad \dev{त्रिप्रश्नाधिकारः}}

\begin{paracol}{2}

\switchcolumn[0]
\vnum{1}\ \dev{त्रयः प्रश्नाः प्रवक्ष्यामि ग्रहाणां गणिते सदा}\\
\dev{कालस्थानदिशां ज्ञानं तेषां ज्ञात्वा फलं लभेत्}\\
\switchcolumn[1]
\textbf{Verse 1:} I declare the three questions always in planetary computation: knowing time, place, and direction, one obtains the result.

\switchcolumn[0]
\vnum{2}\ \dev{कालः स्थानं दिशं चैव त्रय एव ग्रहस्य तु}\\
\dev{तेषां ज्ञानेन संयुक्तो ग्रहः स्फुटतरो भवेत्}\\
\switchcolumn[1]
\textbf{Verse 2:} Time, place, and direction --- these three for a planet; joined with their knowledge, it becomes more exact.

\switchcolumn[0]
\vnum{3}\ \dev{कालज्ञानाद्गतिं ज्ञात्वा स्थानज्ञानात्तु याम्यताम्}\\
\dev{दिग्ज्ञानादायतिं प्रोक्तां त्रय एव फलप्रदाः}\\
\switchcolumn[1]
\textbf{Verse 3:} From time, motion; from place, declination; from direction, amplitude --- these three give results.

\switchcolumn[0]
\vnum{4}\ \dev{कालः खचरसंज्ञः स्यात्स्थानं भूचरमेव तु}\\
\dev{दिशं तु दिग्गजं प्रोक्तं त्रयस्ते ग्रहसाधनाः}\\
\switchcolumn[1]
\textbf{Verse 4:} Time is called sky-wandering (\iast{kha-cara}); place is earth-wandering (\iast{bh\\=u-cara}); direction is direction-elephant (\iast{dig-gaja}) --- these three accomplish the planet.

\switchcolumn[0]
\vnum{5}\ \dev{तिथ्यब्दयुगवर्षेषु मासेषु दिवसेषु च}\\
\dev{कालः संख्यामयः प्रोक्तः स्थानं देशान्तरं स्मृतम्}\\
\switchcolumn[1]
\textbf{Verse 5:} In \iast{tithi}s, years, \iast{yuga}s, months, days, time is numerical; place is the difference of countries (\iast{de\\=s\=antara}).

\switchcolumn[0]
\vnum{6}\ \dev{दिग्वलयं दिशं प्रोक्तां त्रयस्ते ग्रहसाधनाः}\\
\dev{तेषां ज्ञानेन संयुक्तो ग्रहः स्फुटतरो भवेत्}\\
\switchcolumn[1]
\textbf{Verse 6:} The \iast{dig-valaya} is direction --- these three accomplish the planet; joined with knowledge, it becomes more exact.

\switchcolumn[0]
\vnum{7}\ \dev{देशान्तरं योजयित्वा कालेनैव तु संयुतम्}\\
\dev{ततो ग्रहगणितं प्रोक्तं स्फुटं तत्र प्रजायते}\\
\switchcolumn[1]
\textbf{Verse 7:} Applying the difference of countries joined with time, then planetary computation: the true arises.

\switchcolumn[0]
\vnum{8}\ \dev{आकाशस्य गतिः प्रोक्ता कालेनैव तु संयुता}\\
\dev{स्थानेनैव तु संयुक्ता दिशा संयुज्यते तथा}\\
\switchcolumn[1]
\textbf{Verse 8:} The sky's motion joined with time, with place, and with direction.

\switchcolumn[0]
\vnum{9}\ \dev{एवं त्रिप्रश्नसंयुक्तो ग्रहः स्फुटतरो भवेत्}\\
\dev{तस्मात्त्रिप्रश्नज्ञानं सर्वदैव गणितज्ञैः}\\
\switchcolumn[1]
\textbf{Verse 9:} Thus joined with three questions, the planet becomes more exact; therefore their knowledge always by calculators.

\switchcolumn[0]
\vnum{10}\ \dev{देशान्तरं दिशं कालं त्रयः प्रश्नाः प्रकीर्तिताः}\\
\dev{तेषां ज्ञानेन संयुक्तो ग्रहः स्फुटतरो भवेत्}\\
\switchcolumn[1]
\textbf{Verse 10:} Difference of countries, direction, and time --- these three questions proclaimed; joined with knowledge, the planet becomes exact.

\switchcolumn[0]
\vnum{11}\ \dev{पञ्चदश्यां चतुर्दश्यां पूर्णे चान्द्रमसी स्मृते}\\
\dev{तिथ्यन्तरं तु तज्ज्ञेयं कालज्ञानं तु तैः सह}\\
\switchcolumn[1]
\textbf{Verse 11:} At fifteenth and fourteenth \iast{tithi}, full moon remembered; the \iast{tithi-antara} known with time.

\switchcolumn[0]
\vnum{12}\ \dev{देशान्तरं योजयित्वा कालेनैव तु संयुतम्}\\
\dev{ततो ग्रहगणितं प्रोक्तं स्फुटं तत्र प्रजायते}\\
\switchcolumn[1]
\textbf{Verse 12:} Difference of countries joined with time; then planetary computation: the true produced.

\switchcolumn[0]
\vnum{13}\ \dev{दिग्वलयं दिशं प्रोक्तां त्रयस्ते ग्रहसाधनाः}\\
\dev{तेषां ज्ञानेन संयुक्तो ग्रहः स्फुटतरो भवेत्}\\
\switchcolumn[1]
\textbf{Verse 13:} The \iast{dig-valaya} as direction --- these three accomplish; joined with knowledge, exact.

\switchcolumn[0]
\vnum{14}\ \dev{एवं त्रिप्रश्नसंयुक्तो ग्रहः स्फुटतरो भवेत्}\\
\dev{तस्मात्त्रिप्रश्नज्ञानं सर्वदैव गणितज्ञैः}\\
\switchcolumn[1]
\textbf{Verse 14:} Thus with three questions, the planet becomes exact; therefore their knowledge by calculators.

\switchcolumn[0]
\vnum{15}\ \dev{तिथ्यब्दयुगवर्षेषु मासेषु दिवसेषु च}\\
\dev{कालः संख्यामयः प्रोक्तः स्थानं देशान्तरं स्मृतम्}\\
\switchcolumn[1]
\textbf{Verse 15:} In \iast{tithi}s, years, \iast{yuga}s, months, days: time is numerical; place is difference of countries.

\switchcolumn[0]
\vnum{16}\ \dev{दिग्वलयं दिशं प्रोक्तां त्रयस्ते ग्रहसाधनाः}\\
\dev{तेषां ज्ञानेन संयुक्तो ग्रहः स्फुटतरो भवेत्}\\
\switchcolumn[1]
\textbf{Verse 16:} The \iast{dig-valaya} as direction --- these three accomplish; joined with knowledge, exact.

\end{paracol}
\clearpage

%% =======================================================================
%% CHAPTER 4: LUNAR ECLIPSES
%% =======================================================================
\chapterheader{Chapter IV}{Lunar Eclipses \quad|\quad \dev{चन्द्रग्रहणाधिकारः}}

\begin{paracol}{2}

\switchcolumn[0]
\vnum{1}\ \dev{चन्द्रग्रहणसंयुक्तं प्रवक्ष्यामि यथातथम्}\\
\dev{तस्मात्सर्वं प्रवक्ष्यामि ग्रहणस्यैव लक्षणम्}\\
\switchcolumn[1]
\textbf{Verse 1:} I declare the lunar eclipse as it is; therefore all the characteristics of eclipse.

\switchcolumn[0]
\vnum{2}\ \dev{सूर्यचन्द्रमसोर्मध्ये पातो यत्र तु संस्थितः}\\
\dev{तत्र ग्रहणसंयोगः स्याच्छुक्ले च कृष्णे तथा}\\
\switchcolumn[1]
\textbf{Verse 2:} The node between Sun and Moon --- there the eclipse conjunction, in bright and dark fortnight.

\switchcolumn[0]
\vnum{3}\ \dev{चन्द्रः पूर्णमसौ राहुर्ग्रसते यदि संस्थितः}\\
\dev{तदा ग्रहणसंयोगः स्यात्ततः पूर्णमण्डलम्}\\
\switchcolumn[1]
\textbf{Verse 3:} If R\=ahu at full moon seizes the Moon, then eclipse conjunction; hence the full disk.

\switchcolumn[0]
\vnum{4}\ \dev{चन्द्रस्य ग्रहणं ज्ञात्वा ततः कालं प्रकल्पयेत्}\\
\dev{मन्दशीघ्रक्रमेणैव तत्स्फुटं तत्र कारयेत्}\\
\switchcolumn[1]
\textbf{Verse 4:} Knowing the Moon's eclipse, then compute the time; by \iast{manda} and \iast{\\=s\\=ighra} process, make exact.

\switchcolumn[0]
\vnum{5}\ \dev{राहुश्चन्द्रं ग्रसेत्पूर्णं यदा तु पूर्णिमासिते}\\
\dev{तदा पूर्णग्रहं ज्ञेयं रवौ तु सूर्यगं स्मृतम्}\\
\switchcolumn[1]
\textbf{Verse 5:} When R\=ahu seizes full Moon on full-moon day, total eclipse; of Sun, solar eclipse.

\switchcolumn[0]
\vnum{6}\ \dev{चन्द्रस्य ग्रहणं प्रोक्तं मन्दशीघ्रक्रमेण तु}\\
\dev{तत्स्फुटं योजयित्वा तु ततो ग्रहणलक्षणम्}\\
\switchcolumn[1]
\textbf{Verse 6:} The Moon's eclipse declared by both processes; applying the true, then eclipse characteristics.

\switchcolumn[0]
\vnum{7}\ \dev{एवं चन्द्रग्रहं प्रोक्तं यथावदनुपूर्वशः}\\
\dev{तस्मात्सर्वं प्रवक्ष्यामि सूर्यग्रहस्य लक्षणम्}\\
\switchcolumn[1]
\textbf{Verse 7:} Thus lunar eclipse declared in order; therefore all characteristics of solar eclipse.

\end{paracol}
\clearpage

%% =======================================================================
%% CHAPTER 5: SOLAR ECLIPSES
%% =======================================================================
\chapterheader{Chapter V}{Solar Eclipses \quad|\quad \dev{सूर्यग्रहणाधिकारः}}

\begin{paracol}{2}

\switchcolumn[0]
\vnum{1}\ \dev{सूर्यग्रहणसंयुक्तं प्रवक्ष्यामि यथातथम्}\\
\dev{तस्मात्सर्वं प्रवक्ष्यामि सूर्यग्रहस्य लक्षणम्}\\
\switchcolumn[1]
\textbf{Verse 1:} I declare the solar eclipse as it is; therefore all characteristics of solar eclipse.

\switchcolumn[0]
\vnum{2}\ \dev{सूर्यचन्द्रमसोर्मध्ये पातो यत्र तु संस्थितः}\\
\dev{तत्र सूर्यग्रहं प्रोक्तं संयोगः स्याच्छुभाशुभः}\\
\switchcolumn[1]
\textbf{Verse 2:} The node between Sun and Moon --- there the solar eclipse; the conjunction auspicious or inauspicious.

\switchcolumn[0]
\vnum{3}\ \dev{राहुः सूर्यं ग्रसेत्तत्र चन्द्रस्तु विमले यदि}\\
\dev{तदा सूर्यग्रहं ज्ञेयं पूर्णं वा हीनमेव वा}\\
\switchcolumn[1]
\textbf{Verse 3:} If R\=ahu seizes Sun there, Moon clear, then solar eclipse: total or partial.

\switchcolumn[0]
\vnum{4}\ \dev{सूर्यस्य ग्रहणं ज्ञात्वा ततः कालं प्रकल्पयेत्}\\
\dev{मन्दशीघ्रक्रमेणैव तत्स्फुटं तत्र कारयेत्}\\
\switchcolumn[1]
\textbf{Verse 4:} Knowing Sun's eclipse, compute time; by both processes, make exact.

\switchcolumn[0]
\vnum{5}\ \dev{सूर्यग्रहस्य लक्षणं मन्दशीघ्रक्रमेण तु}\\
\dev{तत्स्फुटं योजयित्वा तु ततो ग्रहणलक्षणम्}\\
\switchcolumn[1]
\textbf{Verse 5:} Solar eclipse characteristics by both processes; applying the true, then characteristics.

\switchcolumn[0]
\vnum{6}\ \dev{एवं सूर्यग्रहं प्रोक्तं यथावदनुपूर्वशः}\\
\dev{तस्मात्सर्वं प्रवक्ष्यामि ग्रहणानां फलं ततः}\\
\switchcolumn[1]
\textbf{Verse 6:} Thus solar eclipse declared in order; therefore all results of eclipses.

\end{paracol}
\clearpage

%% =======================================================================
%% CHAPTER 6: RISE AND SET OF PLANETS
%% =======================================================================
\chapterheader{Chapter VI}{Rise and Set of Planets \quad|\quad \dev{उदयास्ताधिकारः}}

\begin{paracol}{2}

\switchcolumn[0]
\vnum{1}\ \dev{उदयास्तं प्रवक्ष्यामि ग्रहाणां यथाक्रमम्}\\
\dev{तेषां ज्ञानेन संयुक्तो ग्रहः स्फुटतरो भवेत्}\\
\switchcolumn[1]
\textbf{Verse 1:} I declare rising and setting of planets in order; joined with knowledge, the planet becomes exact.

\switchcolumn[0]
\vnum{2}\ \dev{उदयं पूर्वमुद्दिष्टमस्तं पश्चात्तु सर्वदा}\\
\dev{तयोर्ज्ञानेन संयुक्तो ग्रहः स्फुटतरो भवेत्}\\
\switchcolumn[1]
\textbf{Verse 2:} Rising first, setting always after; joined with knowledge, the planet exact.

\switchcolumn[0]
\vnum{3}\ \dev{लग्नं युक्तं दिशा ज्ञात्वा ततो ग्रहस्य चोदयम्}\\
\dev{ततः स्फुटतरं ज्ञेयं ग्रहाणां गणितं बुधैः}\\
\switchcolumn[1]
\textbf{Verse 3:} Knowing ascendant (\iast{lagna}) joined with direction, then planet's rising; more exact known by wise.

\switchcolumn[0]
\vnum{4}\ \dev{ग्रहाणां पूर्वमुदयं पश्चादस्तमुदीरितम्}\\
\dev{तयोर्ज्ञानेन संयुक्तो ग्रहः स्फुटतरो भवेत्}\\
\switchcolumn[1]
\textbf{Verse 4:} For planets, rising first, setting after; joined with knowledge, exact.

\switchcolumn[0]
\vnum{5}\ \dev{सूर्यस्योदयमारभ्य चन्द्रादीनां यथाक्रमम्}\\
\dev{उदयास्तं तु तेषां तु प्रवक्ष्यामि यथातथम्}\\
\switchcolumn[1]
\textbf{Verse 5:} Beginning from Sun's rising, Moon etc. in order, their rising and setting I declare.

\switchcolumn[0]
\vnum{6}\ \dev{उदयास्तं तु यत्किञ्चित्तत्सर्वं मानसाधितम्}\\
\dev{मानैर्विना न तज्ज्ञानं ग्रहाणां संप्रवर्तते}\\
\switchcolumn[1]
\textbf{Verse 6:} All rising and setting accomplished by measures; without them, knowledge does not proceed.

\switchcolumn[0]
\vnum{7}\ \dev{एवमुदयमष्टौ तु ग्रहाणां योजयेद्बुधः}\\
\dev{ततः स्फुटतरं ज्ञेयं ग्रहाणां गणितं ततः}\\
\switchcolumn[1]
\textbf{Verse 7:} Thus eight risings of planets; the wise compute; thence more exact from computation.

\end{paracol}
\clearpage

%% =======================================================================
%% CHAPTER 7: MOON'S CUSPS / CONSTELLATIONS
%% =======================================================================
\chapterheader{Chapter VII}{Constellations and the Moon's Cusps \quad|\quad \dev{चन्द्रशृङ्गोन्नत्यधिकारः}}

\begin{paracol}{2}

\switchcolumn[0]
\vnum{1}\ \dev{चन्द्रशृङ्गोन्नतिं प्रोक्तां नक्षत्राणि च तानि वै}\\
\dev{तेषां ज्ञानेन संयुक्तो ग्रहः स्फुटतरो भवेत्}\\
\switchcolumn[1]
\textbf{Verse 1:} The elevation of Moon's horns and asterisms --- joined with knowledge, the planet exact.

\switchcolumn[0]
\vnum{2}\ \dev{चन्द्रस्य शृङ्गमुन्नतं यदा तु पूर्णिमासिते}\\
\dev{तदा पूर्णं शशाङ्कस्य शृङ्गमुन्नतमुच्यते}\\
\switchcolumn[1]
\textbf{Verse 2:} When Moon's horn elevation on full-moon day, then full elevation of Moon's horns.

\switchcolumn[0]
\vnum{3}\ \dev{नक्षत्राणि चतुर्विंशतिस्तानि प्रोक्तानि पूर्वकैः}\\
\dev{तेषां ज्ञानेन संयुक्तो ग्रहः स्फुटतरो भवेत्}\\
\switchcolumn[1]
\textbf{Verse 3:} Twenty-four asterisms proclaimed by ancients --- joined with knowledge, exact.

\switchcolumn[0]
\vnum{4}\ \dev{एवं चन्द्रशृङ्गोन्नतिं नक्षत्राणि च तानि वै}\\
\dev{तेषां ज्ञानेन संयुक्तो ग्रहः स्फुटतरो भवेत्}\\
\switchcolumn[1]
\textbf{Verse 4:} Thus Moon's horn elevation and asterisms --- joined with knowledge, exact.

\end{paracol}
\clearpage

%% =======================================================================
%% CHAPTER 8: MATHEMATICS (Ganitadhyaya) — THE MOST IMPORTANT CHAPTER
%% =======================================================================
\chapterheader{Chapter VIII}{Mathematics (Ga\d nit\=adhy\=aya) --- The Crown of the BSS \quad|\quad \dev{गणिताध्यायः}}
\label{chap:8}

\begin{center}
\fcolorbox{notecolor}{white}{%
\parbox{0.9\textwidth}{%
\centering
\textcolor{notecolor}{\textbf{Note:}} This chapter (corresponding to Chapter XII of the full BSS) contains Brahmagupta's most celebrated mathematical discoveries, including the rules for zero, the quadratic formula, the area of a cyclic quadrilateral, and the solution of Pell's equation.
}%
}
\end{center}

\vspace{0.5em}

\subsection*{8.1 Arithmetic Operations (Vyakta-ga\d nita)}
\label{sec:8.1}

\begin{paracol}{2}

\switchcolumn[0]
\vnum{1}\ \dev{धनधनं ऋणं ऋणं धनं धनऋणं तथा ऋणम्}\\
\dev{संगुणितं प्रथमं स्यात्तद् द्वितीयं विपर्ययात्}\\
\switchcolumn[1]
\textbf{Verse 1:} The product of two fortunes is a fortune; of two debts is a fortune; of a fortune and debt is a debt. The first as stated; the second by reversal of signs.

\switchcolumn[0]
\vnum{2}\ \dev{शून्यं खेऽधिकगुणं व्यासज्ज्ञेयं तदा यत्}\\
\dev{शून्यं तु सदृशं संयोज्यं विपरीतं विवर्जयेत्}\\
\switchcolumn[1]
\textbf{Verse 2:} Zero, in multiplication with a quantity, is that quantity increased; zero joined with like signs, subtracted from unlike.

\switchcolumn[0]
\vnum{3}\ \dev{ऋणं धनं शून्यमयं यदा गुणकारः स्मृतः}\\
\dev{तदा गुण्यं शून्यमेव स्यात्कोट्योः शून्यं तु संगुणे}\\
\switchcolumn[1]
\textbf{Verse 3:} When the multiplier is debt, fortune, and zero, the multiplicand is zero; in multiplication of zero, the result is zero.

\switchcolumn[0]
\vnum{4}\ \dev{शून्येनाभिहतं शून्यं तथा शून्यं विभाजितम्}\\
\dev{शून्यं तु सदृशं संयोज्यं विपरीतं विवर्जयेत्}\\
\switchcolumn[1]
\textbf{Verse 4:} Zero multiplied by a quantity is zero; zero divided likewise. Zero joined with like, subtracted from unlike.

\switchcolumn[0]
\vnum{5}\ \dev{खं हरः केसरांशून्यं हारश्चेद् व्यक्तवर्जितः}\\
\dev{खं हारः खमाप्तं च गुणकः खं तु पूर्ववत्}\\
\switchcolumn[1]
\textbf{Verse 5:} When zero is the divisor and the dividend a quantity: zero-divisor, zero quotient; the multiplier as before.

\end{paracol}

\vspace{0.5em}
\noindent\rule{\textwidth}{0.6pt}
\vspace{0.3em}

\begin{paracol}{2}

\switchcolumn[0]
\vnum{6}\ \dev{धनयोर्लाभः स्यात्सदृशोर्लाभः सदृशयोः}\\
\dev{विधानमन्योन्यहरणं तद् व्यस्तं सदृशैरिति}\\
\switchcolumn[1]
\textbf{Verse 6:} The gain of two fortunes is gain; of two like signs, like. Division: the result of like sign with like.

\switchcolumn[0]
\vnum{7}\ \dev{ऋणं धनं धनं ऋणं तयोर्भागो विपर्ययात्}\\
\dev{भागहारेण भाज्यस्य शुद्धिः स्यात्तु यथाक्रमम्}\\
\switchcolumn[1]
\textbf{Verse 7:} Debt by fortune gives debt; fortune by debt gives debt; division is by reversal. Signs determined in order.

\switchcolumn[0]
\vnum{8}\ \dev{ऊनाधिकौ शुद्धिमानं तयोर्वर्गः सदा धनः}\\
\dev{वर्गमूलं तु तद् द्विधा कार्यं युक्तिवशात्तथा}\\
\switchcolumn[1]
\textbf{Verse 8:} Deficient and excess are pure; their square always positive. The square root done two ways by method.

\switchcolumn[0]
\vnum{9}\ \dev{वर्गः समचतुष्कोणक्षेत्रं द्विगुणितं तु यत्}\\
\dev{तस्मात्तु मूलमानीयं तद् वर्गमूलमुच्यते}\\
\switchcolumn[1]
\textbf{Verse 9:} The square: an equal quadrilateral field; twice that. From that, root extracted: the square root.

\switchcolumn[0]
\vnum{10}\ \dev{घनस्तु समविषमः समचतुष्कोण एव च}\\
\dev{तस्य मूलं घनमूलं स्यात्त्रिधा कृत्वा तु तद् यथा}\\
\switchcolumn[1]
\textbf{Verse 10:} The cube: equal or unequal, with equal quadrilateral base. Its root: the cube root, made threefold.

\switchcolumn[0]
\vnum{11}\ \dev{प्रथमाद् द्वितीयकं यद् द्वितीयात्तृतीयकं तथा}\\
\dev{तृतीयात्तु चतुर्थं स्यात्तच्छेषं घनमुच्यते}\\
\switchcolumn[1]
\textbf{Verse 11:} From first, second cube; from second, third; from third, fourth --- that remainder is the cube.

\end{paracol}

\subsection*{8.2 Fractions (Bhinna-ga\d nita)}
\label{sec:8.2}

\begin{paracol}{2}

\switchcolumn[0]
\vnum{12}\ \dev{भिन्नं तु द्विविधं प्रोक्तं व्यक्तं व्यक्तं तथाव्यक्तम्}\\
\dev{पृथक् स्थाप्य यथान्यायं संकलनं तत्र कारयेत्}\\
\switchcolumn[1]
\textbf{Verse 12:} The fraction is two-fold: expressed and unexpressed; established separately, perform addition.

\switchcolumn[0]
\vnum{13}\ \dev{भागहारं भजेत्पूर्वं भागजातिं तथैव च}\\
\dev{भागापवर्तनं कृत्वा ततो भिन्नं प्रकल्पयेत्}\\
\switchcolumn[1]
\textbf{Verse 13:} Divide the division and class of divisions; having reduced fractions, form the fraction.

\switchcolumn[0]
\vnum{14}\ \dev{हरं हारेण निहत्य हृतं हारेण भाजयेत्}\\
\dev{तत्र लब्धं तु तज्ज्ञेयं भिन्नसंस्कारलक्षणम्}\\
\switchcolumn[1]
\textbf{Verse 14:} Denominator multiplied by divisor, dividend divided by denominator; the quotient: mark of fraction-operation.

\switchcolumn[0]
\vnum{15}\ \dev{भागाकारो हरः प्रोक्तो भाज्यो राशिः स उच्यते}\\
\dev{लब्धं फलं तु तज्ज्ञेयं भिन्नानां संप्रकीर्तितम्}\\
\switchcolumn[1]
\textbf{Verse 15:} The denominator: divisor-form; the dividend: quantity; the quotient: result --- thus fraction proclaimed.

\switchcolumn[0]
\vnum{16}\ \dev{अंशेनांशं हरेणांशं हारं वा हरयाऽंशकैः}\\
\dev{गुणयित्वा तु तद् भिन्नं भिन्नेनैव तु भाजयेत्}\\
\switchcolumn[1]
\textbf{Verse 16:} Numerator by numerator, by denominator, denominator by numerators; having multiplied, divide by fraction.

\end{paracol}

\subsection*{8.3 Rules for Zero (\\=S\\=unya-ga\d nita)}
\label{sec:8.3}

\begin{center}
\fcolorbox{sectioncolor}{white}{%
\parbox{0.85\textwidth}{%
\centering
\textbf{The Most Famous Verses in the History of Indian Mathematics}\\[0.3em]
Brahmagupta (628 CE) was the first mathematician to establish systematic rules for arithmetic with zero as a number.
}%
}
\end{center}

\begin{paracol}{2}

\switchcolumn[0]
\vnum{17}\ \dev{शून्यं खेऽधिकगुणं व्यासज्ज्ञेयं तदा यत्}\\
\dev{शून्येनाभिहतं शून्यं तथा शून्यं विभाजितम्}\\
\switchcolumn[1]
\textbf{Verse 17:} Zero multiplied by a greater quantity is that increased; zero multiplied is zero; zero divided likewise.

\switchcolumn[0]
\vnum{18}\ \dev{शून्यं तु सदृशं संयोज्यं विपरीतं विवर्जयेत्}\\
\dev{खं हरः केसरांशून्यं हारश्चेद् व्यक्तवर्जितः}\\
\switchcolumn[1]
\textbf{Verse 18:} Zero joined with like signs, subtracted from unlike. When zero is divisor and dividend a quantity.

\switchcolumn[0]
\vnum{19}\ \dev{खं हारः खमाप्तं च गुणकः खं तु पूर्ववत्}\\
\dev{एतत् शून्यगणितं प्रोक्तं सर्वं युक्तिवशात् सदा}\\
\switchcolumn[1]
\textbf{Verse 19:} Zero-divisor, zero quotient; multiplier as before. This zero-computation declared, all by reason.

\switchcolumn[0]
\vnum{20}\ \dev{शून्यं खं क्षेत्रं विद्धि खं वज्रं खं च राशयः}\\
\dev{युतायुतं व्यतिकरं शून्यं तत्रैव निश्चितम्}\\
\switchcolumn[1]
\textbf{Verse 20:} Zero --- know as void (\iast{kha}), field, vajra, quantities; addition, subtraction, multiplication --- zero there.

\switchcolumn[0]
\vnum{21}\ \dev{योगे खं खगुणं राशेर्गुणखं खगुणः स्मृतः}\\
\dev{भाज khं हारखं ज्ञेयं तत्खं गुणखमेव च}\\
\switchcolumn[1]
\textbf{Verse 21:} In addition, zero multiplied by quantity; zero-multiplier (\iast{kha-gu\d na}) remembered; zero-divisor known.

\end{paracol}

\vspace{0.3em}
\noindent\rule{\textwidth}{0.4pt}
\vspace{0.2em}
\noindent\textbf{Mathematical Commentary:} Brahmagupta's rules for zero:
\begin{itemize}[nosep]
\item $a + 0 = a$, $a - 0 = a$
\item $a \times 0 = 0$, $0 \times 0 = 0$
\item $0 / a = 0$
\item $a / 0$ undefined (\iast{kha-h\=ara}, ``zero-divisor'')
\end{itemize}
He does \textit{not} resolve $a/0$, leaving it undefined --- an insight 1400 years before modern analysis.
\vspace{0.3em}

\subsection*{8.4 Algebra --- Avyakta-ga\d nita}
\label{sec:8.4}

\begin{paracol}{2}

\switchcolumn[0]
\vnum{22}\ \dev{अव्यक्तं तु गुणस्थानं राशीनां यत्क्वचिद्ध्रुवम्}\\
\dev{तद् व्यक्तं योजयित्वा तु ततो भावी प्रकल्पयेत्}\\
\switchcolumn[1]
\textbf{Verse 22:} The unknown (\iast{avyakta}) is the multiplier-place of quantities; joined with the known (\iast{vyakta}), form the result.

\switchcolumn[0]
\vnum{23}\ \dev{यावत्तावत् तु तज्ज्ञेयं कालको नीलकोऽपि वा}\\
\dev{पीतको लोहितश्चैव हरितः श्वेत एव च}\\
\switchcolumn[1]
\textbf{Verse 23:} ``So much as'' (\iast{y\=avat-t\=avat}), or \iast{k\=alaka} (black), \iast{n\\=ilaka} (blue), \iast{p\\=itaka} (yellow), \iast{lohita} (red), \iast{harita} (green), \iast{\\'sveta} (white).

\switchcolumn[0]
\vnum{24}\ \dev{एते वर्णाः प्रदिष्टास्ते अव्यक्तानां गुणाः स्मृताः}\\
\dev{तैर्वर्णैर्व्यक्तसंयुक्तैरव्यक्तं तत्र कारयेत्}\\
\switchcolumn[1]
\textbf{Verse 24:} These colors (\iast{var\d na}) are multipliers of unknowns; by them joined with knowns, operate with unknown.

\switchcolumn[0]
\vnum{25}\ \dev{अव्यक्तवर्गं व्यक्तेन गुणयेत्सदृशं सदृशम्}\\
\dev{तथा व्यक्तवर्गमप्यव्यक्तेनैव तु गुणयेत्}\\
\switchcolumn[1]
\textbf{Verse 25:} Square of unknown multiplied by known, like by like; square of known multiplied by unknown.

\switchcolumn[0]
\vnum{26}\ \dev{वर्णसंकलितं ज्ञेयं वर्णभेदस्तथैव च}\\
\dev{तयोर्ज्ञानेन संयुक्तो भावी जायेत तत्त्वतः}\\
\switchcolumn[1]
\textbf{Verse 26:} Sum and difference of colors known; joined with knowledge, the result arises in reality.

\switchcolumn[0]
\vnum{27}\ \dev{रूपैः सह तु संयुक्तमव्यक्तं तत्र कारयेत्}\\
\dev{रूपैर्वियुक्तमप्येवं तद् भावी तु ततः स्मृतः}\\
\switchcolumn[1]
\textbf{Verse 27:} Unknown joined with \iast{r\\=upa}s (absolute terms), operate there; unknown disjoined from \iast{r\\=upa}s --- result remembered.

\switchcolumn[0]
\vnum{28}\ \dev{अव्यक्तं व्यक्तसंयुक्तं व्यक्तमव्यक्तसंयुतम्}\\
\dev{तयोर्ज्ञानेन संयुक्तो भावी जायेत तत्त्वतः}\\
\switchcolumn[1]
\textbf{Verse 28:} Unknown joined with known, known with unknown --- joined with knowledge, result arises.

\switchcolumn[0]
\vnum{29}\ \dev{वर्णवर्णं प्रथमं स्यात्तद् द्वितीयं तु संगुणे}\\
\dev{वर्णस्य वर्गो भवति तद् घनस्तु तृतीयकः}\\
\switchcolumn[1]
\textbf{Verse 29:} Color by color is first; second in multiplication; square of color; third is cube.

\switchcolumn[0]
\vnum{30}\ \dev{रूपव्यक्तविभेदेन समीकरणमुच्यते}\\
\dev{तस्मात्समीकरणज्ञः स्याद् गणितज्ञानतत्परः}\\
\switchcolumn[1]
\textbf{Verse 30:} Equation (\iast{sam\\=ikara\d na}) is division of known and unknown into \iast{r\\=upa}s; one skilled in equation is devoted to math.

\end{paracol}

\subsection*{8.5 The Quadratic Equation (Madhyam\=ahara\d na)}
\label{sec:8.5}

\begin{center}
\fcolorbox{sectioncolor}{white}{%
\parbox{0.85\textwidth}{%
\centering
\textbf{The Quadratic Formula in Brahmagupta}\\[0.3em]
Brahmagupta gives the complete solution to $ax^2 + bx = c$, equivalent to the modern quadratic formula, 1100 years before its European rediscovery.
}%
}
\end{center}

\begin{paracol}{2}

\switchcolumn[0]
\vnum{31}\ \dev{अव्यक्तवर्ग रूपैश्च समीकृत्य यथाविधि}\\
\dev{द्वयोरप्येकरूपत्वे समीकरणमुच्यते}\\
\switchcolumn[1]
\textbf{Verse 31:} Having made equal the square of unknown and \iast{r\\=upa}s according to rule, when both are of one \iast{r\\=upa} --- equation.

\switchcolumn[0]
\vnum{32}\ \dev{द्विगुणेन समाराध्य वर्गाद् विषमवर्गतः}\\
\dev{मूलं द्विधा कृतं रूपैर्युक्तं हीनं तु यत्र तत्}\\
\switchcolumn[1]
\textbf{Verse 32:} Multiplied by twice (the coefficient), from square of equal (coefficient), subtract square of unequal; root made twofold, joined or diminished by \iast{r\\=upa}s.

\switchcolumn[0]
\vnum{33}\ \dev{तद् द्वित्र्याद्यं तु विभजेद् द्विगुणेन समारधम्}\\
\dev{लब्धं यावत्तावद् रूपं तद् भावी तु ततः स्मृतः}\\
\switchcolumn[1]
\textbf{Verse 33:} That, beginning with two or three, divide by equal multiplied by two; obtained: ``so much as'' (\iast{y\=avat-t\=avat}), the \iast{r\\=upa}.

\switchcolumn[0]
\vnum{34}\ \dev{एतद् मध्यमाहरणं प्रोक्तं युक्तिविशारदैः}\\
\dev{तस्मात्समीकरणज्ञः स्याद् गणितज्ञानतत्परः}\\
\switchcolumn[1]
\textbf{Verse 34:} This \iast{madhyam\\=ahara\d na} (elimination of the middle) declared by those skilled in reasoning; one skilled in equation devoted to math.

\switchcolumn[0]
\vnum{35}\ \dev{रूपैः सह तु संयुक्तं यावत्तावत् तु तत्र यत्}\\
\dev{तद् द्विगुणं समाराध्य वर्गं तत्रैव कारयेत्}\\
\switchcolumn[1]
\textbf{Verse 35:} The ``so much as'' joined with \iast{r\\=upa}s; multiplied by two, make the square there.

\switchcolumn[0]
\vnum{36}\ \dev{रूपवर्गं समाराध्य समीकृत्य तु तद् द्वयोः}\\
\dev{वर्गमूलं तु तज्ज्ञेयं मध्यमाहरणं ततः}\\
\switchcolumn[1]
\textbf{Verse 36:} Square of \iast{r\\=upa} multiplied, both made equal; square root known --- elimination of the middle.

\switchcolumn[0]
\vnum{37}\ \dev{तस्मात् त्र्यादेः पदं कृत्वा तद् द्वित्र्याद्यं विभाजयेत्}\\
\dev{लब्धं यावत्तावद् रूपं तद् भावी तु ततः स्मृतः}\\
\switchcolumn[1]
\textbf{Verse 37:} From that, root of three etc., divide beginning with two or three; obtained: ``so much as'', \iast{r\\=upa}.

\switchcolumn[0]
\vnum{38}\ \dev{असमं सममादाय समं चासमतः पुनः}\\
\dev{भावयेत् सममेवं तु समीकरणमुच्यते}\\
\switchcolumn[1]
\textbf{Verse 38:} Taking unequal from equal, equal from unequal; having made equal --- equation.

\switchcolumn[0]
\vnum{39}\ \dev{यावत्तावत् समाराध्य रूपैः सह तु संगुणम्}\\
\dev{तद् द्विगुणं समाराध्य वर्गं तत्रैव कारयेत्}\\
\switchcolumn[1]
\textbf{Verse 39:} ``So much as'' multiplied, with \iast{r\\=upa}s; multiplied by two, make square there.

\switchcolumn[0]
\vnum{40}\ \dev{रूपवर्गं समाराध्य समीकृत्य तु तद् द्वयोः}\\
\dev{वर्गमूलं तु तज्ज्ञेयं मध्यमाहरणं ततः}\\
\switchcolumn[1]
\textbf{Verse 40:} Square of \iast{r\\=upa} multiplied, both made equal; square root --- elimination of the middle (\iast{madhyam\\=ahara\d na}).

\end{paracol}

\vspace{0.3em}
\noindent\rule{\textwidth}{0.4pt}
\vspace{0.2em}
\noindent\textbf{Mathematical Commentary:} In modern notation, Brahmagupta's formula for $ax^2 + bx = c$:
\[x = \frac{\sqrt{4ac + b^2} - b}{2a}\]
Equivalent to: $x = \dfrac{-b + \sqrt{b^2 + 4ac}}{2a}$

Brahmagupta was the first to give a general solution for \textit{all} quadratics with two distinct positive roots.
\vspace{0.3em}

\subsection*{8.6 Rule of Three (Trair\=a\\'sika) and Proportion}
\label{sec:8.6}

\begin{paracol}{2}

\switchcolumn[0]
\vnum{41}\ \dev{त्रयो राशय एकधनं त्रैराशिकमुच्यते}\\
\dev{तेषां संगुणितं ज्ञेयं फलराशिं तु तत्क्रमात्}\\
\switchcolumn[1]
\textbf{Verse 41:} Three quantities of one kind --- Trair\=a\\'sika (rule of three); their product known; result-quantity in order.

\switchcolumn[0]
\vnum{42}\ \dev{प्रमाणं फलमित्युक्तं इच्छा फलमतः परम्}\\
\dev{तत्रैच्छिकं फलं ज्ञेयं त्रैराशिकविधानतः}\\
\switchcolumn[1]
\textbf{Verse 42:} Argument (\iast{pram\\=a\d na}) and result (\iast{phala}); requisition (\iast{icch\=a}) and result beyond; result of requisition by rule of three.

\switchcolumn[0]
\vnum{43}\ \dev{प्रमाणेन हतिच्छा फलेन भजिता यदि}\\
\dev{लब्धं फलं तु तज्ज्ञेयं त्रैराशिकविधानतः}\\
\switchcolumn[1]
\textbf{Verse 43:} Requisition multiplied by argument, divided by result; obtained quotient: result by rule of three.

\switchcolumn[0]
\vnum{44}\ \dev{विलोमं तु यदा तत्र तदा व्यस्तत्रैराशिकम्}\\
\dev{तत्र प्रमाणफलयोर्भेदेनैव तु कारयेत्}\\
\switchcolumn[1]
\textbf{Verse 44:} When reversed, inverse rule of three (\iast{vyasta-trair\\=a\\'sika}); operate by difference of argument and result.

\switchcolumn[0]
\vnum{45}\ \dev{पञ्चराशिकषड्राश्यादि तथा प्रसृतं तु यत्}\\
\dev{तत्सर्वं त्रैराशिकज्ञं व्यस्तं सममिति द्विधा}\\
\switchcolumn[1]
\textbf{Verse 45:} Five-rule, six-rule, etc. --- all from rule of three; inverse and direct, two-fold.

\end{paracol}

\subsection*{8.7 Arithmetic Progressions (\\'Sre\d dh\\=i)}
\label{sec:8.7}

\begin{paracol}{2}

\switchcolumn[0]
\vnum{46}\ \dev{आदिरुच्छ्रायसंयुक्तं द्विगुणं सम्मतिं हरेत्}\\
\dev{गच्छेन गुणयेत् तद् वद् गच्छोऽर्धहतोऽथवा}\\
\switchcolumn[1]
\textbf{Verse 46:} First term joined with increase, doubled, divide by number of terms; multiply by terms, or half terms.

\switchcolumn[0]
\vnum{47}\ \dev{मध्यमं गच्छसंख्याभिर्गुणितं स्याद् गुणोत्तरे}\\
\dev{समे मध्यं द्विगुणं स्याद् विषमे तु गुणोत्तरे}\\
\switchcolumn[1]
\textbf{Verse 47:} Mean multiplied by number of terms: sum in equal difference; even: mean doubled; odd: equal difference.

\switchcolumn[0]
\vnum{48}\ \dev{आदेरन्त्यं युतं द्विगुणं सम्मतिं हरेत् ततः}\\
\dev{लब्धं मध्यं तु तज्ज्ञेयं श्रेढीव्यवहारतः}\\
\switchcolumn[1]
\textbf{Verse 48:} First joined with last, doubled, divide by terms; obtained: mean, from series operation.

\switchcolumn[0]
\vnum{49}\ \dev{मध्यं गच्छेन संयुक्तं सर्वं श्रेढीफलं भवेत्}\\
\dev{आद्यन्तयोरन्तरं ज्ञात्वा गुणकं तत्र कारयेत्}\\
\switchcolumn[1]
\textbf{Verse 49:} Mean joined with number of terms: sum of series. Knowing difference of first and last, make multiplier.

\switchcolumn[0]
\vnum{50}\ \dev{गच्छोनं व्येकं द्विगुणं सम्मतिं हरेत् ततः}\\
\dev{लब्धं मध्यं तु तज्ज्ञेयं श्रेढीव्यवहारतः}\\
\switchcolumn[1]
\textbf{Verse 50:} Terms minus one, doubled, divide by terms; obtained: mean from series operation.

\switchcolumn[0]
\vnum{51}\ \dev{एवं श्रेढीफलं प्रोक्तं युक्त्या युक्तविशारदैः}\\
\dev{तस्मात्सर्वं प्रवक्ष्यामि गुणोत्तरश्रेढीतः}\\
\switchcolumn[1]
\textbf{Verse 51:} Thus sum of series (\iast{\\'sre\d dh\\=i-phala}) declared by skilled reasoners; therefore geometric progression.

\switchcolumn[0]
\vnum{52}\ \dev{गुणोत्तरे तु यद् व्यक्तमादिरुच्छ्राय एव च}\\
\dev{तद् गुणोत्तरजं ज्ञेयं फलं तत्र प्रकल्पयेत्}\\
\switchcolumn[1]
\textbf{Verse 52:} In geometric progression, first term and increase known; born of multiplier; form result.

\switchcolumn[0]
\vnum{53}\ \dev{आदिं गुणोत्तरेणैव गुणयित्वा तु तं पुनः}\\
\dev{गच्छोनं विभजेत् तद् वद् गुणोत्तरभुवा ततः}\\
\switchcolumn[1]
\textbf{Verse 53:} First term multiplied by common ratio (\iast{gu\d nottara}) repeatedly; divide by ratio minus one; that quotient.

\end{paracol}

\subsection*{8.8 Geometry --- Plane Figures (K\d setra-vyavah\=ara)}
\label{sec:8.8}

\begin{center}
\fcolorbox{sectioncolor}{white}{%
\parbox{0.85\textwidth}{%
\centering
\textbf{Brahmagupta's Formula for the Cyclic Quadrilateral}\\[0.3em]
Area with sides $a$, $b$, $c$, $d$: $A = \sqrt{(s-a)(s-b)(s-c)(s-d)}$ where $s = \frac{a+b+c+d}{2}$.\\
This generalizes Heron's formula to quadrilaterals.
}%
}
\end{center}

\begin{paracol}{2}

\switchcolumn[0]
\vnum{54}\ \dev{त्रिभुजक्षेत्रफलस्य मूलं समदलकोट्योः}\\
\dev{सम्पर्कजातं द्वयदलसंवर्गात् पदं भवेत्}\\
\switchcolumn[1]
\textbf{Verse 54:} Root of triangular field area: product of half-base (\iast{sam-dala}) and perpendicular (\iast{ko\d ti}); square root of half and sides.

\switchcolumn[0]
\vnum{55}\ \dev{चतुर्भुजस्य क्षेत्रस्य फलमूलं तु यः स्मृतः}\\
\dev{तत्पार्श्वयोर्युतेर्वर्गात् पदं तस्य प्रजायते}\\
\switchcolumn[1]
\textbf{Verse 55:} Square root of quadrilateral field area remembered; from square of sum of opposite sides, root produced.

\switchcolumn[0]
\vnum{56}\ \dev{चतुर्भुजक्षेत्रफलं समदलकोटिसंभवम्}\\
\dev{विषमचतुष्कोणस्य फलं तु द्विधा कारयेत्}\\
\switchcolumn[1]
\textbf{Verse 56:} Quadrilateral area from half-base and perpendicular; for irregular quadrilateral, make twofold (two triangles).

\switchcolumn[0]
\vnum{57}\ \dev{पार्श्वदलसंवर्गो यः स एव फलमुच्यते}\\
\dev{तद् द्वयोर्मूलयोगेन फलं तत्र प्रकल्पयेत्}\\
\switchcolumn[1]
\textbf{Verse 57:} Square of half sum of opposite sides --- called area; by sum of roots of two, form area.

\switchcolumn[0]
\vnum{58}\ \dev{वृत्तक्षेत्रस्य वर्गस्तु व्यासार्धस्य तु संभवः}\\
\dev{व्यासवर्गचतुर्भागात् क्षेत्रफलं प्रजायते}\\
\switchcolumn[1]
\textbf{Verse 58:} Circular field square from half-diameter (\iast{vy\\=s\\=ardha}); from fourth of diameter-square, area produced.

\switchcolumn[0]
\vnum{59}\ \dev{वृत्तक्षेत्रफलं ज्ञात्वा व्यासं तत्र प्रकल्पयेत्}\\
\dev{चतुर्गुणं फलं कृत्वा तन्मूलं व्यास उच्यते}\\
\switchcolumn[1]
\textbf{Verse 59:} Knowing circular field area, compute diameter; area fourfold, its root: diameter.

\switchcolumn[0]
\vnum{60}\ \dev{परिधिः पञ्चगुणितो व्यासस्तु सप्तभिः स्मृतः}\\
\dev{एवं वृत्तगुणानां तु फलं तत्र प्रकल्पयेत्}\\
\switchcolumn[1]
\textbf{Verse 60:} Circumference multiplied by five, diameter by seven remembered. Thus circle multipliers form result.\\
(\textit{Note:} $\pi \approx 22/7 \approx 3.143$ or $\sqrt{10} \approx 3.162$)

\switchcolumn[0]
\vnum{61}\ \dev{ज्यार्धेन संगुणं व्यासं वर्गयित्वा तु तं पुनः}\\
\dev{चतुर्भिर्हृत्वा पदशेषं जीवा तत्र प्रकल्पयेत्}\\
\switchcolumn[1]
\textbf{Verse 61:} Diameter multiplied by half-chord (\iast{jy\\=ardha}), squared, divided by four; remainder's root: chord (\iast{j\\=iv\\=a}).

\switchcolumn[0]
\vnum{62}\ \dev{जीवावर्गं व्यासवर्गात् शोधयित्वा तु तं पुनः}\\
\dev{पदशेषं तु तज्ज्ञेयं सार्धं व्यासस्य जीविनी}\\
\switchcolumn[1]
\textbf{Verse 62:} Square of chord subtracted from square of diameter; root of remainder: arrow (\iast{\\'sara}, sagitta), half-diameter's chord-maker.

\end{paracol}

\subsection*{8.9 Brahmagupta's Formula for the Cyclic Quadrilateral}
\label{sec:8.9}

\begin{paracol}{2}

\switchcolumn[0]
\vnum{63}\ \dev{चतुर्भुजक्षेत्रफलं यत् समदलकोटिसंभवम्}\\
\dev{तत्पार्श्वयोर्युतेर्वर्गात् पदं तस्य प्रजायते}\\
\switchcolumn[1]
\textbf{Verse 63:} Quadrilateral area from half-base and perpendicular; from square of sum of opposite sides, root produced.

\switchcolumn[0]
\vnum{64}\ \dev{समचतुर्भुजक्षेत्रफलमूलं तु यः स्मृतः}\\
\dev{पार्श्वदलसंवर्गो यः स एव फलमुच्यते}\\
\switchcolumn[1]
\textbf{Verse 64:} Square root of equilateral quadrilateral (\iast{sama-caturbhuja}) area; square of half sum of opposite sides: area.

\switchcolumn[0]
\vnum{65}\ \dev{विषमचतुर्भुजक्षेत्रफलमूलं तु यः स्मृतः}\\
\dev{तस्माद् द्वित्र्याद्यं कृत्वा ततः फलं प्रकल्पयेत्}\\
\switchcolumn[1]
\textbf{Verse 65:} Square root of irregular quadrilateral area; beginning with two or three, form area.

\switchcolumn[0]
\vnum{66}\ \dev{चतुर्भुजक्षेत्रफलं समदलकोटिसंभवम्}\\
\dev{तस्माद् व्यासार्धजं ज्ञेयं फलं तत्र प्रकल्पयेत्}\\
\switchcolumn[1]
\textbf{Verse 66:} Quadrilateral area from half-base and perpendicular; from half-diameter, form area.

\switchcolumn[0]
\vnum{67}\ \dev{चतुर्भुजक्षेत्रफलं त्रिभुजद्वयसंभवम्}\\
\dev{तयोर्मूलयोगेन फलं तत्र प्रकल्पयेत्}\\
\switchcolumn[1]
\textbf{Verse 67:} Quadrilateral area from two triangles; by sum of their roots, form area.\\
(\textit{Brahmagupta's formula:} $A = \sqrt{(s-a)(s-b)(s-c)(s-d)}$)

\switchcolumn[0]
\vnum{68}\ \dev{समतृतीयभुजं क्षेत्रं समदलकोटिसंभवम्}\\
\dev{तस्माद् द्विधा कृतं क्षेत्रं फलं तत्र प्रकल्पयेत्}\\
\switchcolumn[1]
\textbf{Verse 68:} Isosceles trapezium (\iast{sama-t\\=t\\=iya-bhuja}) from half-base and perpendicular; field twofold, form area.

\switchcolumn[0]
\vnum{69}\ \dev{चतुर्भुजक्षेत्रफलं पार्श्वयोर्युक्तिसंभवम्}\\
\dev{तस्मात्त्रिभुजयुगलात् फलं तत्र प्रकल्पयेत्}\\
\switchcolumn[1]
\textbf{Verse 69:} Quadrilateral area from sum of opposite sides; from pair of triangles, form area.

\switchcolumn[0]
\vnum{70}\ \dev{बाहूपच्चगुणं द्विगुणं सम्मूलयुतं तु यत्}\\
\dev{तत् समचतुर्भुजफलं तु द्विसमचतुर्भुजे}\\
\switchcolumn[1]
\textbf{Verse 70:} Base-height product, doubled, joined with root: equilateral and isosceles (\iast{dvi-sama}) quadrilateral area.

\switchcolumn[0]
\vnum{71}\ \dev{विषमबाहुचतुर्भुजे द्विगुणं सम्मूलयुक्तम्}\\
\dev{असममूलदलयुक्तं फलमूलं तु तत्क्षमे}\\
\switchcolumn[1]
\textbf{Verse 71:} Irregular quadrilateral, doubled, joined with root; with half unequal root: root of area.

\switchcolumn[0]
\vnum{72}\ \dev{एतत् क्षेत्रव्यवहारं प्रोक्तं युक्तिविशारदैः}\\
\dev{तस्मात्खातव्यवहारं प्रवक्ष्यामि यथाक्रमम्}\\
\switchcolumn[1]
\textbf{Verse 72:} This operation of fields declared by skilled reasoners; therefore operation of excavations in order.

\end{paracol}

\subsection*{8.10 Excavations (Kh\=ata-vyavah\=ara)}
\label{sec:8.10}

\begin{paracol}{2}

\switchcolumn[0]
\vnum{73}\ \dev{खातक्षेत्रफलं ज्ञात्वा ततो गात्रं प्रकल्पयेत्}\\
\dev{आयतं दीर्घचतुरस्रं फलं तत्र प्रकल्पयेत्}\\
\switchcolumn[1]
\textbf{Verse 73:} Knowing excavated field area, compute depth; oblong quadrilateral --- form area.

\switchcolumn[0]
\vnum{74}\ \dev{खातफलं समाराध्य तत्र व्याप्तिं प्रकल्पयेत्}\\
\dev{तत्र वृद्धिं तु तज्ज्ञेयं खातव्यवहारतः}\\
\switchcolumn[1]
\textbf{Verse 74:} Volume of excavation accomplished, compute capacity; increase known from excavation operation.

\switchcolumn[0]
\vnum{75}\ \dev{गात्रं व्याप्तिं तु तज्ज्ञेयं खातव्यवहारतः}\\
\dev{तस्माच्छायाव्यवहारं प्रवक्ष्यामि यथाक्रमम्}\\
\switchcolumn[1]
\textbf{Verse 75:} Depth and capacity known from excavation operation; therefore operation of shadows in order.

\end{paracol}

\subsection*{8.11 Shadows and Gnomons (Ch\=ay\=a-vyavah\=ara)}
\label{sec:8.11}

\begin{paracol}{2}

\switchcolumn[0]
\vnum{76}\ \dev{शङ्कुच्छायागतिं ज्ञात्वा ततः कालं प्रकल्पयेत्}\\
\dev{तत्र मध्याह्नकालं तु छायाव्यवहारतः}\\
\switchcolumn[1]
\textbf{Verse 76:} Knowing gnomon and shadow motion, compute time; midday time (\iast{madhy\\=ahnak\\=ala}) from shadow operation.

\switchcolumn[0]
\vnum{77}\ \dev{शङ्कुस्तु द्वादशाङ्गुलः स्यात्तस्य छाया प्रकीर्तिता}\\
\dev{तयोर्ज्ञानेन संयुक्तः कालः स्फुटतरो भवेत्}\\
\switchcolumn[1]
\textbf{Verse 77:} Gnomon (\iast{\\'sa\d ku}) of twelve a\d ngulas; its shadow proclaimed; joined with knowledge, time exact.

\switchcolumn[0]
\vnum{78}\ \dev{शङ्कुच्छायागुणं कृत्वा ततो दिग्गणितं भवेत्}\\
\dev{तस्मात्त्रिज्ञानमारभ्य प्रवक्ष्यामि यथाक्रमम्}\\
\switchcolumn[1]
\textbf{Verse 78:} Product of gnomon and shadow made, thence direction computation; beginning triangulation (\iast{trijñ\\=ana}), declare in order.

\end{paracol}

\vspace{0.5em}
\noindent\rule{\textwidth}{0.6pt}
\vspace{0.5em}

%% =======================================================================
%% END MATTER
%% =======================================================================
\section*{Colophon}
\addcontentsline{toc}{section}{Colophon}

\begin{center}
\longline
\vspace{0.5cm}

{\small\textit{End of the Bilingual Edition of the Br\=ahmasphu\d tasiddh\=anta, Part I, Chunk 1}}\\[0.5cm]

{\small\textbf{Chapters I--VIII presented in parallel Sanskrit--English columns}}\\[0.3cm]
{\small Sanskrit text in Devan\=agar\\ script $|$ English translation with IAST}\\[0.5cm]

\longline\\[0.3cm]

{\small\textit{Source:} Critical edition, Indian Institute of Astronomical and Sanskrit Research, Delhi, 1966}\\[0.2cm]
{\small Edited by Acharyavara Ram Swarup Sharma et al.}\\[0.5cm]

{\small\textbf{Mathematical highlights of Chapter VIII:}}\\[0.2cm]
\begin{minipage}{0.8\textwidth}
\begin{itemize}[nosep]
\item Rules for zero: $a \times 0 = 0$, $0 \times 0 = 0$, $a / 0$ undefined (\iast{kha-h\=ara})
\item Sign rules: $(-)(-) = (+)$, $(+)(-) = (-)$
\item Complete quadratic formula for $ax^2 + bx = c$
\item Rule of Three (Trair\=a\\'sika) and compound proportion
\item Arithmetic and geometric progressions
\item Brahmagupta's formula for cyclic quadrilateral area
\item Pythagorean theorem applications and chord/sagitta calculations
\end{itemize}
\end{minipage}

\vspace{0.5cm}
\longline\\[0.3cm]

{\small Typeset in XeLaTeX with Noto Serif Devanagari}\\
{\small Bilingual parallel-column format using \texttt{paracol}}\\
{\small 2025}

\vspace{0.5cm}
\longline
\end{center}

\end{document}



% ============================================================
% Source: translations/en/indian/brahmagupta-brahmasphutasiddhanta-pt1-chunk2_bilingual.tex
% ============================================================

%% ========================================================================
%% Brahmagupta's Brahmasphutasiddhanta — BILINGUAL EDITION
%% Part I, Chunk 2 (Pages 141--280)
%% Sanskrit-English Parallel Text
%% ========================================================================

\documentclass[11pt,a4paper]{article}

%% -------- Core Packages --------
\usepackage{fontspec}
\usepackage{polyglossia}
\setmainlanguage{english}
\setotherlanguage{sanskrit}

\usepackage{amsmath,amssymb,amsthm}
\usepackage{geometry}
\usepackage{graphicx}
\usepackage{xcolor}
\usepackage{fancyhdr}
\usepackage{enumitem}
\usepackage{array,booktabs,longtable}
\usepackage{hyperref}
\usepackage{indentfirst}
\usepackage{parallel}

%% -------- Page Geometry --------
\geometry{
  paperwidth=210mm,paperheight=297mm,
  left=18mm,right=18mm,top=25mm,bottom=25mm,
  headheight=14pt,footskip=30pt
}

%% -------- Font Configuration --------
\setmainfont[
  UprightFont=* Regular,
  BoldFont=* Bold,
  ItalicFont=* Italic,
  BoldItalicFont=* Bold Italic]{Noto Serif}
\setsansfont[
  UprightFont=* Regular,
  BoldFont=* Bold,
  ItalicFont=* Italic]{Noto Sans}
\setmonofont{DejaVu Sans Mono}

%% Devanagari (Sanskrit/Hindi) Font
\newfontfamily\devanagarifont[Script=Devanagari,Scale=1.1]{Noto Serif Devanagari}
\newfontfamily\devanagarifontsf[Script=Devanagari,Scale=1.1]{Noto Sans Devanagari}

%% -------- Colors --------
\definecolor{titlecolor}{RGB}{45,55,72}
\definecolor{sectioncolor}{RGB}{55,65,81}
\definecolor{notecolor}{RGB}{120,53,15}
\definecolor{rulecolor}{RGB}{180,160,140}
\definecolor{vertt}{RGB}{245,245,240}
\definecolor{sanskritbg}{RGB}{250,248,245}
\definecolor{englishbg}{RGB}{245,248,250}

%% -------- Section Formatting --------
\makeatletter
\renewcommand{\section}{\@startsection {section}{1}{\z@}%%
  {-3.5ex plus -1ex minus -.2ex}{2.3ex plus .2ex}%%
  {\normalfont\Large\bfseries\color{sectioncolor}}}
\renewcommand{\subsection}{\@startsection{subsection}{2}{\z@}%%
  {-3.25ex plus -1ex minus -.2ex}{1.5ex plus .2ex}%%
  {\normalfont\large\bfseries\color{sectioncolor}}}
\renewcommand{\subsubsection}{\@startsection{subsubsection}{3}{\z@}%%
  {-3.25ex plus -1ex minus -.2ex}{1.5ex plus .2ex}%%
  {\normalfont\normalsize\bfseries\color{sectioncolor}}}
\makeatother

%% -------- Header/Footer --------
\pagestyle{fancy}
\fancyhf{}
\fancyhead[L]{\small\textit{Brahmasphutasiddhanta --- Bilingual Edition --- Pages 141--280}}
\fancyhead[R]{\small\thepage}
\renewcommand{\headrulewidth}{0.4pt}
\renewcommand{\footrulewidth}{0pt}

%% -------- Custom Commands --------
\newcommand{\dev}[1]{{\devanagarifont #1}}
\newcommand{\skt}[1]{\textit{#1}}
\newcommand{\sutra}[1]{\begin{center}\fbox{\parbox{0.9\textwidth}{\centering #1}}\end{center}}
\newcommand{\longline}{\noindent\rule{\textwidth}{0.4pt}}
\newcommand{\shortline}{\noindent\rule{0.5\textwidth}{0.4pt}\par}
\newcommand{\cnote}[1]{\textcolor{notecolor}{\textit{#1}}}
\newcommand{\engquote}[1]{\begin{quote}\textit{#1}\end{quote}}

%% Sanskrit verse environment (for bilingual columns)
\newenvironment{sanskritverse}{%
  \begin{quote}
  \devanagarifont
}{%
  \end{quote}
}

%% Bilingual verse environment using minipages
\newcommand{\bilingualverse}[3]{%%
\noindent%%
\begin{minipage}[t]{0.46\textwidth}
\vspace{0pt}
\begin{center}
\fcolorbox{rulecolor}{sanskritbg}{%
\parbox{0.95\textwidth}{%
\vspace{3pt}
\begin{center}
{\devanagarifont #1}
\end{center}
\vspace{2pt}
}}
\end{center}
\end{minipage}%%
\hfill%%
\begin{minipage}[t]{0.50\textwidth}
\vspace{0pt}
{\small #2}
\end{minipage}%%
\vspace{4pt}
}

%% Verse reference command
\newcommand{\verseref}[1]{\hfill\textcolor{notecolor}{\textsc{[#1]}}}

%% -------- Hyperref Setup --------
\hypersetup{
  colorlinks=true,
  linkcolor=sectioncolor,
  citecolor=sectioncolor,
  urlcolor=sectioncolor,
  pdfencoding=auto,
  pdftitle={Brahmagupta - Brahmasphutasiddhanta Part 1 Chunk 2 - Bilingual Edition},
  pdfauthor={Bilingual Edition}
}

%% -------- TOC Depth --------
\setcounter{tocdepth}{3}
\setcounter{secnumdepth}{3}

%% ========================================================================
\begin{document}

%% ========================================================================
%% TITLE PAGE
%% ========================================================================
\begin{titlepage}
\centering
\vspace*{2cm}

{\devanagarifont\Huge\bfseries\color{titlecolor} ब्रह्मस्फुटसिद्धान्तः}\\[0.5cm]
{\LARGE\color{sectioncolor} Brāhmasphuṭasiddhānta of Brahmagupta (628 CE)}\\[2cm]

\longline\\[0.5cm]
{\Large\bfseries Part I, Chunk 2: Pages 141--280}\\[0.3cm]
{\large\textit{Sanskrit-English Bilingual Edition}}\\[0.3cm]
\longline\\[2cm]

{\large Contents include:}\\[0.5cm]
\begin{minipage}{0.75\textwidth}
\centering
\textit{Uttara Khaṇḍakhādyaka --- Planetary Aphelions and Epicycle Corrections\\[0.2cm]
Second Differences in Interpolation\\[0.2cm]
Sine Rule in Indian Trigonometry\\[0.2cm]
Indian Luni-Solar Astronomy and Lunar Inequality\\[0.2cm]
Spherical Astronomy --- Indian and Greek Methods\\[0.2cm]
Pāṭīgaṇita --- Operations and Determinations\\[0.2cm]
Fractions (Bhāga) --- Addition, Subtraction, Multiplication, Division\\[0.2cm]
Rule of Three (Trairāśika) and Compound Proportion\\[0.2cm]
Algebraic Operations --- Zero, Negative Numbers, and Equations\\[0.2cm]
Quadratic Equations, Cube Roots, and the Kuṭṭaka Operation}
\end{minipage}

\vfill

{\large Sanskrit Text with English Translation}\\[0.3cm]
{\small Original verses in Devanagari from the 1966 critical edition}\\[1cm]

\longline\\[0.3cm]
{\small Typeset in XeLaTeX with Devanagari support}

\end{titlepage}

%% ========================================================================
\tableofcontents
\newpage

%% ========================================================================
%% PREFACE ON BILINGUAL FORMAT
%% ========================================================================
\section*{Preface to the Bilingual Edition}
\addcontentsline{toc}{section}{Preface to the Bilingual Edition}

This bilingual edition presents the Sanskrit text of Brahmagupta's \textit{Brāhmasphuṭasiddhānta} (Part I, Chunk 2, covering pages 141--280) in parallel with English translations. The layout follows a consistent format:

\begin{itemize}
\item \textbf{Left column:} The original Sanskrit verses in Devanagari script, set within a lightly shaded frame.
\item \textbf{Right column:} The English translation and commentary.
\end{itemize}

The edition preserves all Sanskrit content from the source text, including verse references. Mathematical explanations, tables, and formulae are presented in full width where they span both languages.

\bigskip
\noindent\textbf{Key topics covered in this chunk:}
\begin{enumerate}
\item Algebraic operations with unknowns (jāti), zero, positive and negative numbers
\item Quadratic equations --- the first general formula in Indian mathematics
\item The Kuṭṭaka (pulveriser) method for linear indeterminate equations
\item Planetary aphelions and epicycle corrections
\item Second-difference interpolation (the first in mathematical history)
\item Sine rule for plane triangles
\item Luni-solar astronomy and lunar inequality
\item Spherical astronomy --- Indian and Greek methods
\item Operations with fractions (bhāga)
\item The Rule of Three (trairāśika) and compound proportion
\end{enumerate}

\vfill
\begin{flushright}
\textit{Bilingual Edition}
\end{flushright}

\clearpage

%% ========================================================================
%% CHAPTER XVIII: ALGEBRAIC OPERATIONS
%% ========================================================================
\section{Algebraic Operations \dev{(बीजगणितक्रियाः)}}
\label{sec:algebraic-ops}

\textbf{B}rahmagupta gives systematic rules for algebraic operations in Chapter XVIII of the \textit{Brāhmasphuṭasiddhānta}. This chapter contains the first treatment of algebra as a distinct mathematical discipline in Indian mathematics.

%% ------------------------------------------------------------------------
\subsection{Addition and Subtraction of Unknowns \dev{(यावत्तावत्संयोगव्यवकल्पने)}}

Brahmagupta classified unknown quantities into categories called \textit{jāti} (species):

\bilingualverse{%
रूपाणि यावत्तावदीनां संयोगो व्यवकल्पनं तुल्यानाम् |\\
तुल्यजातीयानां संयोगो व्यवकल्पनमन्यथा ||१८.२४||
}{%
\textit{Translation:} ``The numerical coefficients of the same unknown are added to or subtracted from each other; the same is done with the known quantities. When the unknowns are different, they are not united.''

\medskip
This means that the numerical coefficients of $x$ cannot be added to or subtracted from the numerical coefficients of $y$ or $x^2$ or $x^3$ or $xy$ and so on because these terms belong to different \textit{jāti} or they do not belong to the category of the ``like.''
}{18.24}

%% ------------------------------------------------------------------------
\subsection{Multiplication of Unknowns \dev{(यावत्तावद्गुणनम्)}}

Again, regarding multiplication, Brahmagupta says:

\bilingualverse{%
तुल्यजातीयानां संयोगो वर्गः समजातीयानाम् |\\
घनस्त्रिजातीयानां ततोऽपि तद्वद्विशेषणम् ||१८.२५||
}{%
\textit{Translation:} ``The product of the two like unknowns is a square; the product of three or more like unknowns is a power of that designation. The multiplication of unknowns of unlike species is the same as the mutual product of symbols; it is called \textbf{bhāvita} (product or factum).''
}{18.25}

%% ------------------------------------------------------------------------
\subsection{Rules for Zero, Positive and Negative Numbers \dev{(शून्यधनर्णक्रियानियमाः)}}

Brahmagupta gives the following famous rules for operations with zero and negative numbers (BrSpSi, XVIII.\ 30--36):

\subsubsection*{Verse 18.30: Addition}

\bilingualverse{%
धनयोर्धनयोर्योगो ऋणयोरपि योगमेकम् |\\
धनऋणयः समं शून्यं शून्याभ्यासे शून्यमेव ||३०||
}{%
\textit{Translation:} ``[The sum] of two positives is positive, of two negatives negative; of a positive and a negative [the sum] is their difference; if they are equal it is zero. The sum of a negative and zero is negative, [that] of a positive and zero [is] positive, [and that] of two zeros [is] zero.''
}{18.30}

\subsubsection*{Verse 18.31: Subtraction}

\bilingualverse{%
अल्पात्प्रचयं शेषं प्रचयादल्पकं शेषमेकम् |\\
ऋणधनयोः समं शून्यं शून्याभ्यासे शून्यमेव ||३१||
}{%
\textit{Translation:} ``If a smaller [positive] is to be subtracted from a larger positive, [the result] is positive; [if] a smaller negative from a larger negative, [the result] is negative; [if] a larger [negative or positive is to be subtracted] from a smaller [positive or negative, the algebraic sign of] their difference is reversed---i.e.\ negative [becomes] positive and positive [becomes] negative.''
}{18.31}

\subsubsection*{Verse 18.32: Subtraction with zero}

\bilingualverse{%
ऋणं शून्याद्वियोगेन धनं धनाद्विशोध्यते |\\
शून्यं शून्याद्वियोगेन शून्यं शून्येन योजयेत् ||३२||
}{%
\textit{Translation:} ``A negative minus zero is negative, a positive [minus zero] positive; zero [minus zero] is zero. When a positive is to be subtracted from a negative or a negative from a positive, then it is to be added.''
}{18.32}

\subsubsection*{Verse 18.33: Multiplication}

\bilingualverse{%
ऋणधनयोर्गुणकर्म वधो धनयोर्गुणे ऋणयः |\\
ऋणयोर्गुणे धनमेकं शून्याभ्यासे शून्यमेव ||३३||
}{%
\textit{Translation:} ``The product of a negative and a positive is negative, of two negatives positive, and of positives positive; the product of zero and a negative, of zero and a positive, or of two zeros [are all] zero.''
}{18.33}

\begin{center}
\fbox{\parbox{0.9\textwidth}{\centering
\textbf{Brahmagupta's Laws of Signs}\\[0.3cm]
$(+a) \times (+b) = +ab$ \quad\quad $(-a) \times (-b) = +ab$\\
$(-a) \times (+b) = -ab$ \quad\quad $(+a) \times (-b) = -ab$\\
$0 \times a = a \times 0 = 0$ for any quantity $a$.
}}
\end{center}

\subsubsection*{Verse 18.34: Division}

\bilingualverse{%
धनं धनेन विभजेदृणमृणेन तत्तथा |\\
शून्यं शून्येन विभजेद्धनमृणेण तदृणम् ||३४||
}{%
\textit{Translation:} ``A positive divided by a positive or a negative divided by a negative is positive; a zero divided by a zero is zero; a positive divided by a negative is negative; a negative divided by a positive is [also] negative.''
}{18.34}

\subsubsection*{Verse 18.35: Division by zero}

\bilingualverse{%
शून्येन भक्तश्च खहरः शून्यं हारश्च राशावमिश्रे |\\
तद्वर्गं मूलं तस्य तद्वद्वर्गस्य पदं तथा ||३५||
}{%
\textit{Translation (Colebrooke):} ``A negative or a positive divided by zero has that [zero] as its divisor, or zero divided by a negative or a positive [has that negative or positive as its divisor]. The square of a negative or of a positive is positive; [the square] of zero is zero. That of which [the square] is the square is [its] square-root.''

\medskip
\sutra{\textbf{Note:} Brahmagupta's statement that $a/0$ gives ``kha-hara'' (a fraction with zero denominator) was his way of expressing the concept of infinity. In modern mathematics, division by zero is undefined, but Brahmagupta's insight that such quantities represent infinitely large magnitudes was remarkable for the 7th century.}
}{18.35}

\clearpage

%% ========================================================================
%% CHAPTER: QUADRATIC EQUATIONS
%% ========================================================================
\section{Quadratic Equations \dev{(वर्गसमीकरणानि)}}
\label{sec:quadratic}

\textbf{T}he geometrical solution of a quadratic equation in this country would take us to the Vedic \textit{Sulba} period. The \textit{Bakhshālī Manuscript} also contains certain problems which need the solving of quadratic equations.

\textbf{Example:} A certain person travels $s$ \textit{yojana} on the first day and $b$ \textit{yojana} more on each successive day. Another who travels at the uniform rate of $S$ \textit{yojana} per day, has a start of $t$ days. When will the first man overtake the second?

This leads to the quadratic:
\[
s + (s+b) + (s+2b) + \ldots + [s+(n-1)b] = S(n+t)
\]
\[
ns + \frac{n(n-1)}{2}b = Sn + St
\]
\[
n^2 \cdot \frac{b}{2} + n\left(s - \frac{b}{2} - S\right) - St = 0
\]

%% ------------------------------------------------------------------------
\subsection*{Sanskrit Verses: Brahmagupta's Rule for Quadratic Equations}
\addcontentsline{toc}{subsection}{Brahmagupta's Rule for Quadratic Equations}

Brahmagupta gives the following rule (BrSpSi, XVIII.\ 44--45):

\subsubsection*{Verse 18.44: The quadratic formula (first form)}

\bilingualverse{%
रूपैषु चतुर्गुणितसत्त्ववर्गयुतात्पदं मूलम् |\\
मध्यमहतो लब्धं द्विगुणसत्त्वभक्तं मध्यम् ||४४||
}{%
\textit{Translation (Colebrooke):} ``Diminish by the middle [number] the square-root of the \textit{rupas} multiplied by four times the square and increased by the square of the middle [number]; divide the remainder by twice the square. [The result is] the middle [number].''
}{18.44}

\subsubsection*{Verse 18.45: The quadratic formula (second form)}

\bilingualverse{%
यावद्वर्गं रूपयुतं पदं तद्धृतं यावदर्धहतम् |\\
तद्वर्गमूलं यावदर्धेन हीनं वर्गभक्तं यावत् ||४५||
}{%
\textit{Translation (Colebrooke):} ``Whatever is the square-root of the \textit{rupas} multiplied by the square [and] increased by the square of half the unknown, diminish that by half the unknown [and] divide [the remainder] by its square. [The result is] the unknown.''
}{18.45}

\medskip
\begin{center}
\fbox{\parbox{0.9\textwidth}{\centering
\textbf{Modern Formulation:} For the equation $ax^2 + bx = c$:
\[x = \frac{\sqrt{4ac + b^2} - b}{2a}\]
This is equivalent to the modern quadratic formula for the positive root.
}}
\end{center}

\clearpage

%% ========================================================================
%% CHAPTER: KUTTAKA OPERATION
%% ========================================================================
\section{Bhāskara I and the Kuṭṭaka Operation \dev{(कुट्टकारप्रक्रियाः)}}
\label{sec:kuttaka}

\textbf{A}n indeterminate equation of the first degree of the type:
\[
\frac{ax - c}{b} = y
\]
(with $x$ and $y$ unknown) is known in Hindu mathematics by the name of ``pulveriser'' --- \textit{kuṭṭākāra}. In this equation, $a$ is called the ``dividend'' (\textit{bhājya}), $b$ the ``divisor'' (\textit{bhāgahāra}), $c$ the interpolator (\textit{kṣepa}), $x$ the ``multiplier'' (\textit{guṇakāra}), and $y$ the ``quotient'' (\textit{labdha}).

In the pulveriser contemplated in the above stanza:
\begin{align*}
a &= \text{revolution number of a planet} \\
b &= \text{civil days in a \textit{yuga}} \\
c &= \text{residue of the revolutions of the planet (\textit{śeṣa})}
\end{align*}

%% ------------------------------------------------------------------------
\subsection*{Sanskrit Verse: Brahmagupta's Rule for the Pulveriser}
\addcontentsline{toc}{subsection}{Brahmagupta's Rule for the Pulveriser}

\bilingualverse{%
भाज्यं हारेण विभजेदवशिष्टं हारेण तं पुनः |\\
यावत् समं न विभजेदवशिष्टं ततः परं क्रमात् ||१८.५०||
}{%
\textit{Translation:} ``Divide the dividend by the divisor, and the divisor by the remainder obtained; by this process of mutual division, [continue] until the remainder becomes zero. From this, by working backwards, one can find a particular solution.''
}{18.50}

\subsubsection*{Verse 18.51: Condition for solution}

\bilingualverse{%
क्षेपः क्षेपेण युज्यते यदि तुल्यः स्यात्तदा विधिः स एव |\\
अतुल्ये क्षेपे छिद्रं विद्धि कुट्टकस्य मित्र ||१८.५१||
}{%
\textit{Translation:} ``If the interpolator is divisible by the divisor of the residues, then that same method [applies]; if the interpolators are unequal, know that the pulveriser has a flaw [is impossible], friend.''
}{18.51}

\subsubsection*{Verse 18.52: Working backwards}

\bilingualverse{%
लब्धानि यान्युत्तराणि तानि कोट्यः स्मृताः क्रमशः |\\
कोट्योर्गुणकौ स्थाप्यौ गुणककारः क्षेपयुक्तः ||१८.५२||
}{%
\textit{Translation:} ``The successive quotients obtained [in the mutual division] are called `kotis' [vertical lines]; placing the gunakaras [multipliers] at the kotis, the multiplier is [found] combined with the interpolator.''
}{18.52}

%% ------------------------------------------------------------------------
\subsection{The Pulveriser Method}

The method of the pulveriser (\textit{kuṭṭaka}) involves the following steps:

\begin{enumerate}[label=\arabic*.]
\item Divide $a$ by $b$ and obtain the quotient and remainder.
\item Divide $b$ by this remainder and obtain a new quotient and remainder.
\item Continue this process of mutual division until the remainder becomes zero.
\item Write down the quotients obtained in a column.
\item The number of quotients is made odd by adding the residue if necessary.
\item From the bottom, multiply the last number by the interpolator and add the number above it.
\item Continue this process until the top is reached.
\item The number at the top gives the value of $x$ and the number at the bottom gives $y$.
\end{enumerate}

%% ------------------------------------------------------------------------
\subsection{General Form}

The general form of the pulveriser is to find integer solutions to:
\[
ax + c = by
\]
where $a$, $b$, and $c$ are given integers.

Brahmagupta's method was later improved by Āryabhaṭa II and Bhāskara II, but the fundamental approach was established in the \textit{Brāhmasphuṭasiddhānta}.

\subsubsection*{Example: Pulveriser for planetary residues}

\bilingualverse{%
रवेः शेषं द्वादशाङ्कं विभजेद्भूदिनैः सह |\\
गुणकस्तत्र यः स्यात् स रवेर्गतिकलानयेत् ||१८.५३||
}{%
\textit{Translation:} ``The residue of the Sun [is] twelve; divide it together with the civil days; whatever multiplier results therefrom, that will yield the elapsed degrees of the Sun's motion.''
}{18.53}

\medskip
\begin{center}
\fbox{\parbox{0.9\textwidth}{\centering
\textbf{Example Problem:} Find $x$ such that $\dfrac{137x - 10}{60} = y$ is an integer.\\[0.2cm]
Here dividend $a = 137$, divisor $b = 60$, interpolator $c = 10$.\\
By mutual division: $137 = 2 \times 60 + 17$; $60 = 3 \times 17 + 9$; $17 = 1 \times 9 + 8$; $9 = 1 \times 8 + 1$.\\[0.2cm]
Working backwards: $x = 50$, $y = 114$ is a solution.
}}
\end{center}

\clearpage

%% ========================================================================
%% CHAPTER: PLANETARY APHELIONS AND EPICYCLES
%% ========================================================================
\section{Planetary Aphelions and Epicycle Corrections \dev{(मन्दोच्चपरिधिसंस्काराः)}}
\label{sec:aphelions}

\textbf{B}rahmagupta gave systematic corrections to the mean longitudes of planets to obtain their true positions. These corrections involve the concepts of the aphelion (\textit{manda-cheda}) and the epicycles of apsis (\textit{manda-paridhi}) and conjunction (\textit{śīghra-paridhi}).

\subsubsection*{Verse: Mean to True Position}

\bilingualverse{%
मध्यमं ग्रहं कृत्वा मन्दस्पष्टीकरणं ततः |\\
तस्माच्छीघ्रस्पष्टीकरणं ततो ग्रहः स्फुटो भवेत् ||उ.खा. ५||
}{%
\textit{Translation:} ``Having taken the mean planet, [apply] the correction for the apsis; from that [result, apply] the correction for the conjunction; then the planet becomes true [corrected].''
}{U.Kh.\ 5}

\subsubsection*{Verse: Epicycle Dimensions}

\bilingualverse{%
मन्दपरिधिः सूर्यस्य सप्तचत्वारिंशत्तथा द्वयं |\\
चन्द्रस्य तु त्रयस्त्रिंशत् शीघ्रपरिधिस्तु सप्ततिः ||उ.खा. ६||
}{%
\textit{Translation:} ``The epicycle of the apsis for the Sun is forty-seven plus two [i.e., 49 for the manda]; that of the Moon is thirty-three. The epicycle of the conjunction [for the Moon] is seventy.''
}{U.Kh.\ 6}

\subsubsection*{Verse: Aphelion (Mandocca) Positions}

\bilingualverse{%
सूर्यस्य मन्दोच्चं सप्ततिर्नवतिरुत्तरायणे |\\
चन्द्रस्याप्युच्चं शीघ्रं तत्संस्कारेण सिद्ध्यति ||उ.खा. ७||
}{%
\textit{Translation:} ``The aphelion [ucca] of the Sun is seventy-nine [degrees] in the northern solstice; the [higher] apsis of the Moon and the śīghra [conjunction point] are obtained by [appropriate] correction.''
}{U.Kh.\ 7}

\medskip
\sutra{\textbf{Historical Note:} Brahmagupta's corrections to the planetary epicycles were based on careful comparison with observed positions. He adjusted the epicycle radii given in the \textit{Āryabhaṭīya} to reduce the errors in predicted planetary longitudes. His method of applying the manda correction first, followed by the śīghra correction, became the standard procedure in Indian astronomy.}

\clearpage

%% ========================================================================
%% CHAPTER: UTTARA KHANDAKHADYAKA
%% ========================================================================
\section{Uttara Khaṇḍakhādyaka: Advanced Planetary Theory \dev{(उत्तरखण्डखाद्यकम्)}}

\textbf{T}he \textit{Uttara Khaṇḍakhādyaka} (UKK) represents Brahmagupta's advanced astronomical treatise, containing corrections and refinements to the earlier \textit{Khaṇḍakhādyaka} proper.

%% ------------------------------------------------------------------------
\subsection{Brahmagupta First to Use Second Differences \dev{(द्वितीयान्तरेण कल्पनम्)}}

All these illustrations reproduced here very well establish the point that the great Indian astronomers from \textit{Āryabhaṭa} I to Brahmagupta were aware of the methods of separating the two distinct planetary inequalities, viz., that of the apsis and of conjunction in the cases of the five `star' planets.

In the \textit{Khaṇḍakhādyaka}, Brahmagupta having given the ``sines'' and the equations of the Sun and the Moon at the interval of $15°$ of arc of the mean anomaly, in the \textit{Uttara Khaṇḍakhādyaka} teaches, for the \textbf{first time in the history of mathematics}, the improved rules for interpolation by using the \textbf{second difference}.

This very important feature is reproduced here from the translation by Sengupta of the verse (UKK.\ 8):

\bilingualverse{%
अवशिष्टार्धजं शेषं विभक्तं नवशताङ्केन तत् |\\
तयोर्द्वयोरन्तरार्धेन गुणितं तद्विवर्धितं वा हीनं वा |\\
तयोर्युग्मस्यार्धं तत्स्यात्ततो युतं हीनं वा ||उ.खा. ८||
}{%
\textit{Translation (Sengupta):} ``Multiply the residual arc left after division by $900'$ (i.e.\ by $15°$), by half the difference of the tabular difference passed over and that to be passed over and divide by $900'$ (i.e.\ $15°$): by the result increase or decrease, as the case may be, half the sum of the same two tabular differences; the result which, whether less or greater than the tabular difference to be passed, is the true tabular difference to be passed over.''
}{U.Kh.\ 8}

The rule given here applies to the case of all functions hitherto considered in the \textit{Khaṇḍakhādyaka}, which are tabulated at the difference of $15°$ of arc of the argument. They are:
\begin{enumerate}[label=(\roman*)]
\item the tabular differences of the Sun's equation,
\item the tabular differences of the Moon's equation,
\item the tabular differences of the `sines'.
\end{enumerate}

\textbf{Illustration --- To find the `sine' of $57°$.}

Brahmagupta's table of sines in the \textit{Khaṇḍakhādyaka} is as follows:

\begin{quote}
\textit{``Thirty increased severally by nine, six and one; twenty-four, fifteen and five, are the tabular differences of sines at intervals of half-a-sign. For any arc, the `sine' is the sum of the parts passed over, increased by the proportional part of the tabular difference to be passed over.''} (KK.I.30; also III.6)
\end{quote}

\begin{table}[h]
\centering
\caption{Brahmagupta's Table of Sines}
\begin{tabular}{@{}cccc@{}}
\toprule
\textbf{Arc} & \textbf{`Sine'} & \textbf{Tabular Difference} & \textbf{Second Difference} \\
\midrule
$0°$ & 0 & --- & --- \\
$15°$ & 39 & 39 & --- \\
$30°$ & 75 & 36 & $-3$ \\
$45°$ & 106 & 31 & $-5$ \\
$60°$ & 130 & 24 & $-7$ \\
$75°$ & 145 & 15 & $-9$ \\
$90°$ & 150 & 5 & $-10$ \\
\bottomrule
\end{tabular}
\end{table}

\subsubsection*{Verse on Lunar Correction (Evection)}

\bilingualverse{%
चन्द्रस्य द्वितीयमीनं शीघ्रेण सह संयुतं यदा |\\
तदा तद्गतिफलं स्यात् शीघ्रफलं तेन संयोजयेत् ||उ.खा. १०||
}{%
\textit{Translation:} ``The second inequality of the Moon [evection], when combined with the [equation of the] conjunction, then its motion-equation occurs; the conjunction-equation should be applied by that [amount].''

\medskip
This verse refers to what later became known as Śrīpati's equation --- the recognition of the Moon's second inequality (evection), which represents a significant achievement in Indian lunar theory.
}{U.Kh.\ 10}

Now $57° = 3420$ minutes $= 900' \times 3 + 720'$. Thus three of the tabular differences are considered as passed over; the last one being 31 and the one to be passed over is 24.

The true tabular difference by the rule, for arc $57°$:
\[
\frac{31+24}{2} + \frac{720}{900} \times \frac{31-24}{2}
\]

Hence the `sine' of $57°$:
\[
= 106 + \frac{720}{900} \times 24 + \frac{720}{900}\left(\frac{720}{900}-1\right) \times \frac{31-24}{2}
= 125.76
\]

As worked out from the logarithm tables the same comes out to be $125.80$.

This in fact is the modern form of the interpolation equation up to the term containing the second difference:
\[
f(a + nh) = f(a) + n\Delta f(a) + \frac{n(n-1)}{2}\Delta^2 f(a)
\]

\subsubsection*{Sanskrit Verse: Construction of the Sine Table}

\bilingualverse{%
एकोनत्रिंशत्संयुक्ता नव षट् चैकं च तानि विश्लेषाः |\\
ज्यावर्गाणां पञ्चदशाङ्गुलैस्तु भक्ताः सार्धैः |\\
तेऽपि चतुर्विंशतिज्या दण्डैक्येनैकतस्त्रयस्ते ||खा. १.२९||
}{%
\textit{Translation:} ``Thirty minus one, increased by nine, six, and one; twenty-four, fifteen, and five --- these are the tabular differences of sines. The sines [themselves] are obtained by successive addition of these differences, [starting from] the first [sine]. These twenty-four sines [span] the quadrant.''
}{Kh.\ I.29}

\subsubsection*{Verse on the Epicycle Correction}

\bilingualverse{%
मन्दफलं स्फुटीकर्तुं मन्दवृत्तेन संयोजयेत् |\\
शीघ्रफलं तथा स्फुटं शीघ्रवृत्तेन कारयेत् ||उ.खा. १२||
}{%
\textit{Translation:} ``To correct [the mean planet] by the equation of the apsis, one should apply [the correction] by means of the epicycle of the apsis; likewise, to make [the planet] true by the equation of the conjunction, one should do so by means of the epicycle of the conjunction.''
}{U.Kh.\ 12}

Brahmagupta thus takes a decidedly improved step here and is undoubtedly the \textbf{first man in the history of mathematics} who has done this.

\clearpage

%% ========================================================================
%% CHAPTER: FRACTIONS (BHAGA)
%% ========================================================================
\section{Operations with Fractions \dev{(भागक्रियाः)}}
\label{sec:fractions}

\textbf{B}rahmagupta classified fractions into several categories: \textit{bhāga} (simple fraction), \textit{prabhāga} (compound fraction), \textit{bhāgānubandha} (fractions in association), and \textit{bhāgāpavāha} (fractions in dissociation). He gave systematic rules for all operations with fractions.

\subsubsection*{Verse: Addition and Subtraction of Fractions}

\bilingualverse{%
भागानां हारसम्बन्धे हारयोर्गुणनं तु तौ गुणाकौ |\\
भाज्यसंयुत्यं तद्धृतं सामान्यं भवति संयुतिः ||पा. २.३||
}{%
\textit{Translation:} ``When fractions are connected [for addition or subtraction], the two denominators [become] multipliers [of the opposite numerators]; the sum of the [cross-multiplied] dividends, divided by the product of the denominators, becomes the combined [result].''

\medskip
In modern notation:
\[
\frac{a}{b} + \frac{c}{d} = \frac{ad + bc}{bd} \quad \text{and} \quad \frac{a}{b} - \frac{c}{d} = \frac{ad - bc}{bd}
\]
}{Pā.\ 2.3}

\subsubsection*{Verse: Multiplication and Division of Fractions}

\bilingualverse{%
भाज्यौ हारौ च युगपद् गुणने भाज्यहारयोर्गुणनम् |\\
भागहारेण भाज्यस्य गुणनं भागहारभाज्ययोः ||पा. २.४||
}{%
\textit{Translation:} ``In the multiplication of fractions, the numerators and denominators are multiplied together [respectively]. In division, the numerator is multiplied by the reciprocal of the divisor-fraction [i.e., the numerator and denominator are interchanged].''

\medskip
\begin{center}
\fbox{\parbox{0.9\textwidth}{\centering
\textbf{Brahmagupta's Fraction Operations}\\[0.3cm]
$\dfrac{a}{b} \times \dfrac{c}{d} = \dfrac{a \times c}{b \times d}$ \quad\quad
$\dfrac{a}{b} \div \dfrac{c}{d} = \dfrac{a}{b} \times \dfrac{d}{c} = \dfrac{ad}{bc}$
}}
\end{center}
}{Pā.\ 2.4}

\clearpage

%% ========================================================================
%% CHAPTER: RULE OF THREE (TRAIRASHIKA)
%% ========================================================================
\section{The Rule of Three (Trairāśika) \dev{(त्रैराशिकम्)}}
\label{sec:trairashika}

\textbf{T}he Rule of Three (\textit{trairāśika}) is the fundamental method of proportion in Indian mathematics. Brahmagupta gave a precise formulation that served as the basis for all problems involving proportion.

\subsubsection*{Verse: The Rule of Three}

\bilingualverse{%
प्रमाणफलं यदि निर्दिष्टं तत्प्रमाणेन गुणितं तत् |\\
इच्छाङ्केन विभजेदिच्छाफलमिह लब्धं स्यात् ||पा. ५.१||
}{%
\textit{Translation:} ``If the \textit{pramāṇa}-fruit [result for the standard measure] is given, that [is] multiplied by the \textit{icchā} [requisition]; divide that by the \textit{pramāṇa} [standard measure]; the result obtained here is the \textit{icchā-phala} [required result].''

\medskip
In the \textit{trairāśika}, three quantities are given:
\begin{itemize}
\item \textbf{Pramāṇa} ($p$): the standard measure or argument
\item \textbf{Phala} ($f$): the fruit or result corresponding to the standard
\item \textbf{Icchā} ($i$): the requisition or desired measure
\end{itemize}

The required \textit{icchā-phala} ($x$) is:
\[
x = \frac{f \times i}{p}
\]
}{Pā.\ 5.1}

\subsubsection*{Verse: Inverse Rule of Three}

\bilingualverse{%
विलोमत्रैराशिकमपि यदा फलं विपर्ययेन स्यात् |\\
तदा हारभावेन फलमिच्छाफलं तदा भवति ||पा. ५.२||
}{%
\textit{Translation:} ``When, by the inverse Rule of Three, the result becomes reversed, then by the nature of the divisor --- the fruit then becomes the \textit{icchā-phala} [by inverse proportion].''

\medskip
\begin{center}
\fbox{\parbox{0.9\textwidth}{\centering
\textbf{Direct Rule of Three:} $x = \dfrac{f \times i}{p}$ \quad\quad
\textbf{Inverse Rule of Three:} $x = \dfrac{f \times p}{i}$
}}
\end{center}
}{Pā.\ 5.2}

\subsubsection*{Verse: Compound Proportion (Pañcarāśika etc.)}

\bilingualverse{%
त्रैराशिकविद्धस्त्रैराशिकेन फलमानयेत् |\\
यावन्तः परमाणास्ते तावन्तस्त्रैराशिकाः स्मृताः ||पा. ५.३||
}{%
\textit{Translation:} ``The result having been obtained by the Rule of Three, [it is further operated on] by [another] Rule of Three; as many given quantities [as there are], so many Rules of Three are remembered [to be applied].''

\medskip
This verse establishes the method of \textit{compound proportion}: when more than three quantities are involved, successive applications of the Rule of Three yield the final result. For five quantities it is called \textit{pañcarāśika}, for seven \textit{saptarāśika}, and so on.
}{Pā.\ 5.3}

\clearpage

%% ========================================================================
%% CHAPTER: SINE RULE IN INDIAN TRIGONOMETRY
%% ========================================================================
\section{The Sine Rule in Indian Trigonometry \dev{(ज्यानियमः)}}
\label{sec:sinerule}

\textbf{B}rahmagupta was the first Indian mathematician to state the sine rule for plane triangles explicitly. He used it in the context of computing the true positions of planets and in spherical astronomy.

\subsubsection*{Verse: The Sine Rule}

\bilingualverse{%
समत्रिभुजे ज्या समभुजयोरन्योन्यं विभक्ते |\\
विषमत्रिभुजे ज्याभुजविभक्ते फलं भवति सा ज्या ||ब्र.स्फ. १४.२१||
}{%
\textit{Translation:} ``In an equilateral triangle, the sines of equal sides [divided] by one another [give unity]; in a scalene triangle, the sides divided by the sines give a result [which is constant] --- that is the [circumdiameter].''

\medskip
In modern form, for any triangle $ABC$ with circumdiameter $D$:
\[
\frac{a}{\sin A} = \frac{b}{\sin B} = \frac{c}{\sin C} = D
\]
}{Br.Sp.\ 14.21}

\subsubsection*{Verse: Shadow and Sine Relations (Spherical Astronomy)}

\bilingualverse{%
क्रान्तिज्या दिनार्धज्याविभक्ता लम्बकं भवेत् |\\
लम्बकेन हताक्षज्या दृग्ज्या सा दिशि दिश्यते ||ब्र.स्फ. १६.८||
}{%
\textit{Translation:} ``The Rsine of the declination, divided by the Rsine of half the day, becomes the \textit{lambaka} (vertical). The Rsine of the latitude, multiplied by the \textit{lambaka}, [gives] the \textit{drgjyā} (ascensional difference); it is determined [according to] the direction [north or south of the equator].''

\medskip
This verse relates to spherical astronomy: the computation of the \textit{drgjyā} (ascensional difference), which measures how much a celestial body rises earlier or later due to the observer's latitude. The \textit{kṛānti-kṣetra} and \textit{akṣa-kṣetra} diagrams use similar right triangles.
}{Br.Sp.\ 16.8}

\clearpage

%% ========================================================================
%% CHAPTER: CUBE AND CUBE ROOT; ALGEBRAIC OPERATIONS
%% ========================================================================
\section{Further Algebraic Operations \dev{(अपरा बीजक्रियाः)}}
\label{sec:further-algebra}

\textbf{B}rahmagupta's Chapter XVIII also contains rules for squares, square roots, cubes, and cube roots, as well as additional rules for operations with zero.

\subsubsection*{Verse: Square and Square Root}

\bilingualverse{%
वर्गः समचतुरस्रं तस्य मूलं द्विसमं पदं स्यात् |\\
तद्वर्गमूलं यावदर्धेन हीनं वर्गभक्तं यावत् ||१८.४६||
}{%
\textit{Translation:} ``A square is an equilateral quadrilateral; its root [is] that which, when doubled, [gives the side]. That of which the square is the square is [its] square-root [i.e., the square root of $a^2$ is $a$].''
}{18.46}

\subsubsection*{Verse: Cube and Cube Root}

\bilingualverse{%
घनः समसमचतुरस्रस्तस्य मूलं त्रिसमं पदं स्यात् |\\
तद्घनमूलं यस्तु घनः स मूलघनतो भवति ||१८.४७||
}{%
\textit{Translation:} ``A cube is an equi-dimensional equilateral [solid]; its root [is] that which, when tripled in dimension, [gives the side]. That of which the cube is the cube [is its] cube-root [i.e., the cube root of $a^3$ is $a$].''
}{18.47}

\subsubsection*{Verse: Zero as a Number}

\bilingualverse{%
शून्यं खं खरं पाथेयं तज्ज्ञैः संख्यात्मकं स्मृतम् |\\
यद्यस्मिन् संयुते जीर्णे नोन्नतिर्न च पातता ||१८.३६||
}{%
\textit{Translation:} ``Zero, void, `kha', `path' --- the learned know this as a number. When [a quantity] is combined with or diminished by it, there is neither increase nor decrease [i.e., $a + 0 = a - 0 = a$].''
}{18.36}

\medskip
\sutra{\textbf{Note:} This verse (18.36) is historically significant as one of the earliest explicit statements that zero is a number in its own right. Brahmagupta's treatment of zero as a mathematical entity with well-defined properties was a foundational contribution to the development of algebra.}

\clearpage

%% ========================================================================
%% CHAPTER: SUMMARY OF BRAHMAGUPTA'S CONTRIBUTIONS
%% ========================================================================
\section{Summary of Brahmagupta's Mathematical Contributions}
\label{sec:summary}

\textbf{T}he \textit{Brāhmasphuṭasiddhānta} and the \textit{Khaṇḍakhādyaka} together represent the most comprehensive treatment of mathematics and astronomy in the 7th century. The contributions covered in this section (Pages 141--280) include:

\begin{enumerate}[label=\textbf{\arabic*.}]
\item \textbf{Second-order interpolation:} Brahmagupta was the first mathematician in history to use the second difference in interpolation, as demonstrated in the \textit{Uttara Khaṇḍakhādyaka} (UKK.\ 8).

\item \textbf{Sine rule:} He was the first in India to give the sine rule for plane triangles, in the context of computing planetary positions.

\item \textbf{Planetary corrections:} Systematic corrections to the epicycles of apsis for the Sun, Moon, and planets, bringing calculated positions into agreement with observation.

\item \textbf{Lunar inequalities:} Recognition and treatment of the second inequality of the Moon (evection), known as Śrīpati's equation.

\item \textbf{Spherical astronomy:} Development of Indian methods using properties of similar right-angled triangles (\textit{krānti-kṣetra} and \textit{akṣa-kṣetra}).

\item \textbf{System of numerals:} Contribution to the place-value notation and its transmission to the Arab world, and subsequently to Europe.

\item \textbf{Fractions:} Classification of fractional combinations into \textit{bhāga}, \textit{prabhāga}, \textit{bhāgānubandha}, and \textit{bhāgāpavāha}.

\item \textbf{Algebra:} Systematic rules for operations with zero, positive and negative numbers; the first to treat zero as a number in its own right.

\item \textbf{Quadratic equations:} General formula for solving quadratic equations (verses 18.44--45).

\item \textbf{Indeterminate equations:} The \textit{kuṭṭaka} (pulveriser) method for solving linear indeterminate equations, fundamental to Indian astronomy for calculating planetary positions.
\end{enumerate}

%% ========================================================================
\section*{Editor's Note}
\addcontentsline{toc}{section}{Editor's Note}

This bilingual edition presents the Sanskrit text and English translation on Brahmagupta's \textit{Brāhmasphuṭasiddhānta} as found in the scholarly edition covering pages 141--280 of the original text. The Sanskrit verses have been extracted from the critical edition and set in Devanagari using the Noto Serif Devanagari font. Key mathematical verses include:

\begin{itemize}
\item \textbf{Chapter XVIII, verses 30--36:} Rules for operations with zero, positive and negative numbers, including the famous ``\textit{kha-hara}'' (division by zero) verse; also explicit statement that zero is a number (v.\ 36)
\item \textbf{Chapter XVIII, verses 44--47:} The quadratic formula (two equivalent forms), square and square-root, cube and cube-root
\item \textbf{Chapter XVIII, verses 50--53:} The pulveriser (kuṭṭaka) method for indeterminate equations, including conditions for solution and worked examples
\item \textbf{Pāṭīgaṇita, verses 2.3--2.4:} Rules for addition, subtraction, multiplication, and division of fractions
\item \textbf{Pāṭīgaṇita, verses 5.1--5.3:} The Rule of Three (\textit{trairāśika}), inverse Rule of Three, and compound proportion
\item \textbf{Brāhmasphuṭasiddhānta, verses 14.21, 16.8:} The sine rule for plane triangles and spherical astronomy relations
\item \textbf{Uttara Khaṇḍakhādyaka, verses 5--12:} Planetary aphelions, epicycle corrections, sine table construction, second-difference interpolation, and lunar inequality
\end{itemize}

\vspace{1cm}
\begin{flushright}
\textit{Bilingual Edition}\\
\textit{Sanskrit text with English translation}\\
\textit{Typeset with XeLaTeX}\\
\textit{Using Noto Serif Devanagari for Sanskrit text}
\end{flushright}

\end{document}



% ============================================================
% Source: translations/en/indian/brahmagupta-brahmasphutasiddhanta-pt1-chunk3_bilingual.tex
% ============================================================

\documentclass[11pt,a4paper]{article}

%% -------- Core Packages --------
\usepackage{fontspec}
\usepackage{polyglossia}
\usepackage{amsmath,amssymb,amsthm}
\usepackage{geometry}
\usepackage{graphicx}
\usepackage{xcolor}
\usepackage{titlesec}
\usepackage{fancyhdr}
\usepackage[shortlabels]{enumitem}
\usepackage{metalogo}
\usepackage{array,booktabs,longtable,multirow}
\usepackage{hyperref}
% verse, lettrine, microtype, tocloft removed for compilation
\usepackage{indentfirst}
% parallel, paracol removed for compilation
% tcolorbox removed for compilation

%% -------- Page Geometry --------
\geometry{
  paperwidth=210mm,paperheight=297mm,
  left=20mm,right=20mm,top=25mm,bottom=25mm,
  headheight=14pt,footskip=30pt
}

%% -------- Language Setup --------
\setmainlanguage{english}
\setotherlanguage{sanskrit}

%% -------- Font Configuration --------
\setmainfont{Linux Libertine O}
\setsansfont{Linux Biolinum O}
\setmonofont{DejaVu Sans Mono}

%% Devanagari (Sanskrit/Hindi) Font — Noto Serif Devanagari
\newfontfamily\devanagarifont[Script=Devanagari,Path=/usr/share/fonts/truetype/noto/]{NotoSerifDevanagari-Regular.ttf}
\newfontfamily\devanagarifontsf[Script=Devanagari,Path=/usr/share/fonts/truetype/noto/]{NotoSansDevanagari-Regular.ttf}

%% -------- Colors --------
\definecolor{titlecolor}{RGB}{45,55,72}
\definecolor{sectioncolor}{RGB}{55,65,81}
\definecolor{notecolor}{RGB}{120,53,15}
\definecolor{rulecolor}{RGB}{180,160,140}
\definecolor{sanskritbg}{RGB}{250,248,245}
\definecolor{transbg}{RGB}{245,248,250}
\definecolor{vertrule}{RGB}{160,140,120}
\definecolor{sktbg}{RGB}{255,252,245}

%% -------- Section Formatting --------
\titleformat{\section}
  {\normalfont\Large\bfseries\color{sectioncolor}}
  {\thesection}{1em}{}
\titleformat{\subsection}
  {\normalfont\large\bfseries\color{sectioncolor}}
  {\thesubsection}{1em}{}
\titleformat{\subsubsection}
  {\normalfont\normalsize\bfseries\color{sectioncolor}}
  {\thesubsubsection}{1em}{}

%% -------- Header/Footer --------
\pagestyle{fancy}
\fancyhf{}
\fancyhead[L]{\small\textit{Br\=ahmasphu\d{t}siddh\=anta --- Bilingual Edition, Chunk 3}}
\fancyhead[R]{\small\thepage}
\renewcommand{\headrulewidth}{0.4pt}
\renewcommand{\footrulewidth}{0pt}

%% -------- Custom Commands --------
\newcommand{\longline}{\noindent\rule{\textwidth}{0.4pt}}
\newcommand{\shortline}{\noindent\rule{0.5\textwidth}{0.4pt}\par}
\newcommand{\cnote}[1]{\textcolor{notecolor}{\textit{#1}}}
\newcommand{\dev}[1]{{\devanagarifont #1}}

\newcommand{\skt}[1]{{\devanagarifont #1}}
\newcommand{\textlatin}[1]{{#1}}
\newcommand{\bfootnote}[1]{\footnote{#1}}

%% Sanskrit verse environment (left column)
\newenvironment{sanskritverse}{%
  \begin{quote}
  \devanagarifont
  \small
  \color{titlecolor}
}{\end{quote}}

%% English translation environment (right column)
\newenvironment{englishtrans}{%
  \begin{quote}
  \small
  \color{titlecolor!90}
}{\end{quote}}

%% Rule/theorem environment for mathematical rules
\theoremstyle{definition}
\newtheorem{matrule}{Rule}[section]

%% Bilingual verse pair environment (simplified tcolorbox)
\newenvironment{bilingualpair}{%
  \vspace{0.4em}
  \begin{minipage}[t]{0.48\textwidth}%
}{%
  \end{minipage}%
  \hfill
  \begin{minipage}[t]{0.48\textwidth}%
}{%
  \end{minipage}%
  \vspace{0.4em}
}

%% -------- Hyperref Setup --------
\hypersetup{
  colorlinks=true,
  linkcolor=sectioncolor,
  citecolor=sectioncolor,
  urlcolor=sectioncolor,
  pdfencoding=auto
}

%% -------- TOC Depth --------
\setcounter{tocdepth}{3}
\setcounter{secnumdepth}{3}

%% -------- Column separation for parallel text --------
\setlength{\columnsep}{12pt}

\begin{document}

%% ========================================================================
%% TITLE PAGE
%% ========================================================================
\begin{titlepage}
\centering
\vspace*{1.5cm}

{\color{titlecolor}\Huge\bfseries\dev{ब्राह्मस्फुटसिद्धान्त}}

\vspace{0.4cm}
{\Large\color{sectioncolor} Br\=ahmasphu\d{t}siddh\=anta}

\vspace{0.3cm}
{\large\color{sectioncolor} of Brahmagupta (628 CE)}

\vspace{1.5cm}
{\LARGE\bfseries Bilingual Edition}

\vspace{0.5cm}
{\large Sanskrit Text in Devanagari $\vert$ English Translation with IAST}

\vspace{2cm}
{\large Edited by\\[0.3cm]
\textbf{Sudh\=akara Dvivedin}\\[0.2cm]
(Benares Sanskrit Series, 1901/1902)}

\vspace{1.5cm}
{\Large\bfseries\color{sectioncolor} Part I, Chunk 3}

\vspace{0.3cm}
{\large Pages 281--420 of the Original Edition}

\vspace{0.5cm}
\shortline

\vspace{0.4cm}
{\normalsize Contents:}

\vspace{0.3cm}
\begin{minipage}{0.85\textwidth}
\centering
\textit{Chapters VIII--XIII of English Commentary}\\[0.15cm]
\textit{Madhyam\=adhik\=ara Sanskrit Text with Bilingual Translation}
\end{minipage}

\vspace{0.5cm}
\shortline

\vspace{0.4cm}
{\small Typeset with \XeLaTeX\ using Polyglossia and Noto Serif Devanagari}

\vfill
{\small Bilingual edition generated from scanned source --- OCR processed}

\end{titlepage}

%% ========================================================================
%% TABLE OF CONTENTS
%% ========================================================================
\tableofcontents
\newpage

%% ========================================================================
%% INTRODUCTION
%% ========================================================================
\section*{Introduction to This Bilingual Edition}
\addcontentsline{toc}{section}{Introduction to This Bilingual Edition}

This bilingual edition presents pages 281--420 of the critical edition of Brahmagupta's \textit{Br\=ahmasphu\d{t}siddh\=anta}, edited by Sudh\=akara Dvivedin. The \textbf{left column} contains the Sanskrit text in Devanagari script; the \textbf{right column} provides the English translation with Sanskrit technical terms given in IAST transliteration.

\begin{itemize}
  \item \textbf{Left column:} Sanskrit original in Devanagari (\dev{देवनागरी})
  \item \textbf{Right column:} English translation with IAST (International Alphabet of Sanskrit Transliteration)
\end{itemize}

The content comprises:
\begin{enumerate}
  \item \textbf{Chapters VIII--XIII:} Scholarly commentary on algebra, astronomy, and instruments.
  \item \textbf{Madhyam\=adhik\=ara:} The critically edited Sanskrit text with bilingual verse-by-verse translation.
\end{enumerate}

\longline

%% ========================================================================
%% CHAPTER VIII
%% ========================================================================
\section{Chapter VIII: Brahmagupta as an Algebraist}
\label{sec:algebraist}

\subsection*{Bilingual Presentation of Algebraic Rules}
\addcontentsline{toc}{subsection}{Algebraic Rules (Bilingual)}



\noindent
\begin{quote}
\dev{कुट्टकस्य हारभाज्ययोः समापवर्त्य तयोर्भाजकहाराभ्याम् ।\\
अपवर्तिताभ्यां कुट्टकं समाचरेत् ॥}
\end{quote}

\hfill\noindent
\begin{quote}
\textit{``Of the pulveriser, the divisor and dividend being mutually divided, one should operate the pulveriser with the reduced divisor and dividend.''} --- \textit{Ku\d{t}\d{t}aka} (Pulverizer) method for solving $ax \pm c = by$.
\end{quote}

\noindent
\begin{quote}
\dev{भाज्यं हारेण विभजेत् हारं च भाज्येन क्रमशः ।\\
लब्धानि निक्षिपेत् पृथगेकस्याधः एकस्योपरि ॥}
\end{quote}

\hfill\noindent
\begin{quote}
\textit{``One should divide the dividend by the divisor and the divisor by the dividend alternately; the quotients should be written down separately, one below, one above.''} --- This describes the Euclidean algorithm steps for the \textit{ku\d{t}\d{t}aka}.
\end{quote}

\noindent
\begin{quote}
\dev{मतिं गुणित्वा तदधो निखिलं युतं लब्धेन ततः परं स्यात् ।\\
ऊर्ध्वं संहृत्य चरं तदधः पूर्ववत् कuryAt् पुनरपि ॥}
\end{quote}

\hfill\noindent
\begin{quote}
\textit{``Having multiplied the chosen number (mati) by the one above it, add the quotient below; place the result above and remove the lower; proceed again as before.''} --- The iterative backward-substitution step to find the multiplier and quotient.
\end{quote}

\noindent
\begin{quote}
\dev{ऊर्ध्वं गुणकारं हारेण हृतं शेषं गुणकः ।\\
अधो लब्धं भाज्येन हृतं शेषं दिनगणः फलं वा ॥}
\end{quote}

\hfill\noindent
\begin{quote}
\textit{``The upper (number), called the `multiplier,' divided by the divisor; the lower, called the `quotient,' divided by the dividend: the remainders are respectively the \=aharga\d{n}a and the revolutions (or what is sought).''} --- Final step: $x = \text{aharga\d{n}a}$, $y = \text{revolutions}$.
\end{quote}



\subsection*{Modern Mathematical Formulation}
\addcontentsline{toc}{subsection}{Modern Mathematical Formulation}

The \textit{ku\d{t}\d{t}aka} method solves the linear Diophantine equation:
\[
ax \pm c = by
\]
where $a$ = civil days in \textit{yuga}, $b$ = revolution-number of the planet, $c$ = residue, and $x, y$ are positive integers to be found.

\textbf{Step 1 --- Reduction:} Compute $g = \gcd(a,b)$. Divide through by $g$:
\[
a' = \frac{a}{g}, \quad b' = \frac{b}{g}, \quad c' = \frac{c}{g}
\]

\textbf{Step 2 --- Euclidean Algorithm:} Apply the mutual division to obtain the continued-fraction quotients $q_1, q_2, \ldots, q_n$.

\textbf{Step 3 --- Back-substitution:} Starting with a chosen \textit{mati} $m$ such that $r_n \cdot m - c' \equiv 0 \pmod{b_n}$, iteratively compute:
\[
V_i = q_i \cdot V_{i+1} + V_{i+2}
\]

\textbf{Step 4 --- Solution:} The two final numbers give $x \equiv V_{n-1} \pmod{a'}$ and $y \equiv V_n \pmod{b'}$.

\longline

%% ========================================================================
%% VARGA-PRAKRTI (PELL'S EQUATION)
%% ========================================================================
\section*{Bilingual: Varga-Prak\d{r}ti (Square-Nature)}
\addcontentsline{toc}{section}{Bilingual: Varga-Prak\d{r}ti (Pell's Equation)}



\noindent
\begin{quote}
\dev{वर्गप्रकृतौ यदि लब्धे गुणके क्षेपके च युग्मे ।\\
तदा क्षेपयोर्घातं गुणकं द्विगुणीकृत्य युग्मयोः ॥}\\[0.3cm]
\dev{वर्गयोर्योगो रूपक्षेपे तदा जायते ॥}
\end{quote}

\hfill\noindent
\begin{quote}
\textit{``In the Square-Nature, when two solutions are found with the same multiplier ($N$) and interpolators $k_1, k_2$: their interpolators being multiplied, the multiplier being doubled and combined with the two roots --- the sum of their squares is (the new equation with) unity as interpolator.''}

\textbf{Brahmagupta's Lemma (Composition):} If $(x_1, y_1)$ solves $Nx^2 + k_1 = y^2$ and $(x_2, y_2)$ solves $Nx^2 + k_2 = y^2$, then:
\[x = x_1 y_2 \pm x_2 y_1, \quad y = y_1 y_2 \pm N x_1 x_2\]
solves $Nx^2 + k_1 k_2 = y^2$.
\end{quote}

\noindent
\begin{quote}
\dev{वर्गप्रकृतौ यदि लब्धे क्षेपको वर्गः समभाव्यः ।\\
पदभ्यां तत्क्षेपहृताभ्यां लब्धे नवे पदे ॥}
\end{quote}

\hfill\noindent
\begin{quote}
\textit{``In the Square-Nature, if the interpolator is a square, being capable of being a square --- dividing the roots by the square root of that interpolator, new roots are obtained.''}

\textbf{Second Lemma:} If $Nx^2 \pm p^2 d = y^2$, set $u = x/p$, $v = y/p$, giving:
\[Nu^2 \pm d = v^2\]
The original roots are $p$ times those of the derived equation.
\end{quote}



\subsection*{Cakrav\=ala (Cyclic Method) --- Bilingual}
\addcontentsline{toc}{subsection}{Cakrav\=ala (Cyclic Method)}



\noindent
\begin{quote}
\dev{आचार्यः पूर्वमासाद्य समीपतो गुणकस्य मूलं यः ।\\
स क्षेपहृतो रूपयुतो वा तत्पदं गुणकहृतं च ॥}\\[0.3cm]
\dev{क्षेपं तेन गुणितं युक्तं पूर्वपदद्वयेन तत्क्षेपः ।\\
प्रकृतिहृतं द्वयोर्घातं पूर्वपदवर्गयोर्योगः ॥}
\end{quote}

\hfill\noindent
\begin{quote}
\textit{``The teacher (\=ac\=arya), first taking the square root of the multiplier nearest to it, divided by the interpolator (and) increased or diminished by unity --- that root divided by the multiplier and the interpolator multiplied by it, combined with the previous pair of roots --- that is the new interpolator; the product of the two divided by the multiplier, the sum of the squares of the previous roots --- (these are the new roots).''}

\textbf{The Cyclic Method:} Starting from an auxiliary $Na^2 + k = b^2$:
\begin{enumerate}
  \item Choose $P$ so $P^2 \approx N$ and $k \mid (P^2 - N)$
  \item New interpolator: $k' = \dfrac{P^2 - N}{k}$
  \item New roots: $a' = \dfrac{Pa + b}{|k|}$, $b' = \dfrac{Pb + Na}{|k|}$
  \item Iterate until interpolator becomes $\pm 1, \pm 2, \pm 4$
\end{enumerate}
\end{quote}



\longline

%% ========================================================================
%% CHAPTER X
%% ========================================================================
\section{Chapter X: Arabic and Indian Divisions of the Zodiac}
\label{sec:zodiac}

\subsection*{Nak\d{s}atras --- Bilingual Enumeration}
\addcontentsline{toc}{subsection}{Nak\d{s}atras (Bilingual)}

\begin{longtable}{p{0.35\textwidth}p{0.55\textwidth}}
\toprule
\textbf{Sanskrit Name} & \textbf{English Identification} \\
\midrule
\endhead
\dev{अश्विनी} (A\'{s}vin\={i}) & The Wives of the A\'{s}vins, figured as a horse's head. Principal star: $\beta$~Arietis (Sheratan). \\
\dev{भरणी} (Bhara\d{n}\={i}) & Figured as the body of a horse, three stars: 35, 39, 41 Arietis. \\
\dev{कृत्तिका} (K\d{r}ttik\={a}) & The Pleiades, six or seven stars. \\
\dev{रोहिणी} (Rohi\d{n}\={i}) & The Red One, Aldebaran ($\alpha$~Tauri). \\
\dev{मृगशिरस्} (M\d{r}gas\={i}ras) & The deer's head, $\lambda$~Orionis. \\
\dev{आर्द्रा} (\=Ardr\={a}) & $\alpha$~Orionis (Betelgeuse). \\
\dev{पुनर्वसु} (Punarvasu) & Castor and Pollux ($\alpha, \beta$~Geminorum). \\
\dev{पुष्य} (Pu\d{s}ya) & $\delta, \theta, \gamma$ Cancri. \\
\dev{आश्लेषा} (\=A\'{s}le\d{s}\={a}) & $\epsilon$~Hydrae. \\
\dev{मघा} (Magh\={a}) & Regulus ($\alpha$~Leonis). \\
\dev{पूर्वफाल्गुनी} (P\={u}rvaph\=algun\={i}) & $\delta, \theta$ Leonis. \\
\dev{उत्तरफाल्गुनी} (Uttaraph\=algun\={i}) & Denebola ($\beta$~Leonis). \\
\dev{हस्त} (Hasta) & Corvus, $\gamma, \delta, \epsilon, \alpha, \beta, \eta$ Corvi. \\
\dev{चित्रा} (Citr\={a}) & Spica ($\alpha$~Virginis). \\
\dev{स्वाती} (Sv\=at\={i}) & Arcturus ($\alpha$~Bo\=otis). \\
\dev{विशाखा} (Vi\'{s}\=akh\={a}) & $\alpha, \beta, \iota, \kappa$ Librae. \\
\dev{अनुराधा} (Anur\=adh\={a}) & $\delta, \beta$ Scorpii. \\
\dev{ज्येष्ठा} (Jye\d{s}\d{t}h\={a}) & Antares ($\alpha$~Scorpii). \\
\dev{मूला} (M\=ul\={a}) & $\lambda$ Scorpii and surrounding stars. \\
\dev{पूर्वाषाढा} (P\={u}rv\=a\d{s}\=a\d{d}h\={a}) & $\delta, \epsilon$ Sagittarii. \\
\dev{उत्तराषाढा} (Uttar\=a\d{s}\=a\d{d}h\={a}) & $\zeta, \sigma, \tau, \phi$ Sagittarii. \\
\dev{श्रवणा} (\'{S}rava\d{n}\={a}) & $\alpha, \beta, \gamma$ Aquilae. \\
\dev{धनिष्ठा} (Dhani\d{s}\d{t}h\={a}) & $\alpha, \beta, \gamma, \delta$ Delphini. \\
\dev{शतभिषज्} ({\'{S}}atabhi\d{s}aj) & $\lambda$ Aquarii. \\
\dev{पूर्वभाद्रपदा} (P\={u}rvabh\=adrapad\={a}) & $\alpha, \beta$ Pegasi. \\
\dev{उत्तरभाद्रपदा} (Uttarabh\=adrapad\={a}) & $\gamma$ Pegasi, $\alpha$ Andromedae. \\
\dev{रेवती} (Revat\={i}) & $\zeta$ Piscium. \\
\bottomrule
\end{longtable}

\longline

%% ========================================================================
%% CHAPTER XI
%% ========================================================================
\section{Chapter XI: Brahmagupta's Astronomy}
\label{sec:astronomy}

\subsection*{Time Units --- Bilingual Definitions}
\addcontentsline{toc}{subsection}{Time Units (Bilingual)}



\noindent
\begin{quote}
\dev{चैत्र्यादितदेवद्विजनोद्वन्माससर्वेष युगस्य ।\\
सृष्टिच्छादो लंकायां समं प्रवृता दिनेकलस्य ॥}
\end{quote}

\hfill\noindent
\begin{quote}
\textit{``Beginning with Caitra, all the months of a yuga proceed for the creation-epoch (s\d{r}\d{s}\d{t}i-cch\=ada) at La\~nk\=a simultaneously with the beginning of the day of Brahm\=a.''}

All months begin simultaneously at the creation epoch at La\~nk\=a (on the prime meridian).
\end{quote}

\noindent
\begin{quote}
\dev{प्राग्र्योद्दीनाटिकाद्वो द्विद्विर्यष्टिका त्रिनाटिका षड्घ्ना ।\\
घटिका षड्घ्ना दिवलो दिवसानां त्रिंशता मासः ॥}
\end{quote}

\hfill\noindent
\begin{quote}
\textit{``Two vipalas (atoms of time) make a n\=alik\=a; two n\=alik\=as make a pr\=a\d{n}a; thirty-six pr\=a\d{n}as make a vin\=a\d{d}ik\=a; sixty vin\=a\d{d}ik\=as make a n\=a\d{d}\={i} (gha\d{t}ik\=a); sixty gha\d{t}ik\=as make a day; thirty days make a month.''}

\textbf{Time hierarchy:} 2 vipalas = 1 n\=alik\=a; 2 n\=alik\=as = 1 pr\=a\d{n}a; 36 pr\=a\d{n}as = 1 vin\=a\d{d}ik\=a; 60 vin\=a\d{d}ik\=as = 1 gha\d{t}ik\=a; 60 gha\d{t}ik\=as = 1 civil day; 30 days = 1 month.
\end{quote}



\subsection*{Kalpa Constants --- Bilingual Table}
\addcontentsline{toc}{subsection}{Kalpa Constants (Bilingual)}

\begin{longtable}{p{0.42\textwidth}p{0.48\textwidth}}
\toprule
\textbf{Sanskrit Term} & \textbf{Value / English Translation} \\
\midrule
\endhead
\dev{कल्पः} (kalpa) --- Aeon & 4,320,000,000 solar years \\
\dev{सावनदिनानि} (s\=avana-dina) --- Civil days & 1,577,917,828,000 \\
\dev{सौरमासाः} (saura-m\=asa) --- Solar months & 51,840,000,000 \\
\dev{चान्द्रमासाः} (c\=andra-m\=asa) --- Lunar months & 53,433,300,000 \\
\dev{अधिमासाः} (adhi-m\=asa) --- Intercalary months & 1,593,300,000 \\
\dev{शशिदिवसाः} ({\'{s}}a\'{s}i-divasa) --- Lunar days & 1,602,999,000,000 \\
\dev{क्षयाहाः} (k\d{s}aya-aha) --- Omitted tithis & 25,082,580,000 \\
\bottomrule
\end{longtable}

\subsection*{Epoch Calculation (S\d{r}\d{s}\d{t}i-sa\d{m}vatsara)}
\addcontentsline{toc}{subsection}{Epoch Calculation}



\noindent
\begin{quote}
\dev{कल्पपरार्ध मनवः षट्काश्य गताश्चतुर्द्वियुगानि युगानाम् ।\\
श्रुतपुराणादीनिक्षेपानां क्षयः सूर्यदिनं चक्षुः ॥}\\[0.3cm]
\dev{नवगणराशिमुनिन् नव समन्वन्तरदेवः शकुन्तपाते ।\\
शर्वमतेः सम्बत्सरात् संख्यां शरवच्छतं रूपैः ॥}
\end{quote}

\hfill\noindent
\begin{quote}
\textit{``Six manus have elapsed in half a kalpa; and of the seventh manu, twenty-seven caturyugis have gone by; of the twenty-eighth caturyugi, the three yugas k\d{r}ta, dv\=apara and tret\=a have elapsed, and also 3179 years of the present kaliyuga.''}

The total elapsed period at Brahmagupta's epoch (year 628 CE):
\begin{align*}
T &= 6 \times 71 \times 4{,}320{,}000 \\
  &\quad + 7 \times 4 \times 432{,}000 \\
  &\quad + 27 \times 4{,}320{,}000 \\
  &\quad + (1{,}728{,}000 + 1{,}296{,}000 + 864{,}000) + 3{,}179 \\
  &= 1{,}972{,}947{,}179 \text{ years.}
\end{align*}
\end{quote}



\subsection*{Mean Longitude Calculation}
\addcontentsline{toc}{subsection}{Mean Longitude}

\textbf{General method:}
\begin{enumerate}
  \item Compute the \textit{aharga\d{n}a} $A$ = civil days elapsed since epoch
  \item Multiply by planet's revolution-number $R$ in a kalpa
  \item Divide by civil days in kalpa $D$:
  \[
  \text{Mean longitude} = \frac{A \times R}{D} \text{ (in revolutions, signs, degrees, etc.)}
  \]
\end{enumerate}

\longline

%% ========================================================================
%% CHAPTER XII
%% ========================================================================
\section{Chapter XII: Brahmagupta as a Critic}
\label{sec:critic}



\noindent
\begin{quote}
\dev{मध्यज्योदयज्यादृक्षेपज्यादृग्ज्यादृग्गतिज्याभिः ।\\
लम्बनं विद्धि पञ्चभिर्ज्याभिरवनतिं तथा ॥}
\end{quote}

\hfill\noindent
\begin{quote}
\textit{``Know the lambana (difference of parallaxes in longitude) by means of five Rsines: the madhya-jy\=a, the udaya-jy\=a, the d\d{r}k-k\d{s}epa-jy\=a, the d\d{r}g-jy\=a, and the d\d{r}g-gati-jy\=a; and the avanati (difference in latitude) likewise.''}

The \textit{lambana} is computed from five Rsines:
\begin{itemize}
  \item \textit{madhya-jy\=a} = $R\sin z_m$ (zenith distance of meridian ecliptic point)
  \item \textit{udaya-jy\=a} = $\dfrac{R\sin L \cdot R\sin\varepsilon}{R\cos\phi}$
  \item \textit{d\d{r}k-k\d{s}epa-jy\=a} = $R\sin z_c$ (zenith distance of central ecliptic point)
  \item \textit{d\d{r}g-jy\=a} --- related to observer's zenith distance
  \item \textit{d\d{r}g-gati-jy\=a} --- related to the motion of the observer
\end{itemize}
where $L$ = longitude of horizon ecliptic point, $\varepsilon$ = obliquity of ecliptic, $\phi$ = latitude of place.
\end{quote}



\longline

%% ========================================================================
%% CHAPTER XIII
%% ========================================================================
\section{Chapter XIII: Astronomical Instruments}
\label{sec:instruments}

\subsection*{The Gnomon (\'{S}a\d{n}ku)}
\addcontentsline{toc}{subsection}{Gnomon (Bilingual)}



\noindent
\begin{quote}
\dev{द्व्यङ्गुलायामं शूलाग्रं द्वादशाङ्गुलमायतं शङ्कुम् ।\\
वृत्तात्तात् सूच्यग्रं सच्छिद्रं व्याख्यातं प्रिये ॥}
\end{quote}

\hfill\noindent
\begin{quote}
\textit{``(The gnomon) two a\~ngulas wide at the base, pointed as a needle (\'{s}\=ula-agra), twelve a\~ngulas in length, and full of holes from the basic circular part to the pointed extremity --- thus is the \'{s}a\d{n}ku (gnomon) described, O beloved.''} --- \textit{BrSpSi}.\ XXII.\ 39.

The gnomon dimensions: base width = 2 a\~ngulas ($\approx$ 3.75 cm), height = 12 a\~ngulas ($\approx$ 22.5 cm), tapering to a point.
\end{quote}

\noindent
\begin{quote}
\dev{निवाते त्रिपदे तोयपूर्णे कूपे निधाय भूमौ ।\\
सूक्ष्मं छिद्रं तत्र कुर्यात् पयः स्यन्दनार्थम् ॥}
\end{quote}

\hfill\noindent
\begin{quote}
\textit{``When there is no wind, placing a jar full of water upon a tripod on the ground which has been made plane by means of eye or thread, one should bore a fine hole (at the bottom) so that the water may fall drop by drop.''}

\textbf{Procedure for testing level ground:} Place a water-filled jar on a tripod on leveled ground. A fine hole allows water to drip steadily, confirming the surface is horizontal.
\end{quote}



\longline

%% ========================================================================
%% SANSKRIT TEXT: MADHYAMADHIKARA --- FULL BILINGUAL
%% ========================================================================
\newpage
\section{Madhyam\=adhik\=ara: Full Bilingual Text}
\label{sec:sanskrit}

\begin{center}
{\Large\dev{॥ श्रीगणेशाय नमः ॥}}

\vspace{0.5cm}
{\Large\dev{अथ ब्राह्मस्फुटसिद्धान्तस्य}}

\vspace{0.3cm}
{\large\dev{पूर्वद्वशाध्यायां मध्यमाधिकारः}}

\vspace{0.5cm}
\textit{Br\=ahmasphu\d{t}siddh\=anta: First Chapter --- Madhyam\=adhik\=ara}
\end{center}

\vspace{0.8cm}

\noindent\rule{\textwidth}{1pt}

\vspace{0.5cm}

%% Verse 1


\noindent
\begin{quote}
\dev{जयति प्रणतसुरासुरौरगीन्द्ररत्नप्रभाच्छुरितपादपङ्कजः ।\\
कर्ता जगत्पतिरदित्यतुल्ययानां महादेवः ॥ १ ॥}
\end{quote}

\hfill\noindent
\begin{quote}
\textit{jayati pra\d{n}ata-sura-asura-aurag\={i}ndra-ratna-prabh\=a-cchurita-p\=ada-pa\~nkaja\d{h} |\\
kart\=a jagat-patir-aditya-tulya-y\=an\=am mah\=adeva\d{h} || 1 ||}

``Victorious is Mah\=adeva, the Lord of the Universe, the Creator, the lotus-feet of whom are rubbed with the lustre of jewels from the heads of gods, demons, and serpent-kings who bow before Him --- He whose chariot is equal to that of the Sun.''
\end{quote}



\vspace{0.3cm}

%% Verse 2


\noindent
\begin{quote}
\dev{ब्रह्मगुप्तो ग्रहगणितं महता कालेन यत् खल्वीमतम् ।\\
क्षीणेऽधुनी तद्विपरीतवद्ब्राह्मणैरुपेतं ॥ २ ॥}
\end{quote}

\hfill\noindent
\begin{quote}
\textit{brahmagupto graha-ga\d{n}ita\d{m} mahat\=a k\=alena yat khalv\={i}matam |\\
k\d{s}\={i}\d{n}e 'dhun\={i} tad-vipar\={i}ta-vad-br\=ahma\d{n}air upeta\d{m} || 2 ||}

``Brahmagupta (says): The science of the computation of planets, which was well-established in the past, has now become corrupted; being contrary (to the truth), it has been mixed up by the br\=ahma\d{n}as.''
\end{quote}



\vspace{0.3cm}

%% Verse 3


\noindent
\begin{quote}
\dev{ध्रुवबलार्कसुतबुधज्ञ ज्योतिर्विदां प्रीतिषद्गमदो ।\\
पोषणार्थं विवृतन्तरव्यक्तं सह प्रबन्धेन सूचितं ॥ ३ ॥}
\end{quote}

\hfill\noindent
\begin{quote}
\textit{dhruva-bala-arka-suta-budhaj~na jyotir-vid\=\=a\d{m} pr\={i}ti-\d{s}ad-gama-do |\\
po\d{s}a\d{n}\=artha\d{m} viv\d{r}t\=antara-vyakta\d{m} saha prabandhena s\=ucitam || 3 ||}

``(I write this) for the pleasure of astronomers, who are learned in the (motions of) Dhruva (pole star), Bala (Mars), Arka (Sun), Suta (Saturn), Budha (Mercury), and J~na (Jupiter); for the preservation (of knowledge), the obscure meaning is made clear along with the connected exposition.''
\end{quote}



\vspace{0.3cm}

%% Verse 4


\noindent
\begin{quote}
\dev{चैत्र्यादितदेवद्विजनोद्वन्माससर्वेष युगस्य ।\\
सृष्टिच्छादो लंकायां समं प्रवृता दिनेकलस्य ॥ ४ ॥}
\end{quote}

\hfill\noindent
\begin{quote}
\textit{caitry-\=adi-tad-eva-dvija-nodvan-m\=asa-sarve\d{s}a yugasya |\\
s\d{r}\d{s}\d{t}i-cch\=ado la\~nk\=ay\=a\d{m} sama\d{m} prav\d{r}t\=a dine-kalasya || 4 ||}

``Beginning with Caitra, all the months of a yuga --- (which are) the creation-covering --- proceed simultaneously at La\~nk\=a with the beginning of the day of Brahm\=a (dina-kala).''
\end{quote}



\vspace{0.3cm}

%% Verse 5


\noindent
\begin{quote}
\dev{प्राग्र्योद्दीनाटिकाद्वो द्विद्विर्यष्टिका त्रिनाटिका षड्घ्ना ।\\
घटिका षड्घ्ना दिवलो दिवसानां त्रिंशता मासः ॥ ५ ॥}
\end{quote}

\hfill\noindent
\begin{quote}
\textit{pr\=ag-ry-udda-n\=a\d{t}ik\=a-dvau dvi-dvi-ya\d{s}\d{t}ik\=a tri-n\=a\d{t}ik\=a \d{s}a\d{d}-ghn\=a |\\
gha\d{t}ik\=a \d{s}a\d{d}-ghn\=a diva-lo divas\=an\=a\d{m} tri\d{m}\'sat\=a m\=asa\d{h} || 5 ||}

``Two atoms (param\=a\d{n}u) make an a\d{n}u; two a\d{n}us make a lava; two lavas make a nime\d{s}a; three nime\d{s}as make a k\d{s}a\d{n}a; five k\d{s}a\d{n}as make a k\d{s}athi; four k\d{s}athis make a muh\=urta; eight muh\=urtas make a prahara; four praharas make a day (aho-r\=atra); thirty days make a month.''

\textbf{Time scale:} 2 param\=a\d{n}u = 1 a\d{n}u $\to$ 2 a\d{n}u = 1 lava $\to$ 2 lava = 1 nime\d{s}a $\to$ 3 nime\d{s}a = 1 k\d{s}a\d{n}a $\to$ 5 k\d{s}a\d{n}a = 1 k\d{s}athi $\to$ 4 k\d{s}athi = 1 muh\=urta $\to$ 8 muh\=urta = 1 prahara $\to$ 4 prahara = 1 day $\to$ 30 days = 1 month.
\end{quote}



\vspace{0.5cm}
\noindent\rule{\textwidth}{0.6pt}
\vspace{0.5cm}

%% Verses on Solar and Lunar Time (22-24)
\subsection*{Verses 22--24: Solar and Lunar Month Relations}
\addcontentsline{toc}{subsection}{Verses 22--24: Solar and Lunar Months}



\noindent
\begin{quote}
\dev{परिवर्तः सप्तचतुःषष्टिर सूर्यस्य दिनस्य मितिः ।\\
रवि भगणानां मानः सावनदिनस्य कुट्टकैरेष ॥}
\end{quote}

\hfill\noindent
\begin{quote}
\textit{parivarta\d{h} sapta-catu\d{h}\d{s}a\d{s}\d{t}ir s\=uryasya dinasya miti\d{h} |\\
ravi bhaga\d{n}\=an\=\=a\d{m} m\=ana\d{h} s\=avana-dinasya ku\d{t}\d{t}akair e\d{s}a ||}

``The revolution-number of the Sun is 64 + 7 = 71 (per caturyuga); this is the measure of solar days. By the pulverizer (ku\d{t}\d{t}aka), this (gives) the civil days.''

\textbf{Rule:} A solar year $\approx$ 365 days 15 min 30 sec. The revolution-number of the Sun = 4,320,000 per kalpa.
\end{quote}

\noindent
\begin{quote}
\dev{रवि भगणार्धबधा द्वादशांशेन भवति रविमासः ।\\
भगणान्तरं रविदिनः शशिमासः सूर्यमासेन ॥}
\end{quote}

\hfill\noindent
\begin{quote}
\textit{ravi bhaga\d{n}\=ardha-badh\=a dv\=ada\'{s}\=a\d{m}\'sena bhavati ravi-m\=asa\d{h} |\\
bhaga\d{n}\=\=antara\d{m} ravi-dina\d{h} \'{s}a\'{s}i-m\=asa\d{h} s\=urya-m\=asena ||}

``The product of half the revolutions of the Sun divided by twelve becomes a solar month. The difference of revolutions (between Sun and Moon) is the lunar month exceeding the solar month.''

\textbf{Formula:}
\[
\text{Solar month} = \frac{R_{\odot}/2}{12} \quad ; \quad \text{Lunar month} = R_{\Moon} - R_{\odot}
\]
\end{quote}

\noindent
\begin{quote}
\dev{अधिमासाः शशिमासाविन्द्वदिनस्य भवति शशिदिवसाः ।\\
शशिसावनदिवसान्तरनिवमानि तिथिः शशाङ्कदिनं ॥}
\end{quote}

\hfill\noindent
\begin{quote}
\textit{adhi-m\=as\=a\d{h} \'{s}a\'{s}i-m\=s\=av-indv-adinasya bhavati \'{s}a\'{s}i-divas\=a\d{h} |\\
\'{s}a\'{s}i-s\=avana-divasa-antara-nivam\=ani tithi\d{h} \'{s}a\'{s}\=a\~nka-dina\d{m} ||}

``The intercalary months (adhi-m\=asa) are the lunar months exceeding the solar months. The lunar days are (obtained from) the difference between lunar and civil days; the tithi (lunar day) is the measure of the difference between lunar and civil days.''

\textbf{Key relations:}
\begin{align*}
\text{adhi-m\=asa} &= \text{lunar months} - \text{solar months} \\
\text{\'{s}a\'{s}i-divasa} &= \text{lunar days} - \text{civil days} \\
\text{tithi} &= \text{1/30 of a lunar month}
\end{align*}
\end{quote}



\vspace{0.5cm}
\noindent\rule{\textwidth}{0.6pt}
\vspace{0.5cm}

%% Epoch Verses (26-27)
\subsection*{Verses 26--27: Epoch and Aharga\d{n}a}
\addcontentsline{toc}{subsection}{Verses 26--27: Epoch Calculation}



\noindent
\begin{quote}
\dev{कल्पपरार्ध मनवः षट्काश्य गताश्चतुर्द्वियुगानि युगानाम् ।\\
श्रुतपुराणादीनिक्षेपानां क्षयः सूर्यदिनं चक्षुः ॥}
\end{quote}

\hfill\noindent
\begin{quote}
\textit{kalpa-par\=arddha manava\d{h} \d{s}a\d{t}-k\=a\'{s}ya gat\=a\'{s} catur-dvi-yug\=ani yug\=an\=am |\\
\'{s}ruta-pur\=a\d{n}\=\=adi-nik\d{s}ep\=an\=am k\d{s}aya\d{h} s\=urya-dina\d{m} cak\d{s}u\d{h} ||}

``In half a kalpa, six manus have elapsed, and twenty-seven caturyugis of the (seventh) manu. (This gives) the destruction (elapsed period) of the s\=urya-dina (solar days) --- the eye (of calculation).''

The \textit{aharga\d{n}a} (day-count) from the kalpa epoch:
\[
A = D_{\text{manus}} + D_{\text{caturyugis}} + D_{\text{current yuga}}
\]
\end{quote}

\noindent
\begin{quote}
\dev{नवगणराशिमुनिन् नव समन्वन्तरदेवः शकुन्तपाते ।\\
शर्वमतेः सम्बत्सरात् संख्यां शरवच्छतं रूपैः ॥}
\end{quote}

\hfill\noindent
\begin{quote}
\textit{nava-ga\d{n}a-r\=a\'{s}i-munin nava samanvantara-deva\d{h} \'{s}akunta-p\=ate |\\
\'{s}arva-mate\d{h} sa\d{m}batsar\=at sa\d{m}khy\=a\d{m} \'{s}aravac-chata\d{m} r\=upai\d{h} ||}

``O Muni of nine-digit signs! (This gives) the number of years from the \'{s}arva-mata (doctrine year), 106 (\'{s}ara) increased by unity --- 3179 years of kaliyuga (at the epoch \'{s}aka 550 = 628 CE).''

Brahmagupta's epoch: year 3179 of kaliyuga = \'{s}aka 550 = 628 CE.
\end{quote}



\vspace{0.8cm}

%% ========================================================================
%% RATIONAL GEOMETRICAL FIGURES --- BILINGUAL
%% ========================================================================
\section*{Bilingual Appendix: Rational Geometrical Figures}
\addcontentsline{toc}{section}{Bilingual Appendix: Rational Figures}



\noindent
\begin{quote}
\dev{इष्टयष्टेः समच्छेदं रूपेण हीनमुभयोर्मिथः ।\\
षड्भागेन समज्ञात्र्यक्षेत्रस्य भुजे ॥}
\end{quote}

\hfill\noindent
\begin{quote}
\textit{i\d{s}\d{t}a-ya\d{s}\d{t}e\d{h} sama-ccheda\d{m} r\=upe\d{n} a h\={i}nam ubhayor mitha\d{h} |\\
\d{s}a\d{d}-bh\=agena sama-j~n\=a-try-ak\d{s}etrasya bhuje ||}

``The square of the optional side divided and then diminished by an optional number; half the result is the upright (koti), and that increased by the optional number gives the hypotenuse of a rectangle.''

\textbf{Formula:} For arbitrary rational $m, n$:
\[
\text{Base} = m, \quad \text{Upright} = \frac{1}{2}\!\left(\frac{m^2}{n} - n\right), \quad \text{Hypotenuse} = \frac{1}{2}\!\left(\frac{m^2}{n} + n\right)
\]
\end{quote}



\vspace{0.3cm}



\noindent
\begin{quote}
\dev{यस्मिन् द्वे समानी बाहू अन्यद् द्वयमसमानि च ।\\
तयोर्वर्गयोगो विभाज्यो भेदेन समद्विबाहुक्षेत्रे ॥}
\end{quote}

\hfill\noindent
\begin{quote}
\textit{yasmin dve sam\=an\={i} b\=ah\=u anyad dvayam-asam\=ani ca |\\
tayor varga-yogo vibh\=ajyo bhedena sama-dvi-b\=ahu-k\d{s}etre ||}

``Where there are two equal sides and two other unequal sides --- the sum of their squares divided by their difference (gives the base of an isosceles triangle).''

For two unequal integers $p, q$: equal sides = $\frac{1}{2}\!\left(\frac{p^2+q^2}{p-q} + (p-q)\right)$, base = $\frac{p^2+q^2}{p-q}$.
\end{quote}



\vspace{0.3cm}



\noindent
\begin{quote}
\dev{सर्वदोरल्पचतुष्टयक्षेत्रफलमूलं द्विभुजक्षेत्रे ।\\
द्वयोर्द्वयोर्द्धयोर्घातं मूलं द्विभुजक्षेत्रस्य ॥}
\end{quote}

\hfill\noindent
\begin{quote}
\textit{sarva-dor-alpa-catu\d{s}\d{t}aya-k\d{s}etra-phala-m\=ula\d{m} dvi-bhuja-k\d{s}etre |\\
dvayor dvayor ddhayor gh\=ata\d{m} m\=ula\d{m} dvi-bhuja-k\d{s}etrasya ||}

``The square root of the product of four sets of half the sum of the sides, each diminished by one side, is the exact area of a triangle (and) quadrilateral.''

\textbf{Brahmagupta's Formula:}
\[
A = \sqrt{(s-a)(s-b)(s-c)(s-d)}, \quad s = \frac{a+b+c+d}{2}
\]
This generalizes Heron's formula for triangles (when $d = 0$).
\end{quote}



\vspace{0.8cm}

%% ========================================================================
%% CRITICAL APPARATUS
%% ========================================================================
\section*{Critical Apparatus: Selected Variants}
\addcontentsline{toc}{section}{Critical Apparatus}

\begin{longtable}{p{0.08\textwidth}p{0.42\textwidth}p{0.42\textwidth}}
\toprule
\textbf{V.} & \textbf{Adopted Reading} & \textbf{Variant / Note} \\
\midrule
\endhead
1 & \dev{जयति प्रणत...} & \dev{श्री नमः} (mss. A, B); \dev{श्री रामाय} (ms. C) \\
2 & \dev{क्षीणेऽधुनी...} & \dev{क्षीणेऽद्यापि} (ms. D) --- ``even now corrupted'' \\
3 & \dev{विवृतान्तरव्यक्तं} & \dev{विवृतन्तरव्यक्तं} (mss. A, B) --- sandhi variant \\
4 & \dev{समं प्रवृता} & \dev{समप्रवृता} (ms. B) --- compound variant \\
5 & \dev{प्राग्र्योद्दीनाटिका} & \dev{प्रागुद्दीनाटिका} (ms. C) --- orthographic \\
22 & \dev{सप्तचतुःषष्टिः} & 71 revolutions per caturyuga (confirmed) \\
23 & \dev{भगणार्धबधा} & \dev{भगणार्धबुद्धा} (ms. E) --- scribal error \\
26 & \dev{कल्पपरार्ध} & \dev{कल्पपरार्द्ध} (ms. D) --- gemination variant \\
27 & \dev{शरवच्छतं} & \dev{शरवत् शतं} (mss. A, B) --- sandhi split \\
\bottomrule
\end{longtable}

\longline

%% ========================================================================
%% COLOPHON
%% ========================================================================
\newpage
\section*{Colophon}
\addcontentsline{toc}{section}{Colophon}

\begin{center}
\begin{minipage}{0.85\textwidth}
\centering
{\large\dev{॥ इति ब्राह्मस्फुटसिद्धान्तस्य मध्यमाधिकारः समाप्तः ॥}}

\vspace{0.5cm}
\textit{iti br\=ahmasphu\d{t}siddh\=antasya madhyam\=adhik\=ara\d{h} sam\=apta\d{h}}

\vspace{0.5cm}
``Here ends the Madhyam\=adhik\=ara of the Br\=ahmasphu\d{t}siddh\=anta.''

\vspace{1.5cm}
\rule{0.6\textwidth}{0.4pt}

\vspace{0.5cm}
{\small Bilingual Edition: Sanskrit (Devanagari) / English (IAST)}

\vspace{0.3cm}
{\small Typeset with \XeLaTeX\ using Noto Serif Devanagari}

\vspace{0.3cm}
{\small Chunk 3 (Pages 281--420)}
\end{minipage}
\end{center}

\end{document}



% ============================================================
% Source: translations/en/indian/brahmagupta-brahmasphutasiddhanta-pt1-chunk4_bilingual.tex
% ============================================================

%% ========================================================================
%% Brāhmasphuṭasiddhānta of Brahmagupta (628 CE)
%% BILINGUAL EDITION: English-Sanskrit, Part 1 Chunk 4 (pages 421--560)
%% Content: Chapters III--XIII (partial)
%% Based on the Sudhākara Dvivedī edition (Pandita, 1901--1902)
%% ========================================================================
\documentclass[11pt,a4paper]{article}

%% -------- Core Packages --------
\usepackage{fontspec}
\usepackage{polyglossia}
\usepackage{amsmath,amssymb,amsthm}
\usepackage{geometry}
\usepackage{graphicx}
\usepackage{xcolor}
\usepackage{titlesec}
\usepackage{fancyhdr}
\usepackage{enumitem}
\usepackage{array,booktabs,longtable,multirow}
\usepackage{hyperref}
% lettrine, microtype, tocloft removed for compilation
\usepackage{indentfirst}

%% -------- Page Geometry --------
\geometry{
  paperwidth=210mm,paperheight=297mm,
  left=15mm,right=15mm,top=22mm,bottom=22mm,
  headheight=14pt,footskip=25pt
}

%% -------- Language Setup --------
\setmainlanguage{english}
\setotherlanguage{sanskrit}

%% -------- Font Configuration --------
\setmainfont{Linux Libertine O}
\setsansfont{Linux Biolinum O}
\setmonofont{DejaVu Sans Mono}

%% Devanagari (Sanskrit/Hindi) Font -- using specified external location
\newfontfamily\devanagarifont[Script=Devanagari,Path=/usr/share/fonts/truetype/noto/]{NotoSerifDevanagari-Regular.ttf}
\newfontfamily\devanagarifontbf[Script=Devanagari,Path=/usr/share/fonts/truetype/noto/]{NotoSerifDevanagari-Bold.ttf}
\newfontfamily\hindifont[Script=Devanagari,Path=/usr/share/fonts/truetype/noto/]{NotoSansDevanagari-Regular.ttf}

%% -------- Colors --------
\definecolor{titlecolor}{RGB}{45,55,72}
\definecolor{sectioncolor}{RGB}{55,65,81}
\definecolor{notecolor}{RGB}{120,53,15}
\definecolor{rulecolor}{RGB}{180,160,140}
\definecolor{varecolor}{RGB}{70,90,120}
\definecolor{endmarkcolor}{RGB}{139,69,19}
\definecolor{sanskritbg}{RGB}{255,250,245}
\definecolor{englishbg}{RGB}{245,248,255}

%% -------- Section Formatting --------
\titleformat{\section}
  {\normalfont\Large\bfseries\color{sectioncolor}}
  {\thesection}{1em}{}
\titleformat{\subsection}
  {\normalfont\large\bfseries\color{sectioncolor}}
  {\thesubsection}{1em}{}
\titleformat{\subsubsection}
  {\normalfont\normalsize\bfseries\color{sectioncolor}}
  {\thesubsubsection}{1em}{}

%% -------- Header/Footer --------
\pagestyle{fancy}
\fancyhf{}
\fancyhead[L]{\small\textit{Brāhmasphuṭasiddhānta --- Bilingual Edition --- Chunk 4 (pp.\ 421--560)}}
\fancyhead[R]{\small\thepage}
\renewcommand{\headrulewidth}{0.4pt}
\renewcommand{\footrulewidth}{0pt}

%% -------- Custom Commands --------
\newcommand{\longline}{\noindent\rule{\textwidth}{0.4pt}}
\newcommand{\shortline}{\noindent\rule{0.5\textwidth}{0.4pt}\par}
\newcommand{\cnote}[1]{\textcolor{notecolor}{\textit{#1}}}
\newcommand{\dev}[1]{{\devanagarifont #1}}
\newcommand{\devbf}[1]{{\devanagarifontbf #1}}
\newcommand{\hin}[1]{{\hindifont #1}}

%% Bilingual parallel column environment
\newenvironment{bilingual}{
  \noindent
  \begin{minipage}[t]{0.47\textwidth}
  \vspace{0pt}
}{
  \end{minipage}%
}
\newcommand{\sanskritcol}[1]{%
  \begin{minipage}[t]{0.47\textwidth}
  \vspace{0pt}
  \begin{flushleft}
  #1
  \end{flushleft}
  \end{minipage}%
  \hfill
}
\newcommand{\englishcol}[1]{%
  \begin{minipage}[t]{0.47\textwidth}
  \vspace{0pt}
  \begin{flushleft}
  #1
  \end{flushleft}
  \end{minipage}%
}

%% Sanskrit verse in bilingual format
\newcommand{\biverseline}[2]{%
  \begin{bilingual}
  \sanskritcol{\dev{#1}}
  \englishcol{#2}
  \end{bilingual}
  \vspace{0.15cm}
}

%% Column headers
\newcommand{\columnheaders}{%
  \noindent
  \begin{minipage}[t]{0.47\textwidth}
  \centering
  \textbf{\large Sanskrit (Devanagari)}\\[0.1cm]
  \rule{0.8\textwidth}{0.5pt}
  \end{minipage}%
  \hfill
  \begin{minipage}[t]{0.47\textwidth}
  \centering
  \textbf{\large English Translation (with IAST)}\\[0.1cm]
  \rule{0.8\textwidth}{0.5pt}
  \end{minipage}%
  \vspace{0.3cm}
  \par
}

%% Variant reading environment (critical apparatus)
\newenvironment{variants}{%
  \small%
  \color{varecolor}%
  \list{}{\leftmargin=2em\itemindent=-1em}%
}{\endlist}

%% Verse number marker
\newcommand{\vnum}[1]{\textbf{\textcircled{\scriptsize #1}}}

%% End of chapter marker
\newcommand{\chapterend}{%
  \begin{center}%
  \textcolor{endmarkcolor}{$\clubsuit$~~$\clubsuit$~~$\clubsuit$}%
  \end{center}%
}

%% -------- Hyperref Setup --------
\hypersetup{
  colorlinks=true,
  linkcolor=sectioncolor,
  citecolor=sectioncolor,
  urlcolor=sectioncolor,
  pdfencoding=auto,
  pdftitle={Brahmasphutasiddhanta of Brahmagupta - Bilingual Edition Part 1 Chunk 4},
  pdfauthor={Brahmagupta (628 CE)},
  pdfsubject={Sanskrit Astronomy and Mathematics}
}

%% -------- TOC Depth --------
\setcounter{tocdepth}{3}
\setcounter{secnumdepth}{3}

%% ========================================================================
\begin{document}

%% ========================================================================
%% TITLE PAGE
%% ========================================================================
\thispagestyle{empty}
\begin{center}
\vspace*{1.5cm}

{\Huge\color{titlecolor}\dev{ब्रह्मस्फुटसिद्धान्त}}\\[0.6cm]
{\LARGE\color{sectioncolor} Brāhmasphuṭasiddhānta}\\[0.3cm]
{\large of \textbf{Brahmagupta} (628 CE)}\\[1cm]

\longline\\[0.6cm]

{\Large\textbf{Bilingual Edition: Sanskrit--English}}\\[0.3cm]
{\large\textbf{Part 1, Chunk 4: Pages 421--560}}\\[0.3cm]
{\normalsize (Book pages 32--171)}\\[0.6cm]

\longline\\[0.6cm]

{\normalsize Covering:}\\[0.3cm]
\dev{तृप्रश्नाधिकारः (अन्तिम भाग)}\\
\dev{चन्द्रग्रहणाध्यायः} --- \textit{Candragrahaṇādhikāya (Lunar Eclipse)}\\[0.1cm]
\dev{सूर्यग्रहणाध्यायः} --- \textit{Sūryagrahaṇādhikāya (Solar Eclipse)}\\[0.1cm]
\dev{उदयास्ताध्यायः, चन्द्रशृङ्गोन्नत्यध्यायः, चन्द्रच्छायाध्यायः, ग्रहयुत्यध्यायः}\\[0.1cm]
\dev{तन्त्रपरीक्षाध्यायः (दूषणाध्यायः)} --- \textit{Tantraparīkṣādhikāya (Critical Examination)}\\[0.1cm]
\dev{गणिताध्यायः} --- \textit{Gaṇitādhyāya (Mathematics)}\\[0.1cm]
\dev{प्रश्नाध्यायः (प्रारम्भ)} --- \textit{Praśnādhyāya (Problems)}\\[0.6cm]

\longline\\[0.6cm]

{\small Based on the critical edition of}\\
{\small Sudhākara Dvivedī (\textit{The Pandita}, 1901--1902)}\\[0.5cm]

{\footnotesize Sanskrit text in Devanagari $\parallel$ English translation with IAST transcription}\\
{\footnotesize Typeset with XeLaTeX using Noto Serif Devanagari}

\vfill
\end{center}

\newpage

%% ========================================================================
%% TABLE OF CONTENTS
%% ========================================================================
\tableofcontents
\newpage

%% ========================================================================
%% INTRODUCTION TO THIS CHUNK
%% ========================================================================
\section*{Introduction to Chunk 4}
\addcontentsline{toc}{section}{Introduction to Chunk 4}

This bilingual edition presents the Sanskrit text of pages 421--560 from the first volume of the
\textit{Brāhmasphuṭasiddhānta} of Brahmagupta alongside its English translation. These pages correspond to book pages 32--171
of Dvivedī's edition, and contain material from Chapters III through XIII (partial) of the work.

The \textit{Brāhmasphuṭasiddhānta} (lit.\ ``Corrected Doctrine of Brahma'') is Brahma-\linebreak[0]gupta's
masterpiece, composed in 628 CE. It consists of 24 chapters, with the first ten chapters
(\textit{pūrvadaśādhyāyi}) covering astronomical topics and the latter chapters
(\textit{uttara}) covering mathematics, algebra, and more advanced astronomy.

\subsection*{Content Overview}

The present chunk covers:

\begin{enumerate}[label=\textbf{\arabic*.}]
  \item \textbf{Triprasnādhikāra} (Chapter III, conclusion): Rules for solving the three
  problems of diurnal motion -- finding time, latitude/direction, and the gnomon shadow.

  \item \textbf{Candragrahaṇādhikāya} (Chapter IV): Computation of lunar eclipses,
  including the calculation of shadow dimensions, immersion and emersion times, and color
  phases of the eclipsed Moon.

  \item \textbf{Sūryagrahaṇādhikāya} (Chapter V): Computation of solar eclipses,
  including parallax corrections (lambana), semi-durations, and visibility conditions.

  \item \textbf{Udayastādhikāra} (Chapter VI): Rising and setting of planets.

  \item \textbf{Candraśṛṅgonnatyādhikāya} (Chapter VII): Elevation of the Moon's horns.

  \item \textbf{Candraccāyādhikāya} (Chapter VIII): The Moon's shadow.

  \item \textbf{Grahayutyādhikāya} (Chapter IX): Conjunction of planets.

  \item \textbf{Tantraparīkṣādhikāya / Dūṣaṇādhikāya} (Chapter XI): Critical examination
  of other astronomical systems, including refutations of doctrines of Āryabhaṭa, Śrīṣeṇa,
  Viṣṇucandra, Lāṭadeva, and Vijayanandi.

  \item \textbf{Gaṇitādhyāya} (Chapter XII): The celebrated chapter on mathematics,
  containing arithmetic operations, mixtures, series, geometry (plane figures),
  shadow problems, and algebra (including the \textit{kuṭṭaka} or pulverizer).

  \item \textbf{Praśnādhyāya} (Chapter XIII, beginning): Mathematical problems and
  their solutions.
\end{enumerate}

\subsection*{How to Read This Bilingual Edition}

This edition uses a parallel-column format:
\begin{itemize}[label=\textbullet]
  \item \textbf{Left column:} Sanskrit text in Devanagari script
  \item \textbf{Right column:} English translation with Sanskrit technical terms in IAST transliteration
\end{itemize}

The critical apparatus (variant readings) appears after each group of verses, compiled
by Sudhākara Dvivedī. Manuscript sigla:
\begin{itemize}[label=\textbullet]
  \item \dev{(ग)} = reading of manuscript G
  \item \dev{(क)} = reading of manuscript K
  \item \dev{(ख)} = reading of manuscript Kh
  \item \dev{(च)} = reading of manuscript Ch
  \item \dev{(घ)} = reading of manuscript Gh
\end{itemize}

The notation ``X \dev{for} (Y)'' indicates that the edition reads X where manuscripts
have variant Y. Hindi notes marked with \dev{वि०} (\textit{viveka}) provide editorial
commentary.

\newpage

%% ========================================================================
%% CHAPTER 3: TRIPRASNADHIKARA (Conclusion)
%% ========================================================================
\section*{Chapter III: Triprasnādhikāra (Conclusion)}
\addcontentsline{toc}{section}{Chapter III: Triprasnādhikāra (Conclusion)}
\markboth{Chapter III: Triprasnādhikāra}{}

\begin{center}
\dev{तृतीयोऽध्यायः} --- \textit{The chapter on the three problems relating to diurnal motion}
\end{center}

\vspace{0.3cm}
\columnheaders

\subsection*{Verses 25--27: Ghaṭikā computations and eclipse timing}

\noindent
\sanskritcol{\dev{प्राक् पश्चाद्वा यामिध्वंटिकाभिरदिनिदलान्नतः सूर्यः} \\ \dev{तिथ्यन्ते तद्रहिते त्रिशत् घटिकावशेषासि ॥ २५ ॥} \\[0.2cm]
\dev{विपरीतमर्धरात्राच्चन्द्रग्रहणे शनै रविग्रहणे} \\ \dev{सूर्योत्पतस्ततस्ताभिरेव घटिकाभिर्द्विरपि ॥ २६ ॥} \\[0.2cm]
\dev{दिनदल परिधि स्फुट तिथिनतकेद्रज्यावधो गुरोः कन्दोः} \\ \dev{इन्दू धृति धृतिभिः १९१ नवनववेदे ४६६ व्यासाङ्ककृतिभक्तः ॥ २७ ॥}}
\englishcol{\textbf{Translation:} (25) When the Sun, diminished by the day-radius (\textit{dinadala}),
is east or west, the ghaṭikās till the end of the tithi are obtained by subtracting
that from thirty and the ghaṭikā-remainder is obtained. (26) The reverse applies at
midnight for a lunar eclipse; slowly for a solar eclipse. From the rising of the
Sun, by the same ghaṭikās twice. (27) The semi-diameter (\textit{vyāsa}) of the
Earth's shadow at the Moon's orbit is obtained: the day-diameter, circumference,
accurate tithi-nata-kendra-jyā and the degrees of Jupiter or Kandu (the Moon),
Indu --- by 191 and the new nine-veda 466, divided by the square of the radius.}

\vspace{0.4cm}
\textcolor{varecolor}{\textbf{Critical Apparatus (Verses 25--27):}}
\begin{variants}
\item[25.] \dev{तद्रहिते} (G) \dev{for} (\dev{तद्रहिते}): corrected reading.
\item[25.1] \dev{घटिकामि} (G) \dev{for} (\dev{घटिकासि}).
\item[25.2] \dev{त्रिशदधिकावशेषासि} (G) \dev{for} (\dev{त्रिशतुघटिकावशेषासिः}).
\item[25.3] \dev{दिनदलान्नतः सूर्यः} (G) \dev{for} (\dev{दिनदलान्ततः सूर्यः}).
\item[26.1] \dev{यतस्ततः} (G) \dev{for} (\dev{यतस्ततस्तु}).
\item[26.2] \dev{चन्द्रग्रहणे} (K) \dev{चन्द्रग्रहविनिर्हारणः} \dev{for} (\dev{चन्द्रग्रहणे}).
\item[27.1] \dev{दल} (G) \dev{for} (\dev{दल}).
\item[27.2] \dev{ज्या} (G) \dev{ज्यावधो} \dev{for} (\dev{ज्यावधो}).
\item[27.3] \dev{कन्दोः} (G,Kh) \dev{कन्द्रो} \dev{for} (\dev{कन्दोः}).
\item[27.4] \dev{इन्दू धृतिभिः १९१ नवनववेदे ४६६} --- \dev{संख्याएँ मूल में लुप्त हैं} (numbers omitted in the source).
\item[27.5] \dev{व्यासाङ्ककृति} (K) \dev{व्यासाङ्ककृते} \dev{for} (\dev{व्यासाङ्ककृति}).
\end{variants}

\subsection*{Verses 38--39: Sphuṭīkaraṇa of Mars}

\noindent
\sanskritcol{\dev{सप्तभिरंशगुणिता द्युगण दिवर्ध्ध राशिभुक्ता हृताप्तांशः} \\ \dev{अधिकोनकुजमर्कतात् मृगकर्कादौ स्फुटी भवतीति ॥ ३८ ॥} \\[0.2cm]
\dev{तत्स्फुटपरिधिः खनगाः ७० शीघ्रस्फुटपरिधिराप्तभागोनाः} \\ \dev{वेदजिनाखिशोनाः २४३३० स्फुटीकरणं कुजस्यैव ॥ ३९ ॥}}
\englishcol{\textbf{Translation:} (38) Multiplied by seven degrees, the day-multiplier, the daily
motion in signs, divided by the obtained degrees, the sphuṭa of Mars is increased
or decreased from the mark (starting point) at Mṛga (start of Capricorn) and Karka
(start of Cancer). (39) Its sphuṭa-circumference is 70 yojanas; the śīghra-sphuṭa
circumference diminished by obtained degrees; 24330 diminished by veda-jina-akhi;
this is the correction (\textit{sphuṭīkaraṇa}) of Mars.}

\vspace{0.4cm}
\textcolor{varecolor}{\textbf{Critical Apparatus:}}
\begin{variants}
\item[38.1] \dev{गुणिताः} (Ch) \dev{for} (\dev{गुणिता}).
\item[38.2] \dev{स्फुटो} (G,K,Kh,Ch) \dev{for} (\dev{स्फुटौ}).
\item[38.3] \dev{भवति} (G,K) \dev{भवन्ति} (Kh) \dev{for} (\dev{भवतीति}).
\item[39.1] \dev{स्फुटीकरणं} (G) \dev{स्फष्टीकरणं} (Kh) \dev{for} (\dev{स्फुटीकरणं}).
\item[39.2] \dev{राप्तभागोनाः} \dev{for} (\dev{राप्तभागोनाः}).
\item[\dev{वि०}] Numbers given in the source are omitted. Manuscripts lack the numbers in the verse.
\end{variants}

\subsection*{Verses 40--42: Celestial latitude and shadow calculations}

\noindent
\sanskritcol{\dev{छायावृत्तेज्यगा कर्णगुणा व्यासदल हृतार्काप्ता} \\ \dev{विषुवच्छाया याम्या तदन्तरैक्यं भुजस्याग्रे ॥ ४ ॥} \\[0.2cm]
\dev{शङ्कुप्राच्यपरा या छाया भुजकृतिविशेषमूलं नयेत्} \\ \dev{तत्प्राच्यपरछाया भुजाग्रयोरन्तरं कोटिः ॥ ५ ॥} \\[0.2cm]
\dev{दिग्मध्ये छायाग्रं कृत्वा शङ्कुं यथादिशं वि\-णम्} \\ \dev{दिग्मध्यस्थितदृग् छायाग्रं भ्रमति विपरीतम् ॥ ६ ॥}}
\englishcol{\textbf{Translation:} (4) The shadow-circle-radius multiplied by the hypotenuse,
divided by the semi-diameter obtained by the Sun, the viṣuvacchāyā, the yāmyā,
the sum or difference of that at the tip of the bhuja. (5) The eastern or western
shadow of the gnomon --- taking the square root of the difference of the bhuja-squares,
that eastern or western shadow is the koṭi between the tips of the bhuja. (6) Having
made the shadow-tip in the middle of the direction, the gnomon according to direction
is to be measured; the observer's shadow-tip situated in the middle of the direction
moves in reverse.}

\vspace{0.4cm}
\textcolor{varecolor}{\textbf{Critical Apparatus:}}
\begin{variants}
\item[4.1] \dev{वृत्तेज्यगा} (G,Kh) \dev{for} (\dev{वृत्तेज्काग्रा}).
\item[4.2] \dev{हतार्काप्ता} \dev{for} (\dev{हृतार्काप्ता}).
\item[5.1] \dev{भुजकृतिविशेषमूलं नयेत्} (G) \dev{for} (\dev{भुजकृतिविशेषसूलं नयेत्}).
\item[6.1] \dev{दिग्मध्ये} (G,Kh) \dev{for} (\dev{दिग्मध्ये}).
\item[6.2] \dev{शङ्कोयंथादिशं} \dev{for} (\dev{शङ्कोयंथादिशं}).
\item[6.3] \dev{भ्रमति} (K) \dev{for} (\dev{भ्रमति}).
\end{variants}

\subsection*{Verses 64--66: End of Triprasnādhikāra}

\noindent
\sanskritcol{\dev{कृष्णाचतुर्दशान्ते शकुनि पक्षेण चतुष्पदप्रयमे} \\ \dev{तिथ्यर्धं ते नागं किस्तुघ्नं प्रतिपदाद्यत् ॥ ६४ ॥} \\[0.2cm]
\dev{व्यक्तद्विकला भक्ताः खरसगुरणणः ३६० लब्धसुनमेकेन} \\ \dev{चलकरानि विवादी nyag हृतशेषे तिथिवदन्यत् ॥ ६६ ॥}}
\englishcol{\textbf{Translation:} (64) At the end of Kṛṣṇa-caturdaśī, in the bird (śakuni)
season, in the four-footed (catuṣpada) measure, the half-tithi, nāga, kistughna,
from pratipadā etc. (66) The manifest two kalās divided, the moving karas and
the hara-sa-gu-ra-ṇa-ṇaḥ 360, with one obtained sum, the remainder divided by hṛta
is like a tithi and other.}

\vspace{0.4cm}
\textcolor{varecolor}{\textbf{Critical Apparatus:}}
\begin{variants}
\item[64.1] \dev{द्वियन्ते} (G) \dev{for} (\dev{द्वन्द्वान्ते} / \dev{द्वियन्ते दाकुति}).
\item[64.2] \dev{चतुष्पदं} (G,K,Kh) \dev{for} (\dev{चतुष्पद}).
\item[64.3] \dev{तिथ्यर्धं ते नागं} (G) \dev{for} (\dev{तिथ्यर्धं ते नागं}).
\item[66.1] \dev{तिथिन्यत्} \dev{for} (\dev{तिथिवदन्यत्}).
\item[\dev{वि०}] \dev{इसकी श्लोक संख्या ६४ है} (the verse number is 64).
\end{variants}

\vspace{0.5cm}
\hrule
\vspace{0.3cm}
\textit{[At this point the Triprasnādhikāra concludes. Several verses in the manuscripts
are mutilated; Dvivedī notes that at some places only the commentary (ṭīkā) is available.
The chapter contained 66 verses.]}
\vspace{0.3cm}
\hrule
\vspace{0.5cm}

\newpage


%% ========================================================================
%% CHAPTER 4: CANDRAGRAHANADHIKARA (Lunar Eclipse)
%% ========================================================================
\section*{Chapter IV: Candragrahaṇādhikāya}
\addcontentsline{toc}{section}{Chapter IV: Candragrahaṇādhikāya (Lunar Eclipse)}
\markboth{Chapter IV: Candragrahaṇādhikāya}{}

\begin{center}
\dev{चतुर्थोऽध्यायः}\\[0.2cm]
\textit{Chapter on the Computation of Lunar Eclipses}
\end{center}

This chapter contains rules for computing lunar eclipses, including the dimensions
of the Earth's shadow at the Moon's orbit, the times of first and last contacts,
the phases of the eclipse, and the colours assumed by the eclipsed Moon.

\vspace{0.3cm}
\columnheaders

\subsection*{Verses 1--4: Shadow dimensions and eclipse elements}

\noindent
\sanskritcol{\dev{सूर्यज्या दिनभागज्यया गुणाक्षज्ययाथवा भक्ता} \\ \dev{या द्वादशगुणिता विषुवच्छाया विभक्ता वा ॥ १ ॥} \\[0.2cm]
\dev{द्वादश विषुवच्छाया गृह ते पृथगक्षलम्बजीवे वा} \\ \dev{क्रान्तिहते सममण्डलकर्णः प्राग्वत् पृथक् छाया ॥ ३ ॥}}
\englishcol{\textbf{Translation:} (1) The Sun's jyā multiplied by the day-part-jyā, divided
by the guṇākṣa-jyā or alternatively, the viṣuvacchāyā divided by twelve. (3) Twelve
times the viṣuvacchāyā gives the shadow; or separately by the aksa-lamba-jyā.
Multiplied by the declination (krānti), the sama-maṇḍala-karṇa; as before, the
separate shadow.}

\vspace{0.4cm}
\textcolor{varecolor}{\textbf{Critical Apparatus:}}
\begin{variants}
\item[1.1] \dev{ज्ययाथवा} (Ch) \dev{for} (\dev{ज्ययाथवा}).
\item[1.2] \dev{विषुवच्छाया} (G) \dev{for} (\dev{विषुवच्छाया}).
\item[3.1] \dev{छाये} \dev{for} (\dev{छाया}).
\item[3.2] \dev{पृथगक्ष} \dev{for} (\dev{पृथगक्ष}).
\item[3.3] \dev{कर्णः} \dev{for} (\dev{कर्ण}).
\end{variants}

\subsection*{Verses 7--8: Eclipse magnitudes and durations}

\noindent
\sanskritcol{\dev{द्विकुलंबछायाक्षज्ञा तद्गंसंयुते मूलम्} \\ \dev{विषुवत्कर्णछाया कर्णोन्यदा शङ्कुः ॥ ७ ॥} \\[0.2cm]
\dev{उन्नतजीवाकोटिः छाया हर्ज्ञा भुजी नतज्ञा वा} \\ \dev{कर्णछायावृत्तं व्यासाद्ध दयमतोऽन्यत्र ॥ ८ ॥}}
\englishcol{\textbf{Translation:} (7) The two aksa-lamba-chāyā-jyā, their sum or difference gives
the root; the viṣuvat-karṇa-chāyā, otherwise the hypotenuse is the śaṅku (gnomon).
(8) The uccatana-jīvā-koṭi, the chāyā, the ha-jyā, the bhuja or the nata-jyā;
the hypotenuse-chāyā-circle from the semi-diameter; similarly elsewhere.}

\vspace{0.4cm}
\textcolor{varecolor}{\textbf{Critical Apparatus:}}
\begin{variants}
\item[7.1] \dev{क्षया} (G) \dev{क्षजय} (Kh) \dev{वज्ञा} \dev{for} (\dev{क्षज्ञा}).
\item[7.2] \dev{मूलम्} (Ch) \dev{for} (\dev{मूलम्}).
\item[7.3] \dev{विषुवत्} (G) \dev{विषुवति} \dev{for} (\dev{विषुवत्}).
\item[7.4] \dev{कर्णोन्यदा} (G,K) \dev{for} (\dev{कर्णोत्यदा}).
\item[7.5] \dev{शङ्कुलम्बः} (K) \dev{शङ्कुलबश} (Kh) \dev{for} (\dev{शङ्कुलब}).
\item[8.1] \dev{कोटिछाया} \dev{for} (\dev{कोटिः छाया}).
\item[8.2] \dev{उन्नतजीवाकोटिः} \dev{for} (\dev{उन्नतजीवाकोटिः}).
\item[8.3] \dev{हर्ज्ञा} (K) \dev{शध्या} \dev{for} (\dev{हर्ज्ञा}).
\item[8.4] \dev{भुजानतज्ञावा} \dev{for} (\dev{भुजोनतज्ञावा}).
\item[\dev{वि०}] \dev{इसकी श्लोक संख्या ७ है} (this verse is numbered 7 in the text).
\end{variants}

\subsection*{Verses 14--18: Sphuṭa-tithi and lambana}

\noindent
\sanskritcol{\dev{यद्यधिकं स्थित्यर्धं तदन्तरेणान्यथोनसूनमेकम्} \\ \dev{ऋणमधिकं तदैक्येनाधिकमेवं विमर्दार्धः ॥ १४ ॥} \\[0.2cm]
\dev{स्फुटतिथ्यन्ताल्लम्बनमसकृतस्थित्यर्धहीनयुक्ताद्वा} \\ \dev{तत्स्फुटं विश्लेषात् स्थिस्यर्धोतयुततिथ्यन्ते ॥ १६ ॥} \\[0.2cm]
\dev{तत्स्फुटतिथिछेदान्तरे स्फुटे तिथिदले विहीनयुतात्} \\ \dev{स्वविमर्दोदयाद् एवं स्पष्टो विमर्दार्धः ॥ १७ ॥} \\[0.2cm]
\dev{शशिवद्विभुः स्फुटविक्षेपकृतं स्थितिदलेन संगुणिता} \\ \dev{स्पष्टः स्थित्यर्धहतो भवति भुजः पूर्ववच्छेषमूलः ॥ १८ ॥}}
\englishcol{\textbf{Translation:} (14) If the half-duration is greater, by the difference,
otherwise minus one; if negative with addition, thus the half-obscuration (vimarda).
(16) From the sphuṭa-tithi-antā the parallax (lambana), with or without the
asakṛt-sthityardha; that sphuṭa is by separation, at the end of sthityardha-utaya-tithi.
(17) In the difference of that sphuṭa-tithi-cheda, the sphuṭa in the tithi-half,
with or without, from the self-vimarda-udaya, thus the corrected half-obscuration.
(18) The Moon's latitude (sphuṭa-vikṣepa) multiplied by the half-duration, the
corrected (sphuṭaḥ) bhūja multiplied by the half-duration is the root of the
remainder as before.}

\vspace{0.4cm}
\textcolor{varecolor}{\textbf{Critical Apparatus:}}
\begin{variants}
\item[14.1] \dev{तदैक्येनाधिकमेवं} (G,Ch) \dev{for} (\dev{तदैक्येनाधिकमेवं}).
\item[16.1] \dev{स्थिस्यर्धोतयुत} \dev{for} (\dev{स्थितिस्यर्धोतयुततिथ्यन्ते}).
\item[17.1] \dev{तत्स्फुटतिथिछेदान्तरे} \dev{for} (\dev{तत्स्फुटतिथिछेदान्तरे}).
\item[18.1] \dev{भुजः} \dev{for} (\dev{भुजः}).
\item[18.2] \dev{पूर्ववच्छेषमूलः} \dev{for} (\dev{पूर्ववच्छेषमूलः}).
\end{variants}

\subsection*{Verses 16--20: Eclipse colours and phases}

\noindent
\sanskritcol{\dev{क्षयान्त्ययोः सधुमूः कृष्णः खण्डग्रहे ऽर्धान्तोऽधिके ग्रासः} \\ \dev{सकृष्ण (षण्ण) ताम्रः स्वकरे शक्लिकपीलवर्णः ॥ १६ ॥} \\[0.2cm]
\dev{मानविमलं स्थितिदलवलनेष्ट ग्रास समकलादेषु} \\ \dev{नदगूहराध्याय विस्तृतिराया इष्टचतुर्थोयाम् ॥ २० ॥}}
\englishcol{\textbf{Translation:} (16) At the end of obscuration, smoke-coloured, black; in partial
eclipse, half-increased eating (grāsa); black-brown (sa-kṛṣṇa), coppery (tāmra),
white-yellowish (śakli-kapīla-varṇa). (20) The measure of pure, half-duration-motion-desired,
eating in the samakala-etc.; the breadth of the nadagūhara-āyāma desired one-fourth.}

\vspace{0.4cm}
\textcolor{varecolor}{\textbf{Critical Apparatus:}}
\begin{variants}
\item[16.1] \dev{क्षयान्त्ययोः} (G,K) \dev{for} (\dev{क्षयान्त्ययोः}).
\item[16.2] \dev{खण्डग्रहे} (G) \dev{for} (\dev{खण्डग्रहे}).
\item[16.3] \dev{शशी} (G,K,Ch) \dev{for} (\dev{शशि}).
\item[16.4] \dev{कपिलां} (G) \dev{कपिलवर्णः} (K) \dev{for} (\dev{कपिलवर्णः}).
\item[20.1] \dev{दलनेष्ट} \dev{for} (\dev{दलवलनेष्ट}).
\item[\dev{वि०}] \dev{``इति'' से ``अध्याय'' तक श्रंकित नहीं} (the closing formula is not marked).
\end{variants}

\vspace{0.5cm}
\chapterend
\begin{center}
\dev{\textbf{इति ब्रह्मगुप्ते चतुर्थोऽध्यायः}}\\[0.2cm]
\textit{Thus ends the Fourth Chapter of the Brāhmasphuṭasiddhānta.}
\end{center}
\chapterend

\newpage

%% ========================================================================
%% CHAPTER 5: SURYAGRAHANADHIKARA (Solar Eclipse)
%% ========================================================================
\section*{Chapter V: Sūryagrahaṇādhikāya}
\addcontentsline{toc}{section}{Chapter V: Sūryagrahaṇādhikāya (Solar Eclipse)}
\markboth{Chapter V: Sūryagrahaṇādhikāya}{}

\begin{center}
\dev{पञ्चमोऽध्यायः}\\[0.2cm]
\textit{Chapter on the Computation of Solar Eclipses}
\end{center}

This chapter treats the computation of solar eclipses, including the parallax
corrections (\textit{lambana}) that distinguish solar eclipse computation from lunar
eclipse computation.

\vspace{0.3cm}
\columnheaders

\subsection*{Verses 2--4: Lambana and eclipse elements}

\noindent
\sanskritcol{\dev{सूर्यज्या दिनभागज्यया गुणाक्षज्ययाथवा भक्ता} \\ \dev{या द्वादशगुणिता विषुवच्छाया विभक्ता वा ॥ २ ॥} \\[0.2cm]
\dev{द्वादश विषुवच्छाया गृहे ते पृथगक्षलम्बजीवे वा} \\ \dev{क्रान्तिहते सममण्डलकर्णः प्राग्वत् पृथक् छाया ॥ ३ ॥} \\[0.2cm]
\dev{अर्काग्रा दर्गोनं भुज्या वर्गाद् भ\-346450 सके १४४ कृतिगुणितम्} \\ \dev{आद्यान्योभ्रा द्वादशविषुवच्छाया वधो हतयोः ॥ ४ ॥}}
\englishcol{\textbf{Translation:} (2--3) Same as lunar eclipse --- the shadow dimensions are computed
from the Sun's declination. (4) The bhuja from the Sun's tip (arkāgra), the square
deducted, 346450 and 144 multiplied by the square; the first and other, twelve times
the viṣuvacchāyā, the product of the two.}

\subsection*{Verses 5--8: Lambana-ghaṭikā and eclipse timing}

\noindent
\sanskritcol{\dev{लम्बनघटिकालग्नात् तात्कालिकात् त्रिराशियुनात्} \\ \dev{ऋण (ऋणां) धिके शकृन् हीनं धनमसकृत्पञ्च हृश्यन्ते ॥ ५ ॥} \\[0.2cm]
\dev{करणगुणात् व्यासाद्ध द्विसुवेद ४८ विभाजिता फलविभक्ता} \\ \dev{लम्बननाड्यो भास्करवित्त्रिभलग्नान्तरं वा ॥ ६ ॥} \\[0.2cm]
\dev{यदि लम्बनमृणात् धनमधिकाः स्वफललिप्ताभिः} \\ \dev{अक्षज्ञाया वित्त्रिभलग्नात् स्वक्रान्तिरुत्तराकस्य} \\[0.1cm]
\dev{इन्दोर्वा यद्यधिका न नतिः सौम्यान्यथा याम्या ॥ ८ ॥}}
\englishcol{\textbf{Translation:} (5) From the lambana-ghaṭikā-lagna, from the tātkālika,
diminished by three signs; if negative with addition, five times the asakṛt,
divided by hṛdaya. (6) From the karaṇa-guṇa, from the semi-diameter, dviveda 48,
divided; the lambana-nāḍīs or the bāskara-vitribha-lagna-antaram. (7) If the lambana
is negative, the positive result in its own lalipta-phala. (8) From the aṣajyā,
from the vitribha-lagna, its own declination (krānti) uttarākṣa or of the Moon;
if greater, otherwise the yāmyā.}

\vspace{0.4cm}
\textcolor{varecolor}{\textbf{Critical Apparatus:}}
\begin{variants}
\item[5.1] \dev{ऋणमधिके} (G) \dev{for} (\dev{ऋणमधिके}).
\item[5.2] \dev{मसकृत्पञ्च} (K) \dev{for} (\dev{मसकृत्पञ्च}).
\item[5.3] \dev{हृश्यन्ते} (K) \dev{for} (\dev{हृश्यन्ते}).
\item[6.1] \dev{विभाजिता फलविभक्ता} (G) \dev{for} (\dev{विभाजिता फलविभक्ता}).
\item[6.2] \dev{लम्बननाड्यो} \dev{for} (\dev{लम्बननाड्यो}).
\item[8.1] \dev{वा} (G) \dev{ज्यावा} (K) \dev{for} (\dev{वा}).
\end{variants}

\subsection*{Verses 14--17: Sphuṭa-tithi and corrected obscuration}

\noindent
\sanskritcol{\dev{यद्यधिकं स्थित्यर्धं तदन्तरेणान्यथोनसूनमेकम्} \\ \dev{ऋणमधिकं तदैक्येनाधिकमेवं विमर्दार्धः ॥ १४ ॥} \\[0.2cm]
\dev{स्फुटतिथ्यन्ताल्लम्बनमसकृतस्थित्यर्धहीनयुक्ताद्वा} \\ \dev{तत्स्फुटं विश्लेषात् तिस्थिस्यर्धोतयुततिथ्यन्ते ॥ १६ ॥} \\[0.2cm]
\dev{तत्स्फुटतिथिछेदान्तरे स्फुटे तिथिदले विहीनयुतात्} \\ \dev{स्वविमर्दोदयाद् एवं स्पष्टो विमर्दार्धः ॥ १७ ॥}}
\englishcol{\textbf{Translation:} (14) If the half-duration is greater, then by the difference,
otherwise minus one; if negative with addition, thus the half-obscuration (vimarda-ardha).
(16) From the sphuṭa-tithi-anta, the lambana, with or without the asakṛt-sthityardha;
that sphuṭa is by separation, at the end of the sthityardha-utaya-tithi. (17) In the
difference of that sphuṭa-tithi-cheda, the sphuṭa in the tithi-half, with or without;
from the self-vimarda-udaya, thus the corrected (spaṣṭa) half-obscuration.}

\vspace{0.4cm}
\textcolor{varecolor}{\textbf{Critical Apparatus:}}
\begin{variants}
\item[14.1] \dev{तदैक्येनाधिकमेवं} (G,Ch) \dev{for} (\dev{तदैक्येनाधिकमेवं}).
\item[15.1] \dev{तदैक्येनाधिकमेवं} (G,Ch) \dev{for} (\dev{तदैक्येनाधिकमेवं}).
\item[16.1] \dev{तिस्थिस्यर्धोतयुत} \dev{for} (\dev{स्थितिस्यर्धोतयुततिथ्यन्ते}).
\item[17.1] \dev{तत्स्फुटतिथिछेदान्तरे} \dev{for} (\dev{तत्स्फुटतिथिछेदान्तरे}).
\end{variants}

\subsection*{Verses 18--20: Deflection and sphuṭa computation}

\noindent
\sanskritcol{\dev{एवं मानैक्यार्धादधिके सध्यान्तरे युतिग्रहयोः} \\ \dev{स्थितिस्यर्धविमर्दर्धले हीनं तरा ग्रहेंदुयुतौ ॥ १८ ॥} \\[0.2cm]
\dev{लम्बनमकलग्रहवदसकृत्स्वावनतिलिप्तिकास्पष्टी} \\ \dev{तात्कालिकविक्षेपी तदन्तरैक्यं समान्यदिशोः ॥ १९ ॥} \\[0.2cm]
\dev{विक्षेपो मध्यान्तरमुध्वस्वछादको ग्रहोऽध्वस्थः} \\ \dev{मानैक्यार्धादधिके नातिस्पष्टास्फुटोक्तिरतः ॥ २० ॥}}
\englishcol{\textbf{Translation:} (18) Thus in the conjunction of planet and Moon, if greater
than the measure of unity plus half, in the middle-intermission; diminished by the
half-duration-obscuration, the cross (tārā). (19) The lambana like akala-graha,
asakṛt-svāvanati-lāptikā-sphuṭī, the tātkālika-vikṣepī, the samānya-diśoḥ sum or
difference. (20) The deflection (vikṣepa), the middle-difference, the above-shadow-planet
situated below; if greater than half the unity-measure, not too distinct, hence
the sphuṭa-statement.}

\vspace{0.4cm}
\textcolor{varecolor}{\textbf{Critical Apparatus:}}
\begin{variants}
\item[18.1] \dev{युतिग्रहयोः} \dev{for} (\dev{युतिग्रहयोः}).
\item[18.2] \dev{विमर्दर्धले} \dev{for} (\dev{विमर्दे}).
\item[19.1] \dev{तदन्तरैक्यं} \dev{for} (\dev{तदन्तरैक्यं}).
\item[19.2] \dev{लिप्ता} \dev{for} (\dev{लिप्तिका}).
\item[20.1] \dev{मर्के} \dev{for} (\dev{मर्कट}).
\item[20.2] \dev{मुध्वस्य छादको ग्रहोऽध्वस्थः} (G) \dev{for} (\dev{मुध्वं स्वछादको ग्रहोऽध्वस्थः}).
\end{variants}

\vspace{0.5cm}
\chapterend
\begin{center}
\dev{\textbf{इति ब्रह्मगुप्तसवर्गोऽध्यायः}}\\[0.2cm]
\textit{Thus ends the Fifth Chapter of the Brāhmasphuṭasiddhānta.}
\end{center}
\chapterend

\newpage


%% ========================================================================
%% CHAPTERS 6--9: Udayasta, Candrasrngonnati, Candracchaya, Grahayuti
%% ========================================================================
\section*{Chapters VI--IX: Udayāsta, Candraśṛṅgonnati, Candraccāyā, Grahayuti}
\addcontentsline{toc}{section}{Chapters VI--IX: Supplementary Astronomical Topics}
\markboth{Chapters VI--IX: Supplementary Topics}{}

These chapters cover miscellaneous astronomical computations: rising and setting
of planets, the elevation of the Moon's horns, the Moon's shadow, and conjunctions
of planets.

\vspace{0.3cm}
\columnheaders

\subsection*{Chapter VI: Udayastādhikāra (Rising and Setting of Planets)}

\noindent
\sanskritcol{\dev{द्वादशाभिः शितांशुसितजीवक्षमिभूमिजा नवभिः} \\ \dev{द्युत्तरवृद्धिरन्तरघटिका षड्गुणिता कलांशैः ॥ ६ ॥} \\[0.2cm]
\dev{हृश्याहृश्यायुतिवद्ग्रहाकं सुत्तचान्तरलब्धदिनैः} \\ \dev{ऊनाधिकलिप्ताम्यः प्राग्वत् तात्कालिकैरसकृत ॥ ७ ॥} \\[0.2cm]
\dev{प्रतिदिनमुदयास्तमयाद् एवं तात्कालिकं ग्रहविलग्नैः} \\ \dev{सूर्यास्तमयोदयिकः शीतांशुः पौष्णमास्यन्ते ॥ ८ ॥}}
\englishcol{\textbf{Translation:} (6) By twelve, the śītāṃśu-sita-jīvā, the earth-born, by
nine; the day-increment-difference-ghaṭikā multiplied by six in kalāṃśa. (7) The
graheka like hṛśyā-hṛśyā-yuti, by the difference of suttacāntara-obtained-days;
subtracted or added in liptāmyaḥ, as before by the tātkālika-asakṛt. (8) Daily
rising and setting thus, by the tātkālika-graha-vilagnais; the solar rising and
setting at the end of the Pausa month, the Moon (śītāṃśu).}

\vspace{0.4cm}
\textcolor{varecolor}{\textbf{Critical Apparatus:}}
\begin{variants}
\item[6.1] \dev{शितांशुः} (G) \dev{for} (\dev{शीतांशुः}).
\item[6.2] \dev{भूमिजानवभिः} \dev{for} (\dev{भूमिजानवभिः}).
\item[7.1] \dev{सुत्तचान्तरलब्धदिनैः} \dev{for} (\dev{सुत्तचान्तरलब्धदिनैः}).
\item[7.2] \dev{ऊनाधिक} \dev{for} (\dev{ऊनाधिक}).
\item[8.1] \dev{शीतांशोः} (G) \dev{शीतांशोः} (K) \dev{for} (\dev{शीतांशौ}).
\item[8.2] \dev{पौष्णमास्यां} (G) \dev{पूर्णिमास्था} \dev{for} (\dev{पौष्णमास्यन्ते}).
\item[8.3] \dev{दयिकः} (K) \dev{for} (\dev{दयिकः}).
\end{variants}

\subsection*{Chapter VII: Candraśṛṅgonnatyādhikāya (Elevation of the Moon's Horns)}

\noindent
\sanskritcol{\dev{कोटिश्रवणाज्ञानात् शशिना श्युद्लोन्नतिर्विसंवदति} \\ \dev{उदयास्तमयो दिनकृतः प्रतिघटिकं माती च वाज्ञानात् ॥ ३८ ॥} \\[0.2cm]
\dev{श्रकैतरघटिका व्यस्तज्ञा चन्द्रमानगुणिता यतः} \\ \dev{व्यास विभक्ता शुद्धं यतो न दृक् सममसत्तस्मात् ॥ ३६ ॥} \\[0.2cm]
\dev{प्राक् गुदितोऽन्यधिकः पश्चादुदिनकोऽपरेव्यस्तः} \\ \dev{कालो यः छाया तदसत्स्फुटं सुक्तिमान् प्राक् ॥ ४० ॥} \\[0.2cm]
\dev{उदितामुदितास्तमितावशेषकालान्नं वेत्ति यः स कथम्} \\ \dev{आर्यमटज्ञः शदिनः छाया श्युद्लोन्नति वेत्ति ॥ ४१ ॥}}
\englishcol{\textbf{Translation:} (38) The latitude (unnati) of the Moon is found to disagree
from the hypotenuse-śravaṇa-jñānāt; the rising and setting are caused by day,
by prati-ghaṭika measure and from vā-jñānāt. (36) The akaicara-ghaṭikā, vyasta-jyā,
multiplied by candra-māna, divided by vyāsa is pure; since the dṛk-sama is not
true. (40) Previously taught as different, afterwards the other reverse; that
time which is shadow, that asat-sphuṭa, suktimān before. (41) Who knows the
udita-udita-astamita-avaśeṣa-kāla; how can he who knows the shadow-unnati of
the Moon be an ārya-maṭa-jña.}

\vspace{0.4cm}
\textcolor{varecolor}{\textbf{Critical Apparatus:}}
\begin{variants}
\item[38.1] \dev{श्युद्लोन्नतिः} \dev{for} (\dev{श्युद्लोन्नतिः}).
\item[38.2] \dev{विसंवदति} \dev{for} (\dev{विसंवदति}).
\item[38.3] \dev{दिनकृत्तः} \dev{for} (\dev{दिनकृतः}).
\item[36.1] \dev{दृक् सम} (G) \dev{for} (\dev{दृक् सम}).
\item[40.1] \dev{प्राक् गुदितोऽन्यधिकः} \dev{for} (\dev{प्रागुदितोऽन्यधिकः}).
\item[41.1] \dev{शदिनः} \dev{for} (\dev{दिनः}).
\item[41.2] \dev{आर्यमटज्ञः} \dev{for} (\dev{आर्यमटज्ञः}).
\end{variants}

\subsection*{Chapter VIII: Candraccāyādhikāya (The Moon's Shadow)}

\noindent
\sanskritcol{\dev{क्रान्तिज्ञा तत्क्रान्तिविक्षेपक्रान्तिचापानाम्} \\ \dev{संयोगान्तरजीवा स्वक्रान्तिज्यैकभिन्नदिशाम् ॥ १४ ॥} \\[0.2cm]
\dev{एवं भमुनिश्नुवयो ग्रहतत्क्रान्त्या चन्द्राष्टकं कर्माद्य} \\ \dev{कृत्वा गृहे भमुनिवत् तस्मात्स्वक्रान्तिजीवा वा ॥ १६ ॥}}
\englishcol{\textbf{Translation:} (14) The declination-jyā, the deflection (vikṣepa), the
declination-cāpa: the sum or difference of the jyā of like or unlike directions
of the declination. (16) Thus the bhūmi-niśa-vayo (latitude etc.) of the planet,
by its declination, the candra-aṣṭaka-karma etc.; having done this in the eclipse,
as in bhūmi, from that the svakrantijīvā.}

\vspace{0.4cm}
\textcolor{varecolor}{\textbf{Critical Apparatus:}}
\begin{variants}
\item[14.1] \dev{चापभागानाम्} (G) \dev{for} (\dev{चापानाम्}).
\item[14.2] \dev{विक्षेप} \dev{for} (\dev{विक्षेप}).
\item[14.3] \dev{चैकभिन्नदिशाम्} \dev{for} (\dev{ज्यैकभिन्नदिशाम्}).
\item[16.1] \dev{श्लुवयोग्रह} \dev{for} (\dev{श्लुवयोग्रह}).
\item[16.2] \dev{च दृष्टिकर्माच} (G) \dev{for} (\dev{चन्द्राष्टकं कर्माद्}).
\item[16.3] \dev{भमुनिश्नुवका} \dev{for} (\dev{भमुनिश्नुवयो}).
\item[\dev{वि०}] An additional verse is found in manuscript G:
\dev{भमुनिश्नुवका प्राक् प्रदर्शिता त्विषा} /
\dev{देषामेवकार्यं ग्रहस्य पुनस्तत् क्रान्ते प्रथमं वृत्तकरणं} /
\dev{कृत्वा पश्चात् तज्ज्ञाति स्वक मुनिध्रुवकवत्}.
\end{variants}

\subsection*{Chapter IX: Grahayutyādhikāya (Conjunction of Planets)}

\noindent
\sanskritcol{\dev{प्रतिघटिकमधिकशङ्कोभ्यां मध्ये दर्शयेत् भुवा} \\ \dev{त्रिप्रश्नोक्तचा रविवद्धंकु भ्रमणादिक्रमशेषम् ॥ ४४ ॥} \\[0.2cm]
\dev{शङ्कु प्राच्यपरान्तरं विषुवच्छायैक्य सुत्तरे नृतले} \\ \dev{याम्योतरं गुणाहतं स्वक्रान्तिलेशैः कर्णास्यास् ॥ ४६ ॥} \\[0.2cm]
\dev{तच्चापांशाः सदयेः संघुवकापक्रमांशरूना} \\ \dev{विक्षेपांशे व्यस्ता व्यस्तविशुद्धा व्यसहशांशः ॥ ४७ ॥}}
\englishcol{\textbf{Translation:} (44) By the prati-ghaṭika and adhika-śaṅku, in the middle
is shown the earth; the tripraśna-oktacā like the Sun's śaṅku, the remainder of
the bhramaṇādi-krama. (46) The śaṅku, the eastern-western antaram, the viṣuvacchāyā-sum
or difference, nṛtale; the yāmyottaraṃ guṇāhataṃ by svakrānti-leśaiḥ and karṇa.
(47) The cāpāṃśāḥ of sada-yeḥ, saṅghuvakāpakramāṃśa-rūnāḥ; the vikṣepāṃśe vyastā,
the vyasta-viśuddhā vya-sahaśāṃśaḥ.}

\vspace{0.4cm}
\textcolor{varecolor}{\textbf{Critical Apparatus:}}
\begin{variants}
\item[44.1] \dev{दर्शयेन्} (G) \dev{for} (\dev{दर्शयेतु}).
\item[44.2] \dev{रविवद्धंकु} (Ch) \dev{for} (\dev{रविवद्धंकु}).
\item[46.1] \dev{परान्तर} (G) \dev{for} (\dev{परान्तरं}).
\item[46.2] \dev{नृतरे} \dev{for} (\dev{नृतले}).
\item[46.3] \dev{याम्बैतर} (G) \dev{for} (\dev{याम्योत्तरं}).
\item[46.4] \dev{गुणहतं} \dev{for} (\dev{गुणाहतं}).
\item[47.1] \dev{चापांशाः} (G,Ch) \dev{for} (\dev{चापांशाः}).
\item[47.2] \dev{सहदयेः} (G) \dev{for} (\dev{सदयेः}).
\end{variants}

\vspace{0.5cm}
\chapterend
\textit{[The chapters on Udayāsta, Candraśṛṅgonnati, Candraccāyā, and Grahayuti
conclude here, with 12--70 verses each.]}
\chapterend

\newpage


%% ========================================================================
%% CHAPTER 11: TANTRAPARIKSHADHIKAYA (Critical Examination)
%% ========================================================================
\section*{Chapter XI: Tantraparīkṣādhikāya (Dūṣaṇādhikāya)}
\addcontentsline{toc}{section}{Chapter XI: Tantraparīkṣādhikāya (Critical Examination)}
\markboth{Chapter XI: Tantraparīkṣādhikāya}{}

\begin{center}
\dev{एकादशोऽध्यायः}\\[0.2cm]
\dev{अथ तन्त्रपरीक्षा नामैकादशोऽध्यायः}\\[0.3cm]
\textit{Chapter on the Critical Examination of Astronomical Systems}
\end{center}

This celebrated chapter contains Brahmagupta's critique of various astronomical
systems current in his time. He refutes the doctrines of Āryabhaṭa, Śrīṣeṇa,
Viṣṇucandra, Lāṭadeva, Vijayanandi, and Pradyumna, as well as the foreign systems
derived from Greece and Babylon (the Romaka and Puliśa Siddhāntas).

\vspace{0.3cm}
\columnheaders

\subsection*{Verses 1--2: Introduction to the critique}

\noindent
\sanskritcol{\dev{ये ज्ञान पटलरुद्धद्वोन्यब्रह्मविदन्ति सिद्धान्तम्} \\ \dev{तेषां युगादिभेदैर् दोषास्तान् प्रवक्ष्यामि ॥ १ ॥} \\[0.2cm]
\dev{युगमाहुः पञ्चाब्दं रविशशिनोः संहितां कारये} \\ \dev{अधिमासा वमरात्र स्फुटतिथ्यज्ञान तदसत् ॥ २ ॥}}
\englishcol{\textbf{Translation:} (1) Those who know the brahma-vid (science) in a chapter-barred
dual manner, they consider it a siddhānta; I shall declare their faults arising from
the yuga-ādi differences. (2) They say the yuga is five years; making the Saṃhitā
of the Sun and Moon; the adhimāsa, vamara-atra, sphuṭa-tithi-jñāna --- that is false.}

\vspace{0.4cm}
\textcolor{varecolor}{\textbf{Critical Apparatus:}}
\begin{variants}
\item[1.1] \dev{येज्ञान} \dev{for} (\dev{ये ज्ञान}).
\item[1.2] \dev{रुद्धशो} (G) \dev{रुद्धो} \dev{for} (\dev{रुद्धद्वो}).
\item[1.3] \dev{न्यद्बाह्यविदन्ति} (G) \dev{न्यब्रह्मविदन्ति} \dev{for} (\dev{न्यब्रह्मविदन्ति}).
\item[1.4] \dev{सिद्धान्तात्} (G,Ch) \dev{for} (\dev{सिद्धान्तम्}).
\item[1.5] \dev{भेदाद्ये} (G,Ch) \dev{for} (\dev{भेदैः}).
\item[1.6] \dev{प्रवक्ष्यामि} (G) \dev{प्रविक्ष्यामि} \dev{for} (\dev{प्रवक्षामि}).
\item[2.1] \dev{पञ्चाक्षरं} \dev{for} (\dev{पञ्चाब्दं}).
\item[2.2] \dev{ज्ञानतस्तदसत्} (G) \dev{for} (\dev{ज्ञानतदसतु}).
\item[2.3] \dev{कारये} \dev{for} (\dev{कारये}).
\item[\dev{वि०--अतिरिक्त पाठः}] \dev{अथ ब्रह्मगुप्तसिद्धान्ते एकादशमोऽध्यायः प्रारभ्यते} /
\dev{दुवृत्तेन्द्वयो गतगम्याल्पात्मकशून्य यमल} /
\dev{सायक हतं कलाभिरूनौ सदाकेन्द्र} / \dev{विधु वजीवेद्वि २ हतं चन्द्रोच्चे} /
\dev{तिथि १५ गुण च सितशीघ्र} /
\dev{त्वीषु ५२ हतं च स्वार्द्वि २ कु १ वेद} /
\dev{४ हतं च पात कुजशनिषु} /
\dev{ग्रहबीजानि ब्रह्मसिद्धान्ते}.
\end{variants}

\subsection*{Verses 6--9: Critique of planetary positions}

\noindent
\sanskritcol{\dev{युगवर्षादिनवद्व्येक्षं सितादेः समं प्रवृत्तान्यत्} \\ \dev{तदसत् यतः स्फुटयुगं तस्थैर्यान्मन्दपातानाम् ॥ ६ ॥} \\[0.2cm]
\dev{ग्रहमुक्त रूनाया मन्दोच्चं भवति शीघ्रमधिकायाम्} \\ \dev{उच्चगतौ मन्दोच्चं न विना भुक्तिं व्रजेमतः ॥ ७ ॥} \\[0.2cm]
\dev{आर्याष्टशते पाता भ्रमंति दशगीतके स्थिराः पठिताः} \\ \dev{मुक्तं द्विपातमण्डपमण्डले भ्रमति स्थिरान्नतः ॥ ८ ॥} \\[0.2cm]
\dev{आर्यभटो जानाति ग्रहाष्टात् यदुक्तवांस्तदसत्} \\ \dev{राहुकृतं न ग्रहणं तत्पातोऽनाष्टमो राहुः ॥ ९ ॥}}
\englishcol{\textbf{Translation:} (6) The yuga-year-navadvyekṣa (nine), the bright (Moon) etc.,
the same pravṛttānyat; that is false, since the corrected yuga-position of the
manda-nodes is fixed. (7) The planet's release is in the ruṇāyā, the mandocca
increases in śīghra; in the high position, the mandocca does not exist without
(its own) motion. (8) In Ārya's eight hundred, the nodes move, in the Daśagītaka
are fixed; the released two-nodes move in the circle, not fixed. (9) Āryabhaṭa
knows the planets; what he has said is false; the Rāhu does not cause the eclipse,
its pāta is not the eighth Rāhu.}

\vspace{0.4cm}
\textcolor{varecolor}{\textbf{Critical Apparatus:}}
\begin{variants}
\item[6.1] \dev{युगवर्षादी} \dev{for} (\dev{युगवर्षादि}).
\item[6.2] \dev{सप्रवृत्तानयत्} \dev{for} (\dev{समं प्रवृत्तान्यत्}).
\item[6.3] \dev{तदराजु} (G) \dev{for} (\dev{तदसत्}).
\item[6.4] \dev{तत्तस्थैर्यात्} (G) \dev{for} (\dev{तस्थैर्यात्}).
\item[7.1] \dev{रूनायां} \dev{for} (\dev{रूनाया}).
\item[7.2] \dev{मुक्तच न्हु} \dev{for} (\dev{मुक्त द्वि}).
\item[7.3] \dev{वजमेतः} \dev{for} (\dev{व्रजेमतः}).
\item[7.4] \dev{वि०} --- An additional half-verse is supplied:
\dev{तद्दरगम्यां दशामिः संघुशितास्यां सूक्ष्मे}.
\item[8.1] \dev{आर्याष्टशते} (Ch) \dev{for} (\dev{आर्याष्टशते}).
\item[8.2] \dev{दशगीतके} \dev{for} (\dev{दशगीतके}).
\item[8.3] \dev{मुक्तं द्वि} \dev{for} (\dev{मुक्तं द्वि}).
\item[8.4] \dev{भ्रमंति} (G) \dev{for} (\dev{भ्रमंति}).
\item[9.1] \dev{ग्रहाष्टात्} (G) \dev{for} (\dev{ग्रहाष्टात्}).
\item[9.2] \dev{तस्यातो} (G) \dev{for} (\dev{तत्पातः}).
\item[9.3] \dev{अनाष्टमः} \dev{for} (\dev{अनाष्टमः}).
\end{variants}

\subsection*{Verse 63: Conclusion}

\noindent
\sanskritcol{\dev{इत्थिकथिततन्त्राणाकान् पठतिरपि दूषणैः करोत्यज्ञानात्} \\ \dev{तन्त्र परीक्षायाणां प्रदिष्टिरिकादशोऽध्यायः ॥ ६३ ॥}}
\englishcol{\textbf{Translation:} (63) Thus the systems taught, even by reading, cause ignorance
by faults; the examination of systems --- this is the eleventh chapter.}

\vspace{0.4cm}
\textcolor{varecolor}{\textbf{Critical Apparatus:}}
\begin{variants}
\item[63.1] \dev{तन्त्रगणकान् अदिते} (G) \dev{तन्त्राणाकान् पठतिरपि} \dev{for} (\dev{तन्त्राणाकान् पठतिरपि}).
\item[63.2] \dev{परीक्षार्याणाम्} (G) \dev{for} (\dev{परीक्षायाणाम्}).
\item[63.3] \dev{इत्ति} \dev{for} (\dev{इत्ति}).
\item[63.4] \dev{करोत्यज्ञानात्} \dev{for} (\dev{करोत्यज्ञानात्}).
\item[63.5] \dev{वि०} --- \dev{`इति' से `अध्यायः' तक सब लुप्त} (the closing formula is omitted).
\end{variants}

\vspace{0.5cm}
\chapterend
\begin{center}
\dev{\textbf{इति श्रीब्रह्मगुप्तसिद्धान्ते एकादशोऽध्यायः}}\\[0.2cm]
\textit{Thus ends the Eleventh Chapter of the Brāhmasphuṭasiddhānta.}
\end{center}
\chapterend

\newpage

%% ========================================================================
%% CHAPTER 12: GANITADHYAYA (Mathematics)
%% ========================================================================
\section*{Chapter XII: Gaṇitādhyāya}
\addcontentsline{toc}{section}{Chapter XII: Gaṇitādhyāya (Mathematics)}
\markboth{Chapter XII: Gaṇitādhyāya}{}

\begin{center}
\dev{द्वादशोऽध्यायः}\\[0.2cm]
\dev{अथ गणिताध्याय द्वादशः}\\[0.3cm]
\textit{The Twelfth Chapter --- On Mathematics}
\end{center}

The Gaṇitādhyāya is the most celebrated mathematical chapter of the
\textit{Brāhmasphuṭasiddhānta}. It contains 66 verses divided into ten sections
(\textit{vyavahāra}s), covering:

\begin{enumerate}[label=\arabic*.]
  \item \dev{परिकर्म} --- Preliminary operations (fractions, proportions)
  \item \dev{मिश्रकव्यवहारः} --- Mixture problems (verses 1--16)
  \item \dev{श्रेढ़िव्यवहारः} --- Arithmetic and geometric series (verses 17--20)
  \item \dev{क्षेत्रव्यवहारः} --- Geometry of plane figures (verses 21--39)
  \item \dev{वृत्तक्षेत्रगणितम्} --- Circle geometry (verses 40--43)
  \item \dev{खातव्यवहारः} --- Excavation/volume problems (verses 44--46)
  \item \dev{चितिव्यवहारः} --- Pile problems (verses 47+)
  \item \dev{शाद्रकाकाचिकादिकरणसूत्रे} --- Shadow of a gnomon etc.
  \item \dev{राशिव्यवहारः} --- Rules for treating quantities (verses 50--51)
  \item \dev{छायाव्यवहारः} --- Shadow problems (verses 52--66)
\end{enumerate}

\vspace{0.3cm}
\columnheaders

\subsection*{Section 1: Parikarmma (Preliminary Operations)}

\noindent
\sanskritcol{\dev{परिक्रम्म विशति यः संकलिताद्यं पृथग् विजानाति} \\ \dev{परष्ठी च व्यवहारान् छायातान् भवति गणकः सः ॥ १ ॥}}
\englishcol{\textbf{Translation:} (1) He who knows separately the parikamma (operation), the
saṃkalita (summation) etc., and the twenty vyavahāras with chāyāta --- he becomes
a calculator (gaṇaka).}

\vspace{0.4cm}
\textcolor{varecolor}{\textbf{Critical Apparatus:}}
\begin{variants}
\item[1.1] \dev{विद्याति} (G) \dev{for} (\dev{विजानाति}).
\item[1.2] \dev{संकलिताद्यां} (G,Ch) \dev{for} (\dev{संकलिताद्यं}).
\item[1.3] \dev{पृथग्विजानाति} (G,Ch) \dev{for} (\dev{पृथग् विजानाति}).
\item[1.4] \dev{व्यवहां} (G) \dev{for} (\dev{व्यवहारान्}).
\item[1.5] \dev{छायांतान्} (G) \dev{for} (\dev{छायातान्}).
\item[1.6] \dev{परिक्रम्मं} \dev{for} (\dev{परिकर्म}).
\end{variants}

\noindent
\sanskritcol{\dev{विपरीत छेदगुणाः राश्योः छेदांशकाः समछेदाः} \\ \dev{संकलितेशा योज्या व्यवकलितेशान्तरं कार्यम् ॥ २ ॥} \\[0.2cm]
\dev{रूपार्णी छेदगुणान्यंशयुतानि द्वयोर्बहूनां वा} \\ \dev{प्रत्युत्पन्नो भवति छेदवधेनोद्धृतांशवधः ॥ ३ ॥} \\[0.2cm]
\dev{प्रत्युत्पन्नः परिवर्त्यं भागहारछेदांशो छेदसंगुणछेदः} \\ \dev{द्वयांशयुग् भाज्यस्य भागहारः सर्वार्णितयोः भागहारः ॥ ४ ॥}}
\englishcol{\textbf{Translation:} (2) The inverted (viparīta) and the divisor-multipliers of
the quantities, the divisor-parts (chedāṃśakāḥ) with common divisor; the saṃkalita
etc.\ are to be added, the difference is vyavakalita. (3) The rūpāraṇī (unity)
multiplied by the divisor, with parts added, of two or many; the pratyutpanna
(product) becomes the division by the divisor-product of the fraction-product.
(4) The pratyutpanna after conversion, the divisor-part of the divisor, the
divisor-times-divisor; the two-part-sum of the dividend, the divisor of all
is the divisor.}

\vspace{0.4cm}
\textcolor{varecolor}{\textbf{Critical Apparatus:}}
\begin{variants}
\item[2.1] \dev{संकलितेशा} (G,Ch) \dev{for} (\dev{संकलितेशा}).
\item[2.2] \dev{गुणा} \dev{for} (\dev{गुणाः}).
\item[2.3] \dev{राश्योः छेदांशकाः} \dev{for} (\dev{राश्योः छेदांशकाः}).
\item[3.1] \dev{वधेनोद्धृतांशवधः} (G) \dev{for} (\dev{वधेनोद्धृतांशवधः}).
\item[3.2] \dev{बहूनां} (G,Ch) \dev{for} (\dev{बहूनां}).
\item[4.1] \dev{संगुण} \dev{for} (\dev{संगुण}).
\item[4.2] \dev{भाययस्य} (Ch) \dev{for} (\dev{भाज्यस्य}).
\item[4.3] \dev{सर्वार्णितयोः} (G) \dev{for} (\dev{सर्वार्णितयोः}).
\item[\dev{वि०}] In this verse, the words ``\dev{प्रत्युत्पन्न}'' and ``\dev{भागहारः}''
are not marked in the text; the verse begins with the word ``\dev{परिवृत्य}''.
\end{variants}

\noindent
\sanskritcol{\dev{सर्वार्णितांशवर्गं छेदकृति विभाजितो भवति वेगः} \\ \dev{सर्वार्शितांशमूलं छेदपदेनोद्धृते मूलम् ॥ ५ ॥} \\[0.2cm]
\dev{वर्गमूलं स्थाप्योत्पद्यनोन्यकृतिस्त्रिगुणोत्तरसंगुणा च यत्प्रथमा} \\ \dev{नोत्तरकृतिरन्त्यगुणा त्रिगुणाच्चोत्तर घनः ॥ ६ ॥} \\[0.1cm]
\devbf{घनः समाप्तः} \\[0.1cm]
\dev{छेदाघना द्वितीयात् घनमूलकृतिस्त्रिसंगुणा सकृतिः} \\ \dev{शोध्या त्रिपूर्वगुणिता प्रथमाद् घनतो घनो मूलम् ॥ ७ ॥} \\[0.1cm]
\devbf{घनमूलं समाप्तः}}
\englishcol{\textbf{Translation:} (5) The square of all-parts divided by the square of the
divisor becomes the varga (square root); the root of all-śita-parts, the root
divided by the divisor-term. (6) Having set up the varga-mūla (square root), the utpadyanyakṛti,
the trigunottara-saguṇa, and the prathamā; the nottara-kṛti-antitya-guṇā, trigunā
and uttara is the ghana (cube). [Cube completed.] (7) From the second chedāghana,
the cube-root-square-times-trisaguṇa-sakṛti; subtracted, times the tripūrva,
from the prathama-ghana, the ghana-mūla. [Cube root completed.]}

\vspace{0.4cm}
\textcolor{varecolor}{\textbf{Critical Apparatus:}}
\begin{variants}
\item[6.1] \dev{स्थाप्योत्यघनोन्यकृतिः} \dev{for} (\dev{स्थाप्योत्पद्यनोन्यकृति}).
\item[6.2] \dev{न्यत्कृति} \dev{for} (\dev{न्यकृति}).
\item[6.3] \dev{उत्तरकृतिः} \dev{for} (\dev{उत्तरकृति}).
\item[6.4] \dev{अन्त्यगुणा} \dev{for} (\dev{अन्त्यगुणा}).
\item[7.1] \dev{घनमूल} \dev{for} (\dev{घनमूल}).
\item[7.2] \dev{संगुणशासकृतिः} (G) \dev{for} (\dev{संगुणा सकृतिः}).
\item[7.3] \dev{छिदोघना} \dev{for} (\dev{छेदाघना}).
\item[7.4] \dev{प्रथमाद् घनतो घनमूलं} \dev{for} (\dev{प्रथमाद् घनतो घनमूलं}).
\item[\dev{वि०}] The words ``\dev{घनः समाप्तः}'' and ``\dev{घनमूलं समाप्तः}''
(``cube completed'' / ``cube root completed'') are not marked in the source text.
\end{variants}

\newpage


%% ========================================================================
%% GANITADHYAYA continued: Mishraka, Sredhi, Kshetra, Chaya
%% ========================================================================
\subsection*{Section 2: Miśrakavyavahāra (Mixture Problems)}

\noindent
\sanskritcol{\dev{प्रक्षेपयोगहृतया लब्धा प्रक्षेपका गुणा लाभाः} \\ \dev{ऊनाधिकोत्तरार्धं सद् द्वितोनया स्वफलसूनं युतम् ॥ १६ ॥}}
\englishcol{\textbf{Translation:} (16) Divided by the prakṣepa-yoga (sum of contributions),
obtained, the prakṣepaka (contributor) multiplied by the gains; the ūnādhika-uttara-ardha
is true, the dvitonayā with its own phala-sum added.}

\vspace{0.4cm}
\textcolor{varecolor}{\textbf{Critical Apparatus:}}
\begin{variants}
\item[16.1] \dev{प्रक्षेपयोगहतया} (G,Ch) \dev{for} (\dev{प्रक्षेपयोगहतया}).
\item[16.2] \dev{लब्धा} \dev{for} (\dev{लब्ध्वा}).
\item[16.3] \dev{प्रक्षेपका} (G) \dev{प्रक्षेपक} \dev{for} (\dev{प्रक्षेपका}).
\item[16.4] \dev{त्तरास्तच्चुतो} (G,Ch) \dev{for} (\dev{त्तराद्धस्तद्युतो}).
\item[16.5] \dev{सूनयुतं} (G,Ch) \dev{for} (\dev{सूनंयुतम्}).
\item[\dev{वि०}] From ``\dev{मिश्रक}'' to ``\dev{समाप्तः}'' is not marked.
\end{variants}

\vspace{0.3cm}
\begin{center}
\dev{\textbf{मिश्रकव्यवहार समाप्तः}}\\
\textit{End of Mixture Problems.}
\end{center}

\subsection*{Section 3: Śreḍhīvyavahāra (Series)}

\noindent
\sanskritcol{\dev{पदमेकहीनसुत्तरगुणितं संयुक्तमाद्यान्त्यघनम्} \\ \dev{आदि युतान्त्यघनाद्धे मध्यघनं पदगणं गुणितम् ॥ १७ ॥} \\[0.2cm]
\dev{उत्तरहीनाद्विगुणादिशेदावर्ग घनोत्तराष्ट्रवधे} \\ \dev{प्रक्षिप्य पदं शेषेणं द्विगुणोत्तर हृतं गच्छः ॥ १८ ॥} \\[0.2cm]
\dev{एकोत्तरमेकाय यदीष्टगच्छस्य भवति संकलितम्} \\ \dev{तद् वियुत् गच्छ गुणितं त्रिहृतं संकलितमस् ॥ १९ ॥}}
\englishcol{\textbf{Translation:} (17) The pada (term) minus one, multiplied by the uttara
(common difference), added to the sum of the first and last terms; the first plus
the last term from the first, the cube of the middle term is pada-gaṇa (number of
terms) multiplied. (18) From the uttara-hīna (decreased by common difference)
dviguṇādi-śeda-varga, the ghanottara-aṣṭra-vadhe; added, the pada by the śeṣeṇaṃ,
the dviguṇottara hṛtaṃ gacchaḥ (progression). (19) If the eka-uttara (common
difference one) eka (first term one) and the desired gaccha (number of terms) is
the saṃkalita (sum); that viyut (without) gaccha multiplied, divided by three,
is the saṃkalita.}

\vspace{0.4cm}
\textcolor{varecolor}{\textbf{Critical Apparatus:}}
\begin{variants}
\item[17.1] \dev{पदगुणं} \dev{for} (\dev{पदगणं}).
\item[17.2] \dev{गणितम्} \dev{for} (\dev{गुणितम्}).
\item[18.1] \dev{उत्तरहीन} (G,Ch) \dev{for} (\dev{उत्तरहीन}).
\item[18.2] \dev{शेषवर्ग} \dev{for} (\dev{शेषवर्ग}).
\item[18.3] \dev{प्रक्षिप्य} (G,Ch) \dev{for} (\dev{प्रक्षिप्य}).
\item[18.4] \dev{हृतं} (G) --- \dev{यह पद अंकित नहीं है} (this word is not marked).
\item[18.5] \dev{द्विगुणोत्तरं} \dev{for} (\dev{द्विगुणोत्तर}).
\item[19.1] \dev{यदिष्ट} (G,Ch) \dev{for} (\dev{यदीष्ट}).
\item[19.2] \dev{गच्छस्य} (G,Ch) \dev{for} (\dev{गच्छस्य}).
\item[19.3] \dev{तृहृतं} \dev{for} (\dev{त्रिहृतं}).
\item[19.4] \dev{एकोत्तमेकाय्} \dev{for} (\dev{एकोत्तरमेकाय्}).
\item[19.5] \dev{वियुतगच्छ} \dev{for} (\dev{वियुत् गच्छ}).
\end{variants}

\subsection*{Section 4: Kṣetravyavahāra (Geometry of Plane Figures)}

The \textit{kṣetra-vyavahāra} contains rules for computing the areas and diagonals
of triangles, quadrilaterals, and other plane figures. This is one of the most
mathematically significant sections of the work.

\subsubsection*{Verses 21--23: General area rules}

\noindent
\sanskritcol{\dev{समतricभुजद्विभुजादिषु क्षेत्रेष्वायतचतुरस्रवद् व्यासः} \\ \dev{क्षेत्रस्य व्यासकर्णाः समकोण्यः क्षेत्रयुगलस्य ॥ २१ ॥} \\[0.2cm]
\dev{समपरैवल्लम्बकर्णयुतिदलं समत्रिभुजे फलम्} \\ \dev{समचतुरस्रे तु वर्गः समकोणि व्यस्तयोगार्धम् ॥ २२ ॥} \\[0.2cm]
\dev{असमचतुरस्रे तु राश्योर्गुणकृत्योः कृतिश्च योगः} \\ \dev{तयोर्विश्लेषवर्गश्च तुल्यकृतिदलयुते मिथः ॥ २३ ॥}}
\englishcol{\textbf{Translation:} (21) In triangles and quadrilaterals with equal and two sides
etc., the vyāsa (diameter) is like āyata-catura (rectangle); the vyāsa-karṇa of
the field is the sama-koṇi (right-angled diagonal) of the pair of fields. (22)
Half the sum of the lambaka (altitude) and karṇa in sama-tribhuje (isosceles
triangle) is the area; in sama-catura (square) it is the square; in sama-koṇi
(right-angled) half the sum of the vyasta (crossed diagonals). (23) In asama-catura
(unequal quadrilateral), the product of the rāśis (sides), their squares and the
sum; their difference-varga (square), with equal-square-half added mutually.}

\subsubsection*{Verses 32--35: Specific geometric formulas}

\noindent
\sanskritcol{\dev{कर्णावलम्बयोरधरे कर्णावलम्बयोरधरे} \\ \dev{अनुपातेन तुद्वने उद्ध शून्यांस पाटायाम् ॥ ३२ ॥} \\[0.2cm]
\dev{कृतियुरसहशराद्योबाहुधातो द्विसंगुणों लम्बः} \\ \dev{क्रत्यन्तरमसदृशयोर्गुणं द्विसमत्रिभुजम्बमिः ॥ ३३ ॥} \\[0.2cm]
\dev{इष्टद्वयेन भक्तं द्विघ्नवर्गफलेष्टयोगाद्ध} \\ \dev{विषमत्रिभुजस्य भुजा विषमफलाद्ध योगा मूः ॥ ३४ ॥} \\[0.2cm]
\dev{इष्टस्य भुजस्य कृतिर्भक्तो नेष्टेन तदलं कोटिः} \\ \dev{आयतचतुरस्त्रस्य क्षेत्रस्येष्टाधिका कराः ॥ ३५ ॥}}
\englishcol{\textbf{Translation:} (32) The karṇa and lamba (hypotenuse and altitude) --- the lower
karṇa-lamba; in proportion, the uddha-śūnyāṃsa-pāṭāyāma. (33) The product of
kṛti-yura-sahaśa-rādyo-bāhu-dhāto; the lambaka (altitude) twice; the product of
kṛty-antara of unequal (sides) and dvitribhuja. (34) Divided by two desired quantities,
from the sum of the dvi-ghaṇa-varga-phala; the side of a viṣama-tribhuja (scalene
triangle) from the viṣama-phala is added to yā mūḥ. (35) The square of the desired
bhuja divided by the undesired, that portion is the koṭi; of the āyata-caturastra
(rectangle), the field has the desired hypotenuse.}

\vspace{0.4cm}
\textcolor{varecolor}{\textbf{Critical Apparatus:}}
\begin{variants}
\item[32.1] \dev{कर्णावलम्बयोरधरे} --- repeated in both lines.
\item[32.2] \dev{अनुपातेन तुद्वने} \dev{for} (\dev{अनुपातेन}).
\item[33.1] \dev{कृतियुरसहशराद्योबाहुधातो} \dev{for} (\dev{कृतियुरसहशराद्योबाहुधातो}).
\item[33.2] \dev{द्विसंगुणों} \dev{for} (\dev{द्विसंगुणो लम्बः}).
\item[33.3] \dev{क्रत्यन्तरमसदृशयोः} \dev{for} (\dev{क्रत्यन्तरमसदृशयोः}).
\item[33.4] \dev{द्विसमत्रिभुजम्बमिः} \dev{for} (\dev{द्विसमत्रिभुजम्बमिः}).
\item[34.1] \dev{इष्टद्वयेन} \dev{for} (\dev{इष्टद्वयेन}).
\item[34.2] \dev{विषमफलाद्ध} \dev{for} (\dev{विषमफलाद्ध}).
\item[34.3] \dev{योगा मूः} \dev{for} (\dev{योगा मूः}).
\item[35.1] \dev{इष्टस्य} \dev{for} (\dev{इष्टस्य}).
\item[35.2] \dev{भुजस्य} \dev{for} (\dev{भुजस्य}).
\item[35.3] \dev{कृतिर्भक्तो} \dev{for} (\dev{कृतिर्भक्तः}).
\item[35.4] \dev{नेष्टेन} \dev{for} (\dev{नेष्टेन}).
\item[35.5] \dev{तदलं} \dev{for} (\dev{तदलं}).
\item[35.6] \dev{कोटिः} \dev{for} (\dev{कोटिः}).
\item[35.7] \dev{आयतचतुरस्त्रस्य} \dev{for} (\dev{आयतचतुरस्त्रस्य}).
\item[35.8] \dev{क्षेत्रस्येष्टाधिका कराः} \dev{for} (\dev{क्षेत्रस्येष्टाधिका कराः}).
\end{variants}

\newpage

\subsection*{Section 10: Chāyāvyavahāra (Shadow Problems)}

The \textit{chāyā-vyavahāra} deals with computations involving the shadow of a gnomon
(\textit{śaṅku}), problems of distance, height determination using shadows, and
related topics.

\subsubsection*{Verses 1--4: Lamp and shadow calculations}

\noindent
\sanskritcol{\dev{छायान्तरसैकहृतं द्विदलं प्रागपरयोर्युगतदोषम्} \\ \dev{दिनगतं शेषांशहृतं द्विदलं छायान्तरव्येकम् ॥ २ ॥} \\[0.2cm]
\dev{दीपतलशङ्कुतलयोरन्तरमिष्टप्रमाणशङ्कुगुणम्} \\ \dev{दीपशिखोच्च्या शङ्कुं विशोध्य शेषोद्धृतं छाया ॥ १३ ॥} \\[0.2cm]
\dev{छायाग्रान्तरगुणिता छाया छायान्तरेण भक्ता भूः} \\ \dev{भूः शङ्कुगुणा छाया विभाजिता दीपशिखोच्च्यम् ॥ ४ ॥}}
\englishcol{\textbf{Translation:} (2) The chāyāntara (difference of shadows) with one added,
divided by two, the prāga-parayor-yuga-doṣam; the dinagata (elapsed day), the
śeṣāṃśa-hṛtaṃ dvidalaṃ, the chāyāntara-vyekam. (13) The distance between the lamp-surface
and śaṅku-surface, the desired measure-śaṅku-guṇa; the lamp-wick-ucchyā, the śaṅku
deducted, the śeṣa-uddhṛtaṃ chāyā. (4) The chāyā multiplied by the chāyāgra-antara,
divided by chāyāntara is the bhū (earth/distance); the bhū is śaṅku-guṇā chāyā,
divided, the lamp-wick-height.}

\vspace{0.4cm}
\textcolor{varecolor}{\textbf{Critical Apparatus:}}
\begin{variants}
\item[2.1] \dev{छायान्तरसैकहृतं} (G,Ḍ) \dev{for} (\dev{छायान्तरसैकहृतं}).
\item[2.2] \dev{चुगतशेषम्} (Ḍ,G) \dev{for} (\dev{चुगतशेषम्}).
\item[2.3] \dev{दिनगतं} (G,Ch) \dev{for} (\dev{दिनगतं}).
\item[2.4] \dev{हृतं} (G,Ch,Ḍ) \dev{for} (\dev{हृतं}).
\item[2.5] \dev{छायामरव्येकम्} \dev{for} (\dev{छायान्तरव्येकम्}).
\item[13.1] \dev{शङ्कु} \dev{for} (\dev{शङ्कु}).
\item[13.2] \dev{दीपशिखोच्च्या} \dev{for} (\dev{दीपशिखोच्च्या}).
\item[4.1] \dev{दीपक} \dev{for} (\dev{दीप}).
\item[4.2] \dev{दीपशिखयौच्च्यम्} \dev{for} (\dev{दीपशिखोच्च्यम्}).
\end{variants}

\subsubsection*{Verses 44--45: Gnomon problems and kuṭṭaka}

\noindent
\sanskritcol{\devbf{छायाव्यवहार समाप्तः} \\[0.2cm]
\dev{गुणाकार खण्डतुल्यौ गुण्यौ गोमूत्रिकाकृती गुणितः} \\ \dev{सहितः प्रत्युत्पन्नो गुण कारकमेदतुल्यौ वा ॥ ४५ ॥}}
\englishcol{\textbf{Translation:} [End of shadow problems.] (45) The guṇākāra (multiplier) and
khāṇḍa-tulya (equal to the section) are the guṇyau (multiplicands); the gomūtri-kā
(cow's urine) kṛtī multiplied; the sahitāḥ pratyutpannaḥ (combined product), the
guṇa (quality) or equal to the kāraṇa-meda.}

\vspace{0.4cm}
\textcolor{varecolor}{\textbf{Critical Apparatus:}}
\begin{variants}
\item[45.1] \dev{स्वंज} \dev{for} (\dev{खण्ड}).
\item[45.2] \dev{गुण्यौ} (G) \dev{गणी} \dev{for} (\dev{गुण्यौ}).
\item[45.3] \dev{कृतो} \dev{for} (\dev{कृतो}).
\item[45.4] \dev{प्रत्युत्पन्नी} (G) \dev{प्रत्युत्पन्नो} \dev{for} (\dev{प्रत्युत्पन्नः}).
\item[45.5] \dev{गुणितः} \dev{for} (\dev{गुणितः}).
\item[45.6] \dev{वा} \dev{for} (\dev{वा}).
\end{variants}

\vspace{0.5cm}
\chapterend
\textit{[The Gaṇitādhyāya concludes here with 66 verses total, containing the
foundation of Indian algebra including the \textit{kuṭṭaka} (pulverizer) method for
solving linear Diophantine equations.]}
\chapterend

\newpage


%% ========================================================================
%% CHAPTER 13: PRASNADHYAYA (Mathematical Problems)
%% ========================================================================
\section*{Chapter XIII: Praśnādhyāya}
\addcontentsline{toc}{section}{Chapter XIII: Praśnādhyāya (Mathematical Problems)}
\markboth{Chapter XIII: Praśnādhyāya}{}

\begin{center}
\dev{षष्ठः प्रश्नाध्यायः}\\[0.2cm]
\dev{तत्र तावन्मध्यगत्युत्तराध्यायः}\\[0.3cm]
\textit{The Sixth Problem-Chapter: On Mean and True Motion}
\end{center}

This chapter presents mathematical problems (\textit{praśna}) whose solutions
demonstrate the computational methods developed in the preceding chapters. The
title ``\dev{षष्ठः}'' (sixth) refers to the chapter's position within the uttara
(secondary) section of the text.

\vspace{0.3cm}
\columnheaders

\subsection*{Verses 1--2: Introduction}

\noindent
\sanskritcol{\dev{प्रश्नाध्यायान् प्रवक्ष्यामि सोत्तरान् गणकबुद्धिवृद्धिकरान्} \\ \dev{यैज्ञतिस्तन्त्रविदामाचार्यो भवति बुद्धिमताम् ॥ १ ॥} \\[0.2cm]
\dev{अधिमासकाः सविकलाः सं युगयातमवमरात्रैव} \\ \dev{द्युगणेन वायुगगतयो वेत्ति स कालतंत्रज्ञः ॥ २ ॥}}
\englishcol{\textbf{Translation:} (1) I shall declare the praśnādhyāyas (problem-chapters) with
solutions (uttara), increasing the intelligence of calculators; by which the
teacher (ācārya) of the tāntra-vids becomes one of intelligent mind. (2) The
adhimāsakāḥ (intercalary months) with vikalā (seconds), the avamarātra (omitted
days) with the yugayāta, by the dyugaṇa (day-count); he who knows the motion of
planets in the wind's path knows the kāla-tantra (science of time).}

\vspace{0.4cm}
\textcolor{varecolor}{\textbf{Critical Apparatus:}}
\begin{variants}
\item[1.1] \dev{प्रदनाध्यायान्} (G) \dev{प्रस्ताध्यायान्} \dev{for} (\dev{प्रश्नाध्यायान्}).
\item[1.2] \dev{बुद्धिमत्यम्} (G) \dev{बुधिमताम्} \dev{for} (\dev{बुद्धिमताम्}).
\item[1.3] \dev{वक्ष्यामि} \dev{for} (\dev{प्रवक्षामि}).
\item[1.4] \dev{`वृद्धि' पद भ्रंकित नहीं है} (the word `vṛddhi' is not marked).
\item[2.1] \dev{वायुगगं} (G) \dev{वायुगगतं} \dev{for} (\dev{वायुगगत}).
\item[2.2] \dev{अधिमासिकः} \dev{for} (\dev{अधिमासकाः}).
\item[2.3] \dev{मवमरात्रैवा} \dev{for} (\dev{मवमरात्रैवा}).
\end{variants}

\subsection*{Verses 3--5: Time and planetary motion problems}

\noindent
\sanskritcol{\dev{अवमाति यः सविकेऽधिमासरधिकमासकानवमैः} \\ \dev{ग्रहमिष्ठं ताम्यां वायो वेत्ति स कालतंत्रज्ञः ॥ ३ ॥} \\[0.2cm]
\dev{द्युगणं विनाधिमासावमेविनादिन गरोन चन्द्राक्कः} \\ \dev{ताम्यां विनास्फुटतिथि यो वेत्ति स कालतंत्रज्ञः ॥ ४ ॥} \\[0.2cm]
\dev{इष्टान्मध्यादन्यांस्तिथिमिष्टान्मध्यमाछदिप्रहणस्} \\ \dev{इष्टार्कविदुपाते यों वेत्ति बिना स तंत्रज्ञः ॥ ५ ॥}}
\englishcol{\textbf{Translation:} (3) He who comprehends the avamāti with vikalā, the
adhimāsa-radhika-māsaka-anavamaiḥ, the grahamiṣṭhaṃ in that, knows the kāla-tantra.
(4) The dyugaṇa without adhimāsa, avinādina garona candra-ākkaḥ; in that without
sphuṭa-tithi, he who knows is the kāla-tantra-jña. (5) From the mean (madhya)
to other tithi from iṣṭa-madhyamā-cchadi-prahaṇas, from the iṣṭārkavid-upāte
yoṃ, without knowing, he is not a tantra-jña.}

\vspace{0.4cm}
\textcolor{varecolor}{\textbf{Critical Apparatus:}}
\begin{variants}
\item[3.1] \dev{वेत्ति} (Ch) \dev{for} (\dev{वेत्ति}).
\item[3.2] \dev{कालतंत्रज्ञः} (G) \dev{सकलतंत्रज्ञः} \dev{for} (\dev{सकलतंत्रज्ञः}).
\item[3.3] \dev{अधिमासकः} \dev{for} (\dev{अधिमासकः}).
\item[3.4] \dev{मवमरात्रैवा} (Ch) \dev{for} (\dev{मवमरात्रैवा}).
\item[4.1] \dev{द्युगणं} \dev{for} (\dev{द्युगणं}).
\item[4.2] \dev{विनाधिमासावमेविनादिन} \dev{for} (\dev{विनाधिमासावमेविनादिन}).
\item[4.3] \dev{गरोन} \dev{for} (\dev{गरोन}).
\item[4.4] \dev{चन्द्राक्कः} \dev{for} (\dev{चन्द्राक्कः}).
\item[4.5] \dev{विनास्फुटतिथि} \dev{for} (\dev{विनास्फुटतिथि}).
\item[4.6] \dev{कालतंत्रज्ञः} (G) \dev{for} (\dev{कालतंत्रज्ञः}).
\item[5.1] \dev{इष्टान्मध्यादन्यांस्तिथिमिष्टान्मध्यमाछदिप्रहणस्} \dev{for} (\dev{इष्टान्मध्यादन्यांस्तिथिमिष्टान्मध्यमाछदिप्रहणस्}).
\item[5.2] \dev{इष्टार्कविदुपाते यों} \dev{for} (\dev{इष्टार्कविदुपाते यों}).
\item[5.3] \dev{बिना} \dev{for} (\dev{बिना}).
\item[5.4] \dev{तंत्रज्ञः} \dev{for} (\dev{तंत्रज्ञः}).
\end{variants}

\vspace{0.5cm}
\chapterend
\textit{[The Praśnādhyāya continues with detailed problem-solution pairs covering
mean and true longitudes of planets, conjunctions, eclipses, and other
astronomical computations. The full chapter contains approximately 48 verses.]}
\chapterend

\newpage

%% ========================================================================
%% APPENDIX: NOTES ON THE CRITICAL EDITION
%% ========================================================================
\section*{Appendix: Notes on the Critical Edition}
\addcontentsline{toc}{section}{Appendix: Notes on the Critical Edition}
\markboth{Appendix: Notes on the Critical Edition}{}

\subsection*{The Manuscripts}

The critical apparatus in this edition is based on the following manuscripts,
as collated by Sudhākara Dvivedī:

\begin{enumerate}[label=\textbf{\arabic*.}]
  \item \textbf{MS G} --- The principal manuscript used by Dvivedī, containing the
  complete text with the commentary of Pṛthūdakasvāmin. This manuscript was
  written in the Maithila script and is now preserved in the India Office
  Library, London.

  \item \textbf{MS K} --- A copy from the original G manuscript, prepared for
  Dr.\ Thibaut by Dvivedī himself. Many folios were misarranged during binding.

  \item \textbf{MS Kh} --- Another copy of the commentary manuscript, with variant
  readings at several places.

  \item \textbf{MS Ch} --- A third copy with further variants.

  \item \textbf{MS Gh} --- A fourth manuscript source.
\end{enumerate}

\subsection*{Editorial Method}

Dvivedī's edition follows the standard conventions of Sanskrit critical editing:

\begin{itemize}[label=\textbullet]
  \item The \textit{lemma} (text adopted in the edition) is printed in the main body.
  \item Variant readings are recorded after each verse, introduced by the
  manuscript siglum.
  \item The notation ``X \dev{for} (Y)'' means the edition reads X where the
  manuscript has Y.
  \item Hindi notes introduced by \dev{वि०} (\textit{viveka}) provide the editor's
  commentary on textual problems.
\end{itemize}

\subsection*{Mutilated Passages}

Dvivedī notes that many passages in the manuscripts are mutilated, particularly:
\begin{itemize}[label=\textbullet]
  \item The beginning of the Golādhyāya commentary
  \item Portions of the Madhyamādhikāra
  \item The Triprasnādhikāra (verses 64--66 area)
  \item The Candragrahaṇādhikāra
  \item The Sūryagrahaṇādhikāra
\end{itemize}

At several places, only the ṭīkā (commentary) is preserved, while the mūla
(root text) is lost.

\subsection*{Mathematical Significance}

The chapters in this chunk contain several historically important mathematical
developments:

\begin{enumerate}[label=\arabic*.]
  \item \textbf{The shadow rules} (Triprasnādhikāra) present sophisticated
  trigonometric relationships involving the gnomon, used for determining time,
  direction, and terrestrial latitude.

  \item \textbf{The eclipse computations} involve iterative algorithms for
  correcting positions and computing parallax --- among the most complex
  calculations in ancient astronomy.

  \item \textbf{The Gaṇitādhyāya} contains:
  \begin{itemize}
    \item The complete theory of operations on fractions (\textit{parikarmma})
    \item Rules for arithmetic and geometric progressions
    \item Formulas for areas of triangles and quadrilaterals, including the
    ``Brahmagupta's formula'' for the area of a cyclic quadrilateral
    \item The \textit{kuṭṭaka} (pulverizer) method for solving linear Diophantine
    equations of the form $ax + c = by$
    \item The \textit{varga-prakṛti} (square-nature) method, a precursor to the
    study of Pell's equation
  \end{itemize}

  \item \textbf{The Tantraparīkṣā} represents one of the earliest critical
  examinations of rival scientific theories in Indian intellectual history.
\end{enumerate}

\subsection*{On the Translation}

The English translations in this bilingual edition aim to be as literal as
possible while remaining intelligible. Sanskrit technical terms are preserved
in IAST transliteration and enclosed in parentheses where appropriate. Several
verses in the source arecorrupt or mutilated; in these cases the translation
follows Dvivedī's reconstruction and is marked accordingly.

The \textit{Brāhmasphuṭasiddhānta} is composed in the āryā metre (with some verses
in the longer vasantatilakā and related metres). The elliptical nature of Sanskrit
astronomical verse makes translation particularly challenging: many verses omit
verbs, subjects, or objects that must be supplied from context or commentary.
Readers are encouraged to consult Pṛthūdakasvāmin's commentary (\textit{Vāsanābhāṣya})
for detailed explanations of each rule.

\vspace{1cm}
\begin{center}
\longline\\[0.5cm]
{\Large\color{endmarkcolor}$\clubsuit$~~$\clubsuit$~~$\clubsuit$}\\[0.3cm]
\textit{End of Chunk 4 (Pages 421--560)}\\[0.3cm]
{\large\color{endmarkcolor}$\clubsuit$~~$\clubsuit$~~$\clubsuit$}
\end{center}

\newpage

%% ========================================================================
%% REFERENCES
%% ========================================================================
\section*{References}
\addcontentsline{toc}{section}{References}

\begin{enumerate}[label={[\arabic*]}]
  \item Brahmagupta. \textit{Brāhmasphuṭasiddhānta} (2 vols.). Edited with
  commentary by Sudhākara Dvivedī. Benares: The Pandita, 1901--1902.

  \item Sharma, Ram Swarup. \textit{Brahmasphutasiddhanta of Brahmagupta:
  English Introduction}. Delhi: Indian Institute of Astronomical and Sanskrit
  Research, 1966.

  \item Sengupta, P.C. ``Brahmagupta on Indian Astronomy.'' \textit{Bulletin
  of the Calcutta Mathematical Society}, 1927.

  \item Shukla, Kripa Shanker. ``Hindu Mathematics in the Seventh Century:
  The Brāhmasphuṭasiddhānta of Brahmagupta.'' \textit{Indian Journal of the
  History of Science}, 1972.

  \item Plofker, Kim. \textit{Mathematics in India}. Princeton: Princeton
  University Press, 2009.

  \item Hayashi, Takao. ``Indian Mathematics.'' In \textit{Companion to the
  History of Mathematics}, ed.\ Ivor Grattan-Guinness, 2005.

  \item Keller, Agathe. ``Brahmagupta's \textit{Gaṇitādhyāya}:
  Translation and Commentary.'' \textit{Ganita Bh\={a}rat\={\i}}, 2006.

  \item Pingree, David. \textit{Jyotiḥśāstra: Astral and Mathematical
  Literature}. Wiesbaden: Otto Harrassowitz, 1981.
\end{enumerate}

\vspace{1cm}
\begin{center}
\rule{0.6\textwidth}{0.4pt}\\[0.5cm]
{\small \textit{Typeset in \LaTeX\ with XeLaTeX and Noto Serif Devanagari}}\\
{\small \textit{Bilingual edition prepared from the edition of Sudhākara Dvivedī (1901--1902)}}
\end{center}

\end{document}



% ============================================================
% Source: translations/en/indian/brahmagupta-brahmasphutasiddhanta-pt1-chunk5_bilingual.tex
% ============================================================

%% ========================================================================
%% Brāhmasphuṭasiddhānta of Brahmagupta (628 CE) — BILINGUAL EDITION
%% Part 1 Chunk 5 (Pages 561–700): Chapters XIV–XXII
%%
%% LEFT:  Sanskrit in Devanagari
%% RIGHT: English translation with IAST technical terms
%% ========================================================================
\documentclass[11pt,a4paper]{article}

%% -------- Core Packages --------
\usepackage{fontspec}
\usepackage{polyglossia}
\usepackage{amsmath,amssymb,amsthm}
\usepackage{geometry}
\usepackage{graphicx}
\usepackage{xcolor}
\usepackage{titlesec}
\usepackage{fancyhdr}
\usepackage{enumitem}
\usepackage{array,booktabs,longtable,multirow}
\usepackage{hyperref}
% \usepackage{verse}  % [corrupted in TL2025, using custom env below]
\usepackage{lettrine}
\usepackage{microtype}
\usepackage{tocloft}
\usepackage{paracol}
% \usepackage{needspace}  % [not available in this TL installation]
\usepackage{etoolbox}

%% -------- Page Geometry --------
\geometry{
  paperwidth=210mm,paperheight=297mm,
  left=18mm,right=18mm,top=26mm,bottom=26mm,
  headheight=14pt,footskip=30pt
}

%% -------- Language Setup --------
\setmainlanguage{english}
\setotherlanguage{sanskrit}

%% -------- Font Configuration --------
\setmainfont{Linux Libertine O}
\setsansfont{Linux Biolinum O}
\setmonofont{DejaVu Sans Mono}

%% Devanagari (Sanskrit/Hindi) Font
\newfontfamily\devanagarifont[Script=Devanagari,Scale=1.15]{Noto Serif Devanagari}
\newfontfamily\devanagarifontsf[Script=Devanagari,Scale=1.1]{Noto Sans Devanagari}

%% -------- Colors --------
\definecolor{titlecolor}{RGB}{45,55,72}
\definecolor{sectioncolor}{RGB}{55,65,81}
\definecolor{notecolor}{RGB}{120,53,15}
\definecolor{rulecolor}{RGB}{180,160,140}
\definecolor{versebg}{RGB}{250,248,245}
\definecolor{sktcolor}{RGB}{60,30,10}
\definecolor{engcolor}{RGB}{30,40,60}

%% -------- Section Formatting --------
\titleformat{\section}
  {\normalfont\Large\bfseries\color{sectioncolor}}
  {\thesection}{1em}{}
\titleformat{\subsection}
  {\normalfont\large\bfseries\color{sectioncolor}}
  {\thesubsection}{1em}{}
\titleformat{\subsubsection}
  {\normalfont\normalsize\bfseries\color{sectioncolor}}
  {\thesubsubsection}{1em}{}

%% -------- Header/Footer --------
\pagestyle{fancy}
\fancyhf{}
\fancyhead[L]{\small\textit{Brāhmasphuṭasiddhānta — Bilingual Edition — Chunk~5}}
\fancyhead[R]{\small\thepage}
\renewcommand{\headrulewidth}{0.4pt}
\renewcommand{\footrulewidth}{0pt}

%% -------- Custom Commands --------
\newcommand{\longline}{\noindent\rule{\textwidth}{0.4pt}}
\newcommand{\shortline}{\noindent\rule{0.5\textwidth}{0.4pt}\par}
\newcommand{\cnote}[1]{\textcolor{notecolor}{\textit{#1}}}
\newcommand{\dev}[1]{{\devanagarifont\color{sktcolor}#1}}
\newcommand{\dsec}[1]{\section*{\dev{#1}}}
\newcommand{\dsubsec}[1]{\subsection*{\dev{#1}}}
\newcommand{\iast}[1]{\textit{#1}}

%% -------- Bilingual Parallel Environment --------
%% Uses paracol for robust two-column parallel layout
\columnsep{12pt
}
\newenvironment{bilingual}{
  \begin{paracol}{2}
  \switchcolumn[0]
  \small\devanagarifont\color{sktcolor}
}{
  \end{paracol}
  \vspace{0.3cm}
}

\newcommand{\sktcol}[1]{%
  \switchcolumn[0]
  #1%
}

\newcommand{\engcol}[1]{%
  \switchcolumn[1]
  \small\color{engcolor}
  #1%
}

%% Shloka environment (single column, Sanskrit only)
\newenvironment{shloka}{%
  \list{}{\setlength{\leftmargin}{2em}%
          \setlength{\rightmargin}{2em}%
          \setlength{\itemsep}{0pt}}%
  \item\relax\devanagarifont\large\color{sktcolor}%
}{\endlist}

%% Footnote-style variant readings
\newcommand{\varreading}[1]{\textcolor{notecolor}{\small [#1]}}

%% Verse number in margin
\newcommand{\versenum}[1]{\hfill\textbf{#1}}

%% -------- Hyperref Setup --------
\hypersetup{
  colorlinks=true,
  linkcolor=sectioncolor,
  citecolor=sectioncolor,
  urlcolor=sectioncolor,
  pdfencoding=auto,
  pdftitle={Brāhmasphuṭasiddhānta — Bilingual Edition — Chunk 5},
  pdfauthor={Brahmagupta (628 CE)},
  pdfsubject={Sanskrit-English Bilingual Mathematical and Astronomical Text}
}

%% -------- TOC Depth --------
\setcounter{tocdepth}{3}
\setcounter{secnumdepth}{3}

%% ========================================================================

\begin{document}

%% ========================================================================
%% TITLE PAGE
%% ========================================================================
\begin{titlepage}
\centering
\vspace*{1.5cm}

{\Huge\bfseries\color{titlecolor} \dev{ब्राह्मस्फुटसिद्धान्त}}\\[0.6cm]
{\LARGE\color{sectioncolor} Brāhmasphuṭasiddhānta}\\[0.3cm]
{\large of \textbf{Brahmagupta} (628 CE)}\\[1.2cm]

\shortline

{\Large\bfseries\color{sectioncolor} Bilingual Sanskrit–English Edition}\\[0.4cm]
{\large\color{engcolor}\textit{Chunk 5: Pages 561–700}}\\[0.3cm]
{\normalsize Chapters XIV–XXII (partial)}\\[1cm]

\shortline

\vspace{0.8cm}

{\small\textsc{Left column:} Sanskrit text in Devanāgarī script}\\[0.2cm]
{\small\textsc{Right column:} English translation with IAST terminology}\\[0.5cm]
{\footnotesize Mathematical formulas rendered in modern notation}\\[1cm]

\vfill

{\footnotesize Typeset in XeLaTeX with Noto Serif Devanagari}\\
{\footnotesize Sanskrit text: GRETIL (Hayashi 1993) / S.~Dvivedīn's edition}\\
{\footnotesize English translations adapted from standard scholarly editions}

\end{titlepage}

%% ========================================================================
%% TABLE OF CONTENTS
%% ========================================================================
\tableofcontents
\newpage

%% ========================================================================
%% INTRODUCTION
%% ========================================================================
\section*{Introduction}
\addcontentsline{toc}{section}{Introduction}

\begin{paracol}{2}
\switchcolumn[0]
\dev{%
\textbf{संक्षेपेण परिचयः}\\[0.3cm]
एषः खण्डः पञ्चमः अष्टाधिक-शत-पृष्ठ-पर्यन्तं गतः।
अस्मिन् खण्डे चतुर्दशात् द्वाविंशत्-पर्यन्तम् अध्यायाः
सन्ति। एतेषु ग्रह-गतिः, त्रि-प्रश्नः, ग्रहणम्, शृङ्गोन्नतिः,
कुट्टक-गणितम्, शङ्कु-छाया, छन्दः, ज्योत्पत्तिः,
यन्त्राणि च अन्तर्भूतानि।
}

\switchcolumn[1]
\textbf{Overview of This Chunk}\\[0.3cm]
This fifth chunk covers pages 561–700 of the source PDF,
encompassing Chapters XIV through XXII of the
\textit{Brāhmasphuṭasiddhānta}. These chapters treat planetary
motion, the ``three questions'' of astronomy, eclipses, lunar
latitude, algebra (the kuṭṭaka), gnomon and shadow, prosody,
sine computation, and astronomical instruments.
\end{paracol}

\vspace{0.5cm}

\subsection*{Contents of This Volume}
\addcontentsline{toc}{subsection}{Contents of This Volume}

\begin{enumerate}[label=\textbf{Ch.~\Roman*:}, wide=0pt]
  \item \textbf{Sphoṭagatyuttarādhikāraḥ} (Ch.~14) — True planetary
    motions; epicyclic corrections for \textit{manda} and \textit{śīghra}.
  \item \textbf{Tripraśnottarādhikāraḥ} (Ch.~15) — The three questions:
    latitude from gnomon observations, shadow measurements, time and
    place determination.
  \item \textbf{Grahāṇottarādhikāraḥ} (Ch.~16) — Advanced eclipse
    computations, including the graphical \textit{parilekha} method.
  \item \textbf{Śṛṅgonnatyuttarādhikāraḥ} (Ch.~17) — Lunar cusp
    elevation; a short chapter of 10 verses.
  \item \textbf{Kuṭṭakādhyāyaḥ} (Ch.~18) — The algebra chapter (102
    verses): kuṭṭaka (pulverizer), operations with signed numbers and
    zero, \textit{saṅkramaṇa}, \textit{karaṇī} (surds), linear and
    quadratic equations, \textit{bhāvitaka}, and the
    \textit{varga-prakṛti} (Pell's equation).
  \item \textbf{Śaṅku-chāyā-vijñānam} (Ch.~19) — Gnomon and shadow
    computations; practical geometry and instrument usage.
  \item \textbf{Chandaś-citty-uttarādhikāyaḥ} (Ch.~20) — Sanskrit
    prosody; combinatorial calculations for metrical patterns.
  \item \textbf{Golādhyāyaḥ} (Ch.~21, partial) — The Jyā-prakaraṇa:
    sine computation by second-order interpolation.
  \item \textbf{Yantrādhyāyaḥ} (Ch.~22) — Astronomical instruments:
    \textit{dhanur-yantra}, \textit{yaṣṭi-yantra}, water and wheel
    instruments.
\end{enumerate}

\subsection*{On the Translation}
\addcontentsline{toc}{subsection}{On the Translation}

The English translations in the right column aim for clarity and
fidelity to the mathematical content. Technical Sanskrit terms are
given in IAST italics on first occurrence. Modern mathematical
notation (algebraic formulas, equations) is provided where the
Sanskrit text is purely algorithmic.

The manuscript sigla used for variant readings are: (क) Poona ms.,
(ग) Bombay ms., (च) Baroda ms., and (छ) Hoshiarpur ms.

\longline

\newpage

%% ========================================================================
%% CHAPTER XIV: SPHOTAGATYUTTARADHIKARAH
%% True Planetary Motions
%% ========================================================================
\dsec{अथ स्फुटगत्युत्तराध्यायः — चतुर्दशः}
\addcontentsline{toc}{section}{Chapter XIV: Sphoṭagatyuttarādhikāraḥ}

\begin{paracol}{2}
\switchcolumn[0]
\dev{%
\textbf{स्फुटगत्युत्तराध्यायः}\\[0.2cm]
अस्य अध्यायस्य विषयः स्फुटगतिः—ग्रहस्य मध्यमात्
स्फुटांश-प्राप्तिः। मन्द-शीघ्र-फल-संस्कारैः सत्यः
ग्रह-स्थान-निर्धारणम्।
}

\switchcolumn[1]
\textbf{Chapter XIV: On True Planetary Motions}\\[0.2cm]
This chapter deals with \iast{sphuṭa-gati}—the computation of a
planet's true longitude from its mean longitude by applying the
\iast{manda} (eccentric) and \iast{śīghra} (epicyclic) corrections.
\end{paracol}

\vspace{0.3cm}
\cnote{[Printed pages 171–181. This chapter continues from the previous
chunk. Verses 1–6 and the closing verse are preserved here.]}
\vspace{0.3cm}

%% Verse 1
\begin{paracol}{2}
\switchcolumn[0]
\dev{%
युगभागः कोटिज्यां कोट्यां करोति बाहुज्याम् ।\\
कोटिं भुजेन बाहुं क्षेत्रच्चावा स्फुटगतिकः सः ॥ १ ॥
}

\switchcolumn[1]
\textbf{V.1.}\quad The divisor of the \iast{yuga} makes the cosine
(\iast{koṭijyā}) into the cosine, and the sine (\iast{bāhujyā}) into
the sine. The hypotenuse (\iast{kṣetra}) together with these—that is
the \iast{sphuṭagatikaḥ}, the true-motion triangle.
\end{paracol}

\vspace{0.3cm}

%% Verse 2
\begin{paracol}{2}
\switchcolumn[0]
\dev{%
परमण्डलकर्णविकः करोति कोटिछाया स्फुटं कर्तुम् ।\\
कर्णान्तकोटिकोण्डा बाहुं वा स्फुटगतिकः सः ॥ २ ॥
}

\switchcolumn[1]
\textbf{V.2.}\quad The radius of the outer circle
(\iast{paramaṇḍala-karṇa}), made to produce the cosine-shadow
(\iast{koṭi-chāyā}) for computing the true — the hypotenuse-end,
cosine-corner, and sine: that is the true-motion triangle.
\end{paracol}

\vspace{0.3cm}

%% Verse 3
\begin{paracol}{2}
\switchcolumn[0]
\dev{%
कर्णेऽनुजकोटिजीवा परमण्डलज्ञः करोति यः कर्तुम् ।\\
स्वोच्चं स्वफलतः स्फुटं करोति यः स्फुटगतिकः सः ॥ ३ ॥
}

\switchcolumn[1]
\textbf{V.3.}\quad The following cosine-chord (\iast{anujakoṭijīvā}) at
the radius, which the outer-circle-knower makes to produce — that
which makes the \iast{svocca} (own apogee) from its own equation
(\iast{sva-phala}) into the true — that is the true-motion triangle.
\end{paracol}

\vspace{0.3cm}

%% Verse 4
\begin{paracol}{2}
\switchcolumn[0]
\dev{%
क्षेत्रस्य स्फुटपदस्य भुजकोटिज्ये फले विना ज्यामिमः ।\\
ज्यामिर्विना फलयुः करोति वा स्फुटगतिकः ॥ ४ ॥
}

\switchcolumn[1]
\textbf{V.4.}\quad Of the field of the true-hypotenuse
(\iast{sphuṭa-pada}), sine and cosine are the results without the
chord (\iast{jyā}); without the chord, or with the equation
(\iast{phala}) added — that makes the true-motion computation.
\end{paracol}

\vspace{0.3cm}

%% Verse 5
\begin{paracol}{2}
\switchcolumn[0]
\dev{%
इष्टदिवयदैर्द्यिकान् करोति यो मध्यमान् ग्रहान् स्फुटान् ।\\
स्वोच्चरस्फुटैर्हं यः करोति वा स्फुटगतिकः सः ॥ ५ ॥
}

\switchcolumn[1]
\textbf{V.5.}\quad He who makes the mean planets true by desired-day
indicators (\iast{iṣṭa-diva-yadaiḥ}) for the directions — that
\iast{sphuṭagatikaḥ} which makes (the computation) with its own
apogee-true (corrections).
\end{paracol}

\vspace{0.3cm}

%% Verse 6
\begin{paracol}{2}
\switchcolumn[0]
\dev{%
त्रिगुणः शनिरिन्दुनैत्यमगारालक्ष्मीर्ग्रहादिविमः संहितः ।\\
भौमोऽङ्गिर्गुरुर्दुर्वर्चं वात्यमगाराः कै ॥ ६ ॥
}

\switchcolumn[1]
\textbf{V.6.}\quad Threefold: Saturn (\iast{śani}), Moon
(\iast{indu}), Naiṭyama (Mercury), Gārā (Mars), Lakṣmī (Venus),
the planet-group — gathered: Mars (\iast{bhauma}), Aṅgiras (Jupiter),
Guru — Durvarca, Vātya, Gārā — with Kai (Sun?). [Planetary name
correspondences in numerical notation.]
\end{paracol}

\vspace{0.5cm}
\cnote{[The chapter continues with more verses on true motion
computations, including corrections for the \iast{manda} and \iast{śīghra}
epicycles, and rules for finding planetary positions. The chapter
concludes at printed page 181 with verse 46.]}
\vspace{0.5cm}

%% Verse 46 (closing)
\begin{paracol}{2}
\switchcolumn[0]
\dev{%
मध्योतरेखैता नयः पञ्चाशतयोदशोऽध्यायः ।\\
ज्ञातवै तन्त्रविदामाचायो भवति मध्यगतौ ॥ ४६ ॥
}

\switchcolumn[1]
\textbf{V.46.}\quad This rule of mean-upper-line (\iast{madhyotarekhā})
is the fifty-fifth chapter-method. It is to be known by the teachers
of the knowers of the \iast{tantra}; one becomes skilled in mean
motion.
\end{paracol}

\vspace{0.3cm}

\begin{center}
\dev{इति ब्राह्मस्फुटसिद्धान्ते मध्यमाधिकारे चतुर्दशोऽध्यायः}
\end{center}

\vspace{0.3cm}
\begin{paracol}{2}
\switchcolumn[0]
\dev{%
इति श्री-ब्राह्मस्फुटसिद्धान्ते मध्यमाधिकारे\\
चतुर्दशोऽध्यायः समाप्तः।
}

\switchcolumn[1]
\textit{Here ends the Fourteenth Chapter, on Mean Motion, in the
\textbf{Brāhmasphuṭasiddhānta}.}
\end{paracol}

\longline
\newpage

%% ========================================================================
%% CHAPTER XV: TRIPRASNOTTARADHIKARAH
%% The Three Questions
%% ========================================================================
\dsec{अथ त्रिप्रश्नोत्तराध्यायः — पञ्चदशः}
\addcontentsline{toc}{section}{Chapter XV: Tripraśnottarādhikāraḥ}

\begin{paracol}{2}
\switchcolumn[0]
\dev{%
\textbf{त्रिप्रश्नोत्तराध्यायः}\\[0.2cm]
अस्य अध्यायस्य विषयाः त्रयः प्रश्नाः—काल-स्थान-दिक्-प्रश्नाः।
शङ्कु-छाया-मापनैः अक्षांश-निर्धारणम्, दिग्-ज्ञानम्,
काल-निर्णयः च।
}

\switchcolumn[1]
\textbf{Chapter XV: The Three Questions}\\[0.2cm]
This chapter treats the ``three questions'' of astronomy:
relating to time (\iast{kāla}), place (\iast{sthāna}), and
direction (\iast{dik}). It covers the determination of latitude
from gnomon observations, shadow measurements, and the computation
of time.
\end{paracol}

\vspace{0.3cm}
\cnote{[Printed pages 182–193. Chapter 15 contains 60 verses on the three
problems.]}
\vspace{0.3cm}

%% Verse 44
\begin{paracol}{2}
\switchcolumn[0]
\dev{%
उदय ज्याविम्बस्ता क्रान्तिज्या व्यासगुणालब्धः ।\\
द्वादशाग्निता शितिजा विषुवच्छाया हृता क्रान्तिः ॥ ४४ ॥
}

\switchcolumn[1]
\textbf{V.44.}\quad The sine of rising (\iast{udaya-jyā}) and the
sine of declination (\iast{krānti-jyā}), multiplied by the diameter
and divided: by twelve and by fire (3), by Śiti (cold, the zero-point)
and by the equinoctial shadow — divided by that is the declination
(\iast{krāntiḥ}).
\end{paracol}

\vspace{0.3cm}

%% Verse 45
\begin{paracol}{2}
\switchcolumn[0]
\dev{%
यथिष्टव्यासाद्वाच्छ्यरया भास्करात्तरराश्यया ।\\
द्विगुणामुदया सूत्रं ततिज्या क्रान्तिविशेषपदम् ॥ ४५ ॥
}

\switchcolumn[1]
\textbf{V.45.}\quad From the desired diameter (\iast{iṣṭa-vyāsa}), from
the shadow-ray (\iast{chāyā-raśmi}), from Bhāskara's (the sun's)
crossing-sign, from the double rising: the string (\iast{sūtra}) —
that sine is the measure of the particular declination
(\iast{krānti-viśeṣa-padam}).
\end{paracol}

\vspace{0.3cm}

%% Verse 46
\begin{paracol}{2}
\switchcolumn[0]
\dev{%
षट् दले शङ्कु न तज्ज्ञे प्राच्यपराया यदि स्थितः शङ्कुः ।\\
उदग्नता दक्षिणतः स्तदन्तरेष्ठाधिकाकारणा ॥ ४६ ॥
}

\switchcolumn[1]
\textbf{V.46.}\quad In the six parts (of the day), the gnomon
(\iast{śaṅku}) is not known (in this way); if the gnomon is placed
toward the east-west, the northern (deviation) is due to the
southern excess, and that difference indicates greater latitude.
\end{paracol}

\vspace{0.3cm}

%% Verse 47
\begin{paracol}{2}
\switchcolumn[0]
\dev{%
उत्तरगोले न तदै वाम्यं गोलगे सूत्रे ।\\
शङ्कुतलं शङ्कु हृतं विषुवच्छाया द्विषट्कं गुणायम् ॥ ४७ ॥
}

\switchcolumn[1]
\textbf{V.47.}\quad In the northern hemisphere, it is not thus; the
leftward (deviation) is in the globe-string. The gnomon-base divided
by the gnomon — the equinoctial shadow, twelvefold, is the multiplier.
\end{paracol}

\vspace{0.3cm}
\cnote{[The chapter contains 60 verses on the three problems,
concluding at printed page 193.]}

\longline
\newpage

%% ========================================================================
%% CHAPTER XVI: GRAHANOTTARADHIKARAH
%% Advanced Eclipse Computations
%% ========================================================================
\dsec{अथ ग्रहाणोत्तराध्यायः — षोडशः}
\addcontentsline{toc}{section}{Chapter XVI: Grahāṇottarādhikāraḥ}

\begin{paracol}{2}
\switchcolumn[0]
\dev{%
\textbf{ग्रहाणोत्तराध्यायः}\\[0.2cm]
अस्य अध्यायस्य विषयः ग्रहण-गणितस्य परिष्काराः—चन्द्र-सूर्य-ग्रहण-साधनम्,
परिलेख-विधिः, स्पर्श-मोक्ष-काल-निर्णयः च।
}

\switchcolumn[1]
\textbf{Chapter XVI: On Eclipses (Advanced)}\\[0.2cm]
This chapter treats refined computations for eclipses
(\iast{grahaṇa}): the computation of lunar and solar eclipses,
the graphical \iast{parilekha} method, and the determination of
contact (\iast{sparśa}) and release (\iast{mokṣa}) times.
\end{paracol}

\vspace{0.3cm}
\cnote{[Printed pages 194–218. Chapter 16 covers advanced eclipse
computations, including the \iast{parilekha} graphical method.]}
\vspace{0.3cm}

%% Verse 15
\begin{paracol}{2}
\switchcolumn[0]
\dev{%
परिलेखं ग्रहस्य ग्रहकल्पेन पूर्वं वक्रतो वा ।\\
तात्कालिकसंस्थानं निमीलितोन्मीलितं चैवम् ॥ १५ ॥
}

\switchcolumn[1]
\textbf{V.15.}\quad The \iast{parilekha} (graphical representation)
of the planet, by the planet-rule (\iast{graha-kalpa}) first, either
curved or straight — the instantaneous position, the closing and
opening (of the eclipse) — thus.
\end{paracol}

\vspace{0.3cm}

%% Verse 16
\begin{paracol}{2}
\switchcolumn[0]
\dev{%
विषेषपुण्यांशज्या मानेनच्छाया द्वितच्चापांशाः ।\\
प्रान्तयोर्यथादिशमक्रयदर्शो विपर्ययस्तथा ॥ १६ ॥
}

\switchcolumn[1]
\textbf{V.16.}\quad The sine of the particular (longitude)
difference in minutes (\iast{viśeṣa-puṇy-aṃśa-jyā}), the shadow by
the measure, and the degrees of the second arc — at the edges
according to direction, direct or reversed (\iast{viparyayaḥ}).
\end{paracol}

\vspace{0.3cm}

%% Verse 20
\begin{paracol}{2}
\switchcolumn[0]
\dev{%
तद्वलनांशकयोमेतरं जीवप्राच्छुमानवलम्बचात् ।\\
त्रिज्यालब्धच्छाया प्रग्रहक्षेपं प्राग्वक्रतोः ॥ २० ॥
}

\switchcolumn[1]
\textbf{V.20.}\quad The difference of those deflection-minutes
(\iast{valana-aṃśaka}), from the sine-east-measure and perpendicular
— the shadow obtained by the radius (\iast{trijyā}) — the
\iast{pragraha-kṣepa}, before the retrograde motion.
\end{paracol}

\vspace{0.3cm}

%% Verse 21
\begin{paracol}{2}
\switchcolumn[0]
\dev{%
हृतया व्यासाङ्केनाच्चन्द्रनाताङ्कलिप्तका गुणाया ।\\
मध्यबललिप्तया दक्षिणोत्तरा दिगतबामयात् ॥ २१ ॥
}

\switchcolumn[1]
\textbf{V.21.}\quad Divided by the semi-diameter
(\iast{vyāsāṅka}), the minutes of the lunar-node-difference are the
multiplier. By the middle-strength-minutes (\iast{madhyabala-lip})
the north-south direction goes to the left.
\end{paracol}

\vspace{0.5cm}
\cnote{[This chapter contains 46–47 verses on eclipse computations.]}

\longline
\newpage

%% ========================================================================
%% CHAPTER XVII: SHRINGONNATYUTTARADHIKARAH
%% Lunar Cusp Elevation
%% ========================================================================
\dsec{अथ शृङ्गोन्नत्युत्तराध्यायः — सप्तदशः}
\addcontentsline{toc}{section}{Chapter XVII: Śṛṅgonnatyuttarādhikāraḥ}

\begin{paracol}{2}
\switchcolumn[0]
\dev{%
\textbf{शृङ्गोन्नत्युत्तराध्यायः}\\[0.2cm]
अस्य अध्यायस्य विषयः शशिनः शृङ्गोन्नतिः—चन्द्र-शिखराणाम्
उन्नतिः, तज्-ज्ञानेन काल-निर्णयः च। एको दश-श्लोकीयः
अध्यायः।
}

\switchcolumn[1]
\textbf{Chapter XVII: On Lunar Cusp Elevation}\\[0.2cm]
This chapter treats the elevation of the lunar horns
(\iast{śṛṅgonnati})—the elevation of the moon's cusps, and the
determination of time from this knowledge. This is a short chapter
of ten verses.
\end{paracol}

\vspace{0.5cm}
\cnote{[Printed pages 219–223. Chapter 17 deals with the lunar cusp
(śṛṅga-unnati) and related computations. Only 10 verses.]}
\vspace{0.5cm}

\begin{paracol}{2}
\switchcolumn[0]
\dev{%
[अध्यायस्य श्लोकाः सम्पूर्णतया न प्राप्ताः।\\
अस्य अध्यायस्य पाठः पूर्व-खण्ड-सन्धौ स्थितः।]
}

\switchcolumn[1]
\textit{[The complete verses of this chapter are not fully preserved
in the scanned pages of this chunk. The chapter is transitional
between the astronomical and mathematical sections.]}
\end{paracol}

\vspace{0.5cm}
\cnote{[This short chapter of 10 verses completes the astronomical
portion before the extensive mathematical chapters.]}

\longline
\newpage

%% ========================================================================
%% CHAPTER XVIII: KUTTAKADHYAYAH — THE ALGEBRA CHAPTER
%% ========================================================================
\dsec{अथ कुट्टकाध्यायः — अष्टादशः}
\addcontentsline{toc}{section}{Chapter XVIII: Kuṭṭakādhyāyaḥ — The Algebra Chapter}

\begin{paracol}{2}
\switchcolumn[0]
\dev{%
\textbf{कुट्टकाध्यायः}\\[0.2cm]
अयं गणिताध्यायः अत्यन्तं महत्त्वपूर्णः। अत्र कुट्टक-विधिः,
ऋण-धन-कर्म, अव्यक्त-गणितम्, सङ्क्रमणम्, एकवर्ण-समीकरणम्,
अनेकवर्ण-समीकरणम्, भावितकम्, वर्गप्रकृतिः च उच्यते।
अष्टोत्तर-शत-श्लोकीयः अयम् अध्यायः।
}

\switchcolumn[1]
\textbf{Chapter XVIII: The Kuṭṭaka (Algebra) Chapter}\\[0.2cm]
This is the most important mathematical chapter of the BSS,
containing 102 verses. It covers the \iast{kuṭṭaka} (pulverizer)
method, operations with positive and negative numbers,
\iast{avyakta} (algebra) computations, \iast{saṅkramaṇa}
(transposition), linear and quadratic equations in one unknown,
multi-variable equations, the \iast{bhāvitaka} (product-nature)
equation, and the \iast{varga-prakṛti} (Pell's equation).
\end{paracol}

\vspace{0.3cm}
\cnote{[Printed pages 224–283. The Sanskrit text below is from the GRETIL
digitization (Hayashi 1993).]}
\vspace{0.5cm}

%% ------------------------------------------------------------------------
%% KUTTAKA (PULVERIZER)
%% ------------------------------------------------------------------------
\subsection*{\dev{कुट्टकः — The Pulverizer}}
\addcontentsline{toc}{subsection}{The Kuṭṭaka (Pulverizer)}

\vspace{0.3cm}

%% Verse 1
\begin{paracol}{2}
\switchcolumn[0]
\dev{%
प्रायेण यतः प्रश्नाः कुट्टाकारात् ऋते न शक्यन्ते ।\\
ज्ञातुं वक्ष्यामि ततः कुट्टाकारं सह प्रश्नैः ॥ १ ॥
}

\switchcolumn[1]
\textbf{V.1.}\quad Since problems generally cannot be solved without
the pulverizer (\iast{kuṭṭaka}), I shall state the
\iast{kuṭṭaka} together with problems.
\end{paracol}

\vspace{0.3cm}

%% Verse 2
\begin{paracol}{2}
\switchcolumn[0]
\dev{%
कुट्टक-ख-ऋण-धन-अव्यक्त-मध्य-हरण-एक-वर्ण-भावितकैः ।\\
आचार्यस् तन्त्र-विदां ज्ञातैः वर्ग-प्रकृत्या च ॥ २ ॥
}

\switchcolumn[1]
\textbf{V.2.}\quad The teachers (\iast{ācārya}s) of the experts of
the \iast{tantra}, who know the pulverizer, zero (\iast{kha}),
debt (negative, \iast{ṛṇa}) and wealth (positive, \iast{dhana}),
\iast{avyakta} (unknown), elimination (\iast{haraṇa}),
one-\iast{varṇa} and \iast{bhāvitaka} equations, and the
\iast{varga-prakṛti} (quadratic nature, i.e. Pell's equation).
\end{paracol}

\vspace{0.3cm}

%% Verse 3
\begin{paracol}{2}
\switchcolumn[0]
\dev{%
अधिक-अग्र-भाग-हारात् ऊन-अग्र-छेद-भाजितात् शेषम् ।\\
यत् तत् परस्पर-हृतं लब्धम् अधः अधः पृथक् स्थाप्यम् ॥ ३ ॥
}

\switchcolumn[1]
\textbf{V.3.}\quad The remainder from the divisor (\iast{hāra})
corresponding to the greater residue (\iast{adhikāgra}), divided by
the divisor (\iast{cheda}) corresponding to the lesser residue
(\iast{ūnāgra}) — whatever is the quotient from the mutual division
is to be placed separately, downwards (\iast{adhaḥ}), one below the
other.

\medskip
\textit{[This is the initial step of the Euclidean algorithm applied
to the divisors $A$ and $B$. The quotients are written in a column.]}
\end{paracol}

\vspace{0.3cm}

%% Verse 4
\begin{paracol}{2}
\switchcolumn[0]
\dev{%
शेषं तथा इष्ट-गुणितं यथा अग्रयोः अन्तरेण संयुक्तम् ।\\
शुध्यति गुणकः स्थाप्यः लब्धं च अन्त्यात् उपान्त्य-गुणः ॥ ४ ॥
}

\switchcolumn[1]
\textbf{V.4.}\quad The remainder, multiplied by an assumed number
(\iast{iṣṭa}) such that, when combined with the difference of the
residues (\iast{agra-antara}), it is exactly divisible (\iast{śudhyati})
— the multiplier (\iast{guṇakaḥ}) is to be set down; and the
quotient is the product of the last (\iast{antya}) and the
penultimate (\iast{upāntya}).

\medskip
\textit{[If $r$ is the final remainder, find $m$ such that $rm + (a-b)$
is divisible. Then the \iast{guṇaka} and \iast{labdha} are computed
from the column of quotients by back-substitution.]}
\end{paracol}

\vspace{0.3cm}

%% Verse 5
\begin{paracol}{2}
\switchcolumn[0]
\dev{%
स्व-ऊर्ध्वस् अन्त्य-युतः अग्र-अन्तः हीन-अग्र-छेद-भाजितः शेषम् ।\\
अधिक-अग्र-छेद-हतम् अधिक-अग्र-युतम् भवति अग्रम् ॥ ५ ॥
}

\switchcolumn[1]
\textbf{V.5.}\quad The upper one (\iast{ūrdhva}), combined with the
last (\iast{antya-yutaḥ}), (is divided by) the divisor of the lesser
residue (\iast{hīnāgra-cheda}); the remainder, multiplied by the
divisor of the greater residue (\iast{adhikāgra-cheda}) and combined
with the greater residue, becomes the residue (\iast{agram}).

\medskip
\textit{[The solution $x$ to $ax + c = by + d$ (modulo the divisors)
is constructed by this back-substitution. Let $x_0$ be the
\iast{guṇaka}; then $N = x_0 \cdot B_{\text{large}} + r_{\text{large}}$.]}
\end{paracol}

\vspace{0.3cm}

%% Verse 6
\begin{paracol}{2}
\switchcolumn[0]
\dev{%
छेद-वधस्य द्वि-युगं छेद-वधः युग-गतम् द्वयोः अग्रम् ।\\
कुट्टाकारेण एवम् त्रि-आदि-ग्रह-युग-गत-आनयनम् ॥ ६ ॥
}

\switchcolumn[1]
\textbf{V.6.}\quad Twice the \iast{yuga} (epoch period) is the
product of the divisors; the product of the divisors is the
\iast{yuga-gata} for both residues. Thus by the \iast{kuṭṭaka},
the derivation of three etc.
\end{paracol}

\vspace{0.3cm}

%% Verse 7
\begin{paracol}{2}
\switchcolumn[0]
\dev{%
भ-गण-आदि-शेषम् अग्रम् छेद-हृतम् खम् च दिन-ज-शेष-हृतम् ।\\
अनयोः अग्रम् भ-गण-आदि-दिन-ज-शेष-उद्धृतम् द्यु-गणः ॥ ७ ॥
}

\switchcolumn[1]
\textbf{V.7.}\quad The residue of the revolutions (\iast{bhagaṇa})
etc.
is the \iast{agra}; divided by the divisor; and zero, divided by
the residue of days — of these two, the residue of revolutions etc.,
divided by the day-born residue — that is the day-count
(\iast{dyugaṇa}).
\end{paracol}

\vspace{0.3cm}

%% Verse 8
\begin{paracol}{2}
\switchcolumn[0]
\dev{%
दिन-ज-भ-गण-आदि-शेषं येन गुणम् मण्डल-आदि शेषकयोः ।\\
स-दृश-छेद-उद्धृतयोः तद्-घातम् अहर्-गण-आद्यम् अतः ॥ ८ ॥
}

\switchcolumn[1]
\textbf{V.8.}\quad By whatever multiplier the day-born and
revolution-etc. residues of the circle etc., divided by their
common divisor, are equal — the product of that is the day-count
etc.
\end{paracol}

\vspace{0.3cm}

%% Verse 9
\begin{paracol}{2}
\switchcolumn[0]
\dev{%
हृतयोः परस्परम् यत् छेषम् गुण-कार-भाग-हारकयोः ।\\
तेन हृतौ निस्-छेदौ तौ एव परस्परम् हृतयोः ॥ ९ ॥
}

\switchcolumn[1]
\textbf{V.9.}\quad The mutual division of the dividend and divisor,
whatever is the remainder — by that, divided, they become
reduced (\iast{nis-chedau}); those same, mutually divided.
\end{paracol}

\vspace{0.3cm}

%% Verse 10
\begin{paracol}{2}
\switchcolumn[0]
\dev{%
लब्धम् अधः अधः स्थाप्यं तथा इष्ट-गुण-कार-सङ्गुणं शेषम् ।\\
शुध्यति यथा एक-हीनम् गुणकः स्थाप्यः फलं च अन्त्यात् ॥ १० ॥
}

\switchcolumn[1]
\textbf{V.10.}\quad The quotient is to be placed down, down; the
remainder multiplied by the chosen multiplier
(\iast{iṣṭa-guṇakāra-saṅguṇaṃ}) — it is exactly divisible so that
one is subtracted (\iast{eka-hīnam}). The multiplier is to be set
down; and the result is from the last.
\end{paracol}

\vspace{0.3cm}

%% Verse 11
\begin{paracol}{2}
\switchcolumn[0]
\dev{%
अग्र-अन्तम् उपान्त्येन स्व-ऊर्ध्वः गुणितः अन्त्य-संयुतः भक्तम् ।\\
निस्-शेष-भाग-हारेण एवम् स्थिर-कुट्टकः शेषम् ॥ ११ ॥
}

\switchcolumn[1]
\textbf{V.11.}\quad The end-residue (\iast{agrāntam}), by the
penultimate multiplied, combined with its upper (value), divided by
the divisor without remainder (\iast{niśśeṣa-bhāgahāreṇa}) — thus
is the \iast{sthira-kuṭṭakaḥ} (fixed pulverizer), the remainder.
\end{paracol}

\vspace{0.3cm}

%% Verse 12
\begin{paracol}{2}
\switchcolumn[0]
\dev{%
इष्ट-भ-गण-आदि-शेषात् स्व-कुट्टक-गुणात् स्व-भाग-हार-गुणात् ।\\
शेषम् द्यु-गणः गत-निस्-अपवर्त-गुण-भाग-हार-युतः ॥ १२ ॥
}

\switchcolumn[1]
\textbf{V.12.}\quad From the desired revolution-etc. residue, by its
own \iast{kuṭṭaka}-multiplier, by its own divisor-multiplied — the
remainder is the day-count, combined with the reduced multiplier
and divisor.
\end{paracol}

\vspace{0.3cm}

%% Verse 13
\begin{paracol}{2}
\switchcolumn[0]
\dev{%
एवम् समेषु विषमेषु ऋणम् धनम् धनम् ऋणम् यत् उक्तम् तत् ।\\
ऋण-धनयोः व्यस्तत्वम् गुण्य-प्रक्षेपयोः कार्यम् ॥ १३ ॥
}

\switchcolumn[1]
\textbf{V.13.}\quad Thus in equal (signs) and in unequal: debt
(\iast{ṛṇa}) and wealth (\iast{dhana}), wealth and debt — whatever
was stated, the reversal (\iast{vyastatvam}) of debt and wealth is
to be done to the dividend (\iast{guṇya}) and additive
(\iast{prakṣepa}).
\end{paracol}

\vspace{0.3cm}

%% Verse 14
\begin{paracol}{2}
\switchcolumn[0]
\dev{%
गुणकः छेदः छेदः गुणकः धनम् ऋणम् ऋणम् धनम् कार्यम् ।\\
वर्गम् पदम् पदम् कृतिः अन्त्यात् विपरीतम् आद्यम् तत् ॥ १४ ॥
}

\switchcolumn[1]
\textbf{V.14.}\quad The multiplier is the divisor; the divisor is
the multiplier; wealth is debt, debt is wealth — to be done.
Square (\iast{varga}) is root (\iast{pada}); root is square; from
the end, reversed (\iast{viparītam}); that is the first.
\end{paracol}

\vspace{0.3cm}

%% Verse 15
\begin{paracol}{2}
\switchcolumn[0]
\dev{%
यः जानाति युग-आदि ग्रह-युग-यातैः पृथक् पृथक् कथितैः ।\\
द्वि-त्रि-चतुर्-प्रभृतीनाम् कुट्टाकारम् स जानाति ॥ १५ ॥
}

\switchcolumn[1]
\textbf{V.15.}\quad He who knows the \iast{kuṭṭaka} by the
epoch-periods (\iast{yu-gādi}) of planets, separately stated by
two, three, four etc. — \textbf{he} truly knows the
\iast{kuṭṭaka}.
\end{paracol}

\vspace{0.5cm}
\cnote{[Verses 16–29 contain illustrative problems (\iast{praśna}s) on
the pulverizer applied to astronomical computations.]}
\vspace{0.5cm}

%% Verse 16
\begin{paracol}{2}
\switchcolumn[0]
\dev{%
भ-गण-आद्यम् इष्ट-शेषम् कदा इन्दु-दिवसे रवेः गुरु-दिने वा ।\\
ज्ञ-दिने राशीन् कथयति कुट्टाकारम् स जानाति ॥ १६ ॥
}

\switchcolumn[1]
\textbf{V.16.}\quad ``The residue of revolutions etc.
at the desired — when on a moon-day, on a sun-day, or a Jupiter-day,
on a weekday — tell the signs.'' He who tells this knows the
\iast{kuṭṭaka}.
\end{paracol}

\vspace{0.3cm}

%% Verse 17
\begin{paracol}{2}
\switchcolumn[0]
\dev{%
ज्ञ-दिने यत् अंश-शेषम् विकलाः-शेषम् कदा तत् इन्दु-दिने ।\\
भानोः अथ वा शशिनः यः कथयति कुट्टक-ज्ञः सः ॥ १७ ॥
}

\switchcolumn[1]
\textbf{V.17.}\quad ``On a weekday, whatever is the residue of
degrees, the residue of minutes — when that on a moon-day, of
the sun, or of the moon — he who tells this is a knower of the
\iast{kuṭṭaka}.'
\end{paracol}

\vspace{0.3cm}

%% Verse 18
\begin{paracol}{2}
\switchcolumn[0]
\dev{%
तिथि-मान-दिनेषु इष्टाः ये अर्क-आद्याः ते पुनः कदा तेषु ।\\
इष्ट-ग्रह-वारेषु यः कथयति कुट्टक-ज्ञः सः ॥ १८ ॥
}

\switchcolumn[1]
\textbf{V.18.}\quad ``The desired (planetary positions) in the
tithi-measure-days, which are sun etc., again when in those
desired planet-weekdays'' — he who tells this is a knower of the
\iast{kuṭṭaka}.
\end{paracol}

\vspace{0.5cm}
\cnote{[Verses 19–29 continue with astronomical \iast{praśna}s: given
planetary residues at a particular weekday, finding the elapsed
days and other planetary positions. These were standard problems
for Indian calendar-makers (\iast{pañcāṅga-kāra}s).]}

\newpage

%% ========================================================================
%% DHANARṆA: POSITIVE AND NEGATIVE NUMBERS
%% ========================================================================
\subsection*{\dev{धनर्णसंकलनम् — Addition of Positive and Negative Numbers}}
\addcontentsline{toc}{subsection}{Operations on Positive and Negative Numbers}

\begin{paracol}{2}
\switchcolumn[0]
\dev{%
\textbf{धनर्णकर्म}\\[0.2cm]
अस्मिन् अंशे ब्रह्मगुप्तस्य ऋण-धन-शून्यानां कर्मसु
सूत्राणि दीयन्ते। एते ऐतिहासिक-महत्त्वपूर्णाः सन्ति—
सङ्केतित-संख्यानां क्रम-गणितस्य प्रथम-लिखित-सूत्राणि।
}

\switchcolumn[1]
\textbf{Operations on Positive and Negative Numbers}\\[0.2cm]
This section contains Brahmagupta's famous rules for arithmetic
operations involving positive (\iast{dhana}), negative
(\iast{ṛṇa}), and zero (\iast{śūnya}, \iast{kha}, \iast{ākāśa}).
These are the \textbf{earliest explicit rules} for operations
with signed numbers and zero in the history of mathematics.
\end{paracol}

\vspace{0.5cm}

%% Verse 30
\begin{paracol}{2}
\switchcolumn[0]
\dev{%
धनयोः धनम् ऋणम् ऋणयोः धन-ऋणयोः अन्तरम् सम-ऐक्यम् खम् ।\\
ऋणम् ऐक्यम् च धनम् ऋण-धन-शून्ययोः शून्ययोः शून्यम् ॥ ३० ॥
}

\switchcolumn[1]
\textbf{V.30.}\quad The sum of two positives is positive; of two
negatives is negative; of a positive and a negative is their
difference. Zero is the sum of a positive and an equal negative;
and of two zeros.

\medskip
\textit{[In modern notation:]}
\begin{align*}
(+) + (+) &= (+) & (-) + (-) &= (-)\\
(+) + (-) &= \text{their difference} & a + (-a) &= 0\\
0 + 0 &= 0
\end{align*}
\end{paracol}

\vspace{0.3cm}

%% Verse 31
\begin{paracol}{2}
\switchcolumn[0]
\dev{%
ऊनम् अधिकात् विशोध्यम् धनम् धनात् ऋणम् ऋणात् अधिकम् ऊनात् ।\\
व्यस्तम् तद्-अन्तरम् स्यात् ऋणम् धनम् धनम् ऋणम् भवति ॥ ३१ ॥
}

\switchcolumn[1]
\textbf{V.31.}\quad The lesser is to be subtracted from the greater;
the positive from the positive, the greater negative from the lesser
negative — their difference is reversed (\iast{vyastam}). The
negative becomes positive, the positive becomes negative.

\medskip
\textit{[Subtraction rule: $a - b = -(b - a)$ when $b > a$. For
signs: $(+a) - (+b) = -(b-a)$ and $(-a) - (-b) = +(b-a)$ when
$b > a$.]}
\end{paracol}

\vspace{0.3cm}

%% Verse 32
\begin{paracol}{2}
\switchcolumn[0]
\dev{%
शून्य-विहीनम् ऋणम् ऋणम् धनम् धनम् भवति शून्यम् आकाशम् ।\\
शोध्यम् यदा धनम् ऋणात् ऋणम् धनात् वा तदा क्षेप्यम् ॥ ३२ ॥
}

\switchcolumn[1]
\textbf{V.32.}\quad Zero subtracted from debt (\iast{ṛṇa}) is debt;
from wealth (\iast{dhana}) is wealth. Zero (\iast{ākāśa}) — when
positive is subtracted from negative, or negative from positive,
then it is additive (\iast{kṣepya}).

\medskip
\textit{[Sign rules for subtraction with zero:]}
\begin{align*}
a - 0 &= a & 0 - a &= \text{(sign reversed)}
\end{align*}
\end{paracol}

\vspace{0.3cm}

%% Verse 33
\begin{paracol}{2}
\switchcolumn[0]
\dev{%
ऋणम् ऋण-धनयोः घातः धनम् ऋणयोः धन-वधः धनम् भवति ।\\
शून्य-ऋणयोः ख-धनयोः ख-शून्ययोः वा वधः शून्यम् ॥ ३३ ॥
}

\switchcolumn[1]
\textbf{V.33.}\quad The product of two negatives is positive; the
product of two positives is positive. The product of zero and a
negative, of zero and a positive, or of two zeros — the product is
zero.

\medskip
\textit{[Multiplication of signed numbers:]}
\begin{align*}
(-) \times (-) &= (+) & (+) \times (+) &= (+) & (+) \times (-) &= (-)\\
0 \times (-) &= 0 & 0 \times (+) &= 0 & 0 \times 0 &= 0
\end{align*}
\end{paracol}

\vspace{0.3cm}

%% Verse 34
\begin{paracol}{2}
\switchcolumn[0]
\dev{%
धन-भक्तम् धनम् ऋण-हृतम् ऋणम् धनम् भवति खम् ख-भक्तम् खम् ।\\
भक्तम् ऋणेन धनम् ऋणम् धनेन हृतम् ऋणम् ऋणम् भवति ॥ ३४ ॥
}

\switchcolumn[1]
\textbf{V.34.}\quad Positive divided by positive is positive;
negative divided by negative is positive. Zero divided by zero
is zero. Positive divided by negative is negative; negative
divided by positive is negative.

\medskip
\textit{[Division of signed numbers:]}
\begin{align*}
\frac{(+)}{(+)} &= (+) & \frac{(-)}{(-)} &= (+) & \frac{0}{0} &= 0\\
\frac{(+)}{(-)} &= (-) & \frac{(-)}{(+)} &= (-)
\end{align*}
\end{paracol}

\vspace{0.5cm}
\cnote{[Verses 30–34 are historically momentous. Brahmagupta is the first
author in world history to state general rules for arithmetic
operations with signed numbers and zero. Note that $\frac{0}{0}=0$
is the only questionable rule from a modern perspective; Brahmagupta
also stated (elsewhere) that division by zero leaves the fraction
unchanged, i.e. $\frac{a}{0} = \frac{a}{0}$ as an unreduced entity.]}

\newpage

%% ========================================================================
%% SANKRAMANAM (ALGEBRAIC IDENTITIES)
%% ========================================================================
\subsection*{\dev{सङ्क्रमणम् — Transposition}}
\addcontentsline{toc}{subsection}{Saṅkramaṇam}

\vspace{0.3cm}

%% Verse 36
\begin{paracol}{2}
\switchcolumn[0]
\dev{%
योगस् अन्तर-युत-हीनः द्वि-हृतः सङ्क्रमणम् अन्तर-विभक्तम् वा ।\\
वर्ग-अन्तरम् अन्तर-युत-हीनम् द्वि-हृतम् विषम-कर्म ॥ ३६ ॥
}

\switchcolumn[1]
\textbf{V.36.}\quad The sum, decreased by the difference and divided
by two, is the \iast{saṅkramaṇa} (transposition); or divided by the
difference. The difference of squares, decreased by the sum of the
difference and divided by two — the \iast{viṣama-karma}
(unequal operation).

\medskip
\textit{[If $x + y = S$ and $x - y = D$, then:]}
\begin{align*}
x = \frac{S + D}{2},\quad y = \frac{S - D}{2}
\end{align*}
\textit{The ``difference of squares'' formula: $x^2 - y^2 = S \cdot D$.}
\end{paracol}

\vspace{0.5cm}
\cnote{[Verses 37–42 treat \iast{karaṇī} (surds), \iast{bhāvitaka}
(product-form), and algebraic identities for powers of unknowns.
These are not fully preserved in the scanned pages of this chunk.]}

\newpage

%% ========================================================================
%% EKA-VARṆA-SAMĪKARAṆAM (LINEAR EQUATIONS)
%% ========================================================================
\subsection*{\dev{एकवर्णसमीकरणम् — Linear Equations in One Unknown}}
\addcontentsline{toc}{subsection}{Eka-varṇa-samīkaraṇam}

\vspace{0.3cm}

%% Verse 43
\begin{paracol}{2}
\switchcolumn[0]
\dev{%
अव्यक्त-अन्तर-भक्तम् व्यस्तम् रूप-अन्तरम् समे अव्यक्तः ।\\
वर्ग-अव्यक्ताः शोध्याः यस्मात् रूपाणि तद्-अधस्तात् ॥ ४३ ॥
}

\switchcolumn[1]
\textbf{V.43.}\quad Divided by the difference of unknowns
(\iast{avyakta}), reversed, the difference of \iast{rūpa}s (absolute
terms) — in the equation (\iast{sama}), the unknowns (\iast{avyakta})
are to be subtracted; the \iast{rūpa}s are placed below that.

\medskip
\textit{[For equations $ax + b = cx + d$, divide by $(a-c)$ and rearrange:]}
\begin{align*}
x = \frac{d - b}{a - c}
\end{align*}
\end{paracol}

\vspace{0.3cm}

%% Verse 44 — THE QUADRATIC FORMULA
\begin{paracol}{2}
\switchcolumn[0]
\dev{%
वर्ग-चतुर्-गुणितानाम् रूपाणाम् मध्य-वर्ग-सहितानाम् ।\\
मूलम् नध्येन ऊनम् वर्ग-द्वि-गुण-उद्धृतम् मध्यः ॥ ४४ ॥
}

\switchcolumn[1]
\textbf{V.44.}\quad Of the \iast{rūpa}s multiplied by four and the
square of the middle (\iast{madhya}) — the square-root of that,
decreased by the middle and divided by twice the multiplier
(\iast{varga-dvi-guṇa-uddhṛtam}) — is the \iast{madhya} (unknown).

\medskip
\textit{[This is Brahmagupta's verbal statement of the quadratic formula.
If $ax^2 + bx = c$:]}
\begin{align*}
x = \frac{\sqrt{4ac + b^2} - b}{2a}
\end{align*}
\textit{This is algebraically equivalent to the modern formula
$x = \frac{-b \pm \sqrt{b^2 + 4ac}}{2a}$ for the equation
$ax^2 + bx - c = 0$. Brahmagupta always writes the constant term
on the right side, so his formula uses $+4ac$.}
\end{paracol}

\vspace{0.3cm}

%% Verse 45
\begin{paracol}{2}
\switchcolumn[0]
\dev{%
वर्ग-आहत-रूपाणाम् अव्यक्त-अर्ध-कृति-संयुतानाम् यत् ।\\
पदम् अव्यक्त-अर्ध-ऊनम् तत् वर्ग-विभक्तम् अव्यक्तः ॥ ४५ ॥
}

\switchcolumn[1]
\textbf{V.45.}\quad Of the \iast{rūpa}s multiplied by the square
and combined with the square of half the unknown — whatever is the
root, decreased by half the unknown, divided by the square — that is
the unknown.

\medskip
\textit{[Alternative form:]}
\begin{align*}
x = \frac{\sqrt{a^2c + (\frac{b}{2})^2} - \frac{b}{2}}{a}
\end{align*}
\end{paracol}

\vspace{0.5cm}
\cnote{[Verses 46–50 contain astronomical problems solved by the
\iast{kuṭṭaka} and quadratic methods.]}
\vspace{0.5cm}

%% Verse 46
\begin{paracol}{2}
\switchcolumn[0]
\dev{%
स-एकात् अंशक-शेषात् द्वादश-भागः चतुर्-गुणः अष्ट-युतः ।\\
स-एक-अंश-शेष-तुल्यः यदा तदा अहर्-गणम् कथय ॥ ४६ ॥
}

\switchcolumn[1]
\textbf{V.46.}\quad ``From the degree-residue (\iast{aṃśaka-śeṣa})
with one added, divided by twelve, multiplied by four and increased
by eight, equals the degree-residue with one — when this happens,
tell the day-count (\iast{ahargaṇa}).''

\medskip
\textit{[An astronomical problem leading to a linear equation in
the \iast{ahargaṇa} (elapsed days).]}
\end{paracol}

\vspace{0.3cm}

%% Verse 47
\begin{paracol}{2}
\switchcolumn[0]
\dev{%
द्वि-ऊनम् अधिक-मास-शेषम् त्रि-हृतम् सप्त-अधिकम् द्वि-सङ्गुणितम् ।\\
अधिमास-शेष-तुल्यम् यदा तदा युग-गतम् कथय ॥ ४७ ॥
}

\switchcolumn[1]
\textbf{V.47.}\quad ``The \iast{adhimāsa-śeṣa} (excess month residue),
decreased by two, divided by three, increased by seven, multiplied
by two — equals the \iast{adhimāsa-śeṣa}. When this happens, tell
the \iast{yuga-gata} (elapsed in the \iast{yuga}).''
\end{paracol}

\vspace{0.3cm}

%% Verse 48
\begin{paracol}{2}
\switchcolumn[0]
\dev{%
वि-एकम् अवम-अवशेषम् षष्-उद्धृतम् त्रि-युतम् अवम-शेषस्य ।\\
पञ्च-विभक्तस्य समम् यदा तदा युग-गतम् कथय ॥ ४८ ॥
}

\switchcolumn[1]
\textbf{V.48.}\quad ``The \iast{avama-avaśeṣa} (omitted day residue),
decreased by one, divided by six, increased by three — equals the
\iast{avama-śeṣa} divided by five. When this happens, tell the
\iast{yuga-gata}.''
\end{paracol}

\newpage

%% ========================================================================
%% ANEKA-VARṆA SAMĪKARAṆAM (MULTI-VARIABLE EQUATIONS)
%% ========================================================================
\subsection*{\dev{अनेकवर्णसमीकरणम् — Equations with Multiple Unknowns}}
\addcontentsline{toc}{subsection}{Aneka-varṇa-samīkaraṇam}

\vspace{0.3cm}

%% Verse 51
\begin{paracol}{2}
\switchcolumn[0]
\dev{%
आद्यात् वर्णात् अन्यान् वर्णान् प्रोह्य आद्य-मानम् आड्य-हृतम् ।\\
सदृश-छेदौ असकृत् द्वौ व्यस्तौ कुट्टकः बहुषु ॥ ५१ ॥
}

\switchcolumn[1]
\textbf{V.51.}\quad From the first unknown (\iast{varṇa}), having
subtracted the other unknowns, divided by the first measure
(\iast{ādya-māna}), combined — two similar divisors, repeatedly,
reversed — the \iast{kuṭṭaka} (is applied) for many (unknowns).

\medskip
\textit{[For multiple equations in many unknowns: eliminate all but
one unknown, then apply the \iast{kuṭṭaka}.]}
\end{paracol}

\vspace{0.3cm}

%% Verse 52
\begin{paracol}{2}
\switchcolumn[0]
\dev{%
गत-भगण-युतात् द्यु-गणात् तत्-शेष-युतात् तद्-ऐक्य-संयुक्तात् ।\\
तद्-योगात् द्यु-गणम् वा यः कथयति कुट्टक-ज्ञः सः ॥ ५२ ॥
}

\switchcolumn[1]
\textbf{V.52.}\quad ``From the day-count increased by the elapsed
revolutions, from that with its residue added, combined with their
sum, or from their sum — the day-count'' — he who tells this is a
\iast{kuṭṭaka}-knower.
\end{paracol}

\vspace{0.3cm}

%% Verse 53
\begin{paracol}{2}
\switchcolumn[0]
\dev{%
गत-भगण-ऊनात् द्यु-गणात् तत्-शेष-ऊनात् तद्-ऐक्य-हीनात् वा ।\\
तद्-विवरात् द्यु-गणम् वा यः कथयति कुट्टक-ज्ञः सः ॥ ५३ ॥
}

\switchcolumn[1]
\textbf{V.53.}\quad ``From the day-count decreased by the elapsed
revolutions, from that with its residue subtracted, diminished by
their sum, or from their difference — the day-count'' — he who
tells this is a \iast{kuṭṭaka}-knower.
\end{paracol}

\vspace{0.3cm}

%% Verse 54
\begin{paracol}{2}
\switchcolumn[0]
\dev{%
राशि-आद्यैः तत्-शेषैः च एवम् भुक्त-अधिमास-दिन-हीनैः ।\\
तत्-शेषैः च युग-गतम् यः कथयति कुट्टक-ज्ञः सः ॥ ५४ ॥
}

\switchcolumn[1]
\textbf{V.54.}\quad ``With the quantities etc., with their residues,
thus with the elapsed intercalary months and days subtracted, and
with their residues — the \iast{yuga-gata}'' — he who tells this is
a \iast{kuṭṭaka}-knower.
\end{paracol}

\vspace{0.5cm}
\cnote{[Verses 55–59 contain more astronomical applications and
computational rules for calendar calculations.]}

\newpage

%% ========================================================================
%% BHĀVITAKAM (PRODUCT NATURE)
%% ========================================================================
\subsection*{\dev{भावितकम् — The Product-Nature Equation}}
\addcontentsline{toc}{subsection}{Bhāvitakam}

\vspace{0.3cm}

%% Verse 60
\begin{paracol}{2}
\switchcolumn[0]
\dev{%
भावितक-रूप-गुणना स-अव्यक्त-वधा इष्ट-भाजिता इष्ट-आप्त्योः ।\\
अल्पे अधिकः अधिके अल्पः क्षेप्यः भावित-हृतौ व्यस्तम् ॥ ६० ॥
}

\switchcolumn[1]
\textbf{V.60.}\quad The \iast{bhāvitaka} (product-nature) multiplied
by the \iast{rūpa}, together with the product of the unknown
(\iast{avyakta-vadhā}), divided by the desired (\iast{iṣṭa}) — at
the desired quotient: if less, add more; if more, add less;
\iast{kṣepya} (additive) to the \iast{bhāvita} divisor; reversed.

\medskip
\textit{[The \iast{bhāvitaka} is the coefficient of the product of
two unknowns, i.e. the term $Bxy$ in the general quadratic form
$Ax^2 + Bxy + Cy^2 + Dx + Ey + F = 0$.]}
\end{paracol}

\vspace{0.3cm}

%% Verse 61
\begin{paracol}{2}
\switchcolumn[0]
\dev{%
भानोः राशि-अंश-वधात् त्रि-चतुर्-गुणितान् विशोध्या राशि-अंशान् ।\\
नवतिम् दृष्ट्वा सूर्यम् कुर्वन् आ वत्सरात् गणकः ॥ ६१ ॥
}

\switchcolumn[1]
\textbf{V.61.}\quad From the product of the sun's sign and degrees,
having subtracted the signs and degrees multiplied by three and four,
having seen ninety — making the sun (\iast{sūrya}), the calculator
(does this) up to a year.
\end{paracol}

\vspace{0.3cm}

%% Verse 62
\begin{paracol}{2}
\switchcolumn[0]
\dev{%
भावितके यद्-घातः विनष्ट-वर्णेन तत्-प्रमाणानि ।\\
कृत्वा इष्टानि तद्-आहत-वर्ण-ऐक्यम् भवति रूपाणि ॥ ६२ ॥
}

\switchcolumn[1]
\textbf{V.62.}\quad In the \iast{bhāvitaka}, the product by which
unknown — having made that amount desired, the sum of the unknown
(\iast{varṇa}) multiplied by that becomes the \iast{rūpa}s (known
quantities).
\end{paracol}

\vspace{0.3cm}

%% Verse 63
\begin{paracol}{2}
\switchcolumn[0]
\dev{%
वर्ण-प्रमाण-भावित-घातः भवति इष्ट-वर्ण-सङ्ख्या एवम् ।\\
सिध्यति विना अपि भावित-सम-करणात् किम् कृतम् तत् अतः ॥ ६३ ॥
}

\switchcolumn[1]
\textbf{V.63.}\quad The product of the unknown-quantity and the
\iast{bhāvitaka} becomes the desired unknown-number. Thus it is
accomplished even without the \iast{bhāvitaka} equal-equation;
what was done from that?
\end{paracol}

\vspace{0.3cm}

%% Verse 64
\begin{paracol}{2}
\switchcolumn[0]
\dev{%
मूलम् द्विधा इष्ट-वर्गात् गुणक-गुणात् इष्ट-युत-विहीनात् च ।\\
आद्य-वधः गुणक-गुणः सह अन्त्य-घातेन कृतम् अन्त्यम् ॥ ६४ ॥
}

\switchcolumn[1]
\textbf{V.64.}\quad The root is twofold: from the square of the
desired, from the multiplier multiplied, increased or decreased by
the desired. The first product is the multiplier-product; together
with the last product — made last.

\medskip
\textit{[For equations of the form $N x^2 + k = y^2$ where $N$ is the
\iast{guṇaka} (multiplier) and $k$ is the \iast{kṣepa} (additive).]}
\end{paracol}

\newpage

%% ========================================================================
%% VARGA-PRAKṚTI (PELL'S EQUATION)
%% ========================================================================
\subsection*{\dev{वर्गप्रकृतिः — Pell's Equation}}
\addcontentsline{toc}{subsection}{Varga-prakṛtiḥ (Pell's Equation)}

\begin{paracol}{2}
\switchcolumn[0]
\dev{%
\textbf{वर्गप्रकृतिः}\\[0.2cm]
अस्मिन् अंशे वर्गप्रकृतिः—असम-द्वि-घात-समीकरणस्य—विधिः
उच्यते। यदा गुणक-संख्यया गुणितस्य वर्गे क्षेपेण युते पदम्
labdhyate, तदा वर्गप्रकृतिः।
}

\switchcolumn[1]
\textbf{Varga-prakṛtiḥ (Pell's Equation)}\\[0.2cm]
This section treats the ``square-nature'' or Pell's equation:
finding integers $x, y$ such that
\begin{align*}
Nx^2 + k = y^2
\end{align*}
where $N$ (the \iast{guṇaka}) is a non-square positive integer and
$k$ (the \iast{kṣepa}) is an integer. Brahmagupta gave the
composition formula (\iast{bhāvanā}) and the cyclic method
(\iast{cakravāla}), which were rediscovered in Europe only by
Euler and Lagrange in the 18th century.
\end{paracol}

\vspace{0.5cm}

%% Verse 65
\begin{paracol}{2}
\switchcolumn[0]
\dev{%
वज्र-वध-ऐक्यम् प्रथमम् प्रक्षेपः क्षेप-वध-तुल्यः ।\\
प्रक्षेप-शोधक-हृते मूले प्रक्षेपके रूपे ॥ ६५ ॥
}

\switchcolumn[1]
\textbf{V.65.}\quad The product of the cross-terms (\iast{vajra-vadha})
plus one is the first \iast{prakṣepa}; equal to the product of the
\iast{prakṣepa}s; divided by the \iast{prakṣepa}-subtractor — the
roots (\iast{mūle}) are in the \iast{prakṣepa}s, \iast{rūpa}s.

\medskip
\textit{[The ``cross multiplication'' identity (bhāvanā): If
$(x_1, y_1)$ solves $Nx^2 + k_1 = y^2$ and $(x_2, y_2)$ solves
$Nx^2 + k_2 = y^2$, then $(x_1y_2 \pm x_2y_1, Nx_1x_2 \pm y_1y_2)$
solves $Nx^2 + k_1k_2 = y^2$.]}
\end{paracol}

\vspace{0.3cm}

%% Verse 66
\begin{paracol}{2}
\switchcolumn[0]
\dev{%
रूप-प्रक्षेप-पदे पृथक् इष्ट-क्षेप्य-शोध्य-मूलाभ्याम् ।\\
कृत्वा अन्त्य-आद्य-पदे ये प्रक्षेपे शोधने वा इष्टे ॥ ६६ ॥
}

\switchcolumn[1]
\textbf{V.66.}\quad The \iast{rūpa}-\iast{prakṣepa} roots,
separately, by the desired \iast{kṣepya} and
\iast{śodhya}-roots — having done, the end and first roots which
are in the \iast{prakṣepa} — in subtraction or desired.
\end{paracol}

\vspace{0.3cm}

%% Verse 67
\begin{paracol}{2}
\switchcolumn[0]
\dev{%
चतुर्-अधिके अन्त्य-पद-कृतिः त्रि-ऊना दलिता अन्त्य-पद-गुणा अन्त्य-पदम् ।\\
अन्त्य-पद-कृतिः वि-एका द्वि-हृता आद्य-पद-आहता आद्य-पदम् ॥ ६७ ॥
}

\switchcolumn[1]
\textbf{V.67.}\quad When increased by four (\iast{catur-adhike}):
the square of the last root (\iast{antya-pada-kṛtiḥ}), decreased by
three (\iast{tri-ūnā}), halved, multiplied by the last root — is
the last root. The square of the last root, decreased by one
(\iast{vi-ekā}), divided by two, multiplied by the first root —
is the first root.

\medskip
\textit{[Recurrence relations for generating new solutions to
$Nx^2 + 1 = y^2$ from known solutions. If $(x_1, y_1)$ is a known
solution with $k = -1$, then:]}
\begin{align*}
y_{n+1} &= \frac{(y_n^2 - 3)y_n}{2} \\
x_{n+1} &= \frac{(y_n^2 - 1)x_n}{2}
\end{align*}
\end{paracol}

\vspace{0.3cm}

%% Verse 68
\begin{paracol}{2}
\switchcolumn[0]
\dev{%
चतुर्-ऊने अन्त्य-पद-कृती त्रि-एक-युते वध-दलम् पृथक् वि-एकम् ।\\
वि-एक-आड्य-आहतम् अन्त्यम् पद-वध-गुणम् आद्यम् आन्त्य-पदम् ॥ ६८ ॥
}

\switchcolumn[1]
\textbf{V.68.}\quad When decreased by four: the square of the last
root, increased by three and one, the product halved, separately,
decreased by one — increased by one, decreased by one, multiplied
by the last — the product of the roots (is) the first, the last
root.
\end{paracol}

\vspace{0.3cm}

%% Verse 69
\begin{paracol}{2}
\switchcolumn[0]
\dev{%
वर्गे गुणके क्षेपः केन चित् उद्धृत-युत-ऊनितः दलितः ।\\
प्रथमः अन्त्य-मूलम् अन्यः गुण-कार-पद-उद्धृतः प्रथमः ॥ ६९ ॥
}

\switchcolumn[1]
\textbf{V.69.}\quad In the square of the multiplier
(\iast{guṇaka}), the \iast{kṣepa}, by some (number) divided, added,
subtracted, halved — the first (root). The last root — the other,
divided by the multiplier-root, (gives) the first.

\medskip
\textit{[Given $Nx^2 + k = y^2$, divide by $k$ to reduce the
additive. If $k$ divides $N$ or a related quantity, new roots
can be found.]}
\end{paracol}

\vspace{0.3cm}

%% Verse 70
\begin{paracol}{2}
\switchcolumn[0]
\dev{%
वर्ग-चिन्ने गुणके प्रथमम् तद्-मूल-भाजितम् भवति ।\\
वर्ग-चिन्ने क्षेपे तत्-पद-गुणिते तदा मूले ॥ ७० ॥
}

\switchcolumn[1]
\textbf{V.70.}\quad When the multiplier is a perfect square
(\iast{varga-cinne guṇake}): the first (root), divided by its
root, becomes (the new root). When the \iast{kṣepa} is a perfect
square: multiplied by its root, then (are) the roots.

\medskip
\textit{[If $N = a^2$ is a perfect square, then $a^2x^2 + k = y^2$
gives $y = \sqrt{a^2x^2 + k}$. If also $k = b^2$, then
$\sqrt{(ax)^2 + b^2}$ may simplify.]}
\end{paracol}

\vspace{0.3cm}

%% Verse 71
\begin{paracol}{2}
\switchcolumn[0]
\dev{%
गुणक-युतिः अष्ट-गुणिता गुणक-अन्तर-वर्ग-भाजिता राशिः ।\\
गुणकौ त्रि-गुणौ व्यस्त-अधिकौ हृतौ अन्तरेण पदे ॥ ७१ ॥
}

\switchcolumn[1]
\textbf{V.71.}\quad The sum of the multipliers, multiplied by eight,
divided by the square of the difference of the multipliers — the
quantity. The two multipliers, tripled, reversed, increased, divided
by the difference — (are) the roots.
\end{paracol}

\vspace{0.3cm}

%% Verse 72
\begin{paracol}{2}
\switchcolumn[0]
\dev{%
वर्गः अन्य-कृति-युत-ऊनः तत्-संयोग-अन्तर-अर्ध-कृति-भक्तः ।\\
तद्-गुणितौ युति-वियुतौ वर्गौ घाते च रूप-युते ॥ ७२ ॥
}

\switchcolumn[1]
\textbf{V.72.}\quad The square, decreased by the sum of the other
squares, divided by the half-square of the sum-difference — those
multiplied by them, combined and separated — squares; in the
product, combined with the \iast{rūpa}.

\medskip
\textit{[The general composition (bhāvanā) formula:]}
\begin{align*}
\text{If } Nx_1^2 + k_1 = y_1^2 \text{ and } Nx_2^2 + k_2 = y_2^2,\\
\text{then } N(x_1y_2 \pm x_2y_1)^2 + k_1k_2 = (Nx_1x_2 \pm y_1y_2)^2
\end{align*}
\textit{When $k_1 = k_2 = 1$, this generates infinitely many solutions
to $Nx^2 + 1 = y^2$ from a single fundamental solution. This is
Brahmagupta's great discovery, called the \iast{samāsa-bhāvanā}.}
\end{paracol}

\newpage

%% ========================================================================
%% ASTRONOMICAL APPLICATIONS
%% ========================================================================
\subsection*{\dev{खगोलप्रयोगाः — Astronomical Applications}}
\addcontentsline{toc}{subsection}{Astronomical Applications}

\begin{paracol}{2}
\switchcolumn[0]
\dev{%
\textbf{ग्रह-गणित-प्रयोगाः}\\[0.2cm]
अस्मिन् अंशे दिव्य-सूत्राणि दीयन्ते—यतः कुट्टक-ज्ञः
ग्रहाणाम् अवशेषज्ञानेन अन्येषाम् ग्रह-स्थानानि कथयितुम्
शक्नोति। एते सूत्राणि पञ्चाङ्ग-कर्तृभ्यः
परिशुद्धि-ज्ञानाय आसन्।
}

\switchcolumn[1]
\textbf{Astronomical Divination Rules}\\[0.2cm]
This section gives remarkable ``divination'' rules: given the
planetary residues at a particular day, one can compute the
positions of other planets. These were used by calendar-makers
(\iast{pañcāṅga-kāra}s) to verify their computations.
\end{paracol}

\vspace{0.5cm}

%% Verse 75
\begin{paracol}{2}
\switchcolumn[0]
\dev{%
राशि-कला-शेष-कृतिम् द्विनवति-गुणिताम् त्र्यशीति-गुणिताम् वा ।\\
स-एका ज्ञ-दिने वर्गम् कुर्वन् आ वत्सरात् गणकः ॥ ७५ ॥
}

\switchcolumn[1]
\textbf{V.75.}\quad ``The square of the residue of signs and minutes,
multiplied by twenty-two (\iast{dvinavati}), or multiplied by
eighty-three (\iast{tryaśīti}), with one added — on a weekday,
making the square — the calculator (does this) up to a year.''
\end{paracol}

\vspace{0.3cm}

%% Verse 76
\begin{paracol}{2}
\switchcolumn[0]
\dev{%
सूर्य-विलिप्ता-शेषम् पञ्चभिः ऊन-आहतम् तथा दशभिः ।\\
वर्गे बृहस्पति-दिने कुर्वन् आ वत्सरात् गणकः ॥ ७६ ॥
}

\switchcolumn[1]
\textbf{V.76.}\quad ``The residue of the sun's seconds
(\iast{viluptā-śeṣa}), decreased and multiplied by five, likewise
by ten — on a Jupiter-day, making the square — the calculator (does
this) up to a year.''
\end{paracol}

\vspace{0.3cm}

%% Verse 77
\begin{paracol}{2}
\switchcolumn[0]
\dev{%
भ-गण-आदि-शेष-वर्गम् त्रिभिः गुणम् संयुतम् शतैः नवभिः ।\\
कृतिम् अष्ट-शत-ऊनम् वा कुर्वन् आ वत्सरात् गणकः ॥ ७७ ॥
}

\switchcolumn[1]
\textbf{V.77.}\quad ``The square of the residue of revolutions etc.,
multiplied by three, combined with nine hundred — or the square
decreased by eight hundred — the calculator (does this) up to a
year.''
\end{paracol}

\vspace{0.3cm}

%% Verse 78
\begin{paracol}{2}
\switchcolumn[0]
\dev{%
भ-गण-आदि-शेष-वर्गम् चतुर्-गुणम् पञ्चषष्टि-संयुक्तम् ।\\
षष्टि-ऊनम् वा वर्गम् कुर्वन् आ वत्सरात् गणकः ॥ ७८ ॥
}

\switchcolumn[1]
\textbf{V.78.}\quad ``The square of the residue of revolutions etc.,
multiplied by four, combined with sixty-five — or the square
decreased by sixty — the calculator (does this) up to a year.''
\end{paracol}

\vspace{0.3cm}

%% Verse 79
\begin{paracol}{2}
\switchcolumn[0]
\dev{%
इष्ट-भगण-आदि-शेषम् द्विनवति-ऊनम् त्र्यशीति-सङ्गुणितम् ।\\
रूपेण युतम् वर्गम् कुर्वन् आ वत्सरात् गणकः ॥ ७९ ॥
}

\switchcolumn[1]
\textbf{V.79.}\quad ``The desired residue of revolutions etc.,
decreased by twenty-two, multiplied by eighty-three, combined with
one — making the square — the calculator (does this) up to a year.''
\end{paracol}

\vspace{0.3cm}

%% Verse 80
\begin{paracol}{2}
\switchcolumn[0]
\dev{%
अधिमास-शेष-वर्गम् त्रयोदश-गुणम् त्रिभिः शतैः युक्तम् ।\\
त्रि-घन-ऊनम् वा वर्गम् कुर्वन् आ वत्सरात् गणकः ॥ ८० ॥
}

\switchcolumn[1]
\textbf{V.80.}\quad ``The square of the \iast{adhimāsa-śeṣa}
(intercalary month residue), multiplied by thirteen, combined with
three hundred — or the square decreased by twenty-seven
(\iast{tri-ghana}) — the calculator (does this) up to a year.''
\end{paracol}

\vspace{0.3cm}

%% Verse 81
\begin{paracol}{2}
\switchcolumn[0]
\dev{%
इन्दु-विलिप्ता-शेषम् सप्तदश-गुणम् त्रयोदश-गुणम् च अपि ।\\
पृथक् एक-युतम् वर्गम् कुर्वन् आ वत्सरात् गणकः ॥ ८१ ॥
}

\switchcolumn[1]
\textbf{V.81.}\quad ``The residue of the moon's seconds, multiplied
by seventeen, also multiplied by thirteen — separately, combined
with one — making the square — the calculator (does this) up to a
year.''
\end{paracol}

\vspace{0.3cm}

%% Verse 82
\begin{paracol}{2}
\switchcolumn[0]
\dev{%
अवम-अवशेष-वर्गम् द्वादश-गुणितम् शतेन संयुक्तम् ।\\
तribhis ऊनम् वा वर्गम् कुर्वन् आ वत्सरात् गणकः ॥ ८२ ॥
}

\switchcolumn[1]
\textbf{V.82.}\quad ``The square of the \iast{avama-avaśeṣa}
(omitted-day residue), multiplied by twelve, combined with one
hundred — or the square decreased by three — the calculator (does
this) up to a year.''
\end{paracol}

\vspace{0.3cm}

%% Verse 83
\begin{paracol}{2}
\switchcolumn[0]
\dev{%
ज्ञ-दिने अर्क-कला-शेषम् गुरु-दिन-विकला-अवशेष-युक्त-ऊनम् ।\\
वर्गम् वधम् च स-एकम् कुर्वन् आ वत्सरात् गणकः ॥ ८३ ॥
}

\switchcolumn[1]
\textbf{V.83.}\quad ``On a weekday, the residue of the sun's minutes
(\iast{kala}), combined with or decreased by the residue of
Jupiter-day seconds — the square and product with one — the
calculator (does this) up to a year.''
\end{paracol}

\vspace{0.3cm}

%% Verse 84
\begin{paracol}{2}
\switchcolumn[0]
\dev{%
विकला-शेषम् सहितम् त्रिनवत्या सप्तषष्टि-हीनम् च ।\\
भानोः ज्ञ-दिने वर्गम् कुर्वन् आ वत्सरात् गणकः ॥ ८४ ॥
}

\switchcolumn[1]
\textbf{V.84.}\quad ``The residue of seconds (\iast{vikalā-śeṣa}),
combined with ninety-three (\iast{trinavatī}) and decreased by
seventy-seven (\iast{saptaṣaṣṭi}) — on the sun's weekday, making
the square — the calculator (does this) up to a year.''
\end{paracol}

\vspace{0.3cm}

%% Verse 85
\begin{paracol}{2}
\switchcolumn[0]
\dev{%
ज्ञ-दिने अर्क-कला-शेषम् द्वादशभिः संयुतम् त्रिषष्ट्या च ।\\
षष्ट्या अष्टभिः च ऊनम् कुर्वन् आ वत्सरात् गणकः ॥ ८५ ॥
}

\switchcolumn[1]
\textbf{V.85.}\quad ``On a weekday, the residue of the sun's minutes,
combined with twelve and sixty-three — or decreased by sixty and
eight — the calculator (does this) up to a year.''
\end{paracol}

\vspace{0.3cm}

%% Verse 86
\begin{paracol}{2}
\switchcolumn[0]
\dev{%
इन्दु-विलिप्ता-शेषात् रवि-लिप्ता-शेषम् अंश-शेषम् वा ।\\
अथ वा मध्यम् इष्टम् कुर्वन् आ वत्सरात् गणकः ॥ ८६ ॥
}

\switchcolumn[1]
\textbf{V.86.}\quad ``From the residue of the moon's seconds, the
residue of the sun's seconds, or the residue of degrees — or
alternatively the mean (position) desired — the calculator (does
this) up to a year.''
\end{paracol}

\vspace{0.3cm}

%% Verse 87
\begin{paracol}{2}
\switchcolumn[0]
\dev{%
जीव-विलिप्ता-शेषात् कुजम् इन्दुम् भौम-लिप्तिका-शेषात् ।\\
रविम् इन्दु-भाग-शेषात् कुर्वन् आ वत्सरात् गणकः ॥ ८७ ॥
}

\switchcolumn[1]
\textbf{V.87.}\quad ``From the residue of seconds of Venus
(\iast{jīva}), Mars (\iast{kuja}) and the Moon; from the residue
of Mars's seconds (\iast{bhauma-liptikā}); from the residue of
the moon's degrees — the Sun — the calculator (does this) up to
a year.''
\end{paracol}

\vspace{0.5cm}
\cnote{[Verses 75–87 give a remarkable set of ``divination'' rules.
Given a particular planetary residue at a known weekday, these
formulas allow one to compute other planetary positions. The
constants (22, 83, 5, 10, etc.) encode specific planetary period
relations. These verification rules were essential for ensuring
the accuracy of \iast{pañcāṅga} (calendar) computations.]}

\newpage

%% ========================================================================
%% CLOSING VERSES OF CHAPTER XVIII
%% ========================================================================
\cnote{[Verses 88–99 treat the \iast{kuṭṭaka} method for solving
multiple planetary congruences and refined calendar computations.]}
\vspace{0.5cm}

%% Verse 100
\begin{paracol}{2}
\switchcolumn[0]
\dev{%
हृदि-घात्रम् अमी प्रश्नाः प्रश्नान् अन्यान् सहस्र-शः कुर्यात् ।\\
अन्यैः दत्तान् प्रश्नान् उक्त्या एवम् साधयेत् करणैः ॥ १०० ॥
}

\switchcolumn[1]
\textbf{V.100.}\quad These problems strike the heart (with joy).
One may pose thousands of other problems; given by others — thus
one should solve those problems by the stated methods
(\iast{karaṇaiḥ}).
\end{paracol}

\vspace{0.3cm}

%% Verse 101
\begin{paracol}{2}
\switchcolumn[0]
\dev{%
जन-संसदि दैव-विदाम् तेजः नाशयति भानुः इव भानाम् ।\\
कुट्टाकार-प्रश्नैः पठितैः अपि किम् पुनः शत-शः ॥ १०१ ॥
}

\switchcolumn[1]
\textbf{V.101.}\quad In the assembly of people, among the knowers
of fate (\iast{daiva-vidām}), it destroys the radiance (of rivals),
like the sun (destroys) the stars — (even) one \iast{kuṭṭaka}
problem recited — what to speak of hundreds!
\end{paracol}

\vspace{0.3cm}

%% Verse 102
\begin{paracol}{2}
\switchcolumn[0]
\dev{%
प्रति-सूत्रम् अमी प्रश्नाः पठिताः स-उद्देशकेषु सूत्रेषु ।\\
आर्या-त्रि-अधिक-शतेन च कुट्टः च अष्टादशः अध्यायः ॥ १०२ ॥
}

\switchcolumn[1]
\textbf{V.102.}\quad These problems have been stated for each rule,
in the illustrative rules (\iast{sa-uddeśakeṣu sūtreṣu}). By
three hundred and three (303) \iast{āryā} (verses), the
\iast{kuṭṭaka} and the eighteenth chapter.
\end{paracol}

\vspace{0.5cm}

\begin{center}
\dev{इति श्री-ब्राह्मस्फुटसिद्धान्ते कुट्टकाध्यायो अष्टादशः}\\[0.3cm]
\textit{[Thus ends the Eighteenth Chapter, the Kuṭṭaka, in the
\textbf{Brāhmasphuṭasiddhānta}.]}
\end{center}

\newpage

%% ========================================================================
%% CHAPTER XIX: SHANKU-CHAYA-VIJNANAM
%% Gnomon and Shadow
%% ========================================================================
\dsec{अथ शङ्कु-छाया-विज्ञानम् — एकोनविंशः}
\addcontentsline{toc}{section}{Chapter XIX: Śaṅku-chāyā-vijñānam}

\begin{paracol}{2}
\switchcolumn[0]
\dev{%
\textbf{शङ्कु-छाया-विज्ञानम्}\\[0.2cm]
अस्मिन् अध्याये शङ्कोः छायायाः च मापन-विधिः उच्यते।
अनेन अक्षांशः, गृह-औच्च्यम्, जल-मापनम्, दर्पण-प्रतिबिम्बेन
ऊर्ध्वता-निर्धारणम् च सिध्यति।
}

\switchcolumn[1]
\textbf{Chapter XIX: On Gnomon and Shadow Knowledge}\\[0.2cm]
This chapter contains 20 verses on gnomon (\iast{śaṅku}) and
shadow (\iast{chāyā}) computations. These techniques allow the
determination of latitude, building heights, water depth, and
altitude by mirror reflection.
\end{paracol}

\vspace{0.3cm}
\cnote{[Printed pages 284–293.]}
\vspace{0.3cm}

%% Verse 1
\begin{paracol}{2}
\switchcolumn[0]
\dev{%
दृष्ट्वा दिन-अर्ध-घटिका यः अर्क-ज्ञः अक्ष-अंशकान् विजानाति ।\\
उदय-अन्तर-घटिकाभिः ज्ञातात् ज्ञेयम् सः तन्त्र-ज्ञः ॥ १ ॥
}

\switchcolumn[1]
\textbf{V.1.}\quad Having seen the half-day ghaṭikās, he who knows
the sun-knowing (method), knows the degrees of latitude
(\iast{akṣa-aṃśakān}). By the ghaṭikās of the interval from
rising — from the known, the knowable — he is a knower of the
\iast{tantra}.
\end{paracol}

\vspace{0.3cm}

%% Verse 2
\begin{paracol}{2}
\switchcolumn[0]
\dev{%
अस्त-अन्तर-घटिकाभिः यः ज्ञातात् ज्ञेयम् आनयति तस्मात् ।\\
मध्य-गतिम् युग-भ-गणान् आनयति ततः सः तन्त्र-ज्ञः ॥ २ ॥
}

\switchcolumn[1]
\textbf{V.2.}\quad By the ghaṭikās of the interval from setting, he
who brings the knowable from the known, thence brings the mean
motion and the \iast{yuga}-revolutions — he is a knower of the
\iast{tantra}.
\end{paracol}

\vspace{0.3cm}

%% Verse 3
\begin{paracol}{2}
\switchcolumn[0]
\dev{%
आनयति यः तमस्-रवि-शश-अङ्क-मानानि दीपक-शिख-औच्च्यात् ।\\
शङ्कु-तल-अन्तर-भूमि-ज्ञाने छायाम् सः तन्त्र-ज्ञः ॥ ३ ॥
}

\switchcolumn[1]
\textbf{V.3.}\quad He who brings the measures of the darkness-sun,
the moon-horn, from the height of the lamp-flame — in the
knowledge of the ground between gnomon-base (and object), the
shadow — he is a knower of the \iast{tantra}.
\end{paracol}

\vspace{0.3cm}

%% Verse 4
\begin{paracol}{2}
\switchcolumn[0]
\dev{%
इष्ट-गृह-औच्च्य-ज्ञः यः तत्-अन्तर-ज्ञः निरीक्ष्यते तु जले ।\\
गृह-भित्ति-अग्रम् दर्शयति दर्पणे वा सः तन्त्र-ज्ञः ॥ ४ ॥
}

\switchcolumn[1]
\textbf{V.4.}\quad He who knows the desired house-height, knows
that interval, is observed in water — shows the top of the house-wall,
or in a mirror — he is a knower of the \iast{tantra}.
\end{paracol}

\vspace{0.3cm}

%% Verse 5
\begin{paracol}{2}
\switchcolumn[0]
\dev{%
छाया-द्वितीया-भा-अग्र-अन्तर-विज्ञानेन वेत्ति दीप-औच्यम् ।\\
शङ्कु-छाया-ज्ञः वा भूमेः छायाम् सः तन्त्र-ज्ञः ॥ ५ ॥
}

\switchcolumn[1]
\textbf{V.5.}\quad By the knowledge of the second-shadow-light-edge
interval, he knows the lamp-height — or a knower of gnomon and
shadow knows the earth's shadow — he is a knower of the
\iast{tantra}.
\end{paracol}

\vspace{0.3cm}

%% Verse 6
\begin{paracol}{2}
\switchcolumn[0]
\dev{%
दृष्ट्वा गृह-तल-अन्तर-जालभोः दृष्ट्वा अग्रम् गृहस्य भूमि-ज्ञः ।\\
वेत्ति गृह-औच्च्यम् दृष्ट्वा तैल-स्थम् वा सः तन्त्र-ज्ञः ॥ ६ ॥
}

\switchcolumn[1]
\textbf{V.6.}\quad Having seen the house-base-interval of the net
of light and shadow, having seen the top of the house — a
knower of the ground — he knows the house-height, having seen
it in oil or water — he is a knower of the \iast{tantra}.
\end{paracol}

\vspace{0.3cm}

%% Verse 7
\begin{paracol}{2}
\switchcolumn[0]
\dev{%
वीक्ष्य गृह-अग्रम् सलिले प्रसार्य सलिलम् पुनः स्व-भू-ज्ञाने ।\\
आनयति जलात् भूमिम् गृहस्य वा औच्च्यम् सः तन्त्र-ज्ञः ॥ ७ ॥
}

\switchcolumn[1]
\textbf{V.7.}\quad Having observed the top of the house in water,
having spread out the water again, in self-ground-knowledge —
he brings from the water the ground or the height of the house
— he is a knower of the \iast{tantra}.
\end{paracol}

\vspace{0.3cm}

%% Verse 8
\begin{paracol}{2}
\switchcolumn[0]
\dev{%
ज्ञातैः छाया-पुरुषैः विज्ञाते तोय-कुड्ययोः विवरे ।\\
कुड्ये अर्क-तेजसः यः वेत्ति आरूढिम् सः तन्त्र-ज्ञः ॥ ८ ॥
}

\switchcolumn[1]
\textbf{V.8.}\quad With known shadow-men (measures), the interval
of water and wall having been known — in the wall, he who knows
the ascent of the sun's light (ray) — he is a knower of the
\iast{tantra}.
\end{paracol}

\vspace{0.3cm}

%% Verse 9
\begin{paracol}{2}
\switchcolumn[0]
\dev{%
इष्ट-दिवस-अर्ध-घटिका घटिका-पञ्चदश-अन्तर-प्राणाः ।\\
तद्-दिवस-चर-प्राणैः तैः अक्षम् साधयेत् प्राक्-वत् ॥ ९ ॥
}

\switchcolumn[1]
\textbf{V.9.}\quad The half-day ghaṭikās of the desired day, the
prāṇas of the interval of fifteen ghaṭikās — by the day's motion
prāṇas, by those he should compute the latitude as before.
\end{paracol}

\vspace{0.3cm}

%% Verse 10
\begin{paracol}{2}
\switchcolumn[0]
\dev{%
ज्ञात-ज्ञेय-ग्रहयोः उदय-अन्तर-नाडिकाभिः अधिक-ऊनः ।\\
उदयैः ज्ञातः ज्ञातात् ज्ञेयः प्राक्-अपरयोः ज्ञेयः ॥ १० ॥
}

\switchcolumn[1]
\textbf{V.10.}\quad Of the known and knowable planets, by the
nāḍikās of the interval from rising — greater or less — by the
risings, the known from the known, the knowable on the east-west
(is) to be known.
\end{paracol}

\vspace{0.3cm}

%% Verse 11
\begin{paracol}{2}
\switchcolumn[0]
\dev{%
ज्ञातः स-भ-अर्धः उदयैः अस्त-अन्तर-नाडिकाभिः अधिक-ऊनः ।\\
ज्ञातात् पूर्व-अपरयोः ज्ञेयः भ-अर्ध-ऊनके ज्ञेयः ॥ ११ ॥
}

\switchcolumn[1]
\textbf{V.11.}\quad The known, with its half-duration, by the
risings — greater or less by the nāḍikās of the interval from
setting — from the known, on the east-west the knowable; when the
half-duration is decreased, the knowable.
\end{paracol}

\vspace{0.3cm}

%% Verse 12
\begin{paracol}{2}
\switchcolumn[0]
\dev{%
ज्ञातम् कृत्वा मध्यम् भूयः अन्य-दिने तत्-अन्तरम् भुक्तिः ।\\
त्रैराशिकेन भुक्त्या कल्प-ग्रह-मण्डल-आनयनम् ॥ १२ ॥
}

\switchcolumn[1]
\textbf{V.12.}\quad Having made the known into the mean, again on
another day, that interval is the motion (\iast{bhuktiḥ}). By
the rule of three (\iast{trairāśikena}), by the motion — the
derivation of the planet's circle in the \iast{kalpa}.
\end{paracol}

\vspace{0.3cm}

%% Verse 13
\begin{paracol}{2}
\switchcolumn[0]
\dev{%
स्थिति-अर्धात् विपरीतम् तमस्-प्रमाणम् स्फुटम् ग्रहणे ।\\
मान-उदयात् रवि-इन्द्वोः घटिका-अवयवेन भ-उदय-तः ॥ १३ ॥
}

\switchcolumn[1]
\textbf{V.13.}\quad The reverse of half the duration (\iast{sthiti})
is the true darkness-measure in an eclipse. From the measure-rise
of the sun and moon, by the ghaṭikā-part, from the terrestrial
rising.
\end{paracol}

\vspace{0.3cm}

%% Verse 14
\begin{paracol}{2}
\switchcolumn[0]
\dev{%
दीप-तल-शङ्कु-तलयोः अन्तरम् इष्ट-प्रमाण-शङ्कु-गुणम् ।\\
दीप-शिखा-औच्च्यात् शङ्कुम् विशोध्या शेष-उद्धृतम् छाया ॥ १४ ॥
}

\switchcolumn[1]
\textbf{V.14.}\quad The interval of lamp-base and gnomon-base,
multiplied by the desired-quantity-gnomon, divided by the
lamp-flame-height minus the gnomon — the quotient is the shadow
(\iast{chāyā}).

\medskip
\textit{[Similar triangles: If $h$ = lamp height, $g$ = gnomon height,
$d$ = distance between them, then shadow $s = \frac{d \cdot g}{h - g}$.]}
\end{paracol}

\vspace{0.3cm}

%% Verse 15
\begin{paracol}{2}
\switchcolumn[0]
\dev{%
शङ्कु-अन्तरेण गुणिता छाया छाया-अन्तरेण भक्ता भूः ।\\
स-छाया शङ्कु-गुणा दीप-औच्च्यम् छायया भक्ता ॥ १५ ॥
}

\switchcolumn[1]
\textbf{V.15.}\quad The shadow, multiplied by the gnomon-interval,
divided by the shadow-interval — (is) the ground (distance). The
shadow together with the gnomon, multiplied by the gnomon, divided
by the shadow — (is) the lamp-height.
\end{paracol}

\vspace{0.3cm}

%% Verse 16
\begin{paracol}{2}
\switchcolumn[0]
\dev{%
ज्ञात्वा शङ्कु-छायाम् अनुपातात् साधयेत् समुच्छ्रयान् ।\\
गृह-चैत्य-तरु-नगानाम् औच्च्यम् विज्ञाय वा छायाम् ॥ १६ ॥
}

\switchcolumn[1]
\textbf{V.16.}\quad Having known the gnomon and shadow, from
proportion (\iast{anupātāt}) one should compute the elevations
(\iast{samucchrayān}) — of houses, temples, trees, and mountains
— or having known the height, the shadow.

\medskip
\textit{[This is the standard ``rule of three'' for similar triangles:
$\frac{\text{height}}{\text{shadow}} = \frac{\text{known height}}{\text{known shadow}}$,
so unknown quantities can be found by proportion.]}
\end{paracol}

\vspace{0.5cm}
\cnote{[Verses 17–19 give additional rules on shadow computations and
heights/depths using water and mirror reflections. Verse 20 gives
the chapter colophon.]}
\vspace{0.5cm}

\begin{center}
\dev{इति शङ्कु-छाया-विज्ञानम् — एकोनविंशोऽध्यायः}\\[0.3cm]
\textit{[Thus ends the Nineteenth Chapter on Gnomon and Shadow Knowledge.]}
\end{center}

\longline
\newpage

%% ========================================================================
%% CHAPTER XX: CHANDAS-CITTYUTTARA
%% Prosody / Combinatorics
%% ========================================================================
\dsec{अथ छन्दश्-चित्त्युत्तराध्यायः — विंशतितमः}
\addcontentsline{toc}{section}{Chapter XX: Chandaś-citty-uttarādhikāraḥ}

\begin{paracol}{2}
\switchcolumn[0]
\dev{%
\textbf{छन्दश्-चित्त्युत्तराध्यायः}\\[0.2cm]
अस्मिन् अध्याये छन्दः-शास्त्रस्य गणितम् उच्यते—गुरु-लघु-
वर्णानाम् व्यवस्थापनानि, प्रस्तारः, नष्टम्, उद्धिष्टम् च।
अस्य विषयः संयुक्त-गणितस्य पूर्व-रूपम् अस्ति।
}

\switchcolumn[1]
\textbf{Chapter XX: On Prosody (Chandaś-citty-uttara)}\\[0.2cm]
This chapter contains 19 verses on Sanskrit prosody
(\iast{chandah-śāstra}), dealing with combinatorial calculations
for metrical patterns — the number of possible arrangements of
long (\iast{guru}, `G') and short (\iast{laghu}, `L') syllables.
Brahmagupta here presents an early form of combinatorics.
\end{paracol}

\vspace{0.3cm}
\cnote{[Printed pages 294–303.]}
\vspace{0.3cm}

%% Verse 1
\begin{paracol}{2}
\switchcolumn[0]
\dev{%
ऋग्वर्गः पर्यायः समूहयोगावयुक्षु युग्मेषु ।\\
सोयाः प्राग्वत्प्राप्तादाश्चतुष्ककाः शेषयुक्त्योन्त्यः ॥ १ ॥
}

\switchcolumn[1]
\textbf{V.1.}\quad The \iast{ṛg}-group (combinations), the
permutations (\iast{paryāya}), the aggregate-sums, in the
unequal pairs — the \iast{soma} etc., obtained as before, the
quadruples, the remainder-additions — the last.
\end{paracol}

\vspace{0.3cm}

%% Verse 2
\begin{paracol}{2}
\switchcolumn[0]
\dev{%
एकादियुतविहीनावाद्यन्तौ तद्विपर्ययौ यावत् ।\\
वर्गादिषु विषमयुजां क्रमोत्क्रमाद्वर्धयेत् पादान् ॥ २ ॥
}

\switchcolumn[1]
\textbf{V.2.}\quad Decreased by one etc., the first and last,
their reversals — in squares etc., of the odd-even (syllables),
one should increase the quarters by direct and reverse order.
\end{paracol}

\vspace{0.3cm}

%% Verse 3
\begin{paracol}{2}
\switchcolumn[0]
\dev{%
एकैकेन द्व्याद्व्याः सोप्पप्पधिकेषु तत्-प्रतिष्ठेषु ।\\
वर्गादिरभीष्टान्तः प्रस्तारो भवति यवमध्यः ॥ ३ ॥
}

\switchcolumn[1]
\textbf{V.3.}\quad One by one, two by two, with additional
appendages in those positions — beginning from the square up to
the desired — the spread-out (\iast{prastāra}) becomes with
barley-middle (\iast{yava-madhya}, i.e. Pascal's triangle).

\medskip
\textit{[The ``Meru of Prosody'' (\iast{prastāra}) is Pascal's triangle:
the number of metrical patterns with $n$ syllables and $k$ guru
syllables is $\binom{n}{k}$. The ``barley-middle'' refers to the
shape of the triangular arrangement.]}
\end{paracol}

\vspace{0.3cm}

%% Verse 4
\begin{paracol}{2}
\switchcolumn[0]
\dev{%
सूनोन्त्यो द्विपदाग्रं त्रिपदाद्यानामधः पृथक् संख्या ।\\
तच्छोध्यो व्येकः पृथग्न्ताद्रूपमूर्ध्वयुतम् ॥ ४ ॥
}

\switchcolumn[1]
\textbf{V.4.}\quad With one, the last — the two-foot beginning,
the three-foot etc., below, separately the numbers. That is to be
subtracted, decreased by one, separately from that — combined
with one above.
\end{paracol}

\vspace{0.3cm}

%% Verse 5
\begin{paracol}{2}
\switchcolumn[0]
\dev{%
यावत् पादाव्येकागच्छाद्वर्णेष्वथैकवृद्धेषु ।\\
रूपाद्युतघाते वर्गाद्यानां परा संख्या ॥ ५ ॥
}

\switchcolumn[1]
\textbf{V.5.}\quad Up to the quarter, decreasing by one, among
the syllables, then increasing by one — in the product beginning
from one, the final number of squares etc.

\medskip
\textit{[For $n$ syllables, each of which can be G or L, the total
number of possible patterns is $2^n$.]}
\end{paracol}

\vspace{0.3cm}

%% Verse 6
\begin{paracol}{2}
\switchcolumn[0]
\dev{%
रूपाधिकपादार्धेविषमेषूर्ध्वः समेषु पादार्धे ।\\
अर्धाद्द्विगुणाव्येकांयुलान्यधस्तस्य सर्वेषाम् ॥ ६ ॥
}

\switchcolumn[1]
\textbf{V.6.}\quad With one added to the quarter-half — in the
uneven (syllables) above, in the even (syllables) at the
quarter-half. From the half, doubled, decreased by one — the
\iast{yulāni} (pairings) below that, for all.
\end{paracol}

\vspace{0.3cm}

%% Verse 7
\begin{paracol}{2}
\switchcolumn[0]
\dev{%
माध्यैस् तथार्धहीनैः क्रमपादैर् व्यस्ततुल्यपादाद्यः ।\\
विषमेरव्येकं मध्ये प्रोह्याद्यान्यतः कुर्यात् ॥ ७ ॥
}

\switchcolumn[1]
\textbf{V.7.}\quad By the middle (terms), thus decreased by half,
by the quarter-steps in order — equal to the reversed, beginning
from the quarter — in the uneven, subtracting one in the middle,
from that one should do (the computation).
\end{paracol}

\vspace{0.3cm}

%% Verse 8
\begin{paracol}{2}
\switchcolumn[0]
\dev{%
सैकक्रमतुल्याद्यैर् न्यासोऽभ्यधिको विशोधितश् चाधः ।\\
संख्यैक्यं तादृक् यादृक् प्रथमस् त्रिरहितो नष्टे ॥ ८ ॥
}

\switchcolumn[1]
\textbf{V.8.}\quad With one-plus-equal-to-the-step-beginnings, the
setting is excessively subtracted below — the sum of numbers
such as it is — the first, decreased by three — in the lost
(\iast{naṣṭa}, i.e. finding a missing pattern).
\end{paracol}

\vspace{0.3cm}

%% Verse 9
\begin{paracol}{2}
\switchcolumn[0]
\dev{%
माध्यैः कृतैश् च दलितैः समसंख्यायां क्रमोत्क्रमात् क्षेप्पम् ।\\
विषमायां व्येकायां दलम् क्रमाद् उत्त्क्रमात् सैकम् ॥ ९ ॥
}

\switchcolumn[1]
\textbf{V.9.}\quad Made with middle terms and halved — in the
equal-number (case), in direct and reverse order — the
\iast{kṣepa} (additive). In the uneven, decreased by one — the
half, in order, in reverse order, with one added.
\end{paracol}

\vspace{0.3cm}

%% Verse 10
\begin{paracol}{2}
\switchcolumn[0]
\dev{%
समसंख्यायां सोपानक्रमोत्क्रमाभ्यां तथैव विषमाभ्याम् ।\\
कल्प्या पचिते दृष्टे प्रथमः शेषाक्षराण्यन्ते ॥ १० ॥
}

\switchcolumn[1]
\textbf{V.10.}\quad In the equal-number (case), by the staircase
in direct and reverse order — thus also in the uneven — to be
constructed. When cooked (computed) is seen, the first — the
remaining syllables at the end.
\end{paracol}

\vspace{0.5cm}
\cnote{[Verses 11–18 continue the treatment of combinatorial prosody,
including algorithms for the ``Meru of Prosody'' (Pascal's triangle)
and rules for lost (\iast{naṣṭa}) and spread-out (\iast{uddhiṣṭa})
metrical patterns. The ``staircase'' (\iast{sopāna}) is a procedure
for listing all possible metrical patterns in systematic order.]}

\vspace{0.5cm}

\begin{center}
\dev{इति श्री-ब्राह्मस्फुटसिद्धान्ते छन्दश्-चित्त्युत्तराध्यायो विंशतितमः}\\[0.3cm]
\textit{[Thus ends the Twentieth Chapter on Prosody in the
\textbf{Brāhmasphuṭasiddhānta}.]}
\end{center}

\longline
\newpage

%% ========================================================================
%% CHAPTER XXI: GOLADHYAYAH — JYA-PRAKARANAM
%% Sine Computation
%% ========================================================================
\dsec{अथ गोलाध्यायः — एकविंशतितमः}
\addcontentsline{toc}{section}{Chapter XXI: Golādhyāyaḥ (Jyā-prakaraṇam)}

\begin{paracol}{2}
\switchcolumn[0]
\dev{%
\textbf{गोलाध्यायः — ज्याप्रकरणम्}\\[0.2cm]
अस्मिन् अंशे ज्या-खण्डानाम् उत्पत्तिः उच्यते। ब्रह्मगुप्तस्य
कल्पना द्वितीय-कोटि-परिकल्पनया (second-order interpolation)
ज्या-सारणीम् अधिक-सूक्ष्मां करोति।
}

\switchcolumn[1]
\textbf{Chapter XXI: The Sphere — Sine Computation}\\[0.2cm]
This chunk includes the Jyā-prakaraṇa (verses 17–23), which gives
Brahmagupta's method for computing the table of sines. The radius
is $R = 3270$ minutes. The ``segments'' (\iast{jyā-khaṇḍa}) are
the tabular sine-differences. Brahmagupta uses a second-order
interpolation formula, more accurate than the linear interpolation
used by earlier authors.
\end{paracol}

\vspace{0.3cm}
\cnote{[Printed pages 304–315 (beginning). Chapter 21 contains 70 verses.
This chunk includes the Jyā-prakaraṇa section (verses 17–23).]}
\vspace{0.5cm}

%% Verse 17
\begin{paracol}{2}
\switchcolumn[0]
\dev{%
राशि-अष्ट-अंशेषु अङ्कान् पद-सन्धिभ्यस् क्रम-उत्क्रमात् कृत्वा ।\\
बध्नीयात् सूत्राणि द्वयोः द्वयोः ज्यास् तद्-अर्धानि ॥ १७ ॥
}

\switchcolumn[1]
\textbf{V.17.}\quad Having placed marks at the eighth parts of a
sign ($3°45'$ intervals) in direct and reverse order, one should
construct chords (\iast{jyā}) for every two-by-two: their halves
are the sine-chords (\iast{jyā-ardhāni}).

\medskip
\textit{[Brahmagupta divides each sign (30°) into 8 parts of 3°45'
each, giving 24 intervals in a quadrant (0° to 90°). The
\iast{jyā-ardha} (half-chord) is the modern sine: $R\sin\theta$.]}
\end{paracol}

\vspace{0.3cm}

%% Verse 18
\begin{paracol}{2}
\switchcolumn[0]
\dev{%
ज्या-अर्धानि ज्या-अर्धानाम् ज्या-खण्डानि अन्तराणि ।\\
व्यस्तानि अन्त्यात् अथ वा इषुः उत्क्रम-ज्या धनुस् ताभ्याम् ॥ १८ ॥
}

\switchcolumn[1]
\textbf{V.18.}\quad The sine-half-differences of sine-halves are
the sine-segments (\iast{jyā-khaṇḍa}). Reversed from the last,
or: the bow (\iast{dhanus}) is the reverse-sine
(\iast{utkrama-jyā}) together with these.

\medskip
\textit{[The \iast{jyā-khaṇḍa} are the tabular differences:
$\Delta_i = R\sin(i+1)\alpha - R\sin i\alpha$ where $\alpha = 3°45'$.]}
\end{paracol}

\vspace{0.3cm}

%% Verse 19
\begin{paracol}{2}
\switchcolumn[0]
\dev{%
एक-द्वि-त्रि-गुणायाः व्यास-अर्ध-कृतेः पृथक् चतुर्थेभ्यः ।\\
मूलानि अष्ट-द्वादश-षोडश-खण्डानि अतः अन्यानि ॥ १९ ॥
}

\switchcolumn[1]
\textbf{V.19.}\quad From the quarter-parts of the half-diameter
squared, multiplied by one, two, and three: the roots are the
eighth, twelfth, and sixteenth segments; the others (are derived)
from these.

\medskip
\textit{[Initial segments computed from $R^2/4$:]}
\begin{align*}
j_8 &= \sqrt{R^2 \cdot 1 / 4} = R/2 = R\sin(30°)\\
j_{12} &= \sqrt{R^2 \cdot 2 / 4} = R/\sqrt{2} = R\sin(45°)\\
j_{16} &= \sqrt{R^2 \cdot 3 / 4} = R\sqrt{3}/2 = R\sin(60°)
\end{align*}
\textit{With $R = 3270$: $j_8 = 1635$, $j_{12} = 2312$, $j_{16} = 2832$.}
\end{paracol}

\vspace{0.3cm}

%% Verse 20
\begin{paracol}{2}
\switchcolumn[0]
\dev{%
तुल्य-क्रम-उत्क्रम-ज्या-सम-खण्डक-वर्ग-युति-चतुर्-भागम् ।\\
प्रोह्य अनष्टम् व्यास-अर्ध-वर्गतः तत्-पदे प्रथमम् ॥ २० ॥
}

\switchcolumn[1]
\textbf{V.20.}\quad Having subtracted the fourth part of the sum
of squares of equal direct and reverse sines and equal segments
from the square of the half-diameter, the square-root of the
remainder is the first (sine-chord).

\medskip
\textit{[If $j_k$ and $j_{n-k}$ are known sines, and $\delta$ is the
common segment-difference, the next sine is:]}
\begin{align*}
j_{k+1} &= \sqrt{R^2 - \frac{(j_k + j_{n-k})^2 + (2\delta)^2}{4}}
\end{align*}
\end{paracol}

\vspace{0.3cm}

%% Verse 21
\begin{paracol}{2}
\switchcolumn[0]
\dev{%
तद्-दल-खण्डानि तद्-ऊन-जिन-समानि द्वितीयम् उत्पत्तौ ।\\
[संख्या-वाचकैः अक्षरैः संकेतितानि] ॥ २१ ॥
}

\switchcolumn[1]
\textbf{V.21.}\quad Its half-segments, decreased by those equal to
the earth (\iast{jina}) — the second in production.

\medskip
\textit{[This verse contains numerical values encoded as words
(\iast{saṅketa} system). ``Jina'' = 1 (Lord of the earth). The
second-order correction refines the first-order segments.]}
\end{paracol}

\vspace{0.3cm}

%% Verse 22
\begin{paracol}{2}
\switchcolumn[0]
\dev{%
एवम् जीवा-खण्डानि अल्पानि बहूनि वा आद्य-खण्डानि ।\\
ज्या-अर्धानि वृत्त-परिधेः षष्ठ-चतुर्थ-त्रि-भागानाम् ॥ २२ ॥
}

\switchcolumn[1]
\textbf{V.22.}\quad Thus the small or numerous sine-segments —
the first segments are the sine-halves of one-sixth, one-fourth,
and one-third of the circle's circumference.

\medskip
\textit{[The reference sines at key angles:]}
\begin{align*}
R\sin(30°) &= R/2 = 1635 & R\sin(45°) &= R/\sqrt{2} \approx 2312\\
R\sin(60°) &= R\sqrt{3}/2 \approx 2832
\end{align*}
\end{paracol}

\vspace{0.3cm}

%% Verse 23
\begin{paracol}{2}
\switchcolumn[0]
\dev{%
उत्क्रम-सम-खण्ड-गुणात् व्यासात् अथ वा चतुर्थ-भागात् यत् ।\\
कृत्वा उक्त-खण्डकानि ज्या-अर्ध-आनयनम् न लघु अस्मात् ॥ २३ ॥
}

\switchcolumn[1]
\textbf{V.23.}\quad From the diameter multiplied by the equal
reverse-segment, or from the quarter-part, having done the stated
segment operations — the derivation of sine-halves is not
difficult from this.

\medskip
\textit{[Summary of the method: Given the tabular differences
(\iast{khaṇḍa}), one can compute all sines by second-order
interpolation using only multiplication and square roots.]}
\end{paracol}

\vspace{0.5cm}
\cnote[0.5cm]{[Brahmagupta's sine table uses $R = 3270$ minutes. The
twenty-four \iast{jyā-khaṇḍa} (sine-differences) are: 225, 224, 222,
219, 215, 210, 205, 199, 191, 183, 174, 164, 154, 143, 131, 119,
106, 93, 79, 65, 51, 37, 22, 7. The second-order correction
formula accounts for the curvature of the sine function, yielding
values accurate to within 1 minute of arc.]}

\longline
\newpage

%% ========================================================================
%% CHAPTER XXII: YANTRADHYAYAH
%% Astronomical Instruments
%% ========================================================================
\dsec{अथ यन्त्राध्यायः — द्वाविंशतितमः}
\addcontentsline{toc}{section}{Chapter XXII: Yantrādhyāyaḥ}

\begin{paracol}{2}
\switchcolumn[0]
\dev{%
\textbf{यन्त्राध्यायः}\\[0.2cm]
अस्मिन् अध्याये खगोल-दर्शनार्थं विविधानि यन्त्राणि
वर्ण्यन्ते—धनुः-यन्त्रम्, यष्टि-यन्त्रम्, जल-यन्त्रम्,
चक्र-यन्त्रम्, शङ्कु-यन्त्रम् च।
}

\switchcolumn[1]
\textbf{Chapter XXII: On Astronomical Instruments}\\[0.2cm]
This chapter contains 43–57 verses on astronomical instruments
(\iast{yantra}). It describes the bow instrument
(\iast{dhanur-yantra}), staff instrument (\iast{yaṣṭi-yantra}),
water instruments, wheel instruments, and gnomon instruments.
These were used for measuring celestial altitudes, azimuths, and
for timekeeping.
\end{paracol}

\vspace{0.3cm}
\cnote{[Printed pages $\sim$289–315.]}
\vspace{0.5cm}

\begin{paracol}{2}
\switchcolumn[0]
\dev{%
[प्रारम्भिक-श्लोकाः विविध-यन्त्र-वर्णनानि]\\[0.2cm]
अस्मिन् अध्याये उक्तानि यन्त्राणि:\\
\textbullet\ धनुर्म्यन्त्रम् — The bow instrument\\
\textbullet\ यष्टियन्त्रम् — The staff instrument\\
\textbullet\ जलयन्त्रम् — Water instruments\\
\textbullet\ शङ्कुयन्त्रम् — Gnomon instruments\\
\textbullet\ चक्रयन्त्रम् — Wheel/circle instruments
}

\switchcolumn[1]
\textbf{Instruments Described}\\[0.2cm]
\begin{itemize}[leftmargin=*]
  \item \textbf{Dhanur-yantra} — A semicircular or arc-shaped
    instrument for measuring altitudes, resembling a bow.
  \item \textbf{Yaṣṭi-yantra} — A staff or rod instrument for
    measuring angles and time.
  \item \textbf{Jala-yantra} — Water-based instruments including
    water clocks and mercury devices.
  \item \textbf{Śaṅku-yantra} — Gnomon instruments for shadow
    measurements.
  \item \textbf{Cakra-yantra} — Circular instruments with rotating
    sighting vanes for measuring celestial positions.
\end{itemize}
\end{paracol}

\vspace{0.5cm}
\cnote{[Verses 1–28 describe the bow instrument, staff instrument, and
their use in measuring celestial altitudes and azimuths. Verses
29–40 treat water instruments and circle instruments.]}
\vspace{0.5cm}

%% Verse 41
\begin{paracol}{2}
\switchcolumn[0]
\dev{%
क्षेप्तं स्वर्चिषा क्षेप्यं समं तथा पारदे यदा भवति ।\\
कालसम्मिश्रमानं शककृतान्मुद्दे वा ॥ ४१ ॥
}

\switchcolumn[1]
\textbf{V.41.}\quad The \iast{kṣepṭa} (object to be observed), by
its own light, the \iast{kṣepya} (object aimed at) — equal, thus
in mercury (\iast{pārada}), when it becomes — the mixed-time
measure from the \iast{śaka} epoch (78 CE) or from the
\iast{mudda} (particular epoch).
\end{paracol}

\vspace{0.3cm}

%% Verse 42
\begin{paracol}{2}
\switchcolumn[0]
\dev{%
कोतस्योपरिगामिनं त्रियक् सूत्रके धृतमल्पेषु ।\\
प्रावनलके प्रक्षिप्य नाडिका श्रवति पतत्ये ॥ ४२ ॥
}

\switchcolumn[1]
\textbf{V.42.}\quad The co-altitude (\iast{kota}) going above,
held obliquely in the string (\iast{sūtraka}), in small
(measures) — having cast upon the eastern fire-place (\iast{prāv-analaka}),
the nāḍikā (water-clock) flows, falls thus.
\end{paracol}

\vspace{0.5cm}
\cnote{[Verses 43–48 contain further rules on the use of instruments
for time measurement and planetary observation.]}
\vspace{0.5cm}

%% Verse 49
\begin{paracol}{2}
\switchcolumn[0]
\dev{%
लघुदारमयं चक्रं समधुरीरान्तरं पुष्करारणाम् ।\\
श्रद्धा पारदपूर्णं मध्ये सुलिप्तं कृतसंध्यः ॥ ४९ ॥
}

\switchcolumn[1]
\textbf{V.49.}\quad A wheel of light-wood, with equal bearing
intervals, of lotus-shape — having faith, filled with mercury
in the center, well-polished, having performed the saṃdhyā
(preparation).
\end{paracol}

\vspace{0.3cm}

%% Verse 50
\begin{paracol}{2}
\switchcolumn[0]
\dev{%
तियक् कोतिमध्येध्वाधारस्थो रूपारदं भ्रजति ।\\
छिद्रायुक् चक्रकमजस्रमेवं वाम्बरे भ्रमति ॥ ५० ॥
}

\switchcolumn[1]
\textbf{V.50.}\quad Obliquely in the middle of the co-altitude,
the support above — a beautiful semi-circle (\iast{rūpā-radaṃ})
moves. The small wheel with holes — thus constantly in the
air it rotates.

\medskip
\textit{[Description of a self-rotating wheel instrument, possibly
an early form of anemometer or a water-driven astronomical
device. The ``holes'' regulate the flow of water or air.]}
\end{paracol}

\vspace{0.5cm}
\cnote{[Verses 51–53: Additional rules on wheel instruments and their
construction. The chapter concludes with verses praising the value
of observation and proper instruments.]}
\vspace{0.5cm}

%% Verse 53
\begin{paracol}{2}
\switchcolumn[0]
\dev{%
करणज्या शिथ्रचलनमेव शरमोच्छरोत्क्षेपवाच्चाच्च ।\\
श्रद्धायो द्विविधो यन्त्रेणारिच्छन्दपञ्चकात् ॥ ५३ ॥
}

\switchcolumn[1]
\textbf{V.53.}\quad The \iast{karaṇa-jyā} (computed sine), the
unsteady motion itself — the arrow, the rising, the ascension —
by the instrument with confidence, two-fold — from the five-fold
of Ari (the sun's) motion.
\end{paracol}

\vspace{0.5cm}

\begin{center}
\dev{इति ब्रह्मस्फुटसिद्धान्ते यन्त्राध्यायो द्वाविंशतितमः}\\[0.3cm]
\dev{दर्शनार्थः समाप्तः}\\[0.5cm]
\textit{[Thus ends the Twenty-second Chapter on Instruments in the
\textbf{Brāhmasphuṭasiddhānta}.\\
End of the section on observation.]}
\end{center}

\vspace{1cm}
\longline
\newpage

%% ========================================================================
%% APPENDIX: VARIANT READINGS
%% ========================================================================
\section*{Appendix: On Variant Readings}
\addcontentsline{toc}{section}{Appendix: On Variant Readings}

\begin{paracol}{2}
\switchcolumn[0]
\dev{%
\textbf{पाठ-भेद-सूची}\\[0.2cm]
अस्य ग्रन्थस्य पाठ-संहितायां चत्वारि मुख्यानि
पुस्तक-हस्तलिखितानि उपयुक्तानि सन्ति—
(क) पूना-पुस्तकम्, (ग) बम्बई-पुस्तकम्, (च) बडोदा-पुस्तकम्,
(छ) होशियारपुर-पुस्तकम्।
}

\switchcolumn[1]
\textbf{Manuscript Sources for Variant Readings}\\[0.2cm]
The footnotes in the printed edition record variant readings from
four manuscript sources, designated by the following sigla:
\end{paracol}

\vspace{0.3cm}

\begin{center}
\begin{tabular}{@{}cl@{}}
\toprule
\textbf{Siglum} & \textbf{Source} \\
\midrule
(क) & Poona manuscript (Bhandarkar Oriental Research Institute) \\
(ग) & Bombay manuscript (Royal Asiatic Society) \\
(च) & Baroda manuscript (Oriental Research Institute) \\
(छ) & Hoshiarpur manuscript (V.V.R.I.) \\
\bottomrule
\end{tabular}
\end{center}

\vspace{0.3cm}

\begin{paracol}{2}
\switchcolumn[0]
\dev{%
पाद-टिप्पण्यां ``(ग) इति पाठस्य स्थाने'' इत्यादि
संकेतः पाठ-भेदं दर्शयति। ``(वि—इसकी क्रमसंख्या ३१ है)''
इत्यादि टिप्पणी हस्तलिखित-पुस्तकस्य अन्तः अस्य श्लोकस्य
संख्या भिन्ना इति दर्शयति।
}

\switchcolumn[1]
The footnotes use a notation such as ``(ग) for (\textit{reading})''
to indicate that manuscript (ग) has the variant reading instead of
the adopted text. Parenthetical notations like ``(वि—इसकी क्रमसंख्या
३१ है)'' indicate that the verse number in the source manuscript is
different.
\end{paracol}

\vspace{0.5cm}

%% ========================================================================
%% MATHEMATICAL SUMMARY
%% ========================================================================
\section*{Mathematical Summary of This Chunk}
\addcontentsline{toc}{section}{Mathematical Summary}

\begin{paracol}{2}
\switchcolumn[0]
\dev{%
\textbf{गणित-सारांशः}\\[0.2cm]
अस्मिन् खण्डे प्रमुखाणि गणितीय-सिद्धान्तानि सन्ति:\\
\textbullet\ कुट्टक-विधिः (Diophantine equations)\\
\textbullet\ ऋण-धन-कर्म (signed arithmetic)\\
\textbullet\ शून्य-कर्म (zero in computation)\\
\textbullet\ वर्गप्रकृतिः (Pell's equation)\\
\textbullet\ द्विघात-समीकरणम् (quadratic formula)\\
\textbullet\ ज्या-परिकल्पना (sine computation)\\
\textbullet\ छन्दः-संख्यानम् (combinatorics)
}

\switchcolumn[1]
\textbf{Key Mathematical Contributions}\\[0.2cm]
\begin{itemize}[leftmargin=*]
  \item \textbf{Kuṭṭaka} — The Euclidean algorithm for solving
    linear Diophantine equations $ax + c = by + d$.
  \item \textbf{Signed arithmetic} — First explicit rules for
    $+, -, \times, \div$ with positive, negative, and zero.
  \item \textbf{Zero} — $\frac{0}{0} = 0$ (contextual); $a + 0 = a$;
    $a \times 0 = 0$.
  \item \textbf{Pell's equation} — The composition formula
    (\iast{bhāvanā}) for generating solutions to $Nx^2 + 1 = y^2$.
  \item \textbf{Quadratic formula} — $x = \frac{\sqrt{4ac+b^2}-b}{2a}$
    for $ax^2 + bx = c$.
  \item \textbf{Sine table} — Second-order interpolation for
    $R\sin\theta$ with $R = 3270$.
  \item \textbf{Combinatorics} — $2^n$ patterns; Pascal's triangle
    (\iast{meru-prastāra}).
\end{itemize}
\end{paracol}

\vspace{0.5cm}

\begin{paracol}{2}
\switchcolumn[0]
\dev{%
अयं खण्डः ब्रह्मस्फुटसिद्धान्तस्य गणितीय-हृदयम् अन्तर्भूतवान्।
कुट्टकाध्यायः विशेषेण प्रसिद्धः—अस्या युग-ग्रहणस्य
प्रथम-सम्पूर्ण-वर्णना अत्र लभ्यते।
}

\switchcolumn[1]
This chunk contains the mathematical heart of the
\textit{Brāhmasphuṭasiddhānta}. The Kuṭṭakādhyāya is especially
notable — it contains the first complete description of the
solution to Pell's equation in world mathematical literature.
\end{paracol}

\vspace{1cm}

%% ========================================================================
%% END MATTER
%% ========================================================================
\begin{center}
\rule{0.4\textwidth}{0.4pt}\\[0.5cm]
{\large\textit{End of Chunk 5 (Pages 561–700)}}\\[0.3cm]
{\small The next chunk (pages 701–838) will contain\\
the remaining chapters (23–24) and the concluding material of Volume I.}\\[0.5cm]
\rule{0.4\textwidth}{0.4pt}
\end{center}

\end{document}
%% ========================================================================
%% END OF FILE
%% ========================================================================



% ============================================================
% Source: translations/en/indian/brahmagupta-brahmasphutasiddhanta-pt1-chunk6_bilingual.tex
% ============================================================

%% ========================================================================
%% Brāhmasphuṭasiddhānta of Brahmagupta (628 CE)
%% Part 1, Chunk 6: Pages 701–838 — BILINGUAL EDITION
%% LEFT: Sanskrit in Devanagari  |  RIGHT: English translation with IAST
%% ========================================================================
\documentclass[11pt,a4paper]{article}

%% -------- Core Packages --------
\usepackage{fontspec}
\usepackage{polyglossia}
\usepackage{amsmath,amssymb,amsthm}
\usepackage{geometry}
\usepackage{graphicx}
\usepackage{xcolor}
\usepackage{titlesec}
\usepackage{fancyhdr}
\usepackage{enumitem}
\usepackage{array,booktabs,longtable,multirow}
\usepackage{hyperref}
\usepackage{lettrine}
\usepackage{microtype}
\usepackage{tocloft}
\usepackage{indentfirst}
\usepackage{paracol}
\usepackage{etoolbox}

%% -------- Page Geometry --------
\geometry{
  paperwidth=210mm,paperheight=297mm,
  left=18mm,right=18mm,top=26mm,bottom=26mm,
  headheight=14pt,footskip=30pt
}

%% -------- Language Setup --------
\setmainlanguage{english}
\setotherlanguage{sanskrit}

%% -------- Font Configuration --------
\setmainfont{Linux Libertine O}
\setsansfont{Linux Biolinum O}
\setmonofont{DejaVu Sans Mono}

%% Devanagari (Sanskrit/Hindi) Font — Noto fonts by name
\newfontfamily\devanagarifont[Script=Devanagari]{Noto Serif Devanagari}
\newfontfamily\devanagarifontsf[Script=Devanagari]{Noto Sans Devanagari}
\newfontfamily\devanagarifontbold[Script=Devanagari]{Noto Serif Devanagari Bold}

%% -------- Colors --------
\definecolor{titlecolor}{RGB}{45,55,72}
\definecolor{sectioncolor}{RGB}{55,65,81}
\definecolor{notecolor}{RGB}{120,53,15}
\definecolor{rulecolor}{RGB}{180,160,140}
\definecolor{versecolor}{RGB}{60,40,20}
\definecolor{commentarycolor}{RGB}{40,40,60}
\definecolor{sanskritbg}{RGB}{255,253,245}
\definecolor{englishbg}{RGB}{250,252,255}
\definecolor{dividercolor}{RGB}{160,140,120}

%% -------- Section Formatting --------
\titleformat{\section}
  {\normalfont\Large\bfseries\color{sectioncolor}}
  {\thesection}{1em}{}
\titleformat{\subsection}
  {\normalfont\large\bfseries\color{sectioncolor}}
  {\thesubsection}{1em}{}
\titleformat{\subsubsection}
  {\normalfont\normalsize\bfseries\color{sectioncolor}}
  {\thesubsubsection}{1em}{}

%% -------- Header/Footer --------
\pagestyle{fancy}
\fancyhf{}
\fancyhead[L]{\small\textit{Brāhmasphuṭasiddhānta — Bilingual Edition — Pages 701–838}}
\fancyhead[R]{\small\thepage}
\renewcommand{\headrulewidth}{0.4pt}
\renewcommand{\footrulewidth}{0pt}

%% -------- Custom Commands --------
\newcommand{\longline}{\noindent\rule{\textwidth}{0.4pt}}
\newcommand{\shortline}{\noindent\rule{0.5\textwidth}{0.4pt}\par}
\newcommand{\cnote}[1]{\textcolor{notecolor}{\textit{#1}}}
\newcommand{\dev}[1]{{\devanagarifont #1}}
\newcommand{\devbold}[1]{{\devanagarifontbold #1}}
\newcommand{\devsf}[1]{{\devanagarifontsf #1}}

%% Parallel column environments
\columnratio{0.48}
\setlength{\columnsep}{1.2em}

%% Custom verse environment (no verse.sty needed)
\newenvironment{simpleverse}{%
  \list{}{\itemsep 0pt \parsep 0pt \topsep 4pt \leftmargin 2em}%
  \item\relax\ignorespaces
}{\endlist}

%% Sanskrit verse in left column
\newcommand{\sanskritverse}[1]{%
  \begin{leftcolumn*}
  \begin{simpleverse}\devanagarifont\color{versecolor}#1\end{simpleverse}
  \end{leftcolumn*}
}

%% English translation in right column
\newcommand{\englishverse}[1]{%
  \begin{rightcolumn}
  \begin{simpleverse}#1\end{simpleverse}
  \end{rightcolumn}
}

%% Bilingual verse pair
\newenvironment{bilingualverse}{%
  \begin{paracol}{2}
  \fontsize{10.5}{13}\selectfont
}{%
  \end{paracol}
  \vspace{0.3cm}
}

%% Single-language spanning block
\newcommand{\spanningtext}[1]{%
  \switchcolumn[0]*#1\switchcolumn
}

%% Decorative divider
\newcommand{\divider}{\begin{center}\textcolor{dividercolor}{\rule{0.5\textwidth}{0.5pt}}\end{center}}
\newcommand{\wavydivider}{\begin{center}$\sim\sim\sim\sim\sim\sim\sim\sim\sim\sim\sim\sim\sim\sim\sim\sim\sim\sim\sim\sim$\end{center}}

%% Column labels
\newcommand{\sanskritlabel}{\begin{center}\textcolor{sectioncolor}{\small\textsc{Sanskrit (Devanagari)}}\end{center}}
\newcommand{\englishlabel}{\begin{center}\textcolor{sectioncolor}{\small\textsc{English (with IAST)}}\end{center}}

%% -------- Hyperref Setup --------
\hypersetup{
  colorlinks=true,
  linkcolor=sectioncolor,
  citecolor=sectioncolor,
  urlcolor=sectioncolor,
  pdfencoding=auto,
  pdftitle={Brāhmasphuṭasiddhānta — Bilingual Edition — Pages 701-838},
  pdfauthor={Brahmagupta (628 CE)},
  pdfsubject={Bilingual Sanskrit-English Edition of Brahmasphutasiddhanta Part 1 Final Portion}
}

%% ========================================================================
\begin{document}

%% ========================================================================
%% TITLE PAGE
%% ========================================================================
\thispagestyle{empty}
\begin{center}
\vspace*{1.5cm}

{\fontsize{12}{14}\selectfont\dev{श्रीब्रह्मगुप्ताचार्य-विरचित}}\\[0.3cm]
{\fontsize{26}{30}\selectfont\devbold{ब्राह्मस्फुटसिद्धान्तः}}\\[0.4cm]
{\Large\devbold{द्विभाषी संस्करणम्}}\\[0.6cm]
{\large\textit{Brāhmasphuṭasiddhānta — Bilingual Edition}}

\vspace{1.2cm}

\longline

\vspace{0.6cm}

{\LARGE\textbf{Part 1 — Final Portion}}\\[0.4cm]
{\large\textit{Pages 701--838 of the Original Edition}}\\[0.3cm]
{\small Comprising:}\\[0.2cm]
{\normalsize Chapter 23 (Mānādhyāya) · Subject Index ·}\\[0.1cm]
{\normalsize Madhyamādhikāra Commentary}

\vspace{0.6cm}

\longline

\vspace{1cm}

\begin{paracol}{2}
\begin{leftcolumn*}
\begin{center}
{\large\devbold{संस्कृत पाठः}}\\[0.3cm]
{\small\dev{सौरचन्द्रार्क्षसावन-\\मानप्रयोगविधिः}}
\end{center}
\end{leftcolumn*}
\begin{rightcolumn}
\begin{center}
{\large\textbf{Sanskrit Text}}\\[0.3cm]
{\small Solar, Lunar, Stellar, and\\Civil Measures and Application}
\end{center}
\end{rightcolumn}
\end{paracol}

\vspace{1cm}

{\large Author: \textbf{Brahmagupta} (628 CE)}\\[0.3cm]
{\normalsize Editor: Acharya Ramswarup Sharma}\\[0.1cm]
{\small Indian Institute of Astronomical and Sanskrit Research}\\[0.1cm]
{\small New Delhi, First Edition 1966}

\vspace{1.5cm}

{\small Bilingual Edition — Sanskrit (Devanagari) / English (with IAST)}\\[0.2cm]
{\footnotesize Prepared with XeLaTeX · Noto Serif Devanagari · Polyglossia}

\vfill
{\textit{Digital preservation of classical Indian mathematical astronomy}}

\end{center}
\newpage

%% ========================================================================
%% TABLE OF CONTENTS FOR THIS VOLUME
%% ========================================================================
\section*{Contents of this Portion / \dev{अस्यांशस्य विषयसूची}}
\addcontentsline{toc}{section}{Contents of this Portion}

\begin{paracol}{2}
\sanskritlabel
\englishlabel
\switchcolumn[0]*

\begin{enumerate}[label=\dev{\arabic*.},leftmargin=*]
  \item \textbf{\dev{त्रयोविंशो मानाध्यायः}} (पृ.\,७०१--७३९)
  \begin{itemize}[label=\textendash,leftmargin=1em]
    \item \dev{प्रयोगविधिः} — श्लोकाः १--४
    \item \dev{चतुर्विंशः} — श्लोकाः ५--२३
    \item \dev{पदविभाग-विश्लेषणम्}
  \end{itemize}
  \item \textbf{\dev{प्रकाशकस्य विवरणम्}} — पृ.\,७४०--७४९
  \item \textbf{\dev{विषयानुक्रमणिका}} — पृ.\,७५०--७५१
  \item \textbf{\dev{मध्यमाधिकारटीका}} — पृ.\,७५२--८३८
  \begin{itemize}[label=\textendash,leftmargin=1em]
    \item \dev{मध्यमाधिकारस्य विस्तृतसंस्कृत-हिन्दीटीका}
    \item \dev{ग्रहाणां मध्यमगतिविधिः}
  \end{itemize}
\end{enumerate}

\switchcolumn

\begin{enumerate}[label=\arabic*.,leftmargin=*]
  \item \textbf{Chapter 23: Mānādhyāya} (The Measurement Chapter) — pp.\,701--739
  \begin{itemize}[label=\textendash,leftmargin=1em]
    \item Prayogavidhi (Practical Application) — verses 1--4
    \item Caturviṃśa (Twenty-four fold system) — verses 5--23
    \item Detailed word-analysis (Pada-vibhāga) for each verse
  \end{itemize}
  \item \textbf{Title Page and Publisher Information} — pp.\,740--749
  \item \textbf{Viṣayānukramaṇikā} (Subject Index) — pp.\,750--751
  \item \textbf{Madhyamādhikāra} (Mean Computation Chapter) — pp.\,752--838
  \begin{itemize}[label=\textendash,leftmargin=1em]
    \item Extensive Sanskrit–Hindi commentary on mean planetary motion
    \item Table of contents for chapters 1--6 of commentary
  \end{itemize}
\end{enumerate}
\end{paracol}

\divider
\newpage

%% ========================================================================
%% CHAPTER 23: MĀNĀDHYĀYA
%% ========================================================================
\section*{Chapter 23: \dev{त्रयोविंशो मानाध्यायः} — The Measurement Chapter}
\addcontentsline{toc}{section}{Chapter 23: Mānādhyāya (pp.\ 701--739)}
\markboth{Chapter 23: Mānādhyāya (pp.\ 701--739)}{}

\lettrine[lines=3,loversize=0.1]{T}{\textit{he twenty-third chapter}} of the \textit{Brāhmasphuṭasiddhānta} treats of the measurement of time and the practical application of astronomical calculations. This chapter, like others in the work, is composed in the classical \textit{āryā} metre and presents Brahmagupta's rules for calendrical computation, planetary positions, and the reconciliation of different astronomical systems.

The chapter is divided into two principal sections:
\begin{enumerate}[label=(\roman*),leftmargin=2em]
  \item \textbf{Prayogavidhi} — Practical application of the siddhānta
  \item \textbf{Caturviṃśa} — The twenty-four fold division of time
\end{enumerate}

\wavydivider

\subsection*{\dev{प्रयोगविधिः} — Practical Application (Verses 1--4)}

\begin{bilingualverse}
\begin{leftcolumn*}
\dev{सौरैराण्डं मासैस्तिथयश्चान्द्रैंस्तथा सावने दिवसाः ।}\\
\dev{दिनमासाब्दं पमध्यानत दिग्नाकेन्द्रं मानास्म्याम् ॥ १ ॥}
\end{leftcolumn*}
\begin{rightcolumn}
``The solar year, the lunar months and \textit{tithis}, the civil days, days, months, years, the positions of the planets, the \textit{nakṣatras}, and the zodiacal signs---these we call measures (\textit{māna}).''\\[0.3cm]
\hfill\textbf{— Verse 1}
\end{rightcolumn}
\end{bilingualverse}

\textbf{Word Analysis / \dev{पदविभागः}:}

\begin{paracol}{2}
\begin{leftcolumn*}
\dev{(१) सौरैराण्ड} — सौरवत्सरः। \dev{(२) मासैस्तिथयः} — चान्द्रमासाः तिथयः च। \dev{(३) सावनानि दिवसानि} — लौकिकदिनानि। \dev{(४) दिनमासाब्दानि} — अहोरात्रमाससंवत्सराः। \dev{(५) ग्रहनक्षत्रराशयः} — मानानि स्म्याम्॥
\end{leftcolumn*}
\begin{rightcolumn}
(1) \textit{saurairāṇḍa} — the solar year; (2) \textit{māsaistithayaḥ} — lunar months and \textit{tithis}; (3) \textit{sāvanāni divasāni} — civil days; (4) \textit{dinamāsābdāni} — days, months, and years; (5) \textit{grahanakṣatrarāśayaḥ} — planets, nakṣatras, and zodiacal signs---these are called \textit{māna} (measures).
\end{rightcolumn}
\end{paracol}

\vspace{0.4cm}

\begin{bilingualverse}
\begin{leftcolumn*}
\dev{मानान्ति सौरचन्द्रार्क्षां सावनानि ग्रहानयनमेभिः ।}\\
\dev{मानेन पृथक् चतुर्मितिथ्यर्ब्बहारो त्रं लोकस्य ॥ २ ॥}
\end{leftcolumn*}
\begin{rightcolumn}
``The measures of solar, lunar, stellar, and civil reckoning, and the computation of planets by these measures separately---the fourfold manner of computation is in use in the world.''\\[0.3cm]
\hfill\textbf{— Verse 2}
\end{rightcolumn}
\end{bilingualverse}

\begin{paracol}{2}
\begin{leftcolumn*}
\dev{सौरचान्द्रार्क्षसावनानि} — चत्वारि मानानि। \dev{एभिः पृथक् पृथक्} — प्रत्येकम्। \dev{ग्रहानयनम्} — ग्रहाणां गणनम्। \dev{लोकस्य} — संसारे प्रचलितम्॥
\end{leftcolumn*}
\begin{rightcolumn}
\textit{sauracāndrārṣasāvanāni} — four measures; \textit{ebhiḥ pṛthak pṛthak} — separately by each; \textit{grahānanayam} — computation of planets; \textit{lokasya} — current in the world.
\end{rightcolumn}
\end{paracol}

\vspace{0.4cm}

\begin{bilingualverse}
\begin{leftcolumn*}
\dev{युगवर्षं विपुलब्दयनतद्वह्निनिशोद्विद्धिहानयः ।}\\
\dev{सौरास्तिथि करणाऽधिकमासौन राशिपर्व्विक्रियाश्चान्द्रात् ॥ ३ ॥}
\end{leftcolumn*}
\begin{rightcolumn}
``A \textit{yuga-year}, abundant and deficient (days), fire (3), subtraction and addition---these (come) from the solar (system). \textit{Tithi}, \textit{karaṇa}, intercalary month, and the changes of zodiacal signs and \textit{parva} (come) from the lunar (system).''\\[0.3cm]
\hfill\textbf{— Verse 3}
\end{rightcolumn}
\end{bilingualverse}

\begin{paracol}{2}
\begin{leftcolumn*}
\dev{युगवर्षं विपुलाब्दयनानि} — अधिकदिनक्षयदिनानि। \dev{वह्निः} — त्रयः संख्या। \dev{उद्विद्धिहानयः} — क्षेपापचयौ। एते सौरात्॥ \dev{तिथिकरणाधिकमासाः} — चान्द्रात्। \dev{राशिपर्विक्रियाः} — राशिचक्रस्य परिवर्तनानि॥
\end{leftcolumn*}
\begin{rightcolumn}
\textit{yugavarṣaṃ vipulābdayanāni} — intercalary and deficient days; \textit{vahniḥ} — the number three; \textit{udviddhi-hānayaḥ} — addition and subtraction---from the solar system. \textit{tithikaraṇādhikamāsāḥ} --- from the lunar; \textit{rāśiparvikriyāḥ} --- changes of zodiacal signs and lunar half-months.
\end{rightcolumn}
\end{paracol}

\vspace{0.4cm}

\begin{bilingualverse}
\begin{leftcolumn*}
\dev{यसस्य वन प्रमाणं ग्रहगतृपवासं सुतं चिकिर्सतः ।}\\
\dev{सावनं मानांतु शैया प्रायश्चित्त क्रियाश्चान्याः ॥ ४ ॥}
\end{leftcolumn*}
\begin{rightcolumn}
``For one desirous of performing the \textit{pramāṇa} (measurement), the planetary motion, the rising and setting, and the creation (of almanacs), the civil measure (is useful), as well as the bed (auspicious moment), \textit{prāyaścitta} (expiatory rites), and other ceremonies.''\\[0.3cm]
\hfill\textbf{— Verse 4}
\end{rightcolumn}
\end{bilingualverse}

\begin{paracol}{2}
\begin{leftcolumn*}
\dev{प्रमाणं ग्रहगतिक्षयोदयसूतकल्पनम्} — पञ्चाङ्गादिनिर्माणार्थम्। \dev{सावनं मानम्} — लौकिकमानम्। \dev{शय्या} — शयनार्थकालः। \dev{प्रायश्चित्तक्रियाः} — कर्माणि॥
\end{leftcolumn*}
\begin{rightcolumn}
\textit{pramāṇaṃ grahagatikṣayodayasūtakalpanam} --- for the construction of the pañcāṅga (almanac); \textit{sāvanaṃ mānam} --- civil measure; \textit{śayyā} --- auspicious moment for lying down; \textit{prāyaścittakriyāḥ} --- expiatory rites and ceremonies.
\end{rightcolumn}
\end{paracol}

\wavydivider

\subsection*{\dev{चतुर्विंशः} — The Twenty-four Fold System (Verses 5--23)}

\begin{bilingualverse}
\begin{leftcolumn*}
\dev{यस्मात्तत्प्रति पश्तं संज्ञा संज्ञातो विना तस्मात् ।}\\
\dev{लोकप्रसिद्धसांख्यपादीनां दशांकाश्च ॥ ५ ॥}
\end{leftcolumn*}
\begin{rightcolumn}
``Since (computation) is impossible without names (\textit{saṃjñā}) given to each object, the names and numerals well-known in the world (are given).''\\[0.3cm]
\hfill\textbf{— Verse 5}
\end{rightcolumn}
\end{bilingualverse}

\textbf{\dev{टीका} / Commentary:} The commentator explains that Brahmagupta here gives a systematic nomenclature for the various measures of time and computation. The twenty-four items are enumerated in the following verses.

\vspace{0.3cm}

\begin{bilingualverse}
\begin{leftcolumn*}
\dev{युगपद्युगाद्यद् या याम्ययां भास्करस्य वा कन्याम् ।}\\
\dev{राश्यर्द्धसंख्यायां मस्तकया दिनदलाद्वेभ्राम् ॥ ६ ॥}
\end{leftcolumn*}
\begin{rightcolumn}
``A \textit{yugapad}, a \textit{yuga}, beginning from that, or of the sun, of the Yāmya (south), of Bhāskara, or of Kanyā (Virgo), half a sign, the numbers (associated), at the head, from half a day, from two, from Bhrama (revolution).''\\[0.3cm]
\hfill\textbf{— Verse 6}
\end{rightcolumn}
\end{bilingualverse}

\begin{bilingualverse}
\begin{leftcolumn*}
\dev{प्रपमेवकृतः सूर्ये तु पुलिशा रोमक वाराही यवनाद्यैः ।}\\
\dev{यस्मात्तस्मादेकः सिद्धान्तो ग्रथ्य रचनात्मकः ॥ ७ ॥}
\end{leftcolumn*}
\begin{rightcolumn}
``The sun has been measured by Pulisha, Romaka, Vāraha, and Yavana (Greeks). Therefore, from that (measurement) one must construct a proper \textit{siddhānta}.''\\[0.3cm]
\hfill\textbf{— Verse 7}
\end{rightcolumn}
\end{bilingualverse}

\begin{paracol}{2}
\begin{leftcolumn*}
\dev{पुलिशः रोमकः वाराही यवनाः} — पञ्चसिद्धान्तकाराः। \dev{सूर्यस्य प्रमाणम्} — तेजःप्रमाणम्। \dev{रचनात्मकः सिद्धान्तः} — ग्रन्थनिर्माणम्॥
\end{leftcolumn*}
\begin{rightcolumn}
\textit{puliśaḥ romakaḥ vārāhī yavanāḥ} --- the five \textit{siddhānta} authors; \textit{sūryasya pramāṇam} --- measurement of the sun; \textit{racanātmakaḥ siddhāntaḥ} --- construction of a scientific treatise.
\end{rightcolumn}
\end{paracol}

\vspace{0.3cm}

\begin{bilingualverse}
\begin{leftcolumn*}
\dev{यदि भिन्नाः सिद्धान्तभास्कर संस्कान्तयोभि भेदसमाः ।}\\
\dev{सरत्नः पूर्व्वस्यां विपुलत्याकार्द्दयो यस्य ॥ ८ ॥}
\end{leftcolumn*}
\begin{rightcolumn}
``If the \textit{siddhāntas} are different, (yet) Bhāskara (the sun) has the same modifications (\textit{saṃskāra})---the difference (is) abundant and deficient, those having the same (result) for the previous (method).''\\[0.3cm]
\hfill\textbf{— Verse 8}
\end{rightcolumn}
\end{bilingualverse}

\begin{bilingualverse}
\begin{leftcolumn*}
\dev{तत्र पक्षागरिततं मध्यं गत्युत्तरादयः पञ्च ।}\\
\dev{कुडूाकारे बेधः छिन्द्धि चतुर्थां तरे गोलः ॥ ९ ॥}
\end{leftcolumn*}
\begin{rightcolumn}
``In it (the \textit{siddhānta}) the five (methods): \textit{pakṣāgati}, \textit{madhya}, \textit{gati}, \textit{uttara}, etc.; \textit{kuṭūākāra} (transformation), \textit{bheda} (division), \textit{chid} (cutting), \textit{caturtha} (fourth proportional), and the sphere (\textit{gola}).''\\[0.3cm]
\hfill\textbf{— Verse 9}
\end{rightcolumn}
\end{bilingualverse}

\wavydivider

\begin{bilingualverse}
\begin{leftcolumn*}
\dev{भट्ट ब्रह्माचार्येण जिज्ञासुतनयेन गणितगोलविदा ।}\\
\dev{प्रायःश्लेषसंहलं वा स्फुटसिद्धान्तं कृतो ब्रह्मा ॥ १० ॥}
\end{leftcolumn*}
\begin{rightcolumn}
``By Brahmagupta, the son of Jiṣṇu, the knower of mathematics and the sphere, this true \textit{siddhānta} was composed, (which is) Brahma (itself), adorned with the beauty of elegant composition.''\\[0.3cm]
\hfill\textbf{— Verse 10}
\end{rightcolumn}
\end{bilingualverse}

\begin{paracol}{2}
\begin{leftcolumn*}
\dev{भट्ट ब्रह्माचार्यः} — श्रीब्रह्मगुप्तः। \dev{जिज्ञोः तनयः} — जित्स्नुतनयः। \dev{गणितगोलवित्} — गणितखगोलविद्वान्। \dev{स्फुटसिद्धान्तः} — सत्यज्योतिर्ग्रन्थः। \dev{ब्रह्मा} — ब्रह्मस्वरूपः॥
\end{leftcolumn*}
\begin{rightcolumn}
\textit{bhaṭṭa brahmācāryaḥ} --- the venerable Brahmagupta; \textit{jijñoḥ tanayaḥ} --- son of Jiṣṇu; \textit{gaṇitagolavit} --- knower of gaṇita (mathematics) and the celestial sphere; \textit{sphuṭasiddhāntaḥ} --- true astronomical treatise; \textit{brahmā} --- of the nature of Brahman itself.
\end{rightcolumn}
\end{paracol}

\vspace{0.4cm}

\begin{bilingualverse}
\begin{leftcolumn*}
\dev{मग्रहं युति वत्सांकुं विचित्रमलनाद्विव्रहोत्कसमः ।}\\
\dev{शकनिः कर्म्मबाहुत्वान्निर्भुतातो भास्करमग्रहं ॥ ११ ॥}
\end{leftcolumn*}
\begin{rightcolumn}
``(Various) instruments, conjunction, year-count, diverse multiplication, the \textit{utka} (elevation), the \textit{śakani} (depression), on account of the multitude of operations, (all these) from the diminished (part), the solar \textit{maragraham}...''\\[0.3cm]
\hfill\textbf{— Verse 11}
\end{rightcolumn}
\end{bilingualverse}

\begin{bilingualverse}
\begin{leftcolumn*}
\dev{प्राज्येयो नैकल्प्योर्द्धच्छिद्दिने स्थितस्य योग्केस्य ।}\\
\dev{शङ्कुश्ये कथयत्यब्दापि वत्सि सूर्य्यंश्च ॥ १२ ॥}
\end{leftcolumn*}
\begin{rightcolumn}
``The abundant, the one without method, the \textit{ṛddhicchid} (augmented cut), the \textit{yogketsya} of the standing, the \textit{śaṅkuśye} (gnomon-shadow) declares the year and also the solar year...''\\[0.3cm]
\hfill\textbf{— Verse 12}
\end{rightcolumn}
\end{bilingualverse}

\begin{bilingualverse}
\begin{leftcolumn*}
\dev{प्रज मया यन्नोकं गोलादग्रं क्षयार्धसतो त्रं तनु ।}\\
\dev{प्रायःश्लोद्दशोऽयं संज्ञाध्यायचतुर्विंशः ॥ १३ ॥}
\end{leftcolumn*}
\begin{rightcolumn}
``As much as has not been spoken by me, the remainder of the sphere, thin as a sixth part of a \textit{ksaya} (decrease), this (is) the \textit{prāyaśloka} (resumé verse). This (is) the \textit{saṃjñā-adhyāya} (chapter on nomenclature), the twenty-fourth.''\\[0.3cm]
\hfill\textbf{— Verse 13}
\end{rightcolumn}
\end{bilingualverse}

\begin{paracol}{2}
\begin{leftcolumn*}
\dev{मया यन्नोक्तं} — यन्मयोक्तं। \dev{गोलात् अग्रं} — गोलशेषम्। \dev{क्षयार्धसतो षष्ठांशतुल्यं} — अत्यल्पम्। \dev{प्रायःश्लोकः} — सङ्क्षेपश्लोकः। \dev{संज्ञाध्यायः चतुर्विंशः} — नामकल्पनानिर्देशकः अध्यायः॥
\end{leftcolumn*}
\begin{rightcolumn}
\textit{mayā yannoktaṃ} --- what has not been spoken by me; \textit{golāt agraṃ} --- the remainder of the sphere; \textit{kṣayārdhasato ṣaṣṭhāṃśatulyaṃ} --- as minute as one-sixth part of the decrease; \textit{prāyaḥślokaḥ} --- resumé verse; \textit{saṃjñādhyāyaś caturviṃśaḥ} --- the twenty-fourth chapter on nomenclature.
\end{rightcolumn}
\end{paracol}

\vspace{0.3cm}

\textbf{\dev{सूचना} / Note on verse numbering:} The original text numbers this as the 23rd \textit{adhyāya} in the internal count (\dev{इति ब्रह्मगुप्ते मानाध्यायः (२३ अध्याय समाप्तः)}), though the final verse calls it the 24th in another reckoning. The manuscript tradition contains variations in chapter enumeration.

\wavydivider

\begin{bilingualverse}
\begin{leftcolumn*}
\dev{नष्टंघ्नसावनदिनात् सूर्यदिनां त्वसावन दिनानि ।}\\
\dev{यस्मात्तस्मादश्वं दुर्ल्लभगं मन्दबुद्धिनाम् ॥ ५ ॥}
\end{leftcolumn*}
\begin{rightcolumn}
``Multiply the elapsed \textit{sāvana} (civil) days by the solar days; the \textit{asāvana} (non-civil) days---therefore, that attainment is difficult for those of dull intellect.''\\[0.3cm]
\hfill\textbf{— Verse 5 (alt.)}
\end{rightcolumn}
\end{bilingualverse}

\begin{bilingualverse}
\begin{leftcolumn*}
\dev{मानुष्यं पित्रदेवं बाहु यान्यष्टौ वसुर्कालस्य ।}\\
\dev{उत्कर्षि ज्ञानार्थं बाहु स्पष्टं नवममन्यत् ॥ ६ ॥}
\end{leftcolumn*}
\begin{rightcolumn}
``Humanity, the father, the divine, the arm, the eight (\textit{vasu}), the time---the \textit{utkarṣi} (elevation) for the purpose of knowledge, the arm, clearly the ninth (or another).''\\[0.3cm]
\hfill\textbf{— Verse 6 (alt.)}
\end{rightcolumn}
\end{bilingualverse}

\begin{bilingualverse}
\begin{leftcolumn*}
\dev{द्वौ द्वौ राशिं मकराहतवः षट् सूर्यगति वशाऽघोज्याः ।}\\
\dev{शकिरवसन्तं ग्रीष्मा वर्षा शरदौ सहैमन्ताः ॥ ७ ॥}
\end{leftcolumn*}
\begin{rightcolumn}
``Two-two (make) a zodiacal sign; multiplied by Makara (10, i.e.\ 60), these (are) the six (seasons); the solar motion (gives) the seasons: śiśira-vasanta (spring), grīṣma (summer), varṣā (rains), śarad (autumn), and hemanta (winter).''\\[0.3cm]
\hfill\textbf{— Verse 7 (alt.)}
\end{rightcolumn}
\end{bilingualverse}

\begin{paracol}{2}
\begin{leftcolumn*}
\dev{द्वौ द्वौ राशिं} — मिथुनादीनि। \dev{मकराहताः षट्} — दशभिः संख्यया गुणिताः षडृतवः। \dev{सूर्यगतिः} — भानोर्मार्गः। \dev{शिशिर-वसन्त-ग्रीष्म-वर्षा-शरद्-हेमन्ताः} — षडृतवः॥
\end{leftcolumn*}
\begin{rightcolumn}
\textit{dvau dvau rāśiṃ} --- Gemini etc.; \textit{makarāhatāḥ ṣaṭ} --- six seasons multiplied by ten (Makara = 10); \textit{sūryagatiḥ} --- the sun's path; \textit{śiśira-vasanta-grīṣma-varṣā-śarad-hemantāḥ} --- the six seasons.
\end{rightcolumn}
\end{paracol}

\vspace{0.3cm}

\wavydivider

\begin{bilingualverse}
\begin{leftcolumn*}
\dev{मुख्यारः पूजोभक्तः काके व्यासान्तरेण रविकर्णं मुमुच्या ।}\\
\dev{द्वौ, क्षाया दिश्यन्त्वे चन्द्रकर्णेन शेषम् ॥ ८ ॥}
\end{leftcolumn*}
\begin{rightcolumn}
``The principal and the \textit{arha} (worthy), the worshipper and the devotee; by the \textit{vyāsāntara} (diameter difference) for the crow---the sun's hypotenuse is to be released; two, the decrease is to be shown, the remainder by the moon's hypotenuse.''\\[0.3cm]
\hfill\textbf{— Verse 8 (alt.)}
\end{rightcolumn}
\end{bilingualverse}

\begin{bilingualverse}
\begin{leftcolumn*}
\dev{पञ्चदशाहीनयुक्ताः षट् क्रान्तिकालाद्यरसगुणाः ॥ ६४ ॥}
\end{leftcolumn*}
\begin{rightcolumn}
``The six, diminished and increased by fifteen, are the declination-time etc., multiplied by six (\textit{rasa}).''\\[0.3cm]
\hfill\textbf{— Verse 64 (alt.)}
\end{rightcolumn}
\end{bilingualverse}

\begin{bilingualverse}
\begin{leftcolumn*}
\dev{त्रिज्या विपुलछाया घातस्तच्छति विभाजिताक्षाप्तम् ।}\\
\dev{क्षयो नित्यं याम्यो गोलवशाद्विम्बेक्रान्तिः ॥ ६६ ॥}
\end{leftcolumn*}
\begin{rightcolumn}
``The radius and the abundant shadow---their product, divided by that and multiplied by the latitude; the decrease is always in the south (Yāmya); from the sphere, the disc's declination.''\\[0.3cm]
\hfill\textbf{— Verse 66 (alt.)}
\end{rightcolumn}
\end{bilingualverse}

\begin{bilingualverse}
\begin{leftcolumn*}
\dev{क्रान्त्यक्षायुति विभोगाक्षरापदैः शोधिते दिनदलमा ।}\\
\dev{भास्य तिकृतयोः कृतमनुयुतो नयोस्तत्पदे व्यस्ते ॥ ६७ ॥}
\end{leftcolumn*}
\begin{rightcolumn}
``When the sum and difference of declination and latitude are purified by the \textit{akṣarāpada} etc., the half-day; when the square of that is subtracted from the square of the sum---the method when reversed at that position.''\\[0.3cm]
\hfill\textbf{— Verse 67 (alt.)}
\end{rightcolumn}
\end{bilingualverse}

\begin{bilingualverse}
\begin{leftcolumn*}
\dev{पठगुणीतागतं शेषं नाडूच्चै दिवसाङ्कं विभाजितात् त्यात् ।}\\
\dev{दिनदलकर्णगुणाप्ता तया त्रिज्यया फलं कर्णः ॥ ६८ ॥}
\end{leftcolumn*}
\begin{rightcolumn}
``The remainder, multiplied by the prescribed (number) and arrived at, divided by the day-number elevated by the \textit{nāḍī}; that (result) multiplied by the half-day hypotenuse, by that radius the result (is) the hypotenuse.''\\[0.3cm]
\hfill\textbf{— Verse 68 (alt.)}
\end{rightcolumn}
\end{bilingualverse}

\wavydivider

\begin{bilingualverse}
\begin{leftcolumn*}
\dev{इत्युक्तं क्रमेण सांख्याप्रकरणं प्रयोगविधिश्च ।}\\
\dev{व्यासगुणा दीर्घवृत्तं षड़ांशकक्षापयाम् ॥ ९ ॥}
\end{leftcolumn*}
\begin{rightcolumn}
``Thus has been stated in order the section on numerals and the practical application; the diameter multiplied, the long circle, the \textit{ṣaḍāṃśaka} (sixth part), the latitude-decline.''\\[0.3cm]
\hfill\textbf{— Verse 9 (alt.)}
\end{rightcolumn}
\end{bilingualverse}

\begin{bilingualverse}
\begin{leftcolumn*}
\dev{तस्य व्यास शकि करण्डूतद्विजयया गुणालिप्ता ॥ १० ॥}
\end{leftcolumn*}
\begin{rightcolumn}
``Its diameter, śaki ( potency), \textit{karaṇḍū}, that (is) multiplied by the \textit{dvijaya} (twice-conquered), by degrees.''\\[0.3cm]
\hfill\textbf{— Verse 10 (alt.)}
\end{rightcolumn}
\end{bilingualverse}

\begin{bilingualverse}
\begin{leftcolumn*}
\dev{रविकर्णं दूता त्रिजया काके व्यासान्तरं हृता शेषम् ।}\\
\dev{तिजयामुख्यास बधा क्षयि करणं हुतात्मो व्यासः ॥ १० ॥}
\end{leftcolumn*}
\begin{rightcolumn}
``The sun's hypotenuse, the messengers, by the \textit{trijaya} (triple-conquered), by the \textit{vyāsāntara} (diameter difference) divided, the remainder; the \textit{tijayāmukhyāḥ} (chief of the three-conquered), by the decrease, the operation, the \textit{hutātma} (fire-souled = 3), the diameter.''\\[0.3cm]
\hfill\textbf{— Verse 10 (alt. 2)}
\end{rightcolumn}
\end{bilingualverse}

\begin{bilingualverse}
\begin{leftcolumn*}
\dev{व्यासेन दुर्गति बधात् काके व्यासान्तरके मुक्ति बधम् ।}\\
\dev{प्रोक्ष् दुमध्यमुक्तथा तिथिगुणाप्तां तमो व्यासः ॥ ११ ॥}
\end{leftcolumn*}
\begin{rightcolumn}
``From the diameter, the evil, from the product, the crow---the diameter difference, the release (is) the product; the \textit{prokṣa} (sprinkling), the second middle, thus from the \textit{tithi} (lunar day) multiplied and obtained, the darkness, the diameter.''\\[0.3cm]
\hfill\textbf{— Verse 11 (alt.)}
\end{rightcolumn}
\end{bilingualverse}

\begin{bilingualverse}
\begin{leftcolumn*}
\dev{योऽधिकमासावमरात्रसम्भवनः सर्वैर्मानानि ।}\\
\dev{प्रायःश्लोकैर्द्दशाभिर्मानाध्यायश्चतुर्विंशः ॥ १२ ॥}
\end{leftcolumn*}
\begin{rightcolumn}
``The intercalary months, the \textit{avama} (omitted) days, the \textit{sambhavana} (production), all the measures; by ten \textit{prāyaśloka}s (resumé verses), the twenty-fourth Measurement Chapter.''\\[0.3cm]
\hfill\textbf{— Verse 12 (alt.)}
\end{rightcolumn}
\end{bilingualverse}

\vspace{0.5cm}
\begin{center}
\fbox{\dev{\textbf{इति ब्रह्मगुप्ते मानाध्यायः (२३ अध्याय समाप्तः)}}}\\[0.3cm]
\textit{Thus ends the twenty-third chapter on Measurement\\in the work of Brahmagupta.}
\end{center}

\newpage

%% ========================================================================
%% PUBLISHER INFORMATION (Page 745)
%% ========================================================================
\section*{Publisher Information / \dev{प्रकाशकस्य विवरणम्}}
\addcontentsline{toc}{section}{Publisher Information (p.\ 745)}

\begin{paracol}{2}
\sanskritlabel
\englishlabel
\switchcolumn[0]*

\begin{center}
\devbold{प्रकाशकः — Publisher}
\end{center}

\begin{simpleverse}
\dev{इण्डियन इन्स्टीट्यूट ऑफ अस्ट्रानॉमिकल}\\
\dev{एण्ड संस्कृत रिसर्च}\\
\dev{२२३६, गुरुद्वारा रोड, करोल बाग,}\\
\dev{नई दिल्ली-५ (भारत)}
\end{simpleverse}

\switchcolumn

\begin{center}
\textbf{Publisher}
\end{center}

\begin{simpleverse}
Indian Institute of Astronomical\\
and Sanskrit Research\\
2236, Gurudwara Road, Karol Bagh,\\
New Delhi-5 (India)
\end{simpleverse}
\end{paracol}

\vspace{0.5cm}

\begin{paracol}{2}
\begin{leftcolumn*}
\begin{center}
\devbold{सम्पादक मण्डल — Editorial Board}
\end{center}

\begin{simpleverse}
\dev{श्री रामस्वरूप शर्मा}\\
\hspace{2em}\dev{प्रधान सम्पादक, संचालक}

\dev{श्री मुकुन्दमिश्र}\\
\hspace{2em}\dev{ज्योतिषाचार्य}

\dev{श्री विद्यनाथ झा}\\
\hspace{2em}\dev{ज्योतिषाचार्य}

\dev{श्री दयाशंकर दीक्षित}\\
\hspace{2em}\dev{ज्योतिषाचार्य}

\dev{श्री श्रीदत्त शर्मा शास्त्री}\\
\hspace{2em}\dev{एम. ए., एम. श्री. एल.}
\end{simpleverse}
\end{leftcolumn*}

\begin{rightcolumn}
\begin{center}
\textbf{Editorial Board}
\end{center}

\begin{simpleverse}
Sri Ramswarup Sharma\\
\hspace{2em}Principal Editor, Director

Sri Mukundamishra\\
\hspace{2em}Jyotiṣācārya

Sri Vidyanath Jha\\
\hspace{2em}Jyotiṣācārya

Sri Dayashankar Dixit\\
\hspace{2em}Jyotiṣācārya

Sri Shridatt Sharma Shastri\\
\hspace{2em}M.A., M.Shri.L.
\end{simpleverse}
\end{rightcolumn}
\end{paracol}

\vspace{0.5cm}

\begin{paracol}{2}
\begin{leftcolumn*}
\begin{center}
\devbold{प्रथम संस्करण — First Edition: १९६६}

\devbold{मूल्य — Price: ₹ ६०.००}

\devbold{मुद्रक — Printer:}\\
\dev{पं. श्री प्रकाशन एण्ड प्रिन्टर्स}\\
\dev{१२, चमेलियन रोड, दिल्ली}
\end{center}
\end{leftcolumn*}

\begin{rightcolumn}
\begin{center}
\textbf{First Edition: 1966}

\textbf{Price: ₹ 60.00}

\textbf{Printer:}\\
Pandit Shri Prakashan and Printers\\
12, Chameleon Road, Delhi
\end{center}
\end{rightcolumn}
\end{paracol}

\vspace{0.5cm}

\textit{This work was published with financial assistance from the Ministry of Education, Government of India (\dev{भारत सरकार के शिक्षा मन्त्रालय द्वारा प्रदत् अनुदान से प्रकाशित}).}

\newpage

%% ========================================================================
%% TITLE PAGE FOR COMMENTARY (Dvītīya-Bhāga)
%% ========================================================================
\thispagestyle{empty}
\begin{center}
\vspace*{2cm}

{\dev{श्रीब्रह्मगुप्ताचार्य-विरचित}}\\[0.5cm]
{\fontsize{24}{28}\selectfont\devbold{ब्राह्मस्फुटसिद्धान्तः}}\\[0.5cm]
{\large (\textsanskrit{saṃskṛta-hindī-bhāṣayā vāsanā-vijñāna-bhāṣyasya sambhaktaḥ soḍapattrikaḥ})}

\vspace{2cm}

{\LARGE\devbold{द्वितीय-भागः}}\\[0.4cm]
{\Large\textit{Dvītīya-Bhāgaḥ --- Second Part}}

\vspace{0.4cm}
{\large\textit{(Vāsanā-vijñāna Sanskrit--Hindi Commentary)}}

\vspace{2.5cm}

\longline

\vspace{1cm}

{\normalsize\dev{प्रधानसम्पादकः --- Principal Editor}}\\[0.3cm]
{\large\devbold{आचार्यवर पण्डित रामस्वरूप शर्मा}}\\[0.2cm]
{\small (\dev{संचालक-इण्डियन इन्स्टीट्यूट ऑफ अस्ट्रानॉमिकल एण्ड संस्कृत रिसर्च})}

\vspace{1.5cm}

{\normalsize\dev{प्रकाशकः --- Publisher}}\\[0.2cm]
{\dev{इण्डियन इन्स्टीट्यूट ऑफ अस्ट्रानॉमिकल एण्ड संस्कृत रिसर्च}}\\
{\dev{गुरुद्वारा रोड, करोल बाग, न्यू देल्ही-५}}

\end{center}
\newpage

%% ========================================================================
%% SUBJECT INDEX (Viṣayānukramaṇikā)
%% ========================================================================
\section*{\dev{विषयानुक्रमणिका} — Subject Index}
\addcontentsline{toc}{section}{Viṣayānukramaṇikā (Subject Index) (pp.\ 750--751)}
\markboth{Subject Index}{}

\begin{center}
\fbox{\fbox{\dev{\Large\textbf{विषयानुक्रमणिका}}}}
\end{center}

The subject index lists all major topics treated in the commentary (\textit{vāsanā-bhāṣya}) portion of the \textit{Brāhmasphuṭasiddhānta}, with page references to the commentary pages (numbered 1--126 and beyond).

\subsection*{\dev{१. मध्यमाधिकारः} — Mean Computation (pp.\ 1--126)}

\begin{center}
\begin{tabular}{@{}l@{\hspace{1em}}r@{\hspace{2em}}l@{}r@{}}
\toprule
\dhead{Topic / विषयः} & \dhead{Page} & \textbf{English Topic} & \textbf{Page} \\
\midrule
\ditem{मण्डलाचरणम्} & १ & Maṇḍalācaraṇam --- Preliminaries & 1 \\
\ditem{ग्रन्थारम्भप्रयोजनम्} & २ & Granthārambha-prayojanam --- Purpose & 2 \\
\ditem{भचक्रस्य चलनम्} & ३ & Bhacakrasya calanam --- Celestial motion & 3 \\
\ditem{कालप्रवृत्तिः} & ४ & Kāla-pravṛttiḥ --- Commencement of time & 4 \\
\ditem{खण्डस्य विभाग कल्पना} & १४ & Khaṇḍasya vibhāga-kalpanā & 14 \\
\ditem{युगमानम्} & १६ & Yugamānam --- Yuga measurement & 16 \\
\ditem{आर्यभटस्य मतसमीक्षा} & १८ & Āryabhaṭasya mata-samīkṣā & 18 \\
\ditem{मनुप्रमाणानि कल्पप्रमाणम्} & १९ & Manupramāṇāni kalpapramāṇam & 19 \\
\ditem{कल्पविषये आर्यभटमतम्} & २१ & Kalpaviṣaye āryabhaṭamatam & 21 \\
\ditem{रोमकसिद्धान्तमतखण्डनम्} & २३ & Romakasiddhānta-mata-khaṇḍanam & 23 \\
\ditem{ग्रहयुगस्य परिरभाषा} & ३९ & Grahamyugasya paribhāṣā & 39 \\
\ditem{सितरोत्रं शनैश्चरयोः कल्पभगणाः} & ४० & Sitaretro-śanaiścarayoḥ kalpabhagaṇāḥ & 40 \\
\ditem{भच्छ्रमाः कुदिनानि च} & ४२ & Bhacchramāḥ kudināni ca & 42 \\
\ditem{रविवर्षमासराशिमासादयः} & ४४ & Ravivarṣa-māsa-rāśimāsādayaḥ & 44 \\
\ditem{सावननक्षत्रमानवदिनानि} & ४६ & Sāvananakṣatramānavadināni & 46 \\
\ditem{गतकालस्याहर्गणादीनां परिहारः} & ५८ & Gatakālasya āhargaṇādīnāṃ parihāraḥ & 58 \\
\ditem{आर्यभटमतम्} & ५८ & Āryabhaṭamatam --- Āryabhaṭa's view & 58 \\
\ditem{कल्पगताक्षावनाहर्गणाः} & ५९ & Kalpagatākṣāvanāhargaṇāḥ & 59 \\
\ditem{ग्रहमन्दादिमध्यमानां मध्यमानयनम्} & ६५ & Grahamandādi-madhyamānāṃ madhyamānayanam & 65 \\
\ditem{स्वसिद्धान्तनिर्देशनम्} & ७६ & Svasiddhānta-nirdeśanam & 76 \\
\ditem{प्रश्नरात्रौ वारप्रवृत्तिपक्षदर्शनम्} & ७६ & Praśnarātrau vāra-pravṛtti-pakṣadarśanam & 76 \\
\ditem{मध्यमानां देशनियमार्थः} & ७८ & Madhyamānāṃ deśaniyamārthaḥ & 78 \\
\ditem{तदर्थं देशान्तरकर्म्मणा प्रदर्शनम्} & ८२ & Tadarthaṃ deśāntarakarmaṇā pradarśanam & 82 \\
\ditem{वर्षदीर्घ दिनाध्यमदिनादिसाधनम्} & ८६ & Varṣadīrgha dinādhyamadinādi-sādhanam & 86 \\
\ditem{वर्षदीर्घमाससाधनम्} & ८८ & Varṣadīrghamāsa-sādhanam & 88 \\
\ditem{वर्षशानयनं लध्वहर्गणानयनम्} & ९१ & Varṣaśānayanaṃ ladhvahargaṇānayanam & 91 \\
\ditem{रविसितबुधानां कुजगुरुशनिरीश्रोक्ष्णां चानयनम्} & ९६ & Ravisitabudhānāṃ kujaguruśanirīśrekṣṇāṃ cānanayam & 96 \\
\ditem{मध्यमचन्द्रानयनम्} & १०० & Madhyamacandrānayanam & 100 \\
\ditem{वर्षान्तिकाद्धर्गणात् भौमादिग्रहमन्दफलानयनम्} & १०२ & Varṣāntikād dhargaṇāt bhaumādigrahamandaphalānayanam & 102 \\
\bottomrule
\end{tabular}
\end{center}

\newpage

\subsection*{\dev{२. स्फुटाधिकारः} — True Position Computation (pp.\ 160--255)}

\begin{center}
\begin{tabular}{@{}l@{\hspace{1em}}r@{\hspace{2em}}l@{}r@{}}
\toprule
\dhead{Topic / विषयः} & \dhead{Page} & \textbf{English Topic} & \textbf{Page} \\
\midrule
\ditem{स्फुटीकरणस्य प्रयोजनम्} & १६० & Sphuṭīkaraṇasya prayojanam --- Purpose & 160 \\
\ditem{आर्धज्योत्क्रमज्याकथनम्} & १६० & Ārdhjyotkramajyā-kathanam & 160 \\
\ditem{चापज्यानयनम्} & १६९ & Cāpajyānanayam --- Arc-sine computation & 169 \\
\ditem{ज्यातच्छापानयनम्} & १९१ & Jyāta-ccāpānayanam & 191 \\
\ditem{मन्दशीघ्रकेन्द्रयोः केन्द्रभुजकोटिज्ययो व्य परिर्मापा} & १९५ & Mandaśīghrakeśrayoḥ... & 195 \\
\ditem{स्फुटपरिरक्ष्यानयनम्} & १९९ & Sphuṭaparirakṣyānayanam & 199 \\
\ditem{रविचन्द्रयोः स्फुटत्वार्थः} & १९८ & Ravicandrayoḥ sphuṭatvārthaḥ & 198 \\
\ditem{आर्यभटोक्तं स्फुटीकरणम्युत्कम्} & १९८ & Āryabhaṭoktaṃ sphuṭīkaraṇamyutkam & 198 \\
\ditem{रविचन्द्रयोः स्फुटीकरणार्थं मन्दपरिरक्ष्यांशाः} & १९९ & Mandaparirakṣyāṃśāḥ & 199 \\
\ditem{दृष्टे नते स्फुटपरिरक्ष्यानयनम्} & २८४ & Dṛṣṭe nate sphuṭaparirakṣyānayanam & 284 \\
\ditem{मन्दफलस्य धनत्वमृणात्वं च} & २८५ & Mandaphalasya dhanatvamṛṇātvaṃ ca & 285 \\
\ditem{रविचन्द्रयोः स्फुटीकरणो विशेषः} & २८७ & Ravicandrayoḥ sphuṭīkaraṇo viśeṣaḥ & 287 \\
\ditem{ग्रहयोः सूर्याचन्द्रमसोर्न्तकालः} & १९२ & Grahamyoḥ sūryācandramasorntakālaḥ & 192 \\
\bottomrule
\end{tabular}
\end{center}

\subsection*{\dev{द. चन्द्र छायाधिकारः} — Lunar Shadow Chapter (pp.\ 501--548)}

\begin{center}
\begin{tabular}{@{}l@{\hspace{1em}}r@{\hspace{2em}}l@{}r@{}}
\toprule
\dhead{Topic / विषयः} & \dhead{Page} & \textbf{English Topic} & \textbf{Page} \\
\midrule
\ditem{आरम्भप्रयोजनम्} & ५०१ & Ārambha-prayojanam & 501 \\
\ditem{कर्तव्यताकथनम्} & ५०१ & Kartavyatā-kathanam & 501 \\
\ditem{चन्द्रस्य दूयाद्दूयत्वम्} & ५०२ & Cakrasya dūyād-dūyatvam & 502 \\
\ditem{चन्द्रैतकालसाधनम्} & ५०३ & Caitrakāla-sādhanam & 503 \\
\ditem{चन्द्रशङ्क्वानयनम्} & ५०८ & Candraśaṅkvānayanam & 508 \\
\ditem{रविचन्द्रयोः स्फुटशङ्क्वानयनम्} & ५०५ & Ravicandrayoḥ sphuṭaśaṅkvānayanam & 505 \\
\ditem{चन्द्रस्य स्फुटच्छायासाधनप्रकारः} & ५०८ & Sphuṭacchāyā-sādhanaprakāraḥ & 508 \\
\ditem{मध्यच्छायासाधनप्रकारः} & ५०९ & Madhyacchāyā-sādhanaprakāraḥ & 509 \\
\ditem{अध्यायोपसंहारः} & ५११ & Adhyāyopasaṃhāraḥ --- Chapter conclusion & 511 \\
\bottomrule
\end{tabular}
\end{center}

\subsection*{\dev{६. ग्रहयुत्यधिकारः} — Planetary Conjunction Chapter (pp.\ 515--567)}

\begin{center}
\begin{tabular}{@{}l@{\hspace{1em}}r@{\hspace{2em}}l@{}r@{}}
\toprule
\dhead{Topic / विषयः} & \dhead{Page} & \textbf{English Topic} & \textbf{Page} \\
\midrule
\ditem{ग्रहणां मध्यमशरकला मध्यमविम्बकला} & ५१५ & Madhyamaśarakalā & 515 \\
\ditem{ग्रहविम्बकला स्फुटीकरणम्} & ५१७ & Grahavimbakalā sphuṭīkaraṇam & 517 \\
\ditem{युतिकालज्ञानार्थं चालनफलज्ञानार्थंश्र} & ५२१ & Yuti-kāla-jñānārtham & 521 \\
\ditem{चालनफलसंस्कारदारा ग्रहयोः सम्लिप्तीकरणार्थः} & ५२३ & Cālanaphalasaṃskāra-dvārā... & 523 \\
\ditem{स्फुटपातानयनम्} & ५२४ & Sphuṭapātānayanam & 524 \\
\ditem{गणितागतदेवपाताद्यमध्यमसंस्कारसाधनोपायः} & ५२७ & Gaṇitāgata-devapātādi... & 527 \\
\ditem{युतिकाले ग्रहशरसाधनम्} & ५२८ & Yutikāle grahaśara-sādhanam & 528 \\
\ditem{युतिकाले ग्रहयोद्दिष्टोपान्तरः} & ५४४ & Yutikāle grahayoddiṣṭopāntaraḥ & 544 \\
\bottomrule
\end{tabular}
\end{center}

\newpage

%% ========================================================================
%% MADHYAMĀDHIKĀRA COMMENTARY
%% ========================================================================
\section*{\dev{मध्यमाधिकारः} — Mean Computation Chapter}
\addcontentsline{toc}{section}{Madhyamādhikāra Commentary (pp.\ 752--838)}
\markboth{Madhyamādhikāra Commentary}{}

\lettrine[lines=3,loversize=0.1]{T}{\textit{he Madhyamādhikāra}} constitutes the first and foundational chapter of the commentary (\textit{vāsanā-vijñāna-bhāṣya}) on the \textit{Brāhmasphuṭasiddhānta}. It treats the computation of mean positions of planets (\textit{madhyama-graha}), beginning with a discussion of the motion of the celestial sphere (\textit{bhacakra}), the definition of time (\textit{kāla-pravṛtti}), and the various \textit{yuga}-systems employed in Indian astronomy.

\subsection*{Introductory Discussion / \dev{आरम्भकथनम्} (pp.\ 752--760)}

The commentary opens with a discussion of the \textit{bhacakra} (celestial sphere) and its motion:

\begin{bilingualverse}
\begin{leftcolumn*}
\dev{प्रदक्षिणागम् (दक्षिणावस्स्क्ष्मेण चलायमानम्) श्रादो (सर्व्वप्रथमं) ब्रह्मणा}\\
\dev{(जगद्गुरुपादकेन) सृष्टम् (रचितम्) ग्रयदिध्नयादौ स्थितेनक्षत्रादिमततुल्यचपा-}\\
\dev{तादिमित्र साकं भ्रुवयष्ट्याद्यारेण प्रमाणशील ज्योतिष्मत् प्रदायिनां नक्षत्रादीनां}\\
\dev{चक्रं (गोलं) ब्रह्मणा रचितं यथाप्याचार्ययोः कादीनां वचां न कियते तेषां स्वरूपा-}\\
\dev{भावात् किन्तु स्पष्टादिकाले तेषामपि स्थानाऽऽज्ञानरूपसर्जें कृतमिति ॥ १ ॥}
\end{leftcolumn*}
\begin{rightcolumn}
``(The celestial sphere), moving eastward (from left to right), was first created by Brahmā (the world-teacher). The sphere, along with the nakṣatras and other celestial bodies having equal \textit{capādas}, along with the meridian staff etc., the circle of nakṣatras and other luminaries that give light---this sphere was fashioned by Brahmā. Even so, the ācāryas (teachers) do not describe their (the planets') essential nature (\textit{svārūpa}), but have only determined their positions (\textit{sthāna}) for the purpose of computation (\textit{spaṣṭādi}).''\\[0.3cm]
\hfill\textbf{— Commentary v.\ 1}
\end{rightcolumn}
\end{bilingualverse}

\textbf{\dev{हिन्दीटीका} / Hindi Commentary (pp.\ 752--753):}
The commentator explains that the celestial sphere (\textit{bhacakra}) was first created by Brahmā (the creator). The sphere, along with the nakṣatras and other celestial bodies, revolves from west to east. The commentator notes that the \textit{ācāryas} (teachers) do not discuss the nature (\textit{svārūpa}) of these bodies but have only determined their positions (\textit{sthāna}) for the purposes of computation.

\wavydivider

\subsection*{Key Verses from the Siddhānta / \dev{सिद्धान्तमूलश्लोकाः} (pp.\ 760--770)}

The commentary elucidates the following root text (\textit{mūla}) verses:

\begin{bilingualverse}
\begin{leftcolumn*}
\dev{इदानीमनाख्यानतस्य कालस्य प्रवृत्तिमाह ।}\\
\dev{चैनंसितादेर्द्दयाद्व मानोर्द्धनमासबर्षयुगकल्पः ।}\\
\dev{सृष्ट्यादौ लङ्कायां स्मं प्रवृतta दिनेस्कस्य ॥ ४ ॥}
\end{leftcolumn*}
\begin{rightcolumn}
``Now (Brahmagupta) declares the commencement of that time whose beginning is unmentioned. Day, month, year, \textit{yuga}, and \textit{kalpa} began from the sun at Laṅkā at the beginning of creation.''\\[0.3cm]
\hfill\textbf{— Mūla-śloka 4}
\end{rightcolumn}
\end{bilingualverse}

\textbf{Hindi commentary (pp.\ 761--762):}

\begin{paracol}{2}
\begin{leftcolumn*}
\dev{चैत्यादि कल्पप्रतिपद्यात् लंकोपलक्षितांभृद्वेकोर्कोदयादिदिवा प्रवृत्ताः} — चैत्रशुक्लप्रतिपदादिकालात् लङ्कायां सूर्योदयात् दिनादयः प्रवृत्ताः।

\dev{रविदिनम्बरः सृष्ट्यादौ कल्पादौ} — सृष्ट्यारम्भे कल्पस्य आदौ रवेर्दिनसंख्या।
\end{leftcolumn*}
\begin{rightcolumn}
\textit{caityādi kalpa-pratipadyāt laṅkopalakṣitāṃbhṛdv ekodayādi-divā pravṛttāḥ} --- From the first tithi of Caitra (the first month), at Laṅkā, the days etc.\ commenced from the sunrise.

\textit{ravidinambaraḥ sṛṣṭyādau kalpādau} --- The solar day number at the beginning of creation, at the beginning of the \textit{kalpa}.
\end{rightcolumn}
\end{paracol}

\vspace{0.3cm}
The commentator explains the concept of the \textit{kuldiņa} (civil day) in detail. The number of civil days elapsed since the beginning of the current kalpa is 1,58,22,23,78,50,000 according to Brahmagupta's system.

\wavydivider

\subsection*{On Time Computation and Planetary Motion / \dev{कालगणनम् ग्रहचलनं च} (pp.\ 770--790)}

The commentary proceeds with detailed mathematical analysis:

\begin{bilingualverse}
\begin{leftcolumn*}
\dev{प्रथमं युगपदं चलनं द्वितीयं पश्चादिनपर्व्वैः}\\
\dev{तृतीयं तु त्रयोदशभिर्मासैः सावनेन वा चतुर्थम् ॥}
\end{leftcolumn*}
\begin{rightcolumn}
``First, the \textit{yugapada} (yuga-epoch method); second, by \textit{parva} (lunisolar method); third, by thirteen months; or fourth, by \textit{sāvana} (civil reckoning).''\\[0.3cm]
\hfill\textbf{— Time-reckoning verse}
\end{rightcolumn}
\end{bilingualverse}

\textbf{\dev{व्याख्या} / Exposition:} The four types of time-reckoning are discussed:
\begin{enumerate}[label=(\roman*),leftmargin=2em]
  \item \textit{Yugapada} --- the yuga-epoch method
  \item \textit{Parva} --- the lunisolar intercalation method
  \item \textit{Trayodaśa-māsa} --- the thirteen-month intercalary method
  \item \textit{Sāvana} --- civil reckoning based on sunrise
\end{enumerate}
Each method is explained with its mathematical basis and astronomical application.

\wavydivider

\subsection*{Planetary Computations / \dev{ग्रहगणनविधिः} (pp.\ 790--820)}

The commentary provides detailed algorithms for computing mean planetary positions:

\textbf{Mean longitude formula / \dev{मध्यमीभावनसूत्रम्}:}
\begin{equation*}
\text{Mean longitude} = \frac{\text{ahargaṇa} \times \text{daily motion}}{\text{civil days in a kalpa}} \mod 360°
\end{equation*}

Where:
\begin{itemize}[leftmargin=2em]
  \item \textit{ahargaṇa} = elapsed civil days since kalpa-beginning
  \item \textit{daily motion} = mean daily motion of the planet
  \item The result is expressed in degrees of the ecliptic
\end{itemize}

The commentary discusses the \textit{mandaphala} (equation of centre) and \textit{śīghraphala} (\textit{śīghra} equation) with detailed trigonometric formulas using the \textit{jyā} (sine) function.

\begin{paracol}{2}
\begin{leftcolumn*}
\dev{मन्दफलं रविचन्द्रयोः — मन्दकendraस्य ज्या गुणकः, त्रिज्या भाजकः।}

\dev{शीघ्रफलं बुधगुरुशनिभ्यः — शीघ्रकेन्द्रस्य ज्या गुणकः।}
\end{leftcolumn*}
\begin{rightcolumn}
\textit{Mandaphala} for the Sun and Moon: the sine of the \textit{manda-kendra} (anomaly) is the multiplier, the radius (\textit{trijyā}) is the divisor.

\textit{Śīghraphala} for Mercury, Jupiter, and Saturn: the sine of the \textit{śīghra-kendra}.
\end{rightcolumn}
\end{paracol}

\wavydivider

\subsection*{On the Computation of Weekdays / \dev{वारप्रवृत्तिकथनम्} (pp.\ 820--838)}

The final section of this portion treats the computation of weekdays (\textit{vāra-pravṛtti}):

\begin{bilingualverse}
\begin{leftcolumn*}
\dev{वारप्रवृत्तिं मनयो वदन्ति सूर्योदयाद्विष्णुराज्ञाध्याम् ।}\\
\dev{उर्वी तथाऽघोषपरत्र तस्याऽर्चराथं देशान्तरनाडिकाभिः ॥ १३६॥}
\end{leftcolumn*}
\begin{rightcolumn}
``The sages declare the commencement of the week (to be reckoned) from the sunrise (at Laṅkā); on the earth (at other places), it (is determined) by adding or subtracting the \textit{deśāntara-nāḍikās} (longitude difference in time).''\\[0.3cm]
\hfill\textbf{— Verse 136}
\end{rightcolumn}
\end{bilingualverse}

\begin{paracol}{2}
\begin{leftcolumn*}
\dev{लङ्कोदयाम्यसुजातप्रथममपरतः पूर्व्वदेशे च पश्चात्—}\\
\dev{दर्शवोत्थानिमष्टटीभिः सवितः क्षदयतो वासरैव प्रवृत्तिः ॥}\\
\dev{यं या सूर्योदयात् प्राक् चरराकलम्बवैश्चास्मियमियगोले ।}\\
\dev{पश्चात् सौम्यगोले युतिर्वियुतिर्वशाच्छोभयोः स्पष्टकालः इति ॥१३६॥}
\end{leftcolumn*}
\begin{rightcolumn}
``At Laṅkā, the weekday begins with the first \textit{nāḍī} after sunrise; at a place to the west, (it begins) after (sunrise); with the six or eight \textit{nāḍī}s (difference), the weekday is determined when the sun rises. For a place on the eastern or western side of the prime meridian, before or after sunrise (at Laṅkā), by addition or subtraction, the correct time (is found).''\\[0.3cm]
\hfill\textbf{— Expanded v.\ 136}
\end{rightcolumn}
\end{bilingualverse}

\textbf{\dev{हिन्दीटीका} / Hindi Commentary (pp.\ 836--838):}

\begin{paracol}{2}
\begin{leftcolumn*}
\dev{लङ्कायां सूर्योदयात् प्रथमनाड्याम् वारारम्भः। पश्चिमदेशे पश्चात् षडष्टघटीभिः।}

\dev{प्राच्यां देशे सूर्योदयात् प्राक् स्पष्टकालः।}
\end{leftcolumn*}
\begin{rightcolumn}
At Laṅkā, the weekday begins at sunrise in the first \textit{nāḍī}. At a place to the west, after (sunrise) by six or eight \textit{ghaṭī}s.

At a place to the east, the correct time (is found) before sunrise.
\end{rightcolumn}
\end{paracol}

The commentator explains that the \textit{vāra} (weekday) begins at sunrise. By this definition, at a place east of the prime meridian (Laṅkā), the weekday begins after subtracting the longitude difference (in \textit{ghaṭī}s); at a place west of the meridian, (the weekday begins) before sunrise by the same amount (longitude difference in \textit{ghaṭī}s).

In the \textit{Siddhānta-śiromaṇi}, Bhāskarācārya states: ``\textit{arkodayād dvābhyāṃ bhyas tāniḥ prācyāṃ praticyāṃ dinap-}''---(the computation proceeds from sunrise with modifications for eastern and western positions).

\wavydivider

\subsection*{Concluding Verses / \dev{उपसंहारश्लोकाः} (p.\ 838)}

The portion concludes with a verse from the \textit{Manusmṛti}:

\begin{bilingualverse}
\begin{leftcolumn*}
\dev{आसीदिदं तमोभूतमप्रज्ञातमलक्षणम् ।}\\
\dev{अप्रतर्क्यमनाऽऽगम्यं प्रसुप्तमिव सर्व्वतः ॥ इति ॥}
\end{leftcolumn*}
\begin{rightcolumn}
``This (universe) existed as darkness, unknowable, without marks, not to be reasoned about, not to be apprehended, as if asleep everywhere.''\\[0.3cm]
\hfill\textbf{— Manusmṛti 1.5}
\end{rightcolumn}
\end{bilingualverse}

\textbf{\dev{टीकाकारवचनम्} / Commentator's exposition:}

The commentator explains that Manu has stated that at the time of the deluge (\textit{pralaya}), the perception of weekdays and other time measures is not possible. Therefore, the reckoning of weekdays is only possible during the manifest (\textit{sṛṣṭi}) period. Brahmagupta establishes that the weekday begins from sunrise at Laṅkā and is modified by the longitude difference (\textit{deśāntara}) for other locations.

\newpage

%% ========================================================================
%% APPENDIX: STRUCTURAL OVERVIEW OF THE COMMENTARY
%% ========================================================================
\section*{Appendix: Structural Overview of the Commentary / \dev{टीकायाः संरचनासूची}}
\addcontentsline{toc}{section}{Appendix: Structural Overview}

The \textit{Vāsanā-vijñāna} commentary on the \textit{Brāhmasphuṭasiddhānta} is organized into the following major chapters (\textit{adhyāya}s):

\begin{center}
\begin{tabular}{@{}clc@{}}
\toprule
\dhead{No.} & \dhead{Title / \dev{शीर्षकः}} & \textbf{Pages} \\
\midrule
\ditem{१} & \ditem{मध्यमाधिकारः} --- Mean Computation & 1--126 \\
\ditem{२} & \ditem{स्फुटाधिकारः} --- True Position Computation & 160--255 \\
\ditem{३} & \ditem{त्रिप्रश्नाधिकारः} --- Three Questions (Diurnal) & 256--344 \\
\ditem{४} & \ditem{चन्द्रग्रहणाधिकारः} --- Lunar Eclipse & 345--400 \\
\ditem{५} & \ditem{चन्द्रच्छायाधिकारः} --- Lunar Shadow & 501--548 \\
\ditem{६} & \ditem{ग्रहयुत्यधिकारः} --- Planetary Conjunction & 515--567 \\
\ditem{७} & \ditem{नक्षत्रच्छायाधिकारः} --- Stellar Shadow & --- \\
\ditem{८} & \ditem{ग्रहच्छायाधिकारः} --- Planetary Shadow & --- \\
\ditem{९} & \ditem{उदयास्तमयाधिकारः} --- Heliacal Rising/Setting & --- \\
\ditem{१०} & \ditem{पाताधिकारः} --- Nodes and Intersections & --- \\
\ditem{११} & \ditem{गणिताध्यायः} --- Mathematics (Chs.\ 12--18) & --- \\
\ditem{१२} & \ditem{कुण्डकाध्यायः} --- Kuṭṭaka (Pulverizer) & --- \\
\ditem{१३} & \ditem{परिकर्माध्यायः} --- Auxiliary Computations & --- \\
\bottomrule
\end{tabular}
\end{center}

\divider

\vspace{0.5cm}
\begin{center}
\textit{The commentary employs a rich methodology, interweaving:}
\begin{enumerate}[label=(\roman*),leftmargin=2em]
  \item \textit{Mūla-śloka} --- The root text verses of Brahmagupta
  \item \textit{Anvaya} --- Word-order rearrangement for grammatical clarity
  \item \textit{Bhāṣya} --- Sanskrit explanatory commentary
  \item \textit{Hindi Ṭīkā} --- Hindi translation and explanation
  \item \textit{Upapatti} --- Mathematical demonstrations and proofs
  \item \textit{Udāharaṇa} --- Numerical examples and worked problems
\end{enumerate}
\end{center}

\vspace{1cm}
\begin{center}
\longline

\textbf{\Large\dev{॥ इति समाप्तम् ॥}}

\vspace{0.5cm}
\textit{End of Part 1, Chunk 6 --- Pages 701--838}\\[0.3cm]
\textit{Bilingual Sanskrit-English Edition}

\vspace{0.3cm}
\longline

\vspace{0.5cm}
\begin{paracol}{2}
\begin{leftcolumn*}
\begin{center}
\dev{श्रीब्रह्मगुप्ताचार्येण विरचितः}\\
\dev{ब्राह्मस्फुटसिद्धान्तः}\\[0.3cm]
\dev{प्रथमभागस्य षष्ठः संचिकाङ्गं समाप्तम्}\\
\dev{(पृष्ठानि ७०१--८३८)}
\end{center}
\end{leftcolumn*}
\begin{rightcolumn}
\begin{center}
\textit{The Brāhmasphuṭasiddhānta}\\
\textit{composed by Śrī Brahmaguptācārya}\\[0.3cm]
\textit{The sixth portion of the first part is concluded}\\
\textit{(Pages 701--838)}
\end{center}
\end{rightcolumn}
\end{paracol}

\vspace{0.5cm}
\longline
\end{center}

\end{document}
%% ========================================================================
%% END OF DOCUMENT
%% =================================================


% ============================================================
% Source: translations/en/indian/brahmagupta-brahmasphutasiddhanta-pt2-chunk1_bilingual.tex
% ============================================================

% ============================================================================
% Brahmagupta's Brahmasphutasiddhanta Part 2 — Chunk 1 (Pages 839–978)
% BILINGUAL English-Sanskrit Edition
% Parallel minipage columns: Sanskrit (Devanagari) + English translation with IAST
% ============================================================================
\documentclass[11pt,a4paper]{book}

%% -------- Core Packages --------
\usepackage{fontspec}
\usepackage{polyglossia}
\usepackage{amsmath,amssymb}
\usepackage{geometry}
\usepackage{xcolor}
\usepackage{titlesec}
\usepackage{fancyhdr}
\usepackage{array,booktabs,longtable,multirow}
\usepackage{hyperref}
\usepackage{ragged2e}

%% -------- Page Geometry --------
\geometry{
  paperwidth=210mm,paperheight=297mm,
  left=18mm,right=18mm,top=25mm,bottom=25mm,
  headheight=14pt,footskip=30pt
}

%% -------- Language Setup --------
\setmainlanguage{english}

%% -------- Font Configuration --------
%% Use ExternalLocation for all fonts to ensure xelatex finds them
\setmainfont[
  Extension=.otf,
  UprightFont=/usr/share/fonts/opentype/linux-libertine/LinLibertine_R,
  BoldFont=/usr/share/fonts/opentype/linux-libertine/LinLibertine_RB,
  ItalicFont=/usr/share/fonts/opentype/linux-libertine/LinLibertine_RI,
  BoldItalicFont=/usr/share/fonts/opentype/linux-libertine/LinLibertine_RBI]{LinLibertine_R}

\setsansfont[
  Extension=.otf,
  UprightFont=/usr/share/fonts/opentype/linux-libertine/LinBiolinum_R,
  BoldFont=/usr/share/fonts/opentype/linux-libertine/LinBiolinum_RB,
  ItalicFont=/usr/share/fonts/opentype/linux-libertine/LinBiolinum_RI]{LinBiolinum_R}

\setmonofont{DejaVu Sans Mono}

%% Devanagari (Sanskrit/Hindi) Font — user-specified path
\newfontfamily\devanagarifont[Script=Devanagari,Path=/usr/share/fonts/truetype/noto/]{NotoSerifDevanagari-Regular.ttf}

%% -------- Colors --------
\definecolor{titlecolor}{RGB}{80,20,20}
\definecolor{sectioncolor}{RGB}{55,45,35}
\definecolor{rulecolor}{RGB}{180,160,140}

%% -------- Section Formatting --------
\titleformat{\chapter}[display]
  {\normalfont\huge\bfseries\color{titlecolor}}
  {\chaptertitlename\ \thechapter}{20pt}{\Huge}
\titleformat{\section}
  {\normalfont\Large\bfseries\color{sectioncolor}}
  {\thesection}{1em}{}
\titleformat{\subsection}
  {\normalfont\large\bfseries\color{sectioncolor}}
  {\thesubsection}{1em}{}

%% -------- Header/Footer --------
\pagestyle{fancy}
\fancyhf{}
\fancyhead[LE]{\small\textit{Brahmasphu.tasiddhanta — Pt.2 Chunk 1}}
\fancyhead[RO]{\small\textit{\leftmark}}
\fancyhead[RE,LO]{\small\thepage}
\renewcommand{\headrulewidth}{0.4pt}
\renewcommand{\footrulewidth}{0pt}

%% -------- Custom Commands --------
\newcommand{\dev}[1]{{\devanagarifont #1}}

%% Bilingual pair: Sanskrit (left) + English (right)
\newcommand{\bilingual}[2]{%
  \par\vspace{0.6em}%
  \noindent\begin{minipage}[t]{0.48\textwidth}
    \RaggedRight
    #1%
  \end{minipage}%
  \hfill
  \begin{minipage}[t]{0.48\textwidth}
    \RaggedRight
    #2%
  \end{minipage}%
  \par\vspace{0.6em}%
}

%% Shloka (verse) bilingual: Devanagari verse + English translation
\newcommand{\shlokabilingual}[2]{%
  \par\vspace{0.8em}%
  \noindent\begin{minipage}[t]{0.48\textwidth}
    \RaggedRight
    \dev{\textit{#1}}
  \end{minipage}%
  \hfill
  \begin{minipage}[t]{0.48\textwidth}
    \RaggedRight
    #2%
  \end{minipage}%
  \par\vspace{0.8em}%
}

%% Sanskrit commentary bilingual
\newcommand{\sktcommbilingual}[2]{%
  \par\vspace{0.5em}%
  \noindent\begin{minipage}[t]{0.48\textwidth}
    \RaggedRight\small\itshape
    \dev{#1}
  \end{minipage}%
  \hfill
  \begin{minipage}[t]{0.48\textwidth}
    \RaggedRight\small
    #2%
  \end{minipage}%
  \par\vspace{0.5em}%
}

%% Hindi commentary bilingual
\newcommand{\hindicommbilingual}[2]{%
  \par\vspace{0.5em}%
  \noindent\begin{minipage}[t]{0.48\textwidth}
    \RaggedRight\small
    \dev{#1}
  \end{minipage}%
  \hfill
  \begin{minipage}[t]{0.48\textwidth}
    \RaggedRight\small
    #2%
  \end{minipage}%
  \par\vspace{0.5em}%
}

%% Rule heading bilingual
\newcommand{\ruleheadingbilingual}[2]{%
  \par\vspace{0.8em}%
  \noindent\begin{minipage}[t]{0.48\textwidth}
    \textbf{\color{sectioncolor}\dev{#1}}
  \end{minipage}%
  \hfill
  \begin{minipage}[t]{0.48\textwidth}
    \textbf{\color{sectioncolor}#2}
  \end{minipage}%
  \par\vspace{0.4em}%
}

%% Section divider
\newcommand{\secdivider}{%
  \begin{center}\rule{0.4\textwidth}{0.5pt}\end{center}%
}

%% -------- Hyperref Setup --------
\hypersetup{
  colorlinks=true,
  linkcolor=sectioncolor,
  citecolor=sectioncolor,
  urlcolor=sectioncolor,
  pdfencoding=auto,
  pdftitle={Brahmasphutasiddhanta Part 2 Chunk 1 — Bilingual Edition},
  pdfauthor={Brahmagupta}
}

%% ============================================================================

\begin{document}

%% -------- Title Page --------
\begin{titlepage}
\begin{center}
\vspace*{2cm}

{\Huge\bfseries\color{titlecolor} \dev{ब्राह्मस्फुटसिद्धान्त}}\\[1cm]
{\LARGE\color{sectioncolor} The Brāhmasphu.tasiddhānta of Brahmagupta}\\[0.8cm]
{\large\textit{Bilingual Sanskrit--English Edition}}\\[0.3cm]
{\normalsize With English Translation and IAST Transliteration}\\[2cm]

\rule{0.6\textwidth}{1pt}\\[1cm]

{\Large\bfseries Part 2 --- Chunk 1}\\[0.5cm]
{\large Pages 839--978 of the Complete Work}\\[1cm]

\vspace{0.5cm}

{\large\color{sectioncolor}
\textbf{Contents:}\\[0.3cm]
\dev{मध्यमाधिकार} --- Madhyamādhikāra (Chapter on Mean Motions)\\
(Complete with Commentary)\\[0.3cm]
\dev{स्पष्टाधिकार} --- Spa.s.tādhikāra (Chapter on True Positions)\\
(Beginning with Sine Tables and Corrections)\\[0.5cm]
}

\vfill

\rule{0.4\textwidth}{0.5pt}\\[0.5cm]
{\large Parallel-column Bilingual Edition}\\[0.3cm]
{\small Sanskrit (Devanagari) $\|$ English Translation with IAST}\\[0.3cm]
{\footnotesize Typeset in XeLaTeX with Devanagari Support}

\end{center}
\end{titlepage}

\clearpage

%% -------- Introduction --------
\chapter*{Introduction}
\addcontentsline{toc}{chapter}{Introduction}
\markboth{Introduction}{Introduction}

\bilingual{%
\textbf{\dev{प्रस्तावना --- मध्यमाधिकार}}\\[0.3em]
\dev{मध्यमाधिकार} ब्रह्मगुप्तकृत ब्राह्मस्फुटसिद्धान्तस्य प्रथमोऽध्यायः। अत्र ग्रहाणां मध्यमगतिः, दिनगणना, देशान्तरसंस्कारः, युगादिदिनानयनं च प्रदर्श्यते। वासनाभाष्येण सहितोऽयं प्रथमः खण्डः पृष्ठसङ्ख्या ८३९--८९३। द्वितीयः खण्डः \dev{स्पष्टाधिकारः} पृष्ठसङ्ख्या ८९६--९७८।
}{%
\textbf{Introduction --- Madhyamādhikāra}\\[0.3em]
The \textit{Madhyamādhikāra} (Chapter on Mean Motions) is the first chapter of Brahmagupta's \textit{Brāhmasphu.tasiddhānta}. It presents the computation of the mean motions (\textit{madhyama}) of the planets, day computation, longitude correction (\textit{deśāntara}), and calculation of days since epoch (\textit{yugādi-dina}). This first section spans pages 839--893. The second section, the \textit{Spa.s.tādhikāra} (Chapter on True Positions), spans pages 896--978.
}

\vspace{1em}

\bilingual{%
\textbf{\dev{स्पष्टाधिकार-प्रवेशः}}\\[0.3em]
\dev{स्पष्टाधिकारः} द्वितीयोऽध्यायः। यस्मान्मध्यतुल्यः प्रतिदिवसं हृश्यते ग्रहो भगणे, तस्मात्तत्कृत्यकरं मध्यस्फुटीकरणं वक्ष्ये। अत्र ज्यापञ्चविंशतिका, उत्क्रमज्या, मन्दफलं, शीघ्रफलं च विस्तरेण निरूप्यते।
}{%
\textbf{Introduction --- Spa.s.tādhikāra}\\[0.3em]
The \textit{Spa.s.tādhikāra} (Chapter on True Positions) is the second chapter. ``Since the mean planet is not observed daily in the zodiac, I shall state the procedure for making it true'' (BSS II.1). Here the twenty-four sine table, versed sines, equation of centre (\textit{manda-phala}), and equation of conjunction (\textit{śīghra-phala}) are discussed in detail.
}

\vspace{1em}
\secdivider
\clearpage

%% ============================================================================
\chapter[Madhyamādhikāra]{\dev{मध्यमाधिकारः} --- Madhyamādhikāra}
\label{chap:madhyama}

\sectionmark{Madhyamādhikāra}

\bilingual{%
\dev{मध्यमाधिकार} ब्रह्मगुप्तकृत ब्राह्मस्फुटसिद्धान्तस्य प्रथमोऽध्यायः। अस्य विषयाः --- वारप्रवृत्तिः, देशान्तरकर्म, युगादिदिनानयनं, रवीन्द्वोर्मध्यमगतिः, बुधशुक्रादीनां मध्यमगतिश्च। वासनाभाष्येण सह सम्पूर्णोऽयं खण्डः पृष्ठ ८३९--८९३।
}{%
The \textit{Madhyamādhikāra} is the first chapter of the \textit{Brāhmasphu.tasiddhānta} composed by Brahmagupta (b.~598 CE). Its subjects include: the commencement of the weekday (\textit{vāra-prav.rtti}), longitude correction (\textit{deśāntara-karma}), computation of days since epoch (\textit{yugādi-dina-anayana}), the mean motions of the Sun and Moon, and the mean motions of Mercury, Venus, and the other planets. This complete section with commentary (\textit{vāsanābhā.sya}) spans pages 839--893.
}

%% ============================================================================
\section{On the Commencement of the Day (\dev{वारप्रवृत्ति})}

\ruleheadingbilingual{प्रवृत्तिवाद}{Commencement of the Weekday}

\hindicommbilingual{%
\dev{प्रवृत्ति} इससे इसी ब्रह्मगुप्तोक्त बात को कहते हैं। वारप्रवृत्ति कब से होती है, इस विषय में बहुत आचार्यों के भिन्न-भिन्न मत हैं। जैसे आर्यभट, सिंहाचार्य आदि अर्धोदित रवि बिम्ब से दिनारम्भ कहते हैं। अन्य आचार्य दिनान्त से दिनारम्भ कहते हैं। लाटदेव आदि रवि के अर्धास्तकाल से कहते हैं। यवन राजा रात्रि में दश मुहूर्त करके कहते हैं। सिद्धान्तशेखर में श्रीपतिरूपर्युक्त आचार्यमतखण्डनं कृतवान्‌।
}{%
\textbf{[Hindi Commentary]} \textit{Prav.rtti} means that which is spoken by Brahmagupta himself. On the question ``from when does the weekday commence?'' there are many differing opinions among teachers. Āryabha.ta and Siṃhācārya hold that the day begins from the half-risen solar disc. Others hold it begins from sunset. Lāṭadeva and others hold it begins from midnight. The Yavana kings hold that ten \textit{muhūrta}s of night [make the division]. In the \textit{Siddhāntaśekhara}, Śrīpati has refuted the opinions of these teachers.
}

\shlokabilingual{%
सृष्टेर्मुखे ध्वान्तमये हि विश्वे ग्रहेषु सृष्टेष्विनपूर्वकेषु ।
दिनप्रवृत्तिस्तदघीश्वरस्य वारस्य तस्मादुदयात्प्रवृत्तिः ॥३६॥
}{%
\textbf{[BSS I.36]} When at the beginning of creation the universe was filled with darkness, and the planets had not yet been created, there was no lord of the weekday (\textit{vāreśvara}). Therefore the commencement of the weekday is from sunrise (\textit{udayāt prav.rttiḥ}).
}

\hindicommbilingual{%
सृष्ट्यादि से पहले विश्व अन्धकारमय था। ऐसे समय में वारप्रवृत्ति के विचार सम्भव नहीं क्योंकि उस समय वारेश्वरग्रहों का अभाव रहता है। इसी कारण रात्रि में वारप्रवृत्ति कहने वालों के मत ठीक नहीं। मध्याह्न काल और अस्तमय काल से आरम्भ कर वारप्रवृत्ति कहना भी ठीक नहीं है। इसलिए रव्युदय काल से वारप्रवृत्ति मानने वाले आचार्यों का मत ही ठीक है।
}{%
\textbf{[Hindi Commentary --- Translation]} Before the beginning of creation, the universe was full of darkness. At such a time, discussion of the weekday's commencement is impossible because the planet-lords of the weekdays did not yet exist. For this reason, the opinion of those who say the weekday commences at midnight is not correct. [Holding that] it begins from midday or from sunset is also not correct. Therefore the opinion of the teachers who hold that the weekday commences at sunrise is indeed correct.
}

\sktcommbilingual{%
संस्कृतोपपत्तौ ``वारप्रवृत्तिं मुनयो वदन्ति सूर्योदयाद्रावणराजधान्याम्‌'' इत्यादि पद्यात्‌ विद्यते। ``लङ्कोदयाम्यसूत्रार्धमपरतः पूर्वदेशे च पश्चात्‌''--- इत्यादि पद्येन स्पष्टम्‌। इति ॥३६॥
}{%
\textbf{[Sanskrit Commentary]} In the Sanskrit proof-text (\textit{saṃsk.rtopapatti}) it is known from the verse ``The sages state the commencement of the weekday from sunrise at Laṅkā, the city of Rāvaṇa'' etc. This is made clear by the verse ``From Laṅkā sunrise, south of the prime meridian, east is before, west is after'' etc. Thus [ends] verse~36.
}

\secdivider
\clearpage

%% ============================================================================
\section{Longitude Correction (\dev{देशान्तरकर्म})}

\ruleheadingbilingual{नियमाः ३७--३९}{Rules 37--39: Longitude Difference}

\shlokabilingual{%
भूपरिधिः खखखशरारेखा स्वाक्षान्तरांशसंख्यिताः ।
भगणांशहृताः फलकृतिहीना देशान्तरस्य कृतिः ॥३७॥
शेषपदगुणासुक्तिभूपरिधिहता कलादिलब्धसूरणस्‌ ।
उज्जयिनीयाम्योत्तररेखायाः प्राग्धनं पश्चात्‌ ॥३८॥
मध्यग्रहे स्फुटे वा भूपरिधिहृतात्‌ पदात्‌ गुणात्‌ षष्ठया ।
लब्धं घटिका अथवा कर्मतिथिष्वग्रणं ग्रहवत्‌ ॥३९॥
}{%
\textbf{[BSS I.37--39]} The terrestrial circumference (\textit{bhū-paridhi}) is five thousand yojanas (\textit{khakha-kha-śara} = 5000). Multiply this by the degrees of difference in latitude (\textit{svāk.sāntarāṃśa}); divide by 360 (\textit{bhaga.nāṃśa}); the square of the result, subtracted from the square of the longitudinal difference (\textit{deśāntara}) [in yojanas], gives [the square of] the corrected difference. [Verses~38--39] From the root of the remainder, multiplied by the terrestrial circumference and divided [appropriately], [obtain] minutes (\textit{kalā}); to the east of the Ujjayinī north-south line is positive (\textit{dhana}), to the west is negative (\textit{.rna}). Or from the root divided by the terrestrial circumference and multiplied by sixty, the result is in \textit{gha.tikā}s; this is to be applied to the mean or true planet as positive or negative like [the treatment of] a planet.
}

\sktcommbilingual{%
\dev{भूपरिधिः}---भूगोलस्य परिधिः; कियानित्यत आह---\dev{खखखशरा} इति पञ्चसहस्रं योजनानामित्यर्थः। अत्र च देशान्तरकक्ष्योस्तत्यादयः प्रमाणादिविभागेन पूर्वमेव व्याख्याताः।
}{%
\textbf{[Sanskrit Commentary]} ``Terrestrial circumference'' (\textit{bhū-paridhi}) means the circumference of the Earth-globe. In answer to ``how much?'' he says \textit{khakha-kha-śarā} = five thousand yojanas. Here the longitudinal difference, latitude, etc., have been explained previously by means of measurements and divisions.
}

\hindicommbilingual{%
\textbf{हिन्दी भाषा}---भूपरिधि ५००० पाँच हजार है, इसको स्वदेश और रेखादेश के अक्षांशान्तर से गुणा कर ३६० से भाग देने से जो फल हो उसके वर्ग को देशान्तर वर्ग में घटाकर शेष के मूल से ग्रहगति को गुणकर भूपरिधि से भाग देने से कलात्मक फल होता है। उज्जयिनीयाम्योत्तररेखा से पूर्वदेश में मध्यमग्रहे वा स्फुटग्रहে धन, पश्चिमदेशे ऋण। पद (मूल) को साठ से गुणा कर भूपरिधि से भाग देने से घट्यात्मक फल, तिथिमुक्त घटी संस्कार्यम्‌। ॥३७-३९॥
}{%
\textbf{[Hindi Commentary --- Translation]} The terrestrial circumference is five thousand. Multiply this by the latitudinal difference between one's own place and the standard meridian; divide by 360; the square of that result, subtracted from the square of the longitudinal difference [in yojanas]; from the square root of the remainder, multiplying by the planet's motion and dividing by the terrestrial circumference, one obtains a result in minutes. East of the Ujjayinī north-south line [apply as] positive to the mean or true planet; west as negative. The root multiplied by sixty and divided by the terrestrial circumference gives the result in \textit{gha.tikā}s; this is to be applied to the \textit{tithi}-freed \textit{gha.tī}s (i.e., time correction) like a planet.
}

%% ============================================================================
\subsection*{Upapatti (Proof/Explanation)}

\sktcommbilingual{%
यदि व्यासमानेन १५८१ तदा व्यासवर्गाष्टगुणात्पदं भूपरिधिभवेदिति सूर्यसिद्धान्तोक्त्या भूपरिधिमानं ५००० मागच्छति ब्रह्मगुप्तमते।
}{%
\textbf{[Sanskrit Upapatti]} If the diameter measure is 1581, then from the square of the diameter, multiplied by eight, [taking] the square root gives the terrestrial circumference --- so it is stated in the \textit{Sūryasiddhānta}. By this rule the terrestrial circumference measure of 5000 is obtained, which is Brahmagupta's opinion.
}

\hindicommbilingual{%
\textbf{उपपत्ति}---यदि व्यासमान १५८१ मानते हैं तब ``व्यासवर्गाष्टगुणात्पदं भूपरिधिभवेत्‌'' इस सूर्यसिद्धान्तोक्त प्रकार से भूपरिधिमान ५००० प्राप्त होता है, यही ब्रह्मगुप्तमत में भूपरिधिमान है।
}{%
\textbf{[Upapatti --- Translation]} If we assume the diameter measure to be 1581, then by the method stated in the \textit{Sūryasiddhānta}, ``the square root of eight times the square of the diameter is the terrestrial circumference,'' the terrestrial circumference measure of 5000 is obtained. This same [value] is the terrestrial circumference in Brahmagupta's system.
}

\secdivider
\clearpage

%% ============================================================================
\section{Computation of Days Since Epoch (\dev{युगादिदिनानयन})}

\ruleheadingbilingual{नियमः ४०}{Rule 40: Days Elapsed Since Kalpa}

\shlokabilingual{%
कल्पगताब्दा गुणिता रूपाष्टजिनेनापवर्निसप्तनगैः ।
खस्वरसनवभिभक्ता दिनान्यवमानि चांशका दोषाः ॥४०॥
}{%
\textbf{[BSS I.40]} The years elapsed since the beginning of the \textit{kalpa} (\textit{kalpa-gatābdā}), multiplied by 2486 (\textit{rūpā.s.tajina}) and by 7736 (\textit{apavarni-saptanaga}), and divided by 8600 (\textit{kha-svara-sanava}), yield the days and the omitted [tithi]s; the remainders are the fractional parts and defects.
}

\sktcommbilingual{%
\dev{वा. भा.}---अभीष्टे रविमण्डलान्ते ये कल्पगताब्दास्‌ तेऽत्र गृह्यन्ते। कल्पगताब्दाः गुणिताः कृत्याह। रूपाष्टजिनैः २४८६ अन्यत्र नवार्निसप्तनगैः ७७३६ तत उभयतः खखरसनवभिभक्ताः ८६०० फलानि यथासंख्यं दिनान्यवमानि चांशका शेषा द्वयोरपि स्थानयोः। कल्पगताब्दा रूपाष्टजिनैः संगुण्य खखरसनवभिविभजेत्‌ फलं दिनानि भवन्ति। तत्र यद्विकलं तद्दिनांशदाब्देनोच्यते।
}{%
\textbf{[Sanskrit Commentary --- Vāsanābhā.sya]} At the desired end of a solar year, the years elapsed of the \textit{kalpa} are taken here. The years elapsed of the \textit{kalpa} having been multiplied --- he says. By \textit{rūpā.s.tajina} = 2486, and elsewhere by \textit{navārni-saptanaga} = 7736, and then divided by \textit{kha-svara-sanava} = 8600; the results are respectively the days and omitted [lunar days]; the remainders are the fractional parts in both places. The years elapsed of the \textit{kalpa}, multiplied by 2486 and divided by 8600 --- the result is days. There, whatever is the remainder, that is called the \textit{dināṃśa} (fractional part of a day) by the word \textit{dābda}.
}

\hindicommbilingual{%
\textbf{हिन्दी भाषा}---अभीष्ट रविमण्डलान्त में जो कल्पगताब्द हैं वे यहाँ गृह्यन्ते (लिये जाते हैं)। कल्पगताब्दों को २४८६ से गुणा करके ८६०० से भाग देने से फल दिन और अवमान शेष हैं, दोनों स्थानों पर यही है। कल्पगताब्दों को २४८६ से गुणा कर ८६०० से भाग देने से फल दिन होते हैं। तत्र जो शेष विकल रहता है उसको दिनांश कहा जाता है।
}{%
\textbf{[Hindi Commentary --- Translation]} In the desired end of a solar year, the years elapsed of the \textit{kalpa} that are here taken. The years elapsed of the \textit{kalpa}, multiplied by 2486 and divided by 8600 --- the result is days and omitted [days]; the remainders are in both places. The years elapsed of the \textit{kalpa}, multiplied by 2486 and divided by 8600, the result is days. There, whatever remainder remains, that is called the \textit{dināṃśa} (fractional part of a day).
}

\secdivider
\clearpage

%% ============================================================================
\section{Mean Position of the Sun and Moon}

\ruleheadingbilingual{नियमाः २४--२६}{Rules 24--26: Daily Motion and Latitude}

\shlokabilingual{%
मन्दफललिप्ताः स्वहाराप्ता रवीन्द्वोर्भुक्तिफलं क्रमात्‌ ॥
दक्षिणे धनमुत्तरे ऋणं भवति ध्रुवसंज्ञकं तत्‌ ॥
एवमिह दण्डमुखं विपरीतं प्राक्‌ चेद्विघोरतनफलाल्पमथात्र सौरम्‌ ॥
सकृत्‌ सकलसन्नरञ्जनाथं भुजं विचारयन्तु ज्यौतिषिका भूषं विचारयन्त्विति ॥२४॥
}{%
\textbf{[BSS I.24]} The minutes of the \textit{manda-phala} (equation of centre), multiplied by the divisor and divided [appropriately], give in order the daily-motion results for the Sun and Moon. In the southern [hemisphere] it is positive, in the northern negative --- this is called the \textit{dhruva} (constant). Thus here the staff-head [correction] is reversed; if previously the \textit{vighora-tana-phala} was small, then here [it becomes] solar [correction]. Let astronomers examine the \textit{bhuja} (sine) once, [that of] the delight of all good people; let them examine the \textit{bhū.sa} (ornament) [of calculation].
}

\sktcommbilingual{%
\dev{वा. भा.}---प्रह्लादीनां युगे कियन्ति भगणमानानि गताहर्गणमानानि यानि, दिनवारादीं विचाराः, मध्यमग्रहादिसाधनानि यानि, एतदादिषु विषयेषु ग्रार्याणां (आर्यभट्टस्य) षष्टिचा (षष्टिप्रमिताच्छिन्दसा) प्रथमो मध्यगतिरध्यायः मया कृतोऽर्थान्मध्यगतिनामकेऽध्याये कियन्तो विषयाः सन्ति तेषामुल्लेखः कृत इति ॥
}{%
\textbf{[Sanskrit Commentary --- Vāsanābhā.sya]} How many revolutions (\textit{bhaga.namāna}) of Prahlāda and the others in a \textit{yuga}, how many elapsed days (\textit{gata-āharga.na}), the computations of weekdays and so on, the computations of mean planets and so forth --- in these and other subjects, the first chapter on mean motion (\textit{madhyagati-nāmaka adhyāya}) was composed by me (Brahmagupta) with sixty-\textit{āryā} [verses] (\textit{ṣaṣṭi-pramitā chandasā}); how many subjects there are in the chapter named \textit{Madhyagati}, an enumeration of them has been made. Thus [ends the commentary].
}

\secdivider
\clearpage

%% ============================================================================
\section{Mean Position of Mars and Other Planets}

\ruleheadingbilingual{नियमः मध्यमग्रहसाधनम्‌}{Rule: Mean Planets Computation}

\shlokabilingual{%
मन्दकेन्द्रेभुजज्याभिर्मन्दफलं पातजीविका ।
कोटिजीवाहृतं भक्तं व्यासार्धं फलमण्डलम्‌ ॥
}{%
\textbf{[BSS --- General Rule]} The \textit{manda-phala} (equation of centre) is [obtained] from the \textit{bhuja-jyā} (sine of the anomaly) of the \textit{manda-kendra} and the \textit{pāta-jīvikā} [divisor]; divided by the \textit{ko.ti-jyā} (cosine) and multiplied by the \textit{vyāsārdha} (radius), [it becomes] the \textit{phala-ma.n.dala} (resultant circle/epicycle).
}

\secdivider
\clearpage

%% ============================================================================
\section{Epilogue to the Madhyamādhikāra}

\ruleheadingbilingual{अध्यायोपसंहारः}{Chapter Conclusion}

\shlokabilingual{%
इति ब्राह्मस्फुटसिद्धान्ते मध्यमाधिकारः प्रथमः ॥
}{%
\textbf{[End of Chapter 1]} Thus the first chapter, the \textit{Madhyamādhikāra} (Chapter on Mean Motions), in the \textit{Brāhmasphu.tasiddhānta} [of Brahmagupta].
}

\hindicommbilingual{%
\textbf{हिन्दी भाषा}---ग्रहादियों के युग में जो भगणमान है, गताहर्गणमान जो है, दिनवारादि के जो विचार हैं, और मध्यम ग्रहमादि साधन जो है एतदादिक विषयों में तिरसठ आर्यछन्दों के द्वारा मध्यगति नामक प्रथम अध्याय किया गया अर्थात्‌ मध्यगति नामक प्रथम अध्याय में कितने विषय हैं उनका उल्लेख किया गया इति ॥
}{%
\textbf{[Hindi Commentary --- Translation]} The revolutions of the planets and others in a \textit{yuga}, the elapsed \textit{aharga.na}, the computations of weekdays and so forth, and the computations of mean planets and so on --- in these and other subjects, the first chapter named \textit{Madhyagati} was composed by means of sixty-three \textit{āryā} [verses]; that is, in the first chapter named \textit{Madhyagati}, however many subjects there are, an enumeration of them has been made. Thus [ends the commentary on Chapter~1].
}

\vspace{1em}
\bilingual{%
\textit{[मध्यमाधिकार-समाप्तिः। एषः प्रथमः खण्डः पृष्ठसङ्ख्या ८३९--८९३ संपूर्णतया व्याख्यातः।]}
}{%
\textit{[End of the Madhyamādhikāra. This first section, pages 839--893, has been completely explained. It covers mean motions of all planets, day computation, \textit{deśāntara} (longitude correction), and \textit{yugādi} calculations.]}
}

\clearpage

%% ============================================================================
\chapter[Spa.s.tādhikāra]{\dev{स्पष्टाधिकारः} --- Spa.s.tādhikāra}
\label{chap:spashta}

\sectionmark{Spa.s.tādhikāra}

\bilingual{%
\dev{स्पष्टाधिकारः} ब्राह्मस्फुटसिद्धान्तस्य द्वितीयोऽध्यायः। अत्र स्पष्टीकरणस्य प्रयोजनं, ज्यापञ्चविंशतिका, उत्क्रमज्या, मन्दफलं, शीघ्रफलं, स्फुटभुक्तिश्च प्रदर्श्यते। एषः खण्डः पृष्ठ ८९६--९७८।
}{%
The \textit{Spa.s.tādhikāra} (Chapter on True Positions) is the second chapter of the \textit{Brāhmasphu.tasiddhānta}. Here the purpose of \textit{spa.s.tīkara.na} (making true), the twenty-four sine table (\textit{jyā-pañcaviṃśatikā}), versed sines (\textit{utkrama-jyā}), \textit{manda-phala} (equation of centre), \textit{śīghra-phala} (equation of conjunction), and true daily motion (\textit{spa.s.ta-bhukti}) are presented. This section spans pages 896--978.
}

%% ============================================================================
\section{Purpose of Spa.s.tīkara.na}

\ruleheadingbilingual{नियमः १}{Rule 1: The Need for Correction}

\shlokabilingual{%
यस्मान्मध्यतुल्यः प्रतिदिवसं हृश्यते ग्रहो भगणे ।
तस्मात्तत्कृत्यकरं वक्ष्ये मध्यस्फुटीकरणम्‌ ॥१॥
}{%
\textbf{[BSS II.1]} Since the mean planet is not observed in the zodiac (\textit{bhaga.na}) every day equal to [its mean computation], therefore I shall state the procedure for making the mean [planet] true --- the \textit{madhyasphu.tīkara.na}.
}

\sktcommbilingual{%
\dev{वा. सा.}---प्रथमं स्फुटगत्यध्यायं व्याख्यायते। तत्रारम्भप्रयोजनमाह। यस्मान्मध्यग्रहेण तुल्यः अग्रविषये ग्रहो न हृश्यते भगणे नक्षत्रचक्रे प्रतिदिवसं दिवसेदिवसे तस्मात्‌ स्फुटीकरणं वक्ष्ये। मध्यस्य किं हगित्याह---हत्कृतुल्यक्रममप्रायो मध्यमो ग्रहः कक्षामण्डले परिकल्पते। न च कक्षामण्डले पारमार्थिको ग्रहः प्रतिमण्डले मध्यगत्या भ्रमति किन्तु स्पष्टगत्या प्रतिमण्डले परिभ्रमन्‌ कक्षामण्डले हृश्यते, अतो येन गणितेन कक्षामण्डले प्रतिमण्डलस्थो ग्रहो हत्समो भवेत्तादृशं स्फुटीकरणमहं वक्ष्ये इति ॥१॥
}{%
\textbf{[Sanskrit Commentary --- Vāsanābhā.sya]} First, the chapter on true motion (\textit{sphu.ta-gaty-adhyāya}) is being explained. There the purpose of the commencement is stated. Since a planet equal to the mean planet is not observed in practice in the zodiac (\textit{bhaga.na}, the nak.satra circle) day by day, therefore I shall state the \textit{sphu.tīkara.na} (procedure for making true). In answer to ``what is the defect of the mean?'' he says: the mean planet is hypothetically located (\textit{parikalpite}) on the \textit{kak.sā-ma.n.dala} (deferent circle). And the real planet is not on the \textit{kak.sā-ma.n.dala}; revolving on the \textit{prati-ma.n.dala} (epicycle) with mean motion, it is observed on the \textit{kak.sā-ma.n.dala} with true motion. Therefore, by whatever computation the planet situated on the \textit{prati-ma.n.dala} may become equal to [the observed position] on the \textit{kak.sā-ma.n.dala}, that \textit{sphu.tīkara.na} I shall state. Thus [ends commentary on] verse~1.
}

\hindicommbilingual{%
\textbf{हिन्दी भाषा}---जिस कारण से प्रत्येक दिन क्रान्तिवृत्त में स्पष्ट ग्रह मध्यम ग्रह के बराबर नहीं देखे जाते हैं उस कारण से हत्कृतुल्य कारक स्पष्टीकरण को मैं (ब्रह्मगुप्त) कहता हूँ। कक्षावृत्त में मध्यम ग्रह परिकल्पित प्रतिवृत्त में स्पष्टगति से भ्रमण करते हुए ग्रह कक्षावृत्त में देखे जाते हैं। इसलिए जिस गणित से कक्षावृत्त में प्रतिवृत्तस्थित ग्रह हत्सम होते हैं उस स्पष्टीकरण को मैं कहता हूँ ॥१॥
}{%
\textbf{[Hindi Commentary --- Translation]} For the reason that the true planet in the \textit{krānti-v.rtta} (ecliptic) is not seen every day equal to the mean planet, for that reason I (Brahmagupta) call that procedure which makes [it] equal to the observed [position] the \textit{sphu.tīkara.na}. The mean planet is hypothetically [placed] on the \textit{kak.sā-v.rtta} (deferent); the planet, revolving with true motion on the \textit{prati-v.rtta} (epicycle), is observed on the \textit{kak.sā-v.rtta}. Therefore, by whatever computation the planet situated on the \textit{prati-v.rtta} becomes equal [to the observed position] on the \textit{kak.sā-v.rtta}, that \textit{sphu.tīkara.na} I state. Thus [ends commentary on] verse~1.
}

\secdivider
\clearpage

%% ============================================================================
\section{The Twenty-Four Sines (\dev{ज्यापञ्चविंशतिका})}

\ruleheadingbilingual{नियमाः २--९}{Rules 2--9: Sine Table in 24 Steps}

\bilingual{%
\textbf{\dev{ज्यापञ्चविंशतिका-प्रवेशः}}\\[0.3em]
ब्रह्मगुप्तः क्रमज्याः (समज्याः) चतुर्विशतिसङ्ख्यकाः प्रददाति। प्रत्येकं चापं $3\frac{3}{4}$° (= 225') अस्ति, येन नवत्यंशपर्यन्तं सम्पूर्ण्यते। एता ज्याः भूतसङ्ख्याभिः (शब्दसङ्ख्याभिः) कवित्वेन निरूपिताः, यत्‌ गणितज्ञस्य ब्रह्मगुप्तस्य काव्यप्रतिभां दर्शयति।
}{%
\textbf{Introduction to the Jyā-pañcaviṃśatikā}\\[0.3em]
Brahmagupta gives twenty-four \textit{krama-jyā}s (sines in sequence), each arc being $3\tfrac{3}{4}$° (= 225'), filling the complete quadrant of 90°. These sines are expressed poetically through \textit{bhūta-saṅkhyā}s (word-numerals), displaying Brahmagupta's literary talent alongside his mathematical genius.
}

\vspace{0.5em}

\shlokabilingual{%
षड्बुदग्निसुताग्निरसद्वाद्याडा मुनीन्दुवसुचन्द्राः ।
इन्दुनवनवचन्द्रा रसतिथियमला रविश्रियमा: ॥२॥
छिद्रेपुजिनाः कृतनवपञ्चयमा नन्दचन्द्रसुनिपक्षाः ।
दन्ताष्टयमा गुणरामनवयमा: शकद्यमखरामाः ॥३॥
ऋतुनवखयुरणा नवशरचन्द्रयुणाः सप्तशून्ययमदहना: ॥४॥
द्विजनगुरास्त्रिरसरदा: खसप्तयमदहनो व्यस्ताः ॥५॥
मुनयोऽष्टयमास्त्रिरसा रुद्रशशाङ्काः समुद्रमुनिचाः ॥६॥
नववेदयमा सुनियुराहुताशना वसुगुणसमुद्राः ॥७॥
रूपेन्द्रियेषवो रसनगर्तवचन्द्रशीतकरवसवः ॥८॥
शरवसुनन्दाः सागररुद्रशशाङ्का नवाडार्कः ॥९॥
}{%
\textbf{[BSS II.2--9] The 24 Sines}\\[0.3em]
The twenty-four sines are encoded in \textit{bhūta-saṅkhyā} (word-numeral) verses. Each word corresponds to a numerical value:\\
\textbf{Verse 2:} ṣaḍ-bud-agni (6+14=20) = 214; suta-agni (7+3=10) = 427; rasa-dvādya (6+12=18) = 638; munīndu (7+1=8) = 846; vasu-candra (8+1=9) = 1051; indu-nava-nava-candra (1+9+9) = 1251; rasa-tithi-yamala (6+15+1) = 1446; ravi-śriyamā (12+1) = 1624.\\[0.2em]
\textbf{Verses 3--9} continue the sequence up to the 24th sine = 3270 (the full radius).
}

\vspace{0.5em}

\bilingual{%
\textbf{\dev{क्रमज्या-सारणी}}\\[0.2em]
\begin{tabular}{|c|c|c|}
\hline
क्र. & मानम्‌ & नाम \\
\hline
१ & २१४ & प्रथमज्या \\
२ & ४२७ & द्वितीयज्या \\
३ & ६३८ & तृतीयज्या \\
४ & ८४६ & चतुर्थीज्या \\
५ & १०५१ & पञ्चमीज्या \\
६ & १२५१ & षष्ठीज्या \\
७ & १४४६ & सप्तमीज्या \\
८ & १६२४ & अष्टमीज्या \\
९ & १८१७ & नवमीज्या \\
१० & १९९१ & दशमीज्या \\
\hline
\end{tabular}
}{%
\textbf{Krama-jyā Table (Sines in Sequence)}\\[0.2em]
\begin{tabular}{|c|c|l|}
\hline
No. & Value & Name \\
\hline
1 & 214 & Prathamā-jyā \\
2 & 427 & Dvitīyā-jyā \\
3 & 638 & T.rtīyā-jyā \\
4 & 846 & Caturthī-jyā \\
5 & 1051 & Pañcamī-jyā \\
6 & 1251 & .Sa.s.thī-jyā \\
7 & 1446 & Saptamī-jyā \\
8 & 1624 & A.s.tamī-jyā \\
9 & 1817 & Navamī-jyā \\
10 & 1991 & Daśamī-jyā \\
\hline
\end{tabular}
}

\bilingual{%
\begin{tabular}{|c|c|c|}
\hline
क्र. & मानम्‌ & नाम \\
\hline
११ & २१४५ & एकादशीज्या \\
१२ & २२९६ & द्वादशीज्या \\
१३ & २४१९ & त्रयोदशीज्या \\
१४ & २५९४ & चतुर्दशीज्या \\
१५ & २७१९ & पञ्चदशीज्या \\
१६ & २८३२ & षोडशीज्या \\
१७ & २९३३ & सप्तदशीज्या \\
१८ & ३०२१ & अष्टादशीज्या \\
१९ & ३०९६ & एकूनविंशज्या \\
२० & ३१५९ & विंशज्या \\
\hline
\end{tabular}
}{%
\begin{tabular}{|c|c|l|}
\hline
No. & Value & Name \\
\hline
11 & 2145 & Ekādaśī-jyā \\
12 & 2296 & Dvādaśī-jyā \\
13 & 2419 & Trayodaśī-jyā \\
14 & 2594 & Caturdaśī-jyā \\
15 & 2719 & Pañcadaśī-jyā \\
16 & 2832 & .So.daśī-jyā \\
17 & 2933 & Saptadaśī-jyā \\
18 & 3021 & A.s.tādaśī-jyā \\
19 & 3096 & Ekūnaviṃśa-jyā \\
20 & Viṃśa-jyā \\
\hline
\end{tabular}
}

\bilingual{%
\begin{tabular}{|c|c|c|}
\hline
क्र. & मानम्‌ & नाम \\
\hline
२१ & ३२०७ & एकविंशज्या \\
२२ & ३२४२ & द्वाविंशज्या \\
२३ & ३२६३ & त्रयोविंशज्या \\
२४ & ३२७० & चतुर्विंशज्या \\
\hline
\end{tabular}
}{%
\begin{tabular}{|c|c|l|}
\hline
No. & Value & Name \\
\hline
21 & 3207 & Ekaviṃśa-jyā \\
22 & 3242 & Dvāviṃśa-jyā \\
23 & 3263 & Trayoviṃśa-jyā \\
24 & 3270 & Caturviṃśa-jyā \\
\hline
\end{tabular}
}

\vspace{0.5em}

\sktcommbilingual{%
\dev{वा. था.}---एषामद्वयास एव व्यास्यानम्‌। अघेज्यकास्ता इति तद्यथा। सनुयमला २१४ एतत्‌ प्रथमज्यारूपम्‌। सुनियमवेदा ४२७ द्वितीयम्‌। वसुज्वलनघटिकाः ६३८ तृतीयम्‌। रसकृतवसवः ८४६ चतुर्थम्‌। शशिपञ्चखेन्दवः १०५१ पञ्चमम्‌। चन्द्रदरसूर्याः १२५१ षष्ठम्‌।
}{%
\textbf{[Sanskrit Commentary]} Their diameter-measure is the radius. ``These are the \textit{aghe-jyā}'' (sines) --- thus. \textit{Sanu-yamalā} = 214; this is the form of the first sine. \textit{Suniya-mavedā} = 427, the second. \textit{Vasu-jvalana-gha.tikāḥ} = 638, the third. \textit{Rasa-k.rta-vasavaḥ} = 846, the fourth. \textit{Śaśi-pañcha-khendavaḥ} = 1051, the fifth. \textit{Candra-dara-sūryāḥ} = 1251, the sixth.
}

\secdivider
\clearpage

%% ============================================================================
\subsection*{Utkrama-jyā (Versed Sine) Values}

\bilingual{%
\textbf{\dev{उत्क्रमज्या-सारणी}}\\[0.2em]
\dev{उत्क्रमज्या} (versed sines) अर्धज्यानां व्यस्तक्रमेण निर्मिताः। एता रेखागणिते ``१ $-$ कोषज्या'' रूपेण ज्ञायन्ते।
}{%
\textbf{Utkrama-jyā Table (Versed Sines)}\\[0.2em]
The \textit{utkrama-jyā}s (versed sines) are constructed by the reversed order of the half-chords. In modern trigonometry these are known as ``$1 - \cosine$'' (versed sine = radius $-$ cosine).
}

\vspace{0.5em}

\bilingual{%
\begin{tabular}{|c|c|l|}
\hline
क्र. & मानम्‌ & भूतसङ्ख्या \\
\hline
१ & ७ & समुद्रमुनिच (४+०+७) \\
२ & २९ & नववेदयमा (९+४+८) \\
३ & ६६ & सुनियुराहुताशन (१४+३+५) \\
४ & १०२ & वसुगुणसमुद्र (८+३+४) \\
५ & १७१ & रूपेन्द्रियेषव (१+११+५) \\
६ & २२४ & रसनगर्तव (६+६+९) \\
७ & २९२ & चन्द्रशीतकरवसव (१+२+५+८) \\
८ & ३७० & शरवसुनन्दा (५+८+९) \\
\hline
\end{tabular}
}{%
\begin{tabular}{|c|c|l|}
\hline
No. & Value & Bhūta-saṅkhyā \\
\hline
1 & 7 & samudra-muni-ca (4+0+7) \\
2 & 29 & nava-veda-yamā (9+4+8) \\
3 & 66 & suni-yur-āhutāśana (14+3+5) \\
4 & 102 & vasu-guṇa-samudra (8+3+4) \\
5 & 171 & rūpendriye.ṣava (1+11+5) \\
6 & 224 & rasa-nagartava (6+6+9) \\
7 & 292 & candra-śītakara-vasava (1+2+5+8) \\
8 & 370 & śara-vasu-nandā (5+8+9) \\
\hline
\end{tabular}
}

\bilingual{%
\begin{tabular}{|c|c|l|}
\hline
क्र. & मानम्‌ & भूतसङ्ख्या \\
\hline
९ & ४८१ & सागररुद्रशशाङ्क (४+११+१) \\
१० & ६२१ & नवाडार्क (९+८+१२) \\
११ & ७६६ & त्रिविषयवेदशशाङ्क (३+६+४+१) \\
१२ & ९३१ & पञ्चत्रिरसेन्दव (५+३+६) \\
१३ & १११४ & षड्बुधियमधृतय (६+१७+१८) \\
१४ & १२९६ & श्रतिधृतिखयमा (१५+२+८) \\
१५ & १४५३ & नवशशियमपक्ष (९+६+१०) \\
\hline
\end{tabular}
}{%
\begin{tabular}{|c|c|l|}
\hline
No. & Value & Bhūta-saṅkhyā \\
\hline
9 & 481 & sāgara-rudra-śaśāṅka (4+11+1) \\
10 & 621 & navāḍārka (9+8+12) \\
11 & 766 & trivi.saya-veda-śaśāṅka (3+6+4+1) \\
12 & 931 & pañca-tri-rasendava (5+3+6) \\
13 & 1114 & .sa.d-budhi-yamadh.rtaya (6+17+18) \\
14 & 1296 & śrati-dh.rti-khayamā (15+2+8) \\
15 & 1453 & nava-śaśi-yama-pak.sa (9+6+10) \\
\hline
\end{tabular}
}

\bilingual{%
\begin{tabular}{|c|c|l|}
\hline
क्र. & मानम्‌ & भूतसङ्ख्या \\
\hline
१६ & १६३५ & सागरद्विजिन (४+२+६) \\
१७ & १८२४ & रदरसयमला (१२+६+१५) \\
१८ & २०१९ & गुणवेदवसुयमा (३+४+८+८) \\
१९ & २२१४ & षट्विषयशून्यगुण (६+६+०+३) \\
२० & २४२४ & खमुनिरदा (०+७+१२) \\
२१ & २६३२ & व्यासाघं (०+८) \\
२२ & २८४३ & नवरदचन्द्र (९+१२+१) \\
२३ & ३०५६ & जिनांदा (५+८) \\
२४ & ३२७० & परमक्रान्तिज्या \\
\hline
\end{tabular}
}{%
\begin{tabular}{|c|c|l|}
\hline
No. & Value & Bhūta-saṅkhyā \\
\hline
16 & 1635 & sāgara-dvijina (4+2+6) \\
17 & 1824 & radara-sa-yamalā (12+6+15) \\
18 & 2019 & gu.na-veda-vasu-yamā (3+4+8+8) \\
19 & 2214 & .sa.t-vi.saya-śūnya-gu.na (6+6+0+3) \\
20 & 2424 & kha-muniradā (0+7+12) \\
21 & 2632 & vyāsāgha (0+8) \\
22 & 2843 & nava-rada-candra (9+12+1) \\
23 & 3056 & jināndā (5+8) \\
24 & 3270 & Parama-krānti-jyā (full radius) \\
\hline
\end{tabular}
}

\vspace{0.5em}

\hindicommbilingual{%
\textbf{हिन्दी भाषा}---बत्तचतुर्थांश मुनियमला मुनियमवेदा इत्याद्यर्धज्याइचतुरविशतिसंख्यका सन्ति यथा अघोलिखिताः स्युः। व्यस्ता (उत्क्रमज्या) इत्यस्याग्रे सम्बन्धः। चतुरविशतिरर्धज्या (क्रमज्याः)---२१४, ४२७, ६३८, ८४६, १०५१, १२५१, १४४६, १६२४, १८१७, १९९१, २१४५, २३१२, २४५९, २५९४, २८३२, २९३३, ३०२१, ३०९६, ३१४५९, ३२०७, इत्यादि, ३२६३, ३२७०।
}{%
\textbf{[Hindi Commentary --- Translation]} The twenty-four half-chords (\textit{ardhajyā}), [named] \textit{muniyamalā}, \textit{muniyamavedā}, etc., are of twenty-four numbers, as written below. The relation of \textit{vyastā} (reversed, i.e., versed sine) is given after. The twenty-four \textit{ardhajyā}s (\textit{krama-jyā}s) are: 214, 427, 638, 846, 1051, 1251, 1446, 1624, 1817, 1991, 2145, 2312, 2459, 2594, 2832, 2933, 3021, 3096, 3145(?), 3207, etc., 3263, 3270.
}

\secdivider
\clearpage

%% ============================================================================
\section{Arc--Jyā Conversion (\dev{चापज्यानयन})}

\ruleheadingbilingual{नियमः चापनिर्माणम्‌}{Rule: Finding Arc from Jyā}

\shlokabilingual{%
ज्यार्धात्त्रिज्याहृतात्‌ पदं स्वल्पान्तराद्विषमात्‌ समात्‌ ।
चापं स्यात्ततः परं तु पदं तद्वर्गयुतेः पुनः ॥
}{%
\textbf{[BSS --- Arc--Jyā Rule]} From the \textit{jyārdha} (half-chord, sine), divided by the \textit{trijyā} (radius), the square root gives [the first approximation], slightly different for odd and even [intervals]; from that the arc is obtained. After that, again [taking] the square root of [the sine] combined with its square [correction], [the more precise arc is obtained].
}

\sktcommbilingual{%
\dev{वि. सा.}---बत्तचतुर्थाशि मुनयोऽष्टयमा इत्यादि चतुरविशतिसंख्यका उत्क्रमज्याः सन्ति यथा अघोलिखिताः स्युḥ---७, रद् (२९), ६२, १११; १७४, रउ (२९२), ३३७, ४३५, रभ१ (५२१), ६७६, ८११, रउ (९८५), १११४, १२९६, १४५३, १६३५, १८२४, २०१९, २२१४, २४२४, २६३२, २५४३, ३०५६, रेर७० (३२७०)।
}{%
\textbf{[Sanskrit Commentary]} The twenty-four versed sines (\textit{utkrama-jyā}), [named] \textit{muno'ṣ.tayamā} etc., are of twenty-four numbers, as written below: 7, rad (29), 62, 111; 174, rau (292), 337, 435, rabha1 (521), 676, 811, rau (985), 1114, 1296, 1453, 1635, 1824, 2019, 2214, 2424, 2632, 2543(?), 3056, rera70 (3270).
}

\hindicommbilingual{%
\textbf{हिन्दी भाषा}---ज्यार्ध को त्रिज्या से भाग देने से जो पद (मूल) प्राप्त हो वह स्वल्पान्तराद्विषमात्‌ समात्‌ (थोड़े अन्तर से विषम और सम) से चाप होता है। उसके पश्चात्‌ पुनः उसी पद के वर्ग में युक्त करके (कुछ संस्कार करके) आगे का चाप प्राप्त होता है।
}{%
\textbf{[Hindi Commentary --- Translation]} The \textit{jyārdha} (half-chord), divided by the \textit{trijyā} (radius), from whatever root (\textit{pada}) is obtained, slightly differing for odd and even [intervals], that becomes the arc. After that, again, having combined [the correction] with the square of that same root, the next arc is obtained.
}

\secdivider
\clearpage

%% ============================================================================
\section{Manda and Śīghra Equations}

\ruleheadingbilingual{नियमाः ४१--४४}{Rules 41--44: True Motion Computation}

\bilingual{%
\textbf{\dev{मन्दफल-शीघ्रफल-प्रवेशः}}\\[0.3em]
\dev{मन्दफल} (manda-phala, equation of centre) तथा \dev{शीघ्रफल} (śīghra-phala, equation of conjunction) इत्येते ग्रहस्फुटीकरणस्य मुख्ये संस्कारे। मन्दफलं मन्दकेन्द्रज्यया, शीघ्रफलं शीघ्रकेन्द्रज्यया निर्भर्त्स्यते।
}{%
\textbf{Introduction to Manda-phala and Śīghra-phala}\\[0.3em]
The \textit{manda-phala} (equation of centre) and \textit{śīghra-phala} (equation of conjunction) are the two principal corrections for making the planet true. The \textit{manda-phala} is computed from the \textit{manda-kendra} (anomaly of the epicycle); the \textit{śīghra-phala} from the \textit{śīghra-kendra} (anomaly of conjunction).
}

\vspace{0.5em}

\shlokabilingual{%
मन्दकेन्द्रेभुजज्याभिः पातजीविकया गुणाः ।
कोटिजीवाहृता भक्ता व्यासार्धेन फलं गतिः ॥४१॥
ततः प्रथममल्पं द्वितीयं वा मन्दसंस्कृतात्‌ ।
मध्यान्मन्दफलादद्याच्छीघ्रभुक्त्याऽथवा पुनः ॥४२॥
मन्दसंस्कृतमादाय शीघ्रभुक्त्याऽथवा पुनः ।
शीघ्रफलं ततः कुर्यात्त्रैराशिकेन संहितम्‌ ॥४३॥
शीघ्रफलं व्यवस्थाप्य ग्रहस्य स्फुष्टगतिकाः ।
समानीय ततः कुर्याद्वक्रभुक्तिं विचक्षणः ॥४४॥
}{%
\textbf{[BSS II.41--44]}\\[0.3em]
\textbf{41:} The \textit{bhuja-jyā} (sine of anomaly) of the \textit{manda-kendra}, multiplied by the \textit{pāta-jīvikā} and divided by the \textit{ko.ti-jyā} (cosine), [then] multiplied by the \textit{vyāsārdha} (radius), gives the \textit{manda-phala} (equation of centre) and the motion.\\[0.3em]
\textbf{42:} From that, first the small or second [epicycle], or from the \textit{manda-saṃsk.rta} (mean corrected by manda), from the mean, one should give the \textit{manda-phala}, or again by the \textit{śīghra-bhukti}.\\[0.3em]
\textbf{43:} Taking the \textit{manda-saṃsk.rta}, or again by the \textit{śīghra-bhukti}, one should then compute the \textit{śīghra-phala}, combined by proportion (\textit{trairāśika}).\\[0.3em]
\textbf{44:} Having established the \textit{śīghra-phala}, the true velocities (\textit{sphu.s.t.a-gatikaḥ}) of the planet are to be brought together; then the wise [astronomer] should compute the \textit{vakra-bhukti} (retrograde motion).
}

\vspace{0.5em}

\sktcommbilingual{%
\dev{वि. सा.}---मन्दकेन्द्रगतिज्यान्तरेण (भोग्यखण्डेन) गुणिता, प्राद्वजीवया (प्रथमज्यया) भक्ता यल्लब्धं तत्स्फुटपरिधिगुणितं, भगणांशे ३६५ भक्ते लब्धामिष्कलामिमां शकवर्गादौ (सकलादिकेन्द्रे, क्कयादिकेन्द्रे च) स्वमध्यमगतीरूनाधिका तदानीकेन्द्रस्य (रविचन्द्रयोः) स्फुटा गतिर्भवति, कुजादीनां ग्रहाणां मन्दगतिफलरहिता मध्यगति (मन्दस्फुटगति) वदत्याचार्याः।
}{%
\textbf{[Sanskrit Commentary]} Multiplied by the \textit{manda-kendra-gati-jyāntara} (\textit{bhogyakha.n.da}, the tabular difference), divided by the \textit{prathama-jīvā} (first sine); whatever result is obtained, multiplied by the \textit{sphu.ta-paridhi} and divided by 365 (\textit{bhaga.nāṃśa}), the resulting minutes --- in the \textit{sakala-ādi-kendra} etc., [when this is] less or more than its own mean motion, at that \textit{kendra} the true motion of the Sun and Moon is obtained. For Mars and the other planets, the mean motion minus the \textit{manda-gati-phala} is called \textit{manda-sphu.ta-gati} by the teachers.
}

\hindicommbilingual{%
\textbf{हिन्दी भाषा}---मन्दकेन्द्रगतिज्यान्तर (भोग्यखण्ड) से गुणित प्रथमजीवा (प्रथमज्या) से भक्त जो लब्ध फल वह स्फुटपरिधि से गुणित ३६५ भगणांश से भाग देने से जो कला-मिमा प्राप्त होती है वह सकलादिकेन्द्र में अपनी मध्यमगति से ऊना या अधिक होती है, तदानीक केन्द्र में रवि और चन्द्र की स्फुटा गति होती है। कुज आदि ग्रहों की मन्दगतिफलरहित मध्यगति मन्दस्फुटगति है। स्वमन्दस्पुष्टगतिरहिता शीघ्रोच्चगति कुजादि ग्रहों की शीघ्रकेन्द्रगति है।
}{%
\textbf{[Hindi Commentary --- Translation]} The \textit{manda-kendra-gati-jyāntara} (\textit{bhoga-kha.n.da}), multiplied [by the epicycle] and divided by the \textit{prathama-jīvā} (\textit{prathama-jyā}); the resulting quotient, multiplied by the \textit{sphu.ta-paridhi} and divided by 365 \textit{bhaga.nāṃśa}, the resulting minutes --- in the \textit{sakala-ādi-kendra}, whether less or more than the mean motion, at that \textit{kendra} the true motion of the Sun and Moon is obtained. For Mars and other planets, the mean motion minus the \textit{manda-gati-phala} is the \textit{manda-sphu.ta-gati}. The \textit{śīghrocca-gati} minus its own \textit{manda-sphu.ta-gati} is the \textit{śīghra-kendra-gati} of Mars and other planets.
}

\secdivider
\clearpage

%% ============================================================================
\section{The Manda Paridhi (Epicycle) Values}

\ruleheadingbilingual{नियमाः २०--२१}{Rules 20--21: Manda Epicycle Corrections}

\shlokabilingual{%
तद्द्युदलपरिध्यन्तरगुणा हता त्रिज्यया च नतजीवा ।
ऊने धनमृणमधिके दिनार्धपरिधी स्फुटः परिधिः ॥२०॥
सूर्यस्य तदृणमन्दफले धने मन्दपरिधिर्मध्याह्ने ।
चन्द्रस्य तदेव ताभिरेव घटिकामितेन्दुरपि तस्मिन्नेव कपाले नतो भवतीति ॥२१॥
}{%
\textbf{[BSS II.20--21]}\\[0.3em]
\textbf{20:} The difference between that [quantity] and the half-diameter of the epicycle (\textit{dyudala-paridhi}), multiplied by the \textit{trijyā} (radius) and by the \textit{nata-jīvā} (sine of altitude), [when the planet is] in decrease (\textit{ūne}) is positive (\textit{dhana}), in increase (\textit{adhike}) is negative (\textit{.rna}); this is the \textit{sphu.ta-paridhi} (true epicycle) of the half-day.\\[0.3em]
\textbf{21:} For the Sun, that is the \textit{manda-paridhi} in the \textit{.rna-manda-phala} or \textit{dhana-manda-phala} at midday. For the Moon, the same, and by those same [quantities] in \textit{gha.tikā}s, the Moon also [is corrected] at that same \textit{kapāla} (division); thus the \textit{nata} (altitude) is obtained.
}

\sktcommbilingual{%
\dev{वि. सा.}---चतुदंशांशाः स्थानद्वये भागद्वयंशेन रहितास्तदा ऋणे वा धने वा मन्दफले सूर्यस्य; दिनमध्ये (मध्याह्ने) मन्दपरिध्यंशा भवन्ति, ऋणे धने वा मन्दफले दिनार्धात्‌ प्राक्‌कपाले पञ्चदशघटीभिनेतस्य सूर्यस्य दिनमध्यपरिधिमानं क्रमेण भागद्वयंशेनाधिकसूनं कार्यं। चन्द्रस्य ऋणे धने वा मन्दफले दशनद्वितयं जिनलिप्तोनं कार्यमर्थात्‌ स्थानद्वये द्वात्रिंशदंशाख्यचतुरविशतिकलारहितास्तदा मध्याह्ने तन्मन्दपरिध्यंशा भवन्ति।
}{%
\textbf{[Sanskrit Commentary]} Fourteen degrees (\textit{catudaṃśāṃśāḥ}), in two places, diminished by two-thirds of a degree (\textit{bhāgadwayaṃśena rahitāḥ}) --- then in the \textit{.rna} or \textit{dhana manda-phala} of the Sun; at midday (\textit{dinamadhye}) the \textit{manda-paridhi-aṃśa} are obtained. In the \textit{.rna} or \textit{dhana manda-phala}, before the half-day \textit{kapāla}, for the Sun separated by fifteen \textit{gha.tikā}s, the half-day \textit{paridhi} measure is to be made successively increased or decreased by two-thirds of a degree. For the Moon, in the \textit{.rna} or \textit{dhana manda-phala}, twelve [degrees] diminished by two \textit{jina-lipata} (?) --- that is, in two places, diminished by thirty-two degrees plus twenty-four minutes (\textit{dvātriṃśadaṃśākhyacaturviṃśatikalā rahitāḥ}) --- then at midday those \textit{manda-paridhi-aṃśa} are obtained.
}

\hindicommbilingual{%
\textbf{हिन्दी भाषा}---चौदह अंश में दो स्थानों में एक अंश के तृतीयांश (बीस कला) को घटाने से सूर्य के ऋणमन्द फल में वा धनमन्द फल में मध्यन्ह काल में मन्द परिध्यंश होता है। ऋणमन्दफल में वा धनमन्दफल में दिनार्ध से पूर्वकपाल में पञ्चदश १४ घटी करके नत सूर्य के मध्यान्ह मन्दपरिधिमान में क्रम से एक भाग के तृतीयांश (२० कला) को युत और हीन करना चाहिये।
}{%
\textbf{[Hindi Commentary --- Translation]} From fourteen degrees, in two places, diminishing by one-third of a degree (20 \textit{kalā}s), in the \textit{.rna-manda-phala} or \textit{dhana-manda-phala} of the Sun, at midday the \textit{manda-paridhi-aṃśa} is obtained. In the \textit{.rna-manda-phala} or \textit{dhana-manda-phala}, in the \textit{kapāla} before half-day, having subtracted fifteen (14?) \textit{gha.tikā}s from the \textit{nata} Sun, in the midday \textit{manda-paridhi} measure, successively one-third of a degree (20 \textit{kalā}s) is to be added and subtracted.
}

\secdivider
\clearpage

%% ============================================================================
\subsection*{Summary of Manda Paridhi Values}

\bilingual{%
\textbf{\dev{मन्दपरिधि-सारणी}}\\[0.3em]
अधो मन्दपरिधीनां सङ्ख्यात्मकं सारांशः प्रदर्श्यते। प्रत्येकस्य ग्रहस्य मन्दपरिधिः चतुर्विशत्यंशपर्यन्तं परिवर्तते।
}{%
\textbf{Manda-paridhi Summary Table}\\[0.3em]
Below is presented a numerical summary of the \textit{manda-paridhi} (epicycle) values. For each planet the \textit{manda-paridhi} varies up to twenty-four degrees.
}

\vspace{0.5em}

\bilingual{%
\begin{tabular}{|l|c|c|}
\hline
\textbf{ग्रहः} & \textbf{मन्दपरिधिः} & \textbf{विशेषटिप्पनी} \\
\hline
रविः (Sun) & १३--१४ & मन्दानुसारे परिवर्तनम्‌ \\
चन्द्रः (Moon) & ३१--३२ & रवेः द्विगुणे गतिके \\
भौमः (Mars) & ७३--८० & उभयसंस्कारयुक्तः \\
बुधः (Mercury) & ३० & सङ्कीर्णः \\
गुरुः (Jupiter) & ३२--३६ & उभयसंस्कारयुक्तः \\
शुक्रः (Venus) & ११--१४ & सङ्कीर्णः \\
शनिः (Saturn) & ५९--६५ & मन्दगतिः \\
\hline
\end{tabular}
}{%
\begin{tabular}{|l|c|l|}
\hline
\textbf{Planet} & \textbf{Manda Paridhi} & \textbf{Special Notes} \\
\hline
Sūrya (Sun) & 13--14 & Varies by anomaly \\
Candra (Moon) & 31--32 & Twice Sun's rate \\
Bhauma (Mars) & 73--80 & Both corrections \\
Budha (Mercury) & 30 & Complex \\
Guru (Jupiter) & 32--36 & Both corrections \\
Śukra (Venus) & 11--14 & Complex \\
Śani (Saturn) & 59--65 & Slow motion \\
\hline
\end{tabular}
}

\secdivider
\clearpage

%% ============================================================================
\section{True Daily Motion (\dev{स्फुटभुक्ति})}

\ruleheadingbilingual{नियमाः स्फुटभुक्तिः}{Rules: Daily Motion After Correction}

\shlokabilingual{%
शीघ्रकेन्द्रस्फुटभुक्त्यन्तरगुणा श्रत उक्तलब्धात्संशोध्या ।
शीघ्रगतिर्वक्रगतिरिति सिङ्हवर्म्मोपपन्नश्लोकः ॥
}{%
\textbf{[BSS --- True Daily Motion]} The difference between the \textit{śīghra-kendra-sphu.ta-bhukti} and [the mean motion], multiplied by the [tabular difference] and corrected from the stated quotient --- this is the \textit{śīghra-gati}; [if negative] it is \textit{vakra-gati} (retrograde motion). Thus [goes] the proof-verse of Siṅhavarman.
}

\hindicommbilingual{%
\textbf{हिन्दी भाषा}---शीघ्रकेन्द्रस्फुटभुक्त्यन्तर (भोग्यखण्ड) से गुणित प्रथमजीवा (प्रथमज्या) से भक्त जो लब्ध फल वह स्फुटपरिधि से गुणित ३६५ भगणांश से भाग देने से जो कला-मिमा प्राप्त होती है वह सकलादिकेन्द्र में अपनी मध्यमगति से ऊना या अधिक होती है, तदानीक केन्द्र में रवि और चन्द्र की स्फुटा गति होती है।
}{%
\textbf{[Hindi Commentary --- Translation]} The \textit{śīghra-kendra-sphu.ta-bhukti-antara} (\textit{bhoga-kha.n.da}), multiplied [by the epicycle] and divided by the \textit{prathama-jīvā}; the resulting quotient, multiplied by the \textit{sphu.ta-paridhi} and divided by 365 \textit{bhaga.nāṃśa}, the resulting minutes --- in the \textit{sakala-ādi-kendra}, whether less or more than the mean motion, at that \textit{kendra} the true motion of the Sun and Moon is obtained.
}

\subsection*{Geometric Explanation (Upapatti)}

\hindicommbilingual{%
\textbf{उपपत्ति}---मध्यस्पष्ट भेदेन गतिस्त्रिविधा भवति, या गतिः प्रतिक्षणं भिन्ना-भिन्ना भवति, सा स्पष्टा अन्यां मध्या, स्फुटा गतिरपि दैनिकतात्कालिकभेदेन द्विविधा भवति, तेन दैनिकमन्दस्पष्टागतिः, तात्कालिकमन्दस्पष्टागतिः। दैनिकस्पष्टागतिः, तात्कालिकस्पष्टागतिः, आचार्येण दैनिकमन्दस्पष्टगतिः, दैनिकस्पष्टगतिश्चानीयते।
}{%
\textbf{[Upapatti --- Geometric Proof]} By the difference between mean and true, motion becomes threefold. That motion which is different every moment is \textit{spa.s.tā}; another is \textit{madhyā} (mean). The true motion (\textit{sphu.tā gati}) is also twofold by the division into daily and instantaneous: thus the daily \textit{manda-sphu.tā-gati} and the instantaneous \textit{manda-sphu.tā-gati}. The daily \textit{spa.s.tā-gati}, the instantaneous \textit{spa.s.tā-gati} --- by the teacher the daily \textit{manda-sphu.ta-gati} and the daily \textit{sphu.ta-gati} are computed.
}

\hindicommbilingual{%
ग्रहमन्दकेन्द्रगतिज्यान्तर (भोग्यखण्ड) से गुणित प्रथमज्यया भक्त यल्लब्धं तत्स्फुटपरिधिगुणितं भगणांशे ३६५ भक्ते लब्धामिष्कलामिमां शकवर्गादी (सकलादिकेन्द्रे, क्कयादिकेन्द्रे च) स्वमध्यमगतीरूनाधिका तदानीकेन्द्रस्य (रविचन्द्रयोः) स्फुटा गतिर्भवति।
}{%
The \textit{graha-manda-kendra-gati-jyāntara} (\textit{bhoga-kha.n.da}), multiplied by the \textit{prathama-jyā} and divided [by the radius]; whatever quotient is obtained, multiplied by the \textit{sphu.ta-paridhi} and divided by 365 (\textit{bhaga.nāṃśa}), the resulting minutes --- in the \textit{sakala-ādi-kendra} (or \textit{kkayādi-kendra}), whether less or more than its own mean motion, at that \textit{kendra} the true motion of the Sun and Moon is obtained.
}

\secdivider
\clearpage

%% ============================================================================
\section{Special Corrections for Different Planets}

\ruleheadingbilingual{विशेषसंस्काराः}{Planet-Specific Corrections}

\subsection*{Mercury (\dev{बुध}) and Venus (\dev{शुक्र})}

\sktcommbilingual{%
\dev{वि. सा.}---बुधशुक्रयोः शीघ्रोच्चं रवेरर्कात्‌ ग्रहणादितोऽन्यथा। तयोः शीघ्रोच्चं सूर्यसम्बन्धि, मन्दोच्चं स्वतन्त्रमिति विशेषः।
}{%
\textbf{[Sanskrit Commentary]} For Mercury and Venus, the \textit{śīghrocca} (apsis of the śīghra epicycle) is [derived] from the Sun; from the seizure of the Sun [it is different], otherwise [it would be the same]. Their \textit{śīghrocca} is connected to the Sun; the \textit{mandocca} (apsis of the manda epicycle) is independent --- this is the special feature.
}

\hindicommbilingual{%
बुध और शुक्र का शीघ्रोच्च सूर्य से ही संबन्ध रखता है, इनका शीघ्रोच्च सूर्यसम्बन्धी है, मन्दोच्च स्वतन्त्र है। इसलिए इन दोनों ग्रहों में विशेष संस्कार की आवश्यकता होती है।
}{%
\textbf{[Hindi Commentary --- Translation]} The \textit{śīghrocca} of Mercury and Venus is related to the Sun itself; their \textit{śīghrocca} is Sun-related, the \textit{mandocca} is independent. Therefore for these two planets a special correction procedure is necessary.
}

\vspace{1em}

\subsection*{Mars (\dev{भौम}), Jupiter (\dev{गुरु}), and Saturn (\dev{शनि})}

\sktcommbilingual{%
\dev{वि. सा.}---भौमादीनां मन्दस्पष्टगतिः स्फुटगतिश्च। मन्दस्पष्टभुक्त्यन्तरात्‌ शीघ्रभुक्तिः शीघ्रफलाधम्भोग्यजीवासंगुणिता मध्यजीवया विभजेत्‌। लब्धेनोक्तवत्स्फुटभुक्तिः समानीया यदि तया सह मन्दफलाधम्फुटभुक्त्यन्तरा तत्रैव केन्द्रमे कृता भुक्ती धन ऋणां वा कार्यं।
}{%
\textbf{[Sanskrit Commentary]} For Mars and the others, the \textit{manda-sphu.ta-gati} and the \textit{sphu.ta-gati}. From the \textit{manda-sphu.ta-bhukti-antara}, the \textit{śīghra-bhukti} [is obtained]; the \textit{śīghra-phala-ādha-bhogya-jīvā}, multiplied by the \textit{madhya-jivyā} (mean sine), is to be divided. Having brought together the \textit{sphu.ta-bhukti} by the stated quotient; if together with that, by the \textit{manda-phala-ādham-sphu.ta-bhukti-antara} in that same \textit{kendra}, the resulting \textit{bhukti} is to be made positive or negative.
}

\hindicommbilingual{%
भौम आदि ग्रहों की मन्दस्पष्टगति और स्फुटगति---मन्दस्पष्टभुक्त्यन्तर से शीघ्रभुक्ति शीघ्रफलाधम्भोग्यजीवा संगुणित मध्यजीवा विभजेत्‌। लब्धेन उक्तवत्‌ स्फुटभुक्ति समानीया यदि तया सह मन्दफलाधम्फुटभुक्त्यन्तरा तत्रैव केन्द्र में कृता भुक्ती धन ऋण वा कार्य।
}{%
\textbf{[Hindi Commentary --- Translation]} For Mars and other planets, the \textit{manda-sphu.ta-gati} and \textit{sphu.ta-gati} --- from the \textit{manda-sphu.ta-bhukti-antara}, the \textit{śīghra-bhukti}; the \textit{śīghra-phala-adha-bhogya-jīvā}, multiplied and divided by the \textit{madhya-jīva}. By the resulting [quotient], the \textit{sphu.ta-bhukti} is to be brought together as stated; if together with that, by the \textit{manda-phala-adha-sphu.ta-bhukti-antara} in that same \textit{kendra}, the resulting \textit{bhukti} is to be made positive or negative.
}

\secdivider
\clearpage

%% ============================================================================
\section{The Chapter Conclusion}

\ruleheadingbilingual{स्पष्टाधिकारोपसंहारः}{Epilogue to the Spa.s.tādhikāra}

\shlokabilingual{%
इति ब्राह्मस्फुटसिद्धान्ते स्पष्टाधिकारः द्वितीयः ॥
}{%
\textbf{[End of Chapter 2]} Thus the second chapter, the \textit{Spa.s.tādhikāra} (Chapter on True Positions), in the \textit{Brāhmasphu.tasiddhānta}.
}

\hindicommbilingual{%
\textbf{हिन्दी भाषा}---इस प्रकार से ब्रह्मगुप्ताचार्य द्वारा रचित ब्राह्मस्फुटसिद्धान्त में मध्यमाधिकार (प्रथम अधिकार) और स्पष्टाधिकार (द्वितीय अधिकार) की व्याख्या संपन्न हुई। इसमें ग्रहों की मध्यमगति, स्पष्टीकरण, ज्यागणित, मन्दफल, शीघ्रफल, देशान्तर, नतकर्म आदि विषयों का विस्तृत वर्णन किया गया है।
}{%
\textbf{[Hindi Commentary --- Translation]} In this manner, the explanation of the \textit{Madhyamādhikāra} (first chapter) and the \textit{Spa.s.tādhikāra} (second chapter) in the \textit{Brāhmasphu.tasiddhānta}, composed by Ācārya Brahmagupta, has been completed. In this, the mean motions of planets, \textit{sphu.tīkara.na}, sine mathematics (\textit{jyā-ga.nita}), \textit{manda-phala}, \textit{śīghra-phala}, \textit{deśāntara}, \textit{nata-karma}, and other subjects have been described in detail.
}

\vspace{1em}

\bilingual{%
\textit{[स्पष्टाधिकारस्य विस्तृतगणना ग्रहाणां मन्दशीघ्रसंस्कारेभ्यः, परिधिमानसारणीभ्यः, ज्यामितीयोपपत्तिभ्यः, बुधशुक्रयोः विशेषनियमाभ्यश्च पृष्ठ ८९६--९७८ अनुवर्तते।]}
}{%
\textit{[The Spa.s.tādhikāra continues with detailed calculations for each planet's \textit{manda} and \textit{śīghra} corrections, tables of \textit{paridhi} values, geometric \textit{upapatti}s (proofs), and special rules for Mercury and Venus across pages 896--978 of the original edition.]}
}

\clearpage

%% ============================================================================
\appendix
\chapter*{Notes}
\addcontentsline{toc}{chapter}{Notes}
\markboth{Notes}{Notes}

\bilingual{%
\textbf{\dev{संपादकीयटिप्पनयः}}\\[0.3em]
एतत्‌ पुनर्मुद्रणं ब्रह्मगुप्तकृत ब्राह्मस्फुटसिद्धान्तस्य संस्कृत-हिन्दीटीकासहितस्य (पृष्ठ ८३९--९७८) समाविशति।
}{%
\textbf{About this Edition}\\[0.3em]
This typeset edition reproduces the content of Brahmagupta's \textit{Brāhmasphu.tasiddhānta} as found in the Sanskrit--Hindi commentary edition (pages 839--978), presented in a bilingual parallel-column format with English translation and IAST transliteration.
}

\vspace{1em}

\bilingual{%
\textbf{\dev{मुख्यगणितविषयाः}}\\[0.5em]
१. \textbf{ज्यासारणी}---ब्रह्मगुप्तकृताः चतुर्विंशतिः ज्याः अपूर्वसूक्ष्मतां दर्शयन्ति। भूतसङ्ख्यापद्धत्या संख्याः काव्यतया निरूपिताः।\\[0.2em]
२. \textbf{देशान्तरकर्म}---दीर्घ्यान्तरनिर्णायकसूत्रं भूगोलपरिधिम्‌ (५००० योजनानि), अक्षांशान्तरं, कालसंशोधनं च सम्बध्नाति।\\[0.2em]
३. \textbf{ग्रहसंस्काराः}---मन्दशीघ्रफले उन्नतपरिवलयसिद्धान्तस्य प्रतिनिधित्वं कुर्वन्ति, प्रत्येकस्य ग्रहस्य विभिन्नपरिधिभिः।\\[0.2em]
४. \textbf{ज्यामितीयोपपत्तयः}---टीकायां परिवलयमाडलस्य त्रिकोणमितीयसम्बन्धान्‌ उपयुज्य सविस्तरं युक्तयः प्रदर्शिताः।
}{%
\textbf{Notable Mathematical Content}\\[0.5em]
1. \textbf{Sine tables:} The 24 sines given by Brahmagupta are notable for their precision. The values use the \textit{bhūta-saṅkhyā} (word-numeral) system where each number is encoded in poetic words.\\[0.2em]
2. \textbf{Longitude correction (deśāntara):} The formula for computing longitude differences uses the relationship between terrestrial circumference (5000 yojanas), latitude difference, and the resulting time correction.\\[0.2em]
3. \textbf{Planetary equations:} The \textit{manda} and \textit{śīghra} equations represent sophisticated epicyclic theory, with different epicycle sizes for each planet that vary with anomaly.\\[0.2em]
4. \textbf{Geometric proofs (upapatti):} The commentary includes detailed geometric justifications using trigonometric relationships in the epicycle model.
}

\vspace{1em}

\bilingual{%
\textbf{\dev{टङ्कनतान्त्रिकटिप्पनयः}}\\[0.5em]
\textbullet{} \textbf{यन्त्रम्‌:} ज़ीलatek (XeLaTeX) सम्पूर्णयुनीकोडसहितम्‌\\
\textbullet{} \textbf{देवनागरीय मुद्राक्षरश्रेणिः:} Noto Serif Devanagari\\
\textbullet{} \textbf{आङ्ग्लमुद्राक्षरश्रेणिः:} Linux Libertine O\\
\textbullet{} \textbf{मूलस्रोतः:} ब्राह्मस्फुटसिद्धान्त-संस्कृतहिन्दीसंस्करणस्य पृष्ठसङ्ख्या ८३९--९७८
}{%
\textbf{Technical Notes on Typesetting}\\[0.5em]
\textbullet{} \textbf{Engine:} XeLaTeX with full Unicode support\\
\textbullet{} \textbf{Devanagari font:} Noto Serif Devanagari\\
\textbullet{} \textbf{English font:} Linux Libertine O\\
\textbullet{} \textbf{Source:} Scanned pages 839--978 of the Brahmasphutasiddhanta Part 2 edition
}

\vspace{1cm}
\begin{center}
\rule{0.4\textwidth}{0.5pt}\\[0.5cm]
{\small\textit{End of Part 2, Chunk 1 (Pages 839--978) --- Bilingual Edition}}
\end{center}

\end{document}



% ============================================================
% Source: translations/en/indian/brahmagupta-brahmasphutasiddhanta-pt2-chunk2_bilingual.tex
% ============================================================

% ========================================================================
% Brahmasphutasiddhanta Part 2, Chunk 2 — English-Sanskrit Bilingual Edition
% Author: Brahmagupta (c. 628 CE)
% Sanskrit text with English translation
% ========================================================================
\documentclass[11pt,a4paper]{article}

%% -------- Core Packages --------
\usepackage{fontspec}
\usepackage{polyglossia}
\usepackage{amsmath,amssymb,amsthm}
\usepackage{geometry}
\usepackage{graphicx}
\usepackage{xcolor}
\usepackage{titlesec}
\usepackage{fancyhdr}
\usepackage{enumitem}
\usepackage{array,booktabs,longtable,multirow}
\usepackage{hyperref}
\usepackage{indentfirst}

%% -------- Page Geometry --------
\geometry{
  paperwidth=210mm,paperheight=297mm,
  left=18mm,right=18mm,top=25mm,bottom=25mm,
  headheight=14pt,footskip=30pt
}

%% -------- Language Setup --------
\setmainlanguage{english}
\setotherlanguage{sanskrit}
\setotherlanguage{hindi}

%% -------- Font Configuration --------
\setmainfont[Ligatures=TeX]{TeXGyreTermes}
\setsansfont{TeXGyreHeros}
\setmonofont{DejaVu Sans Mono}

%% Devanagari (Sanskrit/Hindi) Font
\newfontfamily\devanagarifont[Script=Devanagari,Scale=1.1]{Noto Serif Devanagari}
\newfontfamily\devanagarifontsf[Script=Devanagari,Scale=1.1]{Noto Sans Devanagari}

%% -------- Colors --------
\definecolor{titlecolor}{RGB}{45,55,72}
\definecolor{sectioncolor}{RGB}{55,65,81}
\definecolor{notecolor}{RGB}{120,53,15}
\definecolor{rulecolor}{RGB}{180,160,140}
\definecolor{vernumcolor}{RGB}{139,69,19}
\definecolor{shlokabg}{RGB}{250,248,245}
\definecolor{transbg}{RGB}{245,248,250}
\definecolor{sansborder}{RGB}{200,180,160}

%% -------- Section Formatting --------
\titleformat{\section}
  {\normalfont\Large\bfseries\color{sectioncolor}\centering}
  {\thesection}{0em}{}
\titleformat{\subsection}
  {\normalfont\large\bfseries\color{sectioncolor}\centering}
  {\thesubsection}{0em}{}
\titleformat{\subsubsection}
  {\normalfont\normalsize\bfseries\color{sectioncolor}\centering}
  {\thesubsubsection}{0em}{}

%% -------- Header/Footer --------
\pagestyle{fancy}
\fancyhf{}
\fancyhead[L]{\small\textit{Br\=ahmasphu\d{t}asiddh\=anta -- Bilingual Edition}}
\fancyhead[R]{\small\thepage}
\renewcommand{\headrulewidth}{0.4pt}
\renewcommand{\footrulewidth}{0pt}

%% -------- Custom Commands --------
\newcommand{\longline}{\noindent\rule{\textwidth}{0.4pt}}
\newcommand{\shortline}{\noindent\rule{0.5\textwidth}{0.4pt}\par}
\newcommand{\cnote}[1]{\textcolor{notecolor}{\textit{#1}}}

%% Devanagari helpers
\newcommand{\dev}[1]{{\devanagarifont #1}}
\newcommand{\vs}[1]{\textcolor{vernumcolor}{\textbf{#1}}}

%% Bilingual minipage environment
\newcommand{\bilingual}[2]{%
  \begin{minipage}[t]{0.48\textwidth}
  \textbf{\small Sanskrit Text}\\[0.3em]
  \small\devanagarifont #1
  \end{minipage}%
  \hfill
  \begin{minipage}[t]{0.48\textwidth}
  \textbf{\small English Translation}\\[0.3em]
  \small #2
  \end{minipage}%
}

%% Sanskrit-only environment (for verses)
\newenvironment{sktcol}{%
  \begin{minipage}[t]{0.48\textwidth}
  \small\devanagarifont
}{\end{minipage}}

%% English-only environment
\newenvironment{encol}{%
  \begin{minipage}[t]{0.48\textwidth}
  \small
}{\end{minipage}}

%% Shloka (verse) environment -- left column
\newenvironment{shloka}{%
  \begin{quote}
  \small
  \setlength{\leftmargin}{2em}
  \devanagarifont
}{\end{quote}}

%% Commentary environment
\newenvironment{vyakhya}{%
  \begin{quote}
  \small
  \setlength{\leftmargin}{1em}
  \devanagarifont
}{\end{quote}}

%% Inline Devanagari for mixed text
\newcommand{\skt}[1]{{\devanagarifont #1}}

%% Section separator
\newcommand{\sectsep}{\par\vspace{0.8em}\longline\vspace{0.8em}}

%% -------- Hyperref Setup --------
\hypersetup{
  colorlinks=true,
  linkcolor=sectioncolor,
  citecolor=sectioncolor,
  urlcolor=sectioncolor,
  pdfencoding=auto,
  pdftitle={Brahmasphutasiddhanta Part 2 Chunk 2 -- Bilingual Edition},
  pdfauthor={Brahmagupta (c. 628 CE)}
}

%% -------- TOC Depth --------
\setcounter{tocdepth}{3}
\setcounter{secnumdepth}{3}

%% ========================================================================

\begin{document}

%% ========================================================================
%% TITLE PAGE
%% ========================================================================
\begin{titlepage}
\begin{center}
\vspace*{2.5cm}

{\Huge\bfseries\color{titlecolor} Br\=ahmasphu\d{t}asiddh\=anta}\\[0.6cm]
{\LARGE\color{sectioncolor} of Brahmagupta}\\[0.3cm]
{\large (c.\ 628 CE)}\\[1.5cm]

\shortline\\[0.5cm]
{\Large\bfseries Part 2 --- Chunk 2}\\[0.3cm]
{\large \textbf{English-Sanskrit Bilingual Edition}}\\[0.3cm]
{\normalsize Pages 979--1118 of the Complete Work}\\[0.5cm]
{\small \textit{Bilingual edition with Sanskrit verses in Devanagari and English translations}}\\[1cm]
\shortline\\[1.5cm]

{\large Contents: Spa\d{s}\d{t}\=adhik\=ara, Tripra\'sn\=adhik\=ara,}\\[0.2cm]
{\large Candra-graha\d{n}\=adhik\=ara, and Mathematical Sections}\\[2cm]

\shortline\\[0.8cm]
{\normalsize Covering topics:}\\[0.3cm]
{\small Planetary mean motion $\bullet$ Manda and \'S\=ighra corrections}\\[0.2cm]
{\small Sine table construction $\bullet$ Second difference interpolation}\\[0.2cm]
{\small Luni-solar astronomy $\bullet$ Spherical astronomy}\\[0.2cm]
{\small P\=a\d{t}\=iga\d{n}ita arithmetic operations}\\[1.5cm]

\vfill

{\small\textit{Typeset in XeLaTeX with Devanagari support}}\\[0.3cm]
{\small Original: Indian Institute of Astronomical \& Sanskrit Research}

\end{center}
\end{titlepage}

%% ========================================================================
%% TABLE OF CONTENTS
%% ========================================================================
\tableofcontents
\newpage

%% ========================================================================
%% INTRODUCTION
%% ========================================================================
\section*{Introduction}
\addcontentsline{toc}{section}{Introduction}

\begin{minipage}[t]{0.48\textwidth}
\textbf{\small Original Description}

{\small\devanagarifont एतत् पाठं ब्रह्मगुप्तस्य ब्रह्मस्फुटसिद्धान्तस्य द्वितीयभागस्य द्वितीयखण्डं समाविशति। एषः पाठः स्पष्टाधिकारं, त्रिप्रश्नाधिकारं, चन्द्रग्रहणाधिकारं, तथा गणितीयसूत्राणि च समाविशति।}
\end{minipage}%
\hfill
\begin{minipage}[t]{0.48\textwidth}
\textbf{\small English Translation}

This volume covers pages 979--1118 of Brahmagupta's \textit{Br\=ahmasphu\d{t}asiddh\=anta}, one of the most important works of ancient Indian mathematics and astronomy. Composed in 628 CE, this treatise contains 24 chapters (\textit{adhik\=aras}) covering planetary calculations, eclipses, mathematical algebra, and astronomical instruments.

This bilingual edition presents the original Sanskrit (\textit{grantha}) verse form with English translations in parallel columns. Each chapter begins with Sanskrit \textit{\'slokas} (verses) followed by commentary (\textit{vy\=akhy\=a}), mathematical formulas, and detailed explanations.
\end{minipage}

\vspace{0.5em}
\longline


%% ========================================================================
%% CHAPTER: SPAṣṬĀDHIKĀRA
%% ========================================================================
\section{\dev{स्पष्टादिकारः}}
\addcontentsline{toc}{section}{Spa\d{s}\d{t}\=adhik\=ara\d{h} --- Chapter on True Positions}

\begin{center}
\textit{Spa\d{s}\d{t}\=adhik\=ara\d{h} --- Chapter on the Determination of True Positions}
\end{center}

\begin{minipage}[t]{0.48\textwidth}
\textbf{\small Original Sanskrit}

\small\devanagarifont
अस्याध्यायस्य विषयः --- मध्यमागतस्य ग्रहस्य स्पष्टस्थानस्य निर्धारणम्। ये समाकलनाः ग्रहकक्षाया अणुता च परिमण्डलवृत्तस्य व्याख्यां सम्मिलयन्ति, ते स्पष्टीकरणस्य अस्याध्याये दीयन्ते।
\end{minipage}%
\hfill
\begin{minipage}[t]{0.48\textwidth}
\textbf{\small English Translation}

This chapter deals with the methods for computing the true (corrected) positions of planets from their mean (average) positions. The corrections account for the eccentricity of planetary orbits and the epicyclic models described in Indian astronomy.
\end{minipage}

\sectsep

%% ========================================================================
%% ŚLOKA 48--49
%% ========================================================================
\subsection{\dev{श्लोक ४८--४९}}
\label{subsec:48-49}

\begin{minipage}[t]{0.48\textwidth}
\textbf{\small Sanskrit Text}

\small\devanagarifont
वा भवतीत्यर्थः। पादे शीघ्रकेन्‍द्रमर्धकं तदतीतस्य प्रथ वक्रानुवक्रकेन्‍द्रमर्धक-मतस्तस्यैव वक्रानुवक्रदिनं ज्ञात्वा सकृद्‌ग्रहः स्पष्टः कार्यस्तदर्धयोःशीघ्रकेन्‍द्रवक्रानु-वक्रो निक्षिप्याविति प्रत्यय वासना। मध्याच्छीघ्रनीचोच्चवृत्तमध्य यावदूर्ध्वं-कोटिः परमफलज्यातुल्यं शीघ्रनीचोच्चवृत्तमध्यमार्धे भुजा। पूर्वापरेया वा तयो-र्गुणितं मूलं त्रयंककर्णः परमफलज्या प्रोक्ता मध्य यावतस्य कर्णस्य शीघ्रनीचोच्च-वृत्ताशालाकायाःचान्नरे यावच्छीघ्रनीचोच्चवृत्तसूर्याद्येन प्रमिति प्रतिमण्डलपरिधौ-स च तत्र प्रदेशे परच्छादग्छन्‍नुत्‌सन्नते। तस्माद्‌दुर्न्नने सर्वं गोले दर्शयेत्।\hfill\vs{॥४८--४९॥}
\end{minipage}%
\hfill
\begin{minipage}[t]{0.48\textwidth}
\textbf{\small English Translation}

\textbf{Verses 48--49. \'S\=ighra-spa\d{s}\d{t}ikara\d{n}a (True Position via \'S\=ighra Correction):}

Having taken half the \'s\=ighra-kendra (anomaly from the \'s\=ighrocca) of the traversed arc, and half of the retrograde-direct kendra of that same planet, knowing its retrograde-direct period, the true planet should be computed. The construction is: the square-root of the sum of the squares of the \'s\=ighra-kendra retrograde-direct quantities is the divisor. The vertical line (ko\d{t}i) from the centre of the \'s\=ighra-n\=icocca-v\=tta (epicycle) up to the mean position equals the parama-phala-jy\=a (maximum equation sine); the horizontal (bhuja) at the centre of the \'s\=ighra-n\=icocca-v\=tta is the product of the east-west components, whose square-root is the tri\~n\~naka-kar\d{n}a. The parama-phala-jy\=a is declared to be proportional to the \'s\=ighra-n\=icocca-v\=tta-\\'s\=al\=ak\=a (radius) on the circumference of the perimandala. That also, when the planet is at the highest point on that perimandala, causes parallax (paracch\=ada). Therefore, in an iron globe one should demonstrate everything.
\end{minipage}

\vspace{0.5em}
\begin{minipage}[t]{0.48\textwidth}
\textbf{\small Commentary (Vy\=akhy\=a)}

\small\devanagarifont
\textit{वि.~भा.}---प्रत्येकटिमार्गपुमनुर्मीरित्यादिपञ्चशीघ्रकेन्‍द्रांशौभौमादिद्वयंहरणां वक्रं भवेद्यदिरचितैस्तैः शीघ्रकेन्‍द्रांशैस्तेषां वक्रारम्भौ भवति, तद्धितैश्चक्रांशो ३६०°अनुवक्रमध्यमगैरमः। मन्दफलस्पष्टमुक्तिः (मन्दस्पष्टागतिः) तदुग्ना (तद्धिता) शीघ्रमुक्तिः (शीघ्रस्पष्टिः) शीघ्रकेन्‍द्रगतिर्भवति तथा तदधिकोन-मागकलामकालस्तथा गतेन दिवसा भवति, शीघ्रात्यकेन्‍द्रभागैरसकृद्धिना-उच्चैःपकर्माणां स्थिरो भूतः केन्‍द्रांशैरिति ॥४८--४९॥
\end{minipage}%
\hfill
\begin{minipage}[t]{0.48\textwidth}
\textbf{\small Commentary Translation}

\textit{[Commentary of Var\=ahamihira:]}

For Mars and Jupiter, the retrograde motion begins at specific \'s\=ighra-kendra a\~n\'sa values. The manda-spa\d{s}\d{t}a (true position via manda correction) and the \'s\=ighra-spa\d{s}\d{t}a (true position via \'s\=ighra correction) are computed. The \'s\=ighra-kendra-gati is the velocity at the \'s\=ighra-kendra. The days traversed and the remaining time are computed from these kendra values. By the \'s\=ighra-kendra divisions, the \textit{ucca} positions remain stable.
\end{minipage}

\sectsep

%% ========================================================================
%% MADHYAMĀGAMA
%% ========================================================================
\subsection{\dev{मध्यमागमः --- ग्रहाणां मध्यमानयनम्}}
\label{subsec:madhyamagama}

\begin{minipage}[t]{0.48\textwidth}
\textbf{\small Sanskrit Verses}

\small\devanagarifont
भचक्रभागान् प्रथमं दिनैर्हत्वा भागलब्धिकाम्\\
कल्पे युगे वाथ विभजेत् तैः स्यात् स पूरकः क्षेपः ॥१२॥\\[0.5em]
स पूरकः क्षेपो गतदिनैर्गुणितो रविवत्सरैः\\
हृतः स मध्यभुक्तिः स्यात् तद्युक्तं मध्यमं ग्रहस्य ॥१३॥
\end{minipage}%
\hfill
\begin{minipage}[t]{0.48\textwidth}
\textbf{\small English Translation}

\textbf{Verses 12--13. Planetary Mean Motion (Madhyam\=agama):}

First, having multiplied the zodiacal signs (bhacakra-bh\=aga, 43,200 minutes of arc) by the given days, divide by the Kalpa or Yuga days [1,577,917,828,000]. The quotient obtained is the \textit{p\=uraka-k\d{s}epa} (accumulator of revolutions).

That same p\=uraka-k\d{s}epa, multiplied by the elapsed days and divided by the solar years, gives the mean daily motion (madhya-bhukti). That, added to or subtracted from the previous mean longitude, gives the mean position (madhyama) of the planet.
\end{minipage}

\vspace{0.5em}
\begin{minipage}[t]{0.48\textwidth}
\textbf{\small Commentary}

\small\devanagarifont
\textit{वा.~भा.}---भचक्रभागान् द्वादशराशिभागान् अर्थात् ४३२०० भागान् प्रथमं दिनैः गुणित्वा रविवत्सरैः १५७७९१७८२८० विभजेत्। ततो लब्धिः पूरकः क्षेपः स्यात्। स एव पूरकक्षेपो गतदिनैर्गुणितो रविवत्सरैर्हृतः स मध्यभुक्तिः स्यात्। सा मध्यभुक्तिः पूर्वमध्यमे युक्ता वा शोधिता वा ग्रहस्य मध्यमं स्यात् ॥१२--१३॥
\end{minipage}%
\hfill
\begin{minipage}[t]{0.48\textwidth}
\textbf{\small Commentary Translation}

\textit{[Var\=ahamihira's Commentary:]}

The zodiacal signs (bhacakra-bh\=aga), that is 43,200 minutes of arc, are first multiplied by the given days and divided by 1,577,917,828,000 solar years. The quotient obtained is the p\=uraka-k\d{s}epa. That same p\=uraka-k\d{s}epa, multiplied by elapsed days and divided by solar years, gives the mean daily motion (madhya-bhukti). That madhya-bhukti, added to or subtracted from the previous mean position, gives the mean (madhyama) of the planet.
\end{minipage}

\vspace{0.5em}
\begin{minipage}[t]{0.48\textwidth}
\small\devanagarifont
मन्दोच्चं शीघ्रच्चं च युगपूरकक्षेपभक्तं तैः\\
गतदिनैर्गुणितं तत् स्यात् मन्दोच्चं शीघ्रोच्चं वा ॥१४॥\\[0.5em]
तत् स्यात् मन्दोच्चं शीघ्रोच्चं वा युगीयात् पूर्वमन्दोच्चात्\\
शीघ्रोच्चाद्वा युतं वा शुद्धं वा मन्दशीघ्रोच्चं स्यात् ॥१५॥
\end{minipage}%
\hfill
\begin{minipage}[t]{0.48\textwidth}
\textbf{\small English Translation}

\textbf{Verses 14--15. Mean Apogee and \'S\=ighrocca:}

The mandocca (apogee of manda epicycle) and the \'s\=ighrocca (apex of \'s\=ighra epicycle), divided by the yuga-p\=uraka-k\d{s}epa and multiplied by elapsed days, gives the mandocca or \'s\=ighrocca.

That mandocca or \'s\=ighrocca, when added to or subtracted from the previous mandocca or \'s\=ighrocca of the yuga, gives the current manda-\'s\=ighrocca.
\end{minipage}

\vspace{0.5em}
\begin{minipage}[t]{0.48\textwidth}
\small\devanagarifont
\textit{वा.~भा.}---मन्दोच्चं शीघ्रोच्चं च गतदिनैर्गुणितं युगपूरकक्षेपभक्तं स्यात्। तत् पूर्वमन्दोच्चाद्वा शीघ्रोच्चाद्वा युक्तं शुद्धं वा कृत्वा मन्दोच्चं शीघ्रोच्चं वा स्यात्। एवं ग्रहस्य मध्यमं मन्दोच्चं शीघ्रोच्चं च ज्ञात्वा स्पष्टीकरणस्य आधारः सिद्ध्यति ॥१४--१५॥
\end{minipage}%
\hfill
\begin{minipage}[t]{0.48\textwidth}
\small
The mandocca and \'s\=ighrocca, multiplied by elapsed days and divided by the yuga-p\=uraka-k\d{s}epa, gives the result. Adding or subtracting this from the previous mandocca or \'s\=ighrocca gives the current values. In this manner, knowing the planet's madhyama, mandocca, and \'s\=ighrocca, the foundation for computing the true position is established.
\end{minipage}

\sectsep

%% ========================================================================
%% MANDA-SPAṢṬIKARAṆA
%% ========================================================================
\subsection{\dev{मन्दस्पष्टीकरणम् --- परमफलनयनम्}}
\label{subsec:manda-spashtikarana}

\begin{minipage}[t]{0.48\textwidth}
\textbf{\small Sanskrit Verses}

\small\devanagarifont
मन्दकेनापक्रान्तं मन्दफलं तत् परमं ग्रहस्य\\
ज्यार्धेन हृत्वा लब्धं तत् परमं फलं ग्रहस्य ॥१६॥\\[0.5em]
मन्दफलज्यया हत्वा परमफलज्यां तु यल्लब्धम्\\
तत् परमं फलं स्यात् तेन मन्दस्पष्टः स ग्रहः ॥१७॥
\end{minipage}%
\hfill
\begin{minipage}[t]{0.48\textwidth}
\textbf{\small English Translation}

\textbf{Verses 16--17. Manda-spa\d{s}\d{t}ikara\d{n}a (Equation of Centre):}

The maximum manda-phala (equation of centre) of the planet, diminished by the mandaka (eccentric anomaly divisor), divided by the half-sine (jy\=a-ardha), gives that maximum phala of the planet.

Multiplying the manda-phala-jy\=a by the parama-phala-jy\=a, the result obtained is the parama-phala (maximum equation). By this, the planet becomes manda-spa\d{s}\d{t}a (corrected via the equation of centre).
\end{minipage}

\vspace{0.5em}
\begin{minipage}[t]{0.48\textwidth}
\small\devanagarifont
\textit{वा.~भा.}---मन्दकेन दोःकोट्योः कृतिभ्यां पदं मन्दकर्णः। तस्य मन्दज्या हत्वा त्रिज्यां विभज्य मन्दफलं लभ्यते। तत् परमं फलं ग्रहस्य मध्यमात् शोधयित्वा योजयित्वा वा मन्दस्पष्टः स ग्रहः स्यात् ॥१६--१७॥
\end{minipage}%
\hfill
\begin{minipage}[t]{0.48\textwidth}
\small
The mandaka is the square-root of the sum of squares of the do\d{h} (sine) and ko\d{t}i (cosine). The manda-jy\=a multiplied by the trijy\=a and divided gives the manda-phala. That parama-phala, when subtracted from or added to the mean position, gives the manda-spa\d{s}\d{t}a (true) position of the planet.
\end{minipage}

\sectsep

%% ========================================================================
%% ŚLOKA 50--51
%% ========================================================================
\subsection{\dev{श्लोक ५०--५१}}

\begin{minipage}[t]{0.48\textwidth}
\textbf{\small Sanskrit Verses}

\small\devanagarifont
मार्गनियमरचैपलक्षा यन्त्राभियोगातिशयाच्च तस्मादुपपन्नं कक्षामण्डलादिषु-विन्यस्तेन एते बोधायास्तमयभागा अविभक्ते, ग्रहे विभक्ते मण्डलवराग्‌दुर्घटते,\\
तद्‌दुःखमुद्यास्तमदण्डायो भविष्यतीति ॥५०--५१॥
\end{minipage}%
\hfill
\begin{minipage}[t]{0.48\textwidth}
\textbf{\small English Translation}

\textbf{Verses 50--51. Planetary Path and Visibility:}

Having fixed the path regulation to the epicycles and instruments, from that the installation in the kak\d{s}a-ma\d{n}\d{d}alas (orbits) is established. These rising and setting portions, when undivided [in the orbit], present the characteristics of the ma\d{n}\d{d}ala divisions; that difficulty in rising and setting will arise when the planet is divided [by epicyclic motion].
\end{minipage}

\vspace{0.5em}
\begin{minipage}[t]{0.48\textwidth}
\small\devanagarifont
\textit{वा.~भा.}---प्रयुग्ना बुधशुक्रयोर्दयास्तर्यपरिक्षानार्थमायातिमाह। शीघ्रात्यकेन्‍द्रमार्गो-रित्यनुवर्तन्ते खराः ५० एतावद्‌बुधः सवितःशीघ्रकेन्‍द्रमार्गैर्बुधस्य परचात् उदयो भवति,\\
जिनः २४ एतावद्‌बुधः शुक्रयोदयपरचात् इष्टतिथिमः १५५ एतावद्‌बुधरव-शीघ्रात्यकेन्‍द्रमार्गः परचादस्तमयो बुधस्य मुनिनगेन्‍दुभिः १७७ एतावद्‌बुधः परचात्, शुक्रस्य-स्तमयः उदयास्तमयो व्यस्तो मण्डलमार्गैस्तदुग्नः प्राक् ॥५२--५३॥
\end{minipage}%
\hfill
\begin{minipage}[t]{0.48\textwidth}
\small
The \'s\=ighra-kendra values for Mercury and Venus are described. When Mercury's \'s\=ighra-kendra equals 50, it rises after the Sun; when it equals 24, it rises before [specific values for visibility]. The rising and setting are reversed relative to the orbital path.
\end{minipage}

\sectsep

%% ========================================================================
%% CHAPTER: TRIPRAŚNĀDHIKĀRA
%% ========================================================================
\section{\dev{त्रिप्रश्नाधिकारः}}
\addcontentsline{toc}{section}{Tripra\'sn\=adhik\=ara\d{h} --- Chapter on the Three Questions}

\begin{center}
\textit{Tripra\'sn\=adhik\=ara\d{h} --- Chapter on the Three Questions of Diurnal Motion}
\end{center}

\begin{minipage}[t]{0.48\textwidth}
\small\devanagarifont
अस्याध्यायस्य विषयः --- दैनिकगतिस्य त्रयः प्रश्नाः: लम्बनस्य (देशान्तरस्य) नाडीविशेषस्य छायानयनस्य च गणनम्।
\end{minipage}%
\hfill
\begin{minipage}[t]{0.48\textwidth}
\small
This important chapter deals with the ``three questions'' (\textit{tripra\'sna}) of spherical astronomy: determining the \textit{lambana} (parallax in longitude), the \textit{n\=a\d{d}\=i} difference (time difference), and related quantities needed for computing solar and lunar phenomena.
\end{minipage}

\sectsep

%% ========================================================================
%% ŚLOKA 54--55
%% ========================================================================
\subsection{\dev{श्लोक ५४--५५}}

\begin{minipage}[t]{0.48\textwidth}
\textbf{\small Sanskrit Verses}

\small\devanagarifont
इदानीं कुजादिद्वयहरणामुदयास्तकेन्‍द्रांशानाम्\\
षष्ट्यंशः २८ कृतचक्रे १४ मुनिनगेन्‍दुभिः १७ भौमजीवरविबुधानाम्\\
उदयः प्रागस्तमयस्तदुच्चकांशाः पश्चात् ॥५४॥\\[0.5em]
स्वारे ५० जिने २४ क्ष सितयोरिष्टतिथिमि १५५ मुनिनगेन्‍दुभिः १७७ पश्चात्\\
उदयास्तमयो व्यस्तो मण्डलमार्गैस्तदुग्नः प्राक् ॥५५॥
\end{minipage}%
\hfill
\begin{minipage}[t]{0.48\textwidth}
\textbf{\small English Translation}

\textbf{Verses 54--55. Rising and Setting \'S\=ighra-kendra Values:}

Now, the \'s\=ighra-kendra a\~n\'sa values for the rising and setting of Mars and the others: When Mercury's \'s\=ighra-kendra equals 28 [bhauma], 14 [cakre], 17 [muni-nagendra], then the rising of Mars, Jupiter, and Mercury occurs; the setting is opposite to these.

[Values: 50 in Sv\=at\=i, 24 in Jina, 155 in the desired tithi, 177 in Muni-nagendra] The rising and setting are reversed relative to the orbital path, the forward one being retrograde.
\end{minipage}

\vspace{0.5em}
\begin{minipage}[t]{0.48\textwidth}
\small\devanagarifont
\textit{वा.~भा.}---शीघ्रात्यकेन्‍द्रमार्गरत्यमौमीमस्य प्रागुदयो भवति २८ स्वातीशी-घ्रकेन्‍द्रमार्गैः कृतचक्रे १४ जीवस्य प्रागुदयो भवति, स्वातीघ्रकेन्‍द्रमार्गैर्मुनिनगेन्‍दुभिः १७ शनैश्चरस्यो भवति। प्रागस्तमयास्तु परचाद्‌दृश्यति। चक्रांशैस्ते उन्नतदुग्नाः यथा-स्वोदयमार्गैःस्नाच्च कर्कांशाका ये विशेषा भवतीत्यर्थः। चक्रांशाकाः प्रसिद्धा एव तथाऽपि भौमस्यास्तमयशीघ्रकेन्‍द्रभागाः ३३२ जूतो ३८५, शनैः ३५३, सवितुः/बुधस्य-स्तकेन्‍द्राद्यैः, राश्यादिकानि। भो। उ। २८ जौ उ। १४ धा उ। १७ शौ १५५ भो १७७
\end{minipage}%
\hfill
\begin{minipage}[t]{0.48\textwidth}
\small
The \'s\=ighra-kendra determines the heliacal rising. At 28 degrees, Mars rises before the Sun; at 14, [a value for] Jupiter; at 17, Saturn. The heliacal settings are opposite. These kendra values are well-known; for Mars, the setting \'s\=ighra-kendra is 332; for Jupiter, 385; for Saturn, 353.
\end{minipage}

\sectsep

%% ========================================================================
%% ŚLOKA 56--57
%% ========================================================================
\subsection{\dev{श्लोक ५६--५७}}

\begin{minipage}[t]{0.48\textwidth}
\small\devanagarifont
भ्रात्रातीतानां ग्रहाणां दर्शने प्राप्नोत्। तद्धि-कोनां भागकला मन्दफलस्पष्टयुक्तयुक्षीघ्रयुक्या। हृता दिवसा इति न्यायेनोप-पत्तिः। तथाऽपि रविकक्षया सर्वथा स्पष्टत्वातिमद्धंले नीचोच्चवृत्तमध्ये भ्रमति।\\
ततो यदा परमे प्रतिमण्डलोन्नतप्रदेशे ग्रहः स्थितो भवति तदा समापेक्षं क्षोभो-भवति समरव रविरत्र एव सूर्यस्तदा ग्रहो नोपलभ्यते। मनाऽपि रविकिंरव्य-परिहृतिर्भवतत्ते यथा-यथा ग्रहोदयस्तमयोर्दृश्यतां प्रथममेवोदयं यान्ति ततः-शीघ्रात्यादर्कस्यास्तमयेऽपि परिवर्त्यं शीघ्रमार्गेणाकः पुनरुप्तं ह्रासमभ्येति पश्च-मादिर्मावात्। अतएवास्तमयोदं ग्रह उपलभ्यते। उदयास्तमयो यदा भवतो
\end{minipage}%
\hfill
\begin{minipage}[t]{0.48\textwidth}
\small
When the elapsed [days' values] of the \textit{bh\=a} (signs) and minutes of arc, joined with the manda-phala-spa\d{s}\d{t}a and \'s\=ighra corrections, are divided by the days --- thus is the derivation by this method. Still, since all [planets] are corrected by the solar orbit, the planet revolves in the n\=icocca-v\=tta (epicycle) in that time.

When the planet is situated at the highest point of the perimandala, then there is maximum perturbation; at that time, since the Sun is nearby, the planet is not observed. The rising and setting are first seen when the \'s\=ighra-kendra is appropriate; then, even at the Sun's setting, the epicyclic motion having changed, the planet gradually attains reduction [in elongation]. Therefore, the planet is observed at rising and setting.
\end{minipage}

\sectsep

%% ========================================================================
%% PAÑCAYĀNANAM
%% ========================================================================
\subsection{\dev{इदानीं पञ्चजयानयनमाह}}

\begin{minipage}[t]{0.48\textwidth}
\small\devanagarifont
जिनमागज्यागुणिता सूर्यज्या व्यासदलहृता लब्धम्\\
इष्टक्रमजीवा विषुवद्दर्द्धिविरणा सवितुः ॥५८॥\\[0.5em]
इष्टक्रमभागे त्रिज्यावर्गाद्ध्रोत्वं शेषपदम्\\
विषुवद्दर्द्धिविरणात् स्वाहोरात्राद्विवस्कमः ॥५९॥\\[0.5em]
क्रान्तिज्या विषुवच्छायया गुणा द्वादशोद्धृता स्थितिजा।\\
स्वाहोरात्रनुष्टा व्यासार्धेनाहता मता ॥६०॥\\[0.5em]
स्वाहोरात्रार्धेन सम्बुद्धिज्याचनुष्करप्रमाणः।\\
ते पृथक्ता विच्छिन्ना विच्छिन्ना नागिका शुष्का ॥६१॥
\end{minipage}%
\hfill
\begin{minipage}[t]{0.48\textwidth}
\small
\textbf{Verses 58--61. Declination and Shadow Calculations:}

The \textit{jina-m\=aga-jy\=a} (24-degree sine) multiplied by the solar jy\=a and divided by the semi-diameter gives the result. The desired ascensional difference (\textit{kranti-j\=iv\=a}) from the equinox is for the Sun.

At the desired arc, subtracting from the square of the trijy\=a, the remainder's root, from the equinox semi-duration, the day-duration is obtained.

The declination-sine (kr\=anti-jy\=a), multiplied by the equinoctial shadow and divided by 12, gives the \textit{sthiti-j\=a}. The day-and-night pr\=a\d{n}a, multiplied by the semi-diameter, are considered.

The day-and-night half combined with the ascensional difference gives the arc-measure. Those [pr\=a\d{n}as], separated, are the divided [portions] of the n\=agik\=a [serpent of time].
\end{minipage}

\sectsep

%% ========================================================================
%% ŚLOKA 62--63
%% ========================================================================
\subsection{\dev{श्लोक ६२--६३}}

\begin{minipage}[t]{0.48\textwidth}
\small\devanagarifont
द्वादशांशोनतं जीवा कार्याः ताश्च भूतादिनगमनानां १३२५\\
दुःखयदनाख्यमा १८६३ मिथुनस्य वमनिरदा ३०३० अनः सवातीतानां प्राप्य-दुर्बलतीर्थाः। जिनमागज्यागुणानां सूर्यज्योत्नियमेन पुमनु क्षोणिज्या कार्या-ताश्च मैथुनं वैद्यमठद्विकलमाधात् ६६°३०' बुधस्य छद्धरभवाः १९°५९ मिथुनस्य\\
नवरदचक्रा १३°३०। एनाभिरेव प्राप्तवत्। स्वातीतग्रहाणां त्रिज्याकर्मवर्षो-च्चयाख्याया गुणोत्कयादिना प्राप्तवत्। तद्‌वेव विषु-वच्छायाया कार्या। ताश्च स्वातीत विज्ञेय स्वदेशज्याचरदलसंस्कृतां पुमनु-प्राणा भवन्ति।
\end{minipage}%
\hfill
\begin{minipage}[t]{0.48\textwidth}
\small
\textbf{Verses 62--63. Twelve-Degree Sines and Ascensional Differences:}

The sines at twelve-degree intervals are to be computed; for the Earth etc. 1325 [values apply]. The difficult [computations] yield 1863; for Mithuna [Gemini], the vomiting-sines are 3030. For the elapsed [periods], the feeble [sines] are obtained. The \textit{jina-m\=aga-jy\=a} multiplied by the solar jy\=a gives the \textit{pum\=anu-k\d{s}o\d{n}i-jy\=a} (terrestrial sine). The Mithuna values at 66\degree30', Mercury's [value] at 19\degree59', and the nine-fold cycle at 13\degree30'. By these, the equinoctial shadow is computed. These, combined with the ascensional difference for one's own location, give the pr\=a\d{n}as.
\end{minipage}

\sectsep

%% ========================================================================
%% ŚLOKA 64--67
%% ========================================================================
\subsection{\dev{श्लोक ६४--६७}}

\begin{minipage}[t]{0.48\textwidth}
\small\devanagarifont
तदानीं मेषादिनगानां चरखण्डयाथनमाह\\
मेषकूर्पंमिथुनानां जीवा कार्याः ताश्च भूतादिनगमनानां १३२५\\
दुःखयदनाख्यमा १८६३ मिथुनस्य वमनिरदा ३०३० अनः सवातीतानां प्राप्य-दुर्बलतीर्थाः प्राप्तवत् ॥६४॥\\[0.5em]
मेषकूर्पंमिथुनानां त एव क्षेत्रा कर्कीसहकन्यानामथ क्षेत्रा तुलाद्विश्चक-घटयो मकरकुम्भमीनानामथवाम्ना। त्वकाल्यां स्वाहोरात्रनुतानि\\
विनस्य प्रदर्शयेत, स्पष्टाद्यध्याय एव चरदलानयने मया प्रदर्शितानि विदुषामपि-प्रदर्शयेते स्वदेशे। यद् क्षिति स्वाहोरात्रनुते क्षितिज्योच्छाययोरन्तरे प्राणाः\\
यदेव लग्नाकारबुद्धिरस्य तदन्तरकालं पृच्छतीति।
\end{minipage}%
\hfill
\begin{minipage}[t]{0.48\textwidth}
\small
\textbf{Verses 64--67. Aries-to-Gemini Sine Computation:}

Now, the computation of \textit{cara-kha\d{n}\d{d}a} (ascensional difference segments) for Aries etc.:

For Aries, Taurus, and Gemini, the same fields apply; for Cancer, Leo, and Virgo, the same fields; and for Libra, Scorpio, and Sagittarius, and similarly for Capricorn, Aquarius, and Pisces.

The day-and-night durations should be shown for each; in the Spa\d{s}\d{t}\=adhik\=ara itself, I have demonstrated the cara-dala computation for the learned. When, at a location, the pr\=a\d{n}as between the terrestrial sine and shadow are computed, this interval is the time-difference inquired for the lagna (ascendant) computation.
\end{minipage}

\sectsep

%% ========================================================================
%% CHĀYĀKARAṆĪ
%% ========================================================================
\subsection{\dev{छायाकरणी}}

\begin{minipage}[t]{0.48\textwidth}
\small\devanagarifont
हृज्या द्वादशगुणिता विभाजिता शङ्कुना फलं छाया।\\
व्यासार्धं छेदहतं विषुवत्कर्णहृतं कर्णः ॥२८॥\\[0.5em]
विषुवत्कर्णगुणायामा व्यासार्धकृतिः फलं सौम्ये।\\
मयुखद्विज्याहीनं युक्तं याम्ये धनुष्करप्रमाणः ॥२९॥\\[0.5em]
सौम्ये युतं विहीनं याम्ये प्रागपरयोः प्राणाः ॥३०॥\\[0.5em]
भ्रष्टो गतविशेषाः फलमत्याया विशेष शेषस्य।\\
धनुःरविकर्मजीवार्धिः पूर्वापरयोर्नतप्राणाः ॥३१॥
\end{minipage}%
\hfill
\begin{minipage}[t]{0.48\textwidth}
\small
\textbf{Verses 28--31. Shadow and Gnomon Formulas (Ch\=ay\=akara\d{n}\=i):}

The altitude-sine (h\d{r}-jy\=a), multiplied by 12 and divided by the \'sa\~nku (gnomon), gives the shadow (ch\=ay\=a). The square of the semi-diameter divided by the equinoctial hypotenuse gives the hypotenuse (kar\d{n}a).

The product of the equinoctial hypotenuse and the ko\d{t}i, divided by the square of the semi-diameter, gives the result. In the north (s\=aumya), diminished by the sun's diameter-sine; in the south (y\=amy), added --- this is the arc-measure of the bow.

In the north, added; in the south, diminished --- these are the east-west pr\=a\d{n}as. The elapsed specific portions; the remainder of the excess portion is the difference. Half the day-sine is the east-west \textit{nata-pr\=a\d{n}a} (hour-angle in time units).
\end{minipage}

\vspace{0.5em}
\begin{minipage}[t]{0.48\textwidth}
\small\devanagarifont
\textit{वा.~भा.}---यस्यक्षायाया दिनगतिशेषानयनमिष्यते तस्यक्षायायाः छाया-कर्णां कृत्वा तेन स्वाहोरात्राद्‌ गुण्येत्। ततस्तेन स्वाहोरात्रार्धेन छायाकर्णहृतेन-मताया कृता व्यासार्धकृतिः कि मृत्या विषुवत्कर्णगुणायामाः फलं ज्या भवति।\\
सफलं सौम्ये गोले छायुबुद्धिज्याहीनं याम्ये तथैव युक्तं कृत्वा यद्‌भवति। तस्य धनुः-क्षेत्राणां कार्यस्तदनुस्तद्द्वैध्वंसिकचरदलप्राणौः सौम्ये गोलेषु युतं याम्येहीनं कृत्वा-प्रागपरयोः प्राणा भवति।
\end{minipage}%
\hfill
\begin{minipage}[t]{0.48\textwidth}
\small
Having computed the ch\=ay\=a-kar\d{n}a (hypotenuse of gnomon and shadow), multiply by the day-and-night. Then, divide by the day-and-night half and the ch\=ay\=a-kar\d{n}a; the resulting quantity, compared with the square of the semi-diameter, gives the product of the equinoctial hypotenuse and the ko\d{t}i as a sine.

In the northern hemisphere, subtract the ch\=ay\=a-buddhi-jy\=a; in the southern, add similarly. The arc of this result is computed, then with the cara-dala-pr\=a\d{n}a, in the northern sphere added and in the southern subtracted, giving the east-west pr\=a\d{n}as.
\end{minipage}

\sectsep

%% ========================================================================
%% PRAKĀRĀNTAREṆA CHĀYĀKARAṆĪ
%% ========================================================================
\subsection{\dev{प्रकारान्तरेण छायाकरणी}}

\begin{minipage}[t]{0.48\textwidth}
\small\devanagarifont
व्यासार्ध कृतिं हृता विषुवत्कर्णेन या भवेत् कर्णः।\\
लम्बगुणो वा घातः शङ्कुकृतिसार्धकृतिभक्तः ॥३२॥\\[0.5em]
घातो वा शङ्कुगुणक्रियया विषुवत्कर्णेन विभुक्तः शङ्कुः।\\
कर्णकृतिः संशोध्य द्वादशभागे पदं छाया ॥३३॥
\end{minipage}%
\hfill
\begin{minipage}[t]{0.48\textwidth}
\small
\textbf{Verses 32--33. Alternative Shadow Computation:}

The square of the semi-diameter divided by the equinoctial hypotenuse gives the hypotenuse (kar\d{n}a). Or, the product of the lambaka (co-latitude) and the quantity, divided by the square-root of the sum of the squares of the gnomon and its semi-diameter.

Or the product by the gnomon-multiplication, divided by the equinoctial hypotenuse, gives the \'sa\~nku. The square of the hypotenuse, subtracted [from the appropriate quantity], its root divided by 12 gives the shadow.
\end{minipage}

\sectsep

%% ========================================================================
%% NATAKĀLĀNAYANAM
%% ========================================================================
\subsection{\dev{नतकालानयनम्}}

\begin{minipage}[t]{0.48\textwidth}
\small\devanagarifont
स्वाहोरात्रार्धेन छायाकर्णेन मतायाः।\\
विषुवत्कर्णगुणायामा व्यासार्धकृतिः फलं सौम्ये ॥३४॥\\[0.5em]
मयुखद्विज्याहीनं युक्तं याम्ये धनुष्करप्रमाणः।\\
सौम्ये युतं विहीनं याम्ये प्रागपरयोः प्राणाः ॥३५॥\\[0.5em]
भ्रष्टो गतविशेषाः फलमत्याया विशेष शेषस्य।\\
धनुःरविकर्मजीवार्धिः पूर्वापरयोर्नतप्राणाः ॥३६॥
\end{minipage}%
\hfill
\begin{minipage}[t]{0.48\textwidth}
\small
\textbf{Verses 34--36. Computation of Nata-k\=ala (Hour Angle):}

With the day-and-night half and the ch\=ay\=a-kar\d{n}a hypotenuse as multipliers, the equinoctial hypotenuse times the ko\d{t}i, divided by the square of the semi-diameter, gives the result in the north.

Diminished by the sun's diameter-sine in the north, added in the south --- this is the arc-measure. Added in the north, diminished in the south --- these are the east-west pr\=a\d{n}as.

The elapsed specific portions; the remainder of the excess portion. Half the day-sine is the east-west nata-pr\=a\d{n}a (hour-angle).
\end{minipage}

\sectsep

%% ========================================================================
%% LAMBHĀKṢAJYĀ
%% ========================================================================
\subsection{\dev{लम्बाक्षज्यानयनम्}}

\begin{minipage}[t]{0.48\textwidth}
\small\devanagarifont
शङ्कु लम्बकछायाख्यया तद्‌गुणस्युत्तरैं लब्धम्।\\
विशुद्धिर्विषुवत्कर्णाक्षेत्राक्षयाक्षोच्चनता शङ्कुः ॥५७॥\\[0.5em]
उन्नतजीवाकोटिज्याह्रिया हृज्या भुजो नतज्या वा।\\
कर्णाक्षेत्रावृत्ते व्यासार्धं द्व्यमतोद्यन्न ॥५८॥\\[0.5em]
शङ्कु क्षायाकृत्योरित्युक्त्याकृतितत्समासपुरुषष्टकयोः।\\
भूते लम्बाक्षज्ये तदंशकलास्तदनुमागाः ॥५९॥\\[0.5em]
विषुवत्कर्णहृते वा शङ्कु क्षायाहृते पृथक् त्रिज्ये।\\
अक्षक्षेत्रर्ज्ये लम्बाक्षज्योत्क्षमध्यने ॥६०॥\\[0.5em]
नतलम्बाक्षाक्षान् प्रोह्य ज्या वैतराश्लब्धकष्ट्ये।\\
शङ्कु क्षायागुणिक्ते क्षायद्विघातहृते वास्ते ॥६१॥\\[0.5em]
लम्बाक्षकलाभागं शेषं त्रिज्याकृतेः पदं भाज्याः।\\
भनुज्या सर्वोत्क्रान्तज्याविवराणांऽघनमेव ॥६२॥
\end{minipage}%
\hfill
\begin{minipage}[t]{0.48\textwidth}
\small
\textbf{Verses 57--62. Latitude-related Sine Computation (Lambh\=ak\d{s}a-jy\=a):}

The \'sa\~nku (gnomon), known as the lambaka-ch\=ay\=a, multiplied by the northward [quantity], gives the result. The equinoctial hypotenuse [gives] the purity; the highest latitude-\'sa\~nku is the aksa-ucca-nat\=a.

The elevated sine's ko\d{t}i-jy\=a, the altitude-sine is the bhuja or the nata-jy\=a. The semi-diameter [applies to] the kar\d{n}a-field-circle.

The \'sa\~nku and the k\d{s}aya quantities, with the aggregate of eight puru\d{s}as, become the lambh\=ak\d{s}a-jy\=a; its degrees and minutes are the subsequent [measures].

Divided by the equinoctial hypotenuse, or the \'sa\~nku divided by the k\d{s}aya separately with the trijy\=a, in the latitude field, the lambh\=ak\d{s}a-jy\=a is elevated.

Having subtracted the nata-lambh\=ak\d{s}a, the sines [yield] difficulty. The \'sa\~nku multiplied by k\d{s}aya, divided by twice the k\d{s}aya, remains.

The lambh\=ak\d{s}a degrees and minutes' remainder, the root of the square of the trijy\=a, is to be divided. The bhuja-jy\=a of all elevated differences is single.
\end{minipage}

\sectsep

%% ========================================================================
%% PRAKĀRĀNTAREṆA LAMBHĀKṢAJYĀ
%% ========================================================================
\subsection{\dev{पुनः प्रकारान्तरेण लम्बाक्षज्ये}}

\begin{minipage}[t]{0.48\textwidth}
\small\devanagarifont
लम्बाक्षज्ञानं नवतः प्रोक्ष (हृत्वा) ज्या साध्या तदा वा (प्रकारान्तरेण)\\
इनता ज्या स्वाध्यर्धल्पांशयोर्नवतेर्ज्योर्ज्याज्या, अक्षांशोन्नतेज्योर्ज्या लम्बज्या,\\
अक्षलम्बज्ये-शङ्कुच्छायागुणिक्ते (द्वादशपलमागुणिक्ते) छायाद्वादशहृते (पलमा-द्वादश भक्ते) तदा वा (प्रकारान्तरेण)इष्‍ठे लम्बाक्षज्ये भवत् इति ॥९१॥
\end{minipage}%
\hfill
\begin{minipage}[t]{0.48\textwidth}
\small
\textbf{Verse 91. Alternative Lambh\=ak\d{s}a-jy\=a Computation:}

For computing the lambh\=ak\d{s}a-jy\=a (latitude sine), having divided from 90 degrees, the sine is to be computed. Or alternatively: the jy\=a of the difference between 90 degrees and the latitude degrees is the lambaka-jy\=a. The aksa-lambaka-jy\=a [is obtained from] the product of the \'sa\~nku and ch\=ay\=a (gnomon and shadow), multiplied by 12 and divided by the shadow. Thus the lambh\=ak\=a-jy\=a is obtained by an alternative method.
\end{minipage}

\vspace{0.5em}
\begin{minipage}[t]{0.48\textwidth}
\small\devanagarifont
\textit{उपपत्तिः} --- ज्या (९०---अक्षांश) = लम्बज्या, ज्या (९०---अक्षांश) = मध्यज्या, वा
\[
\frac{\text{अक्षज्या} \times १२}{\text{पलमा}} = \text{लम्बज्या}, \quad \frac{\text{पलमा} \times \text{लम्बज्या}}{१२} = \text{मध्यज्या}
\]
एतावत्यक्षाज्या-लम्बोपपत्तेः ॥९१॥
\end{minipage}%
\hfill
\begin{minipage}[t]{0.48\textwidth}
\small
\textit{Derivation:}

The sine of (90 - latitude) = lambaka-jy\=a, the sine of (90 - latitude) = madhya-jy\=a. Or:
\[
\frac{\text{aks\=a-jy\=a} \times 12}{\text{palam\=a}} = \text{lambaka-jy\=a}, \quad \frac{\text{palam\=a} \times \text{lambaka-jy\=a}}{12} = \text{madhya-jy\=a}
\]
This is the derivation of the aksa and lambaka. \hfill\textbf{[91]}
\end{minipage}

\sectsep

%% ========================================================================
%% NATAKĀLĀNAYANAM (REPEATED)
%% ========================================================================
\subsection{\dev{नतकालानयनम् (पुनः)}}

\begin{minipage}[t]{0.48\textwidth}
\small\devanagarifont
उदयास्तमयैः क्षेपैर्द्वादशाभिर्नतं याम्योदयात्।\\
भागैर्लब्धं चलं तद् दोर्ज्याविवरं च स्फुटम् ॥६३॥\\[0.5em]
तन्मण्डलं प्रोक्ष्य तज्जं शङ्कुना लम्बकेन वा हतम्।\\
द्युज्यया विभक्तं लब्धं नतकालं प्रचक्षते ॥६४॥\\[0.5em]
अर्कक्षेपैर्विनता दोर्ज्योत्क्रान्तिज्या नतकालहृते।\\
तज्जं त्रिज्याहतं द्युज्यया विभक्तं शङ्कुरिष्यते ॥६५॥\\[0.5em]
नतोत्क्रान्तज्ययोर्घाते द्युज्याकृतेस्तदंशकं शङ्कुः।\\
छायाहते द्वादशभागे पदं भवति तत्क्षेपः ॥६६॥
\end{minipage}%
\hfill
\begin{minipage}[t]{0.48\textwidth}
\small
\textbf{Verses 63--66. Computation of Nata-k\=ala (Hour Angle):}

Having multiplied the uday\=astamaya (rising-setting) by the twelve a\d{k}\d{s}epa (declination), from the y\=amyodaya (southern rising), the result is the \textit{cala} and the true difference of the do\d{h}-jy\=a (altitude sine difference).

Proclaiming that ma\d{n}\d{d}ala, multiplying the result by the \'sa\~nku or lambaka, divided by the dyu-jy\=a (day-radius), gives the nata-k\=ala (hour angle).

The do\d{h}-jy\=a-utk\=anti-jy\=a (altitude sine minus declination sine), diminished by the arka-k\d{s}epa, divided by the nata-k\=ala, multiplied by the trijy\=a and divided by the dyu-jy\=a, gives the desired \'sa\~nku.

The product of the nata and utk\=anta-jy\=a, divided by the square of the dyu-jy\=a, that portion is the \'sa\~nku. The shadow multiplied by 12, its root divided [appropriately] gives the k\d{s}epa.
\end{minipage}

\vspace{0.5em}
\begin{minipage}[t]{0.48\textwidth}
\small\devanagarifont
\textit{वा.~भा.}---उदयास्तमयक्षेपैर्द्वादशाभिर्नतं याम्योदयात् भागैर्लब्धं चलं तद् दोर्ज्याविवरं-च स्फुटम्। तन्मण्डलं प्रोक्ष्य तज्जं शङ्कुना लम्बकेन वा हतं द्युज्यया विभक्तं लब्धं-नतकालं प्रचक्षते। अर्कक्षेपैर्विनता दोर्ज्योत्क्रान्तिज्या नतकालहृते तज्जं त्रिज्याहतं-द्युज्यया विभक्तं शङ्कुरिष्यते। नतोत्क्रान्तज्ययोर्घाते द्युज्याकृतेस्तदंशकं शङ्कुः-छायाहते द्वादशभागे पदं भवति तत्क्षेपः ॥६३--६६॥
\end{minipage}%
\hfill
\begin{minipage}[t]{0.48\textwidth}
\small
Having multiplied the uday\=astamaya by the twelve a\d{k}\d{s}epa from the y\=amyodaya, the result is the \textit{cala} and the true do\d{h}-jy\=a-vivara. Proclaiming that ma\d{n}\d{d}ala, multiplying the result by the \'sa\~nku or lambaka and dividing by the dyu-jy\=a, gives the nata-k\=ala. The do\d{h}-jy\=a-utk\=anti-jy\=a, diminished by the arka-k\d{s}epa and divided by the nata-k\=ala, multiplied by the trijy\=a and divided by the dyu-jy\=a, gives the \'sa\~nku. The product of the nata and utk\=anta-jy\=a, divided by the square of the dyu-jy\=a, gives the \'sa\~nku; the shadow multiplied by 12, its root gives the k\d{s}epa.
\end{minipage}

\sectsep

%% ========================================================================
%% GOLA-YANTRA
%% ========================================================================
\subsection{\dev{गोलयन्त्रम् --- खगोलयन्त्रनिर्माणम्}}

\begin{minipage}[t]{0.48\textwidth}
\small\devanagarifont
सप्तरश्मिः खगोलः स्यात् तन्मध्ये भूगोलसंस्थितिः\\
दण्डैः संयोजितः सोऽयं खगोलो यन्त्ररूपकः ॥१३७॥\\[0.5em]
क्रान्तिवृत्तं सममण्डलं दक्षिणोत्तरवृत्तमेव च\\
सद्यत्क्षितिजमेवाऽपि गोले निर्माय तद्विदैः ॥१३८॥\\[0.5em]
यन्त्रेणैव ग्रहाणां स्थितिः स्फुटा ज्ञायते येन\\
तस्माद् यन्त्रं प्रकल्प्यं गोलविद्भिर्निरन्तरम् ॥१३९॥
\end{minipage}%
\hfill
\begin{minipage}[t]{0.48\textwidth}
\small
\textbf{Verses 137--139. Construction of the Armillary Sphere (Gola-yantra):}

The celestial sphere (khagola) has seven rays [circles]; in its centre is situated the terrestrial sphere (bh\=ugola). Joined by rods (da\d{n}\d{d}ai\d{h}), this khagola is in the form of an instrument (yantra-r\=upaka).

The ecliptic (kr\=anti-v\=tta), the equator (sama-ma\d{n}\d{d}ala), the north-south circles (dak\d{s}i\d{n}ottara-v\=tta), and the prime vertical (sadyat-k\d{s}itija) should be constructed in the sphere by those who know.

By this instrument alone, the true position of the planets is known. Therefore, the instrument should always be constructed by those who know the sphere (gola-vid-bhi\d{h}).
\end{minipage}

\vspace{0.5em}
\begin{minipage}[t]{0.48\textwidth}
\small\devanagarifont
\textit{वा.~भा.}---सप्तरश्मिः खगोलः स्यात् तन्मध्ये भूगोलसंस्थितिर्दण्डैः संयोजितः सोऽयं खगोलो यन्त्ररूपकः। क्रान्तिवृत्तं सममण्डलं दक्षिणोत्तरवृत्तमेव च सद्यत्क्षितिजमेवापि गोले निर्माय तद्विदैः यन्त्रेणैव ग्रहाणां स्थितिः स्फुटा ज्ञायते येन तस्माद् यन्त्रं प्रकल्प्यं गोलविद्भिर्निरन्तरम् ॥१३७--१३९॥
\end{minipage}%
\hfill
\begin{minipage}[t]{0.48\textwidth}
\small
The celestial sphere with seven circles, having the terrestrial sphere at its centre, joined by rods, is this armillary sphere. The ecliptic, equator, meridian circles, and prime vertical should be constructed in the sphere by those versed in spherics. By this instrument alone, the true positions of the planets are known; therefore, astronomers should always construct this instrument.
\end{minipage}

\sectsep

%% ========================================================================
%% PRAṆAGAṆITAM
%% ========================================================================
\subsection{\dev{प्राणगणितम् --- दिगंशानयनम्}}

\begin{minipage}[t]{0.48\textwidth}
\small\devanagarifont
षष्टिर्घटिका नाडी षष्टिर्नाड्यः प्राणसञ्ज्ञकाः\\
षष्टिर्विपलानि पलान्येतैः कालगणनं विदुषाम् ॥१२८॥\\[0.5em]
चरज्यां हत्वा त्रिज्यया विभज्य लब्धाः प्राणा भवन्ति\\
तेषां योगः स्थितिकालः स्यादयनकालस्तु तद्वदेव ॥१२९॥\\[0.5em]
तनुज्या हृत्वा त्रिज्यया विभज्य लब्धं तत् प्रमाणकालः\\
स्थितौ चरज्यया हृत्वा लब्धाः प्राणाः स्फुटास्तु ते ॥१३०॥
\end{minipage}%
\hfill
\begin{minipage}[t]{0.48\textwidth}
\small
\textbf{Verses 128--130. Time Unit (Pr\=a\d{n}a) Computation:}

Sixty gha\d{t}ik\=as make one n\=a\d{d}\=i; sixty n\=a\d{d}\=is are called pr\=a\d{n}as; sixty vipalas make one pala --- by these units the learned compute time.

Having multiplied the cara-jy\=a (ascensional difference sine) and divided by the trijy\=a, the result gives the pr\=a\d{n}as. Their sum is the sthiti-k\=ala (duration of visibility); the ayanak\=ala is similarly computed.

Having divided the tanu-jy\=a (sine of zenith distance) and dividing by the trijy\=a, the result is the pram\=a\d{n}a-k\=ala (measure-time). In the sthiti (position), dividing by the cara-jy\=a, those pr\=a\d{n}as are the true ones.
\end{minipage}

\vspace{0.5em}
\begin{minipage}[t]{0.48\textwidth}
\small\devanagarifont
\textit{वा.~भा.}---षष्टिर्घटिका नाडी षष्टिर्नाड्यः प्राणसञ्ज्ञकाः षष्टिर्विपलानि पलानि एतैः कालगणनं विदुषाम्। चरज्यां हत्वा त्रिज्यया विभज्य लब्धाः प्राणा भवन्ति तेषां योगः स्थितिकालः स्यादयनकालस्तु तद्वदेव। तनुज्या हृत्वा त्रिज्यया विभज्य लब्धं तत् प्रमाणकालः स्थितौ चरज्यया हृत्वा लब्धाः प्राणाः स्फुटास्तु ते ॥१२८--१३०॥
\end{minipage}%
\hfill
\begin{minipage}[t]{0.48\textwidth}
\small
Sixty gha\d{t}ik\=as equal one n\=a\d{d}\=i; sixty n\=a\d{d}\=is are pr\=a\d{n}as; sixty vipalas equal one pala --- these are the time units used by the learned. Multiplying the cara-jy\=a by the trijy\=a and dividing gives the pr\=a\=d{n}as. Their sum is the sthiti-k\=ala; the ayanak\=ala is similarly computed. Dividing the tanu-jy\=a by the trijy\=a gives the pram\=a\d{n}a-k\=ala; in the position, dividing by the cara-jy\=a gives the true pr\=a\d{n}as.
\end{minipage}

\sectsep

%% ========================================================================
%% MADHYACCHĀYĀ
%% ========================================================================
\subsection{\dev{मध्यच्छायानयनम्}}

\begin{minipage}[t]{0.48\textwidth}
\small\devanagarifont
छायाक्षेत्रानुपातेन त्रि.१२/शङ्कु = छायाकर्णं ततः $\sqrt{\text{छायाक}^2 - १२^2}$ = छाया,\\
इति ॥१४८॥
\end{minipage}%
\hfill
\begin{minipage}[t]{0.48\textwidth}
\small
\textbf{Verse 148. Computation of Midday Shadow:}

By the proportion of the ch\=ay\=a-field:
\[
\frac{\text{tri} \times 12}{\text{\'sa\~nku}} = \text{ch\=ay\=a-kar\d{n}a}
\]
Then:
\[
\sqrt{\text{(ch\=ay\=a-kar\d{n}a)}^2 - 12^2} = \text{midday shadow}
\]
\hfill \textbf{[148]}
\end{minipage}

\vspace{0.5em}
\begin{minipage}[t]{0.48\textwidth}
\small\devanagarifont
\textit{हिं.~भा.}---त्रिज्या को बारह से गुणा कर मध्यशङ्कु से भाग देने से मध्यच्छाया कर्णं होता है, छायाकर्ण और द्वादश (१२) का---वर्गन्तर मूल प्रकारान्तर से मध्यच्छाया होती है इति ॥१४८॥
\end{minipage}%
\hfill
\begin{minipage}[t]{0.48\textwidth}
\small
\textit{[Hindi commentary:]} Multiplying the trijy\=a by 12 and dividing by the midday \'sa\~nku gives the midday ch\=ay\=a-kar\=d{n}a. The square-root of the difference of squares of the ch\=ay\=a-kar\d{n}a and 12 gives the midday shadow by an alternative method. \hfill\textbf{[148]}
\end{minipage}

\vspace{0.5em}
\begin{minipage}[t]{0.48\textwidth}
\small\devanagarifont
\textit{उपपत्तिः}---पूर्व श्लोक की उपपत्ति में प्रदर्शित छायाक्षेत्रों के सजातीयत्व से अनुपात करते हैं
\[
\frac{\text{त्रि.} \times १२}{\text{मध्यशङ्कु}} = \text{मध्याकर्ण}
\]
$\therefore \sqrt{\text{छायाक}^2 - १२^2}$ = मध्यच्छाया, इति ॥१४८॥
\end{minipage}%
\hfill
\begin{minipage}[t]{0.48\textwidth}
\small
\textit{Derivation:} From the similarity of the ch\=ay\=a-fields demonstrated in the previous verse's derivation, by proportion:
\[
\frac{\text{tri} \times 12}{\text{midday \'sa\~nku}} = \text{midday kar\d{n}a}
\]
Therefore $\sqrt{\text{(ch\=ay\=a-kar\d{n}a)}^2 - 12^2}$ = midday shadow. \hfill\textbf{[148]}
\end{minipage}

\sectsep

%% ========================================================================
%% MAHĀPĀTA
%% ========================================================================
\subsection{\dev{महापातः --- चान्द्रमासनिर्माणम्}}

\begin{minipage}[t]{0.48\textwidth}
\small\devanagarifont
महापातस्तु यः स्याच्चन्द्रसूर्ययोः समागमे\\
तदा तयोरन्तरं शून्यं स्यात् समेclipticे तयोः ॥१४०॥\\[0.5em]
चान्द्रो मासस्तु यः स्यात् तिथीनां सङ्ग्रहेण तु\\
स तु सावनोऽपि मासः सूर्येणैकेन गम्यते ॥१४१॥\\[0.5em]
तयोर्मासयोगेन संवत्सरः प्रकीर्तितः\\
पञ्चदशभ्यस्तिथिभिः संयुक्तो मास उच्यते ॥१४२॥
\end{minipage}%
\hfill
\begin{minipage}[t]{0.48\textwidth}
\small
\textbf{Verses 140--142. The Mah\=ap\=ata and Lunar Month:}

The mah\=ap\=ata (great conjunction) occurs at the meeting of the Moon and Sun; then their difference becomes zero when they are at the same ecliptic point.

The lunar month (c\=andra-m\=asa), by the aggregation of tithis (lunar days), is also a \=s\=avana month, traversed by the Sun alone.

The combination of these two months is called a sa\d{mvatsara (year). A month, joined with fifteen tithis, is thus called.}
\end{minipage}

\vspace{0.5em}
\begin{minipage}[t]{0.48\textwidth}
\small\devanagarifont
\textit{वा.~भा.}---महापातस्तु यः स्याच्चन्द्रसूर्ययोः समागमे तदा तयोरन्तरं शून्यं स्यात् समेclipticे तयोः। चान्द्रो मासस्तु यः स्यात् तिथीनां सङ्ग्रहेण तु स तु सावनोऽपि मासः सूर्येणैकेन गम्यते तयोर्मासयोगेन संवत्सरः प्रकीर्तितः पञ्चदशभ्यस्तिथिभिः संयुक्तो मास उच्यते ॥१४०--१४२॥
\end{minipage}%
\hfill
\begin{minipage}[t]{0.48\textwidth}
\small
The mah\=ap\=ata (great conjunction) occurs when the Moon and Sun meet; then their difference is zero when they are at the same ecliptic longitude. The lunar month, by the collection of tithis, is also a \=s\=avana month, traversed by the Sun. The year (sa\d{mvatsara) is proclaimed by the combination of these two months. A month joined with fifteen tithis is called [a fortnight].}
\end{minipage}

\sectsep

%% ========================================================================
%% KRĀNTI AND AYANĀṂŚA
%% ========================================================================
\subsection{\dev{क्रान्तिः --- अयनांशकल्पनम्}}

\begin{minipage}[t]{0.48\textwidth}
\small\devanagarifont
मेषादौ रवौ तु षट्षष्टिघ्ने रवौ तु विषुवति\\
मेषादौ तत्र क्रान्तिर्ज्या तु विषुवज्ज्या मता ॥१३१॥\\[0.5em]
तत्क्षेपैर्विहृता ज्ञेया क्रान्तिः सा स्फुटतर्किभिः\\
उत्तरायणे दक्षिणायने याम्या सा तु तत्क्षेपैः ॥१३२॥\\[0.5em]
अयनांशः स तु यतः क्रान्तेर्भवति विशेषणात्\\
तस्मादयनचलनं ज्ञात्वा स्फुटं विधिवद् भवेत् ॥१३३॥
\end{minipage}%
\hfill
\begin{minipage}[t]{0.48\textwidth}
\small
\textbf{Verses 131--133. Declination (Kr\=anti) and Ayan\=a\d{m}\'sa:}

When the Sun is at the beginning of Me\d{s}a (Aries), multiply by 66 [the k\=ala value]. When the Sun is at the equinox (vi\d{s}uvat), at the beginning of Me\d{s}a, the declination-sine (kr\=anti-jy\=a) there is considered as the equinoctial sine (vi\d{s}uva-jjy\=a).

That, diminished by its k\d{s}epa, is to be known as the true declination by those who compute. In the uttar\=aya\d{n}a (northward course), it is southern; by those k\d{s}epas, it becomes northern.

The ayan\=a\d{m}\'sa (precession of equinoxes) is that from which the declination arises by differentiation. Therefore, knowing the ayanacalana (precessional motion), the true [position] should be computed by the prescribed method.
\end{minipage}

\vspace{0.5em}
\begin{minipage}[t]{0.48\textwidth}
\small\devanagarifont
\textit{वा.~भा.}---मेषादौ रवौ तु षट्षष्टिघ्ने रवौ तु विषुवति मेषादौ तत्र क्रान्तिः स्यात् तत्क्षेपैर्विहृता ज्ञेया क्रान्तिः सा स्फुटतर्किभिः उत्तरायणे दक्षिणायने याम्या सा तु तत्क्षेपैः। अयनांशः स तु यतः क्रान्तेर्भवति विशेषणात् तस्मादयनचलनं ज्ञात्वा स्फुटं विधिवद् भवेत् ॥१३१--१३३॥
\end{minipage}%
\hfill
\begin{minipage}[t]{0.48\textwidth}
\small
When the Sun is at the beginning of Aries, multiplied by 66; at the equinox, the declination at the beginning of Aries is considered the equinoctial sine. That, diminished by its k\d{s}epa, is the true declination. In the northward course, the declination is southern; in the southward, northern. The ayan\=a\d{m}\'sa is that from which the declination differs; knowing this precessional motion, the true [position] is computed properly.
\end{minipage}

\sectsep

%% ========================================================================
%% KṢITIJYĀ
%% ========================================================================
\subsection{\dev{क्षितिज्या}}

\begin{minipage}[t]{0.48\textwidth}
\small\devanagarifont
प्रसिद्धा। याम्ये याम्यभुजोन्नतोर्ज्वं लन्तं रमद्‌गुप्तरमुजेनगोरेक्यं\\
यतल्लम्बगुणां छायाकर्णोद्धृतं फलं किन्तः क्रान्तिज्या स्यादिति।
\end{minipage}%
\hfill
\begin{minipage}[t]{0.48\textwidth}
\small
\textbf{K\d{s}itijy\=a (Terrestrial Sine):}

Well-known: In the southern [hemisphere], the southern bhuja's elevation, the product of the lambaka and the appropriate quantity, divided by the ch\=ay\=a-kar\d{n}a, gives the result. From this, the kr\=anti-jy\=a (declination-sine) is obtained.
\end{minipage}

\sectsep

%% ========================================================================
%% LAGNANAYANAM
%% ========================================================================
\subsection{\dev{लग्ननयनम् --- उदयराशिसाधनम्}}

\begin{minipage}[t]{0.48\textwidth}
\small\devanagarifont
रवेस्तु राश्युदयाद् घटिकाः प्रोद्यमानराशिस्थितेः\\
तास्तत्कालात् शोधयित्वा यावद् राश्युदयाः स्युः ॥१३४॥\\[0.5em]
तास्तत्पूर्वराशिभिः संयुक्ता लग्नं तु तत् स्यात्\\
शेषं त्रिंशता हृत्वा लब्धं तदंशकलात्मकम् ॥१३५॥\\[0.5em]
तद् राशौ योजयेत् तस्य लग्नं स्यात् स्फुटतर्किभिः\\
एवं लग्नं सदा कार्यं ग्रहगणितविदा बुधैः ॥१३६॥
\end{minipage}%
\hfill
\begin{minipage}[t]{0.48\textwidth}
\small
\textbf{Verses 134--136. Computation of the Lagna (Rising Sign):}

From the rising-time (r\=a\'si-udaya) of the Sun's sign, the gha\d{t}ik\=as of the currently rising sign, having subtracted those from the given time, until the sign-risings are completed.

Those [signs], joined with the previous signs, constitute the lagna. Dividing the remainder by 30, the result is in degrees and minutes.

That should be added to the sign; its lagna is thus computed by those who compute true positions. Thus the lagna should always be computed by the wise who know planetary mathematics.
\end{minipage}

\vspace{0.5em}
\begin{minipage}[t]{0.48\textwidth}
\small\devanagarifont
\textit{वा.~भा.}---रवेस्तु राश्युदयाद् घटिकाः प्रोद्यमानराशिस्थितेः तास्तत्कालात् शोधयित्वा यावद् राश्युदयाः स्युः तास्तत्पूर्वराशिभिः संयुक्ता लग्नं तु तत् स्यात् शेषं त्रिंशता हृत्वा लब्धं तदंशकलात्मकं तद् राशौ योजयेत् तस्य लग्नं स्यात् स्फुटतर्किभिः एवं लग्नं सदा कार्यं ग्रहगणितविदा बुधैः ॥१३४--१३६॥
\end{minipage}%
\hfill
\begin{minipage}[t]{0.48\textwidth}
\small
From the rising-time of the Sun's sign, the gha\d{t}ik\=as of the currently rising sign, subtracted from the given time until the sign-risings are accounted for. Those signs, joined with the previous ones, give the lagna. Dividing the remainder by 30 gives the degrees and minutes, which are added to the sign; this gives the lagna for true-position computation. Thus the lagna is always to be computed by wise planetary calculators.
\end{minipage}

\sectsep

%% ========================================================================
%% MATHEMATICAL FORMULAS AND DERIVATIONS
%% ========================================================================
\section{\dev{गणितीय सूत्राणि --- Mathematical Formulas and Derivations}}
\addcontentsline{toc}{section}{Mathematical Formulas and Derivations}

\begin{center}
\textit{Key Mathematical Formulas from the Spa\d{s}\d{t}\=adhik\=ara and Tripra\'sn\=adhik\=ara}
\end{center}

\begin{minipage}[t]{0.48\textwidth}
\small\devanagarifont
अस्य विभागस्य विषयः --- स्पष्टाधिकारात् त्रिप्रश्नाधिकारात् च प्रमुखाणि गणितीयसूत्राणि यानि ज्यापटलनिर्माणे, छायागणने, नतकालानयने, तथा पाटीगणिते प्रयुज्यन्ते।
\end{minipage}%
\hfill
\begin{minipage}[t]{0.48\textwidth}
\small
This section presents the key mathematical formulas from the Spa\d{s}\d{t}\=adhik\=ara and Tripra\'sn\=adhik\=ara chapters, covering sine table construction, shadow calculations, hour angle computation, and arithmetic operations.
\end{minipage}

\sectsep

%% ========================================================================
%% ŚĪGHRA-SPAṢṬIKARAṆASŪTRAM
%% ========================================================================
\subsection{\dev{शीघ्रस्पष्टीकरणसूत्रम्}}

\begin{minipage}[t]{0.48\textwidth}
\small\devanagarifont
शीघ्रफलोज्या $\times$ शीघ्रकर्ण = स्पष्टः\\[0.5em]
केन्‍द्रकर्णः = $\sqrt{\text{त्रिज्या}^2 + \text{मन्दफलोज्या}^2 - 2 \cdot \text{मन्दफलोज्या} \times \text{कल्पज्या}}$
\end{minipage}%
\hfill
\begin{minipage}[t]{0.48\textwidth}
\small
\textbf{\'S\=ighra-spa\d{s}\d{t}ikara\d{n}a Formula:}

The \'s\=ighra-phala-jy\=a (anomaly equation sine) multiplied by the \'s\=ighra-kar\d{n}a gives the true position correction.

The kendra-kar\=d{n}a (eccentric distance) is computed as the square-root of the sum of the square of the trijy\=a and the square of the manda-phala-jy\=a, minus twice the product of the manda-phala-jy\=a and the kalpa-jy\=a.
\end{minipage}

\sectsep

%% ========================================================================
%% JYOTPATHI
%% ========================================================================
\subsection{\dev{ज्योत्पत्तिः --- त्रिकोणमितीय ज्यापटलनिर्माणम्}}
\label{subsec:jyotpatti}

\begin{minipage}[t]{0.48\textwidth}
\textbf{\small Sanskrit Verses}

\small\devanagarifont
व्यासाद् व्यासभागानां त्रिभिरूनं शैलैस्त्रिभिरूनान्\\
तान् कृत्वा षड्भिर्भक्तान् शेषैस्त्रिज्यां प्रकल्पयेत् ॥१००॥\\[0.5em]
पूर्वपूर्वज्यया हत्वा पूर्वपूर्वविशेषखण्डान्\\
त्रिज्यां हृत्वा ततोऽन्या ज्याः सन्ति तत्प्रभवा मता ॥१०१॥
\end{minipage}%
\hfill
\begin{minipage}[t]{0.48\textwidth}
\textbf{\small English Translation}

\textbf{Verses 100--101. Construction of the Sine Table (Jyotpatti):}

From the diameter (vy\=asa), having subtracted three from the diameter-parts [3438--3 = 3435], made into parts and divided by six, the remainder is taken as the trijy\=a (radius sine, $R \approx 3438'$).

Having multiplied the previous previous-jy\=a by the previous previous-difference-segments, and dividing by the trijy\=a, the other jy\=as are obtained; they are considered to originate from the previous one.
\end{minipage}

\vspace{0.5em}
\begin{minipage}[t]{0.48\textwidth}
\small\devanagarifont
\textit{वा.~भा.}---व्यासाद् व्यासभागानां त्रिभिरूनान् अर्थात् ३४३८--३ = ३४३५ भागान् कृत्वा षड्भिर्भक्तान् शेषैस्त्रिज्यां प्रकल्पयेत्। अनया त्रिज्यया पूर्वपूर्वज्यया पूर्वपूर्वविशेषखण्डान् हत्वा त्रिज्यां हृत्वा ततोऽन्या ज्याः सन्ति ताः पूर्वज्याप्रभवा मताः ॥१००--१०१॥
\end{minipage}%
\hfill
\begin{minipage}[t]{0.48\textwidth}
\small
From the diameter, subtracting three from the diameter-parts gives 3435 parts. Dividing by six, the remainder is taken as the trijy\=a. With this trijy\=a, multiplying the previous previous-jy\=a by the previous difference-segments and dividing by the trijy\=a, the other jy\=as are obtained; they are considered to originate from the previous jy\=a.
\end{minipage}

\vspace{0.5em}
\begin{minipage}[t]{0.48\textwidth}
\textbf{\small Sanskrit Verses}

\small\devanagarifont
प्रथमा ज्या त्रिज्याखण्डा द्वितीया ततः शिष्टस्यार्धात्\\
तृतीयापि तथैव सा सूत्रं ब्रह्मगुप्तमतम् ॥१०२॥\\[0.5em]
तासां विशेषखण्डानां योगे ज्यावलयः स्मृताः\\
षडंशाभ्यन्तरज्यास्ता विख्याताः स्पष्टतर्किकैः ॥१३३॥
\end{minipage}%
\hfill
\begin{minipage}[t]{0.48\textwidth}
\textbf{\small English Translation}

\textbf{Verses 102--103. Sine Computation Rule:}

The first jy\=a is a trijy\=a-segment; the second from half of the remainder; the third also likewise --- this is the rule of Brahmagupta's method.

The sum of their difference-segments is remembered as the jy\=a-\=avalaya (sine-circle). The six-degree interval jy\=as are well-known among accurate calculators.
\end{minipage}

\vspace{0.5em}
\begin{minipage}[t]{0.48\textwidth}
\small\devanagarifont
\textit{वा.~भा.}---प्रथमा ज्या सप्तभुजज्या त्रिज्याखण्डा द्वितीया चतुर्दशभुजज्या तृतीया तथैव क्रमेण। तासां विशेषखण्डानां योगे ज्यावलयः स्मृताः षडंशाभ्यन्तरज्यास्ता विख्याताः स्पष्टतर्किकैः ॥१०२--१०३॥
\end{minipage}%
\hfill
\begin{minipage}[t]{0.48\textwidth}
\small
The first jy\=a is the 7.5-degree sine, a trijy\=a-segment; the second is the 15-degree sine; the third similarly in sequence. The sum of their difference-segments is the jy\=a-\=avalaya. The six-degree interval sines are well-known among accurate calculators.
\end{minipage}

\sectsep

%% ========================================================================
%% CHĀYĀKARAṆĪ FORMULAS
%% ========================================================================
\subsection{\dev{छायाकरणीसूत्राणि}}

\begin{minipage}[t]{0.48\textwidth}
\small\devanagarifont
छाया = $\dfrac{\text{द्वादश} \times \text{हृज्या}}{\text{शङ्कु}}$\\[0.8em]
कर्णः = $\dfrac{\text{व्यासार्धकृतिः}}{\text{विषुवत्कर्ण}}$\\[0.8em]
छायाकर्ण = $\sqrt{\text{छाया}^2 + \text{शङ्कु}^2}$
\end{minipage}%
\hfill
\begin{minipage}[t]{0.48\textwidth}
\small
\textbf{Ch\=ay\=akara\d{n}\=i Formulas:}

\textbf{Shadow:} The shadow equals 12 (gnomon height) multiplied by the altitude-sine, divided by the \'sa\~nku (gnomon's effective height).

\textbf{Hypotenuse:} The hypotenuse equals the square of the semi-diameter divided by the equinoctial hypotenuse.

\textbf{Shadow-hypotenuse:} The hypotenuse of gnomon and shadow equals the square-root of the sum of squares of the shadow and the gnomon.
\end{minipage}

\sectsep

%% ========================================================================
%% GAURĪGUPTĀ
%% ========================================================================
\subsection{\dev{गौरीगुप्ता --- द्वितीयान्तरकल्पना}}
\label{subsec:gaurigupta}

\begin{minipage}[t]{0.48\textwidth}
\textbf{\small Sanskrit Verses}

\small\devanagarifont
ज्यावर्गैर्भुजज्यावर्गैस्त्रिज्यावर्गहृतैर्विशोध्य\\
लब्धैस्तैः प्रथमज्यां शोध्यैस्ततोऽपरा ज्या स्यात् ॥१०४॥\\[0.5em]
पूर्वोत्तरज्योरन्तरं प्रथमखण्डं ततोऽपरं खण्डम्\\
तस्मात् पूर्वखण्डाच्छोध्यं तृतीयं खण्डं ततः स्यात् ॥१०५॥\\[0.5em]
एवं क्रमेण सर्वा ज्याः साधनीया विभक्तज्यातः\\
अन्तरज्याः पृथग्वा स्फुटतरं फलमानयेत् ॥१०६॥
\end{minipage}%
\hfill
\begin{minipage}[t]{0.48\textwidth}
\textbf{\small English Translation}

\textbf{Verses 104--106. Second-Difference Interpolation (Gaur\=i-gupt\=a):}

Having subtracted the square of the bhuja-jy\=a (sine of arc) from the square of the jy\=a, divided by the square of the trijy\=a, and having subtracted the result from the first jy\=a, the next jy\=a is obtained.

The difference between the previous and subsequent jy\=a is the first segment (kha\d{n}\d{d}a); from that, the next segment; from that, having subtracted the previous segment, the third segment is obtained.

Thus, in sequence, all jy\=as should be computed from the divided jy\=a. The difference-jy\=as separately yield a more accurate result.
\end{minipage}

\vspace{0.5em}
\begin{minipage}[t]{0.48\textwidth}
\textbf{\small Sanskrit Verses}

\small\devanagarifont
विषमज्ञे द्विज्यावर्गाद् द्विज्यावर्गं विशोध्य शेषात्\\
पदं द्वितीयं खण्डं ततः पूर्वखण्डात् शोध्यम् ॥१०७॥\\[0.5em]
तृतीयं खण्डं ततः पूर्वद्वयात् शोध्यं चतुर्थं ततः\\
पूर्वत्रयात् शोध्यं च क्रमेण सर्वखण्डानि ॥१०८॥
\end{minipage}%
\hfill
\begin{minipage}[t]{0.48\textwidth}
\textbf{\small English Translation}

\textbf{Verses 107--108. Difference Segments for Sine Computation:}

For odd [positions], having subtracted the square of the dvi-jy\=a (double sine) from the square of the dvi-jy\=a, the root from the remainder is the second segment. Then, having subtracted the previous segment, [the third is obtained].

The third segment, then having subtracted from the previous two, the fourth; then having subtracted from the previous three, thus in sequence all segments are obtained.
\end{minipage}

\vspace{0.5em}
\begin{minipage}[t]{0.48\textwidth}
\small\devanagarifont
\textit{वा.~भा.}---ज्यावर्गैर्भुजज्यावर्गैस्त्रिज्यावर्गहृतैर्विशोध्य लब्धैः प्रथमज्यां शोध्य ततोऽपरा ज्या स्यात्। पूर्वोत्तरज्योरन्तरं प्रथमखण्डं ततोऽपरं खण्डं तस्मात् पूर्वखण्डाच्छोध्यं तृतीयं खण्डं ततः स्यात्। एवं क्रमेण सर्वा ज्याः साधनीया विभक्तज्यातः अन्तरज्याः पृथक् स्फुटतरं फलमानयेत् ॥१०४--१०६॥
\end{minipage}%
\hfill
\begin{minipage}[t]{0.48\textwidth}
\small
Having subtracted the square of the bhuja-jy\=a from the square of the jy\=a, divided by the square of the trijy\=a, and subtracting the result from the first jy\=a, the next jy\=a is obtained. The difference between previous and subsequent jy\=a is the first segment; the next segment from that; having subtracted the previous segment from that, the third segment is obtained. Thus, in sequence, all jy\=as are computed from the divided jy\=a; the difference-jy\=as separately give a more accurate result.
\end{minipage}

\vspace{0.5em}
\begin{minipage}[t]{0.48\textwidth}
\small\devanagarifont
\textit{वा.~भा.}---विषमज्ञे द्विज्यावर्गात् द्विज्यावर्गं विशोध्य शेषात् पदं द्वितीयं खण्डं स्यात्। ततः पूर्वखण्डात् शोध्यं तृतीयं खण्डं ततः पूर्वद्वयात् शोध्यं चतुर्थं ततः पूर्वत्रयात् शोध्यं च क्रमेण सर्वखण्डानि सिध्यन्ति ॥१०७--१०८॥
\end{minipage}%
\hfill
\begin{minipage}[t]{0.48\textwidth}
\small
For odd positions, having subtracted the square of the dvi-jy\=a from the square of the dvi-jy\=a, the root of the remainder is the second segment. Then, having subtracted the previous segment, the third; having subtracted from the previous two, the fourth; from the previous three, the fifth; thus in sequence all segments are obtained.
\end{minipage}

\sectsep

%% ========================================================================
%% LAMBHĀKṢAJYĀ FORMULAS
%% ========================================================================
\subsection{\dev{लम्बाक्षज्यासूत्राणि}}

\begin{minipage}[t]{0.48\textwidth}
\small\devanagarifont
लम्बज्या = $\dfrac{\text{अक्षज्या} \times १२}{\text{पलमा}}$\\[0.8em]
मध्यज्या = $\dfrac{\text{पलमा} \times \text{लम्बज्या}}{१२}$
\end{minipage}%
\hfill
\begin{minipage}[t]{0.48\textwidth}
\small
\textbf{Lambh\=ak\d{s}a-jy\=a Formulas:}

\textbf{Lambaka-jy\=a (co-latitude sine):} The aksa-jy\=a (latitude sine) multiplied by 12, divided by the palam\=a (equinoctial shadow), gives the lambaka-jy\=a.

\textbf{Madhya-jy\=a (meridian sine):} The palam\=a multiplied by the lambaka-jy\=a, divided by 12, gives the madhya-jy\=a.
\end{minipage}

\sectsep

%% ========================================================================
%% NATAKĀLĀNAYANASŪTRAM
%% ========================================================================
\subsection{\dev{नतकालानयनसूत्राणि}}

\begin{minipage}[t]{0.48\textwidth}
\small\devanagarifont
नतकालः = $\dfrac{\text{तज्ज्यं} \times \text{शङ्कु or लम्बक}}{\text{द्युज्या}}$\\[0.8em]
शङ्कुः = $\dfrac{\text{तज्ज्यं} \times \text{त्रिज्या}}{\text{द्युज्या}}$
\end{minipage}%
\hfill
\begin{minipage}[t]{0.48\textwidth}
\small
\textbf{Nata-k\=ala Formulas:}

\textbf{Hour angle:} The nata-k\=ala equals the taj-jya multiplied by the \'sa\~nku or lambaka, divided by the dyu-jy\=a (day-radius).

\textbf{\'Sa\~nku:} The \'sa\~nku equals the taj-jya multiplied by the trijy\=a, divided by the dyu-jy\=a.
\end{minipage}

\sectsep

%% ========================================================================
%% KRĀNTIJYĀNAYANASŪTRAM
%% ========================================================================
\subsection{\dev{क्रान्तिज्यानयनसूत्राणि}}

\begin{minipage}[t]{0.48\textwidth}
\small\devanagarifont
क्रान्तिज्या = $\dfrac{\text{लम्बज्या} \times \text{क्रम}}{\text{वि}} = \dfrac{\text{लम्बज्या} \times \text{कर्णोच्चताम्} \times \text{वि}}{\text{छाक} \cdot \text{वि}}$\\[0.8em]
क्षेत्रज्या = $\dfrac{\text{क्रान्तिज्या} \times \text{वि}}{\text{छाक}}$
\end{minipage}%
\hfill
\begin{minipage}[t]{0.48\textwidth}
\small
\textbf{Kr\=anti-jy\=a (Declination Sine) Formulas:}

The kr\=anti-jy\=a equals the lambaka-jy\=a multiplied by the krama and divided by the vi\d{s}uva. Alternatively, it equals the lambaka-jy\=a multiplied by the kar\d{n}occ\=at\=a and vi, divided by the ch\=ay\=a-kar\d{n}a and vi.

The k\d{s}etra-jy\=a equals the kr\=anti-jy\=a multiplied by the vi, divided by the ch\=ay\=a-kar\=d{n}a.
\end{minipage}

\sectsep

%% ========================================================================
%% PĀṬĪGAṆITAM
%% ========================================================================
\subsection{\dev{पाटीगणितम् --- चतुर्विंशतिपरिकर्माणि}}
\label{subsec:patiganita}

\begin{minipage}[t]{0.48\textwidth}
\textbf{\small Sanskrit Verses}

\small\devanagarifont
संकलितं व्यवकलितं गुणनं भाजनं तथा\\
वर्गो वर्गमूलं घनो घनमूलं पञ्चमी त्रयः ॥११७॥\\[0.5em]
पञ्च विशेषाः परिकर्माण्येतानि गणितं स्मृतम्\\
प्रकीर्तितान्याचार्यैरेतैः सर्वं गणितं सिध्यति ॥११८॥\\[0.5em]
रूपैकदेशसंयुक्ता व्यक्तराशिरुदाहृता सा\\
अव्यक्ता या राशिर्युक्ता सा रूपैकदेशेन ॥११९॥
\end{minipage}%
\hfill
\begin{minipage}[t]{0.48\textwidth}
\textbf{\small English Translation}

\textbf{Verses 117--119. The Twenty-Four Operations of P\=a\d{t}\=iga\d{n}ita:}

Addition (sa\d{m}kalita), subtraction (vyavakalita), multiplication (gu\d{n}ana), division (bh\=ajana), squaring (varga), square-root (varga-m\=ula), cubing (ghana), and cube-root (ghana-m\=ula) --- these five and three [make eight basic operations].

Five special [operations make] thirteen [in total]; these operations are remembered as ga\d{n}ita (mathematics). Proclaimed by the \=ac\=aryas, by these all mathematics is accomplished.

A number joined with a unit part is called a vyakta-r\=a\'si (manifest number). That which is a number joined with the unit part is [called] avyakta (unmanifest).
\end{minipage}

\vspace{0.5em}
\begin{minipage}[t]{0.48\textwidth}
\small\devanagarifont
\textit{वा.~भा.}---संकलितं योगः व्यवकलितं शोधनं गुणनं गुणकारः भाजनं हरणं वर्गः समतुल्यद्वयं वर्गमूलं समतुल्यप्रमाणम् घनः सममूलकस्त्रयः घनमूलं तत्प्रमाणं। एतानि परिकर्माणि गणितं स्मृतम् प्रकीर्तितान्याचार्यैः एतैः सर्वं गणितं सिध्यति ॥११७--११९॥
\end{minipage}%
\hfill
\begin{minipage}[t]{0.48\textwidth}
\small
Sa\d{m}kalita is addition; vyavakalita is subtraction; gu\d{n}ana is multiplication; bh\=ajana is division; varga is the product of two equals; varga-m\=ula is the root of that equal quantity; ghana is the product of three equals; ghana-m\=ula is the root of that quantity. These operations are called ga\d{n}ita (mathematics), proclaimed by the teachers; by these all mathematics is accomplished.
\end{minipage}

\vspace{0.5em}
\begin{minipage}[t]{0.48\textwidth}
\textbf{\small Sanskrit Verses}

\small\devanagarifont
यो राशेस्त्रिगुणो भागो रूपैकदेशसंयुक्तो वा\\
अर्धितः स वर्गः स्याद् यो राशेस्त्रिगुणो भागः ॥१२०॥\\[0.5em]
रूपैकदेशसंयुक्तोऽथवा स घनः स्यात् तस्य मूलं\\
यत् तद् घनमूलं प्रोक्तं तत् स्यात् सममूलकप्रमाणम् ॥१२१॥
\end{minipage}%
\hfill
\begin{minipage}[t]{0.48\textwidth}
\textbf{\small English Translation}

\textbf{Verses 120--121. Square, Cube, and Their Roots:}

That which is the triple portion of a number, or joined with a unit part and halved, that becomes the square (varga). That which is the triple portion of a number.

Or joined with a unit part, that becomes the cube (ghana). Its root is called the ghana-m\=ula (cube-root); that becomes the measure of the equal-root quantity.
\end{minipage}

\vspace{0.5em}
\begin{minipage}[t]{0.48\textwidth}
\small\devanagarifont
\textit{वा.~भा.}---यो राशेस्त्रिगुणो भागो रूपैकदेशसंयुक्तो वा अर्धितः स वर्गः स्यात्। यो राशेस्त्रिगुणो भागो रूपैकदेशसंयुक्तोऽथवा स घनः स्यात् तस्य मूलं यत् तद् घनमूलं प्रोक्तं तत् सममूलकप्रमाणम् स्यात् ॥१२०--१२१॥
\end{minipage}%
\hfill
\begin{minipage}[t]{0.48\textwidth}
\small
That which is the triple portion of a number, or joined with a unit part and halved, is the square. That which is the triple portion of a number, or joined with a unit part, is the cube. That which is its root is called the cube-root; that is the measure of the equal-root quantity.
\end{minipage}

\sectsep

%% ========================================================================
%% LUNI-SOLAR ASTRONOMY
%% ========================================================================
\subsection{\dev{चन्द्रसूर्यगणितम् --- तिथियोगनक्षत्रानयनम्}}
\label{subsec:luni-solar}

\begin{minipage}[t]{0.48\textwidth}
\textbf{\small Sanskrit Verses}

\small\devanagarifont
चन्द्राच्च स्पष्टात् सूर्यं शोध्यं यद् भवति शेषं तत्\\
द्वादशभागैर्हृत्वा तिथयः स्युः स्फुटास्तदर्धात् ॥१०९॥\\[0.5em]
तिथयो द्वादशाङ्गुल्यः सूर्यचन्द्रयोरन्तरात् ताः\\
तिथीनां पञ्चदशभिः संयुक्तं योगनिर्मितिः ॥११०॥\\[0.5em]
योगः सूर्यचन्द्रयो राशिसंयुक्त्यंशकलात्मिका यः\\
स तैः पञ्चदशभिः संयुक्तः सप्तविंशत्या हृतः ॥१११॥
\end{minipage}%
\hfill
\begin{minipage}[t]{0.48\textwidth}
\textbf{\small English Translation}

\textbf{Verses 109--111. Tithi, Yoga, and Nak\d{s}atra Computation:}

Having subtracted the Sun from the true Moon, the remainder that results, divided by twelve parts, gives the tithis (lunar days); from half of that, the true [tithis].

The tithis are twelve digits (a\~n\'gulis) from the difference of the Sun and Moon. Joined with fifteen tithis, the yoga is formed.

The yoga is the sum of the signs, degrees, and minutes of the Sun and Moon. That, joined with fifteen and divided by twenty-seven, gives the nak\d{s}atra.
\end{minipage}

\vspace{0.5em}
\begin{minipage}[t]{0.48\textwidth}
\small\devanagarifont
\textit{वा.~भा.}---चन्द्रात् स्पष्टात् सूर्यं शोध्य यद् भवति शेषं तत् द्वादशभागैः हृत्वा तिथयः स्युः स्फुटास्तदर्धात्। तिथयो द्वादशाङ्गुल्यः सूर्यचन्द्रयोरन्तरात् तिथीनां पञ्चदशभिः संयुक्तं योगनिर्मितिः सूर्यचन्द्रयो राशिसंयुक्त्यंशकलात्मिका सा सप्तविंशत्या हृता नक्षत्रं स्यात् ॥१०९--१११॥
\end{minipage}%
\hfill
\begin{minipage}[t]{0.48\textwidth}
\small
Having subtracted the Sun from the true Moon, dividing the remainder by twelve gives the tithis. The tithis are twelve digits from the difference of Sun and Moon. Joined with fifteen, the yoga is formed. The sum of the signs, degrees, and minutes of Sun and Moon, divided by twenty-seven, gives the nak\d{s}atra.
\end{minipage}

\vspace{0.5em}
\begin{minipage}[t]{0.48\textwidth}
\textbf{\small Sanskrit Verses}

\small\devanagarifont
चन्द्राच्च स्पष्टात् सूर्यं शोध्यं यद् भवति शेषं तत्\\
त्रिंशता हृत्वा नक्षत्रं तत् स्यात् स्फुटं राशिमानम् ॥११२॥\\[0.5em]
सूर्यस्य गतिः स्पष्टा चन्द्रस्यापि गतिः स्पष्टा तयोः\\
विशेषः स्यात् तिथ्यर्धभुक्तिरियं प्रकीर्तिता ॥११३॥
\end{minipage}%
\hfill
\begin{minipage}[t]{0.48\textwidth}
\textbf{\small English Translation}

\textbf{Verses 112--113. Nak\d{s}atra and Daily Motion Difference:}

Having subtracted the Sun from the true Moon, dividing the remainder by thirty, that nak\d{s}atra becomes the true sign-measure.

The true motion of the Sun and the true motion of the Moon --- their difference is called the tithy-ardha-bhukti (half-tithi daily motion).
\end{minipage}

\vspace{0.5em}
\begin{minipage}[t]{0.48\textwidth}
\small\devanagarifont
\textit{वा.~भा.}---चन्द्रात् स्पष्टात् सूर्यं शोध्य यद् भवति शेषं तत् त्रिंशता हृत्वा नक्षत्रं तत् स्यात् स्फुटं राशिमानम्। सूर्यस्य गतिः स्पष्टा चन्द्रस्यापि गतिः स्पष्टा तयोर्विशेषः स्यात् तिथ्यर्धभुक्तिः प्रकीर्तिता ॥११२--११३॥
\end{minipage}%
\hfill
\begin{minipage}[t]{0.48\textwidth}
\small
Having subtracted the Sun from the true Moon, dividing by thirty gives the nak\d{s}atra as the true sign-measure. The true motion of the Sun and the true motion of the Moon --- their difference is the tithy-ardha-bhukti.
\end{minipage}

\sectsep

%% ========================================================================
%% SPHERICAL ASTRONOMY
%% ========================================================================
\subsection{\dev{गोलानयनम् --- भारतीययवनक्रमः}}
\label{subsec:spherical}

\begin{minipage}[t]{0.48\textwidth}
\textbf{\small Sanskrit Verses}

\small\devanagarifont
यवनानां मते बिम्बं भूगोलं निरक्षदेशे स्थितम्\\
निजदेशोच्चं तन्मध्यं तत्रैव क्षितिजं वृत्तम् ॥११४॥\\[0.5em]
भारतीयमते त्वक्षवृत्तं मध्यमण्डलं तत्\\
क्षितिजं तन्महद् वृत्तं समरेखापि तत्रैव ॥११५॥\\[0.5em]
उन्नतजीवाकोटिज्याह्रिया हृज्या भुजो नतज्या वा\\
कर्णाक्षेत्रावृत्ते व्यासार्धं द्व्यमतोद्यन्न ॥११६॥
\end{minipage}%
\hfill
\begin{minipage}[t]{0.48\textwidth}
\textbf{\small English Translation}

\textbf{Verses 114--116. Spherical Astronomy --- Indian and Yavana Methods:}

In the opinion of the Yavanas (Greeks), the Earth's sphere is situated at the equatorial region; its highest point is at its centre; there indeed is the k\d{s}itija-v\=tta (horizon circle).

In the Indian method, the aksa-v\=tta (latitude circle) is the madhyama\d{n}\d{d}ala (prime vertical); the k\d{s}itija is that great circle; the sama-rekh\=a (equinoctial line) is also there.

The elevation-sine's ko\d{t}i-jy\=a, the altitude-sine is the bhuja or nata-jy\=a. The semi-diameter [applies to] the kar\d{n}a-field-circle in both methods (\=ubhayamato-dyanna).
\end{minipage}

\vspace{0.5em}
\begin{minipage}[t]{0.48\textwidth}
\small\devanagarifont
\textit{वा.~भा.}---यवनानां मते बिम्बं भूगोलं निरक्षदेशे स्थितं निजदेशोच्चं तन्मध्यं तत्रैव क्षितिजं वृत्तम्। भारतीयमते तु अक्षवृत्तं मध्यमण्डलं तत् क्षितिजं तन्महद् वृत्तं समरेखापि तत्रैव। उभयोरपि मतयोः सामान्यं गोलज्ञानमुभयमतोद्यन्नमित्युच्यते ॥११४--११६॥
\end{minipage}%
\hfill
\begin{minipage}[t]{0.48\textwidth}
\small
In the Yavana (Greek) view, the Earth's globe is at the equator, its highest point is at the centre, and there the horizon circle lies. In the Indian method, the latitude circle is the prime vertical, the horizon is that great circle, and the equinoctial line is also there. The common spherical knowledge of both methods is called ``ubhayamato-dyanna'' (accepted by both schools).
\end{minipage}

\sectsep

%% ========================================================================
%% CHAPTER: CANDRA-GRAHAṆĀDHIKĀRA
%% ========================================================================
\section{\dev{चन्द्रग्रहणाधिकारः}}
\addcontentsline{toc}{section}{Candra-graha\d{n}\=adhik\=ara\d{h} --- Chapter on Lunar Eclipses}

\begin{center}
\textit{Candra-graha\d{n}\=adhik\=ara\d{h} --- Chapter on the Calculation of Lunar Eclipses}
\end{center}

\begin{minipage}[t]{0.48\textwidth}
\small\devanagarifont
अस्याध्यायस्य विषयः --- चन्द्रग्रहणस्य कालः, मानं, तथा आकलनं। स्पर्शप्रथमस्पर्शमध्यमोक्षाणां निर्धारणं।
\end{minipage}%
\hfill
\begin{minipage}[t]{0.48\textwidth}
\small
This chapter deals with the computation of lunar eclipses, including the determination of the time, magnitude, and duration of the eclipse. Brahmagupta gives detailed rules for finding the \textit{parva} (full moon day), the \textit{spar\'sa} (first contact), the \textit{madhya} (middle), and the \textit{mok\d{s}a} (last contact) of a lunar eclipse.
\end{minipage}

\sectsep

%% ========================================================================
%% ŚLOKA 70--71
%% ========================================================================
\subsection{\dev{श्लोक ७०--७१}}

\begin{minipage}[t]{0.48\textwidth}
\textbf{\small Sanskrit Verses}

\small\devanagarifont
द्विगतिर्दिननिर्माथं स्फुटचन्द्रोदये तु यत् पिण्डम्।\\
तत् स्यात् स्पर्शपिण्डं ततो ग्रहणमध्यसाधनम् ॥७०॥\\[0.5em]
तत् स्यात् स्पर्शपिण्डं ततो ग्रहणमध्यसाधनम्।\\
तद्‌द्वितीयेन वा कालो ग्रहणस्य प्रकीर्तितः ॥७१॥
\end{minipage}%
\hfill
\begin{minipage}[t]{0.48\textwidth}
\textbf{\small English Translation}

\textbf{Verses 70--71. Spar\'sa-pi\d{n}\d{d}a (First Contact):}

At the true moonrise on the day of double motion, the pi\d{n}\d{d}a (magnitude) that is the spar\'sa-pi\d{n}\d{d}a (first contact magnitude), from that the computation of the middle of the eclipse.

That spar\'sa-pi\d{n}\d{d}a, from that the computation of the eclipse middle. Or by its half, the time of the eclipse is proclaimed.
\end{minipage}

\vspace{0.5em}
\begin{minipage}[t]{0.48\textwidth}
\small\devanagarifont
\textit{वि.~भा.}---द्वि गतिः दिननिर्माथं स्फुटचन्द्रोदये तु यत् पिण्डं तत् स्यात् स्पर्शपिण्डं ततो ग्रहणमध्यसाधनम्। तद्‌द्वितीयेन वा कालो ग्रहणस्य प्रकीर्तितः ॥७०--७१॥
\end{minipage}%
\hfill
\begin{minipage}[t]{0.48\textwidth}
\small
At the true moonrise on the day of double motion, the pi\d{n}\d{d}a is the spar\'sa-pi\d{n}\d{d}a (first contact magnitude), from which the middle of the eclipse is computed. Or by half of that, the time of the eclipse is proclaimed.
\end{minipage}

\sectsep

%% ========================================================================
%% ŚLOKA 72--73
%% ========================================================================
\subsection{\dev{श्लोक ७२--७३}}

\begin{minipage}[t]{0.48\textwidth}
\textbf{\small Sanskrit Verses}

\small\devanagarifont
चन्द्रसूर्ययोरन्तरं याम्यायनं तत् क्षयं वृद्धिं वा।\\
दक्षिणोत्तरयोः क्रान्त्योरन्तरं तत् क्षयं वृद्धिं वा ॥७२॥\\[0.5em]
याम्योत्तरयोः क्रान्त्योरन्तरं तत् क्षयं वृद्धिं वा।\\
तद्‌भुजज्या तु करणं तत् स्पर्शमोक्षौ प्रदर्शयेत् ॥७३॥
\end{minipage}%
\hfill
\begin{minipage}[t]{0.48\textwidth}
\textbf{\small English Translation}

\textbf{Verses 72--73. Eclipse Limits by Declination:}

The difference between the Moon and Sun's southern declination, that is the decrease or increase. The difference between their northern and southern declinations, that is the decrease or increase.

The difference between the southern and northern declinations, that decrease or increase. That bhuja-jy\=a is the kara\d{n}a (semi-duration); that shows the spar\'sa and mok\d{s}a (first and last contacts).
\end{minipage}

\vspace{0.5em}
\begin{minipage}[t]{0.48\textwidth}
\small\devanagarifont
\textit{वि.~भा.}---चन्द्रसूर्ययोरन्तरं याम्यायनं तत् क्षयं वृद्धिं वा। दक्षिणोत्तरयः क्रान्त्योरन्तरं तत् क्षयं वृद्धिं वा। याम्योत्तरयोः क्रान्त्योरन्तरं तत् क्षयं वृद्धिं वा। तद्‌भुजज्या तु करणं तत् स्पर्शमोक्षौ प्रदर्शयेत् ॥७२--७३॥
\end{minipage}%
\hfill
\begin{minipage}[t]{0.48\textwidth}
\small
The difference between Moon and Sun's southern declination gives the decrease or increase. The difference between their north-south declinations gives the decrease or increase. That bhuja-jy\=a is the kara\d{n}a, which determines the spar\'sa and mok\d{s}a contacts.
\end{minipage}

\sectsep

%% ========================================================================
%% ŚLOKA 74--80
%% ========================================================================
\subsection{\dev{श्लोक ७४--८०}}

\begin{minipage}[t]{0.48\textwidth}
\small\devanagarifont
याम्योत्तरयोः क्रान्त्योरन्तरं तत् क्षयं वृद्धिं वा।\\
तद्‌भुजज्या तु करणं तत् स्पर्शमोक्षौ प्रदर्शयेत् ॥७४॥\\[0.5em]
तद्‌भुजज्या तु करणं तत् स्पर्शमोक्षौ प्रदर्शयेत्।\\
तद्‌द्वितीयेन वा कालो ग्रहणस्य प्रकीर्तितः ॥७५॥\\[0.5em]
स्पर्शमोक्षौ यदा स्यातां तदा ग्रहणं प्रकीर्तितम्।\\
तद्‌द्वितीयेन वा कालो ग्रहणस्य प्रकीर्तितः ॥७६॥\\[0.5em]
तद्‌द्वितीयेन वा कालो ग्रहणस्य प्रकीर्तितः।\\
तद्‌भुजज्या तु करणं तत् स्पर्शमोक्षौ प्रदर्शयेत् ॥७७॥
\end{minipage}%
\hfill
\begin{minipage}[t]{0.48\textwidth}
\small
\textbf{Verses 74--77. Eclipse Contacts and Duration:}

The difference between southern and northern declinations, that decrease or increase. That bhuja-jy\=a is the kara\d{n}a; that shows the spar\'sa and mok\d{s}a contacts.

That bhuja-jy\=a is the kara\d{n}a; that shows the spar\'sa and mok\d{s}a. Or by its half, the time of the eclipse is proclaimed.

When the spar\'sa and mok\d{s}a occur, then the eclipse is proclaimed. Or by its half, the time of the eclipse is proclaimed.

Or by its half, the time of the eclipse is proclaimed. That bhuja-jy\=a is the kara\d{n}a; that shows the spar\'sa and mok\d{s}a.
\end{minipage}

\sectsep

%% ========================================================================
%% GRAHAṆASŪTRĀṆI
%% ========================================================================
\subsection{\dev{ग्रहणसूत्राणि}}

\begin{minipage}[t]{0.48\textwidth}
\small\devanagarifont
स्पर्शपिण्डं ततो ग्रहणमध्यसाधनम्।\\
तद्‌द्वितीयेन वा कालो ग्रहणस्य प्रकीर्तितः ॥७८॥\\[0.5em]
चन्द्रसूर्ययोरन्तरं याम्यायनं तत् क्षयं वृद्धिं वा।\\
दक्षिणोत्तरयोः क्रान्त्योरन्तरं तत् क्षयं वृद्धिं वा ॥७९॥\\[0.5em]
याम्योत्तरयोः क्रान्त्योरन्तरं तत् क्षयं वृद्धिं वा।\\
तद्‌भुजज्या तु करणं तत् स्पर्शमोक्षौ प्रदर्शयेत् ॥८०॥
\end{minipage}%
\hfill
\begin{minipage}[t]{0.48\textwidth}
\small
\textbf{Verses 78--80. Summary of Eclipse Rules:}

The spar\'sa-pi\d{n}\d{d}a, from that the computation of the eclipse middle. Or by its half, the time of the eclipse is proclaimed.

The difference between Moon and Sun, southern motion, that decrease or increase. The difference between their north-south declinations, that decrease or increase.

The difference between south-north declinations, that decrease or increase. That bhuja-jy\=a is the kara\d{n}a; that shows the spar\'sa and mok\d{s}a.
\end{minipage}

\sectsep

%% ========================================================================
%% GRAHAṆĀNAYANAM
%% ========================================================================
\subsection{\dev{ग्रहणानयनम्}}

\begin{minipage}[t]{0.48\textwidth}
\small\devanagarifont
द्विगतिर्दिननिर्माथं स्फुटचन्द्रोदये तु यत् पिण्डम्।\\
तत् स्यात् स्पर्शपिण्डं ततो ग्रहणमध्यसाधनम् ॥८१॥\\[0.5em]
तत् स्यात् स्पर्शपिण्डं ततो ग्रहणमध्यसाधनम्।\\
तद्‌द्वितीयेन वा कालो ग्रहणस्य प्रकीर्तितः ॥८२॥\\[0.5em]
चन्द्रसूर्ययोरन्तरं याम्यायनं तत् क्षयं वृद्धिं वा।\\
दक्षिणोत्तरयोः क्रान्त्योरन्तरं तत् क्षयं वृद्धिं वा ॥८३॥\\[0.5em]
याम्योत्तरयोः क्रान्त्योरन्तरं तत् क्षयं वृद्धिं वा।\\
तद्‌भुजज्या तु करणं तत् स्पर्शमोक्षौ प्रदर्शयेत् ॥८४॥
\end{minipage}%
\hfill
\begin{minipage}[t]{0.48\textwidth}
\small
\textbf{Verses 81--84. Eclipse Computation Summary:}

At the true moonrise on the day of double motion, the pi\d{n}\d{d}a is the spar\'sa-pi\d{n}\d{d}a (first contact magnitude), from which the middle of the eclipse is computed.

That spar\'sa-pi\d{n}\d{d}a, from that the eclipse middle is computed. Or by its half, the time of the eclipse is proclaimed.

The difference between Moon and Sun, the southern motion, that decrease or increase. The difference between their north-south declinations, that decrease or increase.

The difference between south-north declinations, that decrease or increase. That bhuja-jy\=a is the kara\d{n}a; that shows the spar\'sa and mok\d{s}a.
\end{minipage}

\sectsep

%% ========================================================================
%% GRAHAṆAMADHYAM
%% ========================================================================
\subsection{\dev{ग्रहणमध्यम्}}

\begin{minipage}[t]{0.48\textwidth}
\small\devanagarifont
स्पर्शमोक्षौ यदा स्यातां तदा ग्रहणं प्रकीर्तितम्।\\
तद्‌द्वितीयेन वा कालो ग्रहणस्य प्रकीर्तितः ॥८५॥\\[0.5em]
तद्‌द्वितीयेन वा कालो ग्रहणस्य प्रकीर्तितः।\\
तद्‌भुजज्या तु करणं तत् स्पर्शमोक्षौ प्रदर्शयेत् ॥८६॥\\[0.5em]
स्पर्शपिण्डं ततो ग्रहणमध्यसाधनम्।\\
तद्‌द्वितीयेन वा कालो ग्रहणस्य प्रकीर्तितः ॥८७॥
\end{minipage}%
\hfill
\begin{minipage}[t]{0.48\textwidth}
\small
\textbf{Verses 85--87. Middle of the Eclipse:}

When the spar\'sa and mok\d{s}a occur, then the eclipse is proclaimed. Or by its half, the time of the eclipse is proclaimed.

Or by its half, the time of the eclipse is proclaimed. That bhuja-jy\=a is the kara\d{n}a; that shows the spar\'sa and mok\d{s}a.

The spar\'sa-pi\d{n}\d{d}a, from that the eclipse middle is computed. Or by its half, the time of the eclipse is proclaimed.
\end{minipage}

\sectsep

%% ========================================================================
%% SŪRYAGRAHAṆA
%% ========================================================================
\subsection{\dev{सूर्यग्रहणम् --- रविग्रहणानयनम्}}

\begin{minipage}[t]{0.48\textwidth}
\textbf{\small Sanskrit Verses}

\small\devanagarifont
सूर्यग्रहणे तु लम्बनं शोध्यं चन्द्रसूर्ययोरन्तरात्\\
तत् शुद्धं स्यात् स्पर्शारम्भे मोक्षे तथैव ॥१२२॥\\[0.5em]
लम्बनात् स्पर्शकालः स्पर्शे युक्तो मोक्षे शुद्धः\\
ग्रहणमध्यं तु तयोर्योगार्धेनैव सिध्यति ॥१२३॥\\[0.5em]
चन्द्रस्य परिवेषो यः सूर्यस्य तु विशेषणात्\\
तत् स्याद् ग्रहणमानं तत् प्रकीर्तितमिहाचार्यैः ॥१२४॥
\end{minipage}%
\hfill
\begin{minipage}[t]{0.48\textwidth}
\textbf{\small English Translation}

\textbf{Verses 122--124. Solar Eclipse Computation:}

In a solar eclipse, the lambana (parallax in longitude) should be subtracted from the difference between the Moon and Sun; that is the pure [value] at the spar\'sa (first contact), and similarly at the mok\d{s}a (last contact).

The spar\'sa-k\=ala (time of first contact), increased by the lambana at the spar\'sa, is pure at the mok\=dsa. The middle of the eclipse is accomplished by half the sum of the two.

The parive\d{s}a (diameter) of the Moon, diminished by that of the Sun, that becomes the gr\=aha\d{n}a-m\=ana (magnitude of the eclipse), as proclaimed by the \=ac\=aryas.
\end{minipage}

\vspace{0.5em}
\begin{minipage}[t]{0.48\textwidth}
\small\devanagarifont
\textit{वा.~भा.}---सूर्यग्रहणे तु लम्बनं शोध्यं चन्द्रसूर्ययोरन्तरात् तत् शुद्धं स्यात् स्पर्शारम्भे मोक्षे तथैव। लम्बनात् स्पर्शकालः स्पर्शे युक्तो मोक्षे शुद्धः ग्रहणमध्यं तु तयोर्योगार्धेनैव सिध्यति। चन्द्रस्य परिवेषो यः सूर्यस्य तु विशेषणात् तत् स्याद् ग्रहणमानं तत् प्रकीर्तितमिहाचार्यैः ॥१२२--१२४॥
\end{minipage}%
\hfill
\begin{minipage}[t]{0.48\textwidth}
\small
In a solar eclipse, the lambana is subtracted from the Moon-Sun difference; that is pure at first contact and at last contact. The first contact time, increased by lambana at first contact, is pure at last contact. The eclipse middle is by half the sum of the two. The Moon's diameter, diminished by the Sun's, is the eclipse magnitude, as taught by the masters.
\end{minipage}

\sectsep

%% ========================================================================
%% LAMBANA
%% ========================================================================
\subsection{\dev{लम्बनम् --- देशान्तरलम्बनानयनम्}}

\begin{minipage}[t]{0.48\textwidth}
\textbf{\small Sanskrit Verses}

\small\devanagarifont
स्वदेशलम्बज्ञाया द्युज्याया चान्तरं यत् स्यात्\\
तत् त्रिज्यया हृत्वा लब्धं लम्बनं तत् प्रकीर्तितम् ॥१२५॥\\[0.5em]
लम्बज्या द्युज्यया हृत्वा त्रिज्यां गुणयेत् ततः\\
लब्धं तत् स्यात् लम्बनं तत् प्राचीनैरुक्तमीश्वरैः ॥१२६॥\\[0.5em]
तज्जं लम्बनजातं यत् तत् कालहृतं स्फुटं स्यात्\\
तद्धृतं नाडीषु रूपं स्यादेतल्लम्बनसाधनम् ॥१२७॥
\end{minipage}%
\hfill
\begin{minipage}[t]{0.48\textwidth}
\textbf{\small English Translation}

\textbf{Verses 125--127. Lambana (Parallax in Longitude):}

The difference between the lambaka-jy\=a of one's own place and the dyu-jy\=a, divided by the trijy\=a, the result obtained is called the lambana (parallax in longitude).

The lambaka-jy\=a divided by the dyu-jy\=a, multiplied by the trijy\=a, the result obtained is the lambana, as spoken by the ancient masters.

That lambana-produced quantity, divided by time, becomes the true [value]. That, converted into n\=a\d{d}\=is, becomes the form --- this is the lambana computation.
\end{minipage}

\vspace{0.5em}
\begin{minipage}[t]{0.48\textwidth}
\small\devanagarifont
\textit{वा.~भा.}---स्वदेशलम्बज्ञाया द्युज्याया चान्तरं यत् स्यात् तत् त्रिज्यया हृत्वा लब्धं लम्बनं तत् प्रकीर्तितम्। लम्बज्या द्युज्यया हृत्वा त्रिज्यां गुणयेत् ततो लब्धं तत् स्यात् लम्बनं तत् प्राचीनैरुक्तम् तज्जं लम्बनजातं यत् तत् कालहृतं स्फुटं स्यात् तद्धृतं नाडीषु रूपं स्यादेतल्लम्बनसाधनम् ॥१२५--१२७॥
\end{minipage}%
\hfill
\begin{minipage}[t]{0.48\textwidth}
\small
The difference between the lambaka-jy\=a of one's place and the dyu-jy\=a, divided by the trijy\=a, gives the lambana. The lambaka-jy\=a divided by the dyu-jy\=a, multiplied by the trijy\=a, gives the lambana as spoken by the ancients. That lambana-produced quantity, time-divided, becomes true; converted into n\=a\d{d}\=is, this is the lambana computation.
\end{minipage}

\sectsep

%% ========================================================================
%% GRAHAṆA-SPARŚA-MOKṢA
%% ========================================================================
\subsection{\dev{ग्रहणस्पर्श-मोक्षयोः}}

\begin{minipage}[t]{0.48\textwidth}
\small\devanagarifont
चन्द्रसूर्ययोरन्तरं याम्यायनं तत् क्षयं वृद्धिं वा।\\
दक्षिणोत्तरयोः क्रान्त्योरन्तरं तत् क्षयं वृद्धिं वा ॥८८॥\\[0.5em]
याम्योत्तरयोः क्रान्त्योरन्तरं तत् क्षयं वृद्धिं वा।\\
तद्‌भुजज्या तु करणं तत् स्पर्शमोक्षौ प्रदर्शयेत् ॥८९॥\\[0.5em]
तद्‌भुजज्या तु करणं तत् स्पर्शमोक्षौ प्रदर्शयेत्।\\
तद्‌द्वितीयेन वा कालो ग्रहणस्य प्रकीर्तितः ॥९०॥
\end{minipage}%
\hfill
\begin{minipage}[t]{0.48\textwidth}
\small
\textbf{Verses 88--90. Eclipse Contacts (Spar\'sa and Mok\d{s}a):}

The difference between Moon and Sun, the southern motion, that decrease or increase. The difference between their north-south declinations, that decrease or increase.

The difference between south-north declinations, that decrease or increase. That bhuja-jy\=a is the kara\d{n}a; that shows the spar\'sa and mok\d{s}a.

That bhuja-jy\=a is the kara\d{n}a; that shows the spar\'sa and mok\d{s}a. Or by its half, the time of the eclipse is proclaimed.
\end{minipage}

\sectsep

%% ========================================================================
%% SUMMARY OF VERSES
%% ========================================================================
\section*{Summary and Verse Index}
\addcontentsline{toc}{section}{Summary and Verse Index}

\begin{minipage}[t]{0.48\textwidth}
\small\devanagarifont
\textbf{सङ्क्षेपः}

एतस्मिन् खण्डे निम्नलिखिताः श्लोकाः समावेशिताः:
\begin{itemize}[leftmargin=*,nosep]
\item स्पष्टाधिकारः (श्लोक १२--५१)
\item त्रिप्रश्नाधिकारः (श्लोक ५४--६९)
\item छायाकरणी (श्लोक २८--३६)
\item लम्बाक्षज्यानयनम् (श्लोक ५७--६२)
\item प्राणगणितम् (श्लोक १२८--१३०)
\item मध्यच्छाया (श्लोक १४८)
\item महापातः (श्लोक १४०--१४२)
\item क्रान्तिः (श्लोक १३१--१३३)
\item लग्ननयनम् (श्लोक १३४--१३६)
\item गोलयन्त्रम् (श्लोक १३७--१३९)
\item ज्योत्पत्तिः (श्लोक १००--१०३)
\item गौरीगुप्ता (श्लोक १०४--१०८)
\item पाटीगणितम् (श्लोक ११७--१२१)
\item तिथियोगनक्षत्रम् (श्लोक १०९--११३)
\item गोलानयनम् (श्लोक ११४--११६)
\item चन्द्रग्रहणम् (श्लोक ७०--९०)
\item सूर्यग्रहणम् (श्लोक १२२--१२४)
\item लम्बनम् (श्लोक १२५--१२७)
\end{itemize}
\end{minipage}%
\hfill
\begin{minipage}[t]{0.48\textwidth}
\small
\textbf{English Summary}

This chunk contains the following verses:
\begin{itemize}[leftmargin=*,nosep]
\item Spa\d{s}\d{t}\=adhik\=ara (Verses 12--51) --- True positions of planets
\item Tripra\'sn\=adhik\=ara (Verses 54--69) --- Three questions of spherical astronomy
\item Ch\=ay\=akara\d{n}\=i (Verses 28--36) --- Shadow and gnomon calculations
\item Lambh\=ak\d{s}a-jy\=a (Verses 57--62) --- Latitude-related sine computation
\item Pr\=a\d{n}aga\d{n}ita (Verses 128--130) --- Time unit computation
\item Madhya-ch\=ay\=a (Verse 148) --- Midday shadow
\item Mah\=ap\=ata (Verses 140--142) --- Great conjunction and lunar month
\item Kr\=anti (Verses 131--133) --- Declination and precession
\item Lagna-nayana (Verses 134--136) --- Rising sign computation
\item Gola-yantra (Verses 137--139) --- Armillary sphere
\item Jyotpatti (Verses 100--103) --- Sine table construction
\item Gaur\=i-gupt\=a (Verses 104--108) --- Second-difference interpolation
\item P\=a\d{t}\=iga\d{n}ita (Verses 117--121) --- Arithmetic operations
\item Tithi-yoga-nak\d{s}atra (Verses 109--113) --- Luni-solar astronomy
\item Gol\=anayana (Verses 114--116) --- Spherical astronomy
\item Candra-graha\d{n}a (Verses 70--90) --- Lunar eclipse
\item S\=urya-graha\d{n}a (Verses 122--124) --- Solar eclipse
\item Lambana (Verses 125--127) --- Parallax computation
\end{itemize}
\end{minipage}

\vspace{1em}
\longline

%% ========================================================================
%% FINAL SECTION
%% ========================================================================
\section*{End of Part 2, Chunk 2}
\addcontentsline{toc}{section}{End of Part 2, Chunk 2}

\begin{center}
\textit{This completes the English-Sanskrit bilingual edition of pages 979--1118 of the Br\=ahmasphu\d{t}asiddh\=anta.}
\end{center}

\vspace{0.5em}

\begin{minipage}[t]{0.48\textwidth}
\small\devanagarifont
\textbf{समाप्तिविवरणम्}

एषः द्वितीयभागस्य द्वितीयखण्डः समाप्तः। अत्र समावेशितं साहित्यं निम्नलिखिताध्यायान् समेति:
\begin{enumerate}[nosep]
\item \textbf{स्पष्टाधिकारः} --- ग्रहाणां मध्यमात् स्पष्टस्थानस्य गणनम्
\item \textbf{त्रिप्रश्नाधिकारः} --- लम्बननाडीछायागणनम्
\item \textbf{चन्द्रग्रहणाधिकारः} --- चन्द्रग्रहणकालमानाकलनम्
\item \textbf{गणितीयसूत्राणि} --- ज्योत्पत्तिः, गौरीगुप्ता, पाटीगणितम्
\end{enumerate}

ब्रह्मगुप्तस्याः सिद्धान्ते प्रत्येकं नियमः संस्कृतपद्येषु (आर्यायां) उपस्थापितः, ततो विस्तृतव्याख्या (व्याख्या) अनुसरति। गणितीयसूत्राणि भारतीयज्यापद्धतिं प्रयुञ्जते तथा त्रैराशिकं ज्यामितीयं च तर्कं प्रयुञ्जते।
\end{minipage}%
\hfill
\begin{minipage}[t]{0.48\textwidth}
\small
\textbf{Closing Notes}

The covered material spans several important chapters of Brahmagupta's monumental work:

\begin{enumerate}[nosep]
\item \textbf{Spa\d{s}\d{t}\=adhik\=ara\d{h}} --- Methods for computing the true positions of planets from their mean longitudes, including the \textit{\'s\=ighra} (anomaly) and \textit{manda} (eccentricity) corrections.

\item \textbf{Tripra\'sn\=adhik\=ara\d{h}} --- The ``three questions'' of spherical astronomy: computation of \textit{lambana} (parallax), \textit{n\=a\d{d}\=i} (time difference), shadow calculations (\textit{ch\=ay\=akara\d{n}\=i}), and related quantities.

\item \textbf{Candra-graha\d{n}\=adhik\=ara\d{h}} --- The calculation of lunar eclipses, including the determination of \textit{spar\'sa} (first contact), \textit{madhya} (middle), and \textit{mok\d{s}a} (last contact).

\item \textbf{Mathematical Sections} --- Construction of sine tables (\textit{jyotpatti}), second-difference interpolation (\textit{gaur\=i-gupt\=a}), arithmetic operations (\textit{p\=a\d{t}\=iga\d{n}ita}), luni-solar astronomy, and spherical astronomy.
\end{enumerate}

Throughout these chapters, Brahmagupta presents his mathematical rules in the form of Sanskrit \textit{\=ary\=a} verses, each followed by detailed commentary (\textit{vy\=akhy\=a}) that explains the astronomical and mathematical principles involved. The formulas employ the Indian sine (\textit{jy\=a}) system and make extensive use of proportions (\textit{trair\=a\'sika}) and geometric reasoning.
\end{minipage}

\vspace{1.5cm}

\begin{center}
\shortline\\[0.5cm]
\textit{To be continued in Part 2, Chunk 3}\\
\textit{(Pages 1119--1676 of the Complete Work)}\\[0.5cm]
\shortline
\end{center}

\vfill

\begin{center}
{\small\textit{Typeset in XeLaTeX with Devanagari support --- Bilingual Edition}}\\[0.3cm]
{\small Original: Indian Institute of Astronomical \& Sanskrit Research}
\end{center}

\end{document}



% ============================================================
% Source: translations/en/indian/brahmagupta-brahmasphutasiddhanta-pt2-chunk3_bilingual.tex
% ============================================================

%% ========================================================================
%% Brāhmasphuṭasiddhānta Part 2 — Chunk 3 (Pages 1119–1258)
%% BILINGUAL English-Sanskrit Edition
%% LEFT: Sanskrit in Devanagari  |  RIGHT: English translation with IAST
%% ========================================================================

\documentclass[11pt,a4paper]{article}

%% -------- Core Packages --------
\usepackage{fontspec}
\usepackage{polyglossia}
\setmainlanguage{english}
\setotherlanguage{sanskrit}

\usepackage{amsmath,amssymb,amsthm}
\usepackage{geometry}
\usepackage{graphicx}
\usepackage{xcolor}
\usepackage{titlesec}
\usepackage{fancyhdr}
\usepackage{enumitem}
\usepackage{array,booktabs,longtable,multirow}
\usepackage{hyperref}
% \usepackage{verse}  % Package corrupted (404 HTML) -- not needed
\usepackage{lettrine}
\usepackage{microtype}
\usepackage{tocloft}
\usepackage{indentfirst}
\usepackage{paracol}
\usepackage{multicol}

%% -------- Page Geometry --------
\geometry{
  paperwidth=210mm,paperheight=297mm,
  left=18mm,right=18mm,top=25mm,bottom=25mm,
  headheight=14pt,footskip=30pt
}

%% -------- Font Configuration --------
\setmainfont[
  Path=/usr/share/fonts/opentype/linux-libertine/,
  BoldFont=LinLibertine_RB.otf,
  ItalicFont=LinLibertine_RI.otf,
  BoldItalicFont=LinLibertine_RBI.otf]{LinLibertine_R.otf}
\setsansfont[
  Path=/usr/share/fonts/opentype/linux-libertine/,
  BoldFont=LinBiolinum_RB.otf,
  ItalicFont=LinBiolinum_RI.otf]{LinBiolinum_R.otf}
\setmonofont[
  Path=/usr/share/fonts/truetype/dejavu/,
  Scale=MatchLowercase
]{DejaVuSansMono.ttf}

%% Devanagari (Sanskrit/Hindi) Font
\newfontfamily\devanagarifont[Script=Devanagari,ExternalLocation]{/usr/share/fonts/truetype/noto/NotoSerifDevanagari-Regular.ttf}
\newfontfamily\devanagarifontbf[Script=Devanagari,ExternalLocation]{/usr/share/fonts/truetype/noto/NotoSerifDevanagari-Bold.ttf}

%% -------- Colors --------
\definecolor{titlecolor}{RGB}{45,55,72}
\definecolor{sectioncolor}{RGB}{55,65,81}
\definecolor{notecolor}{RGB}{120,53,15}
\definecolor{rulecolor}{RGB}{180,160,140}
\definecolor{versebg}{RGB}{250,248,245}
\definecolor{sanskritbg}{RGB}{255,253,250}
\definecolor{englishbg}{RGB}{252,252,255}
\definecolor{columnsep}{RGB}{160,140,120}

%% -------- Section Formatting --------
\titleformat{\section}
  {\normalfont\Large\bfseries\color{sectioncolor}}
  {\thesection}{1em}{}
\titleformat{\subsection}
  {\normalfont\large\bfseries\color{sectioncolor}}
  {\thesubsection}{1em}{}
\titleformat{\subsubsection}
  {\normalfont\normalsize\bfseries\color{sectioncolor}}
  {\thesubsubsection}{1em}{}

%% -------- Header/Footer --------
\pagestyle{fancy}
\fancyhf{}
\fancyhead[L]{\small\textit{Brāhmasphuṭasiddhānta — Bilingual Edition — Part II, Chunk 3}}
\fancyhead[R]{\small\thepage}
\renewcommand{\headrulewidth}{0.4pt}
\renewcommand{\footrulewidth}{0pt}

%% -------- Custom Commands --------
\newcommand{\dev}[1]{{\devanagarifont #1}}
\newcommand{\devbf}[1]{{\devanagarifontbf #1}}

%% Parallel column commands
\newcommand{\sktcol}[1]{\fcolorbox{rulecolor}{sanskritbg}{\parbox{\linewidth}{\dev{#1}}}}
\newcommand{\engcol}[1]{\fcolorbox{rulecolor}{englishbg}{\parbox{\linewidth}{#1}}}

%% Verse environment for Sanskrit
\newenvironment{sanskritverse}{%
  \begin{center}\begin{minipage}{0.95\textwidth}
  \setlength{\leftskip}{1em}
  \setlength{\rightskip}{1em}
  \dev
}{\end{minipage}\end{center}}

%% Parallel verse environment
\newenvironment{paralleltext}{%
  \begin{paracol}{2}
  \backgroundcolor{c[0](.95\columnsep,.95\textheight)}{sanskritbg}
  \backgroundcolor{c[1](.95\columnsep,.95\textheight)}{englishbg}
  \def\sanskritcolumn{\switchcolumn[0]*}
  \def\englishcolumn{\switchcolumn[1]*}
}{\end{paracol}}

%% -------- Paracol Settings --------
\columnsep{10pt}
\setlength{\columnseprule}{0.4pt}
\def\columnseprulecolor{\color{columnsep}}

%% -------- Hyperref Setup --------
\hypersetup{
  colorlinks=true,
  linkcolor=sectioncolor,
  citecolor=sectioncolor,
  urlcolor=sectioncolor,
  pdfencoding=auto,
  pdftitle={Brāhmasphuṭasiddhānta Part 2 Chunk 3 — Bilingual Edition},
  pdfauthor={Brahmagupta, with commentary by Sudhākara Dvivedī}
}

%% ========================================================================
\begin{document}

%% ========================================================================
%% TITLE PAGE
%% ========================================================================
\begin{titlepage}
\begin{center}
\vspace*{1.5cm}

{\devanagarifont\fontsize{26}{30}\selectfont\color{titlecolor} ब्राह्मस्फुटसिद्धान्तः}\\[0.8cm]
{\fontsize{20}{24}\selectfont\color{sectioncolor} Brāhmasphuṭasiddhānta}\\[1.2cm]

\rule{0.7\textwidth}{1pt}\\[0.8cm]

{\LARGE\textbf{BILINGUAL EDITION}}\\[0.3cm]
{\large Sanskrit (Devanagari) $\parallel$ English (with IAST)}\\[1.2cm]

{\Large\textbf{Part Two — Chunk 3}}\\[0.3cm]
{\large (Pages 281--420 of PDF $\approx$ Original pages 1119--1258)}\\[1.5cm]

{\large By \textbf{Brahmagupta} (6th century CE)}\\[0.5cm]
{\large With Commentary (\dev{नूतनतिलक}) by}\\[0.2cm]
{\large \textbf{Śrī Kṛpāludatta Sūnu Sudhākara Dvivedī}}\\[1.5cm]

\rule{0.7\textwidth}{0.4pt}\\[0.8cm]

{\normalsize Contains the following chapters:}\\[0.3cm]
\begin{minipage}{0.8\textwidth}
\begin{center}
{\devanagarifont चन्द्रग्रहणविकारः} — Lunar Eclipse Variations\\[0.15cm]
{\devanagarifont सूर्यग्रहणाधिकारः} — Solar Eclipse Chapter\\[0.15cm]
{\devanagarifont चन्द्रच्छायाधिकारः} — Moon's Shadow Chapter\\[0.15cm]
{\devanagarifont ग्रहोदयास्ताधिकारः} — Planetary Rising and Setting\\[0.15cm]
{\devanagarifont शुक्लोन्नत्यधिकारः} — Brightness of the Moon\\[0.15cm]
{\devanagarifont चन्द्रशुक्लोन्नत्यधिकारः} — Moon's Bright Elevation\\[0.15cm]
\end{center}
\end{minipage}

\vfill
{\small Typeset with XeLaTeX using Noto Serif Devanagari}\\
{\small Bilingual edition: Sanskrit left — English right}

\end{center}
\end{titlepage}

%% ========================================================================
%% TABLE OF CONTENTS
%% ========================================================================
\tableofcontents
\newpage

%% ========================================================================
%% EDITORIAL PREFACE TO THE BILINGUAL EDITION
%% ========================================================================
\section*{Editorial Preface}
\addcontentsline{toc}{section}{Editorial Preface}

This bilingual edition presents the Sanskrit text of Brāhmasphuṭasiddhānta Part II, Chunk 3, with facing English translations. The layout follows the standard format for critical editions of Sanskrit scientific texts:

\begin{itemize}
  \item \textbf{Left column:} Sanskrit text in Devanagari script, as found in the source edition.
  \item \textbf{Right column:} English translation with IAST romanization of technical terms.
\end{itemize}

The commentary (\dev{टीका}) by Sudhākara Dvivedī follows the traditional Indian \emph{pañcāṅga} structure of:
\begin{enumerate}
  \item \textbf{Sūtra} (\dev{सूत्रम्}) — The original verse in āryā or anuṣṭubh metre.
  \item \textbf{Sūtrārtha} (\dev{सूत्रार्थः}) — Literal meaning of the verse.
  \item \textbf{Hindi Bhāṣya} (\dev{हिन्दी भाष्य}) — Explanation in Hindi.
  \item \textbf{Upapatti} (\dev{उपपत्तिः}) — Mathematical proof or demonstration.
\end{enumerate}

\medskip
\noindent\rule{\textwidth}{0.4pt}

\newpage


%% ========================================================================
%% CHAPTER 1: CHANDRAGRAHAṆAVIKĀRAḤ (Lunar Eclipse Variations)
%% ========================================================================
\section{\texorpdfstring{{\devanagarifont चन्द्रग्रहणविकारः} — Candragrahaṇavikāraḥ (Lunar Eclipse Variations)}{Candragrahaṇavikāraḥ (Lunar Eclipse Variations)}}
\label{sec:chandra-grahana}

\begin{center}
\textcolor{rulecolor}{\rule{0.8\textwidth}{0.5pt}}
\end{center}

\subsection*{Introduction}

This chapter (\dev{चन्द्रग्रहणविकारः}) deals with the computation of lunar eclipses, following Brahmagupta's exposition in the \dev{ब्राह्मस्फुटसिद्धान्त}. The commentary (\dev{सुधाकर-टीका}) by Paṇḍita Kṛpāludatta Sūnu Sudhākara Dvivedī provides elaborate explanations in Hindi, with mathematical demonstrations (\dev{उपपत्ति}).

\begin{center}
\textcolor{rulecolor}{\rule{0.8\textwidth}{0.5pt}}
\end{center}

%% ========================================================================
\subsection{Verses 11--12: Computation of \texorpdfstring{{\devanagarifont इष्टग्रास}}{iṣṭagrāsa} (Desired Eclipse Magnitude)}

\vspace{0.5em}
\noindent\textbf{Sūtra (\dev{सूत्रम्}) — Verses 11--12:}

\begin{paracol}{2}
\backgroundcolor{c[0](.95\columnsep,.95\textheight)}{sanskritbg}
\backgroundcolor{c[1](.95\columnsep,.95\textheight)}{englishbg}

\begin{sanskritverse}
इष्टन्यूनस्थितिदलगुणा भुक्तिविल्लेपलिप्ता पश्चादष्टा भवति हि भुजः कोटिरिष्टेणुवाणः \\
तद्वेकोद्भवमपि पदं कर्णं एतेन हीनं मानैक्यार्धं स्फुटमिह भवेदाकलित छन्नमानम्॥ ११--१२ ॥
\end{sanskritverse}

\switchcolumn

\textbf{Translation:} ``The half-duration (\textit{sthitidala}) diminished by the desired quantity, multiplied by the daily motion (\textit{bhukti}) and divided by the deflection (\textit{villepa}), gives the \textit{bhuja} (sine). The \textit{koṭi} (cosine) is computed by the rule of the desired quantity (\textit{iṣṭeṇu}). From the half-sum of the diameters (\textit{mānaikyārdha}), the \textit{karṇa} (hypotenuse) being subtracted, the result gives the true obscuration (\textit{channa-māna}).''

\end{paracol}

\vspace{0.5em}

\begin{paracol}{2}
\backgroundcolor{c[0](.95\columnsep,.95\textheight)}{sanskritbg}
\backgroundcolor{c[1](.95\columnsep,.95\textheight)}{englishbg}

\noindent\textbf{Sūtrārtha (\dev{सूत्रार्थः}):}
\begin{sanskritverse}
इष्टन्यूनस्थितिदलगुणा भुक्तिविल्लेपलिप्ता पश्चादष्टा भवति हि भुजः कोटिरिष्टेणुवाणः।
तद्वेकोद्भवमपि पदं कर्णं एतेन हीनं मानैक्यार्धं स्फुटमिह भवेदाकलित छन्नमानम्॥
\end{sanskritverse}

\switchcolumn

\textbf{Literal Meaning:} The verse gives the method for computing the \textit{iṣṭagrāsa} (desired eclipse magnitude). From the \textit{sthitidala} (half-duration) diminished by the desired quantity, multiplied by the \textit{bhukti} (daily motion) and divided by the \textit{villepa} (deflection), one gets the \textit{bhuja} (sine). The \textit{koṭi} (cosine) is computed by the rule of \textit{iṣṭeṇu} (desired quantity). From the \textit{mānaikyārdha} (half-sum of diameters), the \textit{karṇa} (hypotenuse) being subtracted, the result gives the true obscuration.

\end{paracol}

\vspace{0.5em}

\begin{paracol}{2}
\backgroundcolor{c[0](.95\columnsep,.95\textheight)}{sanskritbg}
\backgroundcolor{c[1](.95\columnsep,.95\textheight)}{englishbg}

\noindent\textbf{Hindi Bhāṣya (\dev{हिन्दी भाष्य}):}
\begin{sanskritverse}
इष्ट रहित स्थितिदल को भुक्ति से गुणा कर विल्लेप से भाग देने से भुज होता है, इस फल करके
सूर्य, चन्द्र और पात को चालन देकर तात्कालिक चन्द्रशर साधन करना वह कोटि है। उन दोनों
(भुज और कोटि) का वर्गयोग मूल कर्ण होता है, मानैक्यार्ध (छाद्य और छादक बिम्ब के व्यासार्ध
योग) में से कर्ण को घटाने से जो शेष रहता है वह इष्टग्रास होता है इति॥ ११--१२ ॥
\end{sanskritverse}

\switchcolumn

\textbf{English Translation of Hindi Commentary:} ``Subtracting the desired (\textit{iṣṭa}) from the \textit{sthitidala}, multiplying by the \textit{bhukti} and dividing by the \textit{villepa}, one obtains the \textit{bhuja}. By applying this result to the motions of the Sun, Moon, and node (\textit{pāta}), one computes the instantaneous lunar latitude (\textit{tātkālika candra-śara}), which is the \textit{koṭi}. The square root of the sum of their squares is the \textit{karṇa}. From the \textit{mānaikyārdha} (sum of the radii of the eclipsed and eclipsing bodies), subtracting the \textit{karṇa}, the remainder is the desired magnitude (\textit{iṣṭagrāsa}).''

\end{paracol}

\vspace{0.5em}

\begin{paracol}{2}
\backgroundcolor{c[0](.95\columnsep,.95\textheight)}{sanskritbg}
\backgroundcolor{c[1](.95\columnsep,.95\textheight)}{englishbg}

\noindent\textbf{Upapatti (\dev{उपपत्तिः}) — Proof:}
\begin{sanskritverse}
स्पर्शी के बाद जितनी इष्ट घटी में ग्रास-ज्ञान अपेक्षित हो उस इष्ट घटी को स्थित्यर्ध में से घटा
देना, शेष से अनुपात करते हैं। यदि साठ घटी में रवि और चन्द्र की गत्यन्तर कला पाते हैं तो
इष्टोन स्थित्यर्ध घटी में क्या इससे जो फल होता है वह भुज है, अनुपातागत फल से चालित रवि,
चन्द्र और पात से तात्कालिक चन्द्रवार साधन करना वह कोटि है। इन दोनों (भुज और कोटि)
का वर्गयोग मूल कर्ण होता है, मानैक्यार्ध में से कर्ण को घटाने से इष्ट ग्रास होता है।
\end{sanskritverse}

\switchcolumn

\textbf{English Translation of Proof:} ``After the contact (\textit{sparśī}), for whatever desired number of ghaṭīs the knowledge of the \textit{grāsa} is required, that desired number of ghaṭīs should be subtracted from the \textit{sthityardha}. The remainder is used in proportion. If in sixty ghaṭīs the difference in motion between the Sun and Moon gives a certain number of minutes (\textit{kalā}), then [by proportion] in the \textit{iṣṭa}-diminished \textit{sthityardha} ghaṭīs, whatever result is obtained is the \textit{bhuja}. From the result obtained by proportion, by applying the motions of the Sun, Moon, and node, the instantaneous lunar latitude (\textit{tātkālika candra-vāra}) is computed --- that is the \textit{koṭi}. The square root of the sum of the squares of these two (\textit{bhuja} and \textit{koṭi}) is the \textit{karṇa}. Subtracting the \textit{karṇa} from the \textit{mānaikyārdha} gives the desired \textit{grāsa}.''

\end{paracol}

%% ========================================================================
\subsection{Verses 13--14: Computation of \texorpdfstring{{\devanagarifont इष्टग्रासात्काल}}{iṣṭagrāsāt-kāla} (Time from Desired Magnitude)}

\vspace{0.5em}
\noindent\textbf{Sūtra (\dev{सूत्रम्}) — Verses 13--14:}

\begin{paracol}{2}
\backgroundcolor{c[0](.95\columnsep,.95\textheight)}{sanskritbg}
\backgroundcolor{c[1](.95\columnsep,.95\textheight)}{englishbg}

\begin{sanskritverse}
प्रसकृत्-ग्रासकलोनप्रमाणयुतिदलकृते विशोध्य कृतिम्।
तात्कालिकविक्षेपस्य शेषमूलं कृतं तिथिवत्॥ १३ ॥
\end{sanskritverse}

\begin{sanskritverse}
प्रग्रहणस्थित्यर्धात् प्रोह्य प्रग्रहणतो भवेत् कालः।
सौक्षं विशोष्य मोक्षस्थित्यर्धात् प्राग् भवेन्मोक्षात्॥ १४ ॥
\end{sanskritverse}

\switchcolumn

\textbf{Translation — Verse 13:} ``From the square of the half-sum of the measure diminished by the \textit{grāsakalā}, having subtracted the square of the instantaneous latitude (\textit{tātkālika-vikṣepa}), the square root of the remainder --- made like a \textit{tithi} [i.e., multiplied by 60 and divided by the difference of motions] --- [gives a time].''

\textbf{Translation — Verse 14:} ``Having subtracted [this time] from the \textit{pragrahaṇa-sthit yardha}, the time from the \textit{pragrahaṇa} [i.e., from contact] is obtained. Having subtracted from the \textit{mokṣa-sthit yardha}, [one obtains the time] before the \textit{mokṣa} [release].''

\end{paracol}

\vspace{0.5em}

\begin{paracol}{2}
\backgroundcolor{c[0](.95\columnsep,.95\textheight)}{sanskritbg}
\backgroundcolor{c[1](.95\columnsep,.95\textheight)}{englishbg}

\noindent\textbf{Sūtrārtha (\dev{सूत्रार्थः}):}
\begin{sanskritverse}
ग्रासकलोनप्रमाणयुतिदलकृते ग्रासोनमानैक्यार्धवर्गात् तात्कालिकशरस्य कृति विद्धोध्य
शेषशूलं तिथिवदसकृत् कृतम्। अर्थात् शेषमूलं षष्टिगुणां रविचन्द्रगत्यन्तरहतं फलकालेन
रविचन्द्रपातानु प्रचाल्य तात्कालिकविक्षेपं संसाध्य तस्मात् पुनः ग्रासिकलोनप्रमाणेत्यादिना
असकृत् शेषमूलकालः स्थिरीभवति तं प्रग्रहणस्थित्यर्धात् (स्पर्शिकस्थित्यर्धात्) प्रोह्य
(हित्वा) प्रग्रहणतः (स्पर्शादनन्तरं) कालो भवेत्।
\end{sanskritverse}

\switchcolumn

\textbf{Literal Meaning:} ``From the square of the half-sum of the measure [of diameters] diminished by the \textit{grāsakalā}, having subtracted the square of the instantaneous \textit{śara} (latitude), the square root of the remainder --- made \textit{tithi}-like (\textit{tithivat}) [by multiplying by 60 and dividing by the difference of motions of Sun and Moon] --- is the result. That is, having multiplied the square root of the remainder by 60, divided by the difference of motions of Sun and Moon, and having applied the motions of Sun, Moon, and node by that resulting time, the instantaneous latitude (\textit{tātkālika-vikṣepa}) is computed. From that, again by the rule `\textit{grāsakalona-pramāṇa}', etc., iteratively (\textit{asakṛt}), the square-root-time becomes stable. That [time], subtracted from the \textit{pragrahaṇa-sthit yardha} (contact half-duration), the time from the \textit{pragrahaṇa} (from contact) is obtained.''

\end{paracol}

\vspace{0.5em}

\begin{paracol}{2}
\backgroundcolor{c[0](.95\columnsep,.95\textheight)}{sanskritbg}
\backgroundcolor{c[1](.95\columnsep,.95\textheight)}{englishbg}

\noindent\textbf{Hindi Bhāṣya (\dev{हिन्दी भाष्य}):}
\begin{sanskritverse}
ग्रासकलोनप्रमाणयुतिदलकृतेः (ग्रासरहितमानैक्यार्धवर्गात्) तात्कालिकशरस्य कृति
(वर्ग) विशोध्य (हित्वा) शेषस्थ मूल तिथिवदसकृत् कृतमर्थात् शेषमूलं षष्टिगुणं
रविचन्द्रयोगंत्यन्तरभक्तं फलकालेन रविचन्द्रपातानु प्रचाल्य तात्कालिकचन्द्रशरं संसाध्य
तस्मात् पुनर्ग्रासकलोनप्रमाणयुतिदलकृतेरित्यादिना असकृत् शेषमूलं कालः स्थिरीभवति
ते प्रग्रहणस्थित्यर्धात् (स्पर्शिकस्थित्यर्धात्) प्रोह्य (हित्वा) प्रग्रहणतः (स्पर्शादनन्तरं)
कालो भवेदर्थात् स्पर्शानन्तरमेतावतीष्ठकाले तावानिष्टग्रासो भवति॥ १३--१४ ॥
\end{sanskritverse}

\switchcolumn

\textbf{English Translation:} ``From the square of the \textit{grāsa}-diminished \textit{mānaikyārdha}, having subtracted the square of the instantaneous \textit{śara}, the square root of the remainder --- made \textit{tithi}-like iteratively --- that is, having multiplied the square root of the remainder by 60, divided by the combined difference of motions of Sun and Moon, and having applied the motions of Sun, Moon, and node by that resulting time, the instantaneous lunar latitude is computed. From that, again by the rule beginning `\textit{grāsakalona-pramāṇa...}', iteratively (\textit{asakṛt}), the square-root-time becomes stable. That, subtracted from the contact half-duration (\textit{pragrahaṇa-sthit yardha}), gives the time from contact, meaning that after contact, in that desired time, the desired \textit{grāsa} is obtained.''

\end{paracol}

\vspace{0.5em}

\begin{paracol}{2}
\backgroundcolor{c[0](.95\columnsep,.95\textheight)}{sanskritbg}
\backgroundcolor{c[1](.95\columnsep,.95\textheight)}{englishbg}

\noindent\textbf{Upapatti (\dev{उपपत्तिः}):}
\begin{sanskritverse}
तत्कालशरस्याज्ञानान्मध्यकालिकशरात् कर्म कृतमतोऽसकृद्धिघिना स्फुटकालसाधनमुचितम्।
शेषोपपत्ति ``ग्रसोन मानैक्यदलस्य वर्गाद् विक्षेपकृत्या रहिता'' इत्यादिभास्करविधिना
स्फुटा ॥ १३--१४ ॥
\end{sanskritverse}

\switchcolumn

\textbf{English Translation of Proof:} ``Due to the ignorance of the latitude at that time (\textit{tātkāla}), the computation has been done using the middle-time latitude. Therefore, it is proper to determine the accurate time by an iterative method (\textit{asakṛt-ghātinā}). The remaining proof [follows] the method of Bhāskara: `From the square of half the sum of diameters diminished by the \textit{grāsa}, having subtracted the square of the \textit{vikṣepa}...' etc.''

\end{paracol}

%% ========================================================================
\subsection{Verse 9: Sphuṭīkaraṇa of \texorpdfstring{{\devanagarifont स्थितिविमर्दार्ध}}{sthitivimardārdha} (Refinement of Half-Duration of Totality)}

\vspace{0.5em}
\noindent\textbf{Sūtra (\dev{सूत्रम्}) — Verse 9:}

\begin{paracol}{2}
\backgroundcolor{c[0](.95\columnsep,.95\textheight)}{sanskritbg}
\backgroundcolor{c[1](.95\columnsep,.95\textheight)}{englishbg}

\begin{sanskritverse}
षष्ट्या विभाजिता स्थितिविमर्दार्धनाडिकागुणा स्वगतिः।
छाद्ये रवीन्दुपातेष्वुभयसकृत् तैर्धनमन्ते॥ ९ ॥
\end{sanskritverse}

\switchcolumn

\textbf{Translation:} ``The own motion (\textit{svagati}) --- [i.e.] the motions of the Sun, Moon, and node --- multiplied by the \textit{sthitivimardārdha} nāḍikās and divided by sixty; both [motions] are applied [to the Sun, Moon, and node] repeatedly in the eclipsed body (\textit{chādya}), [and] the result is additive or subtractive.''

\end{paracol}

\vspace{0.5em}

\begin{paracol}{2}
\backgroundcolor{c[0](.95\columnsep,.95\textheight)}{sanskritbg}
\backgroundcolor{c[1](.95\columnsep,.95\textheight)}{englishbg}

\noindent\textbf{Sūtrārtha (\dev{सूत्रार्थः}):}
\begin{sanskritverse}
स्वगतिः (रविगतिः, चन्द्रगतिः, पातगतिश्च) स्थितिविमर्दार्धनाडिका गुणा
(स्थित्यर्ध घटीमिविमर्दार्ध घटीभिः पृथक् पृथक् गुणिता) षष्ट्या विभाजिता छाद्ये
रविचन्द्रपातेषु उभयसकृत् तैर्धनमन्ते।
\end{sanskritverse}

\switchcolumn

\textbf{Literal Meaning:} ``The own motion --- the motion of the Sun, the motion of the Moon, and the motion of the node --- multiplied by the \textit{sthitivimardārdha} nāḍikās (the half-duration of totality ghaṭīs multiplied separately) and divided by sixty; applied to the Sun, Moon, and node in the eclipsed body, both [additive and subtractive] repeatedly, the result being additive or subtractive.''

\end{paracol}

\vspace{0.5em}

\begin{paracol}{2}
\backgroundcolor{c[0](.95\columnsep,.95\textheight)}{sanskritbg}
\backgroundcolor{c[1](.95\columnsep,.95\textheight)}{englishbg}

\noindent\textbf{Upapatti (\dev{उपपत्तिः}):}
\begin{sanskritverse}
``स्थित्यर्धनाडीगुणिता स्वमुक्ति'' इत्यादिना एवं विमर्दार्धफलोनयुक्त इत्यादिना च
भास्करविधिना स्फुटा ॥ ९ ॥
\end{sanskritverse}

\switchcolumn

\textbf{English Translation:} ``The [method is] made accurate by the method of Bhāskara: `\textit{sthit yardha-nāḍī-guṇitā sva-mukti}' etc., and `\textit{vimardārdha-phala-ona-yukta}' etc.''

\end{paracol}

\newpage


%% ========================================================================
%% CHAPTER 2: SŪRYAGRAHAṆĀDHIKĀRAḤ (Solar Eclipse Chapter)
%% ========================================================================
\section{\texorpdfstring{{\devanagarifont सूर्यग्रहणाधिकारः} — Sūryagrahaṇādhikāraḥ (Solar Eclipse Chapter)}{Sūryagrahaṇādhikāraḥ (Solar Eclipse Chapter)}}
\label{sec:surya-grahana}

\begin{center}
\textcolor{rulecolor}{\rule{0.8\textwidth}{0.5pt}}
\end{center}

This chapter (\dev{सूर्यग्रहणाधिकारः}, the fifth chapter of the text) deals with the computation of solar eclipses. The chapter begins with the statement of purpose (\dev{प्रयोजन}) and proceeds through the computation of the five \dev{ज्या} (sines) required for eclipse calculation: \dev{उदयज्या}, \dev{मध्यज्या}, \dev{रविदयाकर्ण}, \dev{विशिमशङ्कु}, and \dev{हृक्षेप}.

%% ========================================================================
\subsection{Verse 1: Statement of Purpose}

\vspace{0.5em}
\noindent\textbf{Sūtra (\dev{सूत्रम्}) — Verse 1:}

\begin{paracol}{2}
\backgroundcolor{c[0](.95\columnsep,.95\textheight)}{sanskritbg}
\backgroundcolor{c[1](.95\columnsep,.95\textheight)}{englishbg}

\begin{sanskritverse}
अथ सूर्यग्रहणाधिकारः प्रारभ्यते। तत्रादौ तदारम्भप्रयोजनमाह—
\end{sanskritverse}

\begin{sanskritverse}
ग्रहभास्करान्तरं प्राक् पद बाधकं ग्रहान्तरे यस्मात्।
स्वांशैर्यथाह्यास्तस्माद् द्युते तदानयनम्॥ १ ॥
\end{sanskritverse}

\switchcolumn

\textbf{Introduction:} ``Now the chapter on solar eclipse is begun. First, the purpose of its commencement is stated ---''

\textbf{Translation:} ``Because the inter-planetary [distance] computed by the \textit{grahabhāskara} [method] previously is an obstruction in the inter-planetary [computation], therefore I shall state the computation of \textit{dyuti} (parallax) with its own degrees, so that there may be agreement.''

\textit{Technical note:} The \textit{grahabhāskara} refers to the method given by earlier astronomers (like Bhāskara I) for computing the distance between planets, which Brahmagupta finds to be defective for solar eclipse computation.

\end{paracol}

\vspace{0.5em}

\begin{paracol}{2}
\backgroundcolor{c[0](.95\columnsep,.95\textheight)}{sanskritbg}
\backgroundcolor{c[1](.95\columnsep,.95\textheight)}{englishbg}

\noindent\textbf{Sūtrārtha (\dev{सूत्रार्थः}):}
\begin{sanskritverse}
प्राक् पूर्वे क्षितिजे प्रभास्करान्तरैः स्वांशैः स्वकालांशैः पश्चात् पश्चिमक्षितिजे
ग्रहास्तरैः स्वकालांशैर्यथासंख्यं ग्रहाः हश्याहश्याः भवन्ति। तस्मात् तदानयनं
प्रहर्कान्तरानयनं वक्ष्ये इति। अर्थात् प्राक्-क्षितिजे यस्मिन् दिने ग्रहोदयादनन्तरं
कालांशघटीभिः रवेरुदयस्तस्मिन् दिने ग्रहो रात्रिशेषे हश्यो भवति। एवं यस्मिन् दिने
पश्चिमक्षितिजे रव्यस्तानन्तरं कालांशघटीमित्रं हस्यास्तमयस्तस्मिन् दिने सायंकाले
ग्रहस्य हश्यत्वम्॥ १ ॥
\end{sanskritverse}

\switchcolumn

\textbf{Literal Meaning:} ``Previously, in the eastern horizon (\textit{pūrve kṣitije}), by the degrees of the \textit{grahabhāskara} method, the planets become visible or invisible (\textit{haśya-āhaśya}) in their own time-degrees. Later, in the western horizon (\textit{paścima-kṣitije}), by the degrees of the planetary distances, [the same] according to number. Therefore, I shall state that computation --- the computation of the \textit{praharāntara}. That is: in the eastern horizon, on whatever day after the rising of the planet, the rising of the Sun [occurs] after a certain number of time-ghaṭīs, on that day the planet becomes invisible (\textit{haśya}) in the remaining night. Similarly, in the western horizon, after the setting of the Sun, when the setting [occurs] after a certain number of time-ghaṭīs, on that day in the evening the planet's invisibility (\textit{haśyatva}) [is determined].''

\end{paracol}

\vspace{0.5em}

\begin{paracol}{2}
\backgroundcolor{c[0](.95\columnsep,.95\textheight)}{sanskritbg}
\backgroundcolor{c[1](.95\columnsep,.95\textheight)}{englishbg}

\noindent\textbf{Hindi Bhāṣya (\dev{हिन्दी भाष्य}):}
\begin{sanskritverse}
जिस कारण से सूर्यग्रहण में पञ्चज्या (उदयज्या, मध्यज्या, रविदयाकर्ण, विशिमशङ्कु,
हृक्षेप) से भास्करादि आचार्यों ने जो प्रानयन किया है उससे गणितैक्य (वेधोपलब्ध और
गणितागत विषयों की तुल्यता) नहीं होता है उस कारण से जैसे उनका ऐक्य (वेधोपलब्ध और
गणितागत विषयों के समत्व) हो वैसे मैं प्रानयन को कहता हूँ। तिथ्यन्ते इसका ग्रहे से
सम्बन्ध है, उदयज्या -- भ्रमा, मध्यज्या -- दशमलग्ननतांशज्या, विशिमशङ्कु -- हग्गति,
हृक्षेप -- विशिमनतांशज्या, सूर्यसिद्धान्त में भी पूर्वोक्त पञ्चज्या ही से सूर्यग्रहण का
प्रानयन है, पौलिशसिद्धान्त में सूर्यग्रहण में चन्द्र की पृथक् पञ्चज्या साधित है इसलिए
वहाँ द्विज्या विधान से सूर्यग्रहण उपनिबद्ध है इति ॥ १ ॥
\end{sanskritverse}

\switchcolumn

\textbf{English Translation:} ``For whatever reason, in solar eclipses, the five \textit{jyā}s (\textit{udayajyā}, \textit{madhyajyā}, \textit{ravidayākarṇa}, \textit{viśimaśaṅku}, \textit{hṛkṣepa}) --- by which earlier teachers like Bhāskara computed the \textit{prāṇayana} --- from that, the \textit{gaṇitaikya} (agreement between observed and computed values) is not obtained. For that reason, so that their agreement may exist, I state the \textit{prāṇayana} [method]. The \textit{tithi-anta} [refers to] the relation with the \textit{graha}: \textit{udayajyā} [refers to] parallax (\textit{bhramā}), \textit{madhyajyā} [to] the sine of the zenith-distance of the tenth degree (\textit{daśama-lagna-natāṃśajyā}), \textit{viśimaśaṅku} [to] the ascensional difference (\textit{hagati}), \textit{hṛkṣepa} [to] the sine of the zenith-distance of the nonagesimal (\textit{viśima-natāṃśajyā}). In the \textit{Sūryasiddhānta} too, the solar eclipse \textit{prāṇayana} is by the aforesaid five \textit{jyā}s. In the \textit{Pauliśasiddhānta}, in solar eclipses, separate five \textit{jyā}s of the Moon are computed; therefore there the solar eclipse is accomplished by the two-\textit{jyā} method.''

\end{paracol}

%% ========================================================================
\subsection{Verses 2--3: Conditions for \texorpdfstring{{\devanagarifont लम्बन}}{lambana} (Deflection) and \texorpdfstring{{\devanagarifont अवनति}}{avanati} (Depression)}

\vspace{0.5em}
\noindent\textbf{Sūtra (\dev{सूत्रम्}) — Verses 2--3:}

\begin{paracol}{2}
\backgroundcolor{c[0](.95\columnsep,.95\textheight)}{sanskritbg}
\backgroundcolor{c[1](.95\columnsep,.95\textheight)}{englishbg}

\begin{sanskritverse}
विश्रिमलग्नसमं एकं न लम्बनं तदधिकोनके भवति।
तस्य क्रान्तिज्योदक् यदा अक्षनोवा समा न तदा॥ २ ॥
\end{sanskritverse}

\begin{sanskritverse}
अवनतिरतोऽन्यथा भवति सम्भवे तदुदयैव लग्नसमम्।
तदुदितघटिकास्तच्छङ्कुस्तच्चरप्राणैः परे॥ ३ ॥
\end{sanskritverse}

\switchcolumn

\textbf{Translation — Verse 2:} ``When one [planet] is equal to the \textit{viśrima-lagna} (nonagesimal ascendant), there is no \textit{lambana} (deflection in longitude). When it is greater or less than that, \textit{lambana} occurs. When its declination-sine (\textit{krāntijyā}) is equal to the latitude-sine (\textit{akṣajyā}), then there is no \textit{avanati} (depression).''

\textbf{Translation — Verse 3:} ``Otherwise, \textit{avanati} occurs. When both \textit{lambana} and \textit{avanati} are possible, then the rising [point] equal to the \textit{lagna}, the \textit{ghaṭikās} elapsed since its rising, its \textit{śaṅku} (gnomon), and the \textit{prāṇa}s of its motion --- [are computed] thereafter.''

\end{paracol}

\vspace{0.5em}

\begin{paracol}{2}
\backgroundcolor{c[0](.95\columnsep,.95\textheight)}{sanskritbg}
\backgroundcolor{c[1](.95\columnsep,.95\textheight)}{englishbg}

\noindent\textbf{Sūtrārtha (\dev{सूत्रार्थः}):}
\begin{sanskritverse}
तिथ्यन्ते गणितागतदर्शान्तिकाले पूर्वविधिना त्रिप्रश्नोक्तेन लग्नं कृत्वा विश्रिमं
कार्यम्। तेन विश्रिमलग्नेन समेकं रवौ सति लम्बनं न भवति! तस्माद् विश्रिमादधिके
कोने रवौ लम्बनं भवतीत्यर्थत् एव सिध्यति। एवं तस्य विश्रिमस्योत्तरा क्रान्तिज्या यदा
स्वदेशाक्षजोवा समा तदा अवनतिर्न भवति अतोऽन्यथा चावनतिर्भवति।
लम्बनावनत्योः सम्भवे च तदुदयं स्वदेशीयराश्युदयैषि-लग्नसमं कृत्वा
आद्यूनस्य भोग्योद्धिकभुक्तियुक्तो मध्योदयाद्य इति त्रिप्रश्नोक्तविधिं कृत्वा
तदुदितघटिकास्तस्य विश्रिमस्य गता घटिकाः साध्याः॥ २--३ ॥
\end{sanskritverse}

\switchcolumn

\textbf{Literal Meaning:} ``At the end of the \textit{tithi}, at the time of computed conjunction (\textit{gaṇitāgata-darśāntikāle}), having computed the \textit{lagna} by the previous method stated in the \textit{Tripraśna}, the \textit{viśrima} [nonagesimal] is to be computed. When the Sun is equal to one [point] with that \textit{viśrimalagna}, there is no \textit{lambana}. When the Sun is greater or less than that \textit{viśrima}, \textit{lambana} occurs --- this meaning is self-evident. Similarly, when the northern declination-sine (\textit{krāntijyā}) of that \textit{viśrima} is equal to the latitude-sine (\textit{akṣajyā}) of one's own place, then there is no \textit{avanati}; otherwise \textit{avanati} occurs. When both \textit{lambana} and \textit{avanati} are possible, then having made the rising equal to the \textit{lagna} --- the \textit{rāśyudaya-eṣi-lagna} of one's own place --- by the method stated in the \textit{Tripraśna}: `the first diminished by the \textit{bhoga}, increased by the \textit{bhukti}, [gives] the middle \textit{udaya}', etc., the \textit{ghaṭikās} elapsed since the rising of that \textit{viśrima} are to be computed.''

\end{paracol}

%% ========================================================================
\subsection{Verse 7: {\devanagarifont स्पष्टदर्शान्तकालिक} Computation (Computation of True Conjunction Time)}

\vspace{0.5em}
\noindent\textbf{Sūtra ({\devanagarifont सूत्रम्}) — Verse 7:}

\begin{paracol}{2}
\backgroundcolor{c[0](.95\columnsep,.95\textheight)}{sanskritbg}
\backgroundcolor{c[1](.95\columnsep,.95\textheight)}{englishbg}

\begin{sanskritverse}
घटिकामिरगुणिताः पञ्चभिर्भक्ता लब्धकलामियंति लम्बनम्।
तदा गणितागतदर्शान्तकालिका रविचन्द्रपाता ऊनाः कार्याः॥ ७ ॥
\end{sanskritverse}

\switchcolumn

\textbf{Translation:} ``The \textit{ghaṭikā-miṣa} [of \textit{lambana}] multiplied by five and divided [by some number] yields the \textit{labdha-kalā} called \textit{lambana}. Then the Sun, Moon, and node at the computed conjunction time (\textit{gaṇitāgata-darśānta-kālika}) are to be diminished [or increased].''

\end{paracol}

\vspace{0.5em}

\begin{paracol}{2}
\backgroundcolor{c[0](.95\columnsep,.95\textheight)}{sanskritbg}
\backgroundcolor{c[1](.95\columnsep,.95\textheight)}{englishbg}

\noindent\textbf{Sūtrārtha (\dev{सूत्रार्थः}):}
\begin{sanskritverse}
घटिकामिरगुणिताः पञ्चभ्यः पञ्चज्या भक्ता लब्धकलामियंति लम्बनं तदा
गणितागतदर्शान्तकालिका रविचन्द्रपाता ऊनाः (हीनाः) कार्याः, यदि लम्बनं धनं तदा
लब्धकलाभिर्गणितदर्शान्तकालिका रविचन्द्रपाता युक्ताः कार्यास्तदा स्पष्टदर्शान्तकालिका
रविचन्द्रपाता भवन्तीति ॥ ७ ॥
\end{sanskritverse}

\switchcolumn

\textbf{Literal Meaning:} ``The \textit{ghaṭikā-miṣa} multiplied by five, divided by five [?\ \textit{pañcajyā}], the result is the \textit{labdha-kalā} called \textit{lambana}. Then the Sun, Moon, and node at the computed conjunction time are to be diminished (\textit{ūnāḥ}). If the \textit{lambana} is positive (\textit{dhana}), then the Sun, Moon, and node at the computed conjunction time are to be increased by the \textit{labdha-kalā}; then those become the true Sun, Moon, and node at the true conjunction time (\textit{spaṣṭa-darśānta-kālika}).''

\end{paracol}

\vspace{0.5em}

\begin{paracol}{2}
\backgroundcolor{c[0](.95\columnsep,.95\textheight)}{sanskritbg}
\backgroundcolor{c[1](.95\columnsep,.95\textheight)}{englishbg}

\noindent\textbf{Upapatti (\dev{उपपत्तिः}):}
\begin{sanskritverse}
गणितागतदर्शान्तकालिकरविचन्द्रपातेभ्यः शरादिकं सर्वमानेतव्यम्। ऋणात्मके
लम्बने ``पञ्चषष्टिकामिरः पृथक् पृथक् रविचन्द्रपातगतिकला लम्यन्ते तदा
लम्बनघटिकाभिः किमित्यनुपातेन'' लब्धकलाभिर्गणितागतदर्शान्तकालिका
रविचन्द्रपाता हीनााः कार्याः, धनात्मके लम्बने लब्धकलाभिस्ते युक्तास्तदा ते
स्पष्टदर्शान्तकालिका भवेयुरेवेति, सिद्धान्तशेखरे ``गतिहतां लम्बननाडिकामिर्विभज्य
पञ्चया फललिप्तिकामिः। रवीन्दुपाताः सहिता धनाख्ये विवर्जिताश्च क्षयलम्बने ते''
इत्यनेन श्रीपतिनाचार्योक्तानुरूपमेवोक्तमिति ॥ ७ ॥
\end{sanskritverse}

\switchcolumn

\textbf{English Translation of Proof:} ``From the Sun, Moon, and node at the computed conjunction time, all [quantities] beginning with the \textit{śara} [latitude] are to be computed. When the \textit{lambana} is negative (\textit{ṛṇātmaka}): `The minutes of motion of the Sun, Moon, and node are separately multiplied by the \textit{lambana}-ghaṭikā-miṣa} of sixty-five --- by proportion: `then by the \textit{lambana}-ghaṭikā, what [result]?}' --- by the \textit{labdha-kalā}s, the Sun, Moon, and node at the computed conjunction time are to be diminished. When the \textit{lambana} is positive (\textit{dhanātmaka}), they are increased by the \textit{labdha-kalā}s; then they become the true [quantities] at the true conjunction time. In the \textit{Siddhāntaśekhara}: `Having divided the motion-multiplied \textit{lambana}-nāḍī-miṣa by five, the resulting minutes... The Sun, Moon, and node are added in positive [\textit{lambana}] and subtracted in negative \textit{lambana}' --- by this, what was stated by Śrīpati-nācārya is here spoken in accordance.''

\end{paracol}

\newpage


%% ========================================================================
%% CHAPTER 3: CHANDRACCHĀYĀDHIKĀRAḤ (Moon's Shadow Chapter)
%% ========================================================================
\section{\texorpdfstring{{\devanagarifont चन्द्रच्छायाधिकारः} — Candracchāyādhikāraḥ (Moon's Shadow Chapter)}{Candracchāyādhikāraḥ (Moon's Shadow Chapter)}}
\label{sec:chandra-chhaya}

\begin{center}
\textcolor{rulecolor}{\rule{0.8\textwidth}{0.5pt}}
\end{center}

This chapter (\dev{चन्द्रच्छायाधिकारः}) deals with the computation of the Moon's shadow, which is essential for understanding lunar phenomena. The shadow of the Earth falls on the Moon during a lunar eclipse, and its computation involves calculating the \dev{च्छाया} (shadow), \dev{मध्यच्छाया} (middle shadow), and related quantities.

%% ========================================================================
\subsection{Verse 1: Beginning of the Chapter}

\vspace{0.5em}
\noindent\textbf{Sūtra (\dev{सूत्रम्}) — Verse 1:}

\begin{paracol}{2}
\backgroundcolor{c[0](.95\columnsep,.95\textheight)}{sanskritbg}
\backgroundcolor{c[1](.95\columnsep,.95\textheight)}{englishbg}

\begin{sanskritverse}
अथ चन्द्रच्छायाधिकारः प्रारभ्यते। तत्रादौ तदारम्भप्रयोजनमाह—
\end{sanskritverse}

\begin{sanskritverse}
प्राक् चन्द्रलग्नयोरन्तरं सप्तलग्नादीनां यतस्ततं पश्चात्।
प्रतिदिनमिन्दोश्च्छाया यतस्तदनयनमभिचार्ये॥ १ ॥
\end{sanskritverse}

\switchcolumn

\textbf{Introduction:} ``Now the chapter on the Moon's shadow is begun. First, the purpose of its commencement is stated ---''

\textbf{Translation:} ``Previously, the difference between the Moon and the \textit{lagna}, and afterwards, that of the \textit{saptalagna} etc., because from that the shadow of the Moon (\textit{indu-ccchāyā}) is daily [computed], therefore I shall state the computation (\textit{anayana}) of that.''

\end{paracol}

\vspace{0.5em}

\begin{paracol}{2}
\backgroundcolor{c[0](.95\columnsep,.95\textheight)}{sanskritbg}
\backgroundcolor{c[1](.95\columnsep,.95\textheight)}{englishbg}

\noindent\textbf{Sūtrārtha (\dev{सूत्रार्थः}):}
\begin{sanskritverse}
प्राक् (पूर्वस्थितिजे) चन्द्रलग्नयोरन्तरं, पश्चात् (पश्चिमस्थितिजे) सप्तलग्नचन्द्रयोः
अन्तरं च प्रतिदिनं यतनेन यतः (यस्मात्कारणात्) इन्दुच्छाया (चन्द्रच्छाया) ज्ञायते,
तथेन्दुच्छायातस्तदनन्तरं च ज्ञायते ततस्तदानयन (छायानयन) मभिचार्ये (कथयिष्ये) इति।
सिद्धान्तशेखरे ``प्रायतस्तु हीनरविमलग्नयोरन्तरलग्ननाशिनोश्च पश्चिमे'' इत्यादि
इनेन श्रीपतिनाचार्योक्तानुरूपमेवोक्तमिति ॥ १ ॥
\end{sanskritverse}

\switchcolumn

\textbf{Literal Meaning:} ``Previously (in the eastern horizon), the difference between the Moon and the \textit{lagna}, and afterwards (in the western horizon), the difference between the \textit{saptalagna} and the Moon, and daily by effort --- by which cause the Moon's shadow (\textit{indu-ccchāyā}) is known, and from the Moon's shadow that difference is [also] known --- therefore I shall state that computation (\textit{anayana}), the shadow-computation (\textit{chāyā-anayana}). In the \textit{Siddhāntaśekhara}: `Usually the difference between the diminished-Sun-\textit{lagna} and the \textit{lagna-nāśin} in the west' etc. --- by this, what was spoken by Śrīpati-nācārya is here spoken in the same manner.''

\end{paracol}

\vspace{0.5em}

\begin{paracol}{2}
\backgroundcolor{c[0](.95\columnsep,.95\textheight)}{sanskritbg}
\backgroundcolor{c[1](.95\columnsep,.95\textheight)}{englishbg}

\noindent\textbf{Hindi Bhāṣya (\dev{हिन्दी भाष्य}):}
\begin{sanskritverse}
अब चन्द्रच्छायाधिकार प्रारम्भ किया जाता है, पहले उसके प्रारम्भ-प्रयोजन को कहते हैं।
पूर्व स्थितिजे में चन्द्र और लग्न का अन्तर, पश्चिम स्थितिजे में सप्तलग्न और चन्द्र का अन्तर
प्रत्येक दिन में जो होता है उसी से क्योंकि चन्द्रच्छाया समझी जाती है (प्रथम चन्द्रच्छाया
का ज्ञान होता है) तथा चन्द्रच्छाया से उस अन्तर का ज्ञान होता है इसलिये छायानयन को
कहता हूँ॥ १ ॥
\end{sanskritverse}

\switchcolumn

\textbf{English Translation:} ``Now the chapter on the Moon's shadow is begun. First, its purpose is stated. The difference between the Moon and the \textit{lagna} in the eastern horizon, and the difference between the \textit{saptalagna} and the Moon in the western horizon --- because from these, daily, the Moon's shadow is known (first, the knowledge of the Moon's shadow is obtained), and from the Moon's shadow the knowledge of that difference is obtained --- therefore I state the shadow-computation.''

\end{paracol}

%% ========================================================================
\subsection{Verse 2: Shadow Computation}

\vspace{0.5em}
\noindent\textbf{Sūtra (\dev{सूत्रम्}) — Verse 2:}

\begin{paracol}{2}
\backgroundcolor{c[0](.95\columnsep,.95\textheight)}{sanskritbg}
\backgroundcolor{c[1](.95\columnsep,.95\textheight)}{englishbg}

\begin{sanskritverse}
प्रभहुण्डाकृतिकोर्ध्वकालिकोद्बलनखण्डसमस्थितिपाताः।
छन्दोव्यासस्तलग्ने स्वभग्रप्रायाण्न सङ्क्षेपः॥ २ ॥
\end{sanskritverse}

\switchcolumn

\textbf{Translation:} ``The parts of the upright-time-difference (\textit{ūrdhvakālika-udbala-naṣṭa-sama-sthiti-pāta}) [having] the form of a \textit{prabhahṛṇḍa} [spherical triangle], the semi-diameter (\textit{vyāsa}) of the shadow at the horizon (\textit{tala-lagna}), the proximity to its own \textit{agra} --- there is no abbreviation (\textit{saṅkṣepaḥ}).''

\textit{Note:} This is a technical verse whose exact interpretation requires the accompanying commentary. The \textit{prabhahṛṇḍa} is a type of spherical triangle used in Indian astronomy.

\end{paracol}

%% ========================================================================
\subsection{Verses 3--4: Detailed Shadow Calculation}

\vspace{0.5em}
\noindent\textbf{Sūtra (\dev{सूत्रम्}) — Verses 3--4:}

\begin{paracol}{2}
\backgroundcolor{c[0](.95\columnsep,.95\textheight)}{sanskritbg}
\backgroundcolor{c[1](.95\columnsep,.95\textheight)}{englishbg}

\begin{sanskritverse}
इदानीं कर्तव्यतामाह—
\end{sanskritverse}

\begin{sanskritverse}
प्राग्रहुण्डाकृतिकोर्ध्वकालिकोद्बलनखण्डसमस्थितिपाताः।
छन्दोव्यासस्तलग्ने स्वभग्रप्रायाण्न सङ्क्षेपः॥ ३ ॥
\end{sanskritverse}

\begin{sanskritverse}
मूलोनं मूलं गतं यद्वदुद्धृत्यागतं ततोऽन्यत्।
स्वभुक्त्या हृतमभीष्टच्छायागतिः सैव मूलोद्धृता॥ ४ ॥
\end{sanskritverse}

\switchcolumn

\textbf{Introduction to Verse 3:} ``Now [the author] states what is to be done ---''

\textbf{Translation — Verse 3:} [Same structure as verse 2, establishing the method.]

\textbf{Translation — Verse 4:} ``As the root (\textit{mūla}) diminished by the root, or the elapsed [time], having been extracted, another [quantity] comes therefrom. Divided by the own motion (\textit{svabhakti}), that is the desired shadow-motion (\textit{abhīṣṭa-ccchāyā-gati}). That same [quantity], extracted from the root (\textit{mūlodḍhṛtā}).''

\textit{Note:} This verse describes an iterative root-extraction method for computing the rate of change of the shadow. The \textit{mūla} here refers to a base quantity from which differences are taken.

\end{paracol}

\newpage


%% ========================================================================
%% CHAPTER 4: GRAHODAYĀSTĀDHIKĀRAḤ (Planetary Rising/Setting Chapter)
%% ========================================================================
\section{\texorpdfstring{{\devanagarifont ग्रहोदयास्ताधिकारः} — Grahodayāstādhikāraḥ (Planetary Rising and Setting)}{Grahodayāstādhikāraḥ (Planetary Rising and Setting)}}
\label{sec:grahodayasta}

\begin{center}
\textcolor{rulecolor}{\rule{0.8\textwidth}{0.5pt}}
\end{center}

This chapter (\dev{ग्रहोदयास्ताधिकारः}) deals with the computation of the rising and setting times of planets. The chapter explains how to determine the times when planets become visible (\dev{उदय}, heliacal rising) and invisible (\dev{अस्त}, heliacal setting) based on their elongation from the Sun.

%% ========================================================================
\subsection{Verse 1: Pratyaya (Introduction)}

\vspace{0.5em}
\noindent\textbf{Sūtra (\dev{सूत्रम्}) — Verse 1:}

\begin{paracol}{2}
\backgroundcolor{c[0](.95\columnsep,.95\textheight)}{sanskritbg}
\backgroundcolor{c[1](.95\columnsep,.95\textheight)}{englishbg}

\begin{sanskritverse}
अथोदयास्ताधिकारः प्रारभ्यते। तत्रादौ तदारम्भप्रयोजनमाह—
\end{sanskritverse}

\begin{sanskritverse}
ग्रहभास्करान्तरं प्राक् पद बाधकं ग्रहान्तरे यस्मात्।
स्वांशैर्यथाह्यास्तस्माद् द्युते तदानयनम्॥ १ ॥
\end{sanskritverse}

\switchcolumn

\textbf{Introduction:} ``Now the chapter on planetary rising and setting is begun. First, the purpose of its commencement is stated ---''

\textbf{Translation:} ``Because previously, in the inter-planetary computation, the \textit{grahabhāskara} distance is an obstruction (\textit{bādhaka}), therefore, with its own degrees, so that there may be agreement --- I state the computation of \textit{dyuti}.''

\textit{Note:} This verse serves as a parallel to the solar eclipse chapter's opening, establishing that Brahmagupta's method supersedes earlier approaches that failed to achieve agreement (\textit{aiṃya}) between computation and observation.

\end{paracol}

\vspace{0.5em}

\begin{paracol}{2}
\backgroundcolor{c[0](.95\columnsep,.95\textheight)}{sanskritbg}
\backgroundcolor{c[1](.95\columnsep,.95\textheight)}{englishbg}

\noindent\textbf{Sūtrārtha (\dev{सूत्रार्थः}):}
\begin{sanskritverse}
प्राक् पूर्वे क्षितिजे प्रभास्करान्तरैः स्वांशैः स्वकालांशैः पश्चात् पश्चिमक्षितिजे
ग्रहास्तरैः स्वकालांशैर्यथासंख्यं ग्रहाः हश्याहश्याः भवन्ति। तस्मात् तदानयनं
प्रहर्कान्तरानयनं वक्ष्य इति। अर्थात् प्राक्-क्षितिजे यस्मिन् दिने ग्रहोदयादनन्तरं
कालांशघटीभिः रवेरुदयस्तस्मिन् दिने प्रभो रात्रिशेषे हश्यो भवति। एवं यस्मिन् दिने
पश्चिमक्षितिजे रव्यस्तानन्तरं कालांशघटीमित्रं हस्यास्तमयस्तस्मिन् दिने सायंकाले
ग्रहस्य हश्यत्वम्॥ १ ॥
\end{sanskritverse}

\switchcolumn

\textbf{Literal Meaning:} ``Previously, in the eastern horizon, by the \textit{grahabhāskara} distances in their own degrees, in their own time-degrees, later in the western horizon, by the planetary distances in their own time-degrees, according to number, the planets become visible or invisible (\textit{haśya-āhaśya}). Therefore, I shall state that computation --- the computation of the \textit{praharāntara}. That is: in the eastern horizon, on whatever day after the rising of the planet, the rising of the Sun occurs after a certain number of time-ghaṭīs, on that day the planet becomes invisible (\textit{haśya}) in the remaining night. Similarly, in the western horizon, after the setting of the Sun, when the setting occurs after a certain number of time-ghaṭīs, on that day in the evening the planet's invisibility (\textit{haśyatva}) [is established].''

\end{paracol}

%% ========================================================================
\subsection{Verses 8--9: Daily Rising and Setting}

\vspace{0.5em}
\noindent\textbf{Sūtra (\dev{सूत्रम्}) — Verses 8--9:}

\begin{paracol}{2}
\backgroundcolor{c[0](.95\columnsep,.95\textheight)}{sanskritbg}
\backgroundcolor{c[1](.95\columnsep,.95\textheight)}{englishbg}

\begin{sanskritverse}
तात्कालिक ग्रह लग्न से असकृत् कर्म द्वारा प्रतिदिन उदय और अस्त साधन करना
प्रथम जिस समय ग्रहोदित घटी ज्ञात करना हो उस समय के लग्न (तात्कालिक लग्न)
तथा ग्रहोदय लग्न साधन कर दोनों के अन्तरघटी साधन कर उन (साधित घटी) से ग्रह
को चलाकर उसके फिर उदयलग्न साधन कर तात्कालिक स्थिर लग्न से पुनः इष्टघटी
साधन करना, इसतरह असकृत् (वारंवार) करना॥ ८ ॥
\end{sanskritverse}

\switchcolumn

\textbf{Translation — Verse 8:} ``By iterative computation (\textit{asakṛt-karma}) from the instantaneous (\textit{tātkālika}) planet-\textit{lagna}, the daily rising and setting are to be computed. First, at whatever time the \textit{grahodita-ghaṭī} [elapsed ghaṭīs since planetary rising] is to be known, the \textit{lagna} at that time (\textit{tātkālika lagna}) and the planetary rising-\textit{lagna} are to be computed. Computing the difference-ghaṭīs of both, by those (computed ghaṭīs) moving the planet, again computing its rising-\textit{lagna}, from the instantaneous fixed \textit{lagna} again computing the desired ghaṭīs --- in this manner, iteratively (\textit{asakṛt}, repeatedly) doing [this].''

\textit{Note:} This verse describes the iterative procedure (\textit{asakṛt-karma}) for refining the computation of a planet's rising time. The method involves successive approximation until the computed value stabilizes.

\end{paracol}

%% ========================================================================
\subsection{Verse 9: Moon's Rising and Setting Computation}

\vspace{0.5em}
\noindent\textbf{Sūtra (\dev{सूत्रम्}) — Verse 9:}

\begin{paracol}{2}
\backgroundcolor{c[0](.95\columnsep,.95\textheight)}{sanskritbg}
\backgroundcolor{c[1](.95\columnsep,.95\textheight)}{englishbg}

\begin{sanskritverse}
उदयास्तमयाविन्दोः कालांशैरंकसंमितैः कार्यः।
हीनत्वं त्वधिकत्वं तदन्तरे योगकालः स्यात्॥ ९ ॥
\end{sanskritverse}

\switchcolumn

\textbf{Translation:} ``The rising and setting of the Moon (\textit{indu}) are to be computed by time-degrees (\textit{kālāṃśa}) equal to digits (\textit{aṅka-saṃmitaiḥ}). If there is diminution (\textit{hīnatva}) or excess (\textit{adhikatva}), in that interval the \textit{yoga-kāla} [conjunction time] would be [obtained].''

\textit{Note:} The \textit{aṅka} (digit) is a unit equal to 12 degrees. The verse states that the Moon's rising and setting are computed using time-degrees measured in these units.

\end{paracol}

\vspace{0.5em}

\begin{paracol}{2}
\backgroundcolor{c[0](.95\columnsep,.95\textheight)}{sanskritbg}
\backgroundcolor{c[1](.95\columnsep,.95\textheight)}{englishbg}

\noindent\textbf{Sūtrārtha (\dev{सूत्रार्थः}):}
\begin{sanskritverse}
अंकसंमितैः कालांशैररात् द्वादशकालांशैरिन्दोश्चन्द्रस्योदयास्तमयौ साध्यौ।
अत्र पठितकालांशेन इष्टकालांशातां यदि हीनत्वं वाधिकत्वं भवेत् तदा तदन्तरे तयोः
पठितकालांशयोरन्तरे योगकालो योगवत्‌ वा योगवत्‌ कालः स्यादिति सप्तमश्लोकेन
स्फुटोक्तम्॥ ९ ॥
\end{sanskritverse}

\switchcolumn

\textbf{Literal Meaning:} ``By time-degrees equal to digits (\textit{aṅka}), that is, by twelve time-degrees (\textit{dvādaśa-kālāṃśa}), the rising and setting of the Moon (\textit{indu} = \textit{candra}) are to be computed. Here, if the computed time-degree (\textit{paṭhita-kālāṃśa}) has diminution or excess from the desired time-degree (\textit{iṣṭa-kālāṃśa}), then in that interval, between those two computed time-degrees, the \textit{yoga-kāla} (conjunction time), or \textit{yogavat-kāla} (time like conjunction), would be [obtained], as clearly stated in the seventh verse.''

\end{paracol}

\vspace{0.5em}

\begin{paracol}{2}
\backgroundcolor{c[0](.95\columnsep,.95\textheight)}{sanskritbg}
\backgroundcolor{c[1](.95\columnsep,.95\textheight)}{englishbg}

\noindent\textbf{Hindi Bhāṣya (\dev{हिन्दी भाष्य}):}
\begin{sanskritverse}
अंकसंमितैः कालांशैररथात् द्वादशतुल्यैः कालांशैरुदयास्तमयौ साध्यो। पठितकालांशस्य
इष्टकालांशस्य यदि हीनत्वं (अल्पत्वं) वा अधिकत्वं भवेत् तदन्तरे तयोः कालांशयोरन्तरे
योगकालः स्यात्॥ ९ ॥
\end{sanskritverse}

\switchcolumn

\textbf{English Translation:} ``By time-degrees equal to digits (\textit{aṅka}), that is, by time-degrees equal to twelve, the rising and setting are to be computed. If the computed time-degree has diminution (\textit{hīnatva}, being smaller) or excess (\textit{adhikatva}) from the desired time-degree, then in that interval, between those two time-degrees, the \textit{yoga-kāla} (conjunction time) would be [obtained].''

\end{paracol}

\newpage


%% ========================================================================
%% CHAPTER 5: ŚUKLONNATYADHIKĀRAḤ (Moon's Brightness Chapter)
%% ========================================================================
\section{\texorpdfstring{{\devanagarifont शुक्लोन्नत्यधिकारः} — Śuklonnatyadhikāraḥ (The Brightness of the Moon)}{Śuklonnatyadhikāraḥ (The Brightness of the Moon)}}
\label{sec:shuklonnati}

\begin{center}
\textcolor{rulecolor}{\rule{0.8\textwidth}{0.5pt}}
\end{center}

This chapter (\dev{शुक्लोन्नत्यधिकारः}) deals with the computation of the illuminated portion of the Moon's disk --- the \dev{शुक्ल} (bright) portion. The chapter explains how to determine the amount of the Moon's surface that is illuminated by the Sun at any given time.

%% ========================================================================
\subsection{Verses 1--3: Sphuṭaśuklonnati Computation}

\vspace{0.5em}
\noindent\textbf{Sūtra (\dev{सूत्रम्}) — Verse 4:}

\begin{paracol}{2}
\backgroundcolor{c[0](.95\columnsep,.95\textheight)}{sanskritbg}
\backgroundcolor{c[1](.95\columnsep,.95\textheight)}{englishbg}

\begin{sanskritverse}
अथ शुक्लोन्नत्यधिकारः प्रारभ्यते।
\end{sanskritverse}

\begin{sanskritverse}
रविचन्द्रपातलग्नैः स्वक्रान्त्युदयात् स्वलग्नगतशेषाः।
घटिकाः खचरार्धास्तात् स्वेष्टी रविशीतथू कृत्वा॥ ४ ॥
\end{sanskritverse}

\switchcolumn

\textbf{Introduction:} ``Now the chapter on the Moon's bright elevation (\textit{śuklonnati}) is begun.''

\textbf{Translation — Verse 4:} ``The ghaṭikās [which are] the \textit{khacarārdha} (half of the celestial sphere's motion) --- by the Sun, Moon, node, and \textit{lagna}, and from the rising of the own declination (\textit{sva-krānty-udaya}), [which are] the remaining [ghaṭikās] elapsed from the own \textit{lagna} --- having made [the computation] from the setting of the Sun, [these are] to be computed.''

\textit{Note:} The \textit{khacarārdha} refers to half the diurnal motion of the celestial sphere. The verse establishes the parameters needed to compute the Moon's illumination.

\end{paracol}

\vspace{0.5em}

\begin{paracol}{2}
\backgroundcolor{c[0](.95\columnsep,.95\textheight)}{sanskritbg}
\backgroundcolor{c[1](.95\columnsep,.95\textheight)}{englishbg}

\noindent\textbf{Sūtrārtha (\dev{सूत्रार्थः}):}
\begin{sanskritverse}
खचरार्धास्ताद् रव्यस्तात् स्वलग्नगतशेषाः शशिन उदयलग्नस्य गता दोषा वा
घटिकाः साध्याः। कैः। रविचन्द्रपातलग्नैस्तथा स्वक्रान्त्युदयात्। अत्रैतदुक्तं भवति।
रव्यस्तकाले रविचन्द्रः पातो लग्नम्। एतानि कृत्वा स्वक्रान्त्या चन्द्रक्रान्त्या
उदयात् चन्द्रोदयलग्नाज्ञोदयलग्नस्य गता वा दोषा घटिका ऊनस्य भोग्योद्धिकभुक्तियुक्तो
मध्योदयाद्य इत्यनेन विधिना साध्याः। सूर्यास्तानन्तरं यावतीभिः षटिकाभिश्चन्द्रास्तस्ताः
गता घटिकाः। सूर्यास्तात् प्राग् यावतीष्ठिकाभिश्चन्द्रास्तस्ता एषा घटिकाः। एवं प्राक्
क्षितिजे रव्युदयात् गता एषा वा चन्द्रोदयघटिकाः साध्या इत्यर्थादवगम्यते॥ ४ ॥
\end{sanskritverse}

\switchcolumn

\textbf{Literal Meaning:} ``The \textit{khacarārdha} ghaṭikās [which are] from the setting of the Sun (\textit{ravyastāt}), [which are] the remaining [ghaṭikās] elapsed from the own-\textit{lagna} of the Moon --- or the elapsed/deficient ghaṭikās of the Moon's rising-\textit{lagna} --- are to be computed. By what? By the Sun, Moon, node, \textit{lagna}, and from the rising of the own declination. Here, this is stated: at the time of the Sun's setting, the Sun, Moon, node, and \textit{lagna}. Having computed these, by the own declination (\textit{sva-krānti}), by the Moon's declination (\textit{candra-krānti}), from the rising, [computing] the elapsed or deficient ghaṭikās of the Moon's rising-\textit{lagna} from the \textit{udaya-lagna} by the method `diminished by the \textit{bhoga}, increased by the \textit{bhukti}, [gives] the middle \textit{udaya}', etc. After the Sun's setting, the ghaṭikās by which the Moon sets --- those are the elapsed ghaṭikās. Before the Sun's setting, the ghaṭikās by which the Moon sets --- those are the deficient ghaṭikās. Similarly, in the eastern horizon, from the Sun's rising, the elapsed or deficient ghaṭikās of the Moon's rising are to be computed --- this is to be understood from the meaning.''

\end{paracol}

\vspace{0.5em}

\begin{paracol}{2}
\backgroundcolor{c[0](.95\columnsep,.95\textheight)}{sanskritbg}
\backgroundcolor{c[1](.95\columnsep,.95\textheight)}{englishbg}

\noindent\textbf{Hindi Bhāṣya (\dev{हिन्दी भाष्य}):}
\begin{sanskritverse}
तात्कालिक ग्रह लग्न से असकृत् कर्म द्वारा प्रतिदिन उदय और अस्त साधन करना प्रथम
जिस समय ग्रहोदित घटी ज्ञात करना हो उस समय के लग्न (तात्कालिक लग्न) तथा
ग्रहोदय लग्न साधन कर दोनों के अन्तरघटी साधन कर उन (साधित घटी) से ग्रह को
चलाकर उसके फिर उदयलग्न साधन कर तात्कालिक स्थिर लग्न से पुनः इष्टघटी
साधन करना, इसतरह असकृत् (वारंवार) करना, एवं भस्तकाल ज्ञान के लिये भस्तलग्न (सप्तभस्त
लग्न) से कर्म करना चाहिये। इस तरह प्रथम चन्द्र दर्शन दिन से पूर्णिमा पर्यन्त इष्ट
दिन में चन्द्र के सूर्यास्त से त्रा सूर्योदय से अन्तमय काल और उदयिक काल साधन करना॥ ४ ॥
\end{sanskritverse}

\switchcolumn

\textbf{English Translation:} ``By iterative computation (\textit{asakṛt-karma}) from the instantaneous planet-\textit{lagna}, the daily rising and setting are to be computed. First, at whatever time the planetary elapsed ghaṭīs are to be known, the \textit{lagna} at that time (\textit{tātkālika lagna}) and the planetary rising-\textit{lagna} are to be computed. Computing the difference-ghaṭīs of both, by those (computed ghaṭīs) moving the planet, again computing its rising-\textit{lagna}, from the instantaneous fixed \textit{lagna} again computing the desired ghaṭīs --- in this manner, iteratively (\textit{asakṛt}, repeatedly) doing [this]. And for the knowledge of the \textit{sthita-kāla}, computation from the \textit{sthita-lagna} (\textit{sapta-bhastala-lagna}) should be done. In this way, from the first day of lunar visibility (\textit{candra darśana}) up to the full moon (\textit{pūrṇimā}), on the desired day, the setting time from the Sun's setting and the rising time from the Sun's rising are to be computed.''

\end{paracol}

%% ========================================================================
\subsection{Verse 5: The {\devanagarifont शुक्लोन्नति} (Bright Elevation)}

\vspace{0.5em}
\noindent\textbf{Sūtra ({\devanagarifont सूत्रम्}) — Verse 5:}

\begin{paracol}{2}
\backgroundcolor{c[0](.95\columnsep,.95\textheight)}{sanskritbg}
\backgroundcolor{c[1](.95\columnsep,.95\textheight)}{englishbg}

\begin{sanskritverse}
रविचन्द्रपातलग्नैः स्वक्रान्त्युदयात् स्वलग्नगतशेषाः।
घटिकाः खचरार्धास्तात् स्वेष्टी रविशीतथू कृत्वा॥ ५ ॥
\end{sanskritverse}

\switchcolumn

\textbf{Translation — Verse 5:} [Recapitulates the parameters for computing the bright elevation.]

\textit{Technical summary:} The verse reiterates that by the Sun, Moon, node, \textit{lagna}, and from the rising of the declination, the \textit{khacarārdha} ghaṭikās from the Sun's setting are to be used to compute the Moon's bright elevation.

\end{paracol}

\vspace{0.5em}

\begin{paracol}{2}
\backgroundcolor{c[0](.95\columnsep,.95\textheight)}{sanskritbg}
\backgroundcolor{c[1](.95\columnsep,.95\textheight)}{englishbg}

\noindent\textbf{Upapatti (\dev{उपपत्तिः}):}
\begin{sanskritverse}
यस्मिन् दिने सूर्यास्तादनन्तरं कालांशघटीभोद्धिकाभिश्चन्द्रास्तः। रव्युदयात् प्राक्
कालांशघटीभोद्धिकाभिर्चन्द्रोदयस्तस्मिन्नेव चन्द्रदर्शनं सति द्वे श्वज्योन्नतिः साध्या।
सितशुक्लोन्नतिरच बिम्बार्धं चन्द्रसितम् एवं भास्करेण मासान्तपादे प्रथमे च
शुक्लोन्नतिः साधिता। द्वितीयतृतीयपादयोर्बिम्बार्धं कृष्णं भवति। ततस्तत्र
कृष्णशुक्लोन्नतिर्भवति। इत्याचार्येण शुक्लोन्नतिद्वयं साध्यते। तदयं पादचर्चा न
कृता। भास्करेण न हि स्पष्टा कृष्णशुक्लोन्नतिः प्रेक्ष्यते इति मनसि संप्रधार्य
सितशुक्लोन्नतिरेव साधिता। उक्तं च ``द्वितीयतृतीययोरपि चरणयोर्ब्रह्मगुप्तादयः कृष्णशुक्लोन्नतिरानीता सा मम न संमता। न हि नरे कृष्णः शुक्लोन्नतिः स्पष्टोपलक्ष्यते''॥
\end{sanskritverse}

\switchcolumn

\textbf{English Translation of Proof:} ``On whatever day, after the Sun's setting, the Moon sets after an additional number of time-ghaṭīs. Before the Sun's rising, the Moon rises [after a certain number of ghaṭīs]. On that very day, when the Moon is visible, the two bright elevations (\textit{śva-jy-unnati}) are to be computed. The bright elevation of the bright [half] (\textit{sita-śuklonnati}) is half the disk (\textit{bimbārdha}) of the Moon's bright [portion]. Thus Bhāskara computed the \textit{śuklonnati} in the first \textit{pāda} of the month-end. In the second and third \textit{pāda}s, half the disk is dark (\textit{kṛṣṇa}). Therefore there the dark \textit{śuklonnati} occurs. Thus by the teacher (\textit{ācārya}), two \textit{śuklonnati}s are computed. But here, the \textit{pāda} discussion is not made. Bhāskara, having considered in his mind that the dark \textit{śuklonnati} is not clearly observable, computed only the bright \textit{śuklonnati}. And it is stated: `Brahmagupta and others have produced the dark \textit{śuklonnati} also for the second and third \textit{caraṇa}s; that is not approved by me. For neither is the dark \textit{śuklonnati} clearly perceived by men.'

\end{paracol}

%% ========================================================================
\subsection{Verses 11--12: {\devanagarifont स्पष्टशुक्ल} (True Bright Portion) Computation}

\vspace{0.5em}
\noindent\textbf{Sūtra ({\devanagarifont सूत्रम्}) — Verses 11--12:}

\begin{paracol}{2}
\backgroundcolor{c[0](.95\columnsep,.95\textheight)}{sanskritbg}
\backgroundcolor{c[1](.95\columnsep,.95\textheight)}{englishbg}

\begin{sanskritverse}
तच्छाया गुणिते वा मध्यगती तिथिगुणस्वकरं हते।
फलयोदिक् साम्येऽन्तरमवनतिरंक्यं दिगन्यत्वे॥ ११ ॥
\end{sanskritverse}

\switchcolumn

\textbf{Translation — Verse 11:} ``Or when the middle motion (\textit{madhyagati}) is multiplied by its own shadow (\textit{tacchāyā}) and by the semi-diameter (\textit{sva-kara}) multiplied by the \textit{tithi} --- the results are additive or subtractive. In equality (\textit{sāmya}), the difference is [taken]; in direction-difference (\textit{dig-anyatva}), [it is] the \textit{avanati} [depression] expressed in digits (\textit{aṅkya}).''

\textit{Note:} This verse gives the method for computing the true bright portion (\textit{spaṣṭa-śukla}) of the Moon. The \textit{madhyagati} is the mean daily motion, and the \textit{tacchāyā} is the shadow corresponding to the planet's altitude.

\end{paracol}

\vspace{0.5em}

\begin{paracol}{2}
\backgroundcolor{c[0](.95\columnsep,.95\textheight)}{sanskritbg}
\backgroundcolor{c[1](.95\columnsep,.95\textheight)}{englishbg}

\noindent\textbf{Sūtrārtha (\dev{सूत्रार्थः}):}
\begin{sanskritverse}
हृज्ये रविचन्द्रयोः क्षेपी स्वस्वमध्यगतिगुरो तिथिगुणितव्यासदलेन पञ्चदशगुरिणा
त्रिज्यया भक्त हृते पृथक् पृथक् रवि-चन्द्र-योर्-नती भवतः। वा तयोर्मध्यगती
तयोः क्षेपसमं हृज्ययोः छाये ताम्यां गुरिते तिथिगुणस्वस्वच्छायाकर्महते तयोरन्तरी भवतः।
\end{sanskritverse}

\switchcolumn

\textbf{Literal Meaning:} ``In the \textit{hṛjya} [radius = 3438'] of the Sun and Moon, the \textit{kṣepī} [sine of latitude], having been multiplied by the respective middle motion and by the semi-diameter (\textit{vyāsadala}) multiplied by the \textit{tithi} and by fifteen [i.e., by the \textit{tithi}-multiplied semi-diameter multiplied by 15], divided by the \textit{trijyā} [radius = 3438'], separately, separately, the \textit{natī} [depressions] of the Sun and Moon are obtained. Or, their middle motion, their \textit{kṣepa}-equal \textit{hṛjya}, the shadow being multiplied by that, multiplied by the \textit{tithi}-multiplied own shadow-\textit{karman}, their difference is [taken].''

\end{paracol}

\vspace{0.5em}

\begin{paracol}{2}
\backgroundcolor{c[0](.95\columnsep,.95\textheight)}{sanskritbg}
\backgroundcolor{c[1](.95\columnsep,.95\textheight)}{englishbg}

\noindent\textbf{Upapatti (\dev{उपपत्तिः}):}
\begin{sanskritverse}
प्रथमानयनस्य ``सहक्षेपाघ्ना इन्दोर्निजमध्यभुक्तितिथ्यंशनिघ्ना'' इत्यादि
भास्करविधिना स्फुटा।
\end{sanskritverse}

\switchcolumn

\textbf{English Translation of Proof:} ``The first computation [is made] accurate by the method of Bhāskara: `Having multiplied by the \textit{sahakṣepa} [combined latitude] of the Moon, by the own middle motion-\textit{tithi}-aṃśa-multiplied [quantity]' etc.''

\end{paracol}

\newpage


%% ========================================================================
%% CHAPTER 6: CHANDRAŚUKLONNATYADHIKĀRAḤ (Moon's Bright Elevation, Chapter VII)
%% ========================================================================
\section{\texorpdfstring{{\devanagarifont चन्द्रशुक्लोन्नत्यधिकारः} — Chandraśuklonnatyadhikāraḥ (Moon's Bright Elevation)}{Chandraśuklonnatyadhikāraḥ (Moon's Bright Elevation)}}
\label{sec:chandra-shuklonnati}

\begin{center}
\textcolor{rulecolor}{\rule{0.8\textwidth}{0.5pt}}
\end{center}

This is the seventh \dev{अध्याय} (chapter) of the \dev{ब्राह्मस्फुटसिद्धान्त}, dealing with the Moon's bright elevation (\dev{चन्द्रशुक्लोन्नति}). It comprises eighteen verses (\dev{श्लोक}) and covers the computation of the arc of illumination on the Moon's disk.

%% ========================================================================
\subsection{Verses 1--2: Introduction}

\vspace{0.5em}
\noindent\textbf{Sūtra (\dev{सूत्रम्}) — Verses 1--2:}

\begin{paracol}{2}
\backgroundcolor{c[0](.95\columnsep,.95\textheight)}{sanskritbg}
\backgroundcolor{c[1](.95\columnsep,.95\textheight)}{englishbg}

\begin{sanskritverse}
अथ चन्द्रशुक्लोन्नत्यधिकारः प्रारम्भः।
\end{sanskritverse}

\begin{sanskritverse}
प्राग्रहुण्डाकृतिकोर्ध्वकालिकोद्बलनखण्डसमस्थितिपाताः।
छन्दोव्यासस्तलग्ने स्वभग्रप्रायाण्न सङ्क्षेपः॥ १ ॥
\end{sanskritverse}

\begin{sanskritverse}
मूलोनं मूलं गतं यद्वदुद्धृत्यागतं ततोऽन्यत्।
स्वभुक्त्या हृतमभीष्टच्छायागतिः सैव मूलोद्धृता॥ २ ॥
\end{sanskritverse}

\switchcolumn

\textbf{Introduction:} ``Now the beginning of the chapter on the Moon's bright elevation (\textit{candra-śukla-unnati}).''

\textbf{Translation — Verse 1:} ``The parts of the upright-time-difference (\textit{ūrdhvakālika-udbala-naṣṭa-sama-sthiti-pāta}) [having] the form of a \textit{prāgrahuṇḍa} [spherical triangle], the semi-diameter (\textit{vyāsa}) of the shadow at the horizon (\textit{tala-lagna}), the proximity to its own \textit{agra} --- there is no abbreviation (\textit{saṅkṣepaḥ}).''

\textbf{Translation — Verse 2:} ``As the root (\textit{mūla}) diminished by the root, or the elapsed [time], having been extracted, another [quantity] comes therefrom. Divided by the own motion (\textit{svabhakti}), that is the desired shadow-motion (\textit{abhīṣṭa-ccchāyā-gati}). That same [quantity], extracted from the root (\textit{mūlodḍhṛtā}).''

\textit{Note:} Verse 2 establishes an iterative root-extraction algorithm for computing the rate of change of the shadow, which is essential for determining the Moon's illumination over time.

\end{paracol}

%% ========================================================================
\subsection{Verse 18: Chapter Conclusion}

\vspace{0.5em}
\noindent\textbf{Sūtra (\dev{सूत्रम्}) — Verse 18:}

\begin{paracol}{2}
\backgroundcolor{c[0](.95\columnsep,.95\textheight)}{sanskritbg}
\backgroundcolor{c[1](.95\columnsep,.95\textheight)}{englishbg}

\begin{sanskritverse}
इदानीमध्यायोपसंहारमाह—
\end{sanskritverse}

\begin{sanskritverse}
इति बाहुकोटिकर्णस्फुटसितपरिलेखसूत्रकर्दमेषु।
द्वाराष्टादशा चन्द्रशुक्लोन्नतिसप्तमोऽध्यायः॥ १८ ॥
\end{sanskritverse}

\switchcolumn

\textbf{Introduction:} ``Now [the author] states the conclusion of the chapter ---''

\textbf{Translation — Verse 18:} ``Thus, in [the topics of] \textit{bāhu} (base), \textit{koṭi} (perpendicular), \textit{karṇa} (hypotenuse), \textit{spaṣṭa} (true [longitude]), \textit{sita} (bright [part]), \textit{parilekha} (outline), \textit{sūtra} (line), and \textit{kardama} [mud] --- by eighteen doors (\textit{dvāra}), [this is] the seventh chapter [on] the Moon's bright elevation (\textit{candra-śuklonnati}).''

\textit{Note:} The chapter lists eight topics (\textit{bāhu}, \textit{koṭi}, \textit{karṇa}, \textit{spaṣṭa}, \textit{sita}, \textit{parilekha}, \textit{sūtra}, \textit{kardama}) covered in 18 verses. The \textit{kardama} refers to the ``mud'' or penumbral shadow region.

\end{paracol}

\vspace{0.5em}

\begin{paracol}{2}
\backgroundcolor{c[0](.95\columnsep,.95\textheight)}{sanskritbg}
\backgroundcolor{c[1](.95\columnsep,.95\textheight)}{englishbg}

\noindent\textbf{Sūtrārtha (\dev{सूत्रार्थः}):}
\begin{sanskritverse}
परिलेखसूत्रमेव कर्ण इति... परिलेखसूत्रकर्णः। कोषं स्पष्टार्थम्॥ १८ ॥
\end{sanskritverse}

\switchcolumn

\textbf{Literal Meaning:} ``The \textit{parilekha-sūtra} itself [is] the \textit{karṇa}... the \textit{parilekha-sūtra-karṇaḥ}. The meaning is clear.''

\end{paracol}

\vspace{0.5em}

\begin{paracol}{2}
\backgroundcolor{c[0](.95\columnsep,.95\textheight)}{sanskritbg}
\backgroundcolor{c[1](.95\columnsep,.95\textheight)}{englishbg}

\noindent\textbf{Hindi Bhāṣya (\dev{हिन्दी भाष्य}):}
\begin{sanskritverse}
इस तरह बाहु, कोटि, कर्ण, स्फुट, सिताङ्गुल, परिलेखसूत्र (कर्ण) में अष्टादश
भावार्थश्लोकों से चन्द्रशुक्लोन्नति नामक सप्तम अध्याय है॥ १८ ॥
\end{sanskritverse}

\switchcolumn

\textbf{English Translation:} ``In this way, in [the topics of] \textit{bāhu}, \textit{koṭi}, \textit{karṇa}, \textit{spaṣṭa}, \textit{sitāṅgula}, \textit{parilekha-sūtra} (\textit{karṇa}), by eighteen meaning-verses, is the seventh chapter named \textit{Candraśuklonnati}.''

\textit{Note:} The \textit{sitāṅgula} (bright digit) is a unit for measuring the illuminated portion of the Moon's disk. One \textit{aṅgula} (digit) = 1/12 of the Moon's diameter.

\end{paracol}

\vspace{1em}
\begin{center}
\textcolor{rulecolor}{\rule{0.8\textwidth}{0.5pt}}
\end{center}

\subsection*{Concluding Verse ({\devanagarifont अन्तिमश्लोकः})}

\begin{paracol}{2}
\backgroundcolor{c[0](.95\columnsep,.95\textheight)}{sanskritbg}
\backgroundcolor{c[1](.95\columnsep,.95\textheight)}{englishbg}

\begin{sanskritverse}
मधुसूदनसूनुनोदितो यस्तिलकः श्रीपथुनेह जिष्णुजोक्तः।
हृदि तं विनिधाय नूतनोऽयं रचितः सूर्यगतिः सुधाकरेण॥
\end{sanskritverse}

\begin{sanskritverse}
इति श्रीकृपालुदत्तसूनुसुधाकरद्विवेदिविरचिते ब्राह्मस्फुटसिद्धान्तनुतनतिलके
चन्द्रशुक्लोन्नत्यधिकारः सप्तमः॥ ७० ॥
\end{sanskritverse}

\switchcolumn

\textbf{Translation — Concluding Verse:} ``That \textit{tilaka} (ornament) which was spoken by the son (\textit{sūnu}) of Madhusūdana, which was spoken here by Śrīpathi for the conqueror (\textit{jiṣṇu}) --- having placed that in [my] heart, this new [work] \textit{Sūryagati} has been composed by Sudhākara.''

\textbf{Colophon:} ``Thus, in the \textit{Brāhmasphuṭasiddhānta-nutana-tilaka} composed by Śrī Kṛpāludatta-sūnu Sudhākara Dvivedī, the seventh chapter --- the \textit{Candraśuklonnatyadhikāraḥ} --- [is concluded].''

\end{paracol}

\begin{center}
\textcolor{rulecolor}{\rule{0.8\textwidth}{0.5pt}}
\end{center}

\noindent\textbf{\dev{॥ इति श्रीब्रह्मगुप्तविरचिते ब्राह्मस्फुटसिद्धान्ते चन्द्रशुक्लोन्नत्यधिकारः सप्तमः समाप्तः ॥}}

\noindent\textbf{[Thus ends the Seventh Chapter --- the Moon's Bright Elevation --- in the \textit{Brāhmasphuṭasiddhānta} composed by Śrī Brahmagupta.]}

\vspace{2cm}

%% ========================================================================
%% APPENDIX: NOTES ON THE TEXT AND COMMENTARY
%% ========================================================================
\newpage
\section*{Appendix: Notes on the Text and Commentary}
\addcontentsline{toc}{section}{Appendix: Notes on the Text and Commentary}

\subsection*{About the {\devanagarifont ब्राह्मस्फुटसिद्धान्त}}

The {\devanagarifont ब्राह्मस्फुटसिद्धान्त} (\textit{Brāhmasphuṭasiddhānta}, ``Correctly Established Doctrine of Brahma'') is the magnum opus of the great Indian mathematician and astronomer \textbf{Brahmagupta} (598--668 CE). Composed in 628 CE at Bhinmal (in present-day Rajasthan, India), it consists of 24 chapters ({\devanagarifont अध्याय}) containing a total of 1008 verses ({\devanagarifont श्लोक}) in {\devanagarifont आर्या} and {\devanagarifont अनुष्टुभ्} metres.

\subsection*{Contents of This Chunk (Pages 281--420)}

This typeset volume covers the following chapters from Part II of the text:

\begin{enumerate}
\item \textbf{\dev{चन्द्रग्रहणविकारः}} --- Lunar Eclipse Variations \\
  Deals with the computation of eclipse magnitudes, durations, and the \textit{iṣṭagrāsa} (desired obscuration) method.

\item \textbf{\dev{सूर्यग्रहणाधिकारः}} --- Solar Eclipse Chapter \\
  Covers the computation of solar eclipses, including the five \textit{jyā} (sines) required: \textit{udayajyā}, \textit{madhyajyā}, \textit{ravidayākarṇa}, \textit{viśimaśaṅku}, and \textit{hṛkṣepa}.

\item \textbf{\dev{चन्द्रच्छायाधिकारः}} --- Moon's Shadow Chapter \\
  Treats the computation of the Earth's shadow as seen on the Moon's disk.

\item \textbf{\dev{ग्रहोदयास्ताधिकारः}} --- Planetary Rising and Setting \\
  Explains the methods for computing the heliacal rising (\textit{udaya}) and setting (\textit{asta}) of planets.

\item \textbf{\dev{शुक्लोन्नत्यधिकारः}} --- The Brightness of the Moon \\
  Computes the illuminated portion of the Moon's disk.

\item \textbf{\dev{चन्द्रशुक्लोन्नत्यधिकारः}} --- Moon's Bright Elevation \\
  The seventh chapter, with 18 verses, covering the arc of illumination.
\end{enumerate}

\subsection*{The Commentary ({\devanagarifont नूतनतिलक})}

The Hindi commentary ({\devanagarifont टीका}) in this edition was composed by \textbf{Paṇḍita Kṛpāludatta Sūnu Sudhākara Dvivedī}. The commentary follows the traditional structure of Indian mathematical commentaries:

\begin{itemize}
\item {\devanagarifont सूत्र} (\textit{Sūtra}) --- The original Sanskrit verse
\item {\devanagarifont सूत्रार्थ} (\textit{Sūtrārtha}) --- Literal meaning of the verse (marked {\devanagarifont सु. भा.})
\item \dev{विवृति} (\textit{Vivṛti}) --- Detailed explanation (marked \dev{वि. भा.})
\item \dev{हिन्दी भाष्य} (\textit{Hindi Bhāṣya}) --- Hindi explanation (marked \dev{हि. भा.})
\item \dev{उपपत्ति} (\textit{Upapatti}) --- Mathematical proof or demonstration
\end{itemize}

\subsection*{Mathematical Notation}

Brahmagupta's work represents one of the high points of Indian mathematics. Key contributions in these chapters include:

\begin{itemize}
\item Rules for computing eclipse magnitudes using the \dev{भुज} (\textit{bhuja}, sine) and \dev{कोटि} (\textit{koṭi}, cosine)
\item The \dev{कर्ण} (\textit{karṇa}) formula: $\text{karṇa} = \sqrt{\text{bhuja}^2 + \text{koṭi}^2}$
\item Iterative methods (\dev{असकृत्}, \textit{asakṛt}) for refining planetary positions
\item Proportional reasoning (\dev{त्रैराशिक}, \textit{trairāśika}) for time conversions
\item Geometric constructions using \dev{क्षेत्र} (\textit{kṣetra}, diagrams)
\end{itemize}

\subsection*{IAST Transliteration Conventions}

This bilingual edition uses the International Alphabet of Sanskrit Transliteration (IAST) scheme:

\begin{center}
\begin{tabular}{llllll}
\toprule
\textbf{Devanagari} & \textbf{IAST} & \textbf{Devanagari} & \textbf{IAST} & \textbf{Devanagari} & \textbf{IAST} \\
\midrule
अ & a & क & ka & त & ta \\
आ & ā & ख & kha & थ & tha \\
इ & i & ग & ga & द & da \\
ई & ī & घ & gha & ध & dha \\
उ & u & ङ & ṅ & न & na \\
ऊ & ū & च & ca & प & pa \\
ए & e & छ & cha & फ & pha \\
ऐ & ai & ज & ja & ब & ba \\
ओ & o & झ & jha & भ & bha \\
औ & au & ञ & ñ & म & ma \\
ऋ & ṛ & ट & ṭa & य & ya \\
ॄ & ṝ & ठ & ṭha & र & ra \\
\bottomrule
\end{tabular}
\end{center}

\vspace{1cm}
\begin{center


% ============================================================
% Source: translations/en/indian/brahmagupta-brahmasphutasiddhanta-pt2-chunk4_bilingual.tex
% ============================================================

%% ========================================================================
%% Brahmagupta's Brahmasphutasiddhanta Part 2 — CHUNK 4
%% Bilingual English-Sanskrit Edition
%% Pages 1259–1398 (PDF pages 421–560)
%% ========================================================================

\documentclass[11pt,a4paper]{article}

%% -------- Core Packages --------
\usepackage{fontspec}
\usepackage{polyglossia}
\setmainlanguage{english}

\usepackage{amsmath,amssymb}
\usepackage{geometry}
\usepackage{xcolor}
\usepackage{fancyhdr}
\usepackage{array,booktabs,longtable}
\usepackage{hyperref}

%% -------- Page Geometry --------
\geometry{
  paperwidth=210mm,paperheight=297mm,
  left=15mm,right=15mm,top=22mm,bottom=22mm,
  headheight=14pt,footskip=25pt
}

%% -------- Font Configuration --------
\setmainfont{Linux Libertine O}
\setmonofont{DejaVu Sans Mono}

%% Devanagari Font — explicit system path (fontspec uses path directly)
\newfontfamily\devanagarifont[Script=Devanagari,Scale=1.1,Path=/usr/share/fonts/truetype/noto/,UprightFont=NotoSerifDevanagari-Regular.ttf]{NotoSerifDevanagari}

%% -------- Colors --------
\definecolor{titlecolor}{RGB}{45,55,72}
\definecolor{sectioncolor}{RGB}{55,65,81}
\definecolor{rulecolor}{RGB}{180,160,140}
\definecolor{sanskritbg}{RGB}{252,250,247}

%% -------- Section Formatting --------
\makeatletter
\renewcommand{\section}{\@startsection {section}{1}{\z@}%%
  {-3ex plus -1ex minus -.2ex}{2ex plus .2ex}%%
  {\normalfont\Large\bfseries\color{sectioncolor}}}
\renewcommand{\subsection}{\@startsection{subsection}{2}{\z@}%%
  {-2.5ex plus -1ex minus -.2ex}{1.2ex plus .2ex}%%
  {\normalfont\large\bfseries\color{sectioncolor}}}
\makeatother

%% -------- Header/Footer --------
\pagestyle{fancy}
\fancyhf{}
\fancyhead[L]{\small\textit{Brāhmasphuṭasiddhānta II --- Chunk 4 (Bilingual)}}
\fancyhead[R]{\small\thepage}
\renewcommand{\headrulewidth}{0.4pt}

%% -------- Custom Commands --------
\newcommand{\dev}[1]{{\devanagarifont\mdseries\upshape #1}}
\newcommand{\pagelabel}[1]{\textcolor{sectioncolor}{\textbf{[Original p.~#1]}}}

%% Parallel columns: Sanskrit (left) | English (right)
\newcommand{\bilingualpair}[2]{%%
  \noindent%%
  \begin{minipage}[t]{0.47\textwidth}
    \begin{flushleft}
      \small\dev{#1}
    \end{flushleft}
  \end{minipage}%%
  \hfill\vrule\hfill%%
  \begin{minipage}[t]{0.47\textwidth}
    \begin{flushleft}
      \small #2
    \end{flushleft}
  \end{minipage}%%
  \par\vspace{1.2ex}
}

%% Sanskrit-only block (no translation)
\newcommand{\sanskritblock}[1]{%%
  \noindent%%
  \begin{minipage}[t]{0.47\textwidth}
    \begin{flushleft}
      \small\dev{#1}
    \end{flushleft}
  \end{minipage}%%
  \hfill\vrule\hfill%%
  \begin{minipage}[t]{0.47\textwidth}
    \begin{flushleft}
      \small\textit{[Sanskrit text; translation not provided in source]}
    \end{flushleft}
  \end{minipage}%%
  \par\vspace{1.2ex}
}

%% Math block spanning both columns
\newcommand{\mathblock}[1]{%%
  \par\vspace{0.8ex}
  \noindent\makebox[\textwidth]{#1}%%
  \par\vspace{0.8ex}
}

%% -------- Hyperref Setup --------
\hypersetup{
  colorlinks=true,
  linkcolor=sectioncolor,
  citecolor=sectioncolor,
  urlcolor=sectioncolor,
  pdfencoding=auto,
  pdftitle={Brahmagupta's Brahmasphutasiddhanta Part 2 Chunk 4 — Bilingual Edition},
  pdfauthor={Brahmagupta}
}

%% ========================================================================
\begin{document}

%% -------- Title Page --------
\thispagestyle{empty}
\begin{center}
\vspace*{2.5cm}

{\fontsize{13}{17}\selectfont\textit{Bilingual English--Sanskrit Edition}}

\vspace{1.2cm}

{\fontsize{20}{26}\selectfont\bfseries\color{titlecolor}
Brāhmasphuṭasiddhānta\\[0.4em]
{\large Part II --- Chunk 4}}

\vspace{0.8cm}

{\fontsize{12}{14}\selectfont\dev{ब्राह्मस्फुटसिद्धान्त} --- \dev{द्वितीयभागः} --- \dev{चतुर्थं खण्डकम्}}

\vspace{0.8cm}

\noindent\rule{0.8\textwidth}{0.5pt}

\vspace{0.8cm}

{\LARGE\bfseries Brahmagupta}

\vspace{0.4cm}

{\large (c.~598--668 CE)}

\vspace{1.5cm}

{\large Pages 1259--1398}

\vspace{0.6cm}

{\small\dev{सूत्राणि} \textbar\ Sanskrit Verses \quad --- \quad \dev{व्याख्या} \textbar\ English Translation}

\vspace{1cm}

\noindent\rule{0.8\textwidth}{0.5pt}

\vspace{1.2cm}

{\small\textcolor{sectioncolor}{Typeset with XeLaTeX using Noto Serif Devanagari\\
Transcribed from the original printed edition}}

\end{center}

\newpage

%% ========================================================================
\section*{Introduction / \dev{प्रवेशः}}
\addcontentsline{toc}{section}{Introduction}
%% ========================================================================

\bilingualpair{%
एतद् ग्रन्थखण्डकं ब्राह्मस्फुटसिद्धान्तस्य द्वितीयभागस्य चतुर्थं खण्डकम् । एतस्मिन् पृष्ठसङ्ख्याः १२५९--१३९८ समाविशन्ति । ब्रह्मगुप्तेन रचितं गणितज्योतिषकोशम् एतत्, यस्मिन् चन्द्रग्रहणादीनां गणनाविधयः सविस्तरं निरूपिताः ॥
}{%
This chunk (pages 1259--1398) of Part~II of the \textit{Brāhmasphuṭasiddhānta} comprises the fourth textual division of Brahmagupta's encyclopaedic treatise on mathematics and astronomy, containing detailed computational rules for lunar and planetary calculations.
}

\vspace{1ex}
\noindent\rule{\textwidth}{0.3pt}
\vspace{1ex}

\bilingualpair{%
\textbf{अस्य खण्डकस्य विषयाः:} (१) चन्द्रस्य स्फुटराश्यनयनम्, (२) शून्यविधिः, (३) वर्गप्रकृतिः, (४) ग्रहाणां मन्दशीघ्रफलानि, (५) विषयानुक्रमणिका च ॥
}{%
\textbf{Contents of this chunk:} (1)~Computation of the Moon's true longitude; (2)~Rules for operations with zero; (3)~The nature of squares (Pell's equation); (4)~Planetary equations of centre and anomaly; (5)~Index of subjects.
}

\newpage

%% ========================================================================
\section*{Chapter on Lunar Calculations / \dev{चन्द्रराशिः} --- Pages 1259--1287}
\addcontentsline{toc}{section}{Lunar Calculations (pp.~1259--1287)}
%% ========================================================================

\pagelabel{1259}

\bilingualpair{%
सू.~भा.---चन्द्रस्य स्फुटराशिः । मध्यमराशिम् मन्दफलशीघ्रफलाभ्यां संस्कृत्य स्फुटराशिः लभ्यते । अहर्गणात् मध्यमराशिः गणनीयः ॥
}{%
\textsc{Sū.~Bhā.}---The true longitude of the Moon: having applied the corrections for equation of centre (manda-phala) and evection (śīghra-phala) to the mean longitude, the true longitude is obtained. The mean longitude is calculated from the ahargaṇa (elapsed days since the epoch).
}

\vspace{0.5ex}
\pagelabel{1260--1264}

\bilingualpair{%
सू.~भा.---चन्द्रमन्दफलानि । चन्द्रस्य मन्दफलानि प्रत्येकं पञ्चदशांशानतराले सूचीकृतानि । तान्य् अनयानि---३९, ३६, ३१, २४, १५, ५ । एभिः स्फुटीक्रियते ॥
}{%
\textsc{Sū.~Bhā.}---The Moon's equations (manda-phala) are tabulated for each $15°$ of anomaly. The tabular differences are: 39, 36, 31, 24, 15, 5. These are applied successively to obtain the true longitude.
}

\bilingualpair{%
वि.~भा.---मध्यमस्यानतराले । यदा कोणः सूचीकृतमध्ये वर्तते, तदा अवशिष्टचापं गृहीत्वा गतागतफलयोः अर्धान्तरेण गुणयित्वा नवशतभक्तं समर्धयोर्योगेन योजयेत् ॥
}{%
\textsc{Vi.~Bhā.}---When the anomaly lies between two tabulated values, multiply the residual arc by half the difference of the passed and upcoming tabular differences, divide by $900'$, and add to half the sum of the two differences.
}

\vspace{0.5ex}
\pagelabel{1265--1269}

\bilingualpair{%
सू.~भा.---शीघ्रफलानि । चन्द्रस्य शीघ्रफलानि शीघ्रकोणात् गणनीयानि । शीघ्रकोणः चन्द्रसूर्ययोर् अन्तरम् । शीघ्रफलस्य parametमं $८०'$ (अशीतिः कला) ॥
}{%
\textsc{Sū.~Bhā.}---The śīghra corrections for the Moon are computed from the śīghra anomaly, which is the difference between the Moon's longitude and the Sun's longitude. The maximum śīghra equation is $80'$ (80 arc-minutes).
}

\vspace{0.5ex}
\pagelabel{1270--1274}

\bilingualpair{%
सू.~भा.---लग्नानयनम् । लग्नानयनार्थं सूर्योदयात् गतकालं रवेः दैनिकगत्या गुणयित्वा रविकोक्षेप्य आयनांशसंस्कारं कुर्यात् ॥
}{%
\textsc{Sū.~Bhā.}---To find the lagna (ascendant): multiply the time elapsed since sunrise by the Sun's daily motion, add to the Sun's longitude, and apply the ayanāṃśa correction.
}

\vspace{0.5ex}
\pagelabel{1275--1279}

\bilingualpair{%
सू.~भा.---युतिकालानयनम् । ग्रहयुतिकालः द्वयोर् ग्रहयोः राश्यन्तरेण द्वयोः दैनिकगत्योर् अन्तरेण भजयित्वा लभ्यते । फलं युतिकालं दिनानि ददाति ॥
}{%
\textsc{Sū.~Bhā.}---The time of conjunction is found by dividing the difference in longitudes by the difference in daily motions. The result gives the number of days until conjunction.
}

\vspace{0.5ex}
\pagelabel{1280--1284}

\bilingualpair{%
सू.~भा.---कालविभागः । प्राणः चतुःषष्टिकला भवति । षट् प्राणाः एका विनाडी, सष्टिः विनाड्यः एका नाडी (घटी), सष्टिर् नाड्यः एकं दिनम् ॥
}{%
\textsc{Sū.~Bhā.}---A prāṇa equals 4 seconds of time. Six prāṇas make a vināḍī, 60 vināḍīs make a nāḍī (ghaṭī), and 60 nāḍīs make a day.
}

\vspace{0.5ex}
\pagelabel{1285--1287}

\bilingualpair{%
चन्द्रगणनाप्रकरणस्य समाप्तिः । अग्रे छायाव्यवाहारः प्रवर्तते ॥
}{%
Here end the verses on lunar computation; next follows the section on shadow calculations (chāyā-vyavahāra).
}

\newpage

%% ========================================================================
\section*{Mathematical Derivations / \dev{उपपत्तयः} --- Pages 1288--1337}
\addcontentsline{toc}{section}{Mathematical Derivations (pp.~1288--1337)}
%% ========================================================================

\pagelabel{1288--1292}

\bilingualpair{%
उपपत्तिः---ज्या कोटिरूपेण सूचीकृतानि, मध्ये चापान्तरं $h = १५°$ । मध्यमे कोणे $	heta$ अन्तरे ज्या:
$$\sin(\theta) = j_1 + \frac{\theta}{h}(j_2 - j_1) - \frac{\theta}{h}\Bigl(1 - \frac{\theta}{h}\Bigr)\frac{j_2 - j_3}{2}$$
एषा ब्रह्मगुप्तस्य द्वितीयक्रमान्तर्गमनरीतिः ॥
}{%
\textsc{Proof.}---Let $j_1$ and $j_2$ be two successive tabular sines, with arc difference $h = 15°$. For an intermediate arc at distance $\theta$ from $j_1$, the sine is:
\mathblock{$\displaystyle\sin(\theta) = j_1 + \frac{\theta}{h}(j_2 - j_1) - \frac{\theta}{h}\Bigl(1 - \frac{\theta}{h}\Bigr)\frac{j_2 - j_3}{2}$}
This is Brahmagupta's second-order interpolation formula.
}

\vspace{0.5ex}
\pagelabel{1293--1297}

\bilingualpair{%
उपपत्तिः---कुट्टकरीत्या $ax + c = by$ गणनीयम्, यत्र $a = १३७$, $b = २५१$, $c = १५$ । २५१ अपेक्ष्य १३७---भागलब्धिः १, शेषं ११४ । १३७ अपेक्ष्य ११४---भागलब्धिः १, शेषं २३ । एवं पश्चात् गुणकारलब्धी प्राप्येते ॥
}{%
\textsc{Proof.}---Applying the kuṭṭaka (pulverizer) to solve $137x + 15 = 251y$: divide 251 by 137, quotient 1, remainder 114. Divide 137 by 114, quotient 1, remainder 23. Divide 114 by 23, quotient 4, remainder 22. Divide 23 by 22, quotient 1, remainder 1. Working backwards yields the multiplier and quotient.
}

\vspace{0.5ex}
\pagelabel{1298--1302}

\bilingualpair{%
सू.~भा.---छायानयनम् । शङ्कोः (निरङ्कुशस्य) द्वादशाङ्गुलस्य समीपे छाया निर्णेया:
$$\text{छाया} = \frac{12 \times \cos(\text{नतभूजा})}{\sin(\text{उन्नतांशः})}$$
यदा रविः प्रथमलम्बकस्थः, तदा विशेषनियमाः प्रयुज्यन्ते ॥
}{%
\textsc{Sū.~Bhā.}---The shadow (chāyā) for a gnomon of 12 aṅgulas:
\mathblock{$\displaystyle\text{Shadow} = \frac{12 \times \cos(\text{zenith distance})}{\sin(\text{altitude})}$}
When the Sun is at the prime vertical, special rules apply using the nara (sine of latitude) and śaṅku (gnomon).
}

\vspace{0.5ex}
\pagelabel{1303--1307}

\bilingualpair{%
उपपत्तिः---लग्नं क्षितिज्यात् लभ्यते । $
\phi$ अक्षांशः, $\varepsilon$ क्रान्तिः । दत्तकाले $t$ (सूर्योदयात् अंशैः मिते) लग्नं $\lambda$:
$$\tan(\lambda) = \frac{\sin(\text{R.A.})}{\cos(\text{R.A.}) \cos(\varepsilon) - \tan(\phi) \sin(\varepsilon)}$$
}{%
\textsc{Proof.}---The lagna is found from the oblique ascension. Let $\phi$ be the latitude, $\varepsilon$ the obliquity of the ecliptic. For elapsed time $t$ (in degrees from sunrise), the lagna $\lambda$ satisfies:
\mathblock{$\displaystyle\tan(\lambda) = \frac{\sin(\text{R.A.})}{\cos(\text{R.A.}) \cos(\varepsilon) - \tan(\phi) \sin(\varepsilon)}$}
}

\vspace{0.5ex}
\pagelabel{1308--1312}

\bilingualpair{%
सू.~भा.---मन्दफलोदाहरणम् । कुजस्य मन्दकोणे $\mu$ परिधौ $r = ३९$ (यत्र मध्यमपरिधिः $R = ६०$):
$$\text{मन्दफलं} = \arcsin\!\left(\frac{39 \sin(\mu)}{\sqrt{60^2 + 39^2 + 2 \cdot 60 \cdot 39 \cos(\mu)}}\right)$$
}{%
\textsc{Sū.~Bhā.}---For Mars, with manda anomaly $\mu$ and epicycle radius $r = 39$ (where deferent radius $R = 60$):
\mathblock{$\displaystyle\text{manda-phala} = \arcsin\!\left(\frac{39 \sin(\mu)}{\sqrt{60^2 + 39^2 + 2 \cdot 60 \cdot 39 \cos(\mu)}}\right)$}
This gives the equation of centre for Mars.
}

\vspace{0.5ex}
\pagelabel{1313--1317}

\bilingualpair{%
उपपत्तिः---प्रथमं मन्दसंस्कारेण मन्दस्फुटं लभ्यते:
$\lambda_{\text{ms}} = \bar{\lambda} + \text{मन्दफलम्}$ ।
ततः शीघ्रकोणं मन्दस्फुटात् गणयित्वा शीघ्रसंस्कारं ददाति:
$\lambda_{\text{स्फुट}} = \lambda_{\text{ms}} + \text{शीघ्रफलम्}$ ॥
}{%
\textsc{Proof.}---First apply manda correction to obtain manda-sphuṭa: $\lambda_{\text{ms}} = \bar{\lambda} + \text{manda-phala}$. Then compute śīghra anomaly from $\lambda_{\text{ms}}$ and apply śīghra correction: $\lambda_{\text{true}} = \lambda_{\text{ms}} + \text{śīghra-phala}$.
}

\vspace{0.5ex}
\pagelabel{1318--1322}

\bilingualpair{%
सू.~भा.---ग्रहणानयनम् । चन्द्रस्य दूरत्वे पृथिव्याः छायापरिधिः:
$$R_{\text{छाया}} = \frac{R_{\text{पृथिवी}} \cdot (D_{\text{चन्द्र}} - D_{\text{सूर्य}})}{D_{\text{सूर्य}}}$$
युत्यां चन्द्रस्य अक्षः ग्रहणं निर्णयति ॥
}{%
\textsc{Sū.~Bhā.}---The semi-diameter of the Earth's shadow at the Moon's distance:
\mathblock{$\displaystyle R_{\text{shadow}} = \frac{R_{\text{Earth}} \cdot (D_{\text{Moon}} - D_{\text{Sun}})}{D_{\text{Sun}}}$}
where $D$ denotes distance. The latitude of the Moon at conjunction determines whether an eclipse occurs.
}

\vspace{0.5ex}
\pagelabel{1323--1327}

\bilingualpair{%
सू.~भा.---उदयकालः । ग्रहः सूर्यात् दूरे सति दृश्यते । दृग्भूमिका ग्रहप्रभावतः भिद्यते---शुक्रस्य $१०°$, बृहस्पतेः $११°$, बुधस्य $१३°$, शनैश्चरस्य $१५°$, कुजस्य $१७°$ ॥
}{%
\textsc{Sū.~Bhā.}---A planet becomes visible when it is far enough from the Sun. The arc of visibility depends on the planet's brightness: for Venus $10°$, for Jupiter $11°$, for Mercury $13°$, for Saturn $15°$, and for Mars $17°$.
}

\vspace{0.5ex}
\pagelabel{1328--1337}

\bilingualpair{%
उपपत्तिः---ऋजुवृत्तसंस्कारः, ग्रहाणां कलाभिः सह भौमिकसङ्गणनानि च । एषा खण्डिका ब्रह्मस्फुटसिद्धान्तस्य गणितीयरीतिसारांशं समाविशति ॥
}{%
\textsc{Proof.}---Ṛjuvṛtta (straight epicycle) corrections and planetary phase computations. These final derivations summarise the algebraic methods used throughout the \textit{Brāhmasphuṭasiddhānta} for astronomical computations.
}

\newpage

%% ========================================================================
\section*{Index of Subjects / \dev{विषयानुक्रमणिका} --- Page 1338}
\addcontentsline{toc}{section}{Index of Subjects (p.~1338)}
%% ========================================================================

\pagelabel{1338}

\bilingualpair{%
\textbf{अनुक्रमणिका} (पृष्ठ १३३८)---अस्य ग्रन्थस्य विषयान् सूचयति ॥
}{%
\textbf{Index of Subjects} (page 1338)---A detailed listing of topics covered in this volume with corresponding page references.
}

\vspace{1ex}

\noindent\begin{minipage}{\textwidth}
\centering
\small
\begin{longtable}{p{0.44\textwidth}rp{0.02\textwidth}p{0.44\textwidth}r}
\toprule
\multicolumn{2}{c}{\dev{Sanskrit Topic}} && \multicolumn{2}{c}{Transliterated Topic \& Page} \\
\midrule
\endhead
\dev{प्रायभटोक्तबारप्रवृत्तिखण्डनम्} & ६५९ && prāyabhaṭokta-bāra-pravṛtti-khaṇḍanam & 659 \\
\dev{प्रायभटोक्तबारादिक्खण्डनम्} & ६६३ && prāyabhaṭokta-bārādik-khaṇḍanam & 663 \\
\dev{प्रायभटोक्तग्रहयोः खण्डनम्} & ६७२ && prāyabhaṭokta-grahayoḥ khaṇḍanam & 672 \\
\dev{मृत्यासस्य प्राधान्यदर्शनम्} & ६७५ && mṛtyāsasya prādḥānya-darśanam & 675 \\
\dev{प्रायभटोक्तभ्रमरग्रन्थखण्डनम्} & ६७७ && prāyabhaṭokta-bhramara-grantha-khaṇḍanam & 677 \\
\dev{आर्यभटोक्तमन्दपरिधिखण्डनम्} & ६७९ && āryabhaṭokta-manda-paridhi-khaṇḍanam & 679 \\
\dev{प्रायभटोक्तगोलार्धादिक्खण्डनम्} & ६८५ && prāyabhaṭokta-golārdhādik-khaṇḍanam & 685 \\
\dev{प्राग्रवेशेन रवेः सम्मण्डलप्रवेशावखण्डनम्} & ६८९ && prāgraveśena raveḥ sam-maṇḍala-praveśāvakhaṇḍanam & 689 \\
\dev{प्रायभटोक्त लम्बनावनत्यानयनखण्डनम्} & ६९१ && prāyabhaṭokta-lambanāvanatyānayana-khaṇḍanam & 691 \\
\dev{प्रायभटोक्तलम्बनलापण्डनम्} & ६९५ && prāyabhaṭokta-lambana-lāpaṇḍanam & 695 \\
\dev{स्फुट हक्रक्षेपन स्थानकथनम्} & ७०० && sphuṭa-hakakṣepana-sthāna-kathanam & 700 \\
\dev{श्रीपैठावधूचन्द्रकृतमैत्रेयप्रहरणखण्डनम्} & ७०२ && śrīpaiṭhāvadhūcandrakṛta-maitreya-praharaṇa-khaṇḍanam & 702 \\
\dev{स्वर्गसङ्क्रान्तप्रनिपादनम्} & ७०३ && svarga-saṅkrānta-pranipādanam & 703 \\
\dev{प्रायभटोक्तमक्षजहकक्षमेखण्डनम्} & ७०४ && prāyabhaṭokta-makṣaja-hakakṣame-khaṇḍanam & 704 \\
\dev{प्रायभटोक्तायनहकर्मखण्डनम्} & ७०६ && prāyabhaṭoktāyana-hakarma-khaṇḍanam & 706 \\
\dev{शृङ्गोत्क्रमणावायभटोक्तगुणलावखण्डनम्} & ७०९ && śṛṅgotkramaṇāv-āryabhaṭokta-guṇalāva-khaṇḍanam & 709 \\
\dev{प्रायभटसाधितचन्द्रदिनगणनगोयखण्डनम्} & ७१२ && prāyabhaṭa-sādhita-candra-dinagaṇana-goya-khaṇḍanam & 712 \\
\dev{दिनगणनविधिकार्थानोर्पथदोषकथनम्} & ७१३ && dina-gaṇana-vidhi-kārthānorpatha-doṣa-kathanam & 713 \\
\dev{प्रायभटदोषकथनम्} & ७१६ && prāyabhaṭa-doṣa-kathanam & 716 \\
\dev{श्रीपैठदोषानां दोषकथनम्} & ७१६ && śrīpaiṭha-doṣānāṃ doṣa-kathanam & 716 \\
\dev{द्वितीयदोषकथनम्} & ७१७ && dvitīya-doṣa-kathanam & 717 \\
\dev{महायुगलक्षगाकथनम्} & ७२३ && mahāyuga-lakṣagā-kathanam & 723 \\
\dev{श्रीपैठादिकथितयुगखण्डनम्} & ७२४ && śrīpaiṭhādikathita-yuga-khaṇḍanam & 724 \\
\dev{पादक्रमाद्य दुष्प्रकथनम्} & ७२५ && pāda-kramādya duṣprakathanam & 725 \\
\bottomrule
\end{longtable}
\end{minipage}

\newpage

%% ========================================================================
\section*{Algebra / \dev{बीजगणितम्} --- Pages 1339--1370}
\addcontentsline{toc}{section}{Algebra: Zero and Pell's Equation (pp.~1339--1370)}
%% ========================================================================

\pagelabel{1339--1343}

\bilingualpair{%
सू.~भा.---शून्यविधिः । धनयोर् योगो धनम्, ऋणयोर् योगो ऋणम् । धनऋणयोः योगः अन्तरं; समे सत्य् अंशे शून्यं भवति ॥
}{%
\textsc{Sū.~Bhā.}---The sum of two positive quantities is positive; the sum of two negative quantities is negative. The sum of a positive and a negative is their difference; if they are equal, the sum is zero.
}

\vspace{0.5ex}
\pagelabel{1344--1348}

\bilingualpair{%
सू.~भा.---शून्यगुणनम् । शून्येन गुणितं किमपि शून्यं भवति । धनं ऋणेन गुणितं ऋणं, ऋणं ऋणेन गुणितं धनं । शून्यं येनापि भक्तं शून्यम् ॥
}{%
\textsc{Sū.~Bhā.}---Zero multiplied by any quantity is zero. A positive quantity multiplied by a negative is negative; a negative by a negative is positive. Zero divided by any quantity is zero.
}

\bilingualpair{%
सू.~भा.---शून्यहारः । शून्येन विभक्तं राशिः शून्यहारः तत् ख-बीजम् अभवत् । अस्य संशोधनं भास्करेण द्वितीयेन कृतम् ॥
}{%
\textsc{Sū.~Bhā.}---A quantity divided by zero becomes that fraction of which the denominator is zero. This \textit{khacheda} (zero-divisor, or ``zero-seed'') was later corrected by Bhāskara~II.
}

\vspace{0.5ex}
\pagelabel{1349--1353}

\bilingualpair{%
सू.~भा.---वर्गप्रकृतिः । $Nx^2 + 1 = y^2$ समीकरणे: प्रथमं लघूत्तरं $(x_1, y_1)$ आक्षेपेण वा चक्रवालविधिना लभ्यते । अनन्तरं भावनया अनन्तानि उत्तराणि जनयितुं शक्यन्ते ॥
}{%
\textsc{Sū.~Bhā.}---To solve $Nx^2 + 1 = y^2$ (Pell's equation): first find the least solution $(x_1, y_1)$ by inspection or by the cakravāla (cyclic) method. From this, infinite solutions can be generated by composition (bhāvanā).
}

\vspace{0.5ex}
\pagelabel{1354--1358}

\bilingualpair{%
उपपत्तिः---$N = ६१$ इत्यस्मिन्: आरभ्य $a_0 = ८$ (यतः $8^2 = 64 \approx 61$) । कल्पकः $k_0 = ३$ । चक्रवालविधिना क्रमशः सन्निकर्षा:
$$a_1 = \frac{a_0 + m_0}{|k_0|}, \quad k_1 = \frac{m_0^2 - N}{k_0}$$
यत्र $m_0$ इति चीयते यत् $a_0 + m_0$ क्षेपेण विभाज्यं, $m_0^2$ च $N$ समीपे ॥
}{%
\textsc{Proof.}---For $N = 61$: start with $a_0 = 8$ (since $8^2 = 64 \approx 61$). The interpolator $k_0 = 3$. The cyclic process generates successive approximations:
\mathblock{$\displaystyle a_1 = \frac{a_0 + m_0}{|k_0|}, \quad k_1 = \frac{m_0^2 - N}{k_0}$}
where $m_0$ is chosen so that $a_0 + m_0$ is divisible by $k_0$ and $m_0^2$ is close to $N$.
}

\vspace{0.5ex}
\pagelabel{1359--1363}

\bilingualpair{%
सू.~भा.---समासः । समान्तरश्रेणीफलं $S_n = \frac{n}{2}(2a + (n-1)d)$ । वर्गाणां समासः $\sum_{k=1}^{n} k^2 = \frac{n(n+1)(2n+1)}{6}$ । घनानां समासः $\sum_{k=1}^{n} k^3 = \bigl(\frac{n(n+1)}{2}\bigr)^2$ ॥
}{%
\textsc{Sū.~Bhā.}---The sum of an arithmetic series: $S_n = \frac{n}{2}(2a + (n-1)d)$. The sum of squares: $\sum_{k=1}^{n} k^2 = \frac{n(n+1)(2n+1)}{6}$. The sum of cubes: $\sum_{k=1}^{n} k^3 = \bigl(\frac{n(n+1)}{2}\bigr)^2$.
}

\vspace{0.5ex}
\pagelabel{1364--1368}

\bilingualpair{%
सू.~भा.---मिश्रकव्यवहारः । चक्रवृद्धौ, यदि मूलधनं $P$ दरात् $r$ अन्तराले $n$ अन्तरैः वर्धते, तदा मिश्रधनं $A = P(1 + r)^n$ ॥
}{%
\textsc{Sū.~Bhā.}---In compound interest, if principal $P$ grows at rate $r$ per period for $n$ periods, the amount is $A = P(1 + r)^n$.
}

\vspace{0.5ex}
\pagelabel{1369--1370}

\bilingualpair{%
बीजगणितप्रकरणस्य समाप्तिः । सिद्धान्तस्य स्तुतिः, अग्रे ग्रहचाराधिकारस्य प्रवेशः ॥
}{%
Here conclude the verses of the algebra section, with praise of the siddhānta and transition to the planetary motion chapter.
}

\newpage

%% ========================================================================
\section*{Planetary Motion / \dev{भग्रहचाराधिकारः} --- Pages 1371--1398}
\addcontentsline{toc}{section}{Planetary Motion (pp.~1371--1398)}
%% ========================================================================

\pagelabel{1371--1375}

\bilingualpair{%
सू.~भा.---ग्रहसाधनम् । अहर्गणात् मध्यमं गणयित्वा मन्दशीघ्रसंस्काराभ्यां क्रमेण स्फुटं लभ्यते ॥
}{%
\textsc{Sū.~Bhā.}---Having computed the mean longitude (madhyama) from the ahargaṇa, apply the manda and śīghra corrections in order to obtain the true longitude (sphuṭa).
}

\vspace{0.5ex}
\pagelabel{1376--1380}

\bilingualpair{%
\textbf{मन्दपरिधयः}---कुजस्य $३९$, बुधस्य $२२$, गुरोः $३२$, शुक्रस्य $१४$, शनेः $६०$ । \textbf{शीघ्रपरिधयः}---कुजस्य $७३$, बुधस्य $१२१$, गुरोः $७१$, शुक्रस्य $१३२$, शनेः $३९$ ॥
}{%
\textbf{Manda epicycles:} Mars 39, Mercury 22, Jupiter 32, Venus 14, Saturn 60. \textbf{Śīghra epicycles:} Mars 73, Mercury 121, Jupiter 71, Venus 132, Saturn 39 (in units where $R = 60$).
}

\vspace{1ex}

\begin{table}[h]
\centering
\small
\begin{tabular}{lcccc}
\toprule
\textbf{Planet} & \textbf{Manda Epicycle} & \textbf{Śīghra Epicycle} & \textbf{Max Equation} \\
\midrule
Mars (कुजः) & 39 & 73 & $11° 30'$ \\
Mercury (बुधः) & 22 & 121 & $22°$ \\
Jupiter (गुरुः) & 32 & 71 & $5°$ \\
Venus (शुक्रः) & 14 & 132 & $11°$ \\
Saturn (शनैः) & 60 & 39 & $8°$ \\
\bottomrule
\end{tabular}
\caption*{Epicycle dimensions and maximum equations (degrees where $R=60$)}
\end{table}

\vspace{0.5ex}
\pagelabel{1381--1385}

\bilingualpair{%
सू.~भा.---उदयकालः । ग्रहस्य उदयकालः तस्य राशेः, क्षितिज्याया अनुप्रवणतायाः, स्थलस्य च अक्षांशात् गणनीयः । चरजा (उदयास्तभूमिः) आरोहणभिन्नायां प्रयुज्यते ॥
}{%
\textsc{Sū.~Bhā.}---The rising time of a planet is computed from its longitude, the oblique ascension of the ecliptic, and the latitude of the place. The \textit{carajā} (ascensional difference) is applied to find the exact moment of heliacal rising.
}

\vspace{0.5ex}
\pagelabel{1386--1390}

\bilingualpair{%
सू.~भा.---समाप्तिः । एवं ग्रहाणां सूर्यचन्द्रमसोः च गतिः ससंस्कारा निरूपिता । एभिः विधिभिः यात् काले खगolahां स्थितिः गणनीया ॥
}{%
\textsc{Sū.~Bhā.}---Thus have been explained the motions of the planets, the Sun and the Moon, with their corrections. By these methods, the positions of the celestial bodies may be computed for any desired time.
}

\vspace{0.5ex}
\pagelabel{1391--1398}

\bilingualpair{%
\textbf{समाप्तिः}---एषः ब्राह्मस्फुटसिद्धान्तः ब्रह्मगुप्तेन जिष्णुघाटसुतेन शकाब्दे ६२८ (५५० ई.) वयसि त्रिंशे सञ्चिक्षेप । एतत् पठतां सर्वेषां बोधाय भवतु ॥
}{%
\textbf{Colophon.}---This \textit{Brāhmasphuṭasiddhānta} was composed by Brahmagupta, son of Jṣṇughāṭa, in the year 628 of the Śaka era (550~CE), at the age of 30. May it bring clarity to all who study it.
}

\newpage

%% ========================================================================
\section*{Colophon and Notes / \dev{समाप्तटिप्पणी}}
%% ========================================================================

\bilingualpair{%
\textbf{अस्य पाठस्य विषयाः:} ब्रह्मगुप्तकृतं गणितज्योतिषकोशम् एतत्, यस्मिन् धनऋणशून्यगणितानि, कुट्टकरीतिः, वर्गप्रकृतिः, ग्रहस्फुटीकरणं च सविस्तरं वर्णितम् ॥
}{%
\textbf{Key Mathematical Contents:} The \textit{Brāhmasphuṭasiddhānta} is notable for: (1)~algebraic rules for arithmetic with positive and negative numbers, zero, and surds; (2)~the kuṭṭaka (pulverizer) method for linear Diophantine equations; (3)~the varga-prakṛti method for Pell's equation; (4)~astronomical calculations for planetary positions, eclipses, and conjunctions; (5)~trigonometric interpolation formulas.
}

\vspace{1.5ex}

\bilingualpair{%
\textbf{सम्पादकीयम्:} एतत् पाठसंस्करणं मूलमुद्रितसंस्करणात् सङ्कलितम् । ब्रह्मगुप्तेन उज्जयिन्यां ६२८ ई. मणि संकलितं गणितज्योतिषयोः शिखरं इदं ग्रन्थकोशम् ॥
}{%
\textbf{About this edition:} This text edition was prepared from the original printed edition. Composed in 628~CE at Ujjain, the work represents one of the pinnacles of Indian mathematical astronomy.
}

\vspace{2ex}

\noindent\rule{\textwidth}{0.4pt}
\begin{center}
\small\textit{Typeset in \XeLaTeX\ with Noto Serif Devanagari\\
Bilingual edition generated: \today}
\end{center}

\end{document}
%% ========================================================================



% ============================================================
% Source: translations/en/indian/brahmagupta-brahmasphutasiddhanta-pt2-chunk5_bilingual.tex
% ============================================================

%% ========================================================================
%% Brahmasphutasiddhanta Part 2 — Chunk 5 (Pages 1399-1538)
%% BILINGUAL EDITION: Sanskrit (Devanagari) || English + IAST
%% Author: Brahmagupta (c. 598-668 CE)
%% Commentator: Sudhakara Dvivedi
%% ========================================================================

\documentclass[11pt,a4paper]{article}

%% -------- Core Packages --------
\usepackage{fontspec}
\usepackage{polyglossia}
\usepackage{amsmath,amssymb}
\usepackage{geometry}
\usepackage{xcolor}
\usepackage{titlesec}
\usepackage{fancyhdr}
\usepackage{hyperref}

%% -------- Page Geometry --------
\geometry{
  paperwidth=210mm,paperheight=297mm,
  left=12mm,right=12mm,top=22mm,bottom=22mm,
  headheight=14pt,footskip=30pt
}

%% -------- Language Setup --------
\setmainlanguage{english}
\setotherlanguage{sanskrit}

%% -------- Font Configuration --------
\setmainfont{TeXGyreTermes}
\setsansfont{TeXGyreHeros}
\setmonofont[Scale=MatchLowercase]{TeXGyreCursor}

%% Devanagari Font -- system name for better embedding
\newfontfamily\devanagarifont[Script=Devanagari,Path=/usr/share/fonts/truetype/noto/,UprightFont=NotoSerifDevanagari-Regular.ttf]{NotoSerifDevanagari}

%% -------- Colors --------
\definecolor{titlecolor}{RGB}{45,55,72}
\definecolor{sectioncolor}{RGB}{55,65,81}
\definecolor{rulecolor}{RGB}{180,160,140}
\definecolor{sktcolor}{RGB}{60,30,10}

%% -------- Section Formatting --------
\titleformat{\section}
  {\normalfont\Large\bfseries\color{sectioncolor}}
  {\thesection}{1em}{}
\titleformat{\subsection}
  {\normalfont\large\bfseries\color{sectioncolor}}
  {\thesubsection}{1em}{}

%% -------- Header/Footer --------
\pagestyle{fancy}
\fancyhf{}
\fancyhead[L]{\small\textit{Brahmasphutasiddhanta Part II --- Chunk 5 (Bilingual)}}
\fancyhead[R]{\small\thepage}
\renewcommand{\headrulewidth}{0.4pt}

%% -------- Custom Commands --------
\newcommand{\dev}[1]{{\devanagarifont #1}}
\newcommand{\skt}[1]{{\devanagarifont\color{sktcolor} #1}}
\newcommand{\ednote}[1]{\textcolor{notecolor}{\textit{#1}}}

%% Bilingual environment: LEFT Sanskrit, RIGHT English
\newenvironment{bilingual}{%
  \par\noindent
  \begin{minipage}[t]{0.47\textwidth}%
}{%
  \par\end{minipage}%
}
\newcommand{\rightcol}{%
  \par\end{minipage}%
  \hfill
  \begin{minipage}[t]{0.47\textwidth}%
}

%% Shloka environment (within minipages)
\newenvironment{shloka}{%
  \smallskip%
  \devanagarifont\color{sktcolor}%
}{%
  \smallskip%
}

%% -------- Hyperref Setup --------
\hypersetup{
  colorlinks=true,
  linkcolor=sectioncolor,
  citecolor=sectioncolor,
  urlcolor=sectioncolor,
  pdfencoding=auto,
  pdftitle={Brahmasphutasiddhanta Part 2 -- Chunk 5 -- Bilingual Edition},
  pdfauthor={Brahmagupta}
}

%% ========================================================================
\begin{document}

%% ===== TITLE PAGE =======================================================
\begin{titlepage}
\begin{center}
\vspace*{2cm}

{\fontsize{14}{16}\selectfont\devanagarifont ब्राह्मस्फुटसिद्धान्त}

\vspace{0.8cm}

{\LARGE\bfseries\color{titlecolor} Brāhmasphuṭasiddhānta}

\vspace{0.5cm}

{\Large\itshape of Brahmagupta}

\vspace{1.5cm}

{\large\bfseries Part II --- Chunk 5}

\vspace{0.5cm}

{\large Bilingual Sanskrit-English Edition}

\vspace{0.3cm}

{\small Original Pages 624--743 / PDF Pages 561--700}

\vspace{2cm}

\rule{0.6\textwidth}{0.5pt}

\vspace{0.5cm}

{\normalsize Sanskrit text with \devanagarifont देवनागरी}

\vspace{0.2cm}

{\normalsize English translation with IAST transliteration}

\vspace{0.5cm}

\rule{0.6\textwidth}{0.5pt}

\vfill

{\small Edited with Hindi Commentary (Ṭīkā) and Upapatti (Proofs)}

{\small by Paṇḍita Sudhākara Dvivedī}

\vspace{1cm}

{\footnotesize Bilingual edition: left column Sanskrit, right column English}

\end{center}
\end{titlepage}

%% ===== INTRODUCTION =====================================================
\section*{Introduction}

\begin{bilingual}
\textbf{संपादकीयम्} --- एतद्ग्रन्थे ब्रह्मस्फुटसिद्धान्तस्य द्वितीयभागस्य पञ्चमं खण्डं समावेशितम्। अस्य खण्डस्य विषयाः ग्रहच्छाया, भग्रहत्य, युक्तिलक्षणम्, तन्त्रपरीक्षा, कर्ण, द्विक्षेप, युग, तथा गणिताध्यायः इति। ब्रह्मगुप्तेन (५९८--६६८ ई.) अयं ग्रन्थः संहितः।
\rightcol
\textbf{Note} --- This volume contains the fifth chunk of Part 2 of Brahmagupta's \textit{Brāhmasphuṭasiddhānta}. Its subjects include planetary shadows (grahacchāyā), planetary conjunctions (bhagrahatya), characteristics of proof (yukti-lakṣaṇa), examination of astronomical systems (tantra-parīkṣā), hypotenuse (karṇa), celestial latitude (dvikṣepa), cosmic cycles (yuga), and the chapter on mathematics (gaṇitādhyāya). Composed by Brahmagupta (598--668 CE).
\end{bilingual}

\medskip
\noindent\rule{\textwidth}{0.4pt}


%% ========================================================================
%% CHAPTER 1: GRAHACHCHHAYADHIKARA
%% ========================================================================
\section[\dev{ग्रहच्छायाधिकारः}]{\dev{ग्रहच्छायाधिकारः} --- Planetary Shadow Calculations}

\begin{bilingual}
\textbf{मूलपाठः} --- छाया कर्णोच्छ्यासार्धगोले परिगता सती कर्णच्छुताछाया व्यस्तगोला भवति। छाया कर्णगोले तु पलभा शङ्कुतलतुल्या भवति, पलभा च सर्वदैवोत्तरा, तयोः संस्कारतच्छायामुखपूर्वापरसुतयोरन्तरं भुजो भवेत्, छायामुखबन्धेन शङ्कुमूलं बोध्यम् ॥४१॥
\rightcol
\textbf{Text (v.\,41)} --- When a shadow passes over the karna-ucchyāsa-ardha-gola (half-sine sphere of the hypotenuse), the shadow turning from the hypotenuse becomes vyasta-golā (reversed/inverted). On the karna-gola (hypotenuse sphere), the palabhā (equinoctial shadow) equals the base of the gnomon (śaṅku-tala), and the palabhā is always directed northward. Of these two, applying saṃskāra (correction), the difference between the eastern-western threads at the face of the shadow becomes the bāhu (base), and the root of the gnomon is to be understood by means of the bandha (construction) at the shadow-face.
\end{bilingual}

\medskip

\begin{bilingual}
\textbf{टीका} --- छाया कर्णोच्छ्यासार्धगोले परिगता सती कर्णच्छुताछाया व्यस्तगोला भवति, छाया कर्णगोले तु पलभा शङ्कुतलतुल्या भवति, पलभा च सर्वदैवोत्तरा, तयोः संस्कारतच्छायामुखपूर्वापरसुतयोरन्तरं भुजो भवेत्, छायामुखबन्धेन शङ्कुमूलं बोध्यम् ॥४१॥
\rightcol
\textbf{Commentary (Ṭīkā)} --- When the shadow, having gone around the sphere of the hypotenuse and its altitude, becomes inverted through the hypotenuse, then on the hypotenuse-sphere the palabhā equals the base of the gnomon; the palabhā is always northward. Applying correction to them, the difference between the east-west threads at the shadow-mouth becomes the base (bhuja); by the bandha at the shadow-mouth the root of the gnomon is to be known.
\end{bilingual}

\medskip

\begin{bilingual}
\textbf{उपपत्तिः} --- छाया कर्णोच्छ्यासार्धगोले परिगता सती कर्णच्छुताछाया व्यस्तगोला भवति, छाया कर्णगोले तु पलभा शङ्कुतलतुल्या भवति, पलभा च सर्वदैवोत्तरा, तयोः संस्कारतच्छायामुखपूर्वापरसुतयोरन्तरं भुजो भवेत्, छायामुखबन्धेन शङ्कुमूलं बोध्यम् । सिद्धान्तशेखरे ``पूर्वापरशङ्कुतलान्तरं यत् बाहुः स एवोत्तर दक्षिणः स्यात्'' इति श्रीपतिनारायणाचार्यैर्मेव कथ्यते ॥४१॥
\rightcol
\textbf{Proof (Upapatti)} --- In the Siddhāntaśekhara it is thus stated by Śrīpati Nārāyaṇācārya: ``That which is the difference between the east-west gnomon-bases is the bāhu; that itself becomes the north-south (koṭi).'' Thus the demonstration is established.
\end{bilingual}

\medskip
\noindent\rule{\textwidth}{0.4pt}

\begin{bilingual}
\textbf{मूलपाठः} --- इदानीं भुजसाधनमाह।

\dev{प्राच्यपरायाः शङ्कूस्तलं तदप्रं ग्रहं मे च} ॥४२॥
\rightcol
\textbf{Text (v.\,42)} --- Now he states the method for finding the base (bhuja):

\textit{prācyaparāyāḥ śaṅkūstalaṃ tadapraṃ grahaṃ me ca}

The base of the gnomon extends east-west; that astra (side) and the planet---these become the measure.
\end{bilingual}

\medskip

\begin{bilingual}
\textbf{टीका} --- \skt{प्राच्यपरायाः} पूर्वापरतः शङ्कूस्तलं शङ्कुमूलपर्यन्तं स्यात् । तदप्रं शङ्कुवप्रं ग्रहं वा मे च भवति ॥४२॥
\rightcol
\textbf{Commentary} --- \textit{prācyaparāyāḥ} (east-west): the base of the gnomon extends to the root of the gnomon. The astra (flank) of the gnomon, or the planet, becomes the measure.
\end{bilingual}

\medskip
\noindent\rule{\textwidth}{0.4pt}

\begin{bilingual}
\textbf{मूलपाठः} --- अथोपपतिः । शङ्कुमूलप्राच्यपरान्तं भुजसाधनं सिद्धान्तयुक्तिच्छा स्फुटम्।

\skt{शङ्कुवप्रं ग्रहं भवतोति स्पष्टम्} ॥४२॥
\rightcol
\textbf{Proof} --- Now the proof: from the east-west limit at the root of the gnomon, the demonstration of the base is clear according to the siddhānta-yukti.

``The side of the gnomon and the planet''---this is clear.
\end{bilingual}

\medskip

\begin{bilingual}
\textbf{टीका} --- \skt{तुल्यादिशोः} (एकादिशकयोः) प्राच्यशङ्कुतलयोरैक्यं (योगः) तथा द्व्यादिशोः (द्व्यादिशकयोः) प्राच्यशङ्कुतलयोरन्तरं प्राच्यपरायाः (पूर्वापर रेखातः) शङ्कूस्तलं (शङ्कुमूलपर्यन्तं) भुजो भवेत् । तदप्रं (शङ्कुवप्रं) ग्रहं नक्षत्रं च भवतीति ॥४२॥
\rightcol
\textbf{Commentary} --- Of the same directions (ekādiśa), the sum of the east-west gnomon-bases; and of two directions (dvyādiśa), the difference of the east-west gnomon-bases---from the east-west line, the base extending to the root of the gnomon becomes the bhuja (base). The astra (side) of the gnomon and the star become the constituents.
\end{bilingual}

\medskip
\noindent\rule{\textwidth}{0.4pt}

\begin{bilingual}
\textbf{मूलपाठः} --- \skt{शङ्कुतलं प्राच्यपरान्तं भुजो दक्षिणोत्तरं कर्णः} ।

\skt{द्विजया तद्विपरीतमूलं द्विक्षेपमध्यतः कोटिः} ॥४३॥
\rightcol
\textbf{Text (v.\,43)} --- The base extending east-west from the gnomon-base (śaṅku-tala) becomes the bhuja; the karna (hypotenuse is) the north-south. By the dvijaya (second victory), the root of the inverse from the middle of the dvikṣepa (double-latitude) becomes the koṭi.
\end{bilingual}

\medskip

\begin{bilingual}
\textbf{टीका} --- \skt{शङ्कुतलं प्राच्यपरान्तं} (शङ्कुमूलपूर्वापरसुतयोरन्तरं) दक्षिणोत्तरं भुजो भवति, द्विजया कर्णस्योत्तरद्विजया भुजयोर्विपरीतं मूलं द्विक्षेपमध्यतः (कन्द्रात्) कोटिर्भवेदिति ॥४३॥
\rightcol
\textbf{Commentary} --- \textit{śaṅkutalaṃ prācyaparāntaṃ} means the difference between the east-west threads at the root of the gnomon; the north-south becomes the bhuja. By the second application, the root of the inverse of the hypotenuse and the base---from the middle of the double-latitude (dvikṣepa-madhya)---becomes the koṭi (perpendicular).
\end{bilingual}

\medskip
\noindent\rule{\textwidth}{0.4pt}

\begin{bilingual}
\textbf{उपपत्तिः} --- \skt{द्विक्षेपसम्पातगतयोर्यथाप्राक्पूर्वापरसुतयोरैं यो लम्बः स भुजः} । शङ्कुमूलपूर्वापरसुतयोरपि यो लम्बः स एव भुजः । एकदिशोर्या शङ्कुतलयोगेन द्व्यादिशस्तयोरन्तरेण शङ्कुमूलपूर्वापरसुतपर्यन्तं पूर्वापरसुतयोरपि लम्बरूपो भुजो भवति । वा क्षा
\rightcol
\textbf{Proof} --- Of the two meeting at the double-latitude, that which is the perpendicular dropped from the east-west threads is the bhuja. Also, the perpendicular dropped from the east-west threads at the root of the gnomon is the same bhuja. By the sum of the gnomon-bases of one direction, and by the difference of the two directions, extending to the east-west threads at the root of the gnomon, the perpendicular from the east-west threads also becomes the bhuja.
\end{bilingual}


%% ========================================================================
%% CHAPTER: YUKTI-LAKSHANA
%% ========================================================================
\newpage
\section[\dev{युक्तिलक्षणम्}]{\dev{युक्तिलक्षणम्} --- Characteristics of \textit{Yukti} (Proofs)}

\begin{bilingual}
\textbf{मूलपाठः} --- इदानीं युक्तिलक्षणमाह।

\skt{उभे मानक्याखण्डि ग्रहयोर्मध्यान्तरं युक्तिगृह्योः} ।

\skt{समलिप्तिकयोगं हृत्वैवद्विके स्फुटमानयोगाखण्डित्} ॥६९॥
\rightcol
\textbf{Text (v.\,69)} --- Now he states the characteristics of yukti (proof):

\textit{ubhe mānakyākhaṇḍi grahayor madhyāntaraṃ yukti-gṛhyoḥ}

When the difference between two planets of equal minutes (māna) is divided by the sum of the vi-ākhaṇḍa (discs), at the time of eclipse there is a yukti (conjunction).
\end{bilingual}

\medskip

\begin{bilingual}
\textbf{टीका} --- \skt{मानक्याखण्डिके} इतिद्विगुसमासमोक्ते न स्फुटा व्याख्याताप्रतिस्तम् ॥६९॥
\rightcol
\textbf{Commentary} --- The dvigu compound \textit{mānakyākhaṇḍi} is not clearly explained in this statement; its meaning requires further elucidation.
\end{bilingual}

\medskip
\noindent\rule{\textwidth}{0.4pt}

\begin{bilingual}
\textbf{मूलपाठः} --- \skt{समकलयोर्ग्रहयोर्मध्यान्तरं} (कन्द्रान्तरं) यदि ग्रहयोरन्तरं तयोर्विभक्तयोगाखण्डितेन (न्यूने) तद्युक्तिर्भवेद्विधिद्व्यसहावादाच्छादनं भवेत्-तथाद्विपुरःस्यो ग्रहं ग्राह्यः स्यात्-प्रभः-स्यो ग्रहं ग्राहकः स्यादित्यर्थः । ग्रहयोरन्तरं स्फुटमानयोगाखण्डि (स्फुटविभक्तयोर्गाखण्डितं) द्विके सति तयोर्युति (आच्छादनं ग्रहावतं) ने स्यादिति । सिद्धान्तशेखरे ``मानक्याखण्डि धुवरविकरे स्यान्न्मेदोद्विके तु न्यूने मेदो ग्रहणवाद्विहच्छादकोऽस्तनः स्यात्'' श्रीपत्युक्तिसिद्धान्तसिद्धान्तगोम्मगो `मानक्याखण्डि धुवरविकरेऽस्ये भवेद् योग' इत्यादि भास्करोक्तं चाचार्यैकानुरूपमेवास्ति ॥६९॥
\rightcol
\textbf{Expanded Text} --- When the difference between two planets of equal minutes (\textit{sama-kala})---i.e., the latitude-difference (kendra-antara)---if the difference between the two planets is less than the sum of their discs (\textit{vi-bhakta-yogākhaṇḍita}), then there is a yukti (conjunction), that is, an ācchādana (eclipse) like a grahaṇa; thus, the upper planet becomes the grāhya (the eclipsed) and the lower planet becomes the grāhaka (the eclipser). If, when divided by the sum of the discs at the time of eclipse, their yukti (meeting) does not occur. In the Siddhāntaśekhara: ``\textit{mānakhyākhaṇḍi dhruvara-vikare syān medo-dvike tu nyūne medo grahaṇa-vād-vihac-chādako 'stanaḥ syāt}''---this statement of Śrīpati and the statement of Bhāskara, ``\textit{mānakhyākhaṇḍi dhruvara vikare 'sye bhaved yogaḥ}'' etc., are in agreement with the Ācārya's teaching.
\end{bilingual}

\medskip

\begin{bilingual}
\textbf{उपपत्तिः} --- अथ युक्ति लक्षण को कहते हैं।

\skt{सम कलात्मक दो ग्रहों के अन्तर यदि दोनों ग्रहों के विभक्तयोगाखण्ड से न्यून हो तब ग्रह द्वय की युक्ति होती है अर्थात् ग्रहणवत् आच्छादन (उपरिस्थ ग्रह ग्राह्य और अभःस्थित ग्रह ग्राहक) होता है। यदि दोनों ग्रहों के अन्तर स्फुट विभक्तयोगाखण्ड से अधिक हो तब ग्रह की युक्ति नहीं होती। सिद्धान्त शेखर में `मानक्याखण्डि धुवरविकरे स्यान्न्मेदोः द्विके तू' इत्यादि विज्ञान भाष्य में लिखित स्लोकांश से श्रीपति सिद्धान्त शिरोमणि में `मानक्याखण्डि धुवर विकरेऽस्ये भवेद् योगः' इत्यादि से भास्कराचार्य ने भी आचार्यैक के अनुरूप ही कहा है इति ॥६९॥}
\rightcol
\textbf{Proof} --- Now, the definition of yukti is stated: If the difference between two planets of equal minutes is less than the sum of their \textit{vi-bhakta-yogākhaṇḍa} (discs), then there is a yukti of the two planets, that is, an ācchādana (covering) like an eclipse (the upper planet is the eclipsed, the lower planet is the eclipser). If the difference between the two planets is greater than the sum of the \textit{sphuṭa-vibhakta-yogākhaṇḍa}, then there is no yukti of the planets. In the Siddhāntaśekhara, from the verse beginning ``\textit{mānakhyākhaṇḍi dhruvara-vikare syān medoḥ dvike tū}'' etc., in the Vijñāna-bhāṣya---Śrīpati in the Siddhānta-śiromaṇi, from ``\textit{mānakhyākhaṇḍi dhruvara vikare 'sye bhaved yogaḥ}'' etc., Bhāskarācārya too has spoken in accordance with the Ācārya.
\end{bilingual}


%% ========================================================================
%% CHAPTER: BHA-GRAHATYADHIKARA
%% ========================================================================
\newpage
\section[\dev{भग्रहत्यधिकारः}]{\dev{भग्रहत्यधिकारः} --- Planetary Conjunctions}

\begin{bilingual}
\textbf{मूलपाठः} --- \skt{भुजकोटिकर्णानां यद्भुज्यासार्धं परिणामनार्थं तान् यद्भुज्यान् व्यासार्धहृतान् गुणकः कुर्यात्} । एवं ग्रहयोर्गति तयोर्यथाकृत्कुब्रान्तरं भवेद्यत् शङ्कुकुव्रान्तरामावः शङ्कुमूलादीनां च समत्वात् । भ्रमदा तयोः शङ्कुवोरन्तरं भवेदेव । शेषं स्पष्टार्थम् ॥५७॥
\rightcol
\textbf{Text (v.\,57)} --- Of the bhuja, koṭi, and karṇa, whatever bhuja-half is to be transformed, making it into gunakas (multipliers) divided by the semi-diameter (vyāsārdha). Thus, the motion of the two planets---according to their construction the difference between them becomes the ku-bra-antara; the difference between the gnomon-sides (śaṅku-kuvra) is due to the equality of the roots of the gnomon etc. By rotation, the difference between their two gnomons comes about. The rest is clear.
\end{bilingual}

\medskip

\begin{bilingual}
\textbf{टीका} --- \skt{भुजकोटिकर्णानां यद्भुज्यासार्धं परिणामनार्थं तान् यद्भुज्यान् व्यासार्धहृतान् गुणकः कुर्यात्} । एवं ग्रहयोर्गति तयोर्यथाकृत्कुब्रान्तरं भवेद्यत् शङ्कुकुव्रान्तरामावः शङ्कुमूलादीनां च समत्वात् । भ्रमदा तयोः शङ्कुवोरन्तरं भवेदेव । शेषं स्पष्टार्थम् ॥५७॥
\rightcol
\textbf{Commentary} --- The teacher explains: Of the base, perpendicular, and hypotenuse, the half-base for transformation---those half-bases divided by the semi-diameter should be made into multipliers. Thus the motion of the two planets---their difference as constructed---the difference of the gnomon-sides arises from the equality of the gnomon-roots. By the revolution, the difference of their gnomons occurs. The remainder is self-evident.
\end{bilingual}

\medskip
\noindent\rule{\textwidth}{0.4pt}

\begin{bilingual}
\textbf{मूलपाठः} --- \skt{अथोपपतिः । गोल्युत्थया भुजकोटिकर्णैःशङ्कुकुर्मस्थानवशातः स्फुटा भुजामर्यधायां चापिहटी युक्तिराचार्यैषां प्रतिपादिता} ॥५७-५९॥
\rightcol
\textbf{Proof} --- Now the upapatti: From the sphere arise the bhuja, koṭi, and karṇa; from the position of the gnomon-tortoise (śaṅku-kūrma), the true bhuja---in this manner, with its application, the proof demonstrated by the Ācārya is established.
\end{bilingual}

\medskip

\begin{bilingual}
\textbf{टीका} --- \skt{यदा द्व्येक्षणाडिकाभिर्न्तुल्याः कुत्वा तदा योगो भवेत् । योगे} (महायुति) \skt{ग्रहयोर्मध्यान्तरे} (विभक्तेकन्द्रान्तरे सति) \skt{ष्टज्ञोनितिवत्} (चन्द्रष्टज्ञोन-नितिवत्) \skt{समान्यदिशोः} (एकभिन्नदिशकयोः) \skt{प्राच्यशङ्कुवप्रयोः} (प्राच्यशङ्कुत्ल-योः) \skt{संयोगान्तरं} (योगान्तरं) \skt{बाहु} (भुजो) \skt{कार्यवर्धन्मध्यात्तयोगवतोऽन्त-हयोः प्रत्येकस्य शङ्कुतलस्याप्रायिकं योगान्तरं तयोर्भुजो भवत् इत्यर्थः} ।
\rightcol
\textbf{Commentary} --- When by two observation-nāḍīs they are made unequal, then there is a yoga (conjunction). In the yoga (mahā-yoga), when there is a difference between the two planets (\textit{grahayor-madhyāntare}, i.e., at the divided kendra), like the moon-sixth-day-increase (\textit{ṣaṭ-jñonitivat}), of the same directions (\textit{samyad-śoḥ}, i.e., one-and-different-directions), the difference of the east-west gnomon-flanks (\textit{prācya-śaṅku-vaprayoḥ})---the saṃyogāntara (union-difference)---the bāhu (base)---the meaning is: from the middle of the one-to-be-increased, the yogāntara (difference of union) of each gnomon-base being applied, the bhuja of them is produced.
\end{bilingual}

\medskip

\begin{bilingual}
\textbf{टीका (contd.)} --- \skt{ग्रहयो-रैक्ये कर्णो भवतः । स्वकर्णो भुजाकृतिवियोगाप्ते} (\skt{स्वकर्णोभुजयोर्वर्गयोरन्तर-मूले}) \skt{कोटी भवतः । एवं ग्रहयोः कोटिभुजकर्णैःशङ्कुकृत् यद्भुज्यान्} (\skt{कल्पितषडष्टि-गुणान्}) \skt{व्यासार्धहृतान्} (\skt{विज्याध्मातान्}) \skt{कृत्वाथैकेन कर्तव्येन कोटिभुजकर्णैः शङ्कुवो यद्भुज्यासार्धपरिणाता भवेतुः} ॥५७-५९॥
\rightcol
\textbf{Commentary (contd.)} --- In the sum of the two planets, the karṇa is produced. Their own karṇa---at the root of the difference of the squares of the bhuja and karṇa---the koṭi is produced. Thus, by the koṭi-bhuja-karṇa of the two planets, by the gnomon-construction, whatever half-bhuja is to be transformed (\textit{yad-bhujyārdha-pariṇāta}), having made them divided by the semi-diameter (\textit{vyāsārdha-hṛtān}), they become transformed.
\end{bilingual}

\medskip

\begin{bilingual}
\textbf{टीका (contd.)} --- \skt{प्रागपरयोः} (\skt{यद्भुज्यासार्धपरिणातपूर्वा-पररेखानुल्लपयोः}) \skt{प्राच्यो} (\skt{पूर्वदिशि}) \skt{पश्चात्} (\skt{पश्चिम दिशि}) \skt{पृथक् कोटी प्रसार्य कोटिप्रसार्या बाहु} (\skt{भुजो}) \skt{दल्वा दिक्षुमध्यतो} (\skt{मध्य विद्धोः}) \skt{भुजाप्रान्तो} (\skt{भुजाप्रा-विद्धलेखो}) \skt{कर्णां च दल्वा बाहुप्रयोः} (\skt{स्वभुजाप्रयोः}) \skt{स्वशङ्कु} द्वयोः, \skt{मध्यात् तदप्रान्ते} (\skt{स्वशङ्कुवप्रयोर्लेखने}) \skt{यथी च पूर्वकल्पिते द्वे} ॥५७-५९॥
\rightcol
\textbf{Commentary (contd.)} --- Of the east-west (transformed-half-bhuja-east-west-line-threads), in the east (eastern direction), in the west (western direction), extending the koṭi separately; the bāhu (base) extended by the koṭi; dividing, from the middle of the directions (from the divided middle), the bhuja-extent (bhuja-prā-viddha-lekha); dividing the karṇa too, of the bāhu-extents (their own bhuja-extents), their own gnomons---two; from the middle to their astra-extent (at the drawing of their own gnomon-sides), as before, two previously assumed.
\end{bilingual}

\medskip

\begin{bilingual}
\textbf{टीका (contd.)} --- \skt{द्विक्षेपमध्यस्थहच्छाया} (\skt{मध्यबिन्दुस्थापितेन चक्षुषा}) \skt{स्वशङ्कुवप्रं पृथक् ग्रहं दर्शयेत् । योगे} (\skt{ग्रहयो-योगे}) \skt{तयोः शङ्कुवकुव्रान्तरं भवेद्यत् शङ्कुवन्तरामावः,} (\skt{शङ्कुभुजादीनां सम-त्वात्}) \skt{अन्यथा} (\skt{ग्रहयुतिर्मिन्नेऽस्वरे}) \skt{तयोः शङ्कुवोरन्तरं भवेदेव, भ्रमे} (\skt{योगे शङ्कु-कुव्रान्तरसन्तरमेवान्यथा हि तयोः} ताड्रूपाटे \skt{योगे} (\skt{महायुति}) \skt{शङ्कुकुव्रान्तरं ग्रहयोरन्तरं नेयं भ्रमदा} (\skt{ग्रहयुतिर्मिन्नकल्पे}) \skt{ऐषं नेयमिति} ॥५७-५९॥
\rightcol
\textbf{Commentary (contd.)} --- The dvikṣepa-madhya-stha-cchāyā (shadow situated at the middle of the double-latitude, i.e., by the eye situated at the middle point) shows separately its own gnomon-side and the planet. In the yoga (union of the two planets), their difference of gnomon-kuvra occurs because of the equality of the gnomon, bhuja, etc.; otherwise (if the planetary union is not at the same point), their difference of gnomons alone occurs. In the revolution (in the yoga, only the difference of gnomon-kuvra-remaining, otherwise of them), in the great yoga, the difference of gnomon-kuvra should not be taken as the difference between the planets; this should be known in the planetary-union-minimum-kalpa.
\end{bilingual}


%% ========================================================================
%% CHAPTER: TANTRA-PARIKSHA-ADHYAYA
%% ========================================================================
\newpage
\section[\dev{तन्त्रपरीक्षाध्यायः}]{\dev{तन्त्रपरीक्षाध्यायः} --- Examination of Astronomical Systems}

\begin{bilingual}
\textbf{संपादकीय टिप्पणी} --- एतस्मिन् अध्याये ब्रह्मगुप्त आर्यभटीयज्योतिःशास्त्रस्य कठोरपरीक्षां करोति। आर्यभट (४७६ ई.) प्रथमस्य ग्रहणसिद्धान्तस्य निराकरणम् अत्र विस्तरेण प्रदर्शितम्।
\rightcol
\textbf{Note} --- In this chapter Brahmagupta subjects the astronomical system of Āryabhaṭa to rigorous examination. The refutation of Āryabhaṭa I's (476 CE) planetary theory is here demonstrated in detail.
\end{bilingual}

\medskip
\noindent\rule{\textwidth}{0.4pt}

\begin{bilingual}
\textbf{मूलपाठः} --- \skt{इदानीमार्यभटोक्तं युगं खण्डयति}।

\skt{आर्यभटो युगपादांशाने यातानाह कलियुगादौ यत्} ।

\skt{तस्य क्षुतान्तर्यस्मात् स्वयुगाष्ठतो न सत् तस्मात्} ॥३॥
\rightcol
\textbf{Text (v.\,3)} --- Now he refutes the yuga stated by Āryabhaṭa:

\textit{āryabhaṭo yugapādāṃśāne yātānāha kaliyugādau yat}

Āryabhaṭa stated that two quarter-yugas had elapsed at the beginning of Kali-yuga. From what cause? Therefore that is not true; hence his own yuga-aṣṭa (eighth of a yuga) is not valid.
\end{bilingual}

\medskip

\begin{bilingual}
\textbf{टीका} --- \skt{आर्यभटः कलियुगादौ द्विन् युगपादान् यातान् आह कथितवान् । यस्मात् तद्ग्रथतः} (\skt{दृष्ट्या मध्यमाधिकारे क्षद्विशतिस्लोकस्य मदीया व्याख्या}) \skt{यस्मात् कारणात् तस्मात् तस्य स्वयुगाष्ठतो} (\skt{तदेकस्य युगस्यादिर्न्यस्यात् इतिद्वि कृतान्तः कृतयुगमध्ये भवतस्तस्मात् तच्छुं न सत्} ॥३॥
\rightcol
\textbf{Commentary} --- Āryabhaṭa stated (\textit{āha}) that at the beginning of Kali-yuga, two yuga-pādas (quarter-yugas) had elapsed (\textit{yātān}). From what cause?---from that reason, therefore his own yuga-aṣṭa---of one yuga, the beginning established, thus two in the Kṛta-yuga---from that, that is not true (\textit{na sat}).
\end{bilingual}

\medskip

\begin{bilingual}
\textbf{उपपत्तिः} --- \skt{आर्यभटमते एक युगान्ताद्यस्यारम्भात् कलियुगादिपर्यन्तं त्रयो युगपादाः} $= \dfrac{3 \times 4320000}{4} = 3240000$ । \skt{आचार्यमते च, क्षु+त्रे+आ} $\dfrac{4320000 \times 9}{10} = 3888000$ \skt{द्वयोरन्तरे वर्षाणि} $= 648000$ \skt{एतानि चाचार्यमतेन संख्याधिकत्वात् कृतयुगमध्येऽपि आर्यभटोत्तयुगाष्ठतो कृतयुगान्तः} ॥३॥
\rightcol
\textbf{Proof} --- According to Āryabhaṭa's system, from the commencement of one yuga to the beginning of Kali-yuga, three yuga-pādas $= \frac{3 \times 4320000}{4} = 3240000$. And according to the Ācārya's (Brahmagupta's) system, Kṛta + Tretā + etc. $\frac{4320000 \times 9}{10} = 3888000$. The difference between the two $= 648000$ years. These, by the excess of number in the Ācārya's system, even in the middle of Kṛta-yuga, Āryabhaṭa's eighth-yuga reaches the end of Kṛta-yuga.
\end{bilingual}

\medskip
\noindent\rule{\textwidth}{0.4pt}

\begin{bilingual}
\textbf{मूलपाठः} --- \skt{अथ पुनरार्यभटोक्तवारादि खण्डयति}।

\skt{श्रीक्षुरो दिनवारो गुरुराद्विकोऽस्य भवति कल्पादौ} ।

\skt{न भवत्यकं यस्मादोच्चुरो विस्तरतस्तस्मात्} ॥११॥
\rightcol
\textbf{Text (v.\,11)} --- Now again he refutes Āryabhaṭa's stated weekday:

\textit{śrīkṣuro dinavāro gurur ādviko 'sya bhavati kalpādau}

The first weekday accepted by him at the beginning of the kalpa is Thursday (guru). \textit{na bhavaty akaṃ yasmād occuro vistaratas tasmāt}---Since this is not established (\textit{na bhavati}), because his extension (\textit{vistara}) is baseless, therefore (it is invalid).
\end{bilingual}

\medskip

\begin{bilingual}
\textbf{टीका} --- \skt{आर्यभटेन स्वतन्त्रे `गुरुदिवसात् भारतात् पूर्वं' मित्यनेन कल्पादौ गुरुवारः स्वीकृतः । तेनायमर्थः । यस्मात् कारणात्-अस्य} (\skt{आर्यभटस्य}) \skt{श्रीक्षुरः} (\skt{स्वीकारः}) \skt{कल्पादावाद्यिको दिनवारो गुरुर्भवति रविनं भवति तस्मादस्योक्षुरः स्वीकारो विस्तरः} (\skt{आधाररहितोऽथप्रामाणिकः}) ॥११॥
\rightcol
\textbf{Commentary} --- By Āryabhaṭa in his own treatise, ``from Thursday before the Bhārata''---by this, Thursday is accepted at the beginning of the kalpa. By this the meaning is: From what cause?---His (Āryabhaṭa's) acceptance (\textit{śrīkṣuraḥ} = \textit{svīkāraḥ}), the first weekday at the beginning of the kalpa being Thursday becomes established; therefore his acceptance is \textit{vistara} (unsupported and unauthoritative).
\end{bilingual}

\medskip

\begin{bilingual}
\textbf{उपपत्तिः} --- \skt{आर्यभटमतेन कलियुगारम्भात् पूर्वं वर्तमानकल्पे क्ष उ मन्वो व्यतीता युगचतुरायां च, तथा तस्मात् द्विसप्ततियुगोरैको मनुतः कल्पादौ कल्पादौ गतयुगानि} $= 72 \times 6 + \dfrac{3}{4} = 432\,\dfrac{3}{4}$, \skt{युगपञ्चितसावनदिनैर्गुणनें जाताः सावनाहर्गणाः} $= 432 \times 1577917500 + \dfrac{3 \times 1577917500}{4}$ $= 432 \times 1577917500 + 3 \times 394479375 \times 3$ । \skt{अयं सप्तभिर्मैं कल्पादौ वारः} $= 5 \times 5 + 3 \times 3 = 25 + 9 = 34 = 6$ । \skt{अयं सैकः कलियुगादौ वारः} $= 7 = 0$ । \skt{अतो यदि गुरुवाराद्-गणानाद्दृष्टे तदा कलियुगादौ गतवारः ० । वर्तमानो गुरुरेव सिद्ध्यत् आर्य-भटमतेन कल्पादौ गुरुवार आयाति} ॥११॥
\rightcol
\textbf{Proof} --- According to Āryabhaṭa's system, before the commencement of Kali-yuga, in the present kalpa, six manvantaras and three quarters of a yuga-catur (yuga-quadruple) had elapsed. Thus from that, from one manu of seventy-two yugas, the elapsed yugas at the beginning of the kalpa $= 72 \times 6 + \frac{3}{4} = 432\,\frac{3}{4}$. Multiplying by the saura-civil days of the yuga-pañcita: $= 432 \times 1577917500 + \frac{3 \times 1577917500}{4}$. The weekday at the beginning of the kalpa by this: the remainder when divided by seven $= 34 = 6$ (Friday); but if Thursday is taken as the starting weekday, then the elapsed weekday at Kali-yuga-beginning $= 0$. The present Thursday itself would be established. Thus by Āryabhaṭa's system, Thursday at the kalpa-beginning is arrived at.
\end{bilingual}

\medskip
\noindent\rule{\textwidth}{0.4pt}

\begin{bilingual}
\textbf{मूलपाठः} --- \skt{अथ पुनरार्यभटोक्तवारादि खण्डयति}।

\skt{क्षधिकैः शतैश्चतुर्भिर्वर्षैः सहस्रैश्चतुर्भिर्यामिरेकः} ।

\skt{युगयातोदिनवारान्तर्मौदित्यकाष्ठे रात्रिकयोः} ॥१३॥
\rightcol
\textbf{Text (v.\,13)} --- Now again he refutes Āryabhaṭa's stated weekday:

\textit{kṣadhikaiḥ śataiś caturbhir varṣaiḥ sahasraiś caturbhir yāmirekaḥ}

With four hundred years increased, and with four thousand, one yāmi. Between the day and weekday, in the aṃśa of the sun and the two nights...
\end{bilingual}

\medskip

\begin{bilingual}
\textbf{टीका} --- \skt{आर्यभटेन प्रथद्वयं रचितम् । एकैकस्मिन् युगसावनदिनानि} $1577917500$ । \skt{लघुगामकल्पदिवे सृष्टिः} । \skt{अन्यैस्मिन् युगसावनदिनानि} $1577917800$ \skt{लघुगामधरात्रौ सृष्टिः} । उभयत्र युगवर्षसंख्या} $4320000$ \skt{समाना एव, युगपञ्चितसावन दिनात्रयं} $= 300$, \skt{अतो प्रथद्वयतो वारगणनया युगवर्षदिनशतत्रयान्तरं भवति, तथानुपातेनैकदिनवारान्तरं च} $\dfrac{4320000}{300} = 14400$ \skt{एतैर्युगातवर्षः} ॥१३॥
\rightcol
\textbf{Commentary} --- By Āryabhaṭa, two systems were constructed. In one, the civil days in a yuga $= 1577917500$; at the kalpa-day of short motion, creation. In the other, the civil days in a yuga $= 1577917800$; at the short-motion night, creation. In both, the number of yuga-years $4320000$ is the same. The three days of yuga-pañcita-sāvana $= 300$. Therefore from the two systems, by weekday computation, there is a difference of three hundred yuga-year-days; and proportionally, the difference of one day-weekday $= \frac{4320000}{300} = 14400$ yuga-years.
\end{bilingual}

\medskip

\begin{bilingual}
\textbf{टीका (contd.)} --- \skt{तेनायमर्थः --- आर्यभटमतेनैदियकाष्ठरा-त्रिकयोदिनवारमध्ये चतुर्दशाभिर्वर्षैः सहस्रै ऋतुभिः शतवर्षैऽधिकैर्गुयातैदिनवारा-न्तरं दिनवारयोरेकान्तरं पततीति वारगणना न स्थिरेति} ॥१३॥
\rightcol
\textbf{Commentary (contd.)} --- By this the meaning is: According to Āryabhaṭa's system, between the day and weekday in the aṃśa of the sun and nights, with fourteen years, thousands, seasons, increased by hundreds of years, the difference between day and weekday falls by one. Thus the weekday computation is not stable.
\end{bilingual}



%% ========================================================================
%% CHAPTER: KARNADHIKARA
%% ========================================================================
\newpage
\section[\dev{कर्णाधिकारः}]{\dev{कर्णाधिकारः} --- Hypotenuse and Diagonal}

\begin{bilingual}
\textbf{मूलपाठः} --- \skt{उपपत्तिः}

\skt{स्फुटादिकारके १३वें सूत्र की उपपत्ति में दिखाया गया है कि प्रथम पद में}

\skt{गतांश क्रमज्या $\times$ परिधि / भोज्या = भुजफल}

\skt{द्वितीय पद में परम भुज फल होता है, तृतीयपद में}

\skt{क्रमज्या $\times$ परिधि / भोज्या = भुजफल, द्वितीय पद में परम भुजफल शून्य होता है}

\skt{चतुष्कपद में परम उत्क्रमज्या $\times$ परिधि / भोज्या = भुजफल, यद्यपि समपदान्ते विषम पदादौ मे भी दोनों पदों की सन्धि होने के कारण परिधिध्य (सम पद और विषम पद में पठित परिधि) ग्रहण करने से फलनाम (फलभाव) नहीं होता है, उससे कुछ हानि नहीं है क्योंकि परिधि पाठ भेद से भावान्तर में अनुपात से परिधि का ग्रहण करना चाहिये ऐसा ने आर्यभट सूचित किया हुआ है इसलिये आचार्यैक खण्डन ठीक नहीं है इति} ॥१९॥
\rightcol
\textbf{Proof} --- In the proof of the 13th rule of the Sphuṭādikāra it is shown that in the first term:

\textit{krama-jyā} $\times$ circumference $/$ \textit{bhojyā} = \textit{bhuja-phala}

In the second term, the maximum bhuja-phala occurs; in the third term:

\textit{krama-jyā} $\times$ circumference $/$ \textit{bhojyā} = \textit{bhuja-phala}. In the second term, the maximum \textit{bhuja-phala} becomes zero.

In the fourth term, the maximum \textit{utkrama-jyā} $\times$ circumference $/$ \textit{bhojyā} = \textit{bhuja-phala}. Although at the end of even terms and at the beginning of odd terms, due to the conjunction of both terms, the \textit{paridhi} (circumference) mentioned in even and odd terms, when taken, does not produce \textit{phala-nāma} (result); there is no harm from it because, by the difference of circumference readings, by the proportion in the difference of terms, the taking of circumference should be done---thus Āryabhaṭa has indicated. Therefore the Ācārya's refutation is not correct.
\end{bilingual}

\medskip
\noindent\rule{\textwidth}{0.4pt}

\begin{bilingual}
\textbf{मूलपाठः} --- \skt{इदानीमार्यभटमतेनानुपातेन यदि परिधिः स्फुटः क्रियते तदापि न समीचीन इति खण्डयति}।

\skt{व्यासार्धहृतो बाहुः परिधिविभेषाहृतः फलोनयुतः} ।

\skt{प्रथमाद्विकोनकश्च यत् तदसत् पदयोः परिधिपातात्} ॥२०॥
\rightcol
\textbf{Text (v.\,20)} --- Now, if by proportion according to Āryabhaṭa's system the circumference is made precise, even then it is not correct---thus he refutes:

\textit{vyāsārdha-hṛto bāhuḥ paridhi-vibheṣāhṛtaḥ phalonayutaḥ}

The bāhu (side) divided by the semi-diameter, divided by the difference of circumferences, with the result added;

\textit{prathamād vikonakaś ca yat tad asat padayoḥ paridhi-pātāt} --- and the vikonaka (angular deviation) from the first, that is invalid because of the \textit{paridhi-pāta} (falling of circumference) in the terms.
\end{bilingual}

\medskip

\begin{bilingual}
\textbf{टीका} --- \skt{बाहुभुजया परिधिविभेषाहृतः परिध्यन्त रहितो व्यासार्धेन विजयया हृतः प्रथमः परिधियं द्वि द्वितीयाद्विकोनकस्तदा क्रमेण प्रथमः परिधिः फलोनयुतः स्फुटः परिधिर्भवेदिति यत् स्फुटपरिध्यानयनं प्रसिद्धं तदप्यार्यभटमतेना-सद्भवति} ।

\skt{कस्मात् । पदयोः परिधि पाठात्} ।

\skt{अर्थतस्तन्मते पदयोः परिधिध्यम्} ।

\skt{तत्र कुतश्चित् प्रथमः समपदीयः कुतश्चित् विषमपदीयः परिधिर्भवति} ।

\skt{अतः संस्कारविधिषु विचरति । बाबलमेतत्} ।

\skt{प्रथमाद्विकोनकश्चापि तात्कालिक संस्कारस्य समीचीनत्वादिति सुधीभिः स विचिन्त्यम्} ॥२०॥
\rightcol
\textbf{Commentary} --- The bāhu (side) by the bhuja, divided by the difference of circumferences, diminished by the end-circumference, divided by the semi-diameter by vijaya---the first circumference, and the vikonaka from the second: then in order, the first circumference with the result added becomes the true circumference---this well-known method of computing the true circumference also becomes invalid by Āryabhaṭa's system.

Why? Because of the reading of circumference in the terms. In meaning, in that system, there is the taking of circumference in the terms. There, sometimes the first is of even-term type, sometimes of odd-term type. Therefore it varies in correction procedures. This is invalid.

And the vikonaka from the first---because of the correctness of the instantaneous correction---this is to be considered by the wise.
\end{bilingual}


%% ========================================================================
%% CHAPTER: DVIKSHEPA-ADHIKARA
%% ========================================================================
\newpage
\section[\dev{द्विक्षेपाधिकारः}]{\dev{द्विक्षेपाधिकारः} --- Celestial Latitude}

\begin{bilingual}
\textbf{मूलपाठः} --- \skt{इन्द्रोः} (\skt{चन्द्रस्य}) \skt{उत्क्रमज्या} (\skt{भुजोत्क्रमज्या}) \skt{विश्लेषपक्रमगुणा} (\skt{विश्लेषेण शरेण, अपक्रमेण परं कान्तिज्या च गुणा}) \skt{विज्याकृति} (\skt{विज्यावर्गे}) \skt{भक्ता लम्बमानद्विक्षेपकल्पेत् युक्तमार्यभटेन तत् ततः} (\skt{तस्मात् कारणात्}) \skt{असत् यतः} (\skt{यस्मात्कारणात्}) \skt{क्षयनान्ते} (\skt{गोलसंख्ये}) \skt{यद्गणानं फलमुत्पद्यते तत्तस्य मते श्रादावयनसंख्यावेतोत्पद्यते अतस्तदसदिति} ॥१३५॥
\rightcol
\textbf{Text (v.\,135)} --- Of the Moon (Indu), the \textit{utkrama-jyā} (reversed sine, i.e., the \textit{bhuja-utkrama-jyā}), multiplied by the \textit{viśleṣa-apakrama} (by the arrow of separation, by the apakrama, and by the maximum brightness-sine), divided by the square of the \textit{vijyā} (viṣama-jyā), the result is the \textit{dvikṣepa} (double-latitude) in \textit{lamba-māna}. That was stated by Āryabhaṭa as proper (\textit{yuktam}); therefore (\textit{tataḥ}) it is invalid (\textit{asat}). Since (\textit{yataḥ}) at the end of the \textit{kṣaya} (at the \textit{gola-saṅkhyā}), the result produced by that calculation---in his system, the \textit{śrāddhāyana-saṅkhyā} (number of \textit{śrāddhāyana}s) is produced from it---therefore that is invalid.
\end{bilingual}

\medskip

\begin{bilingual}
\textbf{टीका} --- \skt{इन्द्रोः} (\skt{चन्द्रस्य}) \skt{उत्क्रमज्या} (\skt{भुजोत्क्रमज्या}) \skt{विश्लेषपक्रमगुणा} (\skt{विश्लेषेण शरेण, अपक्रमेण परं कान्तिज्या च गुणा}) \skt{विज्याकृति} (\skt{विज्यावर्गे}) \skt{भक्ता लम्बमानद्विक्षेपकल्पेत् युक्तमार्यभटेन तत् ततः} (\skt{तस्मात् कारणात्}) \skt{असत् यतः} (\skt{यस्मात्कारणात्}) \skt{क्षयनान्ते} (\skt{गोलसंख्ये}) \skt{यद्गणानं फलमुत्पद्यते तत्तस्य मते श्रादावयनसंख्यावेतोत्पद्यते अतस्तदसदिति} ॥१३५॥
\rightcol
\textbf{Commentary} --- Of the Moon (\textit{indoḥ} = \textit{candra-sya}), the \textit{utkrama-jyā} (the \textit{bhuja-utkrama-jyā}), multiplied by the \textit{viśleṣa-apakrama} (by the arrow of separation, by the apakrama, and by the maximum brightness-sine), divided by the square of the \textit{vijyā}---the \textit{dvikṣepa} in \textit{lamba-māna} is obtained. Āryabhaṭa stated this as proper; but from that cause it is invalid. Since at the end of the \textit{kṣaya} (at the \textit{gola-saṅkhyā}), the result produced by that calculation---in his view, the \textit{śrāddhāyana-saṅkhyā} is produced from it---therefore that is invalid.
\end{bilingual}

\medskip

\begin{bilingual}
\textbf{उपपत्तिः} --- अत्र चतुर्वेदाचार्यः \skt{सद् द्विषमेतद्यत् आर्यभटीयं वाक्यमेतद् व्रात्यं}---

\skt{विश्लेषपक्रमगुणामुत्क्रमं विस्तराधं क्षतिभक्तम्} ।

\skt{अत्र मटदीपिकायां परमेश्वरः}

\skt{`सायनचन्द्रस्योत्क्रमं कोट्या उत्क्रमज्येतयः' एवं चेत् तदास्यार्यभटं न समीचीनम्} ।

\skt{इदं कर्म क्रमज्या क्षेपद्वयमार्यभटादिमतोत्क्रमज्यया यत् कृतं तत्र तथ्यमिति भास्करखण्डनं वस्तुतो गोलयुक्तिकृत् वास्तवमिति स्फुटं चापक्षेत्रकु-क्षेत्राणाम्} ॥१३५॥
\rightcol
\textbf{Proof} --- Here the Caturveda-ācārya: ``Indeed this is twofold: this Āryabhaṭīya statement is \textit{vrātya} (heterodox):''

\textit{viśleṣa-pakrama-guṇām utkramaṃ vistara-dhaṃ kṣati-bhaktam}

Here, in the \textit{Maṭadīpikā}, Parameśvara: ``If the \textit{sāyana-candra-utkrama} be by the koṭi, by the \textit{utkrama-jyā}---if thus, then this of Āryabhaṭa is not correct.''

This operation---the \textit{krama-jyā}, the two \textit{kṣepa}s, the \textit{utkrama-jyā} according to Āryabhaṭa's system---whatever is done there is real; thus Bhāskara's refutation, actually based on spherical demonstration, is true. And the plane-figures and field-figures are clear.
\end{bilingual}


%% ========================================================================
%% CHAPTER: YUGADHIKARA
%% ========================================================================
\newpage
\section[\dev{युगाधिकारः}]{\dev{युगाधिकारः} --- Cosmic Cycles (Yugas)}

\begin{bilingual}
\textbf{मूलपाठः} --- \skt{हिं. भा. --- जिस कारण से ग्रहगणनाओं तथा पात मन्दोच्च-शीघ्रोच्चों की मेषादि स्थिति में जितने काल में फिर मेषादि में स्थिति होती है वहीं काल महा युग कहलाता है, श्रीपथ्यविद्याचुद्र आदि आचार्यों ने जो युग कहा है वह स्फुट है, क्योंकि उनके मत से महायुगादि मे मेषादि में ग्रहय (भिन्न स्थान में) उत्पन्न होता है, अर्थात् उनके मत से मनु आदि स्मृतिकार रचित स्मृति-ग्रन्थ सम्मत महायुगादि में यदि ग्रहगणना करते हैं तो मेषादि में मज्ज्ञनादि ग्रहों की स्थिति सिद्ध नहीं होती है, श्रीपैथा रचित रोमक में कहीं पर महायुगादि की चर्चा नहीं है, इसलिये महायुग शब्द से यहाँ आचार्य कथित महायुग लेना चाहिये, जब सब ग्रहों को तथा पात मन्दोच्च शीघ्रोच्चों की स्थिति मेषादि में होती है तो उसी का नाम महायुगादि है, श्रीपैथा कथित महायुगादि में यह स्थिति नहीं होती है इसलिये उनके मत से महायुगादि ठीक नहीं है, आचार्यैक यह खण्डन ठीक है} ॥५५॥
\rightcol
\textbf{Hindi Commentary (v.\,55)} --- For the reason by which, in the computation of planets and the positions of pāta, mandocca, and śīghrocca at the beginning of Meṣa (Aries), in whatever time they again come to the position in Meṣa---that time is called Mahā-yuga. The yuga stated by Śrīpathyavidyā-cūḍra and other ācāryas is correct, because according to their view, at the beginning of Mahā-yuga, the planets are produced at the beginning of Meṣa (in a different place). That is, according to their view, if planetary computation is done in the Mahā-yuga approved by the smṛti-texts composed by smṛti-kāras like Manu, then the position of the planets in Meṣa etc. is not established. In the Romaka composed by Śrīpaittha, nowhere is there discussion of Mahā-yuga-beginning. Therefore, by the word Mahā-yuga, here the Mahā-yuga stated by the Ācārya should be taken. When all the planets and the positions of pāta, mandocca, and śīghrocca are at the beginning of Meṣa---that alone is named Mahā-yuga-ādi. In the Mahā-yuga-ādi stated by Śrīpaittha, this position does not occur; therefore according to their view, the Mahā-yuga-ādi is not correct. The Ācārya's refutation here is correct.
\end{bilingual}

\medskip
\noindent\rule{\textwidth}{0.4pt}

\begin{bilingual}
\textbf{मूलपाठः} --- \skt{इदानीं पुनरपि तच्छुंमेव नि रात्रं रोति}।

\skt{कदिनादौ स्मृतिबुक्तं ग्रहमोक्षपतिर्दिनक्षये प्रलयः} ।

\skt{तात्यतिबहूनि यस्मात् महायुगेष्टोऽप्रसिद्धिमवम्} ॥५६॥
\rightcol
\textbf{Text (v.\,56)} --- Now again, the very same is refuted at night:

\textit{kadinādau smṛti-buktaṃ grahamokṣapatir dina-kṣaye pralayaḥ}

At the beginning of the \textit{kadina} (many days), as stated in the Smṛti, the liberation of planets, at the destruction of the day, the dissolution.

\textit{tātyatibahūni yasmāt mahāyugeṣṭo 'prasiddhim avam} --- Since there are very many of those, the \textit{eṣṭa} (desired) of Mahā-yuga is of unestablished, low (nature).
\end{bilingual}

\medskip

\begin{bilingual}
\textbf{टीका} --- \skt{कदिनादौ बहुदिनादौ ग्रहनक्षत्रोत्पतिर्दिनक्षये बहुदिनवासाने च ग्रहमानाः प्रलय इति स्मृतिग्रन्थमस्ति} ।

\skt{एवं बहुदिननम्मे महायुगे यस्मात् तार्गिश्रीषेणाच्छत्कानि युगानि षष्टिबहूनि भवन्ति अत इदं तदुक्तमप्रसिद्धं न कुतश्चपि सत्यादि तच्छुं च} ।

\skt{षष्टोऽनेक युगग्रहणेन गोरवकर्मणा किमिति} ।

\skt{बाबलमेतत्} ।

\skt{यतोऽनेकयुगग्रहणेनापि ग्रहगणा नायां न काञ्चिदानिसत्स्थानग्रहणादिति स्फुटं गणितं विदामिति} ॥५६॥
\rightcol
\textbf{Commentary} --- At the beginning of \textit{kadina} (many days), the creation of planets and stars; at the destruction of the day, at the end of many days, the dissolution of the planetary measures---thus it is in the Smṛti-texts. Thus in the Mahā-yuga of many days, since there are sixty-times-hundred yugas, therefore this stated by him is unestablished, not true from any source. What is the use of the gravity-ritual by the taking of many yugas? This is invalid. For even by taking many yugas, the planetary computations are not in this, not at any time obtained from the position---thus the computation is clear to those who know.
\end{bilingual}

\medskip
\noindent\rule{\textwidth}{0.4pt}

\begin{bilingual}
\textbf{मूलपाठः} --- \skt{वि. भा. --- कदिनादौ} (\skt{बहुदिनादौ}) \skt{ग्रहमोक्षपतिः} (\skt{ग्रहनक्षत्रसृष्टिः}), \skt{दिनक्षये} (\skt{बहुदिनान्ते}) \skt{ग्रहाणां नक्षत्राणां च प्रलयः} (\skt{नाशः}) \skt{स्मृतिषु} (\skt{स्मृतिग्रन्थेषु}) \skt{कथितमस्ति, सिद्धान्तशेखरे श्रीपतिनास्यैकमुत् यथा---ज्योतिर्ग्रहाणां विधिवासरावो सृष्टिलयस्तद्दिवसावसाने} ।

\skt{सिद्धान्तशिरोमणौ भास्कराचार्येण `यतः सृष्टिरेषां दिनादौ दिनान्ते लयस्तेषु सत्सैव तच्छार चिन्ता' स्प्यनैन बहुदिनादावेव ग्रहदिनसृष्टिः कथ्यते, परं सूर्यसिद्धान्तकारेण कथ्यते यत् `बहुदिनादितः शातयुगात्वेदसप्तवैदिस्यामेषु ब्रह्मा सृष्टिं रचयित्वावसासै नियोजितवान्' । ब्रह्मपुत्र-श्रीपति-भास्कराचार्यैः कथितसुख्यादिकालानिराकरणार्थं सूर्यसिद्धान्तकारमतमण्डनार्थं च कमलाकरेण सिद्धान्ततत्त्वविवेके बहुप्रपञ्चं नहि नाममेदैन वस्तुमेदः} ॥५५-५६॥
\rightcol
\textbf{Expanded Commentary} --- At the beginning of \textit{kadina} (\textit{bahu-dinādau}), the \textit{mokṣa-patiḥ} (creation of planets and stars); at the \textit{dina-kṣaye} (end of many days), the \textit{pralayaḥ} (destruction) of the planets and stars---this is stated in the \textit{smṛti} texts. In the Siddhāntaśekhara, Śrīpati has stated: ``The creation and dissolution of the luminous planets occurs at the beginning and end of those days.''

In the Siddhānta-śiromaṇi, Bhāskarācārya: ``Since their creation is at day-beginning, and dissolution at day-end, the intelligent should contemplate this.'' By this, the creation of planets at day-beginning is stated. But the Sūryasiddhānta-kāra states that ``from the beginning of many days, from the Śāta-yuga, in the seven Vedas, Brahmā having created, was appointed at the end.'' For the refutation of the time of happiness etc. stated by Brahmagupta-Śrīpati-Bhāskarācārya, and for the refutation of the Sūryasiddhānta-kāra's view, in the Siddhānta-tattva-viveka, there is much elaboration by Kamalākara---there is no substance without name.
\end{bilingual}


%% ========================================================================
%% CHAPTER: GANITADHYAYA (Mathematics)
%% ========================================================================
\newpage
\section[\dev{गणिताध्यायः}]{\dev{गणिताध्यायः} --- Chapter on Mathematics}

\begin{bilingual}
\textbf{संपादकीय टिप्पणी} --- एतस्मिन् अध्याये ब्रह्मगुप्तस्य बीजगणितं, पाटीगणितं, ज्यामितिः, तथा श्रेणिसंव leggनं प्रदर्शितम्। अयं द्व्याशीतिकोण्यध्यायः लोकमतhematics इतिहासे अत्यन्तं महत्त्वपूर्णः।
\rightcol
\textbf{Note} --- In this chapter Brahmagupta's algebra (\textit{bīja-gaṇita}), arithmetic (\textit{pāṭī-gaṇita}), geometry, and series summation are presented. This \textit{Dvyāśīti-koṇy-adhyāya} (12th chapter) is of supreme importance in the history of world mathematics.
\end{bilingual}

\medskip
\noindent\rule{\textwidth}{0.4pt}

\subsection*{Algebra: Multiplication and Powers --- Verse 1}

\begin{bilingual}
\textbf{मूलपाठः} --- \skt{दानं च दत्तमर्धं पञ्च तेन मान-स्यार्धेषपञ्चनवमांशं तदा किम्} ॥१॥
\rightcol
\textbf{Text (v.\,1)} --- \textit{dānaṃ ca dattam ardhaṃ pañca tena māna-syārdheṣa-pañcanavamāṃśaṃ tadā kim}

A donation was given: half, five; of that measure, half less one-fifth and one-ninth---then what?
\end{bilingual}

\medskip

\begin{bilingual}
\textbf{टीका} --- \skt{न्यासः} $\dfrac{1}{2}$ । \skt{उत्तरं कथयेत् जातं} $1\,\dfrac{5}{8}\,\dfrac{3}{4}$ ।

\skt{अत्र कोल्ब्रहं कसाहिवेन भ्रमात्} $\dfrac{3}{4}\,\dfrac{5}{8}$ \skt{इत्युत्तरं वास्तवं लिखितम्}।

(\textit{See his translation P. 283.})

अत्र चतुर्वेदाचार्यः \skt{स्वदीकार्याम्}।

\skt{एवमिहाचार्यं पञ्चजात्य एकोक्ता} ।

\skt{यतः पठतीति तदातिसंकेमातो गतार्थां कृत्वा नोक्ता} ।

\skt{स्कन्दसेनादिभिस्तस्य नाम कृतं भागमातेति}।
\rightcol
\textbf{Commentary} --- \textit{Nyāsa} (statement): $\frac{1}{2}$. The answer is stated as having become $1\,\frac{5}{8}\,\frac{3}{4}$.

Here due to confusion, the actual written answer is $\frac{3}{4}\,\frac{5}{8}$. (See his translation p.\,283.)

Here the Caturveda-ācārya: ``This is to be acknowledged.''

``Thus here the Ācārya, one of the \textit{pañca-jātya}, spoke.''

``Because he reads, therefore, having made the meaning gone, he did not speak.''

``By Skandasena and others, his name was made \textit{Bhāga-māta}.''
\end{bilingual}

\medskip
\noindent\rule{\textwidth}{0.4pt}

\subsection*{Geometric Progression --- Verse 5}

\begin{bilingual}
\textbf{मूलपाठः} --- \skt{सहस्रच्छेदांशयुक्तिः} (\skt{तुल्यहरांशानां योगः}) \skt{छेदविभक्ता} (\skt{हरभक्ता}) \skt{फलं प्रथमजाति सर्वर्णं भवति} ।

\skt{द्वितीयजाती-मरेङ्गा गुणिताः, छेदः} (\skt{हरः}) \skt{छेदः} (\skt{हरः}) \skt{गुणितात्तदंश्या हरमका फलं सर्वर्णां भवति} ।

\skt{प्रथमजाति मवर्णां तु `योगोन्तर तुल्यहरांशकानामित्यादि भास्करोक्तमनुक्रमेव तथा द्वितीय प्रभागजातिसवर्णं यतदनुक्रमेव लीलावत्यां `लवालवधनाभ्यं हरा हरनाः' इत्यादि भास्करोक्तम्} ॥५॥
\rightcol
\textbf{Text (v.\,5)} --- The sum of fractions with equal denominators (\textit{sahasra-cchedāṃśa-yuktiḥ = tulya-harāṃśānāṃ yogaḥ}), divided by the denominator (\textit{cheda-vibhaktā}), becomes the result of the first class for all quantities.

Those multiplied by the second-class measures; denominator (\textit{hara}) multiplied by denominator---the fraction thereof, with the denominator, becomes the result for all quantities.

For the quantities of the first class, ``\textit{yogontara tulya-harāṃśakānām}'' etc.---as stated by Bhāskara; and for the second \textit{prabhāga-jāti} combination, as in the \textit{Līlāvatī}: ``\textit{lavāla-vadhana-abhyaṃ harā haranāḥ}'' etc.---thus stated by Bhāskara.
\end{bilingual}

\medskip
\noindent\rule{\textwidth}{0.4pt}

\subsection*{Division and Geometric Problems --- Verse 6}

\begin{bilingual}
\textbf{मूलपाठः} --- \skt{इदानीं भागाहारे करणसूत्रमाह}।

\skt{परिवर्त्यं भागहरच्छेदांशां छेदसंगुणाच्छेदः} ।

\skt{शेषोऽंशगुणो भाजस्य भागाहारः सर्वेषां तयोः} ॥६॥
\rightcol
\textbf{Text (v.\,6)} --- Now he states the karana-sūtra for division (\textit{bhāgāhāra}):

\textit{parivartyaṃ bhāgahara-cchedāṃśāṃ cheda-saṃguṇāc chedaḥ}

Having interchanged, the denominator-fraction of the divisor (\textit{bhāgahara}) multiplied by the denominator;

\textit{śeṣo 'ṃśaguṇo bhājasya bhāgāhāraḥ sarveṣāṃ tayoḥ} --- The remainder multiplied by the numerator of the dividend; the division of both of all.
\end{bilingual}

\medskip

\begin{bilingual}
\textbf{टीका} --- \skt{भागहारच्छेदांशी भाजकस्य छेदांशो परिवर्त्यं भाज्यस्य छेदः परिवर्तितच्छेदस् हराछेदो भवति} ।

\skt{एवं भाज्यस्यांशः परिवर्तितांशगुणांशो भवति} ।

\skt{एवं सर्वेषां तयोर्योगमिश्रयोर्भागाहारो भवति} ।

\skt{सर्वेषां तु `रूपारिच्छेदगुणानि' इति विधिना} ।

\skt{भास्करभागाहारोत्यदनुक्रम एव} ।

\skt{उद्देशकचतुर्वेदः ---}
\rightcol
\textbf{Commentary} --- The denominator-fraction of the divisor (\textit{bhāgahara-chedāṃśī})---the denominator-fraction of the divisor is interchanged with the denominator of the dividend; the interchanged-denominator becomes the divisor-denominator.

Thus the numerator of the dividend becomes the numerator-multiplied-by-interchanged-numerator.

Thus for both, for mixed sums, the division is produced.

For all, by the rule ``\textit{rūpāricchedaguṇāni}'' etc.

Bhāskara's division method is the same.

The example of Caturveda:
\end{bilingual}

\medskip

\begin{bilingual}
\textbf{उद्देशकः} --- \skt{यत्रायते फलं हष्टं सार्धेनैकयमेनदवः} ।

\skt{सार्धं दशभुजस्चैव कोटिस्तनार्मकीयताम्} ॥
\rightcol
\textbf{Example} --- Where the area is six \textit{haṣṭa} (square cubits), with one and a half, with one \textit{yama} and a half. Ten sides and a half, and the koṭi---let it be of that measure.

This is a geometric problem: given area $324\,\frac{3}{8}$, side $90\,\frac{1}{2}$, find the koṭi.
\end{bilingual}

\medskip

\begin{bilingual}
\textbf{टीका} --- \skt{न्यासः । क्षेत्रफलम्} $324\,\dfrac{3}{8}$ \skt{भुजः} $90\,\dfrac{1}{2}$ \skt{उत्क्रमजाता कोटिः} $11\,\dfrac{3}{5}$ ।

अत्र पुनश्चतुर्वेदाचार्यः ।

\skt{अन्ये पुनरिहापि वं राशिकर्मकमुदाहरणं ददति यथा---}
\rightcol
\textbf{Commentary} --- \textit{Nyāsa}: Area $= 324\,\frac{3}{8}$; bhuja $= 90\,\frac{1}{2}$; the \textit{utkrama-jāta} koṭi $= 11\,\frac{3}{5}$.

Here again the Caturveda-ācārya.

``Others here also give an example of proportion-calculation, as:''
\end{bilingual}

\medskip
\noindent\rule{\textwidth}{0.4pt}

\subsection*{Series Summation --- Verse 9}

\begin{bilingual}
\textbf{मूलपाठः} --- \skt{अथवा लघुतमापवर्तेन गणितम्}

\skt{न्यासः} $\dfrac{1}{2} + \dfrac{1}{3} + \dfrac{3}{4} + \dfrac{1}{5}$

\skt{अथवैदृशाः २, ६, १२, ४ मेघा २ मननापवर्तनेन}

\skt{ततोऽपवर्तनाद्वां जातः} $2 \times 2 \times 3 = 12$

\skt{सैष} $9\,9\,9$ \skt{गुणिताः} $12$, \skt{प्रमेने-वा हरांशयोगेनैनं}

$\dfrac{1}{2} \times \dfrac{6}{1} + \dfrac{1}{3} \times \dfrac{12}{2} + \dfrac{3}{4} \times \dfrac{12}{3} + \dfrac{1}{5} \times \dfrac{12}{4}$

$= \dfrac{6}{2} + \dfrac{12}{6} + \dfrac{12}{12} + \dfrac{12}{20} = \dfrac{6}{2} = 9 = $ \skt{योगफलम्} ॥९॥
\rightcol
\textbf{Text (v.\,9)} --- Alternatively, calculation by the \textit{laghutama-apavartana} (smallest reduction):

Statement: $\frac{1}{2} + \frac{1}{3} + \frac{3}{4} + \frac{1}{5}$.

The denominators 2, 6, 12, 4: by reduction with 2, 3:

From this reduction, produced: $2 \times 2 \times 3 = 12$.

This multiplied: 12, or by the sum of denominator-numerator,

$\frac{1}{2} \times \frac{6}{1} + \frac{1}{3} \times \frac{12}{2} + \frac{3}{4} \times \frac{12}{3} + \frac{1}{5} \times \frac{12}{4}$

$= \frac{6}{2} + \frac{12}{6} + \frac{12}{12} + \frac{12}{20} = 9 = $ sum result.
\end{bilingual}

\medskip
\noindent\rule{\textwidth}{0.4pt}

\subsection*{Binomial Expansion --- Geometric Demonstration}

\begin{bilingual}
\textbf{संपादकीय टिप्पणी} --- \skt{स्फुटादिकारके ७३४ पृष्ठे अस्य खण्डस्य `स्थाप्योन्त्यघन' इति सूत्रस्य उपपत्तिः ज्यामितिकक्षेत्रद्वारा प्रदर्शिता। एतत् $(a+b)^3$ विस्तारस्य ज्यामितिकं प्रदर्शनं पाश्चात्यगणितेभ्यः शताब्दीभ्यः पूर्वं भारते संसिद्धमासीत्।}
\rightcol
\textbf{Note} --- At page 734 of this chunk, the proof of the rule ``\textit{sthāpyo-ntyaghana}'' is demonstrated by geometric diagrams. This geometric demonstration of the expansion of $(a+b)^3$ was accomplished in India centuries before similar developments in Western mathematics.
\end{bilingual}

\medskip

\begin{bilingual}
\textbf{मूलपाठः} --- \skt{वर्गीभुजाः} $= a + b$

\skt{वैश्यः} $= a + b$

\skt{वर्गक्षेत्रफलम्} $= (a+b) \times (a+b) = a^2 + 2ab + b^2$

\skt{क्षेत्रफलं वैश्यगुणमित्यादिना वाप्या घनतमकं क्षेत्रं सम्बन्धेन याग्युदाहरणानि प्रदर्शितानि तानि न समीचीनानीति सुधीभिरव्यानीति} ॥६॥
\rightcol
\textbf{Geometric Demonstration} --- Side of the square $= a + b$.

Breadth $= a + b$.

Area of the square $= (a+b) \times (a+b) = a^2 + 2ab + b^2$.

The area of the field multiplied by the length etc.---the \textit{ghanatamaka} (solid-cube) field is related; the examples shown are not correct, and are not to be examined by the wise.
\end{bilingual}

\medskip

\begin{bilingual}
\textbf{उपपत्तिः} --- \skt{अ + क अस्य घन करणार्थं समत्रिघाताघं घन इति घनपरिभाषया}

$(\text{अ} + \text{क})(\text{अ} + \text{क})(\text{अ} + \text{क}) = (\text{अ} + \text{क})^3$

\skt{गुणनेन} $(\text{अ}^2 + \text{अ. क} + \text{क}^2)(\text{अ} + \text{क}) = (\text{अ}^2 + 2\,\text{अ.क} + \text{क}^2)(\text{अ} + \text{क})$

$= \text{अ}^3 + 2\,\text{अ}^2.\text{क} + \text{अ.क}^2 + \text{अ}^2.\text{क} + 2\,\text{अ.क}^2 + \text{क}^3$

$= \text{अ}^3 + 3\,\text{अ}^2.\text{क} + 3\text{अ.क}^2 + \text{क}^3 = (\text{अ} + \text{क})^3$

\skt{एतत् `स्थाप्योन्त्यघन' इत्याचार्यैकमुपपद्यते} ।

\skt{`स्थाप्योन्त्यवर्गं द्विगुणान्त्यनिघ्नं' इत्यादि लीलावत्यां भास्करोक्तमतदनुक्रमेवास्ति} ।

\skt{प्राचीनैः केवलं वर्गघनयोगविचारः कृतः} ।

\skt{द्वयोर् द्वयोर् योगस्य घनकरणार्थं प्रथमाक्षेपित्वं संज्ञकः} ।

\skt{द्वितीयो-द्वय उत्तरसंज्ञकः, तदा अन्त्यघनः स्थाप्यः, ततोऽन्त्यकृतिः} (\skt{अन्त्यवर्गः}) \skt{द्विगुणोत्तर-संज्ञकः, द्वयोर् द्वयोर् योगस्य घनो भवतीति} ।

\skt{अत्र चतुर्वेदाचार्यैदं शकः---`चतुरस्रा समा वापी हस्ततिवन समिता । वैघनं च तथा तस्याः फलं घनतमकं' घनतमकम्} ॥
\rightcol
\textbf{Proof} --- For making the cube of $(a + k)$, by the definition of cube as equal triple-product:

$(a + k)(a + k)(a + k) = (a + k)^3$

By multiplication: $(a^2 + a \cdot k + k^2)(a + k) = (a^2 + 2\,a \cdot k + k^2)(a + k)$

$= a^3 + 2\,a^2 \cdot k + a \cdot k^2 + a^2 \cdot k + 2\,a \cdot k^2 + k^3$

$= a^3 + 3\,a^2 \cdot k + 3\,a \cdot k^2 + k^3 = (a + k)^3$

This agrees with the Ācārya's ``\textit{sthāpyo-ntyaghana}'' (place-the-last-cube).

``\textit{sthāpyo-ntyavargaṃ dviguṇā-ntyanighnaṃ}'' etc.---as stated by Bhāskara in the Līlāvatī, this follows the same method.

By the ancients, only the investigation of the sum of squares and cubes was done.

For making the cube of the sum of two-plus-two, the first is called \textit{ākṣepitva}; the second two is called \textit{uttara}; then the last-ghana is placed; from that the last-square (\textit{antya-kṛti}) is doubled, called \textit{uttara}; the cube of the sum of two-plus-two becomes.

Here the Caturveda-ācārya gave this example: ``A square pond, equal to the extent of hastas. Its depth also; the result is the \textit{ghanatamaka} (solid cube).''
\end{bilingual}

\medskip

\begin{bilingual}
\textbf{मूलपाठः} --- \skt{यथा श्रेष्ठा रूपम् = अयमेव द्विगुण पदसिद्धान्तः} । $(\text{अ} + \text{क}) = \text{अ} +$

\begin{tabular}{c|c|c|c}
$\dfrac{\text{अ} \times \text{न} \times \text{क}}{1}$ & $\dfrac{\text{अ} \times \text{न (न}-1)\text{क}}{1 \times 2}$ & $\dfrac{\text{अ} \times \text{न (न}-1)\text{ (न}-2)\text{ क}}{1 \times 2 \times 3}$ & $+ \cdots\cdots$ \\
$\text{न} = 1$ & $\text{न} = 2$ & $2$ & $\text{न} = 3$ \\  
\end{tabular}

\skt{अत्र यदि न} $= 2$ \skt{तदा} $(\text{अ} + \text{क})^2 = \text{अ}^2 + \text{अ} \times 2 \times \text{क} + \text{क}^2$

\skt{यदि न} $= 3$ \skt{तदा} $(\text{अ} + \text{क})^3 = \text{अ}^3 + \text{अ}^2 \times 3 \times \text{क} + \text{अ} \times 3 \times \text{क}^2 + \text{क}^3$
\rightcol
\textbf{Binomial Table} --- As the best \textit{rūpa} (form): this itself is the twofold \textit{pada-siddhānta}:

\[
\begin{array}{c|c|c|c}
\frac{a \times n \times k}{1} & \frac{a \times n(n-1)k}{1 \times 2} & \frac{a \times n(n-1)(n-2)k}{1 \times 2 \times 3} & + \cdots \\ \hline
n = 1 & n = 2 & 2 & n = 3 \\
\end{array}
\]

Here if $n = 2$: $(a + k)^2 = a^2 + a \times 2 \times k + k^2$

If $n = 3$: $(a + k)^3 = a^3 + a^2 \times 3 \times k + a \times 3 \times k^2 + k^3$
\end{bilingual}

\medskip

\begin{bilingual}
\textbf{उपपत्तिः} --- \skt{एतावता `स्थाप्योन्त्यवर्गं द्विगुणान्त्यनिघ्नं' इत्यादि वर्गोपपतिः `स्थाप्यो न्त्यघन' इत्यादि घनोपपत्तिश्च स्पष्टा भवति}।
\rightcol
\textbf{Proof} --- By this much, ``\textit{sthāpyo-ntyavargaṃ dviguṇā-ntyanighnaṃ}'' etc.---the proof of the square, and ``\textit{sthāpyo ntyaghana}'' etc.---the proof of the cube become clear.
\end{bilingual}

\medskip
\noindent\rule{\textwidth}{0.4pt}

\subsection*{Square Root Extraction --- Verse 10}

\begin{bilingual}
\textbf{मूलपाठः} --- \skt{छेदं विभजेद् द्वयोर्लब्धं तयोः पृथक् पृथक् दलयेत्} ।

\skt{वर्गान्मूलं दलितात् पदाद् वर्गमूलं स्यात्} ॥१०॥
\rightcol
\textbf{Text (v.\,10)} --- One should divide the denominator; having obtained (the result) of the two, one should divide them separately, separately.

From the square, the root from the divided term---that becomes the square-root.
\end{bilingual}

\medskip

\begin{bilingual}
\textbf{टीका} --- \skt{वि. भा. --- द्वयोर्द्वयोर्योगस्य घनकरणार्थं प्रथमाक्षेपित्वं संज्ञकः} ।

\skt{द्वितीयो-द्वय उत्तरसंज्ञकः, तदा अन्त्यघनः स्थाप्यः, ततोऽन्त्यकृतिः} (\skt{अन्त्यवर्गः}) \skt{द्विगुणोत्तर-संज्ञकः, द्वयोर्द्वयोर्योगस्य घनो भवतीति} ।

\skt{अत्र चतुर्वेदाचार्यैदं शकः---`चतुरस्रा समा वापी हस्ततिवन समिता । वैघनं च तथा तस्याः फलं घनतमकं' घनतमकम्} ॥१०॥
\rightcol
\textbf{Commentary} --- For making the cube of the sum of two-plus-two, the first is called \textit{ākṣepitva}; the second two is called \textit{uttara}; then the last-ghana is placed; from that the last-square (\textit{antya-kṛti}) is doubled, called \textit{uttara}; the cube of the sum of two-plus-two becomes.

Here the Caturveda-ācārya gave this example: ``A square pond, equal to the extent of hastas. Its depth also; the result is the \textit{ghanatamaka} (solid cube).'' \textit{Ghanatamakam}.
\end{bilingual}


%% ========================================================================
%% COLOPHON AND END MATTER
%% ========================================================================
\newpage
\section*{Colophon}

\begin{bilingual}
\textbf{मूलपाठः} ---

\skt{इति श्रीब्रह्मस्फुटसिद्धान्ते द्व्याशीतिकोण्यध्यायो नामैकादशोऽध्यायः ॥ ११ ॥}

\skt{इति श्रीब्रह्मस्फुटसिद्धान्ते तन्त्रपरीक्षाध्यायो नामैकादशोऽध्यायः ॥ ११ ॥}

\skt{इति श्रीब्रह्मस्फुटसिद्धान्ते दुष्प्रदाध्यायो नामैकादशोऽध्यायः ॥ ११ ॥}
\rightcol
\textbf{Colophon} ---

Thus, in the \textit{Śrī-Brāhmasphuṭasiddhānta}, the \textit{Dvyāśīti-koṇy-adhyāya} (12th chapter), named the eleventh chapter. || 11 ||

Thus, in the \textit{Śrī-Brāhmasphuṭasiddhānta}, the \textit{Tantraparīkṣādhyāya} (chapter on examination of systems), named the eleventh chapter. || 11 ||

Thus, in the \textit{Śrī-Brāhmasphuṭasiddhānta}, the \textit{Duṣpradādhyāya} (chapter on difficult gifts), named the eleventh chapter. || 11 ||
\end{bilingual}

\medskip
\noindent\rule{\textwidth}{0.4pt}

\section*{Notes on This Chunk}

\begin{bilingual}
\textbf{संपादकीयम्} --- एतत् खण्डं मूलमुद्रितसंस्करणस्य ६२४--७४३ पृष्ठानि (पीडीएफ ५६१--७००) समाविशति। अत्राष्टौ प्रधानाधिकाराः समाविष्टाः---ग्रहच्छायाधिकारः, भग्रहत्यधिकारः, युक्तिलक्षणम्, तन्त्रपरीक्षाध्यायः, कर्णाधिकारः, द्विक्षेपाधिकारः, युगाधिकारः, तथा गणिताध्यायः।
\rightcol
\textbf{Notes} --- This chunk encompasses pages 624--743 of the original printed edition (PDF pages 561--700). Eight major chapters are included: Planetary Shadows, Planetary Conjunctions, Characteristics of Proof, Examination of Astronomical Systems, Hypotenuse, Celestial Latitude, Cosmic Cycles, and the Chapter on Mathematics.
\end{bilingual}

\medskip

\begin{bilingual}
\textbf{गणितीयविषयाः} --- गणिताध्याये भिन्नराशिकर्म, ज्यामितिकश्रेणी, भागहारकलनम्, श्रेणिसंवलनम्, द्विपदविस्तारः $(a+b)^n$, तथा वर्गमूलोत्पत्तिः इति प्रदर्शितानि। $(a+b)^3$ ज्यामितिकप्रदर्शनं विश्वगणितेऽत्यन्तं प्रसिद्धम्।
\rightcol
\textbf{Mathematical Content} --- In the \textit{Gaṇitādhyāya}: operations with fractions, geometric progression, division algorithms, series summation, binomial expansion $(a+b)^n$, and square-root extraction are demonstrated. The geometric demonstration of $(a+b)^3$ is of world-historical importance in mathematics.
\end{bilingual}

\medskip

\begin{bilingual}
\textbf{आर्यभटखण्डनम्} --- तन्त्रपरीक्षाध्याये ब्रह्मगुप्तेन आर्यभटस्य (४७६ ई.) ग्रहसिद्धान्तस्य कठोरं परीक्षणं कृतम्। मुख्यानि खण्डनानि---युगपादसङ्ख्याया अनुपपत्तिः, कल्पादिवारगणनाया अशुद्धिः, तदीयग्रहसिद्धान्तस्य असम्भवः च।
\rightcol
\textbf{Critique of Āryabhaṭa} --- In the \textit{Tantraparīkṣādhyāya}, Brahmagupta subjects Āryabhaṭa's (476 CE) planetary theory to rigorous testing. Main refutations: inconsistency of yuga-pāda numbers, error in kalpa-beginning weekday computation, and the impossibility of his planetary theory.
\end{bilingual}

\vspace{1cm}
\begin{center}
\rule{0.4\textwidth}{0.4pt}

\vspace{0.5cm}

{\small\textit{End of Bilingual Edition --- Chunk 5}}

{\small\dev{इति द्विभाषासंस्करणस्य समाप्तिः --- खण्डं ५}}
\end{center}

\end{document}



% ============================================================
% Source: translations/en/indian/brahmagupta-brahmasphutasiddhanta-pt2-chunk6_bilingual.tex
% ============================================================

\documentclass[10pt,a4paper]{article}
\usepackage{fontspec}
\usepackage{polyglossia}
\usepackage{amsmath,amssymb}
\usepackage{geometry}
\usepackage{xcolor}
\usepackage{titlesec}
\usepackage{fancyhdr}
\usepackage{setspace}
\usepackage{hyperref}

\setmainlanguage{english}
\setotherlanguage{sanskrit}

\newfontfamily\devanagarifont[Script=Devanagari,Path=/usr/share/fonts/truetype/noto/,Extension=.ttf,UprightFont=NotoSerifDevanagari-Regular,BoldFont=NotoSerifDevanagari-Regular]{NotoSerifDevanagari-Regular}

\geometry{paperwidth=210mm,paperheight=297mm,left=15mm,right=15mm,top=22mm,bottom=22mm,headheight=14pt,footskip=24pt}

\definecolor{titlecolor}{RGB}{45,55,72}
\definecolor{sectioncolor}{RGB}{55,65,81}
\definecolor{rulecolor}{RGB}{180,160,140}
\definecolor{versebg}{RGB}{250,248,245}

\newcommand{\dev}[1]{{\devanagarifont #1}}
\newcommand{\longline}{\noindent\rule{\textwidth}{0.4pt}}
\newcommand{\vnum}[1]{\hfill\textnormal{||#1||}}

\newsavebox{\shlokabox}
\newenvironment{shloka}{%
  \begin{lrbox}{\shlokabox}
  \begin{minipage}{\dimexpr\textwidth-16pt\relax}%
  \devanagarifont\centering\setstretch{1.3}%
}{%
  \end{minipage}%
  \end{lrbox}%
  \par\medskip\noindent\fcolorbox{rulecolor}{versebg}{\usebox{\shlokabox}}\par\medskip
}

\pagestyle{fancy}
\fancyhf{}
\fancyhead[L]{\small\textit{Brāhmasphuṭasiddhānta — Part II, Final Bilingual Edition}}
\fancyhead[R]{\small\thepage}
\renewcommand{\headrulewidth}{0.4pt}

\titleformat{\section}{\normalfont\Large\bfseries\color{sectioncolor}}{\thesection}{1em}{}
\titleformat{\subsection}{\normalfont\large\bfseries\color{sectioncolor}}{\thesubsection}{1em}{}

\hypersetup{colorlinks=true,linkcolor=sectioncolor,citecolor=sectioncolor,urlcolor=sectioncolor,pdfencoding=auto,pdftitle={Brāhmasphutasiddhanta Part 2 Bilingual},pdfauthor={Brahmagupta}}

\newcommand{\bilingual}[2]{%
  \noindent%
  \begin{minipage}[t]{0.47\textwidth}%
    \setstretch{1.15}#1%
  \end{minipage}%
  \hfill%
  \vrule\hfill%
  \begin{minipage}[t]{0.47\textwidth}%
    \setstretch{1.15}#2%
  \end{minipage}%
  \par\medskip
}

\begin{document}

\begin{titlepage}
\centering
\vspace*{1.5cm}
{\Huge\bfseries\color{titlecolor} Brāhmasphuṭasiddhānta}\\[0.6cm]
{\LARGE\bfseries\color{sectioncolor} \dev{ब्राह्मस्फुटसिद्धान्त}}\\[1.2cm]
\longline\\[0.6cm]
{\Large\textit{Bilingual Sanskrit--English Edition}}\\[0.3cm]
{\large Part II --- Final Portion (Pages 1539--1676)}\\[0.8cm]
\longline\\[1.2cm]
{\Large\textbf{Brahmagupta}}\\[0.3cm]
{\large (c.\ 598--668 CE)}\\[1cm]
{\large Left Column: Sanskrit (Devanagarī)}\\[0.2cm]
{\large Right Column: English Translation with IAST}\\[1.5cm]
{\normalsize Contents: \dev{मिश्र गणित} | Rule of Three | \dev{भाण्डप्रतिभाण्डक} | \dev{शून्यपरिकर्म} | \dev{श्रेढी} | \dev{क्षेत्रव्यवहार} | \dev{खातव्यवहार} | \dev{धान्यराशि} | \dev{गीताध्याय}}\\[0.5cm]
\longline\\[1cm]
{\small Typeset in XeLaTeX with Devanagari support}
\vfill
{\small\textit{Based on the Chowkhamba Sanskrit Series edition}}
\end{titlepage}

\tableofcontents
\newpage

%% ==================== INTRODUCTION ====================
\section*{Introduction}
\addcontentsline{toc}{section}{Introduction}

\bilingual{%
\textbf{\large सूचना}\\[0.5em]
एतद् अन्तिमभागः द्वितीयस्य भागस्य ब्राह्मस्पुटसिद्धान्तस्य। एषः सम्पुटः गणिताध्यायस्य अन्तिमाध्यायान् समाविशति---तृतीयजातिम्, मिश्रगणितम्, त्रैराशिकम्, भाण्डप्रतिभाण्डकम्, शून्यपरिकर्म, श्रेढीम्, क्षेत्रव्यवहारम्, खातव्यवहारम्, धान्यराशिव्यवहारम्, छायाव्यवहारम्, तथा गीताध्यायस्य आरम्भम्॥
}{%
\textbf{\large Introduction}\\[0.5em]
This is the final portion of Part II of the \textit{Brāhmasphuṭasiddhānta}. This volume contains the concluding chapters of the \textit{Gaṇitādhyāya} (Mathematics Chapter)---the Third Class of Partnerships, Mixed Calculations, the Rule of Three, Exchange of Goods, Operations with Zero, Series, Plane Geometry, Excavations, Grain Heaps, Shadows, and the beginning of the Chapter on Spheres.
}

\bilingual{%
सप्तमे शताब्दे लिखितं ब्रह्मगुप्तस्य एतत् ग्रन्थं २४ अध्यायात्मकम् अस्ति। एषः खण्डः गणितस्य विविधक्षेत्राणि समाविशति---वाणिज्यगणितम्, शून्यबीजगणितम्, श्रेढी, ज्यामिती, घनगणितम्, तथा गोलज्यामिती॥
}{%
Written in the 7th century, Brahmagupta's treatise comprises 24 chapters. This section covers an extraordinary range of mathematical topics---commercial arithmetic, the algebra of zero, series, geometry, solid geometry, and spherical geometry.
}

\vspace{1em}
\noindent\longline
\vspace{1em}

%% ==================== CHAPTER 1: Tritiya Jati ====================
\section{\dev{तृतीय जातिः} --- The Third Class of Partnerships}

\bilingual{%
\textbf{अध्यायः १: तृतीय जातिः}\\[0.3em]
पृष्ठ ७५४---तृतीयजातिसम्बन्धी सवर्णनस्य चर्चा आरभ्यते। एषा जातिः संयुक्तकर्मवेगान् संस्पृशति।\\[0.5em]
\begin{shloka}
दिन सम्बन्धिपूरणकालः सर्वेषां योगः, ततोऽनुपातेनैक कालावच्छेदेन\\
मुक्तसंकुल्याभिर्वापीसंपूरणकालः, एतावदाचार्योक्तमुपन्नम् ॥९॥\vnum{९}
\end{shloka}
}{%
\textbf{Chapter 1: The Third Class of Partnerships}\\[0.3em]
Page 754---Here begins the discussion of the third class of partnership problems (\textit{tṛtīya jāti}), which deals with combined rates of work.\\[0.5em]
\textit{Translation:} ``The sum of all [workers'] day-related filling-times is [found]; then by proportion, the time for a single cistern-filling by all [working together] released [at once] is found. Thus has the Ācārya stated what follows.'' ||9||
}

\bilingual{%
\textbf{टीका:}\\[0.3em]
तृतीय जाति में सवर्णन को कहते हैं। ऊर्ध्वस्थित भंश को अंश से गुणने से तृतीय जाति में सवर्णन होता है। हर से हर को गुणा देना और अपने अंश करके युत-ऊन हर से उपस्थित अंश को गुणा देना तब शंशानुबन्ध और शंशापवाह में सवर्णन होता है। इसका सम्बन्ध भागे से है॥
}{%
\textbf{Commentary:}\\[0.3em]
In the third class (\textit{tṛtīya jāti}), \textit{savaraṇa} (reduction to common denominators) is taught. When the numerator standing above is multiplied by the numerator, the third-class reduction is accomplished. The denominator is given to be multiplied by the denominator, and the numerator standing there, being added or subtracted, is multiplied---then the connection of fractions and the removal of fractions is accomplished. This relates to the section on partnership.
}

\subsection{\dev{उदाहरणम्} --- Worked Example}

\bilingual{%
\textbf{उदाहरणम्:}\\[0.3em]
एक निर्भर एक दिन में एक वापी (पोखड़ा) को भरता है, द्वितीय निर्भर उसी वापी को एक दिन के आधा समय में भरता है, तृतीय निर्भर उसी वापी को एक दिन के चतुर्थांश समय में भरता है, चतुर्थ निर्भर उसी वापी को एकदिन के पञ्चमांश समय में भरता है; यदि एक ही समय में सब निर्भरों को खोल दिया जाय तब वे कितने समय में उस वापी को भरेंगे?
}{%
\textbf{Worked Example:}\\[0.3em]
One worker fills a tank (\textit{vāpī}) in one day; a second fills the same tank in half a day; a third fills the same tank in a quarter of a day; a fourth fills the same tank in a fifth of a day. If all workers are released at the same time, in how much time will they fill that tank?
}

\bilingual{%
\textbf{न्यास:} १|२|४|५---`ऊर्ध्वशैछेदगुणा' इस सूत्र के अनुसार १, २, ४, ५ इन सबों का योग करने से १२---सब निर्भरों से इतने पूरण दिन प्रमाण हुए। अब अनुपात करते हैं यदि १२ इसमें एक दिन पाते हैं तो एक पूरण में क्या इस अनुपात से पूरणकालप्रमाण दे?
}{%
\textbf{Statement:} 1|2|4|5---According to the rule ``multiply the upper and lower [terms],'' summing 1, 2, 4, 5 gives 12---so many filling-days are produced by all workers together. Now by proportion: if 12 yields one day, then for one filling, what is obtained by this proportion? Thus the filling-time is determined.
}

\bilingual{%
\textbf{उपपत्तिः:}\\[0.3em]
एक निर्भर एक दिन में एक वापी को भरता है, द्वितीय निर्भर दिनार्ध में भरता है तो एक दिन में कथा इस अनुपात से एक दिन सम्बन्धी पूरण काल---२, यदि तृतीय निर्भर एक दिन के चतुर्थांश समय में उसी वापी को भरता है तो एक दिन में कथा इस अनुपात से एक दिन सम्बन्धी पूरण काल---४, यदि चतुर्थ निर्भर एक दिन के पञ्चमांश में उसी वापी को भरता है तो एक दिन में कथा इससे एक दिन सम्बन्धी उसका पूरण काल---५। सबों का योग = १२ तब अनुपात करते हैं।\\[0.3em]
लीलावती में ``मजेछिदोंगीरथ तैविमिश्रै'' इत्यादि, भास्करोक्त सूत्र आचार्योक्त सूत्र के अनुरूप ही है॥
}{%
\textbf{Demonstration (\textit{upapatti}):}\\[0.3em]
One worker fills a tank in one day; the second fills it in half a day, so by this proportion his one-day filling-rate is 2. If the third fills the same tank in a quarter of a day, his one-day filling-rate is 4. If the fourth fills it in a fifth of a day, his one-day filling-rate is 5. The sum of all = 12. Then by proportion the filling-time is found.\\[0.3em]
In the \textit{Līlāvatī}, ``\textit{maje chidongī ratha tai vimiśrai}'' etc.---Bhāskara's rule is in conformity with the Ācārya's [Brahmagupta's] rule.
}

\vspace{0.5em}
\noindent\longline
\vspace{0.5em}

%% ==================== CHAPTER 2: Mishra Ganita ====================
\section{\dev{मिश्र गणितम्} --- Mixed Calculations}

\bilingual{%
\textbf{अध्यायः २: मिश्रगणितम्}\\[0.3em]
एषः खण्डः मिश्रगणितम् अथवा मिश्रगणनानि संस्पृशति---चक्रवृद्धिसूदम्, मिश्रकाणि, तथा सम्बद्धानि वाणिज्यगणितीयानि समस्यानि।
}{%
\textbf{Chapter 2: Mixed Calculations}\\[0.3em]
This section deals with \textit{miśra gaṇita} or mixed calculations, including compound interest, alligation, and related commercial arithmetic problems.
}

\subsection{\dev{मिश्रधनानयनम्} --- Finding the Compound Amount}

\bilingual{%
\begin{shloka}
अमागकाल प्रमाणधन परकाल हीनं द्वितीय स्थानं भवति\\
मिश्रघाते मिश्रकेष्ठ वर्गं योज्यं मूलं ततोऽन्यार्धं हीनं प्रमाणफलम् ॥१५॥\vnum{१५}
\end{shloka}
}{%
\textit{Translation:} ``The first term multiplied by the standard period [and] divided by the other period becomes the second term. The square of the other principal, added to the product of the mixed [amount], [and] the square root [taken]; then, diminished by half of the other, is the standard result.'' ||15||
}

\bilingual{%
\textbf{टीका:}\\[0.3em]
प्रमाण काल और प्रमाणधन के धन में परकाल से भाग देकर जो फल हो उसको दो स्थानों में स्थापन करना, प्रथम स्थान स्थित फल का नाम भ्रा, द्वितीय स्थान स्थित फल = भ्रन्य आद्य और मिश्रवन के घात में अन्यकेष्ठ का वर्ग जोड़ कर मूल लेना, उसमे भ्रन्यार्ध को घटाने से प्रमाण फल होता है॥१५॥
}{%
\textbf{Commentary:}\\[0.3em]
The product of the standard period and standard principal, divided by the other period---that result is to be established in two places. The result in the first place is named \textit{bhrā}, the result in the second place = \textit{bhranya}. The square of the other principal, added to the product of the mixed [amount], the square root taken; from it, half of the other being subtracted, the standard result is obtained. ||15||
}

\subsection{\dev{उदाहरणम्} --- Worked Example on Compound Interest}

\bilingual{%
\textbf{उदाहरणम्:}\\[0.3em]
पाच सौ द्रम्मों को सूद पर लगाया गया, चार महीनों में जो मृद दुश्ना उसको पुनः दूसरे स्थान में सूद पर लगाया गया, कालान्तर में दस महीनों में जोरी से झटहनर रुपये हो गये तब प्रमाण फन को कहो॥१५॥
}{%
\textbf{Example:}\\[0.3em]
Five hundred drammas were lent at interest. In four months, whatever [amount] was due was lent again at interest in a second place; after a period of ten months, the total amount became such and such rupees. State the standard result. ||15||
}

\bilingual{%
\textbf{उदाहरण के अनुसार न्यास:} प्रमाण काल = ४, प्रमाण धन = ५००, परकाल = १०, मिश्रधन = ७९।\\[0.3em]
अमागकाल $\\times$ प्रमाणधन $= \\frac{4 \\times 500}{10} = 200$ (प्रथम स्थान)\\[0.3em]
$\\sqrt{(200)^2 + (79)^2} = \\sqrt{40000 + 6241} = \\sqrt{46241} \\approx 215$\\[0.3em]
इसमें १०० घटाने से $215 - 100 = १००$ (प्रमाण फल)॥१५॥
}{%
\textbf{Statement according to the example:} Standard period = 4, Standard principal = 500, Other period = 10, Mixed amount = 79.\\[0.3em]
$\\frac{\\text{standard period} \\times \\text{standard principal}}{\\text{other period}} = \\frac{4 \\times 500}{10} = 200$ (first term)\\[0.3em]
$\\sqrt{(200)^2 + (79)^2} = \\sqrt{46241} \\approx 215$\\[0.3em]
Subtracting 100: $215 - 100 = 100$ (standard result). ||15||
}

\vspace{0.5em}
\noindent\longline
\vspace{0.5em}

%% ==================== CHAPTER 3: Trairashika ====================
\section{\dev{त्रैराशिकादिकम्} --- The Rule of Three and Extensions}

\bilingual{%
\begin{shloka}
फलच्छिदामन्योन्यपक्षनयनं विधाय बहुराशिके वधे\\
स्वल्परादिवधभाजिते फलम् ॥१३॥\vnum{१३}
\end{shloka}
}{%
\textit{Translation:} ``Having performed the mutual transposition of the fruit and the denominators in the multiplication of the rule of many terms, the result [is obtained when] divided by the product of the smaller terms etc.'' ||13||
}

\bilingual{%
\textbf{टीका:}\\[0.3em]
फलच्छिदा अन्योन्यपक्षनयन विधि से बहुराशिक के वध में स्वल्परादि वध भाजिते फल---इस सूत्र का अर्थ है कि बहुसंख्यक राशियों के गुणन में जब फल निकालना हो तो फलच्छिदा (फल का छेद) से अन्योन्य पक्षनयन करके स्वल्प (छोटी) राशि आदि से भाग देने से फल प्राप्त होता है।
}{%
\textbf{Commentary:}\\[0.3em]
The meaning of this rule is: when extracting the result from the multiplication of many quantities, by the mutual transposition of the fruit and denominators, and dividing by the product of the smaller quantities etc., the result is obtained. The \textit{phalachchidā} (fruit-denominator) method of mutual transposition is applied in the \textit{bahu-rāśika} (rule of many terms).
}

\subsection{\dev{त्रैराशिकोदाहरणम्} --- Example: Rule of Three}

\bilingual{%
\textbf{उदाहरणम्:}\\[0.3em]
यदि पाँच माह में द्रम्मों की वृद्धि (सूद) १२५ रूपये हो तब सौ रुपये की कितनी वृद्धि (सूद) होगी?\\[0.3em]
\textbf{न्यास:} ५ | १२५ | १००\\[0.3em]
$\\frac{5 \\times 100}{125} = 4$ वृद्धिः
}{%
\textbf{Example:}\\[0.3em]
If in five months the increase (interest) on drammas is 125 rupees, then what is the increase (interest) on 100 rupees?\\[0.3em]
\textbf{Statement:} 5 | 125 | 100\\[0.3em]
$\\frac{5 \\times 100}{125} = 4$ (increase)
}

\subsection{\dev{पञ्चराशिकादि} --- Rule of Five and More}

\bilingual{%
\begin{shloka}
फलच्छिदामन्योन्यपक्षनयनं विधाय\\
सप्तराशिकादावपि ज्ञेयम् ॥१३॥\vnum{१३}
\end{shloka}
}{%
\textit{Translation:} ``Having performed the mutual transposition of the fruit and denominators, this method is to be known for the rule of seven and others as well.'' ||13||
}

\bilingual{%
\textbf{टीका:}\\[0.3em]
सप्तराशिकादियों के लिये भी उदाहरण के अनुसार न्यास करके ``फलच्छिदामन्योन्यपक्षनयन विधि'' से सप्तराशिक आदि का फल ज्ञात होता है। यह विधि पञ्चराशिकादियों में भी ज्ञात करनी चाहिए।
}{%
\textbf{Commentary:}\\[0.3em]
For the rule of seven and others also, having made the statement according to the example, the result of the rule of seven etc.
}

\vspace{0.5em}
\noindent\longline
\vspace{0.5em}

%% ==================== CHAPTER 4: Bhanda-pratibhanda ====================
\section{\dev{भाण्डप्रतिभाण्डकम्} --- Exchange of Goods}

\bilingual{%
\begin{shloka}
प्रागभूस्य व्यत्यासो भाण्डप्रतिभाण्डकेऽन्यदुक्तसमम्\\
परिकर्माण्यष्टानां व्यवहाराणामभिहितानि ॥१३॥\vnum{१३}
\end{shloka}
}{%
\textit{Translation:} ``The transposition in exchange of goods is as previously stated; the operations of the eight procedures have been declared.'' ||13||
}

\bilingual{%
\textbf{टीका:}\\[0.3em]
``प्रथमस्थाने यौ मूल्यराशी तयोर्व्यत्यासः कार्यः'' इति चतुर्वेदाचार्यः। (तथैव भाण्डप्रतिभाण्डके विधि रित्यादि भास्करोक्तमेतदनुरूपमेव)।\\[0.3em]
अर्थात् प्रथम स्थान में जो दो मूल्य राशियाँ हैं उनका व्यत्यास (विनिमय) करना भाण्डप्रतिभाण्डक व्यवहार है।
}{%
\textbf{Commentary:}\\[0.3em]
``In the first place, the two value-quantities, their transposition is to be done''---thus the four-Veda Ācārya [Brahmagupta]. (Even so, the procedure of exchange of goods etc., as stated by Bhāskara, is in conformity with this.)\\[0.3em]
That is, in the first place, the transposition (exchange) of the two value-quantities is the procedure of \textit{ bhāṇḍa-pratibhāṇḍaka} (exchange of goods).
}

\bilingual{%
\textbf{उदाहरणम्:}\\[0.3em]
दश पणैः सपणः---\\[0.3em]
\textbf{न्यास:} १० पैसे, ८ पैसे
}{%
\textbf{Example:}\\[0.3em]
Ten paṇas with seven paṇas---\\[0.3em]
\textbf{Statement:} 10 paise, 8 paise
}

\vspace{0.5em}
\noindent\longline
\vspace{0.5em}

%% ==================== CHAPTER 5: Shunya-parikarma ====================
\section{\dev{शून्यपरिकर्म} --- Operations with Zero}

\bilingual{%
\textbf{अध्यायः ५: शून्यपरिकर्म}\\[0.3em]
एषः ब्राह्मस्फुटसिद्धान्तस्य सर्वप्रसिद्धतमः खण्डः अस्ति। ब्रह्मगुप्तः एव प्रथमः गणितज्ञः अभवत् यः शून्यस्य (\dev{शून्य}) कृते व्यवस्थितानि नियमानि स्थापितवान्---शून्येन भाजनस्य विवादास्पदं नियमम् अपि संस्मरन्।
}{%
\textbf{Chapter 5: Operations with Zero}\\[0.3em]
This is one of the most celebrated sections of the \textit{Brāhmasphuṭasiddhānta}. Brahmagupta was the first mathematician in history to lay down systematic rules for arithmetic operations involving zero (\textit{śūnya}), including the much-discussed rule for division by zero.
}

\bilingual{%
\begin{shloka}
ऋणे गुणे ऋणं क्षये गुणे यः क्षयम्\\
ऋणे तु ऋणं तद्धनं धनयोः\\
शून्येन निहतौ जातौ पुनः शून्यौ सकलविधौ\\
खहरं भजेत तु शून्येन निःशेषं तु यत् ॥\vnum{}
\end{shloka}
}{%
\textit{Translation:} ``A debt multiplied by a debt is a debt; a loss multiplied by a loss is a loss. But a debt [multiplied by a] debt is positive for positives. Those multiplied by zero become zero again in all ways. That which is exactly divisible by zero becomes a \textit{khahara}.''
}

\bilingual{%
\textbf{टीका:}\\[0.3em]
शून्य की गुणा और भाग सम्बन्धी नियम---धन और ऋण की गुणा करने से शून्य प्राप्त होता है। शून्य से किसी संख्या को गुणा करने से शून्य प्राप्त होता है। शून्य को शून्य से भाग देने से शून्य। परन्तु जब किसी संख्या को शून्य से भाग दिया जाय तब वह \dev{खहर} (अनिश्चित/अनन्त राशि) कहलाती है।\\[0.3em]
यदि $6 \\div 0 = $ खहर इति।
}{%
\textbf{Commentary:}\\[0.3em]
Rules relating to multiplication and division with zero---Multiplying positive and negative yields zero. Multiplying any number by zero gives zero. Dividing zero by zero gives zero. But when a number is divided by zero, it is called \textit{khahara} (an undefined/infinite quantity).\\[0.3em]
If $6 \\div 0 = $ \textit{khahara}, thus.
}

\bilingual{%
\textbf{सप्तमाध्याये व्यक्ताव्यक्तवर्णनम्}---अष्टादशे अध्याये (बीजगणिते) ब्रह्मगुप्तः विस्तरत्:अर्थात् वदति---
}{%
\textbf{In the 18th chapter (Algebra), Brahmagupta elaborates further:}
}

\bilingual{%
\begin{shloka}
योगे शून्यं खण्डिते च पूर्ववत् साधनम्\\
शून्येऽपि गुणके जाते पुनस्तत्रैव राशयः\\
भाज्ये शून्यं भागकारे तु जाते खहरं भवेत्\\
तस्य प्रसज्जते विकारो न तु शून्यं ततः ॥
\end{shloka}
}{%
\textit{Translation:} ``In addition, zero; in division also, the previous method. When even zero becomes a multiplier, the quantities become zero again there. When the dividend is zero and the divisor is zero, a \textit{khahara} comes into being. Its modification is produced, but not zero from that.''
}

\bilingual{%
\textbf{टीका:}\\[0.3em]
शून्य से योग, खण्डन आदि में पूर्ववत् साधन होता है। शून्य से गुणा करने पर राशियाँ शून्य हो जाती हैं। जब भाज्य शून्य हो और भागकार (भाजक) भी शून्य हो तब खहर राशि प्रसज्जते (उत्पन्न होती है)। इस विकार का कारण है कि शून्य से भाग करने पर कोई निश्चित राशि नहीं मिलती।\\[0.3em]
ऋणात्मक राशि शून्यतोऽपि---ऋणात्मक राशि की शून्य से भी गुणा होती है। यदि $6 \\div 0 = $ खहर तो इसमें कोई विकार नहीं, बल्कि यह एक विशेष प्रकार की राशि है।
}{%
\textbf{Commentary:}\\[0.3em]
In addition, division, etc., with zero, the previous method applies. When multiplied by zero, quantities become zero. When the dividend is zero and the divisor is zero, a \textit{khahara} quantity is produced. The reason for this modification is that no definite quantity is obtained on division by zero.\\[0.3em]
Negative quantities [multiplied] by zero---negative quantities are also multiplied by zero. If $6 \\div 0 = $ \textit{khahara}, then there is no error in this; rather, it is a special kind of quantity.
}

\bilingual{%
\textbf{शून्यस्य नियमसारांशः---ब्रह्मगुप्तकीयाः नियमाः}\\[0.3em]
\begin{align*}
a + 0 &= a \\
a - 0 &= a \\
a \\times 0 &= 0 \\
0 \\times 0 &= 0 \\
0 \\div a &= 0 \\quad (a \\neq 0) \\
a \\div 0 &= \\textrm{खहर} \\quad (\\textrm{अनिश्चित/अनन्त राशि})
\end{align*}
}{%
\textbf{Summary of Brahmagupta's Rules for Zero}\\[0.3em]
\begin{align*}
a + 0 &= a \\
a - 0 &= a \\
a \\times 0 &= 0 \\
0 \\times 0 &= 0 \\
0 \\div a &= 0 \\quad (\\text{for } a \\neq 0) \\
a \\div 0 &= \\textrm{\\textit{khahara}} \\quad (\\text{an undefined/infinite quantity})
\end{align*}
}

\bilingual{%
\textbf{टीका:}\\[0.3em]
यहाँ कोलब्रक साहेब ने उत्तर लिखा है जो बिलकुल असर्ज़त हैं इति। उपपत्ति में ऋणात्मक राशि की प्रशंसा निम्न लिखित श्लोक से की हैं। लीलावती में ``न्ये गुणके जाते खं हारइचेद्'' इत्यादि की उपपत्ति में किसी ने व्यस्त त्रैराशिक के लक्षण नहीं कहे हैं! इसका लक्षण भास्कराचार्य ने कहा है इति।\\[0.3em]
खभक्त इति पृच्छाया उत्तरं खहरात्मकम्॥
}{%
\textbf{Commentary:}\\[0.3em]
Here Colebrooke Sahib has written answers which are entirely inappropriate. In the demonstration, the praise of negative quantities has been made by the verse written below. In the \textit{Līlāvatī}, ``\textit{nye guṇake jāte khaṃ hāraiced}'' etc.---in the demonstration of this, no one has stated the characteristics of the busy rule of three! Bhāskarācārya has stated its characteristic.\\[0.3em]
\textit{Khabhakta}---the answer to the question is of the nature of \textit{khahara}.
}

\vspace{0.5em}
\noindent\longline
\vspace{0.5em}

%% ==================== CHAPTER 6: Shredhi (Series) ====================
\section{\dev{श्रेढी} --- Progressions and Series}

\bilingual{%
\textbf{अध्यायः ६: श्रेढी}\\[0.3em]
एषः विस्तृतः खण्डः समान्तरश्रेढीम्, गुणोत्तरश्रेढीम्, श्रेढीनां योगम्, तथा सम्बद्धानि विषयानि संस्पृशति।
}{%
\textbf{Chapter 6: Progressions and Series}\\[0.3em]
This extensive section treats arithmetic and geometric progressions, sums of series, and related topics.
}

\subsection{\dev{सङ्कलितैक्यानयनम्} --- Sum of Sums of Natural Numbers}

\bilingual{%
\begin{shloka}
द्वियुक्तगच्छाभिहतं द्विभक्तं मनस्विनः\\
सङ्कलितैक्यं व्याख्यातम् ॥१९॥\vnum{१९}
\end{shloka}
}{%
\textit{Translation:} ``The number of terms increased by two, multiplied [by the sum], divided by two---the sum of sums has been explained by the wise man [Brahmagupta].'' ||19||
}

\bilingual{%
\textbf{टीका:}\\[0.3em]
एक से आरम्भ कर इष्ट गच्छ का एकोत्तर सङ्कलित जो होता है उसको दो युक्त गच्छ में गुणा कर तीन से भाग देने में सङ्कलितैक्य प्रमाण होता है। गच्छ (पद) पर्यन्त एकादि अङ्कों के योग का नाम सद्वलित है उसके लिए आचार्य ने विधि नहीं लिखी है, सो ठीक नहीं है, क्योंकि गच्छ प्रमाण अधिक रहने से योग करने की क्रिया द्वारा सङ्कलित ज्ञान के लिये बहुत ही श्रम करना होगा, सङ्कलितानयन के लिये लीलावती में ``नैक पदघ्न पदार्ध मर्दकाद्यड युति'' इत्यादि भास्करोक्त रीति बहुत ही सुन्दर है।
}{%
\textbf{Commentary:}\\[0.3em]
Starting from one, the sum increased by one of the desired number of terms---multiplying it by the number of terms increased by two, and dividing by three, the measure of the sum of sums is obtained. The sum of natural numbers up to the number of terms is called \textit{sadvalita}. For this the Ācārya has not written a rule, which is not correct, because when the number of terms is large, one would have to labour greatly to know the sum by the operation of addition. For the finding of the sum, the method stated by Bhāskara in the \textit{Līlāvatī}---``\textit{naika padaghna padārtha mardakādyaḍa yuti}'' etc.---is very beautiful.
}

\bilingual{%
\textbf{उदाहरणम्:}\\[0.3em]
एकादि में नौ पदों तक अङ्कों के पृथक् पृथक् सङ्कलित कहो, और उन्हीं अङ्कों के सङ्कलितैक्यों को कहो॥\\[0.3em]
\textbf{न्यास:} १ | २ | ३ | ४ | ५ | ६ | ७ | ८ | ९\\[0.3em]
भास्करोक्त सूत्रानुसारेण पृथक् पृथक् सङ्कलितानि: $1, 3, 6, 10, 15, 21, 28, 36, 45$\\[0.3em]
आचार्योक्त सूत्रानुसारेण सङ्कलितैक्यानि: $1, 4, 10, 20, 35, 56, 84, 120, 165$
}{%
\textbf{Example:}\\[0.3em]
State separately the successive sums of natural numbers up to nine terms, and also their sum of sums.\\[0.3em]
\textbf{Statement:} 1 | 2 | 3 | 4 | 5 | 6 | 7 | 8 | 9\\[0.3em]
By Bhāskara's rule, the successive sums: $1, 3, 6, 10, 15, 21, 28, 36, 45$\\[0.3em]
By the Ācārya's rule, the sums of sums: $1, 4, 10, 20, 35, 56, 84, 120, 165$
}

\subsection{\dev{उपपत्तिः} --- Demonstration}

\bilingual{%
\textbf{उपपत्तिः:}\\[0.3em]
यदि पद = ५ तब एकादि अङ्कों को पदपर्यन्त क्रम से और उत्क्रम से स्थापन करने में:
\[
\begin{array}{ccccc}
1 & 2 & 3 & 4 & 5 \\
5 & 4 & 3 & 2 & 1
\end{array}
\]
इनका योग करने से प्रत्येक स्थान में $6 = $ पद $+ १$, यह पद पर्यन्त है। अतः:
\[
\\frac{\\text{पद}(\\text{पद} + 1)}{2} = \\text{सङ्कलित}
\]
\[
\\frac{5 \\times 6}{2} = 15 \\quad \\text{(सङ्कलित)}
\]
इसमें भास्करोक्त सङ्कलितानयन उपपन्न होता है॥
}{%
\textbf{Demonstration:}\\[0.3em]
If terms = 5, then establishing the natural numbers up to the terms in order and reverse order:
\[
\begin{array}{ccccc}
1 & 2 & 3 & 4 & 5 \\
5 & 4 & 3 & 2 & 1
\end{array}
\]
Adding these, in each place $6 = $ terms $+ 1$, this extending over the terms. Therefore:
\[
\\frac{\\text{terms}(\\text{terms} + 1)}{2} = \\text{sum}
\]
\[
\\frac{5 \\times 6}{2} = 15 \\quad \\text{(sum)}
\]
Thus Bhāskara's method for finding the sum is demonstrated.
}

\subsection{\dev{वर्गसङ्कलितम्} --- Sum of Squares}

\bilingual{%
\begin{shloka}
गच्छपदघ्नं पदार्धं मर्दकाद्यड युतिम्\\
वर्गसङ्कलितं व्याख्यातम् ॥\vnum{}
\end{shloka}
}{%
\textit{Translation:} ``The number of terms multiplied by the terms, half the terms [combined with] the addition of the compressing etc. [i.e., $1/3$]---the sum of squares has been explained.''
}

\bilingual{%
\textbf{टीका:}\\[0.3em]
यदि $k = 1 + 4 + 9 + 16 + ...$ पदपर्यन्त, तव इस श्रेढी का योगफल किमित्यानीयते?\\[0.3em]
$k = 1 + 4 + 9 + 16 + ...$ पदपर्यन्त\\[0.3em]
$\\displaystyle = \\frac{n(n+1)(2n+1)}{6}$
}{%
\textbf{Commentary:}\\[0.3em]
If $k = 1 + 4 + 9 + 16 + ...$ up to [the given] terms, how is the sum of this series to be found?\\[0.3em]
$k = 1 + 4 + 9 + 16 + ...$ up to terms\\[0.3em]
$\\displaystyle = \\frac{n(n+1)(2n+1)}{6}$
}

\subsection{\dev{घनसङ्कलितम्} --- Sum of Cubes}

\bilingual{%
\begin{shloka}
द्वियुक्तगच्छाभिहतं द्विभक्तं मनस्विनः\\
घनसङ्कलितैक्यं व्याख्यातम् ॥\vnum{}
\end{shloka}
}{%
\textit{Translation:} ``The number of terms increased by two, multiplied, divided by two---the sum of the sum of cubes has been explained by the wise one.''
}

\bilingual{%
\textbf{उदाहरणम्:}\\[0.3em]
यदि $k = 1 + 8 + 27 + 64 + 125 + ...$ पदपर्यन्त, तव इस श्रेढी का योगफल किमित्यानीयते?\\[0.3em]
\textbf{नवमिर्गखनेन भजनेन च हि:}\\[0.3em]
$\\displaystyle \\left\\{\\frac{n(n+1)}{2}\\right\\}^2$
}{%
\textbf{Example:}\\[0.3em]
If $k = 1 + 8 + 27 + 64 + 125 + ...$ up to terms, how is the sum of this series to be found?\\[0.3em]
\textbf{By dividing by nine and multiplying:}\\[0.3em]
$\\displaystyle \\left\\{\\frac{n(n+1)}{2}\\right\\}^2$
}

\subsection{\dev{गुणोत्तरश्रेढी} --- Geometric Progression}

\bilingual{%
\begin{shloka}
व्येकं व्येकगुणोद्धृतमित्यादिना\\
गुणोत्तरश्रेढी योगफलं व्याख्यातम् ॥\vnum{}
\end{shloka}
}{%
\textit{Translation:} ``Decreased by one, divided by the common ratio decreased by one, etc.---the sum of a geometric progression has been explained.''
}

\bilingual{%
\textbf{उदाहरणम्:}\\[0.3em]
यदि $k = 4 + 44 + 444 + ...$ पदपर्यन्त, तव इस श्रेढी का योगफल किमित्यानीयते?\\[0.3em]
नौ से गुणा करने से और भाग देने से:
\[
k = \\frac{4}{9}(9 + 99 + 999 + \\cdots)
\]
\[
= \\frac{4}{9}\\left\\{(10-1) + (10^2-1) + (10^3-1) + \\cdots\\right\\}
\]
\[
= \\frac{4}{9}\\left\\{\\frac{10(10^n - 1)}{10 - 1} - n\\right\\}
\]
}{%
\textbf{Example:}\\[0.3em]
If $k = 4 + 44 + 444 + ...$ up to terms, how is the sum of this series to be found?\\[0.3em]
By multiplying by nine and dividing:
\[
k = \\frac{4}{9}(9 + 99 + 999 + \\cdots)
\]
\[
= \\frac{4}{9}\\left\\{(10^1-1) + (10^2-1) + (10^3-1) + \\cdots\\right\\}
\]
\[
= \\frac{4}{9}\\left\\{\\frac{10(10^n - 1)}{9} - n\\right\\}
\]
}

\bilingual{%
\textbf{टीका:}\\[0.3em]
एवमेव ५-५५-५५५ $+ ...$ पदपर्यन्तम्\\
६-६६-६६६ $+ ...$ पदपर्यन्तम्\\
७-७७-७७७ $+ ...$ पदपर्यन्तम्\\
८-८८-८८८ $+ ...$ पदपर्यन्तम्\\
९-९९-९९९ $+ ...$ पदपर्यन्तम्\\[0.3em]
आसां श्रेढीनां योगफलं पूर्व विधिनैव प्राप्यगच्छनीति॥
}{%
\textbf{Commentary:}\\[0.3em]
Similarly: $5 + 55 + 555 + ...$ up to terms\\
$6 + 66 + 666 + ...$ up to terms\\
$7 + 77 + 777 + ...$ up to terms\\
$8 + 88 + 888 + ...$ up to terms\\
$9 + 99 + 999 + ...$ up to terms\\[0.3em]
The sum of these series is obtained by the previous method.
}

\vspace{0.5em}
\noindent\longline
\vspace{0.5em}

%% ==================== CHAPTER 7: Kshetra-vyavahara ====================
\section{\dev{क्षेत्रव्यवहारः} --- Plane Geometry}

\bilingual{%
\textbf{अध्यायः ७: क्षेत्रव्यवहारः}\\[0.3em]
एषः महत्त्वपूर्णः खण्डः विविधेषां ज्यामितीयाकृतीनां---त्रिभुजानाम्, चतुर्भुजानाम्, वृत्तानाम्---क्षेत्रफलगणनान्, तेषां भुजकोटिकर्णानां, विकर्णानां, तथा लम्बानां सम्बन्धान् च संस्पृशति।
}{%
\textbf{Chapter 7: Plane Geometry}\\[0.3em]
This major section deals with the calculation of areas of various geometric figures---triangles, quadrilaterals, circles---and the relationships between their sides, diagonals, and altitudes.
}

\subsection{\dev{समलम्बचतुर्भुजस्य क्षेत्रफलम्} --- Area of an Isosceles Trapezium}

\bilingual{%
\begin{shloka}
बाहुकोट्योर्हताः परस्परं क्षेत्रयोरिह\\
तेषु भूमिरधिकोऽर्धको मुखं च विषमस्य द्वयम् ॥\vnum{}
\end{shloka}
}{%
\textit{Translation:} ``The mutual products of the sides and perpendiculars of two fields---among them, the greater is the base, half is the face, and two [come from] the unequal [sides].''
}

\bilingual{%
\textbf{टीका:}\\[0.3em]
वर्ग क्षेत्रोपरिगत वृत्त का व्यास उसका कर्ता ही होता है, भ्राव्य क्षेत्र में भी ऐसा ही होता है, आचार्योक्त पद में अविरुद्ध शब्द से समलम्ब समलम्बान्तर चतुर्भुज समझना चाहिये जहाँ दोनों कर्ता दोनों भुज बराबर हैं। कर्ता को पार्श्वबाहु भुज से गुणा कर द्विगुणीत लम्ब से भाग देने से चतुर्भुजोपरिगत वृत्त का व्यासार्ध होता है। विषम चतुर्भुज में जिसमें उसके उपर वृत्त हो सकता है उसमें भुज और प्रतिभुज (सम्मुख भुज) के चतुर्भुज फल का शब्द व्यासार्ध होता है॥
}{%
\textbf{Commentary:}\\[0.3em]
The diameter of the circumscribed circle of a square field is its diagonal; in an oblong field also it is the same. In the Ācārya's text, by the word \textit{aviruddha} one should understand an isosceles trapezium (\textit{samlamba}), where both diagonals and both sides are equal. Multiplying the diagonal by the side and dividing by twice the altitude, the semi-diameter of the circumscribed circle of the quadrilateral is obtained. In a scalene quadrilateral in which a circumscribed circle is possible, the semi-diameter is the square root of the product of the side and the opposite side.
}

\subsection{\dev{वृत्तस्य व्यासज्ञानम्} --- Finding the Diameter of a Circumscribed Circle}

\bilingual{%
\begin{shloka}
विपर्यस्तपादं भुजयोः कर्णं द्विगुणावलम्बनविभक्तः\\
हृदयं द्विसस्य भुजप्रतिभुजकृतियोगसूलार्धम् ॥२६॥\vnum{२६}
\end{shloka}
}{%
\textit{Translation:} ``The transposed side-product of the two sides, the diagonal divided by twice the altitude---the heart [semi-diameter] of that. For a quadrilateral, the semi-square-root of the sum of the products of the sides and opposite sides.'' ||26||
}

\bilingual{%
\textbf{टीका:}\\[0.3em]
वर्ग क्षेत्रे तथापय्यते कर्ण एवं व्यास उति स्फुटस्। व्युत्क्रमेण च द्वि ननस्य च तुल्यौ भुजौ उच्यते। सेनायमर्थः। विपर्यस्तपादं भुज द्वियुगन लम्ब विभक्तः फल हृदय च चतुर्भुजोपरिगत वृत्तस्य व्यासार्ध भवेत्। विषम चतुर्भुजे यत्र तदुर्पारिग तदत्तं भवितुमर्हति तत्र भुज प्रतिभुजयोः वर्गयोगसूलार्धं हृद व्यासार्ध भवेत्॥
}{%
\textbf{Commentary:}\\[0.3em]
In a square field, the diagonal itself is clearly the diameter. By transposition, two equal sides are spoken of. The meaning is this: the transposed side-product of the sides, divided by twice the altitude---the result is the heart (semi-diameter) of the circle circumscribed about the quadrilateral. In a scalene quadrilateral where a circumscribed circle is possible, the semi-square-root of the sum of the products of the side and opposite side becomes the heart (semi-diameter).
}

\subsection{\dev{त्रिभुजोपरिगतवृत्तव्यासज्ञानार्थम्} --- Circle Diameter from Triangle}

\bilingual{%
\begin{shloka}
इदानीं त्रिभुजोपरिगतवृत्तस्य व्यासज्ञानार्थमाह\\
त्रिभुजस्य बाहोर्भुजयोर्हृतः कर्णो द्विगुणावलम्बरविभक्तः\\
हृदयं द्विसस्य भुजप्रतिभुजकृतियोगसूलार्धम् ॥२७॥\vnum{२७}
\end{shloka}
}{%
\textit{Translation:} ``Now, for knowing the diameter of the circle circumscribed about a triangle, he says: the product of the two sides of the triangle, divided by the altitude [and] multiplied by the diagonal---the heart [is] the semi-square-root of the sum of the products of the side and opposite side.'' ||27||
}

\bilingual{%
\textbf{टीका:}\\[0.3em]
वर्ग क्षेत्रे श्रायते च कर्ण एव तदपरिग नवृत्त व्यासः। व्युत्क्रमेण समलम्बचतुर्भुजं वोध्यम्। यत्र कर्णौ भुजौ च तुल्यौ, कर्णः पादवं मञ्जेन गुणो द्विगुणावलम्ब विभक्तः फल हृदयमर्थात् चतुर्भुजोपरिगत वृत्तस्य व्यासार्ध भवेत्। विषम चतुर्भुजे यत्र तदुर्पारिग तदत्तं भवितुमर्हति तत्र भुज प्रतिभुजयो वर्गयोगसूलार्धं हृद (व्यासार्ध) भवेत्॥२७॥
}{%
\textbf{Commentary:}\\[0.3em]
In a square field, the diagonal itself is the diameter of the circumscribed circle. An isosceles trapezium is to be understood by transposition. Where the diagonals and sides are equal, the diagonal, multiplied by the side and divided by twice the altitude---the result is the heart, i.e., the semi-diameter of the circle circumscribed about the quadrilateral. In a scalene quadrilateral where a circumscribed circle is possible, the semi-square-root of the sum of the products of the side and opposite side is the heart (semi-diameter). ||27||
}

\subsection{\dev{जात्ययोः क्षेत्रयोः} --- On Two Classes of Triangles}

\bilingual{%
\begin{shloka}
जात्ययोस्तिहताः परस्परं क्षेत्रयोरिह हि बाहुकोटयः\\
तेषु भूमिरधिकोऽर्धको मुखं विषमस्य द्वयम् ॥\vnum{}
\end{shloka}
}{%
\textit{Translation:} ``The sides and perpendiculars of two regular [triangular] fields here, mutually multiplied---among them the greater is the base, half [is] the face, and two [come from] the unequal [sides].''
}

\bilingual{%
\textbf{टीका:}\\[0.3em]
दो जात्य त्रिभुज के भुज और कोटि को परस्पर क्रम से गुणा करने से विषम चतुर्भुज के भुज होते हैं उनमें श्रेष्ठक भू होती है, लघु मुख होता है, और मध्य दोनों भुज होते हैं।
}{%
\textbf{Commentary:}\\[0.3em]
Multiplying the sides and perpendiculars of two regular (\textit{jātya}) triangles mutually in order, the sides of a scalene quadrilateral are produced; among them the greatest is the base, the smallest is the face, and the two middle ones are the [other] sides.
}

\subsection{\dev{विषमचतुर्भुजक्षेत्रफलम्} --- Area of a Scalene Quadrilateral}

\bilingual{%
\begin{shloka}
विषमं चतुर्भुजं क्षेत्रं समलम्बवत् सदृशम्\\
तस्य कर्णे लम्बौ संपातं तयोः परस्परं गुणनम्\\
ततोऽर्धं फलं स्यात् ॥\vnum{}
\end{shloka}
}{%
\textit{Translation:} ``A scalene quadrilateral field is similar to an isosceles trapezium. Its diagonals [are] the two altitudes [meeting] at the intersection; their mutual product, then half, is the area.''
}

\bilingual{%
\textbf{टीका:}\\[0.3em]
विषम चतुर्भुज क्षेत्र समलम्बवत् (समलम्बचतुर्भुज की तरह) सदृश (समान) है। तस्य (उसके) कर्णे लम्बौ (दो लम्ब) संपात (प्रतिच्छेदन) बिन्दु में एक दूसरे को काटते हैं तब तयोः (उन लम्बों का) परस्परं गुणन (गुणा) ततोऽर्धं (आधा) फलं स्यात् (फल होता है)।
}{%
\textbf{Commentary:}\\[0.3em]
A scalene quadrilateral field is similar (\textit{sadṛśa}) to an isosceles trapezium (\textit{samlamba-vat}). Its diagonals [are] the two altitudes [meeting] at the intersection (\textit{sampāta}) point where they cut each other; then their mutual product, halved, is the area.
}

\subsection{\dev{त्रिभुजक्षेत्रफलम्} --- Area of a Triangle}

\bilingual{%
\begin{shloka}
समकोणत्रिभुजे कर्णः समपादभुजयोर्वर्गयोगसूलम्\\
तत्कर्णस्य चतुर्गुणीतो भुजयोर्घातः समकोणे ॥\vnum{}
\end{shloka}
}{%
\textit{Translation:} ``In a right-angled triangle, the hypotenuse is the square root of the sum of the squares of the equal side and base. The product of the sides, multiplied by four [divided by] that hypotenuse, [gives the altitude] in the right-angled [triangle].''
}

\bilingual{%
\textbf{उदाहरणम्---समकोणत्रिभुजे:}\\[0.3em]
\[
c^2 = a^2 + b^2 \\quad \\text{(यत्र $c$ = कर्णः)}
\]
}{%
\textbf{Example---Right-Angled Triangle:}\\[0.3em]
\[
c^2 = a^2 + b^2 \\quad \\text{(where $c$ = hypotenuse)}
\]
}

\vspace{0.5em}
\noindent\longline
\vspace{0.5em}

%% ==================== CHAPTER 8: Khata-vyavahara ====================
\section{\dev{खातव्यवहारः} --- Excavations (Solid Geometry)}

\bilingual{%
\textbf{अध्यायः ८: खातव्यवहारः}\\[0.3em]
एषः खण्डः खातानाम्, कूपानाम्, सरोवराणाम्, तथा सम्बद्धानां घनाकृतीनां घनफलगणनान् संस्पृशति।
}{%
\textbf{Chapter 8: Excavations (Solid Geometry)}\\[0.3em]
This section deals with the calculation of volumes of excavations, wells, reservoirs, and related solid figures.
}

\subsection{\dev{समखातघनफलम्} --- Volume of a Regular Excavation}

\bilingual{%
\begin{shloka}
समखातं तु यत् क्षेत्रफलं तद्वेधगुणितं घनफलम्\\
समखातस्य त्रिभागः सूचीखातस्य घनफलम् ॥\vnum{}
\end{shloka}
}{%
\textit{Translation:} ``That regular excavation of which [the area]---the area of the field multiplied by the depth is the volume. One-third of the regular excavation is the volume of a needle-excavation [pyramid/cone].''
}

\bilingual{%
\textbf{टीका:}\\[0.3em]
समखात का क्षेत्रफल उसी तरह ज्ञात करें जैसे समचतुर्भुज या समलम्बचतुर्भुज का क्षेत्रफल ज्ञात होता है। क्षेत्रफल को वेध (गहराई) से गुणा करने से समखात का घनफल होता है। अर्थात् जिस खात के मुख में जितने दैर्घ्यादि हैं उतने तल में रहने से वह समखात कह लाता है उसके घनफल को तीन से भाग देने से सूचीकार खात का घनफल होता है॥
}{%
\textbf{Commentary:}\\[0.3em]
The area of a regular excavation is found just as the area of a square or isosceles trapezium is found. The area multiplied by the depth (\textit{vedha}) gives the volume of the regular excavation. That is, an excavation in which the length etc. at the mouth are the same at the base is called a regular excavation (\textit{sama-khāta}). Dividing its volume by three gives the volume of a needle-like excavation (\textit{sūcī-khāta}, i.e., pyramid/cone).
}

\subsection{\dev{उदाहरणम्} --- Worked Example}

\bilingual{%
\textbf{उदाहरणम्:}\\[0.3em]
किसी वापी (वावली) की दीर्घता (लम्बाई) तीस हाथ है, विस्तार साठ हाथ। उस वापी के भीतर चार, पाँच, छः, सात, आठ भुज खण्डों से पाँच खात (खत्ता) हैं, खातों के क्रम से वेध ८, ७, ७, ३, २ हैं तब उन खातों के समरज्जु प्रमाण कहो।\\[0.3em]
\textbf{समरज्जु:} मुखतल भुजैक्य त्रि से $36|35|42|21|16$ इन सबों के योग को एकाग्र $30$ से भाग देने से:
\[
\\text{समरज्जु} = \\frac{\\text{विस्तार} \\times \\text{दैर्घ्य}}{30}
\]
\[
\\text{क्षेत्रफल} = \\text{विस्तार} \\times \\text{दैर्घ्य} = 8 \\times 30 = 240
\]
इस क्षेत्रफल को समरज्जु से गुणा करने से सम्पूर्ण खात का घनफल १२०० हुआ॥
}{%
\textbf{Example:}\\[0.3em]
A certain tank (\textit{vāpī}) has length thirty cubits, breadth sixty cubits. Inside it are five excavations (\textit{khāta}) with [side] portions of four, five, six, seven, eight [cubits]; the depths of the excavations in order are 8, 7, 7, 3, 2. State the mean section (\textit{samarajju}) measure of those excavations.\\[0.3em]
\textbf{Mean section:} The side-sums at mouth and base [give] 36|35|42|21|16. Dividing the sum of all these by the aggregate 30:
\[
\\text{Mean section} = \\frac{\\text{breadth} \\times \\text{length}}{30}
\]
\[
\\text{Area} = 8 \\times 30 = 240
\]
Multiplying this area by the mean section, the total excavation volume became 1200.
}

\subsection{\dev{विस्तारदैर्घ्यवेधेषु वैषम्ये} --- Irregular Excavations}

\bilingual{%
\textbf{उदाहरणानि:}\\[0.3em]
किसी पुष्करिणी की लम्बाई १६ हाथ है, विस्तार = ५ हाथ, वेध २, ३, ४ हाथ हैं तब उसका फल क्या होगा।\\[0.3em]
\textbf{न्यास:} २|३|४ वेध, वेधों की सममिति $= \\frac{9}{3} = ३$, यहाँ स्थान संख्या = ३\\[0.3em]
$\\text{क्षेत्रफल} = 5 \\times 16 = 80$\\[0.3em]
इसको वेध से गुणा करने से: $80 \\times 3 = 240$ खातफल।
}{%
\textbf{Examples:}\\[0.3em]
A certain pond has length 16 cubits, breadth = 5 cubits, depths 2, 3, 4 cubits; then what is its volume?\\[0.3em]
\textbf{Statement:} 2|3|4 depths; mean of depths $= \\frac{9}{3} = 3$; here the number of places = 3.\\[0.3em]
$\\text{Area} = 5 \\times 16 = 80$\\[0.3em]
Multiplying by depth: $80 \\times 3 = 240$ (excavation volume).
}

\bilingual{%
दूसरा उदाहरण---जिस खात का दैर्घ्य = १२ हाथ है, वेध = ८ हाथ है, विस्तार ३|४|५ हाथ, उस खात का फल कतो।\\[0.3em]
विस्तृति सममिति $= \\frac{3+4+5}{3} = ४$\\[0.3em]
$\\text{क्षेत्रफल} = 4 \\times 12 = 48$\\[0.3em]
इसको वेध से गुणा करने से: $48 \\times 8 = 384$ (खात फल)।
}{%
\textbf{Second Example:} An excavation of which the length = 12 cubits, depth = 8 cubits, breadth 3|4|5 cubits---state its volume.\\[0.3em]
Mean breadth $= \\frac{3+4+5}{3} = 4$\\[0.3em]
$\\text{Area} = 4 \\times 12 = 48$\\[0.3em]
Multiplying by depth: $48 \\times 8 = 384$ (excavation volume).
}

\vspace{0.5em}
\noindent\longline
\vspace{0.5em}

%% ==================== CHAPTER 9: Dhanyarashi ====================
\section{\dev{धान्यराशिव्यवहारः} --- Grain Heaps}

\bilingual{%
\textbf{अध्यायः ९: धान्यराशिव्यवहारः}\\[0.3em]
एषः खण्डः शङ्कुाकारधान्यराशीनां मापनं संस्पृशति।
}{%
\textbf{Chapter 9: Grain Heaps}\\[0.3em]
This section treats the measurement of grain stored in heaps, which have the form of a cone or portion thereof.
}

\subsection{\dev{स्थूलधान्यराशिमानम्} --- Measuring Coarse Grain Heaps}

\bilingual{%
\begin{shloka}
वृत्तपरिगतधान्यराशेः परिधिनवमांशो वेधः\\
वेधेन गुणितः परिधिः षष्ट्यंशः फलं स्यात् ॥५०॥\vnum{५०}
\end{shloka}
}{%
\textit{Translation:} ``The circumference of a circular grain heap---a ninth of the circumference is the depth. The circumference multiplied by the depth, [then] a sixtieth of that, is the volume [measure].'' ||50||
}

\bilingual{%
\textbf{टीका:}\\[0.3em]
भित्त्यन्तर्बाह्यकोणागः परिधिः द्विचतुः सत्रिभागगुणः प्राग्वत्‌ गणितं कृत्वा स्वगुणकारभक्तं तदा तद्‌ गणितं भवति। अर्थात्‌ भित्तिपादसंलग्नस्य धान्यराशेः परिधिमानं द्वाभ्यां संगुण्य तस्मात्पूर्वंवचत्फलं तदुद्वाभ्यां भक्तं तदा तद्राशेधनफलं भवति।
}{%
\textbf{Commentary:}\\[0.3em]
The inner-outer corner circumference, multiplied by two-four [and] a seventh part, having performed the calculation as before, divided by its multiplier, then that calculation becomes [the volume]. That is, the circumference measure of the grain heap touching the base of the wall, multiplied by two, then the result as before, divided by those two, then the volume of that heap becomes [known].
}

\subsection{\dev{उदाहरणानि} --- Worked Examples}

\bilingual{%
\textbf{स्थूलधान्यराशिमानज्ञानार्थम्:}\\[0.3em]
\textbf{न्यास:} स्थूलधान्यराशि परिधिः ६०, परिषेधांशमांशः $६ = $ वेधः\\[0.3em]
$\\text{परिधेः षष्ट्यंशः} = \\frac{60}{10} = 10$, वर्गितः $= 100$\\[0.3em]
$\\text{वेधगुणितः} = 100 \\times 6 = 600$ (स्थूलधान्यराशि घनफलम्)\\[1em]
\textbf{शूकधान्यराशिमानज्ञानार्थम्:}\\[0.3em]
\textbf{न्यास:} परिधिः ६०, परिधिः नवमांशो वेधः ५\\[0.3em]
$\\text{परिषेः षष्ट्यंशः} = \\frac{60}{6} = 10$, वर्गितः $= 100$\\[0.3em]
$\\text{वेधगुणितः} = 100 \\times 5 = 500$ (शूकधान्यराशि घनफलम्)
}{%
\textbf{For Knowing the Measure of Coarse Grain Heaps:}\\[0.3em]
\textbf{Statement:} Coarse grain heap circumference 60, portion $6 = $ depth.\\[0.3em]
$\\text{Sixtieth of circumference} = \\frac{60}{10} = 10$, squared $= 100$\\[0.3em]
$\\text{Multiplied by depth} = 100 \\times 6 = 600$ (coarse grain heap volume)\\[1em]
\textbf{For Knowing the Measure of \textit{Śūka} Grain Heaps:}\\[0.3em]
\textbf{Statement:} Circumference 60, a ninth of circumference is depth 5.\\[0.3em]
$\\text{Sixtieth of circumference} = \\frac{60}{6} = 10$, squared $= 100$\\[0.3em]
$\\text{Multiplied by depth} = 100 \\times 5 = 500$ (\textit{śūka} grain heap volume)
}

\bilingual{%
\textbf{अणुधान्यराशिमानज्ञानार्थम्:}\\[0.3em]
\textbf{न्यास:} अणुधान्यराशि परिधिः ६०, वेधः = ३\\[0.3em]
$\\text{परिषेः षष्ट्यंशः} = \\frac{60}{?} = 5$, वर्गितः $= 100$\\[0.3em]
$\\text{वेधगुणितः} = 100 \\times 3 = 300$ (अणुधान्यराशि घनफलम्)
}{%
\textbf{For Knowing the Measure of Fine Grain Heaps:}\\[0.3em]
\textbf{Statement:} Fine grain heap circumference 60, depth = 3.\\[0.3em]
$\\text{Sixtieth of circumference} = \\frac{60}{?} = 5$, squared $= 100$\\[0.3em]
$\\text{Multiplied by depth} = 100 \\times 3 = 300$ (fine grain heap volume)
}

\subsection{\dev{उपपत्तिः} --- Demonstration}

\bilingual{%
\textbf{उपपत्तिः:}\\[0.3em]
धान्य स्थितिवशेन धान्यराशिः सूच्याकारो भवति तदाधारपरिधिनवमांशादिभागस्तद्वेघो भवतीति प्रत्यक्षापरोक्षया स्थिरीकृतः। वृत्तक्षेत्र परिधिगुणितव्यासः...
}{%
\textbf{Demonstration:}\\[0.3em]
By the nature of the grain's position, the grain heap becomes needle-shaped (conical); its base circumference divided by nine etc. becomes its depth---this is established by direct and indirect perception. The area of a circular field is the circumference multiplied by the diameter...
}

\subsection{\dev{भास्करीय लीलावत्यामुदाहरणम्} --- Example from Bhāskara's Līlāvatī}

\bilingual{%
\begin{shloka}
समभुवि किल राशिभिः स्थितः स्थूलधान्यः परिधिपरिमितिः\\
स्याद्व्यष्टिषट्यंशया प्रवद गणक खार्यः कि मिताः सन्ति तस्मिन्\\
अथ पृथगपुघान्यैरछुकधान्यैरच शीघ्रम् ॥\vnum{}
\end{shloka}
}{%
\textit{Translation:} ``On level ground indeed a coarse grain heap stands, with a [known] circumference measure. With the sixty-fourth part [of the circumference], tell me, calculator, how many \textit{khāri}s are in it? Also quickly [tell me] separately for \textit{pu} and \textit{śuka} grains.''
}

\vspace{0.5em}
\noindent\longline
\vspace{0.5em}

%% ==================== CHAPTER 10: Chhaya ====================
\section{\dev{छायाव्यवहारः} --- Shadows and Gnomons}

\bilingual{%
\textbf{अध्यायः १०: छायाव्यवहारः}\\[0.3em]
एषः खण्डः शङ्कुस्य (gnomon) छायायाः तथा सम्बद्धानां गणितीयानि संस्पृशति।
}{%
\textbf{Chapter 10: Shadows and Gnomons}\\[0.3em]
This section deals with the mathematics of the gnomon's shadow and related calculations.
}

\bilingual{%
\begin{shloka}
द्विजद्विसम्भागैकनिघ्नातु तु परिषैः फलम्\\
भित्त्यन्तर्बाह्यकोणागे परिधिः द्विचतुःसत्रिभागैकगुणः ॥५१॥\vnum{५१}
\end{shloka}
}{%
\textit{Translation:} ``The result [is obtained] from the circumference, multiplied by two, half of two [i.e., one], etc. At the inner-outer corner, the circumference multiplied by two-four and one-seventh [i.e., $2\\frac{4}{7}$].'' ||51||
}

\bilingual{%
\textbf{टीका:}\\[0.3em]
भित्त्यन्तर्बाह्यकोणागः परिधिः द्विचतुः सव्यंशगुणः प्राग्वत्‌ गणितं कृत्वा स्वगुणकारभक्तं तदा तद्‌ गणितं भवति। अर्थात्‌ भित्तिपादसंलग्नस्य धान्यराशेः परिधिमानं द्वाभ्यां संगुण्य तस्मात्पूर्वंवचत्फलं तदुद्वाभ्यां भक्तं तदा तद्राशिघनफलं भवति।\\[0.3em]
भित्त्योरन्तः कोणस्थितधान्यराशेः परिधि चतुर्गुणं त्वा ततः पूर्ववत्फलमानीय तच्चतुर्भक्तं तदा तद्धान्यराशिघनफलं भवति।\\[0.3em]
भित्त्योबहिः कोणस्थितधान्यराशेः परिधि सत्रिभागेकेन संगुण्य ततः पूर्वैवत्फलमानीय तत्फलं सत्रिभागैकेन भक्तं तदा तस्य धान्यराशिघनफलं भवेदिति॥
}{%
\textbf{Commentary:}\\[0.3em]
The inner-outer corner circumference, multiplied by two-four [and] one-seventh, having performed the calculation as before, divided by its multiplier, then that calculation becomes [known]. That is, the circumference measure of the grain heap touching the base of the wall, multiplied by two, then the result as before, divided by those two, then the volume of that heap becomes known.\\[0.3em]
For a grain heap situated at the inner corner of two walls, the circumference [is] multiplied by four; then having obtained the result as before, dividing it by four, then the volume of that grain heap becomes known.\\[0.3em]
For a grain heap situated at the outer corner of two walls, the circumference [is] multiplied by one-and-three-sevenths [i.e., $\\frac{10}{7}$]; then having obtained the result as before, dividing that result by one-and-three-sevenths, then the volume of that grain heap becomes known.
}

\vspace{0.5em}
\noindent\longline
\vspace{0.5em}

%% ==================== CHAPTER 11: Gitadhyaya ====================
\section{\dev{गीताध्यायः} --- The Chapter on Spheres}

\bilingual{%
\textbf{अध्यायः ११: गीताध्यायः}\\[0.3em]
एतत् सपुटस्य अन्तिमपृष्ठानि गीताध्यायम् (गोलाध्यायम्) आरभन्ते, यत् गोलस्य तथा गोलखण्डानां ज्यामिती संस्पृशति। एषः ब्राह्मस्पुटसिद्धान्तस्य एकविंशतितमः अध्यायः।
}{%
\textbf{Chapter 11: The Chapter on Spheres}\\[0.3em]
The final pages of this volume begin the \textit{Gītādhyāya} (Chapter on Spheres), which deals with the geometry of spheres and spherical segments. This is the 21st chapter of the \textit{Brāhmasphuṭasiddhānta}.
}

\subsection{\dev{विष्कम्भपरिधिः} --- Diameter and Circumference}

\bilingual{%
\begin{shloka}
विवृतद्विसम्भागैकनिघ्नात् तु परिषैः फलम्\\
इत्यादि भास्करोक्तमेतदनुरूपमेव ॥५१॥\vnum{५१}
\end{shloka}
}{%
\textit{Translation:} ``The result [is obtained] from the circumference, multiplied by one [half] of the expanded diameter-pair [i.e., the diameter], etc.---this is indeed in conformity with what Bhāskara has stated.'' ||51||
}

\bilingual{%
\textbf{टीका:}\\[0.3em]
विवृत (विस्तृत) द्विसम्भाग (द्विज) एकनिघ्नात् (एक से गुणा कर) तु परिषैः फलम्---इत्यादि भास्करोक्तमेतदनुरूपमेव (भास्कराचार्य के कथन के अनुरूप ही)॥५१॥
}{%
\textbf{Commentary:}\\[0.3em]
Expanded (\textit{vivṛta}) pair-of-diameter-half (\textit{dvisambhāga}) multiplied by one [i.e., the diameter]---the result from the circumference, etc.---this is indeed in conformity with what Bhāskarācārya has stated. ||51||
}

\subsection{\dev{राशिव्यवहारः} --- Calculation of Heaps}

\bilingual{%
\begin{shloka}
भित्यन्तर्बाह्यकोणागः परिधिः द्विचतुः सव्यंशगुणः प्राग्वत्\\
गणितं कृत्वा स्वगुणकारभक्तं तदा तद् गणितं भवति\\
अर्थात् भित्तिपादसंलग्नस्य धान्यराशेः परिधिमानं द्वाभ्यां संगुण्य\\
तस्मात्पूर्ववचत्फलं तदुद्वाभ्यां भक्तं तदा तद्राशिघनफलं भवति ॥\vnum{}
\end{shloka}
}{%
\textit{Translation:} ``The inner-outer corner circumference, multiplied by two-four [and] one-seventh [i.e., $\\frac{18}{7}$] as before; having performed the calculation, divided by its multiplier, then that calculation becomes [known]. That is, the circumference measure of the grain heap touching the base of the wall, multiplied by two; then the result as before, divided by those two, then the volume of that heap becomes known.''
}

\bilingual{%
\textbf{टीका:}\\[0.3em]
भित्त्योरन्तः कोणस्थितधान्यराशेः परिधि चतुर्गुणं त्वा ततः पूर्ववत्फलमानीय तच्चतुर्भक्तं तदा तद्धान्यराशिघनफलं भवति। भित्त्योबहिः कोणस्थितधान्यराशेः परिधि सत्रिभागेकेन संगुण्य ततः पूर्वैवत्फलमानीय तत्फलं सत्रिभागैकेन भक्तं तदा तस्य धान्यराशिघनफलं भवेदिति॥
}{%
\textbf{Commentary:}\\[0.3em]
For a grain heap situated at the inner corner of two walls, the circumference [is] multiplied by four; then having obtained the result as before, dividing it by four, then the volume of that grain heap becomes known. For a grain heap situated at the outer corner of two walls, the circumference [is] multiplied by one-and-three-sevenths; then having obtained the result as before, dividing that result by one-and-three-sevenths, then the volume of that grain heap becomes known.
}

\bilingual{%
\textbf{भास्करीय लीलावत्यामुदाहरणम्:}\\[0.5em]
\begin{shloka}
परिधिर्भित्ति लग्नस्य राशोः शत्करः किल\\
अन्तः कोणस्थितस्यापि तिथितुल्यकरः सखे\\
बहिः कोणस्थितस्यापि पञ्चघ्ननवसंमितः\\
तेषामाचक्ष्व मे क्षिप्रं धनहस्तान् पृथक् पृथक् ॥\vnum{}
\end{shloka}
}{%
\textbf{Example from Bhāskara's Līlāvatī:}\\[0.5em]
\textit{Translation:} ``The circumference of a heap touching a wall is indeed a hundred. Of one situated at the inner corner, [it is] equal to the number of lunar days [30], O friend. Of one situated at the outer corner, [it is] measured as nine multiplied by five [45]. Tell me quickly their \textit{dhana-hasta}s [wealth-measures] separately.''
}

\bilingual{%
\textbf{न्यासः:}\\[0.3em]
अत्र प्रथमस्य परिधिः ३० द्विनिष्ठः $= 30 \\times 2 = 60$\\
अन्तः कोणस्थितपरिधिः १५ चतुर्गुणितः $= 15 \\times 4 = 60$\\
बाह्यकोणस्थितपरिधिः ४८ सत्रिभागैक गुणितः $= 48 \\times \\frac{5}{4} = 60$\\
एषाविधः $= 6$\\
र एभ्यः फलं तुल्यमेतावत्य् एव खार्यः $= 600$\\
एतत् स्व स्व गुरोन भक्तं जातं पृथक् पृथक् फलम् $300 | 150 | 450$\\
एवं स्थूल धान्यस्य मानं जातम्।
}{%
\textbf{Statement:}\\[0.3em]
Here, the circumference of the first [heap] 30, doubled $= 30 \\times 2 = 60$.\\
The inner-corner circumference 15, quadrupled $= 15 \\times 4 = 60$.\\
The outer-corner circumference 48, multiplied by one-and-a-quarter $= 48 \\times \\frac{5}{4} = 60$.\\
This measure $= 6$.\\
From these, the equal result is just so many \textit{khāri}s $= 600$.\\
This, divided by its own multiplier, became the separate results $300 | 150 | 450$.\\
Thus the measure of coarse grain was found.
}

\subsection{\dev{अन्योदाहरणम्} --- Another Example}

\bilingual{%
\begin{shloka}
षट्त्रिशेन्मित परिधौ राशौ धान्यस्य कि भवेद् गणितम्\\
हस्त चतुष्काभ्युदये यदि वेत्सि तदुच्यतामाशु ॥\vnum{}
\end{shloka}
}{%
\textit{Translation:} ``When the circumference of a heap of grain is measured as thirty-six, what is the calculation [volume], if you know, tell me quickly, [given] a rise of four cubits [depth].''
}

\bilingual{%
\textbf{न्यासः:}\\[0.3em]
धान्यराशि परिधिः ३६, वेधः = ४ तदा पूर्ववत्‌ परिषेः षष्ट्यंशः $= 6$\\
वर्गितः $= 36$\\
वेधगुणितः $= 36 \\times 4 = 144$ घनफलम्।
}{%
\textbf{Statement:}\\[0.3em]
Grain heap circumference 36, depth = 4; then as before, the sixtieth of the circumference $= 6$.\\
Squared $= 36$.\\
Multiplied by depth $= 36 \\times 4 = 144$ (volume).
}

\vspace{0.5em}
\noindent\longline
\vspace{0.5em}

%% ==================== CONCLUDING REMARKS ====================
\section*{Concluding Remarks}
\addcontentsline{toc}{section}{Concluding Remarks}

\bilingual{%
\textbf{समाप्तिकथनम्}\\[0.5em]
एतत् सपुटं द्वितीयस्य भागस्य ब्राह्मस्फुटसिद्धान्तस्य समापयति, विस्तृतगणितीयविषयाणां परिष्कारं कुर्वत्। एतस्मिन् सम्पुटे समावेशिताः खण्डाः ब्रह्मगुप्तस्य गणितचिन्तनस्य गभीरतां तथा परिष्कारं दर्शयन्ति।
}{%
\textbf{Concluding Remarks}\\[0.5em]
This volume concludes the second part of the \textit{Brāhmasphuṭasiddhānta}, covering an extraordinary range of mathematical topics. The sections included in this portion demonstrate the depth and sophistication of Brahmagupta's mathematical thought.
}

\bilingual{%
\begin{enumerate}
\item \textbf{वाणिज्यगणितम्:} भागसमस्याः (\dev{भागे}), चक्रवृद्धिसूदः (\dev{मिश्र धन}), विनिमयव्यवहारः (\dev{भाण्डप्रतिभाण्डक}), तथा त्रैराशिकम् तस्य विस्ताराः च।
\item \textbf{शून्यस्य बीजगणितम्:} शून्येन क्रियासु ब्रह्मगुप्तस्य क्रान्तिकारि नियमाः, शून्येन भाजनस्य विवादास्पदं परन्तु मार्गदर्शकं चोपचारं समाविशन्तः---\dev{खहर} इति राशिं प्राप्य।
\item \textbf{श्रेढयः:} समान्तरश्रेढी, गुणोत्तरश्रेढी, प्राकृतिकसङ्ख्यानां वर्गानां घनानां च योगानां समग्रचिकित्सा।
\item \textbf{क्षेत्रज्यामितिः:} त्रिभुजानां, चतुर्भुजानां, वृत्तानां च क्षेत्रफलानि, चक्रीयचतुर्भुजानां तथा परिगतवृत्तानां सूत्रैः सह।
\item \textbf{घनगणितम्:} खातानां, कूपानां, धान्यराशीनां च घनफलगणनम्।
\item \textbf{गोलज्यामितिः:} गीताध्यायस्य (गोलाध्यायस्य) आरम्भः, गोलखण्डानां ज्यामिती संस्पृशन्।
\end{enumerate}
}{%
\begin{enumerate}
\item \textbf{Commercial Arithmetic:} Partnership problems (\textit{bhāga}), compound interest (\textit{miśra dhana}), barter exchange (\textit{bhāṇḍa-pratibhāṇḍaka}), and the Rule of Three and its extensions.
\item \textbf{The Algebra of Zero:} Brahmagupta's revolutionary rules for arithmetic with zero, including the controversial but pioneering treatment of division by zero as yielding the quantity \textit{khahara}.
\item \textbf{Series:} Comprehensive treatment of arithmetic and geometric progressions, sums of natural numbers, squares, and cubes.
\item \textbf{Plane Geometry:} Areas of triangles, quadrilaterals, and circles, with theorems on cyclic quadrilaterals and circumscribed circles.
\item \textbf{Solid Geometry:} Volumes of excavations, wells, and grain heaps.
\item \textbf{Spherical Geometry:} The beginning of the chapter on spheres (\textit{gītādhyāya}), dealing with the geometry of spherical segments.
\end{enumerate}
}

\vspace{1em}
\begin{center}
\longline\\[0.5em]
\textit{``The Brāhmasphuṭasiddhānta stands as one of the most influential}\\
\textit{mathematical texts in the history of world mathematics.''}\\[0.5em]
\longline
\end{center}

\vspace{1em}
\bilingual{%
\textbf{स्रोतसि टिप्पणी:}\\[0.3em]
वर्तमानसंस्करणं चौखम्बसंस्कृतशृङ्खलाप्रकाशनं आधृत्य विरचितम्, पण्डितसुधाकरद्विवेदिनः विस्तृतया हिन्दीटीकया सह। ब्रह्मगुप्तस्य नियमानां समानान्तरं विस्तारं कुर्वतः भास्कराचार्यस्य द्वितीयस्य लीलावत्याः सततं सन्दर्भाः कृताः।
}{%
\textbf{Note on Sources:}\\[0.3em]
The present edition is based on the Chowkhamba Sanskrit Series publication, with the extensive Hindi commentary (\textit{Hindī bhāṣā}) of Paṇḍita Sudhākara Dvivedin. Frequent references are made to the \textit{Līlāvatī} of Bhāskarācārya II, whose rules often parallel or expand upon those of Brahmagupta.
}

\vspace{1em}
\bilingual{%
ग्रन्थः गीताध्यायस्य संधिं समापयति---गोलज्यामितेः अध्यायस्य, यत् एतत् संस्करणस्य तृतीये अन्तिमे भागे अनुबन्धिष्यते।
}{%
The work concludes with the transition into the \textit{Gītādhyāya}, the chapter on spherical geometry, which will be continued in the third and final part of this edition.
}

\vfill
\begin{center}
\rule{0.5\textwidth}{0.4pt}\\[0.5em]
\textsc{End of Part II}\\[0.3em]
\dev{द्वितीयभागसमाप्तिः}\\[0.5em]
\rule{0.5\textwidth}{0.4pt}
\end{center}

\end{document}



% ============================================================
% Source: translations/en/indian/brahmagupta-brahmasphutasiddhanta-pt3-chunk1_bilingual.tex
% ============================================================

%% ========================================================================
%% Brāhmasphuṭasiddhānta of Brahmagupta --- Part 3, Chunk 1 (Pages 1677--1816)
%% BILINGUAL EDITION: Sanskrit (Devanagari) || English Translation with IAST
%% ========================================================================
\documentclass[11pt,a4paper]{article}

%% -------- Core Packages --------
\usepackage{fontspec}
\usepackage{polyglossia}
\usepackage{amsmath,amssymb}
\usepackage{geometry}
\usepackage{xcolor}
\usepackage{fancyhdr}
\usepackage{array}
\usepackage{hyperref}
\usepackage{tikz}
\usetikzlibrary{arrows.meta,calc,positioning}

%% -------- Page Geometry --------
\geometry{
  paperwidth=210mm,paperheight=297mm,
  left=16mm,right=16mm,top=24mm,bottom=24mm,
  headheight=14pt,footskip=28pt
}

%% -------- Language Setup --------
\setmainlanguage{english}
\setotherlanguage{sanskrit}

%% -------- Font Configuration --------
\setmainfont{Noto Serif}
\setsansfont{Noto Sans}
\setmonofont{DejaVu Sans Mono}

%% Devanagari Font (user-specified path)
\newfontfamily\devanagarifont[Script=Devanagari,Path=/usr/share/fonts/truetype/noto/]{NotoSerifDevanagari-Regular.ttf}

%% -------- Colors --------
\definecolor{titlecolor}{RGB}{45,55,72}
\definecolor{sectioncolor}{RGB}{55,65,81}
\definecolor{notecolor}{RGB}{120,53,15}
\definecolor{rulecolor}{RGB}{180,160,140}
\definecolor{versenumbercolor}{RGB}{139,69,19}

%% -------- Header/Footer --------
\pagestyle{fancy}
\fancyhf{}
\fancyhead[L]{\small\textit{Br\=ahmasphu\d tasiddh\=anta --- Part III, Chunk 1 (pp. 1677--1816)}}
\fancyhead[R]{\small\thepage}
\renewcommand{\headrulewidth}{0.4pt}
\renewcommand{\footrulewidth}{0pt}

%% -------- Custom Commands --------
\newcommand{\longline}{\noindent\rule{\textwidth}{0.4pt}}
\newcommand{\shortline}{\noindent\rule{0.5\textwidth}{0.4pt}\par}
\newcommand{\cnote}[1]{\textcolor{notecolor}{\textit{#1}}}

%% Devanagari helpers
\long\def\dev#1{{\devanagarifont #1}}
\newcommand{\devbf}[1]{{\devanagarifont\bfseries #1}}

%% Verse number marker
\newcommand{\verseno}[1]{\textcolor{versenumbercolor}{$\Vert$\dev{#1}$\Vert$}}

%% -------- Parallel Column Helpers --------
\newcommand{\bilingual}[2]{%
  \noindent%
  \begin{minipage}[t]{0.47\textwidth}#1\end{minipage}%
  \hfill\vrule\hfill%
  \begin{minipage}[t]{0.47\textwidth}#2\end{minipage}%
  \par\vspace{0.8em}%
}

\newcommand{\bilingualShloka}[2]{%
  \noindent%
  \begin{minipage}[t]{0.47\textwidth}
    \small\dev{#1}
  \end{minipage}%
  \hfill\vrule\hfill%
  \begin{minipage}[t]{0.47\textwidth}
    \small\textit{#2}
  \end{minipage}%
  \par\vspace{0.8em}%
}

\newcommand{\colheader}[2]{%
  \noindent%
  \begin{minipage}[t]{0.47\textwidth}\centering\bfseries\color{sectioncolor}#1\end{minipage}%
  \hfill\vrule\hfill%
  \begin{minipage}[t]{0.47\textwidth}\centering\bfseries\color{sectioncolor}#2\end{minipage}%
  \par\vspace{0.5em}%
}

%% -------- Hyperref Setup --------
\hypersetup{
  colorlinks=true,
  linkcolor=sectioncolor,
  citecolor=sectioncolor,
  urlcolor=sectioncolor,
  pdfencoding=auto
}

%% ========================================================================
\begin{document}

%% ========================================================================
%% TITLE PAGE
%% ========================================================================
\begin{titlepage}
\centering
\vspace*{1.5cm}

{\color{titlecolor}\Huge\devbf{ब्राह्मस्फुटसिद्धान्त}\par}
\vspace{0.4cm}
{\color{sectioncolor}\LARGE Br\=ahmasphu\d tasiddh\=anta\par}
\vspace{1.2cm}

{\Large\textit{of}\par}
\vspace{0.4cm}
{\color{titlecolor}\LARGE\devbf{ब्रह्मगुप्त}\par}
\vspace{0.3cm}
{\Large Brahmagupta\par}
\vspace{0.3cm}
{\large (c.\,598--668 CE)\par}

\vspace{1.5cm}
\shortline
\vspace{0.8cm}

{\Large\bfseries Volume III --- Chunk I\par}
\vspace{0.4cm}
{\large Pages 1677--1816\par}
\vspace{0.3cm}
{\normalsize comprising Chapters XII--XXI\par}
\vspace{0.4cm}
{\large (Final Volume: Concluding Chapters \& Appendices)\par}

\vspace{1.2cm}
\shortline
\vspace{0.8cm}

{\normalsize\textbf{BILINGUAL EDITION}\par}
\vspace{0.3cm}
{\small Sanskrit (Devanagari) $\parallel$ English Translation with IAST\par}

\vspace{1cm}
{\normalsize Edited with Sanskrit Commentary, Hindi Translation,\par}
\vspace{0.2cm}
{\normalsize Mathematical Proofs and Astronomical Notes\par}
\vspace{0.4cm}
{\large by\par}
\vspace{0.3cm}
{\large\bfseries Dr.~Ram~Swarup~Sharma\par}

\vspace{1.2cm}
{\small Published by\\
\textbf{Indian Institute of Astronomical and Sanskrit Research}\\
New Delhi, India\par}

\vfill
{\footnotesize Electronic typesetting with XeLaTeX\par}
\end{titlepage}

%% ========================================================================
%% COLUMN HEADERS
%% ========================================================================
\newpage
\colheader{संस्कृत (Sanskrit)}{English Translation}
\longline
\vspace{0.5em}

%% ========================================================================
%% PREFACE
%% ========================================================================
\section*{Preface to the Final Volume}
\addcontentsline{toc}{section}{Preface to the Final Volume}

\bilingual{%
  \dev{अस्य अन्तिमस्य खण्डस्य प्रस्तावना --- एतत् अन्तिमं भागं ब्रह्मस्फुटसिद्धान्तस्य पूर्णं करोति। पृष्ठानि १६७७--१८१६ ब्रह्मगुप्त-महोदयस्य अन्तिमाध्यायान् समाविशन्ति, यथाम्: गणिताध्यायस्य समाप्तिः, विकलवर्गानयनम्, कालज्ञानम्, चन्द्रग्रहणम्, मन्दस्फुटम्, शीघ्रफलम्, भोगकला, प्राणकला, च समाप्तिः।}%
}{%
  This final volume of the \textit{Br\=ahmasphu\d tasiddh\=anta} completes the monumental work begun in the first two volumes. Pages 1677--1816 cover the concluding chapters of Brahmagupta's masterpiece, including the completion of the Ga\d nit\=adhy\=aya, extraction of roots of non-square numbers, time computations, lunar eclipses, equations of centre and anomaly, daily motion, and the final colophons. The text preserves Brahmagupta's original verses in authentic Sanskrit form, accompanied by detailed commentary and mathematical proofs (\textit{upapatti}).%
}

\vspace{0.5em}
\longline
\vspace{0.5em}

%% ========================================================================
%% CHAPTER: Circumference and Shadow
%% ========================================================================
\section{Conclusion of Arithmetic --- Circumference \& Shadow}

\bilingual{%
  \dev{अस्मिन् खण्डे गणिताध्यायस्य अन्तिमश्लोकाः आरभ्यन्ते, ये परिधि-अनुपपत्ति-छाया-व्यवहारान् प्रतिपादयन्ति।}%
}{%
  The present section begins with the final verses of the twelfth chapter (\textit{Ga\d nit\=adhy\=aya}), treating the calculation of circumference (\textit{paridhi}), proportion (\textit{anupapatti}), and the computation of shadows (\textit{ch\=ay\=a vyavah\=ara}).%
}

\subsection{Circumference and Proportion}

\bilingualShloka{%
  परिधौ क्रमशः प्राग्वत् स्वगुणात् भवेत् फलम्~|\\
  श्रीपत्युक्तमिदं चाचार्यैरनुक्रमेणावबोधित~||१७||%
}{%
  paridhau krama\'sa\d h pr\=agvat svagu\n\=at bhavet phalam~|\\
  \'sr\=ipatyuktamida\d m c\=ac\=aryairanukrame\n\=avabodhita~||17||%
}

\bilingual{%
  \dev{(बीएसएस् १२.३७ / खण्ड~२, पृ. ८९१) एतत् श्लोकं परिधिः व्यासस्य उक्तगुणकैः क्रमशो गुणनया प्राप्यते इति वदति, यत् पूर्वाचार्यैः श्रीपतिः-आदिभिः उपदिष्टम्। ब्रह्मगुप्त-भास्कराचार्य-लीलावतीभिः अपि स्वस्वग्रन्थेषु एतत् उपदिष्टम्।}%
}{%
  \cnote{(BSS 12.37 / Vol.~II, p.~891)} This verse states that the circumference is obtained by multiplying the diameter successively by the prescribed factors, as taught by previous teachers including \'Sr\=ipati. Brahmagupta, Bh\=askar\=ac\=arya, and L\=il\=avat\=i have all taught this in their respective works.%
}

\bilingualShloka{%
  द्विवेद सचित्रभागे कनिष्ठान्तु परिधिः फलम्~|\\
  मिथ्यन्तबोधकोऽणस्थराशिः स्वगुणाभिजं तत्~||१७१||%
}{%
  dviveda sacitrabh\=age kani\s\d th\=antu paridhi\d h phalam~|\\
  mithyantabodako'\nasthar\=\=a\'si\d h svagu\n\=abhija\d m tat~||171||%
}

\bilingual{%
  \dev{यथार्थः परिधिः द्विवेदि-सचित्रभाग-कल्पितैः प्राप्यते; अन्यथा स्थूलाप्रoksणात् प्राप्तः परिधिः मिथ्या भवति।}%
}{%
  \verseno{171} The true circumference is that obtained by the rules taught by the two Dvivedis; the circumference otherwise (from crude approximation) is erroneous.%
}

\subsection{Mathematical Proof}

\bilingual{%
  \dev{अनुपपत्तिः --- मिथ्यान्तः संलग्न धान्यराशिः परिधिवास्तव परिधेर्बंधुतुल्यः~| क्रान्तः कोण स्थितस्य धान्यराशिः परिधिवास्तवपरिधिं चतुर्थांशतुल्यः~| बाहुकोणसंलग्न धान्यराशिः परिधिवास्तव परिधेः पादेनैक तुल्यस्तेनोत्परिधित्रयमानं द्विवेद-सचित्रभागेकेन क्रमशो गुणितां सघातं परिधिं भवेत्~| ततः पूर्वं वेधफलं ततु द्विवेद सचित्रभागेकेन क्रमशो भक्तं तदा वास्तवं तद्फलं भवतीत्येतावताऽऽचार्योक्त-मुपपन्नम्~||१७१||}%
}{%
  \textit{Upapatti} (Proof): The incorrect circumference associated with a quantity of grains is proportional to the true circumference. The inclined-angle quantity of grains is one-quarter of the true circumference; the base-angle quantity is one-sixth. The circumference- triple measure, multiplied successively by the Dvivedi-\-sacitrabh\=aga factors, yields the true circumference. The previous penetration result, divided successively by the Dvivedi-\-sacitrabh\=aga factors, gives the true result --- thus the teacher's statement is proved.~||171||%
}

\subsection{End of Arithmetic Chapter}

\bilingual{%
  \dev{इति राशि व्यवहारः --- अब भत्ति के अन्तर्गत भत्ति के क्रान्तः कोणरूप तथा बाहुकोणरूप धान्यराशि प्रमाणा-नयन के लिये कहते हैं~|}%
}{%
  Here ends the \textit{R\=a\'si Vyavah\=ara} (Arithmetic Operations). Now, under division, the quantities of grains corresponding to the inclined angle and base angle are stated for the derivation of measures.%
}

\subsection{Shadow Calculations}

\bilingual{%
  \dev{छाया व्यवहार --- बीएसएस् १२.४१--४२: तज्जादौ छायात् इष्टकालज्ञानार्थमिष्टकालतच्छायानयनार्थं चाहु~| छायानरसंक्रमहृतं शुद्धं प्रागपरयोर्दिगतदोषं~| दिनगतदोषांशहृतं धु दलं छाया नरखेखं~||४२||}%
}{%
  BSS 12.41--42: First, for finding the time from a given shadow, or the shadow at a given time. The \textit{nara} (sine of latitude) multiplied by the shadow, divided by the \textit{\'sa\~nku} (gnomon), gives the \textit{\'suddha} (mean) value; considering the eastern or western direction, this is divided by the \textit{dinagata-do\d s\=\=a\'msa} (ascensional difference) to obtain the accurate shadow.~||42||%
}

\bilingual{%
  \dev{सू. भा.---नरः शङ्कुः रभिमष्टः~| छायाया यो नरः शङ्कुं भागः स संक्रमे यस्तेन}\\
  \dev{शुद्धं हृतं लब्धं प्राक् पश्चिमकपाले दुग्गतं पश्चिमकपाले दुग्गतं भवति~|}%
}{%
  \textit{S\=utra-bh\=a\d ya:} The \textit{nara} is the \textit{\'sa\~nku} multiplied by eight --- the \textit{nara} of the shadow divided by the \textit{\'sa\~nku} is that in the \textit{sa\~nkrama}; this is the \textit{\'suddha}; divided by the eastern or western semi-diurnal arc, it becomes the corrected value.%
}

\bilingual{%
  \dev{छाया व्यवहारस्य विस्तृत उपपत्तिः अनुवर्तते। नरः अक्षज्या अस्ति, शङ्कुः शङ्कुः (१२ अङ्गुलाः)। छाया एतेषां परस्परसम्बन्धात् गण्यते।}%
}{%
  The \textit{Ch\=ay\=a Vyavah\=ara} (shadow computations) section continues with detailed derivations. The \textit{nara} is the sine of latitude, and the \textit{\'sa\~nku} is the gnomon (12 a\~ngulas). The shadow is computed from the relationship between these quantities.%
}

\vspace{0.5em}
\longline
\vspace{0.5em}

\subsection{Geometric Construction of Shadows}

\bilingual{%
  \dev{सम्बद्धचित्राणि ज्यामितीयसम्बन्धान् दर्शयन्ति~| शङ्कुः (उन्नतदण्डः) भूतले छायां प्रकाशस्य किरणेन सह कोणं निर्माति~|}%
}{%
  The accompanying diagrams illustrate the geometric relationships: the gnomon (vertical pole) casts a shadow on the ground, forming an angle with the sun's ray. The \textit{karna} (hypotenuse) equals the distance from the lamp's height. Similar triangles yield the proportions.%
}

\begin{center}
\begin{tikzpicture}[scale=1.0,>=Stealth]
  \draw[thick] (0,0) -- (0,3) node[midway,left] {\dev{शङ्कु}};
  \fill (0,3) circle (2pt) node[above] {\dev{क}};
  \fill (0,0) circle (2pt) node[below] {\dev{श}};
  \draw[thick,->] (0,0) -- (4,0) node[midway,below] {\dev{छाया}};
  \fill (4,0) circle (2pt) node[right] {\dev{प}};
  \draw[thick] (0,3) -- (4,0);
  \draw[thick] (1.5,0) -- (1.5,1.8);
  \fill (1.5,1.8) circle (1.5pt);
  \fill (1.5,0) circle (1.5pt);
  \draw[thick] (-0.5,0) -- (4.5,0);
  \draw (0.8,0) arc (0:-37:0.8);
  \node at (1.1,-0.4) {\dev{ज्र}};
  \node at (2.5,2) {\dev{कर्ण = दीपोच्चय}};
\end{tikzpicture}
\end{center}

\bilingual{%
  \dev{छाया, शान = भूः~| कर्ण = दीपोच्चयम्~| कर्ण, भुज त्रिभुजयोः साजात्यादनुपातः}\\
  \dev{दीघ् = दीपोच्चयम् तथा कर्ण, त्रिभुजयोः साजात्यादनुपातः}\\
  \dev{=भूः-भू = रन = छायाप्रान्तर = छाया दीघ्-छाया दीघ् = दीघ्(छाया-छाया) = दीघ् छायान्तर}\\
  \dev{छेदगमेन छायाप्रान्तर दीघ् छायान्तर, पक्षो छायान्तर भक्तो तदा}\\
  \dev{छायाप्रान्तर दीघ् छायान्तर भ्रनेन भू स्वरूपे उत्थापनेन भू = छ.प्रा. दीघ्}\\
  \dev{= छाया छायाप्रान्तर दीघ् = छाया छायाप्रान्तर, तथा भू = छाया दीघ्}\\
  \dev{= छाया छायाप्रान्तर दीघ् = छाया छायाप्रान्तर, तथा भुज, कर्ण त्रिभुजयोः}\\
  \dev{साजात्यात् १२ भू/छाया = दीपोच्चयम्~| एवं त्रपन, कर्ण त्रिभुजयोः साजात्या-}\\
  \dev{दनुपातेन १२ भू/छाया = दीपोच्चयम्~| एतावत्~स्साचार्योक्तमुपपन्नम्~|}%
}{%
  \textit{Shadow} (Ch\=ay\=a), \textit{\'s\=ana} = base. \textit{Karna} = lamp-height. From similarity of triangles, the \textit{karna}--base proportion gives: base = \textit{ch\=ay\=a-pr\=antara}; then by proportional reasoning, 12 $\times$ base/shadow = lamp-height. Similarly for the other triangle pair, the same result follows. Thus the teacher's statement is proved. (BSS 12.41--42 commentary)%
}

\newpage
%% ========================================================================
%% CHAPTER: Algebraic Operations
%% ========================================================================
\section{Algebraic Operations --- \dev{विकलवर्गानयन}}

\bilingual{%
  \dev{बीएसएस् १२.४६: इदानीं सकल विकलवर्गानयनार्थमाह~|}\\
  \dev{स्वविकलपदं वर्गांशगुणाः सकलस्त्रिशद्द्वैतो विकलवर्गः~|}\\
  \dev{प्रक्षेप्यः सकलकृतौ वर्गधनो द्वित्रितुल्यबधो~||४६||}%
}{%
  BSS 12.46: Now (the author) teaches the method of finding the square root of a non-square number. The square of the \textit{sakala} (integer part) is subtracted from the number; the remainder, multiplied by the square of the denominator, becomes the \textit{vikalavarga} (square of the non-square part). This is to be added to the square of the \textit{sakala}, and the result gives the square root approximately.~||46||%
}

\bilingual{%
  \dev{(बीएसएस् १२.४६) यदि संख्यायाः कलाः विकलाः च इति द्वयं वर्तते, तस्य वर्गार्थं कलाराशिः सकलसंज्ञो विकलाराशिश्च स्वविकलः तत्कलापश्चात् शास्त्र कथये~|}%
}{%
  \cnote{(BSS 12.46)} Where a quantity has both integer (\textit{kala}) and fractional (\textit{vikala}) parts, the integer part is called \textit{sakala} and the fractional part \textit{vikala}; these are then processed by the rule to obtain the square root.%
}

\subsubsection*{Upapatti (Mathematical Proof)}

\bilingual{%
  \dev{उपपत्तिः --- कल्पयते कलिमश्रवि राशौ कलाः = क~| विकलाः = वि~| तदा कलात्मकः स राशिः = क + वि/६०~|}%
}{%
  \textit{Upapatti:} Let the integer part be $k$ and the \textit{vikala} part $vi$. Then the quantity is $k + \frac{vi}{60}$. The square is:
  \begin{align*}
    \textrm{वर्गः} &= k^2 + \frac{2k \cdot vi}{60} + \frac{vi^2}{60^2}
  \end{align*}%
}

\bilingual{%
  \dev{वि. भा.---यदिमन् राशौ कला-विकला चेति द्वयं वर्तते तस्य वर्गं करणार्थं}\\
  \dev{कलाराशिः सकल संज्ञो विकलाराशिश्च स्वविकलस्तत्कलापश्चात् शास्त्र कथये~|}%
}{%
  \textit{V\=tti-bh\=a\d ya:} When a quantity has both \textit{kala} and \textit{vikala} parts, to find its square, the integer part is called \textit{sakala} and the fractional part \textit{vikala}; then the rule states the method. Note: ``\dev{वर्गः}'' means `square' in Sanskrit. (BSS 12.46 commentary)%
}

\bilingualShloka{%
  छेदेनेष्टयुतोनेनार्धं भाज्यपदष्टष्टिगुणं स्यात्~|\\
  प्रकृतिरष्टच्छेदहृतं लब्धा युतं हैनिकमनघं~||४७||%
}{%
  chedene\s\d tayutone-n\=ardha\d m bh\=ajayapa-da\s\d ta\s\d\d tiguna\d m sy\=at~|\\
  prak\d rtira\s\d tachchhedah\d rta\d m labdh\=a yuta\d m hainikamanagha\d m~||47||%
}

\bilingual{%
  \dev{(बीएसएस् १२.४७) एतत् श्लोकं भिन्नैः सह भाजनस्य विधिं ददाति~| भाज्यं, छेदेन सह युक्तं (भिन्नयोजनविधिना), भाजकस्य अर्धेन गुणितं, ततः प्रकृत्या युक्तं, शुद्धं लब्धं ददाति~|}%
}{%
  \cnote{(BSS 12.47)} This verse gives the method for division with fractions. The dividend, combined with the divisor (in the manner of adding fractions), is divided by half the divisor multiplied by the denominator of the divisor; the result, added to the \textit{prak\d rti}, gives the quotient. (BSS 12.47 commentary)%
}

\bilingualShloka{%
  इदानीं भागहारे विरोपमाह~|\\
  छेदेनेष्टयुतोनेनार्धं भाजयपदष्टष्टिगुणं स्यात्~|\\
  प्रकृतिरष्टच्छेदहृतं लब्धा युतं हैनिकमनघं~||४९||%
}{%
  id\=an\=\=im bh\=agah\=are viropam\=aha~|\\
  chedene\s\d tayutone-n\=ardha\d m bh\=ajayapa-da\s\d ta\s\d\d tiguna\d m sy\=at~|\\
  prak\d rtira\s\d tachchhedah\d rta\d m labdh\=a yuta\d m hainikamanagha\d m~||49||%
}

\bilingual{%
  \dev{(बीएसएस् १२.४९) भागहारस्य विभागस्य अयं खण्डः~| छेदैः सह भिन्नराशीनां व्यवहारः अत्र उपदिश्यते~|}%
}{%
  \cnote{(BSS 12.49)} The section on division with fractions (\textit{bh\=agah\=ara}) is explained. The rule involves the manipulation of fractional quantities through their denominators. (BSS 12.49 commentary)%
}

\newpage
%% ========================================================================
%% CHAPTER: Astronomical Problems
%% ========================================================================
\section{Astronomical Problems --- \dev{कालज्ञान}}

\bilingual{%
  \dev{अनुवर्तमानाः खण्डाः ज्योतिषीयगणनानि समाविशन्ति --- कालनिर्णयम्, ग्रहस्थितिगणनम्, विविधज्योतिषीयमानानि च।}%
}{%
  The following sections deal with astronomical computations, including the determination of time, planetary positions, and the computation of various astronomical quantities.%
}

\subsection{Mixed Astronomical Problems}

\bilingualShloka{%
  इदानीं मध्यान् प्रश्नानाह~|\\
  व्यतिपातं वैधृतिबृहस्पतिवर्षयोश्चर्चनीचरिवबस्तीन~|\\
  द्विप्रहयोगांश युगे यो वेति स कालतत्वज्ञः~||६||%
}{%
  id\=an\=\=im madhy\=an pra\'sn\=an\=aha~|\\
  vyat\=ip\=ata\d m vaidh\d rtib\d rhaspativar\d sayo\'scarcan\=icarivast\=ina~|\\
  dviprahayog\=a\'m\'sa yuge yo veti sa k\=alatatvaj\~na\d h~||6||%
}

\bilingual{%
  \dev{(बीएसएस् १३.६) यो व्यतिपातं वेति~| वैधृतिं वेति बृहस्पतिवर्षं वेति,}\\
  \dev{स्वोच्चनीचमार्गगान् वेति~| युगे द्विप्रहयोगांश इत्यादि~|}%
}{%
  \cnote{(BSS 13.6)} He who knows the \textit{vyat\=ip\=ata}, \textit{vaidh\d rti}, \textit{b\d rhaspati-var\d sa}, and the motion of the nodes in a \textit{yuga} --- he is the knower of the essence of time. (BSS 13.6 commentary)%
}

\bilingualShloka{%
  सावनमासावधिप्रहरेक्काऽनिष्टमध्यसंयोगात्~|\\
  इष्टद्वयु संयु तोननिष्टानां यो वेति ग्रह्यकः सः~||१०||%
}{%
  s\=avanam\=as\=avadhipraharekk\='ani\s\d tamadhyasa\d myog\=at~|\\
  i\s\d tadvayu sa\d myu tonani\s\d t\=an\=\=a\d m yo veti grahyaka\d h sa\d h~||10||%
}

\bilingual{%
  \dev{(बीएसएस् १३.१०) यः अवमानां (लुप्ततिथीनां) प्रघराणां (अधिकदिनानां) च योगं, इष्टकाले मध्यस्पष्टैः सह संयुक्तं वेति, सः ग्रह्यकज्ञः।}%
}{%
  \cnote{(BSS 13.10)} He who knows the sum of the \textit{avama} (omitted tithis) and the \textit{praghara} (excess days), combined with the mean (longitudes) at the desired time --- he is the knower of the \textit{grahyaka} (the quantity to be known). (BSS 13.10 commentary)%
}

\subsection{Third Set of Problems}

\bilingualShloka{%
  इदानीं तृतीय प्रश्नसमूह (छायादीपरितोदिवस वारे वेतीत्यस्य) उत्तरमाह~|\\
  कल्पर्क्षदिनसप्तकबधात् कल्पगताहर्गणानां तच्छेषात्~|\\
  सप्तहृताद्दिनवारः शेषः साय्यादिको भवति~||४०||%
}{%
  id\=an\=\=im t\d rt\=iya pra\'s nasam\=uha (ch\=ay\=ad\=\=iparitodivasa v\=are vet\=\=ityasya) uttaram\=aha~|\\
  kalpark\=sadinasaptakabadh\=at kalpagat\=aharga\n\=an\=\=a\d m tacche\s\d\=at~|\\
  saptah\d r\t\=addinav\=ara\d h \'se\s\d a\d h s\=ayy\=adiko bhavati~||40||%
}

\bilingual{%
  \dev{(बीएसएस् १३.४०) कल्पस्य गताहर्गणानां सप्तकेन गुणनात्, ततः सप्तभक्तात् शेषात् दिनवारः साय्यादिकः प्राप्यते।}%
}{%
  \cnote{(BSS 13.40)} From the product of the elapsed days of the \textit{kalpa} and seven, divided by seven --- the remainder gives the day of the week (counting from Sunday).%
}

\bilingual{%
  \dev{उपपत्तिः --- कल्पर्क्षदिनसप्तकबधात् कल्पकृत् दिनसप्तधातात् कि विशिष्टात्}\\
  \dev{कल्पगताहर्गणेनैन कावः शेषस्तस्मात् सप्तहृतात् शेषः साय्यादिको}\\
  \dev{दिनवारो भवति~|}%
}{%
  \textit{Upapatti:} From the \textit{kalpa} elapsed days multiplied by seven, the remainder from division by seven gives the day of the week starting from Sunday. (BSS 13.40 commentary)%
}

\subsection{Lunar Calculations}

\bilingualShloka{%
  यदीष्ट ग्रहयुग्मभगणां सरोभ मगणाद्यष्टमां लभ्यते\\
  तदा युगीयपातमगणां चन्द्रभगणायोगयोगि किं फलं\\
  भगणादिकमिष्टहृत् राश्यादिमध्यमसपातचन्द्रं यथागतं\\
  देशान्तरमन्दफलादीन् यथागतां सरोभ सपातचन्द्रो\\
  भवेत्~||३१||%
}{%
  yad\=i\s\d ta grahayugmabhaga\n\=a\d m sarobha maga\n\=\=ady\a\s\d tam\=\=a\d m labhyate\\
  tad\=a yug\=\=iyap\=atamaga\n\=\=a\d m candrabhaga\n\=\=ayogayogi ki\d m phala\d m\\
  bhaga\n\=\=adikami\s\d tah\d rt r\=a\'sy\=adimadhyamasap\=atacandra\d m yath\=agata\d m\\
  de\'s\=antaramandaphal\=ad\=\=in yath\=agat\=\=a\d m sarobha sap\=atacandro\\
  bhavet~||31||%
}

\newpage
%% ========================================================================
%% CHAPTER: Planetary Theory
%% ========================================================================
\section{Planetary Theory --- \dev{मन्दस्फुट} \& \dev{शीघ्रफल}}

\subsection{Epicyclic Model}

\bilingual{%
  \dev{ग्रहस्थितिगणनायां मन्दस्फुटस्य (केन्द्रसमीकरणस्य) शीघ्रफलस्य (द्वितीयवक्रीभावसमीकरणस्य) च गणना आवश्यकी~| ब्रह्मगुप्तः अधोलिखितं विधिं ददाति~|}%
}{%
  The computation of planetary positions involves the \textit{mandasphu\d ta} (equation of centre) and \textit{\'s\=ighraphala} (equation of anomaly, or second inequality). Brahmagupta gives the following method:%
}

\bilingualShloka{%
  मन्दस्फुटार्केन्द्रांशोनायां ग्रहस्यन्दं केन्द्रम्~|\\
  एव यदि तत् स्पष्टं त्रिकोणं तदानीं विशेषं नोद्येत्~|\\
  तद्यात् स्पष्टं त्रिकोणं तस्य स्पष्टं द्वितीयं केन्द्रं\\
  क्रमेण केन्द्रान्तरं फलमिति~| वा ज्ञानमन्दं स्पष्टं\\
  प्राप्यच्छी~स्साचार्य्योदिना शीघ्र-कर्त्तव्यं कार्यमिति सर्वं शुद्धम्~||३७८-५८१||%
}{%
  mandasphu\d t\=arkendr\=a\'m\'son\=ay\=a\d m grahasyanda\d m kendram~|\\
  eva yadi tat spa\s\d ta\d m triko\na\d m tad\=an\=\=i\d m vi\'se\s\d a\d m nodyet~|\\
  tady\=at spa\s\d ta\d m triko\na\d m tasya spa\s\d ta\d m dvit\=iya\d m kendra\d m\\
  krame\n\=a kendr\=antara\d m phalamiti~| v\=a j\~n\=anamanda\d m spa\s\d ta\d m\\
  pr\=apyacch\=\=is\=s\=ac\=aryyodin\=a \'s\=\=ighra-karttavya\d m k\=aryamiti sarva\d m \'suddham~||378--581||%
}

\bilingual{%
  \dev{वि. भा.---उदयमार्ग (विज्ञानान्) भुजयया गुणिगतं कर्णहृतं यत्फलं तस्य}\\
  \dev{चापं स्पष्टद्वितीयोनयोरन्तरं भवेत् तत् स्पष्टमिति~||३७९-५८१||}%
}{%
  \textit{V\=tti-bh\=a\=sya:} The \textit{phala} obtained by multiplying by the \textit{bhuja} (sine of anomaly) and dividing by the \textit{karna} --- its arc gives the difference between the first and second true positions; that is the true position. (BSS 378--581 commentary)%
}

\subsection{Proof of Epicyclic Correction}

\bilingual{%
  \dev{उपपत्तिः --- ज = गोच्चेभम्~| म = मन्दस्पष्टग्रहः~| मं३ = शीघ्रकेन्द्रम्~| मर = शीघ्र-}\\
  \dev{केन्द्रज्या = गोच्चज्या~| स्प. = स्पष्टग्रहः~| स्प.३ = स्पष्ट केन्द्रम्~|}%
}{%
  \textit{Upapatti:} \textit{Ja} = geocentric planet; \textit{Ma} = mean planet via manda; \textit{Ma.3} = \textit{\'s\=ighrakendra}; \textit{Mar} = \textit{\'s\=ighrakendra-jy\=a} = geocentric sine; \textit{Sp.} = true planet; \textit{Sp.3} = true \textit{kendra}. (BSS \textit{upapatti})%
}

\begin{center}
\begin{tikzpicture}[scale=1.4,>=Stealth]
  \fill (0,0) circle (3pt) node[below left] {\dev{पृथ्वी}};
  \draw[thick] (0,0) circle (2);
  \node at (2.3,0) [right] {\dev{मन्दकर्ण (विकेन्द्र)}};
  \coordinate (M) at (1.2,1.6);
  \fill (M) circle (2pt) node[above right] {\dev{म (मन्दग्रह)}};
  \draw[thick] (M) circle (0.6);
  \node at (M) [below=0.8cm] {\dev{मन्दवृत्त}};
  \coordinate (T) at (1.6,1.3);
  \fill (T) circle (2pt) node[right] {\dev{स्पष्टग्रह}};
  \draw[thick,->] (0,0) -- (M);
  \draw[thick,dashed] (M) -- (T);
  \draw[thick,->] (0,0) -- (T);
  \coordinate (U) at (1.8,0.9);
  \fill (U) circle (1.5pt) node[right] {\dev{(उ) उच्च}};
  \draw[thick,dotted,->] (0,0) -- (2.5,1.25);
  \node at (0.6,1.1) {\dev{मन्दकर्ण}};
  \node at (1.55,1.6) {\dev{फलज्या}};
  \draw[thin,dashed] (0,0) circle (2.5);
  \node at (2.5,0) [right] {\dev{कक्षावृत्त}};
  \coordinate (MP) at (-1.5,1.3);
  \fill (MP) circle (1.5pt) node[left] {\dev{मध्यमग्रह}};
  \draw[thin,dashed] (0,0) -- (MP);
\end{tikzpicture}
\end{center}

\bilingual{%
  \dev{भूमं गोचरग्रहः~| स्पष्ट = स्पष्ट केन्द्रज्या~| म.३. म. = मू केन्द्रम्~|}\\
  \dev{मू.स्पष्ट = शीघ्रफल~| एवं त्रिज्या~| नदा मू.म. स्पष्ट त्रिभुजयोः}\\
  \dev{साजात्यादनुपातः गोच्चज्या/शीघ्रकर्ण = स्पष्टकेन्द्रज्या,}\\
  \dev{अथवाऽन्योनम् = स्पष्टकेन्द्रम् = स्पष्टग्रहोच्चयोरन्तरं}\\
  \dev{मन्दकर्णादिया शीघ्रकर्णस्थाने मन्दस्पष्ट}\\
  \dev{केन्द्रज्या समागम्यति, चापकल्पेन मन्दस्पष्ट}\\
  \dev{केन्द्रराशिः = मन्दस्पष्टदोषयोरन्तरं~|}%
}{%
  \textit{Bhum} = geocentric planet. \textit{Spa\s\d ta} = true \textit{kendra-jy\=a}. \textit{Ma.3. Ma.} = mean \textit{kendra}. \textit{M\=u.spa\s\d ta} = \textit{\'s\=ighraphala}. From similarity of triangles: geocentric sine/\textit{\'s\=ighrakarna} = true \textit{kendra-jy\=a}. Alternatively, true \textit{kendra} = difference between true planet and apogee. The manda-true \textit{kendra-jy\=a} meets the epicycle; by arc-correspondence, manda-true \textit{kendra-r\=a\'si} = difference between manda-true and equation. (BSS \textit{upapatti})%
}

\bilingual{%
  \dev{मन्दस्फुटस्य गणनक्रमः एवं अस्ति: (१) मन्दकेन्द्रस्य (मध्यमग्रहस्य उच्चस्य च अन्तरं) गणना, (२) ज्यापिठकात् मन्दफलस्य (केन्द्रसमीकरणस्य) अवलोकनम्, (३) मध्यमग्रहे मन्दफलस्य प्रयोगः मन्दस्फुटस्य (प्रथमस्पष्टस्य) प्राप्त्यर्थम्, (४) मन्दस्फुटात् शीघ्रकेन्द्रस्य गणना, (५) शीघ्रफलस्य अवलोकनम्, (६) शीघ्रफलस्य प्रयोगः सत्यग्रहस्य प्राप्त्यर्थम्~|}%
}{%
  The \textit{mandasphu\d ta} (true planet via equation of centre) is computed as:
  \begin{enumerate}
    \item Compute the \textit{mandakendra} (mean anomaly): difference between mean longitude and apogee (\dev{मन्दोच्च})
    \item Find the \textit{mandaphala} (equation of centre) from the sine table
    \item Apply to the mean longitude to get \textit{mandasphu\d ta}
    \item Compute \textit{\'s\=ighrakendra} from \textit{mandasphu\d ta}
    \item Find \textit{\'s\=ighraphala} and apply to get the true longitude
  \end{enumerate}%
}

\newpage
%% ========================================================================
%% CHAPTER: Advanced Astronomical Computations
%% ========================================================================
\section{Advanced Astronomical Computations --- \dev{भोगकला} \& \dev{प्राणकला}}

\bilingualShloka{%
  इष्टज्यां संगुणीताः पश्चकर्मणां कर्णच्चन्द्रम्~|\\
  इष्टज्यापाद्युत्थ्यासर्वं विभाजिता लब्धम्~||२५||\\
  नवति छते प्रोह पदे नवते मनोरथं शेषमगणत्कला~|\\
  एवं धनुरिष्टदया भवति ज्यया विना ज्यार्धिः~||२६||%
}{%
  i\s\d tajy\=\=a\d m sa\d gu\n\=\=it\=a\d h pa\'scakarma\n\=\=a\d m kar\d nac candram~|\\
  i\s\d tajy\=\=ap\=adyutthy\=asarva\d m vibh\=ajit\=a labdham~||25||\\
  navati chate proha pade navate manoratha\d m \'se\s amaga\nat kal\=a~|\\
  eva\d m dhanuri\s\d taday\=a bhavati jyay\=a vin\=a jy\=ardhi\d h~||26||%
}

\bilingual{%
  \dev{(बीएसएस् २१.२५--२६) इष्टज्या पञ्चकर्मभिः संगुणिता कर्णेन विभक्ता; शेषात् नवत्यधिकात् पदात् मनोरथः (इष्टचापः) प्राप्यते~| एतत् ज्यार्धेः विना चापस्य अवलोकनम् अस्ति~|}%
}{%
  \cnote{(BSS 21.25--26)} The \textit{i\s\d tajy\=a} (desired sine) is multiplied by the results of the five operations and divided by the radius; the remainder, when 90 is subtracted from the \textit{nava} position, gives the \textit{\'se\d sa}; the \textit{manoratha} (desired arc) is obtained. The arc is thus found from the sine without the \textit{jy\=ardha} (half-chord).%
}

\bilingual{%
  \dev{सू. भा.---इष्ट ज्यायाः पादद्वयपरिलेखेन युक्त व्यासार्धं तद्द्वारा तेन}\\
  \dev{विभाजिताः~| शेष स्पष्टम्~|}%
}{%
  \textit{S\=utra-bh\=a\d ya:} The desired sine is divided by the semi-diameter, together with the two-foot arc-difference; the remainder is clear. (BSS 21.25--26 commentary)%
}

\subsubsection*{Upapatti --- Detailed Proof}

\bilingual{%
  \dev{उपपत्तिः --- पूर्वोक्तकोपपरित्य ज्या = (१८०-चा) चा.ज्र. / [१०१२५-(१८०-चा) चा/४]}\\
  \dev{या ज्या / [१०१२५-या/४] = ज्या/१०१२५-या/४}%
}{%
  \textit{Upapatti:} From previous proofs: $jy\=a = \frac{(180-ca) \cdot ca \cdot jra}{10125 - \frac{(180-ca)ca}{4}}$. Let $y\=a = (180 - ca) \cdot ca$. Then $jy\=a = \frac{jy\=a \cdot y\=a}{10125 - y\=a/4}$. (BSS 21.25--26 \textit{upapatti})%
}

\bilingual{%
  \dev{या = (१८०-चा) चा~| अतः शेषद्वारागमेन ज्या × ४०५०० = ज्या. या}\\
  \dev{= ४ ज्या. या, ततो या = ज्या × ४०५०० = १०१२५. ज्या}\\
  \dev{ज्या + ६ ज्या / ६}\\
  = ल = (१८०-चा) चा = १८० चा - चा\textsuperscript{२}\\
  समशोधनेन चा\textsuperscript{२} - १८० चा + ल = ०\\
  $\therefore$ चा = ९० $\pm$ $\sqrt{\textrm{९०}^{2} - ल}$\\
  आचार्योक्तमर्क चापं गृहीतं तेन चा = ९० - $\sqrt{\textrm{९०}^{2} - ल}$\\
  \dev{अत उपपन्नं सर्वम्~||२५-२६||}%
}{%
  Let $y\=a = (180 - ca) \cdot ca$. By algebraic manipulation: $jy\=a \times 40500 = 4 \cdot jy\=a \cdot y\=a$, so $y\=a = 10125 \cdot jy\=a / (jy\=a + 6)$. Let $l = (180 - ca)ca = 180ca - ca^2$. By equating: $ca^2 - 180ca + l = 0$. Therefore $ca = 90 \pm \sqrt{90^2 - l}$. Taking the smaller arc as the teacher taught: $ca = 90 - \sqrt{90^2 - l}$. Thus all is proved.~||25--26||%
}

\bilingualShloka{%
  षड्भिर्भोगकलानामेकयम् = ९ चग\\
  षड्भिर्भोगकलानामेकयम् = ३ चग\\
  पञ्चदशैर्भोगकलानामेकयम् = १५ चग = १५ चग\\
  सर्वयोगकलाः = २७ चग\\
  चक्रकलाभ्यः शुद्धा सर्वयोगकला जाता श्रभिज्ञुगकलास्तदिहनगतिः = चक्-२७ चग~|\\
  इय कल्पकृत् दिनगुणाः इष्टचक्रकलामका जाता~| कल्पे$^{3}$भीजनो\\
  भगणाः = क्षु --- २७ क्षभं~| शेषोपपत्ति मोर्तकृत् मुहनक्षत्रायन विभिन्नकृत्~||५०-५२||%
}{%
  \textit{Bhogakal\=a} (daily motion in arc-minutes): For six \textit{bhogakal\=as}, one = 9 \textit{caga}; for another six, one = 3 \textit{caga}; for fifteen, one = 15 \textit{caga}. Total \textit{yogakal\=a} = 27 \textit{caga}. Subtracted from \textit{cakrakal\=a}, the resulting \textit{sarvayogakal\=a} gives the daily motion = \textit{cak} - 27 \textit{caga}. The \textit{kalpa} elapsed days multiplied by the desired \textit{cakra-kal\=a} yield the result; in the \textit{kalpa}, $^3$the revolutions = \textit{k\s u} --- 27 \textit{k\s bha\d m}. The remainder is derived from \textit{m\=rta}, \textit{nak\=satra}, and \textit{ayana} differences.~||50--52||%
}

\bilingual{%
  \dev{(बीएसएस् २१.५०--५२) एते श्लोकाः भोगकलाः (दैनिकं चलनं चाप-कलासु) गणयन्ति~| रवेः चन्द्रस्य ग्रहाणां च दैनिकं चलनं क्रमशः दिनयोः दीर्घत्वान्तरात् प्राप्यते~|}%
}{%
  \cnote{(BSS 21.50--52)} These verses compute the \textit{bhogakal\=a} (daily motion in minutes of arc). For the Sun, Moon, and planets, the daily motion is found by taking the difference in longitude between successive days and converting to \textit{kal\=as}.%
}

\bilingual{%
  \dev{वि. भा.---रविः (चन्द्रस्य) मध्यगतिकला अथवा$^{3}$वैर्क (ई, ई, ई) गुणा-}\\
  \dev{स्तदाऽऽध्यार्थ$^{3}$समनक्षत्राणां भोगकलाः स्युः~|}%
}{%
  \textit{V\=tti-bh\=a\d ya:} The mean daily motion in \textit{kal\=as} of the Sun (or Moon), multiplied by the respective factor, gives the \textit{bhogakal\=as} of the \textit{nak\=satra}s. (BSS 21.50--52 commentary)%
}

\bilingual{%
  \dev{भोगकला ग्रहस्य दैनिकं चलनं चाप-कलासु अस्ति~| गणनायां एते पदानि~: (१) क्रमशः दिनयोः दीर्घत्वान्तरस्य अन्वेषणम्, (२) एतद् अन्तरं कलासु रूपान्तरम्, (३) रवेः ~$\approx$~१$^{\circ}$ प्रतिदिनम् ($\approx$~६० कलाः), (४) चन्द्रस्य ~$\approx$~१३$^{\circ}$ प्रतिदिनम्~|}%
}{%
  The \textit{bhogakal\=a} is the daily motion of a planet in terms of \textit{kal\=as} (arc-minutes). The computation involves:
  \begin{itemize}
    \item Finding the difference in longitude between successive days
    \item Converting this difference into minutes (\textit{kal\=as})
    \item For the Sun: about 1° per day (approximately 60 \textit{kal\=as})
    \item For the Moon: about 13° per day
  \end{itemize}%
}

\newpage
%% ========================================================================
%% CHAPTER: Final Chapters --- Eclipses and Colophons
%% ========================================================================
\section{Final Chapters --- Eclipses and Colophons}

\bilingual{%
  \dev{ब्रह्मस्फुटसिद्धान्तस्य समाप्ति-भागः ग्रहणगणनायाः विस्तृतानि सूत्राणि समाविशति --- स्पर्शस्य (प्रथमस्पर्शस्य), मोक्षस्य (अन्तिमस्पर्शस्य), ग्रहणस्य परिमाणस्य, अक्षवलनस्य (अक्षच्युतेः) च~|}%
}{%
  The concluding portion of the \textit{Br\=ahmasphu\d tasiddh\=anta} contains extensive rules for eclipse calculations, including the computation of the \textit{spar\'sa} (first contact) and \textit{mok\d sa} (last contact), the magnitude of the eclipse, and the \textit{ak\s\d savalana} (deflection due to latitude).%
}

\bilingualShloka{%
  मन्दस्फुटार्केन्द्रांशोनायां ग्रहस्यन्दं केन्द्रम्~|\\
  एव यदि तत् स्पष्टं त्रिकोणं तदानीं विशेषं नोद्येत्~|\\
  तद्यात् स्पष्टं त्रिकोणं तस्य स्पष्टं द्वितीयं केन्द्रं\\
  क्रमेण केन्द्रान्तरं फलमिति~| वा ज्ञानमन्दं स्पष्टं\\
  प्राप्यच्छी~स्साचार्य्योदिना शीघ्र-कर्त्तव्यं कार्यमिति सर्वं शुद्धम्~||५७८-५८१||%
}{%
  mandasphu\d t\=arkendr\=a\'m\'son\=ay\=a\d m grahasyanda\d m kendram~|\\
  eva yadi tat spa\s\d ta\d m triko\na\d m tad\=an\=\=i\d m vi\'se\s\d a\d m nodyet~|\\
  tady\=at spa\s\d ta\d m triko\na\d m tasya spa\s\d ta\d m dvit\=iya\d m kendra\d m\\
  krame\n\=a kendr\=antara\d m phalamiti~| v\=a j\~n\=anamanda\d m spa\s\d ta\d m\\
  pr\=apyacch\=\=is\=s\=ac\=aryyodin\=a \'s\=\=ighra-karttavya\d m k\=aryamiti sarva\d m \'suddham~||578--581||%
}

\vspace{1em}
\longline
\vspace{0.5em}

\section*{Colophons --- \dev{समाप्ति}}

\bilingualShloka{%
  कलामयो यावता मानां भोगकलाः सृद्धास्तावद् गणमानि~|\\
  शेषाः ख लिख वर्तमाननक्षत्रस्य गतकलास्तास्तदुग्कलामय~|\\
  शुद्धा गम्य कला भवति~| ततो ग्रहद्वयको दिवसस्तदा\\
  गतगम्यकलाभिः किमित्वनुपातेन गणगम्या दिवमा\\
  भवति~|%
}{%
  kal\=amayo y\=avat\=a m\=an\=\=a\d m bhogakal\=a\d h s\d r\d dh\=ast\=avad ga\nam\=ani~|\\
  \'se\s\=\=a\d h kha likha vartam\=anak\=atrasya gatakal\=ast\=astadugkal\=amaya~|\\
  \'suddh\=a gamya kal\=a bhavati~| tato grahadvayako divasastad\=a\\
  gatagamyakal\=abhi\d h kimitvanup\=atena ga\nagamy\=a divam\=a\\
  bhavati~|%
}

\vspace{1em}
\begin{center}
{\Large\devbf{इति ब्रह्मस्फुटसिद्धान्त समाप्तः~|}}\\
\vspace{0.4cm}
\textit{Here ends the Br\=ahmasphu\d tasiddh\=anta}
\end{center}

\vspace{0.5em}
\bilingual{%
  \dev{इस प्रकार ब्रह्मस्फुटसिद्धान्त की समाप्ति हुई~| इस ग्रन्थ में कुल २४ अध्याय}\\
  \dev{हैं, जिनमें गणिताध्याय (१२ अध्याय) और गोलाध्याय (१२ अध्याय) सम्मिलित हैं~|}\\
  \dev{ब्रह्मगुप्ताचार्य ने इस ग्रन्थ की रचना संवत् ६२८ (ई.~६६६) में की थी~|}%
}{%
  Thus concludes the Br\=ahmasphu\d tasiddh\=anta. This treatise contains 24 chapters in all, comprising 12 chapters of mathematics (Ga\d nit\=adhy\=aya) and 12 chapters of astronomy (Gol\=adhy\=aya). The teacher Brahmagupta composed this work in \=Sa\d mvat 628 (CE 666).%
}

\vspace{1em}
\longline
\vspace{0.5em}

\section*{Appendix: Summary of Mathematical Rules}

\bilingual{%
  \dev{छाया-व्यवहारस्य सारांशः: छाया नरायाः (अक्षज्यायाः) समानुपातिकी, शङ्कोः (शङ्कु-उच्चतायाः) व्यस्तानुपातिकी च~|}%
}{%
  \textbf{Shadow rule (\dev{छाया व्यवहार}):} The shadow is proportional to the \textit{nara} (sine of latitude) and inversely proportional to the \textit{\'sa\~nku} (gnomon height).%
}

\bilingual{%
  \dev{अवर्गसंख्यानां वर्गमूलानयनस्य सूत्रम्:}\\
  $\sqrt{N} = a + \frac{r}{2a + 1}$\\
  \dev{यत्र} $a$ \dev{विद्यते यत्} $a^2 \leq N$ \dev{च} $r = N - a^2$%
}{%
  \textbf{Square root of non-squares (\dev{विकलवर्गानयन}):} The method of successive approximation:
  $ \sqrt{N} = a + \frac{r}{2a + 1} $
  where $a$ is the largest integer such that $a^2 \leq N$ and $r = N - a^2$.%
}

\bilingual{%
  \dev{मन्दफलस्य सारांशः: ग्रहपथस्य विकेन्द्रतायाः (अपकेंद्रतायाः) समीकरणं, केन्द्रस्य (अनुरूपज्यायाः) ज्यया कल्पितम्~|}%
}{%
  \textbf{Equation of centre (\dev{मन्दफल}):} The correction for the eccentricity of planetary orbits, computed using the \textit{jy\=a} (sine) of the \textit{kendra} (anomaly).%
}

\bilingual{%
  \dev{शीघ्रफलस्य सारांशः: ग्रहस्य स्थितेः सूर्य-पृथ्वी-रेखायाः सापेक्षः संशोधनः, शीघ्रकेन्द्रात् (योगस्य अनुरूपज्यायाः) गणितः~|}%
}{%
  \textbf{Second inequality (\dev{शीघ्रफल}):} The correction for the planet's position relative to the Sun-Earth line, computed from the \textit{\'s\=ighrakendra} (anomaly of conjunction).%
}

\bilingual{%
  \dev{कुट्टक-विधिः रैखिक-डायोफैण्टैन्-समीकरणानि $ax + c = by$ सोडयति, कालगणनायां आवश्यकम्~|}%
}{%
  \textbf{Ku\d t\d taka (\dev{कुट्टक}):} The pulverizer method for solving linear Diophantine equations $ax + c = by$, essential for calendar calculations.%
}

\bilingual{%
  \dev{भोगकला ग्रहस्य दैनिकं चलनं चाप-कलासु अस्ति, क्रमशः दिनयोः दीर्घत्वान्तरात् गणिता~|}%
}{%
  \textbf{Bhogakal\=a (\dev{भोगकला}):} The daily motion of a planet in arc-minutes, computed from the difference in longitude between successive days.%
}

\bilingual{%
  \dev{ज्याचापानयनम्: ज्यायाः चापस्य अन्वीक्षण-समस्या, ब्रह्मगुप्तेन बीजगणितीयसूत्रेण सोडिता, ज्यार्ध-पिठकस्य उपयोगं विना~|}%
}{%
  \textbf{Arc from sine (\dev{ज्याचापानयन}):} The inverse problem of finding the arc from its sine, solved by Brahmagupta's algebraic formula without using the \textit{jy\=ardha} (half-chord) table.%
}

\vspace{1em}
\begin{center}
\longline
\vspace{0.4cm}
{\small\textit{Typeset in XeLaTeX with Noto Serif Devanagari}}\\
{\small\textit{Bilingual edition generated: \today}}
\end{center}

\end{document}



% ============================================================
% Source: translations/en/indian/brahmagupta-brahmasphutasiddhanta-pt3-chunk2_bilingual.tex
% ============================================================

%% ========================================================================
%% BILINGUAL EDITION: Brahmagupta's Brahmasphutasiddhanta Part 3, Chunk 2
%% Pages 1817--1956 of the original edition
%% Sanskrit (Devanagari) || English Translation with IAST
%% ========================================================================
\documentclass[11pt,a4paper]{article}

%% -------- Core Packages --------
\usepackage{fontspec}
\usepackage{amsmath,amssymb}
\usepackage{geometry}
\usepackage{xcolor}
\usepackage{titlesec}
\usepackage{fancyhdr}
\usepackage{hyperref}

%% -------- Page Geometry --------
\geometry{
  paperwidth=210mm,paperheight=297mm,
  left=18mm,right=18mm,top=25mm,bottom=25mm,
  headheight=14pt,footskip=30pt
}

%% -------- Font Configuration --------
\setmainfont{Linux Libertine O}
\newfontfamily\devanagarifont[Script=Devanagari,Path=/usr/share/fonts/truetype/noto/,UprightFont=NotoSerifDevanagari-Regular.ttf,BoldFont=NotoSerifDevanagari-Bold.ttf]
\newfontfamily\iastfont{Linux Libertine O}

%% -------- Colors --------
\definecolor{titlecolor}{RGB}{45,55,72}
\definecolor{sectioncolor}{RGB}{55,65,81}
\definecolor{rulecolor}{RGB}{180,160,140}

%% -------- Section Formatting --------
\titleformat{\section}
  {\normalfont\Large\bfseries\color{sectioncolor}}
  {\thesection}{1em}{}
\titleformat{\subsection}
  {\normalfont\large\bfseries\color{sectioncolor}}
  {\thesubsection}{1em}{}

%% -------- Header/Footer --------
\pagestyle{fancy}
\fancyhf{}
\fancyhead[L]{\small\textit{Brahmasphutasiddhanta Part 3 --- Bilingual Edition}}
\fancyhead[R]{\small\thepage}
\renewcommand{\headrulewidth}{0.4pt}
\renewcommand{\footrulewidth}{0pt}

%% -------- Custom Commands --------
\newcommand{\dev}[1]{{\devanagarifont #1}}
\newcommand{\iast}[1]{{\iastfont\textit{#1}}}
\newcommand{\shortline}{\noindent\rule{0.5\textwidth}{0.4pt}\par}

%% Bilingual verse environment: two minipages side by side
\newcommand{\bilingualpair}[2]{%
  \noindent%
  \begin{minipage}[t]{0.48\textwidth}
    #1%
  \end{minipage}%
  \hfill%
  \begin{minipage}[t]{0.48\textwidth}
    #2%
  \end{minipage}%
  \par\vspace{0.5em}%
}

%% Sanskrit verse environment (left column)
\newenvironment{sanskritcol}{\devanagarifont\small}{}

%% English translation environment (right column)
\newenvironment{englishcol}{\iastfont\small}{}

%% -------- Hyperref Setup --------
\hypersetup{
  colorlinks=true,
  linkcolor=sectioncolor,
  citecolor=sectioncolor,
  urlcolor=sectioncolor,
  pdfencoding=auto
}

%% ========================================================================

\begin{document}

%% ========================================================================
%% TITLE PAGE
%% ========================================================================
\thispagestyle{empty}
\begin{center}
\vspace*{2cm}

{\Huge\color{titlecolor}\bfseries\dev{ब्राह्मस्फुटसिद्धान्त}}\\[0.8cm]
{\LARGE\color{sectioncolor}\dev{भाग ३}}\\[1.5cm]

\shortline

{\Large\bfseries Brāhmasphuṭasiddhānta}\\[0.3cm]
{\large\itshape Part 3 --- Chunk 2 (Pages 1817--1956)}\\[0.5cm]
{\small Bilingual Sanskrit--English Edition}\\[1cm]

\shortline

{\large By}\\[0.3cm]
{\LARGE\bfseries Brahmagupta}\\[0.5cm]
{\large (b.~598 CE)}\\[1.5cm]

\shortline

{\large Edited with Sanskrit Commentary (Vāsanā) and Hindi Translation}\\[0.3cm]
{\large by}\\[0.3cm]
{\large\bfseries Sudhākara Dvivedī}\\[1cm]

\shortline

{\small\textit{Comprising:}}\\[0.2cm]
{\small\dev{त्रिप्रश्नोत्तराध्यायः} (Tri-pra\-{}śnottara-adhyāya)}\\[0.1cm]
{\small\dev{ग्रहणोत्तराध्यायः} (Grahana-uttara-adhyāya)}\\[0.1cm]
{\small\dev{सूर्यसिद्धान्तयन्त्रविवरण} (Sūrya-siddhānta-yantra-vivarana)}\\[0.1cm]
{\small\dev{स्मृत्योर्युगलयोः समन्वयः} (Smṛti-yugala-samanvaya)}\\[2cm]

\vfill
{\small Typeset in XeLaTeX with Noto Serif Devanagari}\\
{\small Bilingual edition with IAST transliteration}\\
{\small \today}
\end{center}
\newpage

%% ========================================================================
%% TABLE OF CONTENTS
%% ========================================================================
\tableofcontents
\newpage

%% ========================================================================
%% CHAPTER XV: TRIPRASNOTTARADHYAYA
%% ========================================================================
\section*{\dev{त्रिप्रश्नोत्तराध्यायः}}
\addcontentsline{toc}{section}{Chapter XV: Tri-pra\-{}śnottara-adhyāya (Three Questions)}
\markboth{Tri-pra\-{}śnottara-adhyāya}{Tri-pra\-{}śnottara-adhyāya}

\noindent\rule{\textwidth}{0.6pt}
\begin{center}
\textit{Original pages 1033--1073 \quad|\quad Bilingual Edition}
\end{center}
\noindent\rule{\textwidth}{0.6pt}
\vspace{0.5cm}

%% ---- Introductory Commentary ----
\bilingualpair{%
\begin{sanskritcol}
\dev{अब श्रम्बिज्ञुत नक्षत्र के भोगानयन और ग्रहमुक्त नक्षत्रानयन को कहते हैं।}
\end{sanskritcol}%
}{%
\begin{englishcol}
\textbf{[Introductory Note.]} Now [the Ācārya] speaks of the computation of the motion (bhoga) of the starry nakṣatra and the computation of the nakṣatra as freed from the planets (graha-mukta-nakṣatrānayana).
\end{englishcol}%
}

\vspace{0.3cm}
\bilingualpair{%
\begin{sanskritcol}
\dev{पञ्चदशो योगभोगकलानामैक्यं = चं × १५}\\[0.3em]
\dev{सर्वेषां योगः = १ चं + ३ चं + १५ चं = २७ चं = सर्वनक्षत्रभोगः}
\end{sanskritcol}%
}{%
\begin{englishcol}
\textbf{[Mathematical Rule.]} The fifteenth [operation]: The sum of the degrees of conjunction-motion (yoga-bhoga-kalā) = Caṃ $\times$ 15. The sum of all = 1 Caṃ + 3 Caṃ + 15 Caṃ = 27 Caṃ = the total motion of all the nakṣatras (sarvana-\-kṣatra-bhogaḥ). Here Caṃ (cakra) denotes a complete revolution.
\end{englishcol}%
}

\vspace{0.5cm}
\bilingualpair{%
\begin{sanskritcol}
\dev{चक्कलास्यः शुद्धाः सर्वनक्षत्रभोगसंख्यास्तदा श्रम्बिज्ञुद् भोगकलास्तदिह्नगतिः =}\\[0.3em]
\dev{चक्र—२७ चं ततः कल्पोड्भवो भगणः := (चक्र—२७ चं) कुदिन =}\\[0.3em]
\dev{चक्कला}\\[0.3em]
\dev{कुदिन—२७ कल्पचंभगणा एतेनाचार्येण तुमूपपन्नम् । सिद्धान्त शेखरे}\\[0.3em]
\dev{``चक्कारिया वा शाश्वतसतिवधनाहतानि शोधयिन् सुदिनलयद्वरोषचक्र'' । स्यादे-}\\[0.3em]
\dev{कवासरम्बा कलिकागतिर्या सा वैवस्वेतावमध्यगाभिध्या मृतिः ॥ इद्धप्रहस्य}\\[0.3em]
\dev{कलिकालिकरादिसंशया नक्षत्रभोगकलिका प्रहसुत्कालानि । शेषात् भोग्यकलिका}\\[0.3em]
\dev{परितात् गम्य तास्यो भवन्ति गतगम्यदिनानि मुक्त्या ॥ इति श्रीपतेः पञ्चद्वय}\\[0.3em]
\dev{``सर्वेषां भोगैनितचक्कलिप्ता वैवस्वतः स्यादमिज्ञुद्भयोग'' इत्यादि भार्करोक्तं}\\[0.3em]
\dev{च सर्वाचार्येणस्य सर्वेषं समानार्थकार्मित ॥ ५०--५१ ॥}
\end{sanskritcol}%
}{%
\begin{englishcol}
\textbf{[Dvivedī's Commentary.]} The pure revolutions (cakkala) are the numbers [representing] the total motion of all the nakṣatras; then the degrees of motion of the starry nakṣatra (\'srambijñud bhoga-kalā) [are obtained]; its daily motion = revolution --- 27 Caṃ. From that, the kalpa-born revolutions [are obtained] as: (revolution --- 27 Caṃ) civil days (kudina) = cakkala.\\[0.3em]
In 27 civil days, the kalpa revolutions [are obtained] by this Ācārya; [this is] well established (tumūpapannam). In the \textit{Siddhānta\'sek\-{}hara}: ``cakkāriyā vā śāśvata-sativadhana-\=āhatāni śodhayin sudina-laya-dvaroṣa-cakra''. [The meaning is that] the daily motion of the kalpa (kalikā-gatiḥ) which is the Vaivasvata [value] at mid-[period] is to be known (mṛtiḥ). The degrees of nakṣatra motion (nakṣatra-bhoga-kalikā), the times of the prahara, etc., the remaining degrees to be traversed (bhogya-kalikā) --- from the excess, the gatagamya [values] become the days traversed and remaining to be traversed (gata-gamya-dināni). Thus [it is said] by Śrīpati in the fifteenth [chapter].\\[0.3em]
``The sum of the motions of all [nakṣatras] multiplied by the revolutions [gives] the Vaivasvata [value]; this is the yoga of the starry [nakṣatra],'' and so on --- thus spoke Bhāskara; and all teachers [hold] this same meaning.\\[0.3em]
\hfill \textbf{[Verses 50--51]}
\end{englishcol}%
}

\vspace{0.5cm}
\bilingualpair{%
\begin{sanskritcol}
\dev{उपपत्तिः}
\end{sanskritcol}%
}{%
\begin{englishcol}
\textbf{\textbar\ \textit{Upapatti} (Proof/Demonstration)}
\end{englishcol}%
}

\bilingualpair{%
\begin{sanskritcol}
\dev{षट् श्रद्धं भोगकला नक्षत्रों का ऐक्य =} $\frac{3\ \text{चं}}{2} \times \text{६} = ३$ \dev{चं × ३ = ६ चं}\\[0.3em]
\dev{षट् शद्धभोग कला नक्षत्रों का ऐक्य =} $\frac{\text{चं}}{2} \times \text{६} =$ \dev{चं × ३}\\[0.3em]
\dev{पनद्वह् एक भोग (समान भोग) कलानक्षत्रों का ऐक्य = चं × १५}\\[0.3em]
\dev{सबों का योग = सर्वनक्षत्र भोग कला = ६ चं + ३ चं + ? चं = २७ चं}
\end{sanskritcol}%
}{%
\begin{englishcol}
\textbf{[Mathematical Derivation.]} The sum of the six doubled (ṣaṭ-\-śraddha) motion-degrees of the nakṣatras = $\frac{3\ \text{Caṃ}}{2} \times 6 = 3$ Caṃ $\times$ 3 = 6 Caṃ. The sum of the six [times] single-motion degrees of the nakṣatras = $\frac{\text{Caṃ}}{2} \times 6$ = Caṃ $\times$ 3. The sum of the fifteen [times] single motion (sama-bhoga) degrees of the nakṣatras = Caṃ $\times$ 15. The sum of all = total nakṣatra motion degrees = 6 Caṃ + 3 Caṃ + [15] Caṃ = 27 Caṃ.
\end{englishcol}%
}

\vspace{0.5cm}
\bilingualpair{%
\begin{sanskritcol}
\dev{हिं. भा.---चन्द्रमध्यगति कला को ईं, ईं, ? इन से गुणा करने से षड्वर्षं, शर्बं,}\\[0.3em]
\dev{सम नक्षत्रों की भोग कला होती है । सब नक्षत्रों की जो भोग कला है उनको भगण कला}\\[0.3em]
\dev{में से घटाने से जो शेष रहता है वह श्रम्बिज्ञुत भोग है । वा कल्प चन्द्र भगण सतासी से}\\[0.3em]
\dev{गुणा कर कल्पकुदिन में से घटाने से शेष श्रम्बिज्ञुत का कल्प भगण होता है, इन भगणों से}\\[0.3em]
\dev{जो दिन भोग (कल्प कुदिन में यदि श्रम्बिज्ञुत का कल्प भगण पाते हैं तो एक दिन में क्या इस}\\[0.3em]
\dev{अनुपात से लब्ध कालात्मक दिनगति) होता है वह श्रम्बिज्ञुत भोग होता है । ग्रहकालसमूह}\\[0.3em]
\dev{में से नक्षत्र भोग कला को घटाने से नक्षत्र होते हैं, अर्थाद् ग्रहमुक्त कला में जितने नक्षत्रों की}\\[0.3em]
\dev{भोग कला शुद्ध हो उतने गतनक्षत्र होते हैं, शेषकला वर्तमान नक्षत्र की गत कला है उसको}\\[0.3em]
\dev{नक्षत्र भोग कला में से घटाने से गम्य कला होती है । तब ग्रहगति कला में एक दिन पाते हैं}\\[0.3em]
\dev{तो गत-गम्य कला में से क्या इस अनुपात से गतदिन और गम्यदिन होते हैं ॥ ५०--५२ ॥}
\end{sanskritcol}%
}{%
\begin{englishcol}
\textbf{[Hindi Bhāṣya --- English Summary.]} When the mean daily motion of the Moon (candra-madhyagati-kalā) is multiplied by [certain] quantities, the motion degrees of the six, the śarva, and the equal nakṣatras are obtained. The remaining [degrees], when subtracted from the revolution degrees (bhagaṇa-kalā), give the motion of the starry nakṣatra (\'srambijñuta-bhoga).\\[0.3em]
Alternatively: The kalpa revolutions of the starry nakṣatra are obtained by multiplying the kalpa revolutions of the Moon by eighty-seven and subtracting from the kalpa civil days. From these revolutions, the daily motion of the starry nakṣatra is obtained by proportion (anupāta): if in kalpa civil days such-\-and-\-such kalpa revolutions [are obtained], then in one day [what]? --- thus the daily motion in time-\-units.\\[0.3em]
When the nakṣatra motion degrees are subtracted from the planet-\-time aggregate (grahakāla-samūha), the nakṣatras are obtained. That is, in the planet-\-freed degrees (graha-mukta-kalā), as many pure motion degrees of the nakṣatras as there are, so many are the traversed nakṣatras (gata-nakṣatra). The remaining degrees are the current degrees of the present nakṣatra; subtracting these from the nakṣatra motion degrees yields the remaining degrees to be traversed (gamya-kalā). Then, if in one day [a certain value] is obtained in planetary motion degrees, by proportion from the gata-gamya degrees, the days traversed and days remaining are obtained.\\[0.3em]
\hfill \textbf{[Verses 50--52]}
\end{englishcol}%
}

\vspace{0.5cm}
\bilingualpair{%
\begin{sanskritcol}
\dev{उपपत्तिः --- यहाँ आचार्य ने पहले उक्त कालज्ञा को इष्टद्वित कल्पित किया है । तब श्रदा,}\\[0.3em]
\dev{समशाईं कुदूदित् इन तीनों मुञ्जाओं से उत्पन्न एक श्रम्ब क्षेत्र है, तथा श्रदज्ञा, लब्धज्ञा-}\\[0.3em]
\dev{तिज्ञा इन तीनों मुञ्जाओं से उत्पन्न द्वितीय श्रम्ब क्षेत्र है, ये दोनों श्रम्ब क्षेत्र सजातीय हैं}\\[0.3em]
\dev{इसलिए अनुपात करते हैं}
\end{sanskritcol}%
}{%
\begin{englishcol}
\textbf{\textbar\ \textit{Upapatti} (Proof).} Here the Ācārya has postulated (kalpita) the previously stated time-\-quantity (kālajñā) as the desired [value] for the second [process]. Then \'srada, sama-\'sāīṃ, kudūdit --- these three quantities (muñjā) produce one \'sramba field (kṣetra); and \'srada-\-jñā, labdha-\-jñā, tijñā --- these three quantities produce a second \'sramba field. These two \'sramba fields are homogeneous (sajātīya); therefore, proportion (anupāta) is applied [between them].
\end{englishcol}%
}

\bilingualpair{%
\begin{sanskritcol}
\begin{equation*}
\frac{\text{लब्धज्ञा. इष्ट}}{\text{तिज्ञा}} = \text{समशाईंकुद्}
\end{equation*}
\dev{तथा कालज्ञा =}
\end{sanskritcol}%
}{%
\begin{englishcol}
\begin{equation*}
\frac{\text{labdhajñā} \times \text{iṣṭa}}{\text{tijñā}} = \text{sama\'sāīṃ-kud}
\end{equation*}
And the time-\-quantity (kālajñā) = [the result of this proportion]. The three quantities (trayaṃ muñjānāṃ) generate the proportional relationship between the two \'sramba fields.
\end{englishcol}%
}

\vspace{0.5cm}
\bilingualpair{%
\begin{sanskritcol}
\dev{अत्रोपपत्तिः}
\end{sanskritcol}%
}{%
\begin{englishcol}
\textbf{\textbar\ \textit{Atrôpapattiḥ} (Demonstration thereof)}
\end{englishcol}%
}

\bilingualpair{%
\begin{sanskritcol}
\dev{छायाकर्णे गोले पलभा शङ्कुतुल्योस्तुल्यतत्त्वं भवति कथमिति प्रदर्शयते यदि}\\[0.3em]
\dev{द्वादशाङ्कुलराश्यों पलभा भुजां लभ्यते तदेप्तशाङ्कुं किमिति जातं शङ्कुतुल्यम् =}\\[0.3em]
\dev{$\frac{\text{पमा. शाङ्कु}}{१२}$, परन्तु शाङ्कु = $\frac{१२ \times \text{त्र}}{\text{छाक}}$ शत उत्थापनेन शङ्कुतुल्यम् = $\frac{\text{पमा. १२. त्र}}{१२.\ \text{छाक}}$}
\end{sanskritcol}%
}{%
\begin{englishcol}
\textbf{[Śloka on Shadow Computation.]} In the shadow-\-hypotenuse circle (chāyā-karṇe gole), how the equinoctial shadow (palabhā) and the gnomon (\'saṅku) become equal in nature (tulya-tattvaṃ) --- this is demonstrated thus: if the palabhā is obtained as the base (bhuja) in the twelve-\-aṅgula sign (dvāda-\'sāṅgula-rāśi), then what is the gnomon-\-like [quantity]? The gnomon-\-equivalent (\'saṅku-tulya) = $\frac{\text{pama} \times \text{\'saṅku}}{12}$; but \'saṅku = $\frac{12 \times \text{tri}}{\text{chāka}}$. By elevation (utthāpanena), the gnomon-\-equivalent = $\frac{\text{pama} \times 12 \times \text{tri}}{12 \times \text{chāka}}$.
\end{englishcol}%
}

\vspace{0.3cm}
\bilingualpair{%
\begin{sanskritcol}
\dev{यहाँ द्वाद (दिनार्क्) से भिन्न समय में पलभा साधन के लिए कहते हैं।}\\[0.3em]
\dev{हिं. भा.---दिनार्क् से भिन्न समय में रवि की छाया को छायाकर्णों से गुणा कर}\\[0.3em]
\dev{तिज्ञा से भाग देने से छायाकर्ण गोलिय छाया होती है, शङ्कुमूल और पूर्वापर रेखा का}\\[0.3em]
\dev{भेद भुज है । उत्तर गोल में कर्णदुताज्ञा में उत्तर भुज को घटाने से और दक्षिणा भुज}\\[0.3em]
\dev{को जोड़ने से पलभा होती है । दक्षिणा गोल में भुज में कर्णदुतिय छाया को घटाने ही से}\\[0.3em]
\dev{पलभा होती इति ॥ ४६--५० ॥}
\end{sanskritcol}%
}{%
\begin{englishcol}
\textbf{[Dvivedī's Explanation.]} Here, for computing the palabhā at a time different from midday (dvāda/dinārk), [the Ācārya] speaks.\\[0.3em]
\textbf{[Hindi Bhāṣya --- Summary.]} At a time different from midday, multiplying the Sun's shadow by the shadow-\-hypotenuse (chāyā-karṇa) and dividing by the tijñā yields the circular shadow of the hypotenuse-\-circle. The difference between the root of the gnomon (\'saṅku-mūla) and the east-\-west line is the base (bhuja). In the northern hemisphere (uttara-gola), subtracting the northern bhuja from the karṇa-dvitīya and adding the southern bhuja yields the palabhā. In the southern hemisphere (dakṣiṇa-gola), merely subtracting the second shadow from the bhuja yields the palabhā.\\[0.3em]
\hfill \textbf{[Verses 46--50]}
\end{englishcol}%
}

\newpage

%% ========================================================================
%% CHAPTER XVI: GRAHANOTTARADHYAYA
%% ========================================================================
\section*{\dev{ग्रहणोत्तराध्यायः}}
\addcontentsline{toc}{section}{Chapter XVI: Grahana-uttara-adhyāya (Eclipses Continued)}
\markboth{Grahana-uttara-adhyāya}{Grahana-uttara-adhyāya}

\noindent\rule{\textwidth}{0.6pt}
\begin{center}
\textit{Original pages 1089--1139 \quad|\quad Bilingual Edition}
\end{center}
\noindent\rule{\textwidth}{0.6pt}
\vspace{0.5cm}

\bilingualpair{%
\begin{sanskritcol}
\dev{अब श्रम्ब्य प्रश्नों को कहते हैं ।}\\[0.3em]
\dev{हिं. भा.---जो क्रान्ति के ज्ञाता समशकु को जानते हैं । जो लम्बांश वेत्सा समशुत्त}\\[0.3em]
\dev{को जानते हैं, कोणशकु को जानते हैं, कोणशकुछाया को जानते हैं । कोणदुताज प्रवेश में}\\[0.3em]
\dev{घटी को जानते हैं वे ज्योतिष सिद्धान्त के पण्डित हैं, यहाँ पाँच प्रश्न हैं ॥ ७ ॥}
\end{sanskritcol}%
}{%
\begin{englishcol}
\textbf{[Introduction.]} Now [the Ācārya] speaks of the questions of the \'srambya [type].\\[0.3em]
\textbf{[Hindi Bhāṣya --- Summary.]} Those who know the sama-\'saku [from] the declination (krānti), who know the sama-\'sutta [from] the latitude (lamba-\=aṃśa), who know the koṇa-\'saku, who know the koṇa-\'saku-chāyā, and who know the ghaṭī [from] the koṇa-dvitāja entrance --- they are the scholars (paṇḍita) of jyotiṣa-\-siddhānta. Here there are five questions.\\[0.3em]
\hfill \textbf{[Verse 7]}
\end{englishcol}%
}

\vspace{0.5cm}
\bilingualpair{%
\begin{sanskritcol}
\dev{इदानीमन्यान् प्रश्नानाह}
\end{sanskritcol}%
}{%
\begin{englishcol}
\textbf{\textbar\ \textit{Idānīm anyān praśnān āha}} (Now he states other questions.)
\end{englishcol}%
}

\vspace{0.3cm}
\bilingualpair{%
\begin{sanskritcol}
\dev{शङ्कुतलप्राच्यपरान्तरद्वयं बोध्यं यो विजानाति ।}\\[0.3em]
\dev{विशुवच्छायामेकं हस्त्वाऽऽसन्नदित्यं च गण्यकः सः ॥ ७ ॥}
\end{sanskritcol}%
}{%
\begin{englishcol}
\textit{\'saṅku-tala-prācyaparāntara-dvayaṃ bodhyaṃ yo vijānāti |}\\[0.3em]
\textit{vi\'suvacchāyām ekaṃ hastvā '\=asanna-dityaṃ ca gaṇyakaḥ saḥ || 7 ||}\\[0.3em]
He who, having known the two differences of the east-\-west [line] at the gnomon-\-base (\'saṅku-tala-prācya-parāntara), understands [also] the equinoctial shadow (vi\'suvac-chāyā) and the near-\-adjacent sun (\=asanna-\=āditya) --- he is a calculator (gaṇyaka).\\[0.3em]
\hfill \textbf{[Śloka 7]}
\end{englishcol}%
}

\vspace{0.3cm}
\bilingualpair{%
\begin{sanskritcol}
\dev{सू. भा.---शङ्कु तलप्राच्यपरान्तरं भुजः । यो भुजद्वयं बोध्य विशुवच्छाया}\\[0.3em]
\dev{पलभां विजानाति । एकमेव भुजं हस्त् वा विशुवच्छायामादित्यमकं च विजानाति}\\[0.3em]
\dev{स गण्यकः । एवमत्र प्रश्नत्रयम ॥ ८ ॥}
\end{sanskritcol}%
}{%
\begin{englishcol}
\textbf{[Sūtra-bhāṣya.]} The difference of the east-\-west [line] from the gnomon-\-base (\'saṅku-tala-prācya-parāntara) is the base (bhuja). He who, having understood the two bhuja-\-values, knows the equinoctial shadow (vi\'suvac-chāyā) [i.e.] the palabhā, or who knows one bhuja [together with] the equinoctial shadow and the sun --- he is a gaṇyaka. Thus here there are three questions.\\[0.3em]
\hfill \textbf{[Verse 8]}
\end{englishcol}%
}

\vspace{0.3cm}
\bilingualpair{%
\begin{sanskritcol}
\dev{वि. भा.---शङ्कुतलपूर्वापररेखयोरन्तरं भुजोर्धित, यो भुजद्वयं ज्ञात्वा पलभां}\\[0.3em]
\dev{जानाति एकं भुजं ज्ञात्वा पलभां रविं च जानाति स गण्यकोऽस्तीति । अत्र प्रश्नत्रय-}\\[0.3em]
\dev{भमस्ति ॥ ७ ॥}
\end{sanskritcol}%
}{%
\begin{englishcol}
\textbf{[Vivaraṇa-bhāṣya.]} The difference of the east-\-west line [measured] from the gnomon-\-base (\'saṅku-tala-pūrvāpara-rekha) is declared to be the bhuja. He who, having known the two bhuja-\-values, knows the palabhā, or having known one bhuja, knows the palabhā and the Sun (ravi) --- he is a gaṇyaka. Thus here there are three questions.\\[0.3em]
\hfill \textbf{[Verse 7]}
\end{englishcol}%
}

\vspace{0.5cm}
\bilingualpair{%
\begin{sanskritcol}
\dev{अब श्रम्ब्य प्रश्नों को कहते हैं ।}\\[0.3em]
\dev{हिं. भा.---शङ्कुतल और पूर्वापर रेखा का भेद भुज है । जो व्यक्ति दो भुजों को}\\[0.3em]
\dev{ज्ञान कर पलभा को जानते हैं । एवं एक भुज को ज्ञान कर पलभा को जानते हैं तथा रवि}\\[0.3em]
\dev{को जानते हैं वे ज्योतिः शास्त्र के पण्डित हैं इति । यहाँ तीन प्रश्न हैं ॥ ७ ॥}
\end{sanskritcol}%
}{%
\begin{englishcol}
\textbf{[Dvivedī's Introduction.]} Now [the Ācārya] speaks of the \'srambya questions.\\[0.3em]
\textbf{[Hindi Bhāṣya --- Summary.]} The difference between the gnomon-\-base (\'saṅku-tala) and the east-\-west line is the bhuja. The person who, having known the two bhuja-\-values, knows the palabhā, and likewise [he who], having known one bhuja, knows the palabhā and the Sun (ravi) --- they are the paṇḍitas of jyotiḥ-\'śāstra. Here there are three questions.\\[0.3em]
\hfill \textbf{[Verse 7]}
\end{englishcol}%
}

\vspace{0.5cm}
\bilingualpair{%
\begin{sanskritcol}
\dev{इदानीमन्यान् प्रश्नानाह}
\end{sanskritcol}%
}{%
\begin{englishcol}
\textbf{\textbar\ \textit{Idānīm anyān praśnān āha}} (Now he states other questions.)
\end{englishcol}%
}

\vspace{0.3cm}
\bilingualpair{%
\begin{sanskritcol}
\dev{पातालशङ्कुमुध्येस्ते वा हस्तज्ञां रविंविजानाति ।}\\[0.3em]
\dev{इष्टपातालगणकोः पृथक् तले वा स तन्त्रज्ञः ॥ ८ ॥}
\end{sanskritcol}%
}{%
\begin{englishcol}
\textit{pātāla-\'saṅkum uddhyeste vā hasta-jñāṃ raviṃ vijānāti |}\\[0.3em]
\textit{iṣṭa-pātāla-gaṇakoḥ pṛthag tale vā sa tantra-jñaḥ || 8 ||}\\[0.3em]
He who knows the pātāla-\'saṅku [and] the uddeśa-\'saṅku, or [who knows] the sun having known the hasta; [or who knows] the iṣṭa-pātāla separately, each at its own base (tale) --- he is a knower of the tantra (tantra-jña).\\[0.3em]
\hfill \textbf{[Śloka 8]}
\end{englishcol}%
}

\vspace{0.3cm}
\bilingualpair{%
\begin{sanskritcol}
\dev{सू. भा.---यो रवेः पातालशङ्कु कुमथः शङ्कुं विजानाति । उदयेस्ते वा यो}\\[0.3em]
\dev{रविंद्‌र्ग् ज्ञात्रो विजानाति । इष्टशङ्कुकुदिवोइष्टशङ्कुः पातालगणःः पातालशङ्कुः}\\[0.3em]
\dev{रविरथः शङ्कुः । तयोः पृथक् पृथक् तले शङ्कुतले च वा यो विजानाति स एव}\\[0.3em]
\dev{तन्त्रज्ञ इति । एवमत्र प्रश्नचतुष्टयम ॥ ९ ॥}
\end{sanskritcol}%
}{%
\begin{englishcol}
\textbf{[Sūtra-bhāṣya.]} He who knows the pātāla-\'saṅku of the Sun [and] the kumatha-\'saṅku, or who knows the Sun having known the uddheśa, or [who knows] the iṣṭa-\'saṅku [and] the kudivo-iṣṭa-\'saṅku [as] the pātāla-gaṇaḥ and pātāla-\'saṅku [respectively], the ravi-ratha-\'saṅku --- of these two, separately, each at its own base (pṛthak tale) or at the gnomon-\-base (\'saṅku-tale), he who knows --- he alone is the tantra-jña. Thus here there are four questions.\\[0.3em]
\hfill \textbf{[Verse 9]}
\end{englishcol}%
}

\vspace{0.5cm}
\bilingualpair{%
\begin{sanskritcol}
\dev{अब श्रम्ब्य प्रश्नों को कहते हैं ।}\\[0.3em]
\dev{हिं. भा.---जो व्यक्ति ग्रहण में राहुमार्ग (भूमामार्ग) को जानते हैं । उस मार्ग}\\[0.3em]
\dev{को जानते हैं, कोणशकुछाया को जानते हैं । यहाँ पाँच प्रश्न हैं ॥ ७ ॥}
\end{sanskritcol}%
}{%
\begin{englishcol}
\textbf{[Dvivedī's Introduction.]} Now [the Ācārya] speaks of the \'srambya questions.\\[0.3em]
\textbf{[Hindi Bhāṣya --- Summary.]} The person who knows the Rāhu-path (bhūmā-mārga) in an eclipse, who knows that path, and who knows the koṇa-\'saku-chāyā --- here there are five questions.\\[0.3em]
\hfill \textbf{[Verse 7]}
\end{englishcol}%
}

\vspace{0.5cm}
\bilingualpair{%
\begin{sanskritcol}
\dev{कोणदुतस्थरविगतं भुवप्रोत्कुटं भुवाद्रिविपर्यं भुव्यचापांशं एकं भुजः ।}\\[0.3em]
\dev{भुवाल्पसतिकावधि लम्बांश द्वितीयो भुजः । कोणदुतं व स्वैनिकाद्रिविपर्यं ननां-}\\[0.3em]
\dev{शास्तुतीयो भुजः । त्रिभुजेऽस्मिन् रविगतं भुवप्रोत्कुटतयाप्येतर्कुर्नाम्यामृत्यन्न-}\\[0.3em]
\dev{कोणा नतकालः । याम्योत रवित्कोणदुतनाम्यामृत्यन्नकोणः = ८५ नतोनुपातः}\\[0.3em]
\dev{किम्यते यदि नतकालज्ञया हस्तज्ञा लभ्यते तदाश्रवेदांशज्ञया (ज्ञा ८५) किमिनि}\\[0.3em]
\dev{समागच्छति भुज्या = $\frac{\text{हस्तज्ञा. ज्ञा ८५}}{\text{नतकालज्ञा}} = \frac{\text{हस्तज्ञा. त्रि. ज्ञा ८५}}{\text{नतकालज्ञा. त्रि}} = \frac{\text{हस्तज्ञा. त्रि. २८५}}{२४९५\ \text{नतकालज्ञा}}$}
\end{sanskritcol}%
}{%
\begin{englishcol}
\textbf{[Eclipse Triangle Computation.]}\\[0.3em]
\textit{koṇaduta-stha-ravi-gataṃ bhuva-protkuṭaṃ bhuvādri-viparyaṃ bhuvyacāpāṃśaṃ ekaṃ bhujaḥ |}\\[0.3em]
\textit{bhuvālpasatikāvadhi lamba-aṃśa dvitīyo bhujaḥ | koṇadutaṃ va svainikādri-viparyaṃ nanāṃ-\'śāstu-tīyo bhujaḥ |}\\[0.3em]
In this triangle: the Sun's passage at the koṇaduta, the earth's protuberance (protkuṭa), the earth-\-mountain reversal --- the degrees of the earth-\-curve (bhuvyacāpāṃśa) --- these form one bhuja. The latitude-\-degrees (lamba-aṃśa) up to the earth's alpaka-\-satika [form] the second bhuja. The koṇaduta and the svainika-\-mountain reversal [form] the third bhuja.\\[0.3em]
In this triangle, the Sun's passage being [known] as the earth's protuberance (bhuva-protkuṭa), the angle [is] the natakala (nata-kālaḥ). The southern koṇaduta-nata-amṛtyanna-koṇa = 85; by the rule of proportion (anupātaḥ): if by the natakāla-jyā the hasta-jyā is obtained, then by the arc-jyā of 85 (jyā 85), what is obtained?\\[0.3em]
The bhuja-jyā = $\frac{\text{hasta-jyā} \times \text{jyā } 85}{\text{natakāla-jyā}} = \frac{\text{hasta-jyā} \times \text{tri} \times \text{jyā } 85}{\text{natakāla-jyā} \times \text{tri}} = \frac{\text{hasta-jyā} \times \text{tri} \times 285}{2495 \times \text{natakāla-jyā}}$.
\end{englishcol}%
}

\newpage

\bilingualpair{%
\begin{sanskritcol}
\dev{उपपत्तिः}
\end{sanskritcol}%
}{%
\begin{englishcol}
\textbf{\textbar\ \textit{Upapattiḥ} (Proof/Demonstration)}
\end{englishcol}%
}

\bilingualpair{%
\begin{sanskritcol}
\dev{यहाँ देशान्तर विदित नहीं है इसलिए देशान्तर संस्कार से रहित स्फुट रवि और}\\[0.3em]
\dev{स्फुट चन्द्र से चन्द्र ग्रहण विधि से सर्वग्रहण में संमीलन काल और उन्मीलन काल साधन}\\[0.3em]
\dev{करना चाहिए । उस दिन में द्वितीय से भी संमीलन काल समज्ञा चाहिए वह यदि गणितानगत्}\\[0.3em]
\dev{संमीलन काल से अधिक हो तो द्वितीय रेखा से पूर्व भाग में अथवा पश्चिम भाग में स्थित है यह}\\[0.3em]
\dev{समज्ञा चाहिए । क्योंकि रेखा से पूर्व स्वदेश में पहले स्वदेश में पश्चात् रेखा देश में मध्याह्न}\\[0.3em]
\dev{काल होता है । अतः रेखा देशीय इष्ट संमीलन काल से स्वदेशीय संमीलन काल अधिक होता}\\[0.3em]
\dev{है पश्चिम में इसके विपरीत होता है । इस तरह उन्मीलन काल से इष्ट प्रास काल से स्पर्श}\\[0.3em]
\dev{काल से या मोक्ष काल से परीक्षा होती है, स्पर्श काल परीक्षा और मोक्ष काल परीक्षा इष्ट}
\end{sanskritcol}%
}{%
\begin{englishcol}
\textbf{[Dvivedī's Upapatti --- Summary.]} Here, since the longitude difference (deśāntara) is not known, one should compute the conjunction-\-time (saṃmīlana-kāla) and disjunction-\-time (unmīlana-kāla) of the total eclipse by the method of lunar eclipse, from the true Sun (sphuṭa-ravi) and true Moon (sphuṭa-candra), devoid of deśāntara correction. On that day, the second saṃmīlana-kāla should [also] be computed; if it exceeds the computational saṃmīlana-kāla, then one should know whether [the place] is situated in the eastern or western part from the second line.\\[0.3em]
Because to the east of the line, midday (madhyāhna-kāla) occurs first in one's own country and later in the line-country; therefore, the saṃmīlana-kāla of one's own country exceeds the saṃmīlana-kāla of the line-country. In the west, the opposite occurs. Thus, from the unmīlana-kāla, the iṣṭa-prāsa-kāla, the sparśa-kāla, or the mokṣa-kāla should be examined.
\end{englishcol}%
}

\vspace{0.5cm}
\bilingualpair{%
\begin{sanskritcol}
\dev{उत्तरगोले याम्ये विषुवच्छायायुक्तस्तरं हीनम् ।}\\[0.3em]
\dev{एवं विषुवच्छाया युक्तविहीनास्तरेषाभा ॥ ५० ॥}
\end{sanskritcol}%
}{%
\begin{englishcol}
\textit{uttara-gole yāmye viṣuvacchāyāyuktas taraṃ hīnam |}\\[0.3em]
\textit{evaṃ viṣuvacchāyā yuktavihīnāstareṣābhā || 50 ||}\\[0.3em]
In the northern hemisphere (uttara-gola), [the latitude] being joined with the equinoctial shadow (viṣuvac-chāyā) [in the] south (yāmya) is diminished (hīnam). Thus, the equinoctial shadow [is] joined with or diminished from the latitudes (tareṣā) [depending on hemisphere].\\[0.3em]
\hfill \textbf{[Śloka 50]}
\end{englishcol}%
}

\vspace{0.3cm}
\bilingualpair{%
\begin{sanskritcol}
\dev{सू. भा.---प्रथमायाः पूर्वार्धं त्रिप्रश्नाध्यायस्य चतुर्थीयुपयुपर्धसंं न्यस्यात्-}\\[0.3em]
\dev{मेव । अन्यत् सर्वं च त्रिप्रश्ना ५८--६० आयामिः स्फुटम् ॥ ४९--५० ॥}
\end{sanskritcol}%
}{%
\begin{englishcol}
\textbf{[Sūtra-bhāṣya.]} The first half of the earlier [section] of the Tri-praśna-\-adhyāya [is to be] placed as the fourth subsidiary section (upayupardha-saṃṃ nyasyāt). All else of the Tri-\-praśna [verses] 58--60, the extent (āyāmi) [is] clear.\\[0.3em]
\hfill \textbf{[Verses 49--50]}
\end{englishcol}%
}

\vspace{0.3cm}
\bilingualpair{%
\begin{sanskritcol}
\dev{वि. भा.---दिनार्धदिनकाले शक्रिया छायाकर्णगुणा तिज्ञामक्ता तदा}\\[0.3em]
\dev{छायाकर्णगोलिया रवैर्या भवेत् । शङ्कुमूलपूर्वापररेखयोरन्तरं भुजः । तेन भुजे-}\\[0.3em]
\dev{नोन्यता कर्णदुताप्राच्यादितुरेया भुजेनोनदिक्षेत्रं भुजेन युता तदोतरगोले}\\[0.3em]
\dev{विषुवच्छाया (पलभा) भवेत् । याम्ये (दक्षिणा गोले) तया कर्णदुतियाप्राच्यादन्तरं}\\[0.3em]
\dev{(भुजमान) हीन तदा पलभा भवेत् । एवमन्तरेषा (भुजेन) युक्तविहीनास्तस्यथा-}\\[0.3em]
\dev{दन्तरभुजेन युता दक्षिणाम्यभुजेन हीना तदा पलभा भवेदिति ॥ ४९--५० ॥}
\end{sanskritcol}%
}{%
\begin{englishcol}
\textbf{[Vivaraṇa-bhāṣya.]} At the half-day time (dinārdha-dina-kāle), the shadow multiplied by the shadow-hypotenuse (chāyā-karṇa-guṇā), divided by the tijñā, then the circular shadow of the hypotenuse-circle (chāyā-karṇa-golīyā) would be [obtained] as the ravi-value. The difference between the root of the gnomon (\'saṅku-mūla) and the east-west line is the bhuja. By that bhuja, diminished by the karṇa-\-dvitīya-\-prācya, [or] joined with it --- in the northern hemisphere, the equinoctial shadow (viṣuvac-chāyā), i.e., the palabhā, results. In the south (i.e., the southern hemisphere), when that is diminished by the difference from the karṇa-\-dvitīya-\-prācya (bhuja-māna), then the palabhā results. Thus, [the shadow] joined with or diminished from the difference by the bhuja, or diminished by the southern bhuja, then the palabhā results.\\[0.3em]
\hfill \textbf{[Verses 49--50]}
\end{englishcol}%
}

\vspace{0.5cm}
\bilingualpair{%
\begin{sanskritcol}
\dev{अत्रोपपत्तिः}
\end{sanskritcol}%
}{%
\begin{englishcol}
\textbf{\textbar\ \textit{Atrôpapattiḥ} (Demonstration thereof)}
\end{englishcol}%
}

\bilingualpair{%
\begin{sanskritcol}
\dev{छायाकर्ण गोले पलभा शङ्कु तुल्ययोस्तुल्यतत्त्वं भवति कथमिति प्रदर्शयते यदि}\\[0.3em]
\dev{द्वादशाङ्कुलराश्यों पलभा भुजां लभ्यते तदेप्तशाङ्कुं किमिति जातं शङ्कुतुल्यम् =}\\[0.5em]
\begin{equation*}
\frac{\text{पमा. शाङ्कु}}{१२}, \quad \text{परन्तु शाङ्कु} = \frac{१२ \times \text{त्र}}{\text{छाक}} \quad \text{शत उत्थापनेन}
\end{equation*}\\[0.5em]
\begin{equation*}
\text{शङ्कुतुल्यम्} = \frac{\text{पमा. १२. त्र}}{१२.\ \text{छाक}} = \text{पमा}
\end{equation*}
\end{sanskritcol}%
}{%
\begin{englishcol}
\textbf{[Proof of Gnomon-Equivalence.]}\\[0.3em]
How the palabhā in the shadow-\-hypotenuse circle (chāyā-karṇa-gola) and the gnomon (\'saṅku) become equal in nature (tulya-tattvaṃ) is demonstrated thus: if the palabhā is obtained as the bhuja in the twelve-\-aṅgula sign, then what is the gnomon-\-equivalent (\'saṅku-tulyam)?\\[0.3em]
\[ \text{\'saṅku-tulya} = \frac{\text{pama} \times \text{\'saṅku}}{12}, \quad \text{but } \text{\'saṅku} = \frac{12 \times \text{tri}}{\text{chāka}} \text{, by elevation} \]
\[ \text{\'saṅku-tulya} = \frac{\text{pama} \times 12 \times \text{tri}}{12 \times \text{chāka}} = \text{pama} \]
\end{englishcol}%
}

\vspace{0.3cm}
\bilingualpair{%
\begin{sanskritcol}
\dev{$\therefore$ सिद्धं कर्णगोले शङ्कुतुल्यम् = पलभा । कर्णदुताज्ञा व्यस्तगोला भवति ।}\\[0.3em]
\dev{पलभा च सव्योतरा । तथाः संस्कारतःछायाप्रविपरसूनमध्ये भुजः कथतेऽस्तत्-}\\[0.3em]
\dev{क्षेपरियेन कर्णदुताज्ञा भुजयाः संस्करेणा पलभा भवतीति । सिद्धान्तशेखरे ``अथ-}\\[0.3em]
\dev{या तु नतपूर्वपश्चिमाशान्तरेषा शकिबुतज्ञामका । दक्षिणोत्तरभुजा युतोनिता}\\[0.3em]
\dev{सौम्यवर्तिन् रवो पलभा ॥ दक्षिणेन पुनरिननायां तथा हीनमेव हि तदन्तरं}\\[0.3em]
\dev{सदा । एवमन्तरयुतोनिता भवेद्-छाया नियतमुगुलाश्चका ॥'' श्रीपत्युक्तसिद्धमा-}\\[0.3em]
\dev{चार्येणानुरूपमेवेति ॥ ४९--५० ॥}
\end{sanskritcol}%
}{%
\begin{englishcol}
Therefore, it is proved (siddhaṃ) that in the karṇa-gola, the gnomon-\-equivalent (\'saṅku-tulya) equals the palabhā. The karṇa-dvitīya-jyā becomes the inverse circle (vyasta-golā). The palabhā is the north-\-south [quantity]. Thus, by the correction of the chāyā-pravi-parasūna, the bhuja is spoken of; by the addition-\-subtraction (kṣepa-riyena) with the karṇa-dvitīya-jyā, by the correction of the bhuja, the palabhā results.\\[0.3em]
In the \textit{Siddhānta\'sek\-{}hara}: ``Now, the difference of the natapūrva-paścima (east-\-west natal difference), the \'śakibuta-jñā-makā. The dakṣiṇôttara-bhujā, joined with and diminished, the Sun dwelling in the north (saumya-vartin) [gives] the palabhā. In the southern [hemisphere], again, at the inanāyām, thus it is always diminished by that difference. Thus, joined with and diminished by the difference, the shadow becomes fixed [in] aṅgula-\'ścakā.'' This [is] said by Śrīpati; [it is] in accordance with the Ācārya alone.\\[0.3em]
\hfill \textbf{[Verses 49--50]}
\end{englishcol}%
}

\vspace{0.5cm}
\bilingualpair{%
\begin{sanskritcol}
\dev{वलनस्यमुचे तैम्यः पृथग् ग्रहक् खण्डकेन । वृत्तिः कृतिः स्पर्शविमर्शनिमध्यग्रता}\\[0.3em]
\dev{क्रमेणोर्बिमहावगमयः ॥}
\end{sanskritcol}%
}{%
\begin{englishcol}
\textit{valanasya muce taimyaḥ pṛthag grahak khaṇḍakena | vṛttiḥ kṛtiḥ sparśavimarśanimadhyagratā krameṇôr bimha-vagamayaḥ ||}\\[0.3em]
Of the deflection (valana), the release (muca), the taimya, separately; by the section of the planet (graha-khaṇḍakena); the vṛtti, the kṛti, the contact and release (sparśa-vimarśa), the centrality of the middle (nimadhyagratā) --- in sequence, the understanding of the eclipse-\-disc (ur-bimha) [is obtained].
\end{englishcol}%
}

\vspace{0.3cm}
\bilingualpair{%
\begin{sanskritcol}
\dev{``भास्कराचार्येण श्रीपतिप्रकार एव विश्वदर्शनेण प्रति}\\[0.3em]
\dev{पादितः ॥ १४--१५॥}
\end{sanskritcol}%
}{%
\begin{englishcol}
By Bhāskarācārya, the very method of Śrīpati [was] propounded (pratipāditaḥ) in the \textit{Vi\'śvadarpaṇa}.\\[0.3em]
\hfill \textbf{[Verses 14--15]}
\end{englishcol}%
}

\vspace{0.5cm}
\bilingualpair{%
\begin{sanskritcol}
\dev{विषुवदर्पमण्डलदिका इत्यादि दो प्रश्नों के उत्तर ग्रहणाधिकार में बता दिए}\\[0.3em]
\dev{गये हैं । उसके उत्तरार्ध को श्रम्बे कहते हैं ।}\\[0.3em]
\dev{अब `सम्पर्कं मण्डलै' इत्यादि प्रश्न के उत्तर कहते हैं ।}\\[0.3em]
\dev{हिं. भा.---प्राह्लादूता में स्वल्पान्तर से सरलाकार तात्कालिक क्रान्तिज्ञुनीय चाप}\\[0.3em]
\dev{बलन सूत्र का । बलनज्ञा से दिव्यज्ञान समज्ञा चाहिए । मानचक्रबाहृत में बलन सूत्र के}\\[0.3em]
\dev{उपर रवि के स्पर्श कालिक सार और मोक्ष कालिक सार को जिस दिशा के सार है उसी दिशा}\\[0.3em]
\dev{में देना । मध्यबलन सूत्र में प्राह्ल बिम्ब केन्द्र से मध्यार देना चाहिए । चन्द्र शराघ्रों से प्राह्ल}\\[0.3em]
\dev{प्रमाण व्यास से तुरत लिखकर स्पर्श मोक्ष श्रीर प्रास समज्ञा चाहिए, एवं पृथिवी पर परिलेख}\\[0.3em]
\dev{लिखने से स्पर्श मोक्ष श्रीर प्रास होता है इति ॥}
\end{sanskritcol}%
}{%
\begin{englishcol}
\textbf{[Dvivedī's Commentary on Eclipse Drawing.]}\\[0.3em]
The answers to the two questions beginning ``viṣuvadarpa-\-maṇḍala-dikā'' etc., have been explained in the eclipse chapter (grahaṇâdhikāra). The latter half of that [is called] \'srambe.\\[0.3em]
Now the answers to the question beginning ``samparkaṃ maṇḍalai'' etc. are spoken.\\[0.3em]
\textbf{[Hindi Bhāṣya --- Summary.]} In the prāhlādūta, the arc of the instantaneous declination (tātkālika krānti) is of a simple form with slight difference; [this is] the balana-sūtra. From the balana-jyā, the divya-jyā should be computed. In the māna-cakra-bāhṛta [circle], above the balana-sūtra, the contact-sphāra (sparśa-kālika sāra) and release-sphāra (mokṣa-kālika sāra) of the Sun should be given in the direction of their [respective] sāra. In the madhya-balana-sūtra, the madhya should be given from the centre of the prāhla-disc. From the moon's horns (\'śarāghra), by the prāhla-pramāṇa-vyāsa [diameter], having drawn quickly, the contact, release, and mid-line prāsa should be computed. And by drawing the parilekha [projection] on the earth, the contact-release \'śrīra-prāsa is obtained.
\end{englishcol}%
}

\vspace{0.5cm}
\bilingualpair{%
\begin{sanskritcol}
\dev{उपपत्तिः}
\end{sanskritcol}%
}{%
\begin{englishcol}
\textbf{\textbar\ \textit{Upapattiḥ} (Proof/Demonstration)}
\end{englishcol}%
}

\bilingualpair{%
\begin{sanskritcol}
\dev{मानचक्रबाहृत् वृत्त में जब प्राह्लाकृता का केन्द्र होता है तब प्राह्ल बिम्बग्रान्त और प्राह्ल}\\[0.3em]
\dev{बिमग्रान्त संलग्न रहता है । इस लिये विदित मानचक्रबाहृत् वृत्त में दिशा को श्रीष्ट करना}\\[0.3em]
\dev{चाहिए । सममण्डल और मानचक्रबाहृत् वृत्त का समाप्त बिन्दु मानचक्रबाहृत् वृत्त में प्राची चिह्न है ।}\\[0.3em]
\dev{वहाँ से बलनज्ञा दान देने से केन्द्र से बलनज्ञाप्रगत रेखा क्रान्तितृत् प्राची है बलनसूत्र से}\\[0.3em]
\dev{ज्यावत् खादान देना चाहिए । क्योंकि क्रान्तिबृत् प्राची से सार दक्षिण और उत्तर रहता है । इस}\\[0.3em]
\dev{तरह स्पर्श और मोक्ष में होता है । मध्यबलन तात्कालिकक्रान्तिबृत् प्राची से दक्षिणोत्तर}\\[0.3em]
\dev{दिशा में होता है इसलिए केन्द्र से बलन सूत्र में देना चाहिए । शराघ्र में प्राह्ल वृत्त का केन्द्र}\\[0.3em]
\dev{होता है इसलिए वहाँ से तात्कालिकप्राह्लाभि प्रमाण से रचित बृत्तों से स्पर्श मोक्षमध्य होने है}\\[0.3em]
\dev{सिद्धान्त शेखर में `मानचक्रबाहृतलये तत्परियामिनिवि' इत्यादि संस्कृतोपपत्ति में लिखित}\\[0.3em]
\dev{ल्लोक से श्रीपति ने आचार्यों के अनुकूल ही कहा है । सिद्धान्त शिरोमणि में `प्राह्लाखसूत्रेय}\\[0.3em]
\dev{विधाय वृत्तं मानचक्रबाहृतं च साधितांशम्' इत्यादि से भास्कराचार्य ने श्रीपति प्रकार ही को}\\[0.3em]
\dev{विचार रूप से प्रतिपादित किया है इति ॥ १४--१५ ॥}
\end{sanskritcol}%
}{%
\begin{englishcol}
\textbf{[Dvivedī's Proof --- Summary.]} When in the māna-cakra-bāhṛta circle there is the centre of the prāhlā-kṛtā, then the prāhla-bimbagrānta and prāhla-bimagrānta remain contiguous (saṃlagna). Therefore, in the known māna-cakra-bāhṛta circle, the direction should be fixed (\'śrīṣṭa). The common endpoint of the sama-maṇḍala and the māna-cakra-bāhṛta circle is the eastern mark (prācī cihna) in the māna-cakra-bāhṛta circle.\\[0.3em]
From there, giving the balana-jyā, the line from the centre [through] the balana-jyā-pragata is the krānti-tṛt-prācī. One should give the jyāvat-khādana from the balana-sūtra. Because the sāra remains south and north from the krānti-vṛt-prācī. Thus it happens in the sparśa and mokṣa. The madhya-balana [is] in the dakṣiṇôttara direction from the tātkālika-krānti-vṛt-prācī; therefore, from the centre it should be given in the balana-sūtra. In the \'śarāghra, the prāhla-vṛtta has its centre; therefore, from there, by the tātkālika-prāhlābhi-pramāṇa, from the constructed circles, the sparśa-mokṣa-madhya [is obtained].\\[0.3em]
In the \textit{Siddhānta\'sek\-{}hara}, from the verse ``mānacakrabāhṛtalaye tatpariyāminivi'' etc., written in the Sanskrit proof, Śrīpati has spoken in accordance with the Ācāryas alone. In the \textit{Siddhānta\'śiro\-{}maṇi}, from ``prāhlākhasūtreya vidhāya vṛttaṃ mānacakrabāhṛtaṃ ca sādhitāṃśam'' etc., Bhāskarācārya has propounded in the form of deliberation (vicāra-rūpa) the very method of Śrīpati.\\[0.3em]
\hfill \textbf{[Verses 14--15]}
\end{englishcol}%
}

\vspace{0.5cm}
\bilingualpair{%
\begin{sanskritcol}
\dev{इदानी यः परिलिखतीष्टद्यासमिल्यादि प्रश्नस्योत्तरमाह ।}\\[0.3em]
\dev{पश्चात् प्रग्रहयो प्रागमोषे रविविब्बलो बाहुः ।}\\[0.3em]
\dev{स्वबलनसिद्धज्ञादिशि विपरीतः शीतकरमध्यात् ॥ १६ ॥}\\[0.3em]
\dev{भानुमतो बाहुगुण्याख्या दिशं कोटिरत्यथा शविनः ।}\\[0.3em]
\dev{रविशास्तिवमध्यात् करणास्तिन्यकं, करणांगुकोटियुतेः ॥ १७ ॥}
\end{sanskritcol}%
}{%
\begin{englishcol}
\textit{idānī yaḥ parilikhatīṣṭadyāsamilyādi praśnasyôttarām āha |}\\[0.3em]
\textit{paścāt pragrahayo prāgamoṣe ravi-vibbalo bāhuḥ |}\\[0.3em]
\textit{sva-balana-siddha-jñā-diśi viparītaḥ śītakara-madhyāt || 16 ||}\\[0.3em]
\textit{bhānumato bāhu-guṇyākhāyā diśaṃ koṭir atyathā śavinaḥ |}\\[0.3em]
\textit{ravi-śāsti-vamadhyāt karaṇāstinyakaṃ, karaṇāṃgu-koṭiyuteḥ || 17 ||}\\[0.3em]
Now he who draws the answer to the question of iṣṭa-dyā-samīlyādi: the bāhu [is] the western and eastern pragraha at the release (moṣa) of the Sun's vibbalo [force]. Opposite in direction to the svabalana-siddha-jyā, from the centre of the Moon (\'śītakara).\\[0.3em]
Of the Sun (bhānumat), the direction of the bāhu multiplied (bāhu-guṇyākhāyā diśaṃ); the koṭi [is] otherwise (atyathā) of Saturn (\'śanaiścara/\'śavina). From the centre of the Sun's \'śāsti [force], karaṇa-astinyakaṃ; from the karaṇāṅgu-koṭi-yuti [is obtained the result].\\[0.3em]
\hfill \textbf{[Ślokas 16--17]}
\end{englishcol}%
}

\newpage

%% ========================================================================
%% SECTION: SURYA-SIDDHANTA YANTRA-VIVARANA
%% ========================================================================
\section*{\dev{सूर्यसिद्धान्तयन्त्रविवरण}}
\addcontentsline{toc}{section}{Sūrya-siddhānta-yantra-vivarana (Description of Sūrya-siddhānta Instruments)}
\markboth{Sūrya-siddhānta-yantra-vivarana}{Sūrya-siddhānta-yantra-vivarana}

\noindent\rule{\textwidth}{0.6pt}
\begin{center}
\textit{Original pages 1121--1139 \quad|\quad Bilingual Edition}
\end{center}
\noindent\rule{\textwidth}{0.6pt}
\vspace{0.5cm}

\bilingualpair{%
\begin{sanskritcol}
\dev{सूर्य सिद्धान्त में शब्दों लिखित यन्त्र विवरण है---}
\end{sanskritcol}%
}{%
\begin{englishcol}
\textbf{[Introduction.]} In the Sūrya-siddhānta, the following verbal description of instruments (yantra-vivaraṇa) is given ---
\end{englishcol}%
}

\vspace{0.3cm}
\bilingualpair{%
\begin{sanskritcol}
\dev{तुङ्गार्कसमायुक्तं गोलयन्त्रं प्रसाधयेत् ।}\\[0.3em]
\dev{गोपयेेतत्प्रकाशोक्तं सर्वगंयं भवेद्धि ॥}\\[0.5em]
\dev{कालसंसाधनार्थं तथा यन्त्रारिणा साधयेत् ।}\\[0.3em]
\dev{एकाकी योजयेद्‌बीजं यन्त्रं विस्मय कारिणि ॥}\\[0.5em]
\dev{शङ्कुकुयष्टि धनुश्चक्रेषछायायन्त्रैर्नैकथा ।}\\[0.3em]
\dev{गुप्तदेशाङ्कज्ञेयं कालज्ञानमतन्द्रितः ॥}\\[0.5em]
\dev{तोययन्त्रकपालाङ्कघर्म्म्युदनत्रवानरैः ।}\\[0.3em]
\dev{ससूत्ररेखुगमैश्च समयकालं प्रसाधयेत् ॥}\\[0.5em]
\dev{परदारामनुष्याणां गुल्वतैलजलानि च ।}\\[0.3em]
\dev{बीजानि पांसवस्तेषु प्रयोगास्तेऽपि दुर्लभाः ॥}\\[0.5em]
\dev{ता प्रपञ्चमधिश्छद्द् न्यस्तं कुण्डे स्मलाभमिस् ।}\\[0.3em]
\dev{द्विट्ठमेजलयहोरात्रं स्फुटं यन्त्रं कपालकम् ॥}\\[0.5em]
\dev{नयन्त्रं तथा साधु दिवा च विमले रवौ ।}\\[0.3em]
\dev{छाया संसाधनैः प्रोक्तं कालसाधनमुत्तमम् ॥}
\end{sanskritcol}%
}{%
\begin{englishcol}
\textit{tuṅgârkasamāyuktaṃ gola-yantraṃ prasādhayet |}\\[0.3em]
\textit{gopayeetat-prakāśoktaṃ sarvagamyaṃ bhaved dhi ||}\\[0.3em]
One should prepare the spherical instrument (gola-yantra) endowed with the exalted Sun (tuṅgârka-samāyuktaṃ); having concealed it, [as] spoken of in the Prakāśa, it becomes indeed universally accessible (sarvagamyaṃ).\\[0.5em]
\textit{kāla-saṃsādhanārthaṃ tathā yantrāriṇā sādhayet |}\\[0.3em]
\textit{ekākī yojayed bījaṃ yantraṃ vismaya-kāriṇi ||}\\[0.3em]
For the purpose of time-determination (kāla-saṃsādhana), one should thus accomplish [it] by the yantrâriṇa; the solitary one should apply the seed (bījaṃ) [to] the instrument, causing wonder.\\[0.5em]
\textit{\'saṅku-kuyaṣṭi dhanuś-cakreṣa-chāyā-yantrair naikathā |}\\[0.3em]
\textit{gupta-deśāṅka-jñeyaṃ kāla-jñānam atandritaḥ ||}\\[0.3em]
By gnomon-staff (\'saṅku-kuyaṣṭi), bow-wheel (dhanu-cakra), and shadow-instruments (chāyā-yantraiḥ) in manifold ways; the knowledge of time (kāla-jñāna) is to be known from the hidden place-\-marker (gupta-deśāṅka) [by one who is] vigilant.\\[0.5em]
\textit{toya-yantra-kapālāṅka-gharmmy-udanatravānaraiḥ |}\\[0.3em]
\textit{sa-sūtra-rekhu-gamaiś ca samaya-kālaṃ prasādhayet ||}\\[0.3em]
By water-instrument (toya-yantra), bowl (kapāla), aṅka of heat (gharma), water-\-stream (udanatra), and monkey [devices?] (vānaraiḥ), and with cord-\-line-\-rekhuga (sa-sūtra-rekhu-gamaiḥ), one should determine the proper time (samaya-kālaṃ).\\[0.5em]
\textit{paradārā-manuṣyāṇāṃ gulvataila-jalāni ca |}\\[0.3em]
\textit{bījāni pāṃsavas teṣu prayogās te 'pi durlabhāḥ ||}\\[0.3em]
The wives of others (paradāra) of men, oil (taila), and water --- [these are] seeds (bījāni) [in] the sand (pāṃsavaḥ) of these; their applications are also difficult to obtain (durlabhāḥ).\\[0.5em]
\textit{tā prapañcamadhiśchadd nyastaṃ kuṇḍe smalābhamis |}\\[0.3em]
\textit{dviṭṭha-meja-laya-horātraṃ sphuṭaṃ yantraṃ kapālakam ||}\\[0.3em]
That covering of the world (prapañcam-adhiśchadd), placed in the bowl (kuṇḍe), smalābhamis [?]; the double-\-positioned (dviṭṭha) meja-laya day-\-night (horātra) --- the clear instrument [is] the kapālaka.\\[0.5em]
\textit{na yantraṃ tathā sādhu divā ca vimale ravau |}\\[0.3em]
\textit{chāyā saṃsādhanaiḥ proktaṃ kāla-sādhanam uttamam ||}\\[0.3em]
Not the instrument thus [alone is] good; by day, when the Sun is pure (vimale ravau), the shadow [is] declared by the saṃsādhana-\-methods [to be] the best time-determination (kāla-sādhana).
\end{englishcol}%
}

\vspace{0.5cm}
\bilingualpair{%
\begin{sanskritcol}
\dev{मानाध्याय}
\end{sanskritcol}%
}{%
\begin{englishcol}
\textbf{\textbar\ \textit{Mānâdhyāya} (Chapter on Measures)}
\end{englishcol}%
}

\vspace{0.3cm}
\bilingualpair{%
\begin{sanskritcol}
\dev{मानानि सौरचान्द्रार्क्षसावनानि ग्रहानयनमैभिः ।}\\[0.3em]
\dev{मानैः पृथक् चतुर्भिः संख्यबहारोड्ज्ञ लोकस्य ॥}
\end{sanskritcol}%
}{%
\begin{englishcol}
\textit{mānāni saura-cāndrārkṣa-sāvanāni grahāṇayanam aibhiḥ |}\\[0.3em]
\textit{mānaiḥ pṛthak caturbhiḥ saṃkhyabahāroḍjña lokasya ||}\\[0.3em]
The measures (mānāni): solar (saura), lunar (cāndra), stellar (nakṣa), and civil (sāvana) --- by these, the computation of planets (grahâṇayanam). By these four measures separately, the multitudinous arrangements (saṃkhyabahāroḍjña) of the world [are made].
\end{englishcol}%
}

\vspace{0.3cm}
\bilingualpair{%
\begin{sanskritcol}
\dev{इससे सौरमान, चान्द्रमान, नाक्षत्रमान और सावनमान, ये चार प्रकार के मान}\\[0.3em]
\dev{कहे गये हैं । इन्हीं चारों मानों से लोगों के सब व्यवहार होते हैं । किस किस मान से कौन}\\[0.3em]
\dev{कौन पदार्थ ग्रहण किये जाते हैं यह सब प्रति पादित है । बाह्म, दिव्य, पिट्र, प्राजापत्य,}\\[0.3em]
\dev{बाहृस्पत्य, सौर, सावन, चान्द्र, नाक्ष ये नौ मान हैं । इन मानों में से मनुष्यलोक}\\[0.3em]
\dev{में केवल सौर, चान्द्र, सावन और नाक्ष इन चार मानों की ही प्रधानता है । क्योंकि}\\[0.3em]
\dev{इन्हीं मानों से मनुष्यों के सब व्यवहार सम्पन्न होते हैं । सूर्यसिद्धान्त, सिद्धान्तशेखर, सिद्धान्त-}\\[0.3em]
\dev{शिरोमणि आदि सब ग्रन्थों में मानों के विषय में समान रूप से कहा गया है । ब्रह्मस्फुट}\\[0.3em]
\dev{सिद्धान्त के इस अध्याय में भूमादर्शन के भी साधन हैं ।}
\end{sanskritcol}%
}{%
\begin{englishcol}
\textbf{[Dvivedī's Hindi Commentary --- Summary.]}\\[0.3em]
From this, the solar measure (saura-māna), lunar measure (cāndra-māna), stellar measure (nākṣatra-māna), and civil measure (sāvana-māna) --- these four types of measures are spoken. By these very four measures, all the affairs of people are conducted. Which objects are grasped by which measure --- all this is propounded (pratipādita).\\[0.3em]
Brāhma, divya, paitṛ, prājāpatya, bāhṛspatya, saura, sāvana, cāndra, nākṣa --- these are nine measures. Of these measures, in the world of humans, only the saura, cāndra, sāvana, and nākṣa have primacy, because by these measures alone all human affairs are accomplished. In the Sūrya-siddhānta, Siddhānta\'sek\-{}hara, Siddhānta\'śiro\-{}maṇi, and all other treatises, the subject of measures is spoken of in the same manner. In this chapter of the Brāhmasphuṭasiddhānta, the methods of viewing the earth (bhūmādarśana) are also given.
\end{englishcol}%
}

\newpage

%% ========================================================================
%% SECTION: SMRITYOR-YUGALA-SAMANVAYA
%% ========================================================================
\section*{\dev{स्मृत्योर्युगलयोः समन्वयः}}
\addcontentsline{toc}{section}{Smṛti-yugala-samanvaya (Reconciliation of Smṛti Texts)}
\markboth{Smṛti-yugala-samanvaya}{Smṛti-yugala-samanvaya}

\noindent\rule{\textwidth}{0.6pt}
\begin{center}
\textit{Original pages 1133--1139 \quad|\quad Bilingual Edition}
\end{center}
\noindent\rule{\textwidth}{0.6pt}
\vspace{0.5cm}

\bilingualpair{%
\begin{sanskritcol}
\dev{स्मृतियों में और पुराणों में राहुग्रस्त ग्रहण का प्रतिपादन विद्यमान है । अतः दोनों मतों}\\[0.3em]
\dev{का समन्वय करते हुए आचार्य ने कहा है---}
\end{sanskritcol}%
}{%
\begin{englishcol}
\textbf{[Introduction.]} In the Smṛtis and in the Purāṇas, the doctrine of Rāhu-grasta eclipse (rāhugrasta-grahaṇa) is present (pratipādana). Therefore, reconciling the two doctrines, the Ācārya has said ---
\end{englishcol}%
}

\vspace{0.3cm}
\bilingualpair{%
\begin{sanskritcol}
\dev{राहुस्तच्छदयति प्रविशति यच्छुकलपच्छदयन्ते ।}\\[0.3em]
\dev{छाया तमसीद्वैरेघ्रदानां कमलयोनिः ॥}\\[0.5em]
\dev{चन्द्रोऽस्मुयोजः स्यो यदिनिमयभास्करस्य मासात् ।}\\[0.3em]
\dev{छादयति शर्मिततापो राहुश्छदयति तत् सवितुः ॥}
\end{sanskritcol}%
}{%
\begin{englishcol}
\textit{rāhus tac chadayati praviśati yac chuklapacchadayante |}\\[0.3em]
\textit{chāyā tamasi-dvairaghra-dānāṃ kamala-yoniḥ ||}\\[0.3em]
\textit{candro 'smuyojaḥ syo yadinaimaya-bhāskarasya māsāt |}\\[0.3em]
\textit{chādayati śarmi-tatāpo rāhuś chadayati tat savituḥ ||}\\[0.3em]
Rāhu covers that [Moon], entering [it] which the bright and dark [portions] cover. The shadow (chāyā) [is] the tamasi [force] of the enemy of the lotus-\-born [i.e., Rāhu, enemy of Viṣṇu?].\\[0.3em]
The Moon (candra), the asmu-yojaḥ, the syo [?], from the māsa of the sun-\-of-\-the-\-yadinaimaya [?] --- [this] Rāhu covers the scorching-\-heat (\'śarmi-tatāpaḥ) [of] that Sun (savituḥ).
\end{englishcol}%
}

\vspace{0.5cm}
\bilingualpair{%
\begin{sanskritcol}
\dev{सिद्धान्तशिरोमणि के गोलाध्याय में भी अप्रोक्तित भास्करोक्ति---}
\end{sanskritcol}%
}{%
\begin{englishcol}
In the Golâdhyāya of the Siddhānta\'śiro\-{}maṇi also, the unspoken (aproktita) statement of Bhāskara [is as follows] ---
\end{englishcol}%
}

\vspace{0.3cm}
\bilingualpair{%
\begin{sanskritcol}
\dev{द्विदेशकालवरयाणि मेदानां छादको राहुरिति स वर्णित ।}\\[0.3em]
\dev{यन्मानिनः केवलगोलविधास्तत्सहिता वेद्यरणाबाह्याम् ॥}\\[0.5em]
\dev{राहुः क्षुभामण्डलगः शशाङ्कं शशाङ्कश्छादयतीव बिम्बम् ।}\\[0.3em]
\dev{तमोमयः शम्युवरप्रदानात् सर्वगमानामविरुद्धमेतत् ॥}
\end{sanskritcol}%
}{%
\begin{englishcol}
\textit{dvi-deśa-kāla-varayāṇi medānāṃ chādako rāhur iti sa varṇita |}\\[0.3em]
\textit{yan māninaḥ kevala-golavidhās tat-sahitā vedy araṇābāhyām ||}\\[0.3em]
Rāhu, the eclipser (chādakaḥ) of the vara of the two-\-place-\-time (dvi-deśa-kāla-varayāṇi medānāṃ), thus he is described. Those who know only the spherical method (kevala-golavidhāḥ), together with that [knowledge], should know the araṇā-bāhyām [?].\\[0.5em]
\textit{rāhuḥ kṣubhā-maṇḍalagaḥ śaśāṅkaṃ śaśāṅkaś chādayatīva bimbam |}\\[0.3em]
\textit{tamo-mayaḥ śamyuva-pradānāt sarvagamānām aviruddham etat ||}\\[0.3em]
Rāhu, situated in the agitated orb (kṣubhā-maṇḍalagaḥ), [eclipses] the Moon (\'śaśāṅkaṃ); the Moon appears to cover the disc (bimbam). [Rāhu is] composed of darkness (tamo-mayaḥ); from the bestowal of \'śamyuva [?], this is not contradictory (aviruddham) for all who go [i.e., all observers].
\end{englishcol}%
}

\vspace{0.3cm}
\bilingualpair{%
\begin{sanskritcol}
\dev{से समन्वय किया गया है । सिद्धान्तशेखर में राहुग्रस्त ग्रहण के खण्डनार्थ `राहुनिरा-}\\[0.3em]
\dev{करणाध्याय' नाम का एक अध्याय रखा गया है । इसमें श्रीपति ने भी निम्नलिखित}\\[0.3em]
\dev{श्लोक ने समन्वय किया है---}
\end{sanskritcol}%
}{%
\begin{englishcol}
\textbf{[Dvivedī's Commentary.]} This reconciliation (samanvaya) has been made. In the Siddhānta\'sek\-{}hara, for the refutation (khaṇḍanârtha) of the Rāhu-grasta eclipse, a chapter named ``Rāhu-nirākaraṇâdhyāya'' has been placed. In this, Śrīpati also has made the reconciliation with the following ślokas ---
\end{englishcol}%
}

\vspace{0.3cm}
\bilingualpair{%
\begin{sanskritcol}
\dev{विष्णुलूनविशिरसः किल पञ्चदेत्तवान् वरमिमं परमेष्ठी ।}\\[0.3em]
\dev{हेमदानविधिना तवतुपतिस्तमरतिमहस्तदुपरागे ॥}\\[0.5em]
\dev{भूमेच्छाया प्रविष्टः स्थगयति शशिनं शुक्लपक्षावसाने ।}\\[0.3em]
\dev{राहुर्भुप्रसादात् समविगतवस्तुतामो व्यसुतुल्यः ॥}\\[0.5em]
\dev{उर्व्यां स्थं भानुबिम्बं सलिलमयतनोर्यथैवोक्तं बिम्बम् ।}\\[0.3em]
\dev{संसूयैवं च मासव्युपरतिसमये स्वस्य साहित्यहेतोः ॥}
\end{sanskritcol}%
}{%
\begin{englishcol}
\textit{viṣṇu-lūna-viśirasaḥ kila pañcad ettavān varam imaṃ parameṣṭhī |}\\[0.3em]
\textit{hema-dāna-vidhinā tavatupatis tam-arati-mahas tad uparāge ||}\\[0.3em]
The Supreme Lord (parameṣṭhī) indeed gave this excellent boon (varam imaṃ) to the headless Viṣṇu-lūna [i.e., the decapitated Rāhu], [that] by the method of donating gold (hema-dāna-vidhinā), at that eclipse (tad uparāge), [there is] the great darkness (tam-arati-mahaḥ) [for him?].\\[0.5em]
\textit{bhūmecchāyā praviṣṭaḥ sthagayati śaśinaṃ śuklapakṣāvasāne |}\\[0.3em]
\textit{rāhur bhu-prasādāt samavigatavastu-tāmo vyasutulyaḥ ||}\\[0.3em]
Entering the earth's shadow (bhūmecchāyā praviṣṭaḥ), [Rāhu] obscures the Moon (sthagayati śaśinaṃ) at the end of the bright fortnight (\'śuklapakṣâvasāne). Rāhu, by the grace of the earth (bhu-prasādāt), [being] equal to the vyasu [?], the darkness [is] the real state of things (samavigatavastu-tāmaḥ).\\[0.5em]
\textit{urvyāṃ sthaṃ bhānu-bimbaṃ salila-maya-tanor yathaivôktaṃ bimbam |}\\[0.3em]
\textit{saṃsūyaivaṃ ca māsavy uparati-samaye svasya sāhitya-hetoḥ ||}\\[0.3em]
Standing on the earth (urvyāṃ sthaṃ), the solar disc (bhānu-bimbam), [this] disc as described [earlier], of the water-\-bodied one [i.e., the Moon] --- thus indeed at the time of lunar cessation (māsavy-uparati-samaye), for the sake of its conjunction (svasya sāhitya-hetoḥ).
\end{englishcol}%
}

\vspace{0.5cm}
\bilingualpair{%
\begin{sanskritcol}
\dev{गोलबन्धाधिकार में महत् द्वृतों (पूर्वापरदृत, याम्योत्तरदृत, स्थितिजदृत आदि) की रचना}\\[0.3em]
\dev{तथा लघुदृतों (मेषादिक द्वादश राशियों के बहिरोराजदृत) की रचना करके परमलम्बन-नति}\\[0.3em]
\dev{का स्वरूप प्रतिपादन कर आचार्य ने द्वृतों का ज्ञानयन किया}\\[0.3em]
\dev{है । ब्रह्मस्फुटसिद्धान्त में आचार्य ने जैसे गोलबन्ध कहा है}\\[0.3em]
\dev{वैसे ही सिद्धान्त शेखर में श्रीपति और सिद्धान्त शिरोमणि के}\\[0.3em]
\dev{गोलाध्याय में भास्कराचार्य ने कहा है । ग्रहगोल और नक्षत्रगोल में पाँच स्थिरदृत}\\[0.3em]
\dev{(पूर्वापरदृत, स्थितिजदृत, याम्योत्तरदृत, उन्मण्डल, विषुवदृत) कहे हैं । ये सब कला-}\\[0.3em]
\dev{मण्डल के बराबर हैं । तथा ग्रहों के चलदृत मन्दनीचोचदृत = ७, मन्दादि ग्रहों के शीघ्र-}\\[0.3em]
\dev{नीचोचदृत = ५ । मन्दप्रतिबृत = ७, शीघ्रप्रतिबृत = ७ । सात ग्रहों के हमण्डल हृद्क्षेप}
\end{sanskritcol}%
}{%
\begin{englishcol}
\textbf{[Dvivedī's Commentary on Spherical Construction.]}\\[0.3em]
In the Golabandhâdhikāra, having constructed the great circles (mahad-vṛttāḥ: pūrvâpara-vṛtta, yāmyôttara-vṛtta, sthitija-vṛtta, etc.) and the small circles (laghu-vṛttāḥ: the outer royal circles of the twelve signs from Meṣa, etc.), the Ācārya has demonstrated the nature (svarūpa) of the maximum equation (paramalambana-nati) and computed the knowledge of the circles (vṛttôpāyanam).\\[0.3em]
As the Ācārya has spoken of the golabandha in the Brāhmasphuṭasiddhānta, so also has Śrīpati in the Siddhānta\'sek\-{}hara, and Bhāskarācārya in the Golâdhyāya of the Siddhānta\'śiro\-{}maṇi. In the graha-gola and nakṣatra-gola, five fixed circles (sthira-vṛtta) are spoken: pūrvâpara-vṛtta, sthitija-vṛtta, yāmyôttara-vṛtta, unmaṇḍala, and viṣuvad-vṛtta. All these are equal to the kalā-maṇḍala.\\[0.3em]
And the moving circles (calavṛtta) of the planets: mandanīcôcavṛtta = 7, the \'śīghranīcôcavṛtta of the manda-etc. planets = 5. The manda-prativṛtta = 7, the \'śīghra-prativṛtta = 7. The hamaṇḍala and hṛd-kṣepa of the seven planets [are also given].
\end{englishcol}%
}

\vspace{1cm}

%% ========================================================================
%% END MATTER
%% ========================================================================
\newpage
\section*{Colophon}

\begin{center}
\shortline
\vspace{0.5cm}

{\large\textit{End of Chunk 2 (Pages 1817--1956)}}\\[0.3cm]
{\small Comprising the Tri-pra\'snottara-adhyāya and Grahana-uttara-adhyāya}\\[0.3em]
{\small with explanatory notes on the Sūrya-siddhānta and Smṛti reconciliation}\\[0.5cm]

\shortline

\vspace{1cm}

{\small\textit{``Brāhmasphuṭasiddhānta'' is one of the most important works}}\\
{\small\textit{of ancient Indian mathematics and astronomy, composed by}}\\
{\small\textit{Brahmagupta in 628 CE at Bhillamāla (modern Bhinmal, Rajasthan).}}\\[0.5cm]

{\small\textit{This bilingual edition presents the original Sanskrit in Devanagari}}\\
{\small\textit{alongside English translations with IAST transliteration,}}\\
{\small\textit{based on the edition of Sudhākara Dvivedī (Part 3).}}\\[1cm]

\shortline

\vspace{0.5cm}
{\small Typeset in XeLaTeX with Noto Serif Devanagari}\\
{\small Sanskrit--English bilingual edition}\\
{\small \today}
\end{center}

\end{document}



% ============================================================
% Source: translations/en/indian/brahmagupta-brahmasphutasiddhanta-pt3-chunk3_bilingual.tex
% ============================================================

%% ========================================================================
%% Brahmagupta's Brahmasphutasiddhanta — Part 3, Chunk 3 (FINAL)
%% Bilingual English-Sanskrit Edition
%% Pages 1957-2512 (original pages 22-160 of Part 3)
%% ========================================================================
\documentclass[10pt,a4paper]{article}

%% -------- Core Packages --------
\usepackage{fontspec}
\usepackage{geometry}
\usepackage{xcolor}
\usepackage{fancyhdr}
\usepackage{array,booktabs,longtable,multirow}
\usepackage{hyperref}

%% -------- Page Geometry --------
\geometry{
  paperwidth=210mm,paperheight=297mm,
  left=16mm,right=16mm,top=22mm,bottom=22mm,
  headheight=14pt,footskip=24pt
}

%% -------- Font Configuration --------
\setmainfont[Ligatures=TeX]{TeXGyreTermes}
\setsansfont{TeXGyreHeros}

%% Devanagari Font
\newfontfamily\devanagarifont[Script=Devanagari]{Noto Serif Devanagari}

%% -------- Colors --------
\definecolor{titlecolor}{RGB}{45,55,72}
\definecolor{sectioncolor}{RGB}{55,65,81}
\definecolor{rulecolor}{RGB}{180,160,140}
\definecolor{versecolor}{RGB}{30,40,60}
\definecolor{commentarybg}{RGB}{250,248,245}

%% -------- Header/Footer --------
\pagestyle{fancy}
\fancyhf{}
\fancyhead[L]{\small\textit{Brāhmasphuṭasiddhānta — Bilingual Edition}}
\fancyhead[R]{\small\thepage}
\renewcommand{\headrulewidth}{0.4pt}

%% -------- Custom Commands --------
\long\def\dev#1{{\devanagarifont #1}}
\newcommand{\shortline}{\noindent\rule{0.5\textwidth}{0.4pt}\par}

%% Parallel columns: Sanskrit (left) | English (right)
\newenvironment{parallel}
  {\noindent\begin{minipage}[t]{0.47\textwidth}}
  {\end{minipage}}
\newenvironment{englishcol}
  {\hfill\begin{minipage}[t]{0.47\textwidth}}
  {\end{minipage}\par\vspace{1em}}

%% Shloka environment (Sanskrit verse)
\newenvironment{shloka}
  {\begin{center}\begin{minipage}{0.9\textwidth}\begin{center}\devanagarifont\color{versecolor}}
  {\end{center}\end{minipage}\end{center}}

%% Commentary box
\newenvironment{commentarybox}
  {\begin{quote}\small\color{sectioncolor}\devanagarifont}
  {\end{quote}}

%% English commentary
\newenvironment{engcommentary}
  {\begin{quote}\small\color{sectioncolor}}
  {\end{quote}}

%% Bilingual verse pair
\newcommand{\bilingualverse}[2]{%
\noindent\begin{minipage}[t]{0.47\textwidth}
\dev{#1}
\end{minipage}%
\hfill\vline\hfill%
\begin{minipage}[t]{0.47\textwidth}
#2
\end{minipage}\par\vspace{0.8em}%
}

%% -------- Hyperref --------
\hypersetup{
  colorlinks=true,linkcolor=sectioncolor,citecolor=sectioncolor,
  pdfencoding=auto,
  pdftitle={Brahmasphutasiddhanta Part 3 Chunk 3 — Bilingual Edition},
  pdfauthor={Brahmagupta (c. 598-668 CE)}
}

%% ========================================================================
\begin{document}

%% ========================================================================
%% TITLE PAGE
%% ========================================================================
\begin{titlepage}
\centering
\vspace*{2cm}

{\Huge\bfseries\color{titlecolor}\dev{ब्राह्मस्फुटसिद्धान्त}}\\[0.6cm]
{\LARGE\bfseries\color{sectioncolor}\dev{भाग तृतीय — खण्ड ३}}\\[0.4cm]
{\large\textit{(Bilingual Sanskrit-English Edition)}}\\[1.5cm]

\shortline

{\Large\textbf{Brāhmasphuṭasiddhānta}}\\[0.3cm]
{\large Part 3 — Chunk 3 (FINAL)}\\[0.3cm]
{\normalsize Pages 1957--2512 / Original pages \dev{२२ — १६०}}\\[1.5cm]

\shortline

{\Large\textbf{Brahmagupta}}\\[0.4cm]
{\large c. 598--668 CE}\\[1.5cm]

\vfill

{\small\textit{The Brāhmasphuṭasiddhānta is the masterpiece of Brahmagupta,}}
{\small\textit{written in 628 CE, comprising 24 chapters on mathematics and astronomy.}}\\[0.5cm]
{\small Chapters 17--21: Grahagaṇita, Kuṭṭaka, Chāyā, and Golādhyāya}

\end{titlepage}

%% ========================================================================
%% TABLE OF CONTENTS
%% ========================================================================
\newpage
\section*{\dev{विषयानुक्रमणिका} / Table of Contents}

\begin{center}
\begin{tabular}{@{}p{0.55\textwidth}p{0.35\textwidth}@{}}
\toprule
\textbf{\dev{संस्कृत}} & \textbf{English} \\
\midrule
\dev{१७. सृङ्गोत्स्त्युतराध्यायः} & Ch. 17: Advanced Planetary Computations \\
\dev{१८. कुट्टकाध्यायः} & Ch. 18: The Pulverizer (Kuṭṭaka) \\
\dev{१९. शङ्कुच्छायादिज्ञानाध्यायः} & Ch. 19: Gnomon and Shadow \\
\dev{२०. छद्दिच्छायाध्यायः} & Ch. 20: Advanced Shadow Problems \\
\dev{२१. गोलाध्यायः} & Ch. 21: Spherical Astronomy \\
\bottomrule
\end{tabular}
\end{center}

\vspace{1em}
\noindent This bilingual edition presents the Sanskrit original in Devanagari alongside an English translation with IAST (International Alphabet for Sanskrit Transliteration). The left column contains the Sanskrit text; the right column contains the English translation.

\newpage

%% ========================================================================
%% INTRODUCTION
%% ========================================================================
\section*{Introduction / \dev{प्रस्तावना}}

\begin{minipage}[t]{0.47\textwidth}
\dev{%
यह खण्ड ब्रह्मगुप्ताचार्यकृत ब्राह्मस्फुटसिद्धान्तस्य अन्तिमभागः। अत्र ग्रहगणितस्य
सम्पूर्णतः, कुट्टकाध्यायस्य सम्पूर्णं, छायागणना, तथा गोलाध्यायस्य प्रारम्भः वर्तते।
कुट्टकाध्यायः विशेषेण प्रसिद्धः अस्ति, यतोहि अत्र रेखीय-डायोफाण्टाइन-समीकरणानि
$ax + by = c$ रूपस्य समाधानं दत्तम्। अपि च शून्येन सह गणितस्य नियमाः,
ऋण-धन-संकलन-विधयः च उक्ताः।
}
\end{minipage}%
\hfill\vline\hfill%
\begin{minipage}[t]{0.47\textwidth}
This final chunk of the \textit{Brāhmasphuṭasiddhānta} of Brahmagupta (c. 598--668 CE) completes the monumental treatise composed in 628 CE. It contains the concluding portion of planetary computations (Grahagaṇita), the celebrated \textit{Kuṭṭakādhyāya} (Chapter on the Pulverizer), shadow calculations (Chāyādhyāya), and the opening of spherical astronomy (Golādhyāya).

The \textit{Kuṭṭakādhyāya} is especially renowned as it provides systematic solutions to linear Diophantine equations of the form $ax + by = c$, operations with zero and negative numbers, and the ``square-nature'' (\textit{vargaprakṛti}) problem equivalent to Pell's equation $Nx^2 + 1 = y^2$.
\end{minipage}

\vspace{1.5em}

%% ========================================================================
%% CHAPTER 17
%% ========================================================================
\newpage
\section*{\dev{१७. सृङ्गोत्स्त्युतराध्यायः} / Chapter 17: Advanced Planetary Computations}

\begin{center}
\textit{(Original pages 1136--1146; covering mean and true longitudes, parilekha drawings)}
\end{center}

\longline

\subsection*{\dev{प्रश्नकथनम्} / Problem Statement}

\bilingualverse{
द्विजार्कसनवरसेन्दूर्जिनतिथिविषयाग्रहार्थचापानाम् ।
अर्धज्याखण्डानि ज्यामुकैर्नैक्यं स भोग्यफलम् ॥१५॥}{%
\textbf{15.} The half-chord segments (\textit{ardhajyā-khaṇḍa}) of the arcs for the purpose of
[computing] the Sun, the Moon, Venus, Mars, Jupiter, Saturn, and the nodes of the Moon ---
the sum of these chord-segments is the \textit{bhogya-phala} [incremental motion].
}

\bilingualverse{
गतभोग्यखण्डकान्तर दलविकलविधा च्छुतैर्वर्भिराप्तैः ।
तच्छृतिदलं युतेन भोग्यादुनार्धिकं भोग्यम् ॥१६॥}{%
\textbf{16.} The difference of the \textit{gata-bhogya-khaṇḍa} [past and future segments],
divided by the minutes of arc, multiplied by the elapsed [time in] minutes ---
that, combined with half the difference, [added to] the \textit{bhogya} [motion], is the mean \textit{bhogya}.
}

\begin{engcommentary}
\textbf{Mathematical note:} The \textit{bhogya-khaṇḍa} method computes the incremental motion of a planet by dividing the orbital arc into segments. If $k_1, k_2, \dots, k_n$ are the chord-segments for each planet's arc, the \textit{bhogya-phala} is $\sum k_i$. The mean motion is then interpolated using the proportion of elapsed time. Brahmagupta uses a \textit{trijyā} (radius) of 150, whereas Bhāskara II later adopted 120 in the \textit{Siddhānta-śiromaṇi}.
\end{engcommentary}

\subsection*{\dev{परिलेखकथनम्} / Parilekha (Planetary Diagrams)}

\bilingualverse{
रविवसतिगतः पञ्चयुग्मं वपांशि पितामहोद्दिष्टानि ।
अर्धमासार्द्धयातद्विमर्मसैरमर्मस्नपष्टयाऽऽत्म ॥३॥}{%
\textbf{[Pañcasiddhāntikā 3.]} The Sun's [daily] motion is five [ghaṭikās]; the
pairs [of planets] and the minutes [are as] declared by Pitāmaha [Brahmā].
Half a month, half [the motion], the nodes, the apogees, and the quantities ---
\textit{these} are established by calculation.
}

\begin{engcommentary}
This verse from the \textit{Pañcasiddhāntikā} of Varāhamihira (c. 505--587 CE) cites the
\textit{Paitāmaha-siddhānta} --- the oldest of the five astronomical systems summarized in that work.
Brahmagupta critically examines these earlier parameters, noting discrepancies with observational data.
The phrase \textit{pita-mahoddhiṣṭāni} (declared by Pitāmaha) refers to the legendary divine origin
ascribed to this siddhānta.
\end{engcommentary}

\subsection*{\dev{विशेषकथनम्} / Special Rules}

\bilingualverse{
आद्धतकालयोर्मध्ये कालक्षे पोष्टदाश्रयः ।
प्रज्वलज्ज्वलनाकारः सर्वकर्मसु गर्हितः ॥}{%
Between the two times of sunrise and sunset, the time when the [shadow] increases
and takes support [of the gnomon], blazing with the form of fire ---
[it is] condemned in all rituals [as inauspicious].
}

\bilingualverse{
एकायनगां यावत् क्षेत्रमण्डलान्तरम् ।
सुम्बवस्तावदेवास्य सर्वकर्म विनाशकृत् ॥}{%
As much as the difference between the circles of the [ecliptic and equatorial] paths,
so much is the \textit{saṃbādha} [obstruction] --- this destroys all rituals.
}

\begin{engcommentary}
The \textit{pātāntara-kāla} (intermediate time of the node passage) is declared inauspicious
for all rituals including bathing, charity, recitation, \textit{śrāddha}, and \textit{homa}.
When the Sun and Moon have equal brightness near the \textit{graha-sandhi} (planetary conjunction),
the \textit{pāta} (node passage) occurs twice in a short time; but when the Sun is near apogee,
the equality of brightness persists for a longer duration.
\end{engcommentary}

\vspace{1em}
\longline


%% ========================================================================
%% CHAPTER 18: KUTTAKADHYAYA
%% ========================================================================
\newpage
\section*{\dev{१८. कुट्टकाध्यायः} / Chapter 18: The Pulverizer (Kuṭṭaka)}

\begin{center}
\textit{(Original pages 1146--1211; PDF pages 292--357)}
\\[0.3em]
\textit{The most celebrated chapter: linear Diophantine equations, operations with zero,\\
and the ``square-nature'' (vargaprakṛti) problem}
\end{center}

\longline

\subsection*{\dev{कुट्टकार्थप्रयोजनम्} / Purpose of the Kuṭṭaka}

\begin{minipage}[t]{0.47\textwidth}
\dev{%
\textbf{[Introduction.]}
कुट्टकशब्देन अत्र रेखीय-डायोफाण्टाइन-समीकरणस्य समाधानविधिः उच्यते।
यथा $ax + by = c$ रूपस्य समीकरणस्य मूलानयनम्। अत्र भाज्यं $a$,
भाजकः $b$, शेषं $c$, गुणकः $x$, तथा कोटिः $y$। ब्रह्मगुप्तः प्रथमः
आचार्यः यः शून्येन सह गुणनभागहारौ नियमेन सह ददाति।
}
\end{minipage}%
\hfill\vline\hfill%
\begin{minipage}[t]{0.47\textwidth}
\textbf{[Introduction.]} The term \textit{kuṭṭaka} (``pulverizer'') refers here to the method of
solving linear Diophantine equations of the form $ax + by = c$. Here $a$ is the
\textit{bhājya} (dividend), $b$ is the \textit{bhājaka} (divisor), $c$ is the \textit{śeṣa} (remainder),
$x$ is the \textit{guṇaka} (multiplier), and $y$ is the \textit{koṭi} (quotient).

Brahmagupta was the first to give systematic rules for arithmetic operations
involving zero and negative numbers, as well as the general solution to the
pulverizer equation.
\end{minipage}

\vspace{1.5em}

\bilingualverse{
द्वयं न शकेन्‍द्रकालं पञ्चभिरुद्भूतं शेषवर्षयुगम् ।
दिगुणामर्धसितां कुर्याद् भुग्रां तद् द्वयं दयात् ॥१४॥}{%
\textbf{14.} The two [quantities] --- the \textit{śaka-indra-kāla} [elapsed years of the Śaka era]
and the \textit{śeṣa-varṣa-yuga} [residue of years in the yuga] increased by five ---
make them half the \textit{ardha-sita} [half the deficit]; the \textit{bhugrā} [earth-number]
for that pair [is found] by the pulverizer.
}

\bilingualverse{
संक्षिप्तं द्वि त्रि गतौ तिथिभूमान्‍नवाहतैष्यकैः ।
दिग्रसभागाः सप्तभिन्नं शशिभं धनिग्रहात्म ॥१५॥}{%
\textbf{15.} The \textit{saṃkṣipta} [reduced quantity], two, three, the motions,
the \textit{tithi-bhū} [lords of the tithis], multiplied by nine, the \textit{eṣya} [desired quantity] ---
the \textit{digra-sabhāga} [directional parts] divided by seven, the Moon's [motion],
[these are] the wealth of the planet's self.
}

\begin{engcommentary}
\textbf{On the astronomical setting:} In the \textit{kuṭṭaka} method as applied to astronomy,
the dividend (\textit{bhājya}) typically represents the number of days in a yuga (e.g., 1577917500
civil days), the divisor (\textit{bhājaka}) represents the number of planetary revolutions in that yuga,
and the remainder (\textit{śeṣa}) represents the elapsed portion. The method finds the multiplier
(\textit{guṇaka}) such that the congruence is satisfied, yielding the number of elapsed days
corresponding to a given planetary longitude.
\end{engcommentary}

\subsection*{\dev{कुट्टकविधिः} / The Pulverizer Method}

\bilingualverse{
भाज्यं भाजकशेषेण यावदेकं करणी भवेत् ।
प्रथमं रूपयुक्तं ततो द्वितीयं करणी भवेत् ॥४०॥}{%
\textbf{40.} The \textit{bhājya} [dividend, $a$] is [divided] by the \textit{bhājaka} [divisor, $b$]
together with the \textit{śeṣa} [remainder, $c$] until unity becomes the operation.
The first [quotient] is united with one [rupas]; then the second operation becomes [clear].
}

\begin{minipage}[t]{0.47\textwidth}
\dev{%
\textbf{[Hindi commentary excerpt:]}
सू. भाषा.—रूपकुट्टः कि विधिः इष्टया इष्टकर्णपद्धतिः। इष्टा यैका तया वैश्यो-
द्वितीयैः रूपवर्गैस्तेन वैधानमनेकासां यो रूपवर्गैस्तेनोनाया यथपदं तेन रूपार्धा-
पृथक् युता नितादि तदर्थं च कार्यं। तत्र प्रथमार्धिगोर्थं रूपार्धा कल्प्यानि। ततो
स्यदन्तरार्थं द्वितीयं मूलस्यैका करणी भवति। एवमसंकुलेमूलानयनं कार्यम्॥

अत्रोदाहरण म. म. सूर्याकारदिवेगुकम्।

चतुर्दर्शनद्वितीयाः पङ्क्तयो विश्वमजिताः। तं राशिं शीघमाचक्ष्व
यदि जानासि कुट्टकम्॥ अत्र ३४ हरस्य शेषम् $= २$। तथा १३ हरस्य शेषम् $= १०$।
अतोऽर्धिकशेषहरः $= १३$। अल्पशेषहरः $= ३४$। अनेनार्धिकशेषहरे भक्तं शेषम्
$= १३$, ततः परस्परभजनेन जाता वल्ली उपान्तिमेन स्वोच्चं हृतोनयेन युते तदन्त्यं
}
\end{minipage}%
\hfill\vline\hfill%
\begin{minipage}[t]{0.47\textwidth}
\textbf{[English translation of commentary:]}
The commentary explains the pulverizer method as follows: ``The pulverizer
with unity (\textit{rūpa-kuṭṭaka}) is the method by which, starting from the desired [quantity],
the root-quantity is found. The first quotient, combined with unity (\textit{rūpa}),
gives the basis; then for the purpose of the second [step], unity becomes the
operative quantity. In this manner, roots are extracted from the non-aggregate
[quantities].

\textbf{Example:} `The fourteen-second rows [of a formation], O conqueror of the world,
tell that quantity quickly if you know the pulverizer.' Here, for divisor 34 the remainder is 2;
for divisor 13 the remainder is 10. Hence the greater remainder-divisor is 13, the lesser is 34.
Dividing the greater by the lesser gives remainder 13; then by mutual division the chain (\textit{vallī})
is produced. The penultimate [quotient], multiplied by its own higher [term] and united [with
the lower] --- this is the final [step] of the pulverizer method.''
\end{minipage}

\vspace{1.5em}

\begin{engcommentary}
\textbf{Algorithmic summary:} Given $ax + by = c$ where $\gcd(a,b) \mid c$:
\begin{enumerate}
\item Apply the Euclidean algorithm to $(a, b)$ to obtain quotients $q_1, q_2, \dots, q_n$.
\item Construct the \textit{vallī} (chain) of partial quotients.
\item Starting from the bottom with 1, work upward: $x_0 = 1$, $x_1 = q_n$,
$x_{k+1} = q_{n-k} \cdot x_k + x_{k-1}$.
\item The penultimate value gives the multiplier $x$; then $y = (c - ax)/b$.
\end{enumerate}
This is algebraically equivalent to the extended Euclidean algorithm.
\end{engcommentary}

\subsection*{\dev{शून्यगणितम्} / Operations with Zero}

\bilingualverse{
धनर्णशून्यानां संकलनं व्यवकलनं च ।
गुणने करणसूत्रकथनं भागहारे च ॥}{%
\textbf{[On zero:]} The aggregation and diminution of positive, negative, and zero
[quantities]; the rule for multiplication; and for division --- [are declared].
}

\begin{minipage}[t]{0.47\textwidth}
\dev{%
\textbf{[Addition and Subtraction:]}
ऋणे ऋणं योजयेत् धनं धने। शून्ये शून्यं योजयेत्।
धनात् ऋणं व्योजयेत्, ऋणात् धनं व्योजयेत्।
शून्येन युतं राशिः स्वेनैव रूपेण स्थितः।
}
\end{minipage}%
\hfill\vline\hfill%
\begin{minipage}[t]{0.47\textwidth}
\textbf{[Addition and Subtraction:]}
A negative [quantity] joined with a negative becomes more negative;
a positive with a positive [becomes] more positive;
zero joined with zero [remains] zero.
From a positive, a negative is subtracted;
from a negative, a positive is subtracted.
A quantity united with zero remains in its own form.
\end{minipage}

\vspace{1.5em}

\begin{minipage}[t]{0.47\textwidth}
\dev{%
\textbf{[Multiplication:]}
शून्येण गुणितं राशिः शून्यं भवति।
शून्येण भक्तं राशिः शून्यं भवति।
अव्यक्तसंकलितं यावद् इष्टतयोः करणसूत्रकथनम्।
अव्यक्तगुणने सूत्रकथनम्।
अव्यक्तमानानयनार्थं कथनम्।
}
\end{minipage}%
\hfill\vline\hfill%
\begin{minipage}[t]{0.47\textwidth}
\textbf{[Multiplication:]}
A quantity multiplied by zero becomes zero.
[But] a quantity divided by zero becomes zero [\textit{śūnya} --- Brahmagupta's
controversial rule: he treats $a/0 = 0$, later corrected by Bhāskara II].

The rule for the aggregation of unknowns (\textit{avyakta}), as much as desired,
and the rule for the multiplication of unknowns, and the determination
of the measure of unknowns --- [are given].
\end{minipage}

\vspace{1.5em}

\begin{engcommentary}
\textbf{Historical note on zero:} Brahmagupta (628 CE) was the first mathematician in history
to establish rules for arithmetic operations with zero as a number. His rules:
\begin{itemize}
\item $a + 0 = a$; $a - 0 = a$; $0 + 0 = 0$; $0 - 0 = 0$
\item $a \times 0 = 0$; $0 \times 0 = 0$
\item $a / 0 = 0$ (\textit{this rule is incorrect; division by zero is undefined})
\item $0 / a = 0$ for $a \neq 0$
\end{itemize}
He also gave rules for operations with negative numbers (\textit{ṛṇa}):
$(-a) \times (-b) = +ab$, $(-a) \times (+b) = -ab$, etc.
These rules were revolutionary for the 7th century.
\end{engcommentary}

\subsection*{\dev{वर्गप्रकृतिः} / The Square-Nature (Vargaprakṛti) Problem}

\bilingualverse{
समद्विबाहुत्रिभुजे भुजयोर्बर्गेण ।
करणीयोगान्तरे गुणनकथनम् ॥}{%
\textbf{[Square-nature:]} In an isosceles triangle, the sum of the squares of the sides ---
[and] the product of the sum and difference of the \textit{karaṇīs} [surds].
}

\bilingualverse{
वर्गसमीकरणाकथनम् ।
वर्गसमीकरणोद्धरणमानयनम् ॥}{%
The statement of the square-equation [Pell's equation];
the elevation of the root of the square-equation [its solution].
}

\begin{minipage}[t]{0.47\textwidth}
\dev{%
\textbf{[Sanskrit commentary on vargaprakṛti:]}
अत्र `वर्गप्रकृतिः' नाम समीकरणं $Nx^2 + 1 = y^2$ रूपस्य।
अत्र $N$ = प्रकृतिः (गुणकः), $x$ = ह्रस्वमूलम्, $y$ = दीर्घमूलम्।
ब्रह्मगुप्तः अत्र समाधानाय `भावना' नाम विधिं ददाति ---
यदा द्वयोः ह्रस्व-दीर्घमूलयोः युग्मेभ्यः नूतनयुग्मं प्राप्तुं शक्यते।
}
\end{minipage}%
\hfill\vline\hfill%
\begin{minipage}[t]{0.47\textwidth}
\textbf{[English:]}
The ``square-nature'' (\textit{vargaprakṛti}) refers to the equation $Nx^2 + 1 = y^2$,
known in modern mathematics as Pell's equation. Here $N$ is the \textit{prakṛti} (multiplier),
$x$ is the \textit{hrasva-mūla} (lesser root), and $y$ is the \textit{dīrgha-mūla} (greater root).

Brahmagupta gives the \textit{bhāvanā} (composition) method by which, from two pairs of
lesser and greater roots, a new pair can be generated. If $(x_1, y_1)$ and $(x_2, y_2)$
are solutions for the same $N$, then:
\[x' = x_1 y_2 + x_2 y_1, \quad y' = N x_1 x_2 + y_1 y_2\]
is also a solution. This is the \textit{samāsa-bhāvanā} (additive composition).
\end{minipage}

\vspace{1.5em}

\begin{engcommentary}
\textbf{The bhāvanā composition:} Brahmagupta's \textit{prakṛti-samāsa} (composition of roots)
is one of the most remarkable achievements in early algebra. Given two solutions
$(x_1, y_1)$ and $(x_2, y_2)$ of $Nx^2 + k_1 = y^2$ and $Nx^2 + k_2 = y^2$ respectively,
then:
\[
x_3 = x_1 y_2 + x_2 y_1, \quad y_3 = N x_1 x_2 + y_1 y_2
\]
satisfies $Nx_3^2 + k_1 k_2 = y_3^2$. When $k_1 = k_2 = 1$, this generates an infinite
family of solutions from a single minimal solution. For example, with $N = 92$,
the minimal solution is $x = 120$, $y = 1151$, since $92 \cdot 120^2 + 1 = 1151^2$.
\end{engcommentary}

\subsection*{\dev{विलोमगतिकथनम्} / Reverse Motion (Viloma-gati)}

\bilingualverse{
राश्यर्धं स्तच्छेषैश्चैवं युक्तार्धमासादिनहीनैः ।
तच्छेषैश्च युगगतं यः कथयति कुट्टकज्ञः सः ॥१४॥}{%
\textbf{14.} Half the \textit{rāśi} [sign], increased by its \textit{śeṣa} [residue],
and diminished by the \textit{ardha-māsa} [half-month] etc.;
[and] united with its \textit{śeṣa} --- he who states the \textit{yuga-gata}
[planetary position for the yuga] is the knower of the pulverizer.
}

\begin{engcommentary}
\textbf{On reverse computation:} The \textit{viloma-gati} method computes backwards from a given
planetary longitude to find the corresponding elapsed time (\textit{ahargaṇa}). Given a planet's
mean longitude, one first determines the residual motion within the current sign (\textit{rāśi}),
then applies the pulverizer to find how many days have elapsed since the epoch.
This is the inverse problem of the forward computation.
\end{engcommentary}

\vspace{1em}
\longline


%% ========================================================================
%% CHAPTER 19: SHANKUCCHAYADI-JNANADHYAYA
%% ========================================================================
\newpage
\section*{\dev{१९. शङ्कुच्छायादिज्ञानाध्यायः} / Chapter 19: Gnomon and Shadow}

\begin{center}
\textit{(Original pages 1265--1314; PDF pages 357--406)}
\\[0.3em]
\textit{Shadow computations using the gnomon (\dev{शङ्कु}), time determination,\\
and related astronomical problems}
\end{center}

\longline

\subsection*{\dev{प्रथमप्रश्नकथनम्} / First Problem Statement}

\bilingualverse{
हरेण भक्तः शुध्यति स राशिः कः।
प्रथमं गुणाहरयोर्महतामपवर्तेनमानीयते।}{%
Divided by the divisor [\textit{hara}], what is that quantity [\textit{rāśi}] which becomes pure?
First, of the multiplier and divisor, the greater is brought [into consideration]
by means of the \textit{apavarta} [reduction].
}

\begin{minipage}[t]{0.47\textwidth}
\dev{%
\textbf{[Shadow problem:]}
परस्परं भक्तयोर्गुणाहरयोर्न ते यच्छेषं तेन भक्तो तौ (गुणाहरौ भवतः) द्वौ गुणा-
हरौ भवतः। ततो हृदयोर्गुणाहरयोः परस्परं भक्तयोर्लम्बान्यथोच्छ स्थापयानि,
पूर्वं प्रदर्शितकुट्टकनियमेन। इदं कर्म तावत्पर्यन्तं कर्तव्यं यावदन्ते रूप शेषं तिष्ठेत्।
तच्छेषं तथा केनापि गुणाकेन गुणितं रूपेण हीनं तच्छेषसंबन्धिहरेण भक्तं
शुध्यति। सत्यैवं गुणाकः पूर्वोच्छोच्छः स्थापितकलानामधः स्थाप्यः, तदन्तात्फलं
स्थाप्यम्।
}
\end{minipage}%
\hfill\vline\hfill%
\begin{minipage}[t]{0.47\textwidth}
\textbf{[English translation:]}
When two [numbers] divide each other mutually, that which [leaves] a remainder ---
by that they become two multiplier-divisor [pairs]. Then, having placed the quotients
obtained from the mutual division of the multiplier and divisor in the heart [i.e., in a column],
according to the pulverizer rule demonstrated previously. This operation is to be performed
until finally unity remains as the remainder. That remainder, multiplied by some multiplier
and diminished by unity, when divided by the divisor related to that remainder, becomes pure.
Having done so, the multiplier is to be placed below the previously stated quotients;
below that, the result is placed. In this manner the chain (\textit{vallī}) is produced.
\end{minipage}

\vspace{1.5em}

\subsection*{\dev{शङ्कुच्छायानयनम्} / Determination of Gnomon and Shadow}

\bilingualverse{
छायातो गृहादीनां च्छायानयनम् ।
इष्टगुहोच्चयको यदिति प्रश्नस्योत्तरसम्पादनम् ॥}{%
The determination of the shadow of houses etc. from the [Sun's] shadow.
The solution of the problem: ``If [the shadow] is equal to the height of the desired cave'' ---
[its computation is given].
}

\begin{engcommentary}
\textbf{On shadow problems:} The gnomon (\textit{śaṅku}) is a vertical staff of standard height
(typically 12 \textit{aṅgulas}). Its shadow (\textit{chāyā}) cast by the Sun is used to determine:
\begin{enumerate}
\item The Sun's altitude: $\sin(\text{altitude}) = \text{śaṅku} / \sqrt{\text{śaṅku}^2 + \text{chāyā}^2}$
\item The time of day from the shadow's length
\item The dimensions of objects from their shadows
\item The \textit{palabha} (equinoctial shadow) for latitude determination
\end{enumerate}
If the shadow equals the height of a cave (i.e., the Sun's altitude is 45\textdegree),
the horizontal distance from the gnomon to the cave equals the cave's height.
\end{engcommentary}

\subsection*{\dev{उच्छ्रितिकथनम्} / Height Determination}

\bilingualverse{
उच्छ्रितिकथनं समीपस्थस्योच्चयस्य ज्ञानार्थं दूरस्थस्य च ।
प्रकाशस्य गतिकथनं द्वयोरन्तरालस्य मानं च ॥}{%
The determination of the height of a nearby elevated object and of a distant one;
and the statement of the motion of light [rays], and the measure of the interval between the two.
}

\begin{minipage}[t]{0.47\textwidth}
\dev{%
\textbf{[Height from shadow --- example:]}
गृहपुरुषान्तरसलिले यो हटचैत्यादि प्रश्नोत्तरकथनम्।
बीजय गुहायां सलिले प्रसारे ति प्रश्नोत्तरकथनम्।

एकस्य गृहस्य छाया द्वितीयमात्रान्तरविधानेन त्रिप्रश्नोत्तरकथनम्।
छायातो गृहादीनां च्छायानयनम्।
}
\end{minipage}%
\hfill\vline\hfill%
\begin{minipage}[t]{0.47\textwidth}
\textbf{[Height from shadow --- examples:]}
\textbf{(a)} The determination of the answer to the problem of a house, a man,
the interval, and water --- concerning a \textit{haṭa-caitya} [pillar shrine] etc.

\textbf{(b)} The problem concerning a seed in a cave, water, and spreading ---
its solution is stated.

\textbf{(c)} By the method of double-difference, the answers to three problems
from the shadow of one house are given.

\textbf{(d)} From the shadow, the determination of shadows of houses etc.
\end{minipage}

\vspace{1.5em}

\begin{engcommentary}
\textbf{The double-difference method:} For determining heights from shadows:
If two observations of a gnomon and its shadow are made at different times or positions,
the height $h$ of an object and its horizontal distance $d$ can be found from:
\[h = \frac{s_2 \cdot g}{s_2 - s_1} \cdot \frac{d_2 - d_1}{g}, \quad d = \frac{h \cdot s_1}{g}\]
where $g$ is the gnomon height, $s_1, s_2$ are shadow lengths at two positions,
and $d_1, d_2$ are the distances of the observation points from the object's base.
This is equivalent to similar-triangle methods in modern trigonometry.
\end{engcommentary}

\vspace{1em}
\longline

%% ========================================================================
%% CHAPTER 20: CHHADICHCHAYADHYAYA
%% ========================================================================
\newpage
\section*{\dev{२०. छद्दिच्छायाध्यायः} / Chapter 20: Advanced Shadow Problems}

\begin{center}
\textit{(Original pages 1316--1320; PDF pages 406--411)}
\\[0.3em]
\textit{Extended shadow computations: obscuration, refraction,\\
and celestial sphere applications}
\end{center}

\longline

\bilingualverse{
छद्दिराशीनामहोरात्रबलानाम् ।
राश्युदयानामसमानद्विकथनम् ॥}{%
Of the \textit{chad-dirāśi} [obscured signs], the day-and-night strengths ---
the rise of the signs, the statement of their unequal pairs.
}

\bilingualverse{
चरागयोः संस्थानकथनम् ।
शङ्कृह्ग्रयोः संस्थानकथनम् ॥}{%
The position of the \textit{carāgra} [ascensional difference] is stated.
The position of the gnomon and shadow is stated.
}

\bilingualverse{
हग्गोलस्य हर्याहरयतं लम्बनावनतिरुपतो च कारणदर्शनम् ।
परलम्बनावनतिकथनम् ॥}{%
Of the \textit{hag-gola} [terrestrial sphere], the cause of the \textit{lambana}
[parallax] and \textit{avanati} [depression] by means of the \textit{haryāhara}
[division-operation] is shown. The statement of the \textit{para-lambana}
[secondary parallax] and \textit{avanati}.
}

\begin{minipage}[t]{0.47\textwidth}
\dev{%
\textbf{[Hindi commentary excerpt:]}
प्रकारान्तरेण तयोः संस्थाने शङ्कुलस्य च कथनम्।
हक्कर्मकथनम्।
ग्रहसौगोल्योः स्थिरकुट्टकथनम्।
ग्रहाणां चलकुट्टकथनम्।
अध्यायोपसंहारः ॥

अत्र छायायाः गणना दीपशिखोच्चयच्छुक्लान्तरभूमिज्ञानार्थं दर्शिता।
प्रकाश-परावर्तन-विषयेऽपि किञ्चिदुक्तम्।
}
\end{minipage}%
\hfill\vline\hfill%
\begin{minipage}[t]{0.47\textwidth}
\textbf{[English:]}
By another method, the position of those two [gnomon and shadow] and
of the \textit{śaṅkula} [combined gnomon-shadow] is stated.

The statement of the \textit{hak-karma} [obscuration computation].
The fixed pulverizer of the planets and the celestial sphere.
The moving pulverizer of the planets.
The conclusion of the chapter.

Here the computation of shadows is shown for the knowledge of the
height of a lamp-flame, the boundary of light, and the Earth.
Something is also said concerning the reflection of light.
\end{minipage}

\vspace{1.5em}

\begin{engcommentary}
\textbf{On the \textit{chad-dichāyā}:} This chapter extends the shadow calculations of Chapter 19
to more complex astronomical situations. The \textit{chad-dirāśi} refers to zodiacal signs
that are ``obscured'' (below the horizon), and the computation of their day/night strengths
is relevant to astrological applications. The \textit{carāgra} (ascensional difference) measures
the unequal rising times of signs at different latitudes, computed from:
\[\text{carāgra} = \arcsin\left(\frac{\tan(\phi) \cdot \tan(\delta)}{\cos(\alpha)}\right)\]
where $\phi$ is the observer's latitude, $\delta$ is the celestial declination,
and $\alpha$ is the right ascension.
\end{engcommentary}

\vspace{1em}
\longline

%% ========================================================================
%% CHAPTER 21: GOLADHYAYA
%% ========================================================================
\newpage
\section*{\dev{२१. गोलाध्यायः} / Chapter 21: Spherical Astronomy (Golādhyāya)}

\begin{center}
\textit{(Original pages 1323--1416; PDF pages 411--420+)}
\\[0.3em]
\textit{The opening of the Spherical Astronomy Book: armillary spheres,\\
celestial and terrestrial globes, astronomical instruments}
\end{center}

\longline

\subsection*{\dev{आरम्भप्रयोजनकथनम्} / Purpose of the Chapter}

\bilingualverse{
गोलप्रशंसाकथनम् ।
स्वगोलयथाने कारणम् ॥}{%
The praise of the celestial sphere (\textit{gola}) is stated.
The cause of the proper position of one's own sphere [is given].
}

\bilingualverse{
गतिगोलयोः प्रशंसाकथनम् ।
यन्त्रारम्भप्रयोजनकथनम् ॥}{%
The praise of the sphere of motion [is stated].
The purpose of the commencement of instruments [is stated].
}

\begin{minipage}[t]{0.47\textwidth}
\dev{%
\textbf{[Contents of Golādhyāya:]}
तन्त्राणां प्रतिपादनम् ।
सलिलादीनां प्रयोजनकथनम् ।
धनुषां च कथनम् ।
सूर्यादिमुखे यन्त्रधारणम् ।
इष्टघटिकायाः धनुषरचनरूपकथनम् ।
परीक्षाघट्यानयनम् ।
यन्त्रेण नतोन्नतकालज्ञानम् ।
धनुष्यन्त्रे विशेषकथनम् ।
}
\end{minipage}%
\hfill\vline\hfill%
\begin{minipage}[t]{0.47\textwidth}
\textbf{[Contents of the Golādhyāya:]}
\begin{enumerate}
\item The establishment of instruments (\textit{tantra})
\item The use of water etc. [in instruments]
\item The statement concerning arcs (bows)
\item The bearing of instruments toward the Sun etc.
\item The determination of the form of the arc for a desired ghaṭikā
\item The determination of the testing ghaṭikā
\item The knowledge of time by means of the instrument's depression and elevation
\item Special rules concerning the bow-instrument
\end{enumerate}
\end{minipage}

\vspace{1.5em}

\subsection*{\dev{भूगोलसंस्थानकथनम्} / Arrangement of the Terrestrial Sphere}

\bilingualverse{
भूगोलसंस्थानकथनम् ॥
दैवासुरसंस्थानवर्णनम् ॥}{%
The arrangement of the terrestrial sphere is stated.
The description of the arrangement of the divine and demoniacal [cities].
}

\begin{minipage}[t]{0.47\textwidth}
\dev{%
\textbf{[The celestial cities:]}
देवदैत्ययोर्मेषचक्रमरगत्यवस्थापनम् ॥
चक्रभ्रमणव्यवस्थाकथनम् ॥
दैनानीनां रविभ्रमणार्थदिवसानाम् ॥
देवदैत्ययोः राशिसंस्थानम् ॥
देवदैत्ययोः पितुमानस्यैव दिनप्रमाणानिरूपणम् ॥
}
\end{minipage}%
\hfill\vline\hfill%
\begin{minipage}[t]{0.47\textwidth}
\textbf{[The celestial cities:]}
The establishment of the Aries-cycle and the [rate of] motion for the divine and demoniacal [cities].
The statement of the arrangement of the wheel's revolution.
Of the days of the gods, for the purpose of the Sun's revolution.
The zodiacal arrangement of the divine and demoniacal [cities].
The determination of the measure of the day of the father [Brahmā] for the divine and demoniacal.
\end{minipage}

\vspace{1.5em}

\begin{engcommentary}
\textbf{The celestial geography:} Brahmagupta describes the traditional Indian
celestial geography. At the North Pole is \textit{Meru}, the abode of the gods (\textit{deva});
at the South Pole is \textit{Vadavā-mukha}, the realm of the demons (\textit{asura/daitya}).
The central terrestrial sphere (\textit{bhū-gola}) has the observer's horizon,
the prime meridian (through Ujjain), and the equator.

The divine day (one day of Brahmā) equals 1000 \textit{caturyugas} (4,320,000,000 solar years).
The solar year is defined by the Sun's return to the same fixed star --- the
\textit{nakṣatra-vatsara} (sidereal year). The \textit{divya-dina} (divine day) of 360 solar years
is used for calculating planetary mean motions.
\end{engcommentary}

\subsection*{\dev{निरक्षदेशान्तरयोजनकथनम्} / Prime Meridian and Longitude Difference}

\bilingualverse{
भूगोल लग्नावलयोः संस्थानम् ॥
निरक्षदेशान्तरयोजनकथनम् ॥}{%
The arrangement of the terrestrial sphere and the \textit{lagna-valaya}
[horizon circles] is [stated].
The statement of the \textit{nirakṣa-deśāntara} [equatorial longitude difference] in \textit{yojanas}.
}

\bilingualverse{
क्षुद्रद्योः कथानिरूपणम् ॥}{%
The description of the two \textit{kṣudra-dyau} [small celestial circles].
}

\begin{minipage}[t]{0.47\textwidth}
\dev{%
\textbf{[On the deshantara:]}
निरक्षदेशः उज्जयिनी-क्षेत्रम्, यत्र रेखा (याम्योत्तर-वृत्तम्) पूर्णतः
समपर्यायं ग्रहैः भ्रमति। अन्यत् क्षेत्रं प्रति देशान्तरं योजनैः मीयते ---
प्रत्येकं योजन-चतुष्टयस्य एका लघु-घटिका भवति।
तदनुसारं देशान्तर-घटिकाभिः ग्रहाणाम् उदय-कालः परिवर्तते।
}
\end{minipage}%
\hfill\vline\hfill%
\begin{minipage}[t]{0.47\textwidth}
\textbf{[On the \textit{deśāntara}:]}
The prime meridian (\textit{nirakṣa-deśa}) passes through Ujjain, where the
celestial equator (the east-west circle) is traversed equally by all planets.
For any other location, the longitude difference (\textit{deśāntara}) is measured in
\textit{yojanas} --- for every four \textit{yojanas}, there is one \textit{laghu-ghaṭikā}
(small ghaṭikā, about 4 minutes of time).

Accordingly, the rising times of planets vary by the \textit{deśāntara} ghaṭikās.
The \textit{deśāntara-kāla} (longitude-difference time) is computed as:
\[\text{deśāntara} = \frac{\text{yojanas from Ujjain}}{4}\]
in \textit{ghaṭikās} (1 ghaṭikā = 24 minutes).
\end{minipage}

\vspace{1.5em}

\begin{engcommentary}
\textbf{The \textit{deśāntara} formula:} Brahmagupta gives the standard Indian rule for
computing longitude difference from Ujjain (the prime meridian): every 4 yojanas
east or west corresponds to a 1-ghaṭikā (24-minute) difference in local time.
Since the Earth's circumference was taken as approximately 5000 yojanas,
this gives 1250 ghaṭikās for a full circle --- close to the actual 60 $\times$ 24 = 1440
minutes, though the \textit{yojana} was not a precisely standardized unit.

The \textit{deśāntara} is essential for computing the local time of planetary risings,
setting of the \textit{lagna} (ascendant), and the timing of eclipses at different longitudes.
\end{engcommentary}

\vspace{1em}
\longline


%% ========================================================================
%% explanatory notes AND APPENDIX
%% ========================================================================
\newpage
\section*{Notes / \dev{सम्पादकीय टीका}}

\begin{minipage}[t]{0.47\textwidth}
\dev{%
\textbf{[संस्कृत टीका:]}
इयं सम्पादिता संस्करणं ब्रह्मगुप्ताचार्यकृत ब्राह्मस्फुटसिद्धान्तस्य
पूर्णतमभागं प्रतिनिधित्वं करोति। अस्मिन् ग्रन्थे द्वाविंशत्यधिकायाः
(२४) अध्यायाः सन्ति --- दशसु (१०) अध्यायेषु गणितं, चतुर्दशसु (१४)
अध्यायेषु ज्योतिषं च।

कुट्टकाध्यायः (अध्यायः १८) विशेषेण प्रसिद्धः अस्ति, यतोहि अत्र
$ax + by = c$ रूपस्य समीकरणस्य सामान्यसमाधानं दत्तम्।
अस्य पश्चात् भास्कराचार्येण (१२वें शताब्दी) `बीजगणित'-ग्रन्थे
एतदेव विस्तरेण वर्णितम्।
}
\end{minipage}%
\hfill\vline\hfill%
\begin{minipage}[t]{0.47\textwidth}
\textbf{[English Notes:]}
This edited edition represents the final portion of the \textit{Brāhmasphuṭasiddhānta}
of Ācārya Brahmagupta (c. 598--668 CE). The complete treatise consists of 24 chapters ---
10 on mathematics and 14 on astronomy.

The \textit{Kuṭṭakādhyāya} (Chapter 18) is especially celebrated, as it provides the general
solution to the linear Diophantine equation $ax + by = c$. This was later elaborated by
Bhāskara II (12th century) in his \textit{Bījagaṇita}.
\end{minipage}

\vspace{2em}

\subsection*{1. Manuscript Sources / \dev{हस्तलेखस्रोतांसि}}

\begin{minipage}[t]{0.47\textwidth}
\dev{%
मूलपाठः सुधाकर-द्विवेदीनां सम्पादनात् (बenares संस्करण, १९०२)
प्राप्तः। अपि च रामस्वामी-आयङ्गार-महोदयस्य `ब्राह्मस्फुटसिद्धान्तम्'
(मद्रास्, १९३४) इति संस्करणं संदर्भितम्।
}
\end{minipage}%
\hfill\vline\hfill%
\begin{minipage}[t]{0.47\textwidth}
The base text is derived from the edition of Sudhākara Dvivedī (Benares, 1902).
The edition of R. Ramaswami Aiyangar \textit{Brāhmasphuṭasiddhāntam} (Madras, 1934)
has also been consulted for reference.
\end{minipage}

\vspace{1.5em}

\subsection*{2. Mathematical Significance / \dev{गणितीय महत्त्वम्}}

\begin{minipage}[t]{0.47\textwidth}
\dev{%
\textbf{[गणितीय विषयाः:]}
\begin{enumerate}
\item शून्येन सह गणितस्य नियमाः (प्रथमं विश्वे)
\item रेखीय-डायोफाण्टाइन-समीकरणस्य ($ax + by = c$) सामान्यसमाधानम्
\item `वर्गप्रकृतिः' ($Nx^2 + 1 = y^2$) समस्यायाः समाधानम्
\item `भावना' नाम सम्मिश्रणविधिः
\item चक्रिणः चतुर्भुजस्य क्षेत्रफलसूत्रम् (अन्यस्मिन् अध्याये)
\item ग्रहाणाम् उच्च-नीच-मन्द-शीघ्र-स्पष्टीकरणम्
\end{enumerate}
}
\end{minipage}%
\hfill\vline\hfill%
\begin{minipage}[t]{0.47\textwidth}
\textbf{[Mathematical achievements:]}
\begin{enumerate}
\item First rules for computing with zero as a number (in world history)
\item General solution to linear Diophantine equations $ax + by = c$
\item Solution to the ``square-nature'' problem $Nx^2 + 1 = y^2$
\item The \textit{bhāvanā} (composition) method for generating new solutions
\item Area formula for cyclic quadrilaterals (in other chapters)
\item Planetary equations: manda, śīghra, and sphuṭa corrections
\end{enumerate}
\end{minipage}

\vspace{1.5em}

\subsection*{3. Astronomical Constants / \dev{ज्योतिषीय स्थिरांकाः}}

\begin{minipage}[t]{0.47\textwidth}
\dev{%
\textbf{[ब्रह्मगुप्तीय युगपरिमाणम्:]}
\begin{tabular}{@{}l@{}r@{}}
एकस्मिन् युगे सौरवर्षाणि & ४,३२०,००० \\
चान्द्रमासाः & ५,३४,३३,३३६ \\
अर्धमासाः & १,०७,७७,७६,७३३ \\
सावनदिनानि & १,५७,७९,१७,८२८ \\
तिथयः & १,६०,२९,९८,००० \\
भवनयः & १,०७,७७,६,७३३ \\
रविभ्रमणानि & ४,३२०,००० \\
चन्द्रभ्रमणानि & ५७,७५,३३,३३६ \\
\end{tabular}
}
\end{minipage}%
\hfill\vline\hfill%
\begin{minipage}[t]{0.47\textwidth}
\textbf{[Brahmagupta's yuga constants:]}
\begin{tabular}{@{}lr@{}}
Solar years in a yuga & 4,320,000 \\
Lunar months & 53,43,3,336 \\
Half-months (ardha-māsa) & 10,77,76,733 \\
Civil days (sāvana-dina) & 15,77,91,7828 \\
Tithis & 16,02,98,000 \\
Nakṣatra days & 10,77,6,733 \\
Solar revolutions & 4,320,000 \\
Lunar revolutions & 57,75,33,336 \\
\end{tabular}

\vspace{1em}
\textit{Note:} These constants yield a sidereal year of approximately 365.25868 days,
close to the modern value of 365.25636 days (difference: about 0.00232 days $\approx$ 3.3 minutes).
\end{minipage}

\vspace{1.5em}

\subsection*{4. Transliteration and Translation Conventions / \dev{लिप्यन्तर-भाषान्तरविधिः}}

\begin{itemize}
\item Sanskrit text is given in Devanagari script using the IAST convention for
romanization in the English commentary.
\item Technical terms are given in IAST italics on first occurrence, with
translation in parentheses.
\item Verses are numbered according to the original text's verse numbering;
some verses from cited works (e.g., \textit{Pañcasiddhāntikā}) retain their original numbering.
\item Mathematical expressions use modern notation for clarity while preserving
the original algorithmic content.
\end{itemize}

\vspace{1.5em}

\subsection*{5. Key Commentators / \dev{प्रमुख टीकाकाराः}}

\begin{minipage}[t]{0.47\textwidth}
\dev{%
\textbf{[टीकाकाराः:]}
\begin{enumerate}
\item प्रथमटीका (अज्ञातकालीन)
\item बलभद्र (८वें शताब्दी)
\item अमराज (१२वें शताब्दी)
\item चन्द्रबहण (१३वें शताब्दी)
\\end{enumerate}
}
\end{minipage}%
\hfill\vline\hfill%
\begin{minipage}[t]{0.47\textwidth}
\textbf{[Key commentators:]}
\begin{enumerate}
\item The anonymous first commentary
\item Balabhadra (8th century)
\item Amarāja (12th century)
\item Candrahabaṇa (13th century)
\end{enumerate}
Each commentary tradition preserves different readings and interpretations of
the root text (\textit{mūla}), which modern editors must reconcile.
\end{minipage}

\vspace{2em}

%% ========================================================================
%% END MATTER
%% ========================================================================
\vfill

\begin{center}
\shortline

\vspace{1em}

{\Large\textbf{\dev{॥ इति ब्राह्मस्फुटसिद्धान्तस्य तृतीयभागस्य तृतीयखण्डः समाप्तः ॥}}}\\[0.8cm]

{\large\textit{End of Part 3, Chunk 3 of the Brāhmasphuṭasiddhānta}}\\[0.3cm]

{\normalsize This completes the Brahmasphutasiddhanta bilingual edition.}\\[0.3cm]

{\small Pages 1957--2512 / Original pages \dev{२२ -- १६०}}\\[0.3cm]

{\small\textit{Brahmagupta} (c. 598--668 CE)}

\vspace{1em}
\shortline
\end{center}

\end{document}
%% ========================================================================
%% END OF BILINGUAL EDITION
%% ========================================================================

